\documentclass[a4paper,11pt,openany]{ltjsbook}
\usepackage{luatexja}
\usepackage{luatexja-fontspec}
\usepackage{amsmath,amssymb,amsfonts,amsthm}
\usepackage{geometry}
\geometry{top=24mm,bottom=24mm,left=24mm,right=24mm,footskip=10mm}
\usepackage{tcolorbox}
\tcbuselibrary{skins,breakable}
\usepackage{hyperref}
\usepackage{fancyhdr}

\definecolor{primaryblue}{RGB}{20, 60, 110}
\definecolor{darkslate}{RGB}{35, 45, 60}
\definecolor{accentcrimson}{RGB}{160, 30, 45}

\hypersetup{
    colorlinks=true,
    linkcolor=primaryblue,
    citecolor=darkslate,
    urlcolor=accentcrimson,
    pdftitle={純粋数学教程 —— 実数・関数・極限・微積分の厳密解析体系（原著完全翻訳・数学校閲版）},
    pdfauthor={ゴッドフレイ・ハロルド・ハーディ}
}

\usepackage{array}
\usepackage{calc}

% Gutenberg Math & Syllogism macros
\newlength{\TmpLen}
\newenvironment{Syllogism}{$\begin{array}{*3{lc}}}{\end{array}$}

\newcommand{\PadTo}[3][c]{\settowidth{\TmpLen}{$#2$}\makebox[\TmpLen][#1]{$#3$}}
\newcommand{\PadTxt}[3][c]{\settowidth{\TmpLen}{\text{#2}}\makebox[\TmpLen][#1]{#3}}
\newcommand{\Syl}[7]{#1 & (#2)\quad & #3#4 & #5 & #6#7 \\}
\newcommand{\Thus}{\therefore}
\newcommand{\Circ}{\bigcirc}
\newcommand{\Sqr}{\square}
\newcommand{\Tri}{\bigtriangleup}
\newcommand{\Add}[1]{#1}
\newcommand{\etc}{\text{etc}}
\newcommand{\Log}{\operatorname{Log}}
\newcommand{\Squeeze}[2][0.95]{\scalebox{#1}[1]{$#2$}}
\newcommand{\arc}{\text{arc}}
\newcommand{\radius}{\text{radius}}
\newcommand{\PgSep}[1]{}
\newcommand{\PageSep}[1]{}
\newcommand{\Ed}{\text{Ed}}
\newcommand{\Ord}[1]{\textsuperscript{#1}}
\newcommand{\Tag}[1]{\tag*{\ensuremath{#1}}}
\newcommand{\efrac}[2]{\frac{#1}{#2}}
\newcommand{\Title}[1]{\textit{#1}}
\newcommand{\DPad}[1]{\PadTo[l]{9''''}{#1}}
\newcommand{\Z}{\phantom{0}}
\newcommand{\Neg}{\phantom{-}}
\newcommand{\Strut}[1][12pt]{\rule{0pt}{#1}}
\newcommand{\ds}{\displaystyle}
\newcommand{\ArrayCompress}[1][0.6]{\renewcommand{\arraystretch}{#1}}
\newcommand{\Figure}[2][]{}
\newcommand{\Dpt}{\raisebox{1.25ex}{.}}
\newcommand{\CompBox}[1]{\begin{minipage}[t]{0.45\textwidth}\small #1\end{minipage}}
\newcommand{\Comparison}[2]{\par\noindent\vspace{8pt}\CompBox{#1}\hfil\CompBox{#2}\par}
\newcommand{\Eq}[1]{\ensuremath{#1}}
\newcommand{\Pagelabel}[1]{\phantomsection\label{page:#1}}
\newcommand{\PageRef}[2][page]{\hyperref[page:#2]{#1~\pageref*{page:#2}}}
\newcommand{\GLT}[4]{\left.\begin{array}{l}\ds#1 \\ \ds #3\end{array}\right\} >=< \left\{\begin{array}{l}\ds#2 \\ \ds #4\end{array}\right.}










\newcommand{\opp}{\text{opp}}
\newcommand{\adj}{\text{adj}}
\newcommand{\hyp}{\text{hyp}}
\providecommand{\cosec}{\operatorname{cosec}}
\newcommand{\TEntry}[2][]{#2}
\newcommand{\Reason}[2][]{#2}
\newcommand{\BAR}{\rule{0.25in}{1pt}}
\newcommand{\gcm}{\operatorname{gcm}}

\pagestyle{fancy}
\fancyhf{}
\fancyhead[LE,RO]{\thepage}
\fancyhead[RE]{\leftmark}
\fancyhead[LO]{\rightmark}
\setlength{\headheight}{14pt}
\renewcommand{\headrulewidth}{0.4pt}

\title{
    {\huge\bfseries\gtfamily 純粋数学教程 —— 実数・関数・極限・微積分の厳密解析体系（原著完全翻訳・数学校閲版）}\\[14pt]
    {\large 原題：\textit{A Course of Pure Mathematics}}\\[6pt]
    {\normalsize Project Gutenberg eBook \#38769}
}
\author{\textbf{ゴッドフレイ・ハロルド・ハーディ} 著\\[4pt]教育情報分析研究会 訳・数学校閲}
\date{2026年10月 完全決定版}


\definecolor{softblue}{RGB}{238, 244, 252}

% Environments
\newenvironment{Examples}[1]{\par\medskip\noindent\textbf{【例 #1】}\par\begin{itemize}}{\end{itemize}\medskip}
\newenvironment{examples}{\par\medskip\noindent\textbf{【例】}\par\begin{itemize}}{\end{itemize}\medskip}
\newenvironment{Theorem}{\begin{tcolorbox}[colback=softblue,colframe=primaryblue,title=\textbf{定理}]}{\end{tcolorbox}}
\newenvironment{ParTheorem}[1]{\begin{tcolorbox}[colback=softblue,colframe=primaryblue,title=\textbf{定理 (#1)}]}{\end{tcolorbox}}
\newenvironment{Corollary}{\begin{tcolorbox}[colback=white,colframe=primaryblue!60,title=\textbf{系}]}{\end{tcolorbox}}
\newenvironment{Cor}{\begin{tcolorbox}[colback=white,colframe=primaryblue!60,title=\textbf{系}]}{\end{tcolorbox}}
\newenvironment{Definition}{\begin{tcolorbox}[colback=softblue,colframe=darkslate,title=\textbf{定義}]}{\end{tcolorbox}}
\newenvironment{Defn}{\begin{tcolorbox}[colback=softblue,colframe=darkslate,title=\textbf{定義}]}{\end{tcolorbox}}
\newenvironment{Remark}{\begin{tcolorbox}[colback=white,colframe=gray!60,title=\textbf{注釈}]}{\end{tcolorbox}}
\newenvironment{Result}{\begin{tcolorbox}[colback=softblue,colframe=primaryblue,title=\textbf{結果}]}{\end{tcolorbox}}

% Macros
\newcommand{\Item}[1]{\item \textbf{#1}\quad }
\newcommand{\SubItem}[1]{\item[#1] }
\newcommand{\Paragraph}[1]{\par\medskip\noindent\textbf{#1}\par\smallskip}
\newcommand{\SecNo}[2][§]{\textbf{#1#2}}
\newcommand{\SectNo}[2][§]{\textbf{#1#2}}
\newcommand{\Section}[1]{\textbf{§#1}}
\newcommand{\okrickRef}[2]{#1}
\newcommand{\Exs}[1]{\textbf{Exs.~#1}}
\newcommand{\Ex}[1]{\textbf{Ex.~#1}}
\newcommand{\MiscEx}[1]{\textbf{雑題~#1}}
\newcommand{\Page}[1]{}
\newcommand{\PageLabel}[1]{}
\newcommand{\Hang}[1]{#1}
\newcommand{\Loosen}[1]{#1}
\newcommand{\Repeat}[2]{#2}
\newcommand{\Seg}[1]{\overline{#1}}
\newcommand{\Fig}[1]{(図~#1)}
\newcommand{\Cis}{\operatorname{cis}}
\newcommand{\DELTA}{\Delta}
\newcommand{\EPSILON}{\varepsilon}
\newcommand{\DPmod}[2]{#1}
\newcommand{\DPchg}[2]{#2}
\newcommand{\DPtypo}[2]{#2}
\newcommand{\CenterDef}[1]{\begin{center}#1\end{center}}
\newcommand{\CenterLine}[1]{\begin{center}#1\end{center}}
\newcommand{\Topic}[1]{\par\noindent\textbf{#1}\par}
\newcommand{\topic}[1]{\textbf{#1}}
\newcommand{\Emph}[1]{\emph{#1}}
\newcommand{\Ie}{\textit{i.e.}}
\newcommand{\eg}{\textit{e.g.}}
\newcommand{\Real}{\operatorname{Re}}
\newcommand{\Imag}{\operatorname{Im}}
\newcommand{\MathTrip}[1]{#1}
\newcommand{\mathtrip}[1]{#1}
\newcommand{\Construction}[1]{\par\noindent\textbf{【作図 #1】}\par}
\newcommand{\Par}[1]{\par\medskip\noindent\textbf{#1}\par\smallskip}
\newcommand{\btw}{}
\newcommand{\dd}{\mathrm{d}}
\newcommand{\cosech}{\operatorname{cosech}}
\newcommand{\sech}{\operatorname{sech}}
\newcommand{\arcsin}{\operatorname{arcsin}}
\newcommand{\arccos}{\operatorname{arccos}}
\newcommand{\arctan}{\operatorname{arctan}}
\newcommand{\arccot}{\operatorname{arccot}}
\newcommand{\arcsec}{\operatorname{arcsec}}
\newcommand{\arccosec}{\operatorname{arccosec}}
\newcommand{\argsinh}{\operatorname{arsinh}}
\newcommand{\argcosh}{\operatorname{arcosh}}
\newcommand{\argtanh}{\operatorname{artanh}}
\newcommand{\argcoth}{\operatorname{arcoth}}
\newcommand{\bigsqrt}[2][2]{\sqrt[#1]{#2}}
\newcommand{\bigsqrtp}[2][2]{\sqrt[#1]{#2}}
\newcommand{\bigsqrtb}[2][2]{\sqrt[#1]{#2}}
\newcommand{\sqrtp}[1]{\sqrt{#1}}
\newcommand{\sqrtb}[1]{\sqrt{#1}}
\newcommand{\sqrtbr}[1]{\sqrt{#1}}
\newcommand{\tsum}{\sum}
\newcommand{\itemp}{\item}
\newcommand{\pagesepp}[1]{}
\newcommand{\pageseppage}[1]{}
\newcommand{\pagetsep}[1]{}
\newcommand{\MC}[3]{\multicolumn{#1}{#2}{#3}}
\newcommand{\MS}[3]{\multicolumn{#1}{#2}{#3}}

\newcommand{\am}{\operatorname{am}}
\begin{document}
\maketitle
\frontmatter
\tableofcontents
\clearpage
\mainmatter

\chapter{第1章：Chapter I: REAL VARIABLES}

第一章 実変数

1. 有理数。分数$r = p/q$、ここで$p$と$q$は正または負の整数であるとき、そのような分数を有理数と呼ぶ。以下の仮定が成り立つとする：(i) $p$と$q$に共通な因子がないようにする。共通な因子がある場合、それぞれをそれで割ることができる。(ii) $q$は正であるようにする。なぜなら
\[
p/(-q) = (-p)/q,\quad (-p)/(-q) = p/q.
\]
このように定義された有理数に加えて、「有理数$0$」を$p = 0$として導入することができる。

読者は、一般的な有理数の操作に関する算術則について既に馴染みがあると仮定する。以下に示す例は、これ以上の知識が要求されない。

例題 I.
1. $r$と$s$が有理数であるとき、$r + s$, $r - s$, $rs$, および $r/s$も有理数である（最後の場合は$s \neq 0$であることが条件）。

2. $\lambda$, $m$, および $n$が正の有理数であり、$m > n$であるとき、$\lambda(m^{2} - n^{2})$, $2\lambda mn$, および $\lambda(m^{2} + n^{2})$も正の有理数である。よって、すべての辺の長さが有理数である直角三角形を任意に決定する方法を示せ。

3. 終わりのある小数は、分母に$2$または$5$以外の素因数を持たない有理数を表す。逆もまた真であり、このような有理数は一意的に終了小数として表現できる。

[一般的小数理論については第IV章で考慮する。]

4. 正の有理数は以下の単純な順列で並べることができる：
\[
\tfrac{1}{1},\quad
\tfrac{2}{1},\quad
\tfrac{1}{2},\quad
\tfrac{3}{1},\quad
\tfrac{2}{2},\quad
\tfrac{1}{3},\quad
\tfrac{4}{1},\quad
\tfrac{3}{2},\quad
\tfrac{2}{3},\quad
\tfrac{1}{4}, \ \dots.
\]

$p/q$がこの数列の$\left[\frac{1}{2}(p + q - 1)(p + q - 2) + q\right]$番目の項であることを示せ。

この数列では、すべての有理数が無限に繰り返されます。したがって、$1$ は $\frac{1}{1}$, $\frac{2}{2}$, $\frac{3}{3}, \dots$ の形で現れます。当然、より単純な形式で既に出た数を省略することでこの問題を避けることができますが、そうすると $p/q$ の正確な位置を決定する問題が複雑になります。

\par\medskip

\Paragraph{2. 有理数の直線上的表現。} 数学解析の多くの分野では、幾何学的な図示を利用することが便利です。

このような形での幾何学的図示は、分析が幾何に依存していることを意味するものではありません：これらは単なる説明上の例であり、明確性を高めるために使用されるだけです。このようにして、普通の初等幾何学的概念について論理的な分析を行う必要はありません；しかし真実からどれほど遠くにあるかに関係なく、これらの概念が何を意味するかを仮定することができます。

したがって、直線 $\Lambda$、線分 $A_{0}A_{1}$ の長さという概念を仮定します。$A_{0}$ を原点、または点 $0$ と呼び、$A_{1}$ を点 $1$ と呼びます。これらの点は数 $0$ および $1$ を表します。

為了得到表示正有理数 $r = p/q$ 的点，我達选择点 $A_{r}$ 使得
\[
\frac{A_{0}A_{r}}{A_{0}A_{1}} = r,
\]
其中 $A_{0}A_{r}$ 是沿与 $A_{0}A_{1}$ 同方向延伸的线段。当如 \Fig{1} 所示将该线水平画在纸上時，我達假设是从左到右的方向。為了得到表示負有理数 $r = -s$ 的点，自然地可以认為长度是一个具有符号的量，如果沿 $A_{0}A_{1}$ 方向测量，則為正；若反方向测量則為負，因此 $AB = -BA$；我達取表示 $r$ 的点為 $A_{-s}$ 使得
\[
A_{0}A_{-s} = -A_{-s}A_{0} = -A_{0}A_{s}.
\]
\PageSep{3}

これ样，に対して每一个正有理数或負有理数的值 $r$，我達都可以在该线上得到一个点 $A_{r}$，并且満たす
\[
A_{0}A_{r} = r \cdot A_{0}A_{1};
\]
{\Loosen 且如果自然地将 $A_{0}A_{1}$ 取為长度単位，并写成 $A_{0}A_{1} = 1$，則有}
\[
A_{0}A_{r} = r.
\]
我達将点 $A_{r}$ 称作该线上的 \emph{有理点}。

3. 無理数。如果读者在该线上依次标出所有分母為 $1, 2, 3, \dots$ 的有理数対应的点，他将很容易说服自己可以非常紧密地覆盖整个线段。我達可以更精确地表述如下：\emph{如果我達取 $\Lambda$ 上的任意一段 $BC$，我達可以在 $BC$ 上找到任意多个有理点}。

例えば、線分$BC$が区間$A_{1}A_{2}$の内部に含まれる場合を考えます。明らかに、正の整数$k$を選んで
\[
k \cdot BC > 1 \tag{1}
\]
という条件を満たすものがあると仮定すれば、$A_{1}A_{2}$を$k$等分すると、少なくとも一つの分割点（例えば$P$）は$BC$の内部に含まれますが、$B$または$C$とは重複しません。もしこれが成り立たなければ、$BC$は$A_{1}A_{2}$を$k$等分したうちのいずれかの部分に完全に含まれることになり、これは条件\eqref{1}と矛盾します。しかし、$P$は明らかに分母が$k$の有理数に対応しています。したがって、少なくとも一つの有理点$P$が$B$と$C$の間にあることが分かります。すると、さらに別の有理点$Q$を$B$と$P$の間に見つけることができ、また$B$と$Q$の間に別の有理点を見つけることができ、これを無限に繰り返すことができます；つまり、先述した通り、任意に多くの有理点を見つけることが可能です。このような状況を表現するために、「$BC$には無数の有理点が含まれる」と言います。

以下が補足説明です：

このような文句「\emph{無数}」や「\emph{無限}」の意味は、例えば「$BC$には無数の有理点が含まれる」や「$BC$上には無限の有理点がある」といった文脈で考慮されることが詳しく述べられます。\okrickRef{第IV章}{IV}\@。主張「無限の正整数がある」とは、「任意に大きな正整数$n$に対しても、それよりも多くの正整数を見つけることができる」という意味です。これはどのような$n$であっても明らかに真であり、例えば$n = 100,000$や$n = 100,000,000$の場合でも同じです。この主張は「任意に多くの正整数を見つけることができる」という意味と完全に等価です。

読者は、次の主張の真偽を簡単に確認できるでしょう。これは、この節の第二段落で証明されたものと本質的に同等です：任意の有理数$r$と任意の正整数$n$に対して、$r$より左右に別の有理数を見つけられることが示されます。そして、その有理数は$r$から$1/n$未満の差異があります。これはただ異なる表現をするためだけに、$r$からどれほどでも小さい差異を持つ有理数を見つけることができることを述べるのに過ぎません。また、任意の二つの有理数$r$と$s$に対して、それらの間に有理数の連鎖を見つけられ、その連続する項はどれほどでも小さな差異を持ちます。言い換えれば、それらの差異は$n$という任意に事前に指定された正整数で$1/n$未満です。

これらの考慮から、読者は線の性質を十分理解するために、その上にある有理点のみを考えることで十分であると考えるかもしれません。そして確かに、線を単に有理点だけで構成し、他の点（もし存在する場合）をすべて排除した場合、残された図形は直感的な直線の性質の多くを持ち、ざっくりと言えば、直線に非常に似た外観と挙動を持つでしょう。

しかし、少し深く考えると、この見解が我々に深刻な困難をもたらすことが明らかになります。

一時的に直感的な視点で問題を見つめてみましょう。そして、基本幾何学における直線の概念に基づいて、直線が持つ合理的に期待できるいくつかの性質について考えてみます。

直線は点から成り立つものであり、その任意の部分は端点間にあるすべての点で構成される。どのような部分に対しても、その長さという特定のエンティティが関連付けられ、これは任意の標準単位長さに基づいて数値測定可能な量であるべきであり、これらの長さは通常の代数学のルールに従って加算や乗算によって組み合わせられるべきである。

また、任意の2つの与えられた長さの和または積に対応する直線を構成することが可能でなければならない。与えられた直線上の長さ$PQ$が$a$であり、同じ直線上の長さ$QR$が$b$である場合、長さ$PR$は$a + b$でなければならず、これは代数的な加算によるものである。

さらに、一直線上の点$O$から始まる長さ$OP$が$1$、$OQ$が$a$であり、別の直線上の長さ$OR$が$b$の場合、$OP$、$OQ$、$OR$に対する第四比例をEuclidの構成法（エウクリッドの六巻目第12命題）で決定した長さ$OS$は$a \times b$であるべきであり、これは代数的な第四比例である。

これらの定義された和と積は通常の「代数学の法則」を遵守するものでなければならない。すなわち、

\begin{gather*}
a + b = b + a, \quad a + (b + c) = (a + b) + c, \\
ab = ba, \quad a(bc) = (ab)c, \quad a(b + c) = ab + ac.
\end{gather*}

私たちの直線の長さも、等式だけでなく不等式に関する多くの明快な法則を遵守するべきである。例えば、$\Lambda$上から左端に向かって点$A$、$B$、$C$が並んでいる場合、$AB < AC$ という関係が成り立つなどである。

さらに、基本的な直線$\Lambda$上で任意のセグメントや他の直線上の任意のセグメントに対応する点$P$を見つけることが可能でなければならない。これらの性質は私たちの初等幾何学における前提条件に含まれている。

現在、直線を点の連続体として考え、各点が有理数に対応することにより全ての要件を満たすことは非常に簡単でないことがわかります。例えば、長さ \(x\) に対して \(x^2 = 2\) を構成する幾何学的な基本的な構築があります。例えば、等腰直角三角形 \(ABC\) を作ることができます。ここで \(AB = AC = 1\) とします。すると \(BC = x\) のとき \(x^2 = 2\) です。または、Euclidの構成（Euc. VI. 13）を使用して、1 と 2 の間の比例中項を求める方法もあります。図に示すように。

したがって、私たちの要件は、数 \(x\) で測定される長さと、\(\Lambda\) 上の点 \(P\) の存在を要求します。つまり
\[
A_0P = x, \quad x^2 = 2.
\]
ページ分割：6

しかし、容易にわかるように、\emph{平方が 2 となる有理数は存在しない}ことがわかります。実際、さらに言えることは、任意の正分数 \(m/n\)（既約分数）の平方が \(m/n\) であるような有理数も存在しないということです。

仮に
\[
p^2 / q^2 = m / n,
\]
\(p\) と \(q\) の公因数はなく、また \(m\) と \(n\) の公因数もなければ、\(np^2 = mq^2\) となります。\(q^2\) のすべての因子が \(np^2\) を割り切るため、そして \(p\) と \(q\) の公因数がないことから、\(q^2\) のすべての因子が \(n\) を割り切ることになります。したがって \(n = \lambda q^2\)（\(\lambda\) は整数）となります。しかし、これにより \(m = \lambda p^2\) となり、また \(m\) と \(n\) の公因数がないことから、\(\lambda\) は 1 でなければなりません。したがって \(m = p^2, n = q^2\) となります。

特に、\(n = 1\) を取ることにより、整数は平方根が有理数である場合にのみ存在することもわかります。つまり、有理数が自身が整数である場合にのみその平方が整数となることが証明されます。

それから、我々の要求は、既に構築された有理点以外の点Pと数xの存在を含むことが明らかになります。すなわち、$A_{0}P = x$, $x^{2} = 2$；そして（読者には初等代数から覚えているように）私たちは$x = \sqrt{2}$と書きます。

---

次のように、有理数の平方が2であることはないことを別の証明法で示します。

もし可能であれば、既約分数$p/q$を仮定し、$(p/q)^{2} = 2$または$p^{2} = 2q^{2}$を満たす正の分数を考えます。これは簡単に言えば、$(2q - p)^{2} = 2(p - q)^{2}$を導くことになります；そして $(2q - p)/(p - q)$も同じ性質を持つ別の分数となります。明らかに$q < p < 2q$であり、したがって$p - q < q$です。これにより、$p/q$と等しいもう一つの分数があり、分母が小さいことになります。これは$p/q$が既約分数であるという仮定を矛盾させます。

---

以下が例題です：

1. 2の立方根は有理数でないことを示してください。

2. 一般に、最も簡単な形で表された有理数比$p/q$が立方根が有理数であるためには、$p$と$q$がともに完全立方数である必要がありますことを証明してください。

3. ガウスによって提唱され、それ以前の命題を特別な場合として含むより一般的な命題は以下の通りです：「整数係数を持つ代数方程式
\[
x^{n} + p_{1}x^{n-1} + p_{2}x^{n-2} + \dots + p_{n} = 0,
\]
が有理根（ただし非整数）を持つことは不可能である」。

[もし方程式に有理根$a/b$があると仮定すると、$a$と$b$は共通の因子を持たない整数であり、かつ$b$は正です。$x = a/b$を代入し、両辺に$b^{n-1}$を乗じると
\[
-\frac{a^{n}}{b} = p_{1}a^{n-1} + p_{2}a^{n-2}b + \dots + p_{n}b^{n-1}
\]
となり、最も簡単な形で表された分数が整数と等しいことになります。これは矛盾です。したがって$b = 1$であり、根は$a$となります。明らかに$a$は$p_{n}$の約数であることがわかります。]

\par\noindent\textbf{・} {4.} \( p_{n} = 1 \)であり、以下の和がいずれも零でないとき、方程式は有理数の根を持てません：
\[
1 + p_{1} + p_{2} + p_{3} + \dots,\quad
1 - p_{1} + p_{2} - p_{3} + \dots
\]

\par\noindent\textbf{・} {5.} 以下の方程式の有理数の根（存在する場合）を求めよ：
\[
x^{4} - 4x^{3} - 8x^{2} + 13x + 10 = 0.
\]
根は整数であることがわかるため、候補は \( \pm 1, \pm 2, \pm 5, \pm 10 \) であり、これらの値が根であるかを試験によって確認できる。このような方程式の有理数の根も同様にこの方法で決定することができる。

\paragraph{4. 無理数（続）．} 有理数の幾何学的表現は、`数' の概念をさらに新しい種類の数の導入によって拡張すべきであることを示唆します。

同じ結論は幾何学的な言葉を使わなくても得られます。代数学の中心的な問題は方程式の解法であり、例えば以下の方程式：
\[
x^{2} = 1,\quad x^{2} = 2.
\]
最初の方程式は有理数の根 \( \pm 1 \) を持っています。しかし、数の概念が有理数に限定されるなら、第二の方程式は根を持たないとしか言えません；同様に \( x^{3} = 2, \quad x^{4} = 7 \) のような方程式も同じです。これらの事実は、有理数の概念を拡張する必要があることを明らかにします（それが可能である場合）。

方程式 \( x^{2} = 2 \) をより詳しく考えてみましょう。

すでに、有理数$x$が次の方程式を満たすことはないことを確認した。任意の有理数の平方は$2$未満か$2$より大きいいずれかである。従って、有理数を二つのクラスに分けることができる。一つ目のクラスには平方が$2$未満の数が含まれ、もう一つのクラスには平方が$2$より大きい数が含まれる。ここでは正の有理数に限定し、これらの二つのクラスを「クラス$L$」、「下側クラス」、「左側クラス」と呼ぶ一方で、「クラス$R$」、「上側クラス」、「右側クラス」と呼ぶことにする。明らかに、クラス$R$のすべての要素はクラス$L$のすべての要素より大きい。さらに、クラス$L$から平方が$2$未満だが$2$とどれだけでも近い有理数を、クラス$R$から平方が$2$より大きいが$2$とどれだけでも近い有理数を見つけることは容易である。実際、$2$の平方根を通常の算術的手法で計算すると、次の有理数列を得る：
\[
1,\quad 1.4,\quad 1.41,\quad 1.414,\quad 1.4142,\ \dots
\]
これらの平方は
\[
1,\quad 1.96,\quad 1.9881,\quad 1.999396,\quad 1.999961\MS64,\ \dots
\]
すべて$2$未満だが、$2$に近づいていく。そして、計算過程で与えられる十分多くの桁数を使用することで、任意の精度で近似値を得ることができる。さらに、上記の近似値の各最後の桁をそれぞれ$1$増やすと、次の有理数列が得られ：
\[
2,\quad 1.5,\quad 1.42,\quad 1.415,\quad 1.4143,\ \dots
\]
これらの平方は
\[
4,\quad 2.25,\quad 2.0164,\quad 2.00225,\quad 2.000244\MS49,\ \dots
\]
すべて$2$より大きいが、どれだけでも近づいていく。

先行の論理は、おそらく読者に納得させることはできるでしょうが、現代の数学には必要な正確な性質ではありません。形式的な証明を次のように提供できます。まず、$L$ と $R$ のそれぞれから、十分に小さい差異を持つ有理数を一つずつ選びます。第3節で見たように、任意の二つの有理数 $a$ と $b$ に対して、$a$ と $b$ を最初と最後とする有理数の連鎖を作ることができますし、その中のどの二つの連続する数も十分に小さい差異を持つことができます。したがって、$L$ の元 $x$ と $R$ の元 $y$ を選び、$x$ が最初で $y$ が最後となる有理数の連鎖を $x$ と $y$ の間に挿入します。この連鎖は任意に小さな正有理数 $\delta$（たとえば $0.01$ や $0.0001$ や $0.000001$）よりも小さい差異を持つように構築できます。

この連鎖の中で、最後の数が $L$ に属し、最初の数が $R$ に属する必要があります。これらの二つの数は $\delta$ よりも小さい差異を持っています。

次に、$x \in L$ と $y \in R$ を見つけることができることを証明します。すなわち、$2 - x^2$ と $y^2 - 2$ の値が十分に小さくなるような $x$ と $y$ を見つけることができます。$\delta$ の代わりに $\frac{1}{4}\delta$ を用いて先の議論を適用すると、$y - x < \frac{1}{4}\delta$ となる $x$ と $y$ を選べることがわかります。また、明らかに両方とも $2$ より小さいと仮定できます。したがって、

\[
y + x < 4, \quad y^2 - x^2 = (y - x)(y + x) < 4(y - x) < \delta;
\]

となります。そして、$x^2 < 2$ かつ $y^2 > 2$ より、より強力な論証から $2 - x^2$ と $y^2 - 2$ の両方が $\delta$ よりも小さいことがわかります。

また、\( L \)の最大の要素や\( R \)の最小の要素は存在しないこともわかります。なぜなら、\( x \)が\( L \)の任意の要素である場合、\( x^2 < 2 \)です。もし \( x^2 = 2 - \delta \) であれば、\( L \)に属する要素 \( x_1 \) を見つけることができます。このとき、\( x_1^2 \)は2より小さいが、\(\delta\)未満の差で2と異なります。したがって、\( x_1^2 > x^2 \) または \( x_1 > x \) となります。つまり、\( L \)には\( x \)よりも大きい要素があります；そして、\( x \)は\( L \)の任意の要素であるため、\( L \)のどの要素も他のすべての要素より大きくないことがわかります。したがって、\( L \)には最大の要素がなく、同様に\( R \)にも最小の要素はありません。

5. 无理数（续）。我達已经将正有理数分為两类，即 \( L \) 和 \( R \)，満たす以下条件：(i) \( R \) 的任意元素大于 \( L \) 的任意元素；(ii) 我達可以找到 \( L \) 和 \( R \) 中的两个元素，它達之间的差可以任意小；(iii) \( L \) 没有最大值且 \( R \) 没有最小值。我達常识中的直线属性、初等几何和初等代数的要求都要求存在一个数\( x \)，它大于所有 \( L \) 的元素而小于所有 \( R \) 的元素，并且存在相应的点 \( P \) 在直线上，使得 \( P \) 将対应于 \( L \) 中的元素的点与対应于 \( R \) 中的元素的点分开。
\%[図表：図3。]
\fbox{【図版】}

Let us suppose for一時的にある数$x$が存在し、代数学の法則に従って操作できると仮定します。すると例えば$x^{2}$は明確な意味を持ちます。その場合、$x^{2}$は$2$より小さくも大きくもなりません。なぜなら、例えば$x^{2}$が$2$より小さいとするとき、それ以前の議論から、$x^{2}$と$2$の中間にある正の有理数$\xi$を見つけることができます。つまり、$L$の要素の中に$x$より大きいものが存在することになります。これは$x$が$L$の要素を$R$の要素から分離するという仮定に矛盾します。したがって$x^{2}$は$2$より小さくないことが分かりますし、同様に$2$より大きくもありません。従って結論として$x^{2} = 2$であり、$x$は代数学で$\sqrt{2}$と表す数です。もちろんこの数$\sqrt{2}$は有理数ではありません。なぜならどの有理数の平方も$2$にはなりません。これは無理数の最も単純な例であると言えます。

しかし、先ほどの議論は$x^{2} = 2$以外の方程式にもほとんど同じ形で適用できます。例えば$x^{2} = N$の場合、$N$が完全平方でない整数であれば、または
\[
x^{3} = 3,\quad x^{3} = 7,\quad x^{4} = 23,
\]
あるいは後ほど見るように$x^{3} = 3x + 8$の場合でも同様です。これにより無理数$x$と$\Lambda$上の点$P$の存在を信じるようになります。これらの長さは（$\sqrt{2}$のように）基本的な幾何学的手法で構築できない場合でも、このような方程式を満たす$x$が存在することです。

\par\medskip\noindent\textbf{【注釈】}\par
読者は、初等代数の教科書で、方程式 $x^{q} = n$ の根を $\sqrt[q]{n}$ または $n^{1/q}$ で表すことを覚えているでしょう。また、符号
\[
n^{p/q}, \quad n^{-p/q}
\]
は以下のような等式によって意味が付けられています：
\[
n^{p/q} = (n^{1/q})^{p}, \quad n^{p/q} n^{-p/q} = 1.
\]
さらに、これらの定義により、「指数法則」
\[
n^{r} × n^{s} = n^{r+s}, \quad (n^{r})^{s} = n^{rs}
\]
が任意の有理数 $r$ と $s$ に対して拡張されることが思い出されます。
\par\medskip

読者は、今後二つの選択肢のいずれかを選べます。もし好むなら、無理数のような $\sqrt{2}$, $\sqrt{3}, \dots$ の存在と、彼が知っている代数法則に従うことを仮定することができます。\footnote
  {この見解は本書の初版で採用されました。}
このように仮定すれば、次の数節のより抽象的な議論を避けることができ、すぐに \SecNo[§§]{13} 以降へ進むことができます。

一方で、このような素朴な態度に満足しないなら、これらの問題がより詳しく検討される次節を慎重に読むことを推奨します。\footnote
  {これらの節では、Bromwich の \textit{無限級数} の付録 I を自由に引用しています。}

\par\medskip\noindent\textbf{【例 III.】}\par
\Item{1.} $\sqrt{2}$ の近似値として \SecNo[§]{4} で与えられた小数の平方と $2$ の差を求めよ。

\Item{2.} 次の分数の平方と $2$ の差を求めよ：
\[
\tfrac{1}{1}, \quad \tfrac{3}{2}, \quad \tfrac{7}{5}, \quad \tfrac{17}{12}, \quad \tfrac{41}{29}, \quad \tfrac{99}{70}.
\]
\par\medskip

\par\noindent\textbf{・} {3.} \( m/n \)が\(\sqrt{2}\)の良い近似である場合、\( (m + 2n)/(m + n) \)はより良い近似であり、両者の誤差は反対方向にある。この結果を用いて、最後の例で示された近似列を続けよ。

\par\noindent\textbf{・} {4.} \( x \)と\( y \)がそれぞれ\(\sqrt{2}\)の欠陥的な近似と余りの近似であり、かつ \( 2 - x^2 < \delta \), \( y^2 - 2 < \delta \)の場合、\( y - x < \delta \)である。

\par\noindent\textbf{・} {5.} 方程式 \( x^2 = 4 \)は\( x = 2 \)で満たされる。前節の議論がこの方程式にどの程度適用できるかを検討せよ（常に2を4と書き換える）。[もしクラスL, Rを以前のように定義すると、これらにはすべての有理数が含まれない。有理数2は例外であり、2\(^2\)は4より大きいとも小さいとも言えないためである。]

\paragraph{6. 无理数（续）．} 第4節では、正の有理数\( x \)を二つのクラスに分ける特別な方法について議論した。その中で一つのクラスの要素に対して\( x^2 < 2 \)であり、もう一つのクラスの要素に対して\( x^2 > 2 \)であるような分類法があった。このような分類法は問題の数を「分割」と呼ぶ。明らかに、同じように二つのクラスの要素が不等式 \( x^3 < 2 \)と\( x^3 > 2 \)、または\( x^4 < 7 \)と\( x^4 > 7 \)で特徴付けられるような「分割」を作成することも可能である。現在では、このような「分割」の原理を正の有理数全体に対して一般化した形で述べてみよう。

ある正の有理数が持つべき2つの互いに排他的な性質を$P$と$Q$とする。さらに、どのような数においても、$P$を満たす数は$Q$を満たす数よりも小さいという仮定をする。つまり、$P$は「$x^2 < 2$」であり、$Q$は「$x^2 > 2$」である可能性がある。このような性質$P$を持つ数の集合を下クラスまたは左クラス$L$と呼び、$Q$を持つ数の集合を上クラスまたは右クラス$R$と呼ぶ。一般に両方のクラスが存在するが、特定の場合には一方のクラスが存在せず、すべての数がもう片方のクラスに属することもある。例えば、$P$（または$Q$）が有理数であるか、正であるという性質であれば、明らかにこのような状況が生じる。

しかし、現在は両方のクラスが存在する場合に限定し、そのような場合には$\SecNo[§]{4}$で述べたように、$L$から$R$への差が任意に小さくなるような要素をそれぞれ見つけることができる。

$\SecNo[§]{4}$で考慮した特定の場合においては、$L$には最大の要素がなく、$R$にも最小の要素がなかった。クラス内の最大または最小の要素の存在性については極めて重要である。まず最初に述べるが、$L$に最大の要素がある一方で$R$に最小の要素があることはあり得ない。なぜなら、$l$が$L$の最大の要素であり、$r$が$R$の最小の要素であり、かつ$l < r$である場合、$\frac{1}{2}(l + r)$は$l$と$r$の間に存在する正の有理数で、これには属さないという事実から矛盾が生じる。したがって、以下の3つの可能性しかない。すなわち、(i)$L$に最大の要素$l$がある、または(ii)$R$に最小の要素$r$がある、または(iii)$L$に最大の要素がないかつ$R$にも最小の要素がない。

\par\medskip\noindent\textbf{【注釈】}\par
第4節（§4）の部分は最後の可能性の例を与えています。最初の可能性の例は、$P$ を「$x^{2} \leq 1$」、$Q$ を「$x^{2} > 1$」と定めるところから得られます；ここでは $l = 1$ です。もし $P$ を「$x^{2} < 1$」、$Q$ を「$x^{2} \geq 1$」とすると、二番目の可能性の例が得られ、このとき $r = 1$ です。$P$ を「$x^{2} < 1$」、$Q$ を「$x^{2} > 1$」と定める場合では、特別な数 $1$ は分類から逃れることに注意すべきです（参照：例 iii.5）。
\par\medskip

\Paragraph{7. 無理数（続）。}最初の二つの場合において、私たちはセクションが正の有理数 $a$ に対応すると述べます。この $a$ は一つの場合では $l$、もう一つの場合では $r$ です。逆に、このような数 $a$ に対して、対応するセクション $\alpha$ が存在することは明らかです。
足注
：有理数を表す英語の文字に対応するセクションは、その英語の文字から始まるギリシャ文字で表すことに便益があるでしょう。

なぜなら、$P$ と $Q$ をそれぞれ「$x \leq a$」と「$x > a$」または「$x < a$」と「$x \geq a$」とすることができるからです。最初の場合では $a$ は $L$ の最大の要素であり、二番目の場合では $R$ の最小の要素となります。実際には、任意の正の有理数に対してただ二つのセクションが存在します。明確性を避けるために一つを選択することにします；私たちはその中で自身が上クラスに属するセクションを選択することとします。つまり、下クラス $L$ が最大の要素を持たないセクションのみを考えることとします。
ページ番号：13


有理数とその定義により区切られた部分との間の対応があるため、数学的な目的では有理数を部分に置き換えて、公式中のシンボルが部分を表していることとして扱うことも完全に正当です。例えば、$\alpha > \alpha'$ は $\alpha$ と $\alpha'$ が $a$ と $a'$ に対応する部分である場合、$a > a'$ と同じ意味になります。

しかし、有理数をその自身の部分で置き換えたときには、私たちに大きな数の体系の一般化を迫られます。実際には（例えば第4節の部分のように）ある部分はどの有理数にも対応していません。部分全体は正の有理数全体よりも大きい集合であり、すべてのこれらの数に対応する部分とそれ以上のものが含まれています。この事実に基づいて、私たちは数の概念を一般化します。そのため、次の定義を以下に述べますが、次節で修正される可能性があるため、一時的な暫定的なものとして扱うべきです。

\par\medskip\noindent\textbf{【定義】}\par
両方のクラスが存在し、下側のクラスには最大の要素がない有理数の部分を \Emph{正実数} と呼びます。
\par\medskip

\par\medskip\noindent\textbf{【定義】}\par
有理数に対応しない正実数を \Emph{正無理数} と呼びます。
\par\medskip

8. 実数。これまで私たちは正の有理数の一部に限定して議論してきました。これらを一時的に「正の実数」と呼びました。最終的な定義を立てる前に、視点を少し変える必要があります。実数とは、有理数全体（正の有理数だけでなく零も含む）を二つのクラスで分ける場合を考えることです。その後、§6, §7 で述べた有理正数の一部に関するすべての議論を再現し、時として「正」という言葉を使わないようにします。

\par\medskip\noindent\textbf{【定義】}\par
有理数全体の部分集合において、両方のクラスが存在し、下側のクラスに最大値がない場合、その部分集合は「実数」または単に「数」と呼ばれます。
実数が有理数に対応しない場合は、「無理数」と呼ばれます。
\par\medskip

実数が有理数に対応する場合、その実数についても「有理的」という用語を使用します。

定義により、「有理数」という用語は曖昧なものとなります；それは第1章の有理数を指すものか、あるいはその対応する実数を指すものかもしれません。$\frac{1}{2} > \frac{1}{3}$ という主張は、基本的な算術の命題であるか、あるいは有理数の分割に関する命題であるかのいずれかを意味する可能性があります。このような曖昧性は数学において一般的であり、命題間の関係がどの解釈を付与しても変わらないため完全に無害です。$\frac{1}{2} > \frac{1}{3}$ と $\frac{1}{3} > \frac{1}{4}$ より $\frac{1}{2} > \frac{1}{4}$ を推論することは、$\frac{1}{2}$, $\frac{1}{3}$, および $\frac{1}{4}$ が算術的な分数であるか実数であるかの疑問に関係なく同じです。もちろん、（たとえば）`$\frac{1}{2}$' の出現文脈によってその解釈が固定される場合もあります。第9章で述べる $\frac{1}{2} < \sqrt{\frac{1}{3}}$ という主張では、「$\frac{1}{2}$」は実数 $\frac{1}{2}$ を意味するものと理解しなければなりません。

读者はまた、我々が採用した「実数」という概念の定義の正確な形式には特別な論理的重要性はないと注意すべきであることに留意すべきである。我々は「実数」を区間、すなわちクラスのペアとして定義した。同様に、「実数」を下側のクラスや上側のクラスとして定義することができる；実際には、実数のクラスと同じ性質を持つ無限多種類のエンティティのクラスを定義することも容易である。数学において重要なのは、その記号が何かしらの解釈が可能であることです；一般的に、多くの解釈が可能です。そして、数学的な観点から言えば、採用する解釈は関係ありません。

Bertrand Russell氏は「我々が話していることについて何も知らないし、それが真偽に関わらず気にしない」という言葉で、「数学は我々が何を話しているのか知らなくても、それに真偽に関わらず何かを述べることに興味がない科学である」と述べています。この言葉はパラドックスの形で表現されていますが、実際には重要な真理を多く含んでいます。

Russell氏の格言の詳細な意味を分析するのに時間がかかりますが、少なくとも一つの意味は、数学の記号は様々な解釈が可能であり、一般的に好む解釈を選択できるということです。

現在三つの場合を区別する必要があります。負の有理数がすべて下クラスに属し、ゼロと正の有理数がすべて上クラスに属することがあるかもしれません。このセクションは「実数零」と呼びます。あるいは、下クラスに一部の正の数が含まれることがあります。そのようなセクションは「正の実数」と呼びます。さらに、一部の負の数が上クラスに属することがあります。そのようなセクションは「負の実数」と呼びます。

また、すべての数が下クラスまたは上クラスに属するセクションもあります。読者はこれらのセクションも数として定義することを考えてみるかもしれません。これらを「正無限大と負無限大の実数」と呼ぶべきでしょう。

このような手続きには論理的な反対はなく、実際的には不便なことがわかります。加法や乗法の最も自然な定義は適切に機能しないことが明らかになります。さらに、解析の要素を学ぶ初心者にとって最大の困難は「無限」という単語が含まれる文句に対する正確な意味を付与することです。経験から見ると、その数が増えることは混乱を招く可能性が高いようです。

以下に現在定義した正の実数 \(a\) と \SecNo[§]{7} の定義との違いについて述べます。これは下クラスにゼロとすべての負の有理数が加わることです。負の実数の例としては、\SecNo[§]{6} の性質 \(P\) を \(x + 1 < 0\)、\(Q\) を \(x + 1 \geq 0\) として選びます。このセクションは明らかに負の有理数 \(-1\) に対応します。また、\(P\) を \(x^3 < -2\)、\(Q\) を \(x^3 > -2\) として選ぶと、無理数でない負の実数が得られます。

9. 実数間の量的关系。実数概念を拡張した以上、等号、不等号、加法、乗法などの概念も対応する形で拡張せざるを得ないことが明白である。これらの考え方が新しい実数に適用できることを示し、その拡張によって通常の代数学の法則が依然として有効であり、実数に対して一般に代数的操作を行う方法は有理数（§1）に対して同じようにできるかを証明する必要がある。これらすべてを系統的に行うと相当なスペースを占めるため、より系统的な議論の進行方法を概略的に示すことに留まる。

実数をギリシャ文字で表し、その下限クラスと上限クラスの有理数は対応する英文字で表す。つまり、$\alpha$, $\beta$, $\gamma, \dots$ で実数を表し、それぞれの下限クラス$(a)$, $(b)$, $(c), \dots$ と上限クラス$(A)$, $(B)$, $(C), \dots$ の有理数は$a$, $b$, $c, \dots$ および$A$, $B$, $C, \dots$ で表す。

実数$\alpha$と$\beta$について、以下の3つの可能性がある：

(i) すべての$a$が$b$であり、すべての$A$が$B$である場合、このとき$(a)$は$(b)$に等しく、$(A)$も$(B)$に等しい。

(ii) すべての$a$が$b$であり、しかしすべての$A$が$B$ではない場合、このとき$(a)$は$(b)$の真部分集合であり、$(B)$も$(A)$の真部分集合である。ここで$\Ie\ (A)$は$(b)$に含まれるが等しくない。

(iii) すべての$A$が$B$であり、しかしすべての$a$が$b$ではない場合。

これらの3つのケースを図4のようにグラフィカルに示すことができる。

(i) の場合、$\alpha = \beta$と表記する。(ii) の場合、$\alpha < \beta$と表記する。そして(iii) の場合、$\alpha > \beta$と表記する。明らかに、$\alpha$と$\beta$がともに有理数であるとき、これらの定義は最初から仮定した有理数間の等号や不等号の概念と一致し、任意の正数は負数より大きいことがわかる。

この段階で正の数$\alpha$の負の数$-\alpha$を定義します。$(a)$、$(A)$が$\alpha$を構成するクラスである場合、すべての数$-A$を下クラスとし、すべての数$-a$を上クラスとして別の有理数の分割を定義できます。このように定義された実数は明らかに負であり、それを$-\alpha$と表記します。同様に$\alpha$が負または零である場合にも$-\alpha$を定義できます；$\alpha$が負であれば$-\alpha$は正です。また、$-(-\alpha) = \alpha$も明らかです。二つの数$\alpha$と$-\alpha$のうち一つは常に正（$\alpha = 0$の場合は除く）であり、その中で正であるものを$|\alpha|$と表記し、$\alpha$の絶対値またはモジュールと呼びます。

例題 IV.
1. $0 = -0$ を証明せよ。
2. $\alpha = \beta$, $\beta < \alpha$, または $\beta > \alpha$ は $\alpha = \beta$, $\alpha > \beta$, または $\alpha < \beta$ の場合に応じて成り立つことを証明せよ。
3. $\alpha = \beta$ かつ $\beta = \gamma$ のとき、$\alpha = \gamma$ であることを証明せよ。
4. $\alpha \leq \beta$, $\beta < \gamma$ または $\alpha < \beta$, $\beta \leq \gamma$ のとき、$\alpha < \gamma$ であることを証明せよ。
5. $-\beta = -\alpha$, $-\beta < -\alpha$, または $-\beta > -\alpha$ は $\alpha = \beta$, $\alpha < \beta$, または $\alpha > \beta$ の場合に応じて成り立つことを証明せよ。
6. $\alpha > 0$ であるとき$\alpha$が正であり、$\alpha < 0$ であるとき$\alpha$が負であることを証明せよ。
7. $\alpha \leq |\alpha|$ であることを証明せよ。
8. $1 < \sqrt{2} < \sqrt{3} < 2$ であることを証明せよ。
9. 二つの異なる実数$\alpha$と$\beta$について、それらの間に無限に多くの有理数が存在することを証明せよ。

[これらの結果は私たちの定義の直ちなる推論です。]

10. 実数に対する代数的演算. 現在、実数一般に対して基本的な代数的演算（加法など）の意味を定義します。

\par\medskip\noindent\textbf{(i)  加法。}\par\smallskip二つの数 $\alpha$ と $\beta$ の和を定義するために、次の2つのクラスを考える：(i)~すべての和 $c = a + b$ を含むクラス $(c)$, (ii)~すべての和 $C = A + B$ を含むクラス $(C)$。明らかに、全ての場合で $c < C$ である。

また、$(c)$ や $(C)$ のいずれにも属さない有理数は最大でも1つしかない。なぜなら、もし2つ存在すると仮定し、それらを $r$ と $s$ とする。そして $s$ を大きいものとする。このとき、$r$ と $s$ は全ての $c$ よりも大きく、全ての $C$ よりも小さい。したがって、$C - c < s - r$ であることは不可能である。しかし
\[
C - c = (A - a) + (B - b);
\]
であり、$a$, $b$, $A$, $B$ を適切に選べば、両辺ともどれほど小さくてもよいので、これは仮定に矛盾する。

全ての有理数が $(c)$ または $(C)$ に属している場合、クラス $(c)$,~$(C)$ は有理数全体から成るセクションを形成し、すなわちある数 $\gamma$ に対応する。そうでないものがあれば、それを $(C)$ に加える。これにより、セクションまたは実数 $\gamma$ が得られ、これは明らかに有理数であり、$(C)$ の最小元に対応する。どのような場合でも、$\gamma$ を $\alpha$ と $\beta$ の和と呼び、
\[
\gamma = \alpha + \beta
\]
と表す。

\par\medskip\noindent\textbf{【注釈】}\par
もし両方の $\alpha$ と $\beta$ が有理数であるなら、それぞれは上クラス $(A)$ および $(B)$ の最小元である。この場合、$\alpha + \beta$ は $(C)$ の最小元であり、従って我々の加法の定義は以前の考え方に一致する。
\par\medskip

\par\medskip\noindent\textbf{(ii)  減法。}\par\smallskip $\alpha - \beta$ を以下の方程式で定義する：
\[
\alpha - \beta = \alpha + (-\beta).
\]
従って、減法の概念には新たな困難は生じない。

\par\medskip\noindent\textbf{【例 V.】}\par
\Item{1.} $\alpha + (-\alpha) = 0$ を証明せよ。

\Item{2.} $\alpha + 0 = 0 + \alpha = \alpha$ を証明せよ。
\par\medskip

3. α + β = β + α を証明せよ。[これは、クラス (a + b) と (b + a)，または (A + B) と (B + A) が同じであるという事実から即座に導かれる。たとえば、有理数 a および b のとき、a + b = b + a である.]

4. α + (β + γ) = (α + β) + γ を証明せよ。

5. α - α = 0 を証明せよ。

6. α - β = -(β - α) を証明せよ。

7. 減算の定義と上記の例4, 1, 2 より、
\[
(\alpha - \beta) + \beta
  = \{\alpha + (-\beta)\} + \beta
  = \alpha + \{(-\beta) + \beta\}
  = \alpha + 0 = \alpha.
\]
したがって、差 α - β = γ を方程式 γ + β = α で定義することができる。

8. α - (β - γ) = α - β + γ を証明せよ。

9. 加法の前の定義に依存しない減算の定義を挙げよ。[γ = α - β の定義として、c = a - B, C = A - b となるクラス (c), (C) を形成する。この定義が本文で採用した定義と同等であることを容易に示すことができる.]

10. 
\[
\big||\alpha| - |\beta|\big| \leq |\alpha ± \beta| \leq |\alpha| + |\beta|.
\]

11. 実数に対する代数的演算（続）.
(iii) 乗法．乗法に取り組むとき、最初は正の数（0も含む）に限定して考えることが最も便利です。これらは第4節から第7節で考慮した正の有理数の分割集合に一時的に戻ります。その後、加算の場合と同様の道筋をたどることができます。$(c)$ を $(ab)$ 、$(C)$ を $(AB)$ として取り扱います。議論は同じですが、すべての有理数が最多1つを除いて $(c)$ や $(C)$ のいずれかに属属するという証明において異なる点があります。これは加算の場合と同様に、$C - c$ を任意に小さくできる $a$ 、$A$ 、$b$ 、$B$ を選べることを示すことで成立します。ここで使用される恒等式は以下の通りです：
\[
C - c = AB - ab = (A - a)B + a(B - b).
\]

最後に、負の数も定義の範囲内に含めます．$\alpha$ と $\beta$ が正であるとき、以下のようにします：
\[
(-\alpha)\beta = -\alpha\beta,\quad
\alpha(-\beta) = -\alpha\beta,\quad
(-\alpha)(-\beta) = \alpha\beta.
\]

(iv) 分割．分割を定義するために、まず数 $\alpha$（零以外）の逆数 $1/\alpha$ を定義します。最初は正の数と正の有理数の分割集合に限定して、正の数 $\alpha$ の逆数を下クラス $(1/A)$ と上クラス $(1/a)$ を用いて定義します。負の数 $-\alpha$ の逆数は方程式 $1/(-\alpha) = -(1/\alpha)$ によって定義します．最後に、$\alpha/\beta$ を以下の方程式で定義します：
\[
\alpha/\beta = \alpha × (1/\beta).
\]

すべての実数、有理数であれ無理数であれ、初等代数における全ての概念と方法を適用することができます。もちろん、この作業を詳細に行うつもりはありません。より有益で興味深いのは、特に重要ないくつかの無理数のクラスに注意を向けることです。

例 VI.
以下の公式によって表される定理を証明してください：

\begin{minipage}{0.5\textwidth-\parindent}
\Item{1.} $\alpha \times 0 = 0 \times \alpha = 0$.

\Item{2.} $\alpha \times 1 = 1 \times \alpha = \alpha$.

\Item{3.} $\alpha \times (1/\alpha) = 1$.

\Item{4.} $\alpha\beta = \beta\alpha$.
\end{minipage}
\begin{minipage}{0.5\textwidth}
\Item{5.} $\alpha(\beta\gamma) = (\alpha\beta)\gamma$.

\Item{6.} $\alpha(\beta + \gamma) = \alpha\beta + \alpha\gamma$.

\Item{7.} $(\alpha + \beta)\gamma = \alpha\gamma + \beta\gamma$.

\Item{8.} $|\alpha\beta| = |\alpha|\, |\beta|$.
\end{minipage}

12. 数 $\sqrt{2}$。次に、§4--§5 で議論した特定の無理数に戻りましょう。そこで、不等式 $x^{2} < 2$, $x^{2} > 2$ を用いてセクションを構築しました。これは正の有理数のみのセクションでした；しかし、§8 で説明したように、すべての有理数のセクションに置き換えました。このセクションやその定義された数は $\sqrt{2}$ の記号で表されます。

$\sqrt{2}$ 自身を乗算するクラスは次の通りです：(i)~$(aa')$ で、$a$ と $a'$ が平方が $2$ より小さい正の有理数であるとき、(ii)~$(AA')$ で、$A$ と $A'$ が平方が $2$ より大きい正の有理数であるとき。これらのクラスはすべての正の有理数を除き、残るものは $2$ のみです。したがって
\[
(\sqrt{2})^{2} = \sqrt{2}\sqrt{2} = 2.
\]

Again
\[
(-\sqrt{2})^{2}
  = (-\sqrt{2})(-\sqrt{2})
  = \sqrt{2}\sqrt{2}
  = (\sqrt{2})^{2} = 2.
\]
Thus, 方程式 $x^{2} = 2$ の二つの根は $\sqrt{2}$ と $-\sqrt{2}$ である。同様に，方程式 $x^{2} = 3$, $x^{3} = 7, \dots$ を議論し，対応する無理数 $\sqrt{3}$, $-\sqrt{3}$, $\sqrt{7}, \dots$ も検討できる。

13. 二次根式。形が $±\sqrt{a}$ で，$a$ は正の有理数であり，また別の有理数の平方ではないとき，そのような数を純粋な二次根式と呼ぶ。形が $a ± \sqrt{b}$ で，$a$ は有理数であり，$\sqrt{b}$ が純粋な二次根式であるとき，そのような数はしばしば混合二次根式と呼ばれる。

The two numbers $a ± \sqrt{b}$ は，次の二次方程式
\[
x^{2} - 2ax + a^{2} - b = 0
\]
の二つの根である。逆に，$p$ と $q$ が有理数であり，かつ $p^{2} - q > 0$ のとき，方程式 $x^{2} + 2px + q = 0$ の二つの根は二次根式 $-p ± \sqrt{p^{2} - q}$ である。

これらの幾何学的考慮により暗示された唯一の種類の無理数は，純粋な二次根式と混合二次根式であり，さらには重複した平方根の取り出しを含むより複雑な無理数，例えば
\[
\sqrt{2} + \sqrt{2 + \sqrt{2}} + \sqrt{2 + \sqrt{2 + \sqrt{2}}}
\]
のような形で表されるものである。

この形式の任意の数に対応する線分の長さを幾何学的に構築することは容易であり，読者は自分で確認できる。ユークリッド方法（すなわち，コンパスと直尺を使用した幾何学的構成）によってこれらの種類の無理数のみが構築可能であるという性質は，証明のために今後遅らせるべき点である。この二次根式の性質は特に興味深い。

例題 VII.
1. 以下の幾何学的な構築を示せ。
\[
\sqrt{2}, \quad
\sqrt{2 + \sqrt{2}}, \quad
\sqrt{2 + \sqrt{2 + \sqrt{2}}}.
\]

二次方程式 \(ax^{2} + 2bx + c = 0\) は、二つの実数解を持つ（\(\therefore x\) の値が二つ存在する。もし \(b^{2} - ac < 0\) の場合、そのような \(x\) の値は存在しない。初等代数学の本では、この方程式は「複素」な根をもつと述べることがあるが、その意味は第III章で説明される。

\(b^{2} = ac\) のとき、方程式は一つの根を持つ。三つの根が全て有理数である場合、分母の最小公倍数で方程式を乗じることにより、すべての項を整数に置き換えることは問題ない。

方程式の解は \(\left\{-b ± \sqrt{b^{2} - ac}\right\}/a\) であることを忘れないでください。これらの長さを幾何学的に構成することは容易です。まず \(\sqrt{b^{2} - ac}\) を構成します。

もっと洗練された（しかしより直接的ではない）構成法は以下の通りです。
ページ番号の区切り：21

以下が構成法です：

単位円を描き、直径 \(PQ\) とその端点における接線を引きます。\(PP' = -2a/b\) および \(QQ' = -c/2b\) を考慮し、符号に注意してください。（図は \(b\) と \(c\) の符号が同じで \(a\) の符号が異なる場合の例です。他の場合も自分で図を描いてみてください）

\(P'Q'\) を結び、円との交点として \(M\) および \(N\) を取ります。さらに \(PM\) および \(PN\) を引き、それぞれ \(QQ'\) の交点として \(X\) および \(Y\) を取ります。すると \(QX\) および \(QY\) は方程式の解であり、適切な符号をもっています。（この構成法は Kline の『幾何学的要素に関するいくつかの問題についての講義』（フランス語版由来、J. Griess 訳、パリ、1896 年）から引用しています。）

証明は簡単で、読者に任せる。
別の構成法もまた簡便である。以下がその構成法である。

\Construction[]単位長さの線分$AB$を取り、$BC = -2b/a$を$AB$に対して垂直に引く。さらに$CD = c/a$を$BC$に対して垂直かつ$BA$と同じ方向に引く。直径として$AD$を用いて円を描き、この円が$BC$と交わる点を$X$と$Y$とする。すると$BX$と$BY$は根である。

\Item{3.} $ac$が正のとき、$PP'$と$QQ'$は同じ方向に引かれる。$b^{2} < ac$の場合は$P'Q'$は円と交わらないことを確認し、$b^{2} = ac$の場合は接線であることも確認せよ。また、$b^{2} = ac$の場合、第2の構成法による円が$BC$と接することも確認せよ。

\Item{4.} 以下の式を証明せよ。
\[
\sqrtp{pq} = \sqrt{p} × \sqrt{q},\quad
\sqrtp{p^{2}q} = p\sqrt{q}.
\]
\par\medskip

\Paragraph{14. 二次根号に関する定理。}純粋な二次根号が同じ根号の有理倍数で表せるなら、その二次根号は類似であると呼び、そうでない場合は異類であるという。すなわち
\[
\sqrt{8} = 2\sqrt{2},\quad
\sqrt{\frac{25}{2}} = \frac{5}{2}\sqrt{2},
\]
よって$\sqrt{8}$と$\sqrt{\frac{25}{2}}$は類似の根号である。一方、整数$M$と$N$が互いに素でかつどちらも完全平方でないとき、$\sqrt{M}$と$\sqrt{N}$は異類の根号である。もし可能ならば、
\[
\sqrt{M} = \frac{p}{q}\bigsqrt{\frac{t}{u}},\quad
\sqrt{N} = \frac{r}{s}\bigsqrt{\frac{t}{u}},
\]
と仮定すると、すべての文字は整数を表す。

次に、$\DPtypo{\sqrt{MN}}{\sqrtp{MN}}$は明らかに有理数であり、従って(\Ex{ii}.~3)同様に整数である。したがって $MN = P^{2}$ で、$P$は整数である。$a$, $b$, $c, \dots$ を $P$ の素因数とし、
\[
MN = a^{2\alpha} b^{2\beta} c^{2\gamma}\ \dots,
\]
ここで $\alpha$, $\beta$, $\gamma, \dots$ は正の整数である。すると $MN$ は $a^{2\alpha}$ で割り切れるので、または (1) $M$ が $a^{2\alpha}$ で割り切れる、または (2) $N$ が $a^{2\alpha}$ で割り切れる、または (3) $M$ と $N$ がともに $a$ で割り切れる。最後のケースは排除されるべきである、なぜなら $M$ と $N$ には共通の因子がないからである。この議論を各素因数 $a^{2\alpha}$, $b^{2\beta}$, $c^{2\gamma}, \dots$ に対して適用すると、$M$ はこれらの一部が $P_{1}^{2}$ の積で表される因子で割り切れる一方で、$N$ は残りの部分が $P_{2}^{2}$ の積で表される因子で割り切れる。したがって
\[
M = P_{1}^{2},\quad N = P_{2}^{2},
\]
ここで $P_{1}^{2}$ は一部の素因数 $a^{2\alpha}$, $b^{2\beta}$, $c^{2\gamma}, \dots$ の積を表し、$P_{2}^{2}$ は残りの部分の積を表す。従って $M$ と $N$ は両方とも完全平方であり、これは我々の仮定に反する。

\par\medskip\noindent\textbf{【定理】}\par
もし $A$, $B$, $C$, $D$ が有理数で
\[
A + \sqrt{B} = C + \sqrt{D},
\]
であれば、または (i) $A = C$, $B = D$ または (ii) $B$ と $D$ はどちらも有理数の平方である。
\par\medskip

$B - D$ は有理数であり、それゆえ
\[
\sqrt{B} - \sqrt{D} = C - A
\]
も有理数である。もし $B \neq D$ の場合（この場合は明らかに $A = C$）、$\sqrt{B} + \sqrt{D}$ は
\[
(B - D) / (\sqrt{B} - \sqrt{D})
\]
と等しく、これも有理数である。したがって $\sqrt{B}$ と $\sqrt{D}$ は有理数である。

\par\medskip\noindent\textbf{【系】}\par
もし $A + \sqrt{B} = C + \sqrt{D}$ であれば、$A - \sqrt{B} = C - \sqrt{D}$ が成り立つ（$\sqrt{B}$ と $\sqrt{D}$ がどちらも有理数でない限り）。
\par\medskip

\par\medskip\noindent\textbf{【例 VIII.】}\par
\Item{1.} $\sqrt{2}$ と $\sqrt{3}$ は類似の根号であることを最初から証明せよ。

\Item{2.} $a$ が有理数であるとき、$\sqrt{a}$ と $\sqrtp{1/a}$ は類似の根号である（両方が有理数でない限り）。
\par\medskip

3. $a$ と $b$ が有理数であるとき、$\sqrt{a} + \sqrt{b}$ は $\sqrt{a}$ と $\sqrt{b}$ が有理数でない限り有理数にならない。同様に、$a = b$ でない限り $\sqrt{a} - \sqrt{b}$ も有理数にならない。

4. 
\[
\sqrt{A} + \sqrt{B} = \sqrt{C} + \sqrt{D}
\]
のとき、(\ia) $A = C$ かつ $B = D$、または (\ib) $A = D$ かつ $B = C$、または (\ic) $\sqrt{A}$, $\sqrt{B}$, $\sqrt{C}$, $\sqrt{D}$ はすべて有理数かすべて類比的無理数である。[与えられた方程式を二乗し、上記の定理を適用します。]

5. $(a + \sqrt{b})^{3}$ および $(a - \sqrt{b})^{3}$ が有理数でない限り $\sqrt{b}$ は有理数である。

6. $x = p + \sqrt{q}$、$p$ と $q$ が有理数であるとき、任意の整数 $m$ に対して $x^{m}$ を $P + Q\sqrt{q}$ の形で表すことができる。例えば，
\[
(p + \sqrt{q})^{2} = p^{2} + q + 2p\sqrt{q},\quad
(p + \sqrt{q})^{3} = p^{3} + 3pq + (3p^{2} + q)\sqrt{q}.
\]
したがって、任意の有理数係数を持つ多項式 $a_{0}x^{n} + a_{1}x^{n-1} + \dots + a_{n}$（$a_{0}, \dots, a_{n}$ は有理数）も $P + Q\sqrt{q}$ の形で表すことができる。

7. $b$ が完全平方でないとき、$a + \sqrt{b}$ が有理係数を持つ代数方程式の根であるなら、$a - \sqrt{b}$ も同じ方程式の別の根である。

8. $1/(p + \sqrt{q})$ を Ex.6 の形式で表す。[分子と分母に $p - \sqrt{q}$ を掛ける。]

9. Ex.6 および Ex.8 から、任意の形 $G(x)/H(x)$（$G(x)$ と $H(x)$ は有理係数を持つ多項式）も $P + Q\sqrt{q}$ の形で表すことができる。

10. $p$, $q$, および $p^{2} - q$ が正のとき、$\sqrt{p + \sqrt{q}}$ を $\sqrt{x} + \sqrt{y}$ の形で表すことができる。ここで
\[
x = \frac{1}{2}\{p + \sqrt{p^{2} - q}\},\quad
y = \frac{1}{2}\{p - \sqrt{p^{2} - q}\}.
\]

\par\noindent\textbf{・} {11.} $p + \sqrt{q}$ が有理数 $p$ および $q$ のとき、$\sqrtp{p + \sqrt{q}}$ を $\sqrt{x} + \sqrt{y}$ (有理数 $x$ および $y$) の形で表すことができる条件を決定せよ。

\par\noindent\textbf{・} {12.} $a^2 - b$ が正のとき、
\[
\sqrtp{a + \sqrt{b}} + \sqrtp{a - \sqrt{b}}
\]
が有理数であるための必要十分条件は、$a^2 - b$ および $\frac{1}{2}\{a + \sqrtp{a^2 - b}\}$ がいずれも有理数の平方であることである。

\Paragraph{15. 統一連続体。} 全ての実数、有理数および無理数を含む集合は、\Emph{算術的統一連続体}と呼ばれる。

直線 $\Lambda$（§2）は、算術的統一連続体に対応する点から成る点の集合であり、それ以外の点は含まれないとするのが便利である。\footnote
  {この仮定は、(i) 当然の幾何学的目的を達成し (ii) 分析的手法に対する便利な幾何学的説明を提供するための仮定に過ぎない。我々が幾何学的な言葉を使うのは、主として説明のためにであり、幾何学の基礎を研究することは含まれていない。}
直線の点は、これらを含む集合を \Emph{線形統一連続体}と呼ぶことができ、これらの点は算術的統一連続体の便利な像を与える。

我々は有理数や平方根のような実数のいくつかの主要な性質について詳細に考慮した。さらにいくつかの例を挙げて、これらの特定の数のクラスがいかに特別であるか、そしてこれらは無限多様な数の集合、すなわち統一連続体のほんの一部であることを示す。

\par\medskip\noindent\textbf{【注釈】}\par
\Itemp{(i)} 以下のような複雑な根式表現を考えましょう：
\[
z = \sqrt{4 + \sqrt{15}} + \sqrt{4 - \sqrt{15}}.
\]
$z$ の表式が意味を持つとする主張の論理は、次のように考えることができます。まず、§12 で示したように、$y = \sqrt{15}$ という数が存在し、$y^2 = 15$ であることを示します。それから、§10 で述べたように、$4 + \sqrt{15}$, $4 - \sqrt{15}$ の数を定義できます。

次に、$z_1$ についての方程式
\[
z_1^3 = 4 + \sqrt{15}
\]
を考えます。この方程式の右辺は有理数ではありませんが、$x^3 = 2$（または他の任意の有理数）を持つ実数 $x$ の存在を導く同じ論理により、$z_1^3 = 4 + \sqrt{15}$ を満たす数 $z_1$ の存在も導けます。したがって、$z_1 = \sqrt{4 + \sqrt{15}}$ と定義し、同様に $z_2 = \sqrt{4 - \sqrt{15}}$ を定義します。そして、§10 で述べたように、$z = z_1 + z_2$ と定義します。

容易に確認できるのは
\[
z^3 = 3z + 8
\]
です。また、$z^3 = 3z + 8$ を満たす一意な数 $z$ の存在を直接証明することもできました。二つの異なるような数が存在しないことも容易に確認できます。もし $z_1^3 = 3z_1 + 8$ と $z_2^3 = 3z_2 + 8$ である場合、両辺を引いて $z_1 - z_2$ で割ると $z_1^2 + z_1z_2 + z_2^2 = 3$ が得られます。もし $z_1$ と $z_2$ が正である場合、$z_1^3 > 8$, $z_2^3 > 8$ より $z_1 > 2$, $z_2 > 2$ であり、したがって $z_1^2 + z_1z_2 + z_2^2 > 12$ となり、方程式は不可能です。また、$z_1$ や $z_2$ が負であることも容易に確認できます。もし $z_1$ が負で $-\zeta$ と等しい場合、$\zeta$ は正であり $\zeta^3 - 3\zeta + 8 = 0$ または $3 - \zeta^2 = 8/\zeta$ です。したがって $3 - \zeta^2 > 0$ より $\zeta < 2$ であり、これにより $8/\zeta > 4$ となり、$8/\zeta$ は $3 - \zeta^2$（これは $3$ より小さい）に等しくなることは不可能です。
\par\medskip

したがって、$z^3 = 3z + 8$ を満たす $z$ は最多1つであり、また有理数ではない。この方程式の任意の有理根は整数でなければならず、かつ8の約数である（Ex.ii.3）が、1, 2, 4, 8 のいずれも根ではないことが簡単に確認できる。

したがって、$z^3 = 3z + 8$ は最多1つの根を持ち、その根が存在すれば正で有理数ではない。現在、正の有理数 $x$ を2つのクラス $L$, $R$ に分類することができる。$x^3 < 3x + 8$ のとき $L$ クラス、$x^3 > 3x + 8$ のとき $R$ クラスとする。もし $x^3 > 3x + 8$ であり、$y$ が $x$ より大きい任意の数であれば、$y^3 > 3y + 8$ であることも容易に確認できる。なぜなら、もし仮定して $y^3 \leq 3y + 8$ とすると、$x^3 > 3x + 8$ より $y^3 - x^3 < 3(y - x)$ となり、または $y^2 + xy + x^2 < 3$ となるがこれは不可能である；$y$ は正の数であり、かつ $x > 2$（$x^3 > 8$ より）であるから。

同様に、もし $x^3 < 3x + 8$ で $y < x$ のとき、$y^3 < 3y + 8$ であることも示せる。

最後に、クラス $L$ と $R$ は両方存在することが明らかであり、これらは正の有理数または正の実数 $z$ を形成し、方程式 $z^3 = 3z + 8$ を満たす。立方方程式をカルダノ法で解くことができる読者は、方程式から $z$ の明示的な表現を得ることができるだろう。

上記で $x^{3} = 3x + 8$ の方程式に対して用いた直接的な議論は、方程式
\[
x^{5} = x + 16
\]
にも（応用は少し難しいですが）適用可能であり、これにより一意の正実数が存在することに結論付けることができます。しかし、この場合、$x$ のような単純な明示的な表現を得ることは不可能です。実際、方程式の次数が4より大きい根に対してそのような表現を見つけることは（証明は難しいですが）一般には不可能であることが証明できます。したがって、無理数は純粋または混合二次根式や他の根式、あるいはこれらの組み合わせで表されるものに加えて、代数方程式の根でありながらそのような表現で表せないものも存在します。このような表現を見つけることは非常に特別な場合に限られます。

\Itemp{(iii)} しかし、無理数のリストに方程式（例えば $x^{5} = x + 16$ のような）の根を加えても、連続体内に含まれるすべての種類の無理数は尽きません。直径が $A_{0}A_{1}$，すなわち単位長と等しい円を描きます。このような円の周長が数値的に測定可能な長さを持つことを自然に想像するのは正当です。\footnote{証明については \okrickRef{第}{VII} 章をご覧ください。} この長さは通常 $\pi$ で表されます。そして、この数 $\pi$ が任意係数の代数方程式の根ではないことが示されています\footnote{詳細は Hobson の \textit{三角法}（第3版）の pp.~305 以降を参照してください，または同じ著者の \textit{円周の正方形}（カムブリッジ, 1913年）も参照してください。}，例えば
\[
\pi^{2} = n,\quad
\pi^{3} = n,\quad
\pi^{5} = \pi + n,
\]
ここで $n$ は整数です。このようにして、有理数でもなくまたこれまで考慮した無理数のどのクラスにも属さない数を定義することが可能です。そしてこの数 $\pi$ は孤立した例や例外的な例ではありません。他の多くの例も構成できます。実際、このような方程式の根となる特別な種類の無理数は限られているだけでなく、根号で表すことができる数のクラスよりもさらに狭いです。

16. 持続的な実数変数。実数は、二つの視点から考えることができます。それらを集合として、前節で定義された「算術的連続体」として見なすことができますか、または個々のものとして見なしますか。そして、個々に考えるとするとき、特定の指定された数（たとえば $1$, $-\frac{1}{2}$, $\sqrt{2}$, または $\pi$）を思い浮かべるか、あるいは任意の数、未指定の数、数$x$を思い浮かべます。後者の場合が、「$x$は数である」「$x$は長さの測度である」「$x$は有理数または無理数である」といった主張を行うときです。これらの命題における$x$は、「持続的な実数変数」と呼ばれ、個々の数は変数の「値」呼ばれます。

しかし、変数とは必ずしも連続的でなくてもよいです。すべての実数の集合ではなく、その前集合に含まれる部分集合を考慮することができるのです。たとえば、有理数全体の集合や正の整数全体の集合などがあります。最後の場合を考えましょう。任意の正の整数、または未指定の正の整数についての命題（たとえば「$n$は奇数か偶数である」）において、$n$は変数として呼ばれ、「正の整数変数」として定義され、個々の正の整数がその値となります。

自然に `$x$' と `$n$' は変数の例であり、`$x$' がすべての実数からなる「変域」を持つ変数、または正整数からなる「変域」を持つ変数の例です。これらは最も重要な例ですが、しばしば他のケースを考慮する必要があります。小数論においては、例えば `$x$' を任意の数を小数表記で表す際に使用される任意の桁に表します。すると `$x$' は変数ですが、10 つの異なる値しか持たない変数であり、即ち $0$, $1$, $2$, $3$, $4$, $5$, $6$, $7$, $8$, $9$ のいずれかです。読者は異なる変域を持つ他の例を思い浮かべるべきです。普通の生活の中で興味深い例を見つけることができます：巡査 `$x$'、タクシー運転手 `$x$'、年 `$x$'、週の `$x$ 番目の日。これらの変数の値は自然に数ではありません。

17. 実数の「断片」
第4節から第7節では有理数の「断片」について考慮しました。すなわち、有理数（または正の有理数のみ）を二つのクラス $L$ と $R$ に分類する方法で、以下の特性を持っています：

(i) 考慮されるすべての数は二つのクラスのいずれか一つに属します。

(ii) 両方のクラスが存在します。

(iii) $L$ の任意の要素は $R$ の任意の要素よりも小さいです。

これらの考え方は、すべての実数の集合に対して適用可能であり、読者は後述する章でその過程が非常に重要なことを見つけるでしょう。

Let us then suppose that \( P \) and \( Q \) are two properties which are mutually exclusive, and one of which is possessed by every real number. Further let us suppose that any number which possesses \( P \) is less than any which possesses \( Q \). We call the numbers which possess \( P \) the \emph{lower} or \emph{left-hand class} \( L \), and those which possess \( Q \) the \emph{upper} or \emph{right-hand class} \( R \).

\par\medskip\noindent\textbf{【注釈】}\par
Thus \( P \) might be \( x \leq \sqrt{2} \) and \( Q \) be \( x > \sqrt{2} \). It is important to observe that a pair of properties which suffice to define a section of the rational numbers may not suffice to define one of the real numbers. This is so, for example, with the pair `\( x < \sqrt{2} \)' and `\( x > \sqrt{2} \)' or (if we confine ourselves to positive numbers) with `\( x^2 < 2 \)' and `\( x^2 > 2 \)'. Every rational number possesses one or other of the properties, but not every real number, since in either case \( \sqrt{2} \) escapes classification.
\par\medskip

There are now two possibilities. Either \( L \) has a greatest member \( l \), or \( R \) has a least member \( r \). Both of these events cannot occur. For if \( L \) had a greatest member \( l \), and \( R \) a least member \( r \), the number \( \frac{1}{2}(l + r) \) would be greater than all members of \( L \) and less than all members of \( R \), and so could not belong to either class. On the other hand one event must occur.

Let $L_{1}$ と $R_{1}$ を、$L$ および $R$ の有理数のみからなるクラスとする。このとき、$L_{1}$ および $R_{1}$ は有理数の区間を形成する。現在二つの場合がある。

まず、$L_{1}$ が最大の要素 $\alpha$ を持つ場合がある。この場合、$\alpha$ もまた $L$ の最大の要素である。なぜなら、もしそうではないとすると、より大きな要素 $\beta$ が存在する。$\alpha$ と $\beta$ の間に有理数が存在し、これらは $\beta$ より小さいので $L$ に属し、従って $L_{1}$ にも属している。これは明らかに矛盾である。したがって、$\alpha$ は $L$ の最大の要素である。

一方で、$L_{1}$ が最大の要素を持たない場合もある。この場合、$L_{1}$ および $R_{1}$ によって形成される有理数の区間は実数 $\alpha$ となる。この数 $\alpha$ は $L$ または $R$ のいずれかに属する。もし $\alpha$ が $L$ に属すると、以前と同様にして $\alpha$ は $L$ の最大の要素であることが示される。同様に、$\alpha$ が $R$ に属すると、それは $R$ の最小の要素となる。

Thus、いずれの場合でも$L$には最大の要素が存在するか、$R$には最小の要素が存在します。したがって、実数全体に対する任意の分割は、有理数全体に対する分割が時としてしか対応しないように、実数に対する分割が実数に対応するという意味で、しばしばですが必ずしもそうではありません。この結論は非常に重要です；なぜなら、実数全体に対する分割を考慮しても、私たちの数の概念へのさらなる一般化には至らないからです。有理数から出発して、有理数全体に対する分割が新たな数的概念、すなわち実数という概念へと導くことを発見しました。これは有理数よりも広範な概念でした；そして、実数全体に対する分割がさらに一般的な概念を導くと考えられました。しかし、前文の議論はこの考えを否定し、実数全体の集合、すなわち連続体には、有理数全体の集合が欠いていたような種類の完全性があり、技術的な言葉で言えば、連続体は閉じていると表現されます。

我々が証明した結果は次のように述べることができます：

\par\medskip\noindent\textbf{【定理 (デデキントの定理。)】}\par
実数を二つのクラス$L$と$R$に分ける場合、

\Itemp{(i)} すべての数がこの二つのクラスのいずれか一つに属する，

\Itemp{(ii)} 各クラスには少なくとも一つの数が含まれる，

\Itemp{(iii)} $L$の任意の要素は$R$の任意の要素よりも小さい， \\
という条件を満たすとき、ある数$\alpha$が存在し、すべての$\alpha$未満の数は$L$に属し、すべての$\alpha$より大きい数は$R$に属します。この数$\alpha$自体はどちらのクラスにも属する可能性があります。
\par\medskip

\par\medskip\noindent\textbf{【注釈】}\par
すべての数ではなく、ある区間 \(\DPmod{(\beta, \gamma)}{[\beta, \gamma]}\) 内の全ての数 \(x\) に対して考慮される断片がしばしばあります。つまり、\(\beta \leq x \leq \gamma\) の範囲内のすべての数 \(x\) を対象とする断片です。このような数の断片は当然、性質 (i), (ii), および (iii) を満たす二つのクラスに分割されます。このような断片は区間 \(\DPmod{(\beta, \gamma)}{[\beta, \gamma]}\) の実数全体の断片として変換することができます。これには、\(L\) に \(\beta\) より小さいすべての数を追加し、\(R\) に \(\gamma\) より大きいすべての数を追加することで実現します。 Dedekind の定理において述べられた結論は、'区間 \(\DPmod{(\beta, \gamma)}{[\beta, \gamma]}\) の実数'と '実数' を置き換えることで依然として成り立ちますし、この場合の数 \(\alpha\) は不等式 \(\beta \leq \alpha \leq \gamma\) を満たします。
\par\medskip

\Paragraph{18. 積聚点}
任意の方法で定義された実数の系またはそれに対応する直線上の点の系を、実数の集合または点の集合と呼びます。例えば、すべての正整数やすべての有理点が含まれる集合などが考えられます。

ここでは幾何学的な言葉を使用するのが最も便利です。\footnote{読者はこの方法が単なる語彙上の利便性のために採用されていることに気づくでしょう。}
したがって、与えられた点の集合 \(S\) について考えてみましょう。任意の点 \(\xi\) を取ります。これは \(S\) の要素である必要はありません。
その場合、二つの可能性があります。一つは (i) \(\delta\) を正の数として選べる場合で、区間 \(\DPmod{(\xi - \delta, \xi + \delta)}{[\xi - \delta, \xi + \delta]}\) が \(S\) の他の点を含まない、つまりその区間に \(\xi\) 変以外の点がない場合、または (ii) このような \(\delta\) を選べない場合です。

\par\medskip\noindent\textbf{【注釈】}\par
例えば、$S$ がすべての正整数に対応する点からなる集合であるとします。もし $\xi$ が正整数であれば、$\delta$ を $1$ より小さい任意の数として (i) は成り立ちます；もしこの $\xi$ が二つの連続する正整数の中間にある場合、$\delta$ を $\frac{1}{2}$ より小さい任意の数として (i) は成り立ちます。一方で、$S$ がすべての有理数に対応する点からなる集合であるとします。このとき、どんな値の $\xi$ についても (ii) は成り立ちます；これはどのような区間にも無限にある有理数があるためです。
\par\medskip
\PageSep{30}

(ii) が真であると仮定しましょう。その場合、どの長さでも小さい任意の区間 $\DPmod{(\xi - \delta, \xi + \delta)}{[\xi - \delta, \xi + \delta]}$ 内には少なくとも一つの $S$ の点 $\xi_{1}$ が存在し、これは $\xi$ と一致しない；そしてこれは $\xi$ が $S$ の要素であるかどうかに関係なく成り立ちます。この場合、$\xi$ を $S$ の \Emph{累積点} と呼びます。容易に理解できるように、区間 $\DPmod{(\xi - \delta, \xi + \delta)}{[\xi - \delta, \xi + \delta]}$ 内には単なる一つではなく無限にある $S$ の点が存在することです。つまり、$\xi_{1}$ を決定した後でも、区間 $\DPmod{(\xi - \delta_{1}, \xi + \delta_{1})}{[\xi - \delta_{1}, \xi + \delta_{1}]}$ を $\xi$ の周りに選び、$\xi_{1}$ まで達しないようにします。この区間もまた点 $\xi_{2}$（これは $S$ の要素で $\xi$ と一致しない）を含むことがわかります。明らかに、$\xi_{2}$ の位置に $\xi_{1}$ を置き換えて同じ議論を繰り返すことができます；そして無限に続けていけることから、次のような点を任意の数だけ決定できます：
\[
\xi_{1},\quad \xi_{2},\quad \xi_{3},\ \dots
\]
これらはすべて $S$ の要素であり、区間 $\DPmod{(\xi - \delta, \xi + \delta)}{[\xi - \delta, \xi + \delta]}$ 内に存在します。

累積点 $\xi$ は必ずしも $S$ の要素であるとは限らない。以下の例はその可能性を示しています。

\par\medskip\noindent\textbf{【例 IX.】}\par
\Item{1.} $S$が正の整数に対応する点やすべての整数に対応する場合、極限点は存在しない。

\Item{2.} $S$がすべての有理数に対応する場合、直線上の任意の点が極限点である。

\Item{3.} $S$が$\frac{1}{1}$, $\frac{1}{2}$, $\frac{1}{3}, \dots$に対応する場合、極限点は原点のみである。

\Item{4.} $S$がすべての正の有理数に対応する場合、極限点は原点と直線上のすべての正の点である。
\par\medskip

\Paragraph{19. ヴェイエルシュトラスの定理。} 点集合に関する一般理論は解析学の高度な分野において最大限の興味と重要性を持つが、本書のような本には大部分が難しすぎると考えられる。しかし、デデキントの定理から容易に導き出せる1つの基本的な定理があり、後で必要となるため紹介する。

\par\medskip\noindent\textbf{【定理】}\par
点集合$S$が無限個の点を含み、区間$\DPmod{(\alpha, \beta)}{[\alpha, \beta]}$内に完全に存在する場合、少なくとも区間に1つの極限点が存在する。
\par\medskip

直線$\Lambda$上の点を2つのクラスに分類する。点$P$が$L$に属するとき、$P$の右側には無数の$S$の点がある。それ以外の場合、$P$は$R$に属する。すると、デデキントの定理の条件(i)と(iii)が満たされていることが明らかであり、$\alpha$が$L$に属し、$\beta$が$R$に属しているため、条件(ii)も満たされている。

したがって、ある点$\xi$が存在し、どんなに小さな$\delta$でも$\xi - \delta$は$L$に属し、$\xi + \delta$は$R$に属する。つまり、区間$\DPmod{(\xi-\delta, \xi+\delta)}{[\xi - \delta, \xi + \delta]}$には無数の$S$の点が含まれる。したがって、$\xi$は$S$の極限点である。

この点は当然、$\alpha$または$\beta$と一致する可能性があります。例えば、$\alpha = 0$, $\beta = 1$で、$S$が点$1$, $\frac{1}{2}$, $\frac{1}{3}, \dots$からなる場合、そのような場合において、$0$は唯一の累積点となります。

第I章の補足例

\Section{第I章の補足例}

\begin{Examples}
\Item{1.} $ax + by + cz = 0$が次の条件を満たすためには：
(1) $x$, $y$, $z$に対するすべての値；
(2) $\alpha x + \beta y + \gamma z=0$を満たす$x$, $y$, $z$に対するすべての値；
(3) $\alpha x + \beta y + \gamma z = 0$と$Ax + By + Cz = 0$をともに満たす$x$, $y$, $z$に対するすべての値。

\Item{2.} 任意の正有理数は、次の形で一意的に表すことができます：
\[
a_{1} + \frac{a_{2}}{1·2} + \frac{a_{3}}{1·2·3} + \dots + \frac{a_{k}}{1·2·3\dots k},
\]
ここで、$a_{1}$, $a_{2}$, \dots, $a_{k}$は整数であり、
\[
0 \leq a_{1},\quad 0 \leq a_{2} < 2,\quad 0 \leq a_{3} < 3,\ \dots\quad 0 < a_{k} < k.
\]

\Item{3.} 任意の正有理数は、次の単純連分数で一意的に表すことができます：
\[
a_{1} + \cfrac{1}{a_{2} + \cfrac{1}{a_{3} + \cfrac{1}{\dots + \cfrac{1}{a_{n}}}}},
\]
ここで、$a_{1}$, $a_{2}, \dots$は正の整数であり、最初のみが零である可能性があります。

このような連分数に関する理論については、代数学の教科書で見ることができます。有理数および無理性数の表現方法についての詳細な情報については、ホブソン著『実変数関数論』第45-49ページを参照してください。

\Item{4.} $9x^{3} - 6x^{2} + 15x - 10 = 0$の有理根（存在する場合）を求めよ。

\Item{5.} 線分$AB$が点$C$で aurea sectione（黄金分割）に分けられているとします---すなわち、$AB·AC = BC^{2}$。このとき、比$AC/AB$は無理数であることを示せ。

直接的な幾何学的証明については、ブロムウィッチ著『無限級数』§143, p. 363を参照してください。

\Item{6.} \( A \)は無理数である。\( \frac{aA + b}{cA + d} \)，ここで \( a \)，\( b \)，\( c \)，\( d \)が有理数であるとき，この比が有理数になるのはどのような条件か？

\Item{7.} いくつかの基本的な不等式。以下，\( a_{1} \)，\( a_{2} \)，…は正の数（0を含む）であり，\( p \)，\( q \)，…は正の整数であるとする。\( a_{1}^{p} - a_{2}^{p} \)と \( a_{1}^{q} - a_{2}^{q} \)が符号相同であることから，\((a_{1}^{p} - a_{2}^{p}) (a_{1}^{q} - a_{2}^{q}) \geq 0\) であり，または
\[
a_{1}^{p+q} + a_{2}^{p+q} \geq a_{1}^{p} a_{2}^{q} + a_{1}^{q} a_{2}^{p},
\Tag{(1)}
\]
という不等式を得る。この不等式はまた次の形でも表される：
\[
\frac{a_{1}^{p+q} + a_{2}^{p+q}}{2}
  \geq \left(\frac{a_{1}^{p} + a_{2}^{p}}{2}\right)
       \left(\frac{a_{1}^{q} + a_{2}^{q}}{2}\right).
\Tag{(2)}
\]
この公式を繰り返し適用することで，次の不等式を得る：
\[
\frac{a_{1}^{p+q+r+\dots} + a_{2}^{p+q+r+\dots}}{2}
  \geq \left(\frac{a_{1}^{p} + a_{2}^{p}}{2}\right)
       \left(\frac{a_{1}^{q} + a_{2}^{q}}{2}\right)
       \left(\frac{a_{1}^{r} + a_{2}^{r}}{2}\right) \dots,
\Tag{(3)}
\]
特に，
\[
\frac{a_{1}^{p} + a_{2}^{p}}{2} \geq \left(\frac{a_{1} + a_{2}}{2}\right)^{p}.
\Tag{(4)}
\]
\( p = q = 1 \) のとき，式 (1) または \( p = 2 \) のとき，不等式は単なる不等式 \( a_{1}^{2} + a_{2}^{2} \geq 2a_{1} a_{2} \) の異なる表現であり，これは二つの正の数の算術平均がその幾何平均以上であるという事実を表している。

\Item{8.} \Topic{$n$個の数に対する一般化。} $n$個の数 $a_{1}$, $a_{2}$, \dots, $a_{n}$ から形成できる $\frac{1}{2} n(n - 1)$ 個の型 \Eq{(1)} の不等式を書き並べ、その結果を足し合わせると、以下の不等式を得ます：
\[
n \tsum a^{p+q} \geq \tsum a^p \tsum a^q,
\Tag{(5)}
\]
または
\[
\left(\tsum a^{p+q}\right)/n
  \geq \left\{\left(\tsum a^p\right)/n\right\}
       \left\{\left(\tsum a^q\right)/n\right\}.
\Tag{(6)}
\]
したがって、Eq{(3)} の明らかなる拡張形を導くことができ、読者は自分自身で述べることができます。特に以下の不等式があります：
\[
\left(\tsum a^p\right)/n \geq \left\{\left(\tsum a\right)/n\right\}^p.
\Tag{(7)}
\]

\Item{9.} \Topic{算術平均と幾何平均に関する定理の一般形。} 略して異なる性格を持つ不等式は、$a_{1}$, $a_{2}$, \dots, $a_{n}$ の算術平均がそれらの幾何平均以上であるというものです。$a_{r}$ と $a_{s}$ を $a$ の中で最大値と最小値とする（最大値や最小値が複数ある場合は任意に選ぶことができます）とします。そのとき、$G$ を $a_{r}$ と $a_{s}$ の幾何平均とします。$G > 0$ としておきます（$G = 0$ の場合、命題の真偽は明らかです）。そして $a_{r}$ と $a_{s}$ を
\[
a_{r}' = G,\quad a_{s}' = a_{r}a_{s}/G,
\]
に置き換えると、幾何平均の値は変化しません。また，
\[
a_{r}' + a_{s}' - a_{r} - a_{s} = (a_{r} - G)(a_{s} - G)/G \leq 0,
\]
であるため、算術平均は必ずしも増加しないことがわかります。

この議論を繰り返すことで、$a_{1}$, $a_{2}$, \dots, $a_{n}$ の各要素をすべて $G$ に置き換えることができます。最大で $n$ 回の繰り返しが必要です。最終的な算術平均の値が $G$ であるため、初期の値は必ずしも小さくないことが明らかです。
\PageSep{33}

10. シュワルツの不等式。任意の正または負の数 \( a_{1}, a_{2}, \dots, a_{n} \) および \( b_{1}, b_{2}, \dots, b_{n} \) に対して，以下の恒等式が成り立つ：
\[
\left(\sum a_{r} b_{r}\right)^{2}
= \sum a_{r}^{2} \sum a_{s}^{2}
- \sum (a_{r} b_{s} - a_{s} b_{r})^{2},
\]
ここで \( r \) と \( s \) は \( 1, 2, \dots, n \) の値を取る。従って，
\[
\left(\sum a_{r} b_{r}\right)^{2} \leq \sum a_{r}^{2} \sum b_{r}^{2}
\]
という不等式が成り立つ。この不等式は通常シュワルツの不等式と呼ばれ，カuchyによって最初に証明された。

11. \( a_{1}, a_{2}, \dots, a_{n} \) がすべて正であり，\( s_{n} = a_{1} + a_{2} + \dots + a_{n} \) のとき，
\[
(1 + a_{1})(1 + a_{2}) \dots (1 + a_{n})
\leq 1 + s_{n} + \frac{s_{n}^{2}}{2!} + \dots + \frac{s_{n}^{n}}{n!}.
\]
1909年数学オリンピック問題。

12. \( a_{1}, a_{2}, \dots, a_{n} \) および \( b_{1}, b_{2}, \dots, b_{n} \) が正の数であり，これらは大きさで降順に並べられているとき，
\[
(a_{1} + a_{2} + \dots + a_{n})(b_{1} + b_{2} + \dots + b_{n})
\leq n(a_{1}b_{1} + a_{2}b_{2} + \dots + a_{n}b_{n}).
\]

13. \( a, b, c, \dots, k \) および \( A, B, C, \dots, K \) が数であり，最初のセットはすべて正であるとき，
\[
\frac{aA + bB + \dots + kK}{a + b + \dots + k}
\]
は \( A, B, \dots, K \) の代数的に最小値と最大値の間にある。

14. 互いに異なる根号 \(\sqrt{p}\), \(\sqrt{q}\) に対して，\( a + b\sqrt{p} + c\sqrt{q} + d\sqrt{pq} = 0 \)
が成り立つとき，\( a, b, c, d \) が有理数であるなら \( a = 0 \), \( b = 0 \), \( c = 0 \), \( d = 0 \)。

[根号 \(\sqrt{p}\) を有理数 \( M \) と \( N \) を用いて \( M + N\sqrt{q} \) の形に表し，定理14の結果を適用する。]

15. \( a\sqrt{2} + b\sqrt{3} + c\sqrt{5} = 0 \)
が有理数 \( a, b, c \) に対して成り立つとき，\( a = 0 \), \( b = 0 \), \( c = 0 \)。

\Item{16.} $\sqrt{p}$ および $\sqrt{q}$ の任意の多項式（有理数係数を持つ、つまり有限個の項 $A(\sqrt{p})^{m}(\sqrt{q})^{n}$ （$m$ と $n$ は整数、$A$ は有理数）の和）は以下の形で表すことができます：
\[
a + b\sqrt{p} + c\sqrt{q} + d\sqrt{pq},
\]
ここで $a$, $b$, $c$, $d$ は有理数です。

\Item{17.} $\frac{a + b\sqrt{p} + c\sqrt{q}}{d + e\sqrt{p} + f\sqrt{q}}$ （$a$, $b$, 等は有理数）を以下の形で表すことができます：
\[
A + B\sqrt{p} + C\sqrt{q} + D\sqrt{pq},
\]
ここで $A$, $B$, $C$, $D$ は有理数です。

［明らかに
\begin{align*}
\frac{a + b\sqrt{p} + c\sqrt{q}}{d + e\sqrt{p} + f\sqrt{q}}
  &= \frac{(a + b\sqrt{p} + c\sqrt{q})(d + e\sqrt{p} - f\sqrt{q})}
          {(d + e\sqrt{p})^{2} - f^{2}q} \\
  &= \frac{\alpha + \beta\sqrt{p} + \gamma\sqrt{q} + \delta\sqrt{pq}}
         {\epsilon + \zeta\sqrt{p}},
\end{align*}
ここで $\alpha$, $\beta$ 等は有理数で、これらは簡単に求めることができます。必要な簡約化は、分子と分母に $\epsilon - \zeta\sqrt{p}$ をかけることで容易に行えます。例えば，
\[
\frac{1}{1 + \sqrt{2} + \sqrt{3}}
  = \frac{1}{2} + \frac{1}{4}\sqrt{2} - \frac{1}{4}\sqrt{6}.
\]

\Item{18.} $a$, $b$, $x$, $y$ が有理数であり、
\[
(ay - bx)^{2} + 4(a - x)(b - y) = 0,
\]
を満たすとき、以下のいずれかの条件が成り立ちます：
(i) $x = a$, $y = b$ または (ii) $1 - ab$ および $1 - xy$ は有理数の平方です。 \MathTrip{1903.}

\Item{19.} 
\[
ax^{2} + 2hxy + by^{2} = 1,\quad
a'x^{2} + 2h'xy + b'y^{2} = 1
\]
（$a$, $h$, $b$, $a'$, $h'$, $b'$ は有理数）の解がすべて有理数であるとき、以下の式は有理数の平方です：
\[
(h - h')^{2} - (a - a')(b - b'),\quad
(ab' - a'b)^{2} + 4(ah' - a'h)(bh' - b'h).
\]
\MathTrip{1899.}

\Item{20.} $\sqrt{2}$と$\sqrt{3}$は$\sqrt{2} + \sqrt{3}$の立方関数であり、有理係数を持つとともに、$\sqrt{2} - \sqrt{6} + 3$は$\sqrt{2} + \sqrt{3}$の線形関数の比である。1905年。

\Item{21.} 表達式
\[
\sqrt{a + 2m\sqrt{a - m^{2}}} + \sqrt{a - 2m\sqrt{a - m^{2}}}
\]
は、$2m^{2} > a > m^{2}$のとき$2m$に等しくなり、$a > 2m^{2}$のとき$2\sqrt{a - m^{2}}$に等しい。

\Item{22.} $\sqrt{2}$を有理係数を持つ多項式として表すと
\[
a + b\sqrt{2} + c\sqrt{4}
\]
の形で表せる。ここで、$a$, $b$, $c$は有理数である。

より一般に、任意の有理数$p$に対して、$\sqrt[m]{p}$を有理係数を持つ多項式として表すと
\[
a_{0} + a_{1}\alpha + a_{2}\alpha^{2} + \dots + a_{m-1}\alpha^{m-1}
\]
の形で表せる。ここで、$a_{0}$, $a_{1}, \dots$は有理数であり、$\alpha = \sqrt[m]{p}$である。任意のこのような多項式は
\[
b_{0} + b_{1}\alpha + b_{2}\alpha^{2} + \dots + b_{k}\alpha^{k}
\]
の形で表される。ここで、$b$'sは有理数である。もし$k \leq m - 1$であれば、既に必要な形になっている。もし$k > m - 1$であれば、$\alpha^{r}$を$m-1$より高い冪とする。すると$r = \lambda m + s$であり、ここで$\lambda$は整数で$0 \leq s \leq m - 1$；そして$\alpha^{r} = \alpha^{\lambda m + s} = p^{\lambda}\alpha^{s}$である。したがって、すべての$m-1$より高い冪を除去できる。

\Item{23.} $(\sqrt{2} - 1)^{5}$と$(\sqrt{2} - 1)/(\sqrt{2} + 1)$を$a + b\sqrt{2} + c\sqrt{4}$の形で表せ、ここで$a$, $b$, $c$は有理数である。[第2項の分子と分母に$\sqrt{4} - \sqrt{2} + 1$をかけてから計算する。]

\Item{24.} 
\[
a + b\sqrt{2} + c\sqrt{4} = 0
\]
が成り立つとき、$a$, $b$, $c$は有理数であり、$a = 0$, $b = 0$, $c = 0$である。
ページ分割：35

[令$y = \sqrt{2}$とすると、$y^{3} = 2$であり
\[
cy^{2} + by + a = 0.
\]
したがって、$2cy^{2} + 2by + ay^{3} = 0$または
\[
ay^{2} + 2cy + 2b = 0.
\]

Multiplying these二項方程式をそれぞれ$a$と$c$で掛けて引くと、$(ab - 2c^{2})y + a^{2} - 2bc = 0$を得て、または$y = -(a^{2} - 2bc)/(ab - 2c^{2})$となり、有理数であり、これは不可能である。唯一の選択肢は$ab - 2c^{2} = 0$, $a^{2} - 2bc = 0$である。

したがって$ab = 2c^{2}$, $a^{4} = 4b^{2}c^{2}$となる。$a$も$b$もゼロでない場合、第二の式を第一の式で割ると$a^{3} = 2b^{3}$となり：これは不可能であり、$\sqrt{2}$は有理数$a/b$と等しくならないからである。したがって$ab = 0$, $c = 0$となり、元の方程式から$a$,$b$,$c$は全てゼロとなる。

この結果として、もし$a + b\sqrt{2} + c\sqrt{4} = d + e\sqrt{2} + f\sqrt{4}$であれば、$a = d$, $b = e$, $c = f$である。

より一般的には、もし
\[
a_{0} + a_{1}p^{1/m} + \dots + a_{m-1}p^{(m-1)/m} = 0,
\]
$p$が$m$乗根の形でない場合、$a_{0} = a_{1} = \dots = a_{m-1} = 0$であることが証明できるが、その証明はより複雑である。

25. $A + \sqrt{B} = C + \sqrt{D}$であれば、$A = C$, $B = D$または$B$と$D$が両方有理数の立方である場合のみ成立する。

26. $\sqrt{A} + \sqrt{B} + \sqrt{C} = 0$であれば、$A$,$B$,$C$のいずれかがゼロであり、他の二つは等しく反対数であるか、$\sqrt{A}$,$\sqrt{B}$,$\sqrt{C}$は同じ有理数倍の根$\sqrt{X}$である。

27. 有理数$\alpha$,$\beta$を求める：
\[
\sqrtp{7 + 5\sqrt{2}} = \alpha + \beta\sqrt{2}.
\]

28. $(a - b^{3})b > 0$であれば、
\[
\bigsqrtb{a + \frac{9b^{3} + a}{3b}\bigsqrtp{\frac{a - b^{3}}{3b}}} +
\bigsqrtb{a - \frac{9b^{3} + a}{3b}\bigsqrtp{\frac{a - b^{3}}{3b}}}
\]
が有理数である。各立方根の下の数は形
\[
\left\{\alpha + \beta\bigsqrtp{\frac{a - b^{3}}{3b}}\right\}^{3}
\]
で、$\alpha$,$\beta$は有理数である。

29. $\alpha = \sqrt[n]{p}$であれば、$\alpha$の任意の多項式は次数$n$、有理係数を持つ方程式の根となる。

多項式（$x$とします）は次の形で表すことができます：
\[
x = l_{1} + m_{1}\alpha + \dots + r_{1}\alpha^{(n-1)},
\]
ここで、$l_{1}$, $m_{1}, \dots$ は有理数であり、Ex.22において同様です。
\PageSep{36}

同様に
\begin{alignat*}{4}
x^{2} \&= l}{2} \&\&+ m\\textit{{2}\alpha \&\&+ \dots \&\&+ r}{2}\alpha^{(n-1)}, \\
\multispan{9}{\dotfill} \\
x^{n} \&= l${n} &&+ m${n}\alpha \&\&+ \dots \&\&+ r\\textit{{n}\alpha^{(n-1)}.
\end{alignat*}

したがって
\[
L_{1}x + L_{2}x^{2} + \dots +  L_{n}x^{n} = \Delta,
\]
ここで、$\Delta$は行列式
\[
\left| ZZZMATHTOKEN2ZZZ \right|
\]
であり、$L_{1}$, $L_{2}, \dots$ は$l_{1}$, $l_{2}, \dots$ の余因子です。

30. この過程を$x = p + \sqrt{q}$に適用し、§14の定理を導き出します。

31. $y = a + bp^{1/3} + cp^{2/3}$が方程式
\[
y^{3} - 3ay^{2} + 3y(a^{2} - bcp) - a^{3} - b^{3}p - c^{3}p^{2} + 3abcp = 0.
\]
を満たすことを示します。

32. 代数的数について。我々は、$\sqrt{2}$のような無理数が次の型の方程式の根であることを見たことがあります：
\[
a_{0}x^{n} + a_{1}x^{n-1} + \dots + a_{n} = 0,
\]
ここで$a_{0}$, $a_{1}$, \dots, $a_{n}$は整数です。このような無理数は「代数的」な数と呼ばれ、他の無理数（例えば§15の$\pi$）は「超越的」な数と呼ばれます。$x$が代数的数である場合、任意の有理数$k$で$kx$もまた代数的数であり、任意の整数$m$と$n$で$x^{m/n}$もまた代数的数であることを示します。

33. $x$と$y$が代数的数である場合、$x + y$, $x - y$, $xy$ および $x/y$ もまた代数的数です。

We have equations
\begin{alignat*}{4}
a}{0}x^{m} \&+ a\\textit{{1}x^{m-1} \&\&+ \dots \&\&+ a}{m} \&\&= 0, \\
b\\textit{{0}y^{n} \&+ b}{1}y^{n-1} \&\&+ \dots \&\&+ b\_{n} \&\&= 0,
\end{alignat*}
where the $a$'s~and~$b$'s are integers. Write $x + y = z$, $y = z - x$ in the second, and eliminate~$x$. We thus get an equation of similar form
\[
c_{0}z^{p} + c_{1}z^{p-1} + \dots + c_{p} = 0,
\]
satisfied by~$z$. Similarly for the other cases.

34. If
\[
a_{0}x^{n} + a_{1}x^{n-1} + \dots + a_{n} = 0,
\]
where $a_{0}$, $a_{1}$, \dots,~$a_{n}$ are any algebraical numbers, then $x$~is an algebraical
number. [We have $n + 1$~equations of the type
\[
a_{0, r}a_{r}^{m_{r}} + a_{1, r}a_{r}^{m_{r}-1} + \dots + a_{m_{r}, r} = 0\quad
(r = 0,\ 1,\ \dots,\ n),
\]
in which the coefficients $a_{0, r}$, $a_{1, r}, \dots$ are integers. Eliminate $a_{0}$, $a_{1}$, \dots,~$a_{n}$
between these and the original equation for~$x$.]

35. Apply this process to the equation $x^{2} - 2x\sqrt{2} + \sqrt{3} = 0$.

[The result is $x^{8} - 16x^{6} + 58x^{4} - 48x^{2} + 9 = 0$.]
\PageSep{37}

36. Find equations, with rational coefficients, satisfied by
\[
1 + \sqrt{2} + \sqrt{3},\quad
\frac{\sqrt{3} + \sqrt{2}}
     {\sqrt{3} - \sqrt{2}},\quad
\sqrtb{\sqrt{3}+ \sqrt{2}} + \sqrtb{\sqrt{3} - \sqrt{2}},\quad
\sqrt{2} + \sqrt{3}.
\]

37. If $x^{3} = x + 1$, then $x^{3n} = a_{n}x + b_{n} + c_{n}/x$, where
\[
a_{n+1} = a_{n} + b_{n},\quad
b_{n+1} = a_{n} + b_{n} + c_{n},\quad
c_{n+1} = a_{n} + c_{n}.
\]

38. If $x^{6} + x^{5} - 2x^{4} - x^{3} + x^{2} + 1 = 0$ and $y = x^{4} - x^{2} + x - 1$, then $y$~satisfies
a quadratic equation with rational coefficients. \MathTrip{1903.}

[It will be found that $y^{2} + y + 1 = 0$.]
\par\medskip
\PageSep{38}

\chapter{第2章：Chapter II: FUNCTIONS OF REAL VARIABLES}

第2章 実変数の関数

20. 関数の概念。$x$と$y$を2つの連続な実変数とする。これらは固定点$A_{0}$、$B_{0}$から直線$\Lambda$、$\Mu$上に測定される距離$A_{0}P = x$、$B_{0}Q = y$で幾何的に表現できるものと仮定することができる。そして、点$P$と$Q$の位置は独立ではなく、$x$と$y$との間の関係によって結びついているものとする。すなわち、$P$と$x$が分かれば、$Q$と$y$も分かることになる。例えば、$y = x$、または$y = 2x$、または$\frac{1}{2}x$、または$x^{2} + 1$とすることができる。これらのすべての場合において、$x$の値は$y$の値を決定する。あるいはまた、$x$と$y$との間の関係が明示的な式で与えられるものではなく、$P$が分かれば$Q$を決定するための幾何学的構築によって与えられているものとすることができる。

このような状況では、$y$は$x$の関数であると定義される。この1変数に対する他1変数の関数的な依存性という概念は、数学全体の中で最も重要な概念の一つである。読者がこれを明確に理解していることを確認するため、本章では多くの例を用いてその概念を説明する。

しかし、これらの例に進む前に、上記の簡単な関数の3つの特性について注意しておきたい。これらは関数という一般的な概念には必ずしも含まれていないものである：

(1) $y$は任意の$x$の値に対して定義されている；

(2) $y$が与えられた$x$の値に対して、その対応する$y$の値は一つだけである；

(3) $x$と$y$との関係は分析的な式によって表現され、この式を利用して任意の$x$の値に対する$y$の値を直接計算できる。

これらの特定の特性は、多くの重要な関数に属しています。しかし、以下の例を考慮すると、これらが関数にとって必ずしも欠かせないものであるとは言えません。必要なのは、$x$ と $y$ の間に何かしらの関係があることだけです。そして、少なくともある値の $x$ に対応する $y$ の値があります。

\par\medskip\noindent\textbf{【例 X.】}\par
\Item{1.} $y = x$ や $2x$ や $\frac{1}{2}x$ や $x^{2} + 1$ です。このようなケースについては、現在は追加の説明は必要ありません。

\Item{2.} $y = 0$ であり、$x$ の値に関係なくそのようにします。このとき $y$ は $x$ の関数であり、$x$ にどのような値を与えても、対応する $y$ の値（つまり $0$）が既知です。この場合、関数の関係性はすべての $x$ の値に対し同じ $y$ の値を対応させます。$y$ が代わりに $1$ や $-\frac{1}{2}$ や $\sqrt{2}$ であったとしても同様です。このような $x$ の関数は「定数」と呼ばれます。

\Item{3.} $y^{2} = x$ とします。すると、$x$ が正の値である場合、この方程式は各 $x$ の値に対し二つの $y$ の値を定義します（つまり $±\sqrt{x}$）。$x = 0$ のとき、$y = 0$ です。したがって、特定の $x$ の値 $0$ には一つだけの $y$ の値が対応します。しかし、$x$ が負の値である場合、方程式を満たす $y$ の値はありません。つまり、関数 $y$ は負の $x$ の値に対して定義されていません。この関数は特性 (3) を持つものの、特性 (1) や (2) は持ちません。

\Item{4.} 恒温で維持されたガスを密閉したシリンダーに収納し、その上部には滑動式の Piston が設置します。\footnote{この有用な例は H. S. Carslaw の『微分積分入門』から借用しました。}

$A$ を Piston の横截面積とし、$W$ をその重さとします。ガスは Piston によって圧縮され、一定の圧力 $p_{0}$（単位面積あたり）を Piston 上に作用させます。これは重さ $W$ と平衡しており、
\[
W = Ap_{0}
\]
となります。

Let $v_{0}$ がシステムが均衡状態にあるときのガスの体積である。
追加の重さを Piston 上に置くと、Piston は下に押し下げられる。
ガスの体積~($v$) が減少し、Piston 単位面積に対するガスの圧力~($p$) が増加する。ボイルの実験法則は $p$ と $v$ の積が非常に近似して一定であることを主張しており、これが正確であれば、その関係は次のような方程式で表される：
\[
pv = a,
\Tag{(i)}
\]
ここで $a$ は実験によって大まかに定まる数である。

しかし、ボイルの法則はガスが十分に圧縮されない限り、事実に対して適切な近似を与えるだけである。$v$ が減少し、$p$ が一定点を超えて増加すると、それらの関係は方程式~\Eq{(i)} で十分正確には表現されなくなる。既知によれば、真の関係をよりよく近似するものは、`van der Waals の法則' と呼ばれるもので、次のような方程式によって表される：
\[
\left(p + \frac{\alpha}{v^{2}}\right)(v - \beta) = \gamma,
\Tag{(ii)}
\]
ここで $\alpha$、$\beta$、$\gamma$ は実験によって大まかに定まる数である。

もちろん、これらの2つの方程式を合わせても、$p$ と $v$ の関係についての完全な説明にはならない。この関係は現実にはもっと複雑であり、$v$ が変化するにつれてその形が、\Eq{(i)} の形から \Eq{(ii)} の形に近づくなどして変化する。しかし、数学的な観点からは、ある特定の値 $V$ よりも大きいすべての $v$ 値に対して \Eq{(i)} が正確に成り立つ理想状態と、それより小さいすべての $v$ 値に対して \Eq{(ii)} が正確に成り立つ状態を想像することは問題ない。そして、この2つの方程式は $p$ を $v$ の関数として定義するためのものと考えることができる。これは一部の $v$ の値では一つの式で定義され、他の $v$ の値では別の式で定義される関数の例である。

この関数は任意の$v$に対してのみ一つだけ$p$の値が対応し、しかし(1)には該当しない。実際、$v$が負の値を持つとき、$p$は$v$の関数として定義されていない；「負の体積」は意味を持たず、したがって負の$v$の値は考慮されることがない。

5. ある完全に弾性のある球が高さ$\frac{1}{2}g\tau^{2}$から固定された水平面に落下し（回転なし）、継続的に反発するという状況を考えてみよう。初等力学の通常の公式、読者にはおそらく馴染みがあるものであるが、$0 \leq t \leq \tau$ のとき $h = \frac{1}{2}gt^{2}$、$\tau \leq t \leq 3\tau$ のとき $h = \frac{1}{2}g(2\tau - t)^{2}$、一般に $(2n - 1)\tau \leq t \leq (2n + 1)\tau$ のとき
\[
h = \tfrac{1}{2}g(2n\tau - t)^{2}
\]
であり、$h$は時間$t$における球の位置（元の位置から下にどれだけあるか）である。明らかに$h$は正の$t$に対してのみ定義された関数である。

6. $y$を$x$の最大の素因数として定義する。これは特定の$x$の値に対する定義であり、つまり整数値に対するものである。$\frac{11}{3}$や$\sqrt{2}$や$\pi$の最大の素因数は意味を持たず、したがってこれらの$x$の値に対して定義関係は機能しない。したがってこの関数には(1)の特性がない。しかし(2)には該当し、(3)には該当しない。なぜなら$y$を$x$の単純な式で表す簡単な公式がないからである。

7. $y$を$x$を既約分数で表したときの分母として定義する。これは$x$が有理数である場合にのみ定義される関数の例である。つまり、$x = -\frac{11}{7}$のとき$y = 7$であるが、$x = \sqrt{2}$に対しては$y$は定義されず、「$\sqrt{2}$の分母」は意味のない表現である。

\Item{8.} 派出所 \( Cx \) のメトロポリタン警察の警官の 1907 年 8 月 8 日午後 5 時 30 分におけるインチ単位の身長 \( y \) を定義します。すると、\( y \) は特定の整数値の \( x \) に対して定義され、具体的には \( 1, 2, \dots, N \) の各値に対応し、ここで \( N \) はその時刻における分部 \( C \) の警官の総数です。

\Paragraph{21. 関数のグラフ的表現。} 変数 \( y \) が変数 \( x \) の関数であると仮定します。この関係は、\( x \) が \( y \) の関数でもあります。しかし、現在は最初の視点からこの関係を観察することにします。すると、\( x \) を独立変数、\( y \) を従属変数と呼びます。具体的な関数形式が指定されていない場合、それを \( y = f(x) \)（または \( F(x), \phi(x), \psi(x), \dots \)）で表します。

特定の関数の性質は、多くの場合、次のように説明しやすくすることができます。直角に交差した無限延長された線 \( OX, OY \) を引き、\( x \) の値と \( y \) の値をそれぞれ \( OX, OY \) 上の距離で表します。矢印が示す正方向を考慮し、\[ \Fig{6} \] に示すようにします。

任意の \( x \) の値 \( a \) に対して \( y \) が定義され、その値は \( b \)（仮定）であるとします。すると、\( OA = a, OB = b \) とし、長方形 \( OAPB \) を完成させます。図に点 \( P \) を示すと、これは \( x = a \) のときの \( y \) の値が \( b \) であることを示しています。

\( x \) の値 \( a \) に対して複数の \( y \) の値（例えば \( b, b', b'' \））が対応する場合、点 \( P \) ではなく、複数の点 \( P, P', P'' \) を持つことになります。

点 \(P\) を座標系の軸 \(OX\), \(OY\) に対して座標 \((a, b)\) の点と呼びます；\(a\) は点 \(P\) の \(x\) 座標（横座標）を、\(b\) は点 \(P\) の \(y\) 座標（縦座標）を表します。軸 \(OX\) と \(OY\) をそれぞれ \(x\) 軸、\(y\) 軸と呼び、これらを合わせて座標軸と呼びます。また、原点 \(O\) を座標の原点と呼びます。

次に、すべての \(x\) の値（定義されている \(y\) の値に対応する）に対して、\(b\) （または \(b'\), \(b''\), \(\dots\)）という \(y\) の値が決定され、その対応する点 \(P\) （または \(P'\), \(P''\), \(\dots\)）も定められたとします。これらのすべての点の集合を関数 \(y\) のグラフと呼びます。

非常に単純な例として、\(y\) が方程式
\[
Ax + By + C = 0
\]
によって \(x\) の関数として定義されているとしましょう（ここで \(A\), \(B\), \(C\) は任意の固定された数）。\footnote{もし \(B = 0\) の場合、\(y\) は方程式に現れません。この場合、\(y\) は \(x = -C/A\) のみで定義され、すべての値を持つ関数とみなす必要があります。}

このとき、\(y\) はすべての特性（1）, (2), (3) を持つ \(x\) の関数です。容易に示されるように、関数 \(y\) のグラフは一直線であることがわかります。読者はおそらく、解析幾何の教科書で与えられるさまざまな証明の一つまたは複数を知っているでしょう。

また、別の表現法を使用することもあります。\(x\) と \(y\) が方程式 \Eq{(1)} を常に満たすように変化するとき、点 \((x, y)\) の軌跡は一直線であると言います。そして、方程式 \Eq{(1)} をこの軌跡の方程式と呼びます。また、方程式が軌跡を表していると言います。これらの用語 `軌跡', `軌跡の方程式' の使用は非常に一般的であり、\(x\) と \(y\) の関係が解析式で表現できる場合に適用できます。

方程式 $Ax + By + C = 0$ は \emph{一次式の一般形} です。$Ax + By + C$ は $x$ と $y$ の次数が1以下である最も一般的な多項式であり、したがって \emph{一次式の一般形は一直線を表す} ことが分かります。任意の一辺の方程式が一次式であることも同様に証明できます。

さらに興味深い幾何学的図形のいくつかの例を挙げましょう。形式が
\[
(x - \alpha)^{2} + (y - \beta)^{2} = \rho^{2}
\]
または
\[
x^{2} + y^{2} + 2Gx + 2Fy + C = 0,
\]
ここで $G^{2} + F^{2} - C > 0$ の条件下では、円を表します。方程式
\[
Ax^{2} + 2Hxy + By^{2} + 2Gx + 2Fy + C = 0
\]
(\emph{二次式の一般形})は、係数が特定の不等式を満たすと仮定すると、楕円、放物線、双曲線のいずれかを表します。これらの図形についてさらに詳しく議論するには、解析幾何学の本を参照してください。

\paragraph{22. 極座標。}これまでのところ、点$P$の位置は直交座標$OM = x$, $MP = y$の長さで決定しました。もし$OP = r$とし、$\angle MOP = \theta$で$\theta$が$0$から$2\pi$までの角度（正方向で測定）であるとするとき、明らかに
\begin{gather*}
x = r\cos\theta,\qquad y = r\sin\theta, \\
r = \sqrt{x^{2} + y^{2}},\quad \cos\theta : \sin\theta : 1 :: x : y : r,
\end{gather*}
が成り立ち、点$P$の位置は$r$と$\theta$を知ることで同様に決定できます。$r$と$\theta$を極座標（polar coordinates）と呼びます。なお、前者は基本的には正であることに注意してください。\footnote{極座標は時として$r$が正または負の値を持つように定義されることがあります。この場合、$(1, 0)$と$(-1, \pi)$は同じ点に対応します。二つのシステムの違いを示すために、方程式$l/r = 1 - e\cos\theta$（$l > 0$, $e > 1$）を考えます。当時の定義では$r$は正であるため$\cos\theta < 1/e$であり、この方程式は双曲線の一方の枝を表し、もう一方の枝は$-l/r = 1 - e\cos\theta$で与えられます。$r$が負の値を持つシステムでは、方程式は全双曲線を表します。}
\%[図7]
\figure[1.75in]{7}{p043}

もし点$P$がある軌跡上を動くとすると、$r$と$\theta$の間に何らかの関係があることになります。例えば$r = f(\theta)$または$\theta = F(r)$のような関係があります。これをその軌跡の極方程式（polar equation）と呼びます。$(x, y)$座標の方程式から（あるいは逆に）極方程式を導くためには、上記の公式を使用します。

極座標における直線の方程式は以下の形で表されます：
\[
r\cos(\theta - \alpha) = p,
\]
ここで $p$ と $\alpha$ は定数です。方程式 $r = 2a\cos\theta$ は原点を通る円を表し、一般の円の方程式は以下の形で表されます：
\[
r^{2} + c^{2} - 2rc\cos(\theta - \alpha) = A^{2},
\]
ここで $A$, $c$, そして $\alpha$ は定数です。

\#\#\# 23. 関数のさらなる例とそのグラフ的表現
以下の例は、読者に可能な関数の多様な種類についてより良い理解を提供します。

\#\#\#\# A. 多項式
$x$ の \emph{多項式} は以下のような関数です：
\[
a_{0}x^{m} + a_{1}x^{m-1} + \dots + a_{m},
\]
ここで $a_{0}$, $a_{1}$, \dots, $a_{m}$ は定数です。最も単純な多項式は単純なべき乗 $y = x$, $x^{2}$, $x^{3}$, \dots, $x^{m}, \dots$ です。関数 $x^{m}$ のグラフは、$m$ が偶数か奇数かによって異なる2つの種類があります。

まず $m = 2$ の場合を考えます。このとき、グラフ上の3点は $(0, 0)$, $(1, 1)$, $(-1, 1)$ です。他の特別な値を $x$ に代入することで、さらに多くの点が見つかります：例えば
\[
x = \frac{1}{2}, \quad 2, \quad 3, \quad -\frac{1}{2}, \quad -2, \quad 3
\]
はそれぞれ
\[
y = \frac{1}{4}, \quad 4, \quad 9, \quad \frac{1}{4}, \quad 4, \quad 9
\]
を与えます。読者がこれらの点をグラフに描くと、その形が \Fig{8} のようなものであることが予想されます。特別な点を通る曲線を引いて、それをグラフの一部として確認し、新しい $x$ の値を与えて $y$ の対応する値を求めることで、その曲線は描かれた点に非常に近いことが確認できます（描画上の誤差を考慮すると）。この曲線は当然パラボラです。
\%[図表: Fig. 8.]
\begin{figure}[2.25in]{8}{p044}

ただし、現在ではこの根本的な質問に充分な答えを出すことはできません。読者はおそらく、連続曲線とは何かという概念についての何らかの理解を持っているでしょう。つまり、断絶や跳躍がない曲線、図8で概略的に示されているような曲線です。関数$y = x^{2}$のグラフが実際にはそのような曲線であるかどうかについては、単に曲線上のいくつかの孤立した点を構築して証明することはできません。しかし、それらの点を多く構築すればするほど、それがそうであるように思われる可能性は高まります。

この質問は、第V章で適切に議論することができます。その章では、私たちの直感的な連続性の概念が何を意味しているのか、そしてなぜ現在考慮されているグラフや、本章後半でさらに考慮する他のグラフが実際には連続曲線であることをどのように証明できるかについて詳細に検討します。現時点では、読者は直感に基づいてそのような曲線を描くことに満足すればよいでしょう。

また、次のように注意しておきます。
関数$y = x^{2}$は$x$軸に対して常に凸であることが容易に確認できます。図8の点$P_{0}$, $P_{1}$は$(x_{0}, x_{0}^{2})$, $(x_{1}, x_{1}^{2})$です。直線$P_{0}P_{1}$上の点の座標は$x = \lambda x_{0} + \mu x_{1}$, $y = \lambda x_{0}^{2} + \mu x_{1}^{2}$で、$\lambda$と$\mu$は合計が1である正の数です。そして
\[
y - x^{2}
  = (\lambda + \mu)(\lambda x_{0}^{2} + \mu x_{1}^{2})
  - (\lambda x_{0} + \mu x_{1} )^{2}
  = \lambda\mu(x_{1} - x_{0})^{2} \geq 0,
\]
よって直線は曲線の上側に完全に位置します。

曲線 $y = x^4$ は、一般に $y = x^2$ に似ていますが、点 $O$ 近傍では平らで、点 $A$, $A'$ 超過では急傾斜になります（図9参照）。また、$m$ が偶数かつ $4$ より大きいときの $y = x^m$ はさらにその傾向が強まります。$m$ の値が大きくなるにつれて、平らさと急峻さはますます顕著になりますし、最終的には曲線は図中の太線から実質的に区別できなくなります。

次に $m$ が奇数のときの曲線 $y = x^m$ を考えるべきです。偶数の場合と異なる基本的な点は、$m$ が偶数のとき $(-x)^m = x^m$ であり、曲線は $OY$ 軸に関して対称であるのに対し、$m$ が奇数のとき $(-x)^m = -x^m$ で、したがって $x$ が負のとき $y$ も負になります（図10参照）。図10は $y = x$, $y = x^3$ および $m$ の奇数値が大きいときに $y = x^m$ に近づく形を示しています。

これで、任意の多項式のグラフを理論的にはどのように描けるかが分かります。まず、$y = x^m$ のグラフから定数 $C$ を用いて $Cx^m$ のグラフを得ることができます。これは曲線上の各点の縦座標に $C$ を乗じることで得られます。また、$f(x)$ と $F(x)$ のグラフが分かっているとき、$f(x) + F(x)$ のグラフは、それぞれの点に対応する縦座標の和を取ることで求めることができます。

多項式のグラフを描く作業は、後述するより高度な方法を使用することで大幅に容易になりますので、ここではこの問題についてはさらに詳しく触れません。

例題 XI.
1. 曲線 $y = 7x^4$, $y = 3x^5$, $y = x^{10}$ を描きなさい。

読者は慎重に曲線を描き、すべての曲線を一つの図形に描いてください。
\footnote{
$y$軸方向の測定スケールは$x$軸方向よりも大幅に小さく設定することが便利で、図が不自然な大きさにならないようにします。}
そうすると、$x$が大きいほど$x$の高次幂が急速に増加することを理解し、多項式
\[
x^{10} + 3x^{5} + 7x^{4}
\]
（またはさらには $x^{10} + 30x^{5} + 700x^{4}$）においては、$x$が相当大きいとき、最初の項が実質的に優位となることがわかります。したがって、$x = 4$ のとき、$x^{10} > 1,000,000$ ですが、$30x^{5} < 35,000$ で $700x^{4} < 180,000$；また、$x = 10$ のとき、最初の項の優位性はより顕著です。

2. $x = 1$, $10$, $100$ 等の場合、$x^{12}$, $1,000,000x^{6}$, $1,000,000,000,000x$ の相対的な大きさを比較してください。

読者はこのようなタイプの多くの例を作成することをお勧めします。異なる$x$の関数の \emph{相対的な増加率} という概念は、次の章でしばしば扱うものです。

3. $ax^{2} + 2bx + c$ のグラフを描いてください。
\[
y - \frac{(ac - b^{2})}{a} = a\left(x + \frac{b}{a}\right)^{2}
\]
新しい軸は古い軸と並行し、点 $x = -\frac{b}{a}$, $y = \frac{(ac - b^{2})}{a}$ を通ります。新しい方程式は $y' = ax'^{2}$ です。この曲線は放物線です。

4. $y = x^{3} - 3x + 1$, $y = x^{2}(x - 1)$, $y = x(x - 1)^{2}$ の曲線を追跡してください。

\Paragraph{24.} \Topic{\Item{B.} 有理関数。}
次の単純さと重要性において次に来るクラスの関数は、\emph{有理関数}です。有理関数は一つの多項式を別の多項式で割った商です：したがって、$P(x)$, $Q(x)$ が多項式であるとき、一般的な有理関数を
\[
R(x) = \frac{P(x)}{Q(x)}
\]
と表記します。

特定の場合、$Q(x)$が1または他の任意の定数（すなわち$x$を含まない）に簡約されるとき、$R(x)$は多項式になります：従って、有理関数のクラスには多項式のクラスがサブクラスとして含まれます。以下に定義に関する注意すべき点があります。

\par\medskip\noindent\textbf{【注釈】}\par
\Item{(1)} 通常、$P(x)$と$Q(x)$が共通因子$x + a$または$x^{p} + ax^{p-1} + bx^{p-2} + \dots + k$を持つことを仮定し、これらの共通因子は除かれます。

\Item{(2)} しかし、共通因子の除去は通常関数を変更します。例えば関数$x/x$を考えましょう。これは有理関数です。共通因子$x$を取り除くと、$1/1 = 1$となります。しかし、元の関数は常に$1$に等しいわけではありません：$x \neq 0$ の限り$1$に等しく、$x = 0$ のとき無意味な形$0/0$になります。したがって、関数$x/x$は$x \neq 0$のとき$1$に等しく、$x = 0$のときは定義されません。従って、常に$1$に等しい関数$1$とは異なります。

\Item{(3)} 次のような関数
\[
\left(\frac{1}{x + 1} + \frac{1}{x - 1}\right) \bigg/
\left(\frac{1}{x} + \frac{1}{x - 2}\right)
\]
は、通常の代数学の法則により標準形
\[
\frac{x^{2}(x - 2)}{(x - 1)^{2} (x + 1)}
\]
に簡約できます。しかし、再び注意すべき点として、この簡約が常に正当であるとは限らないことを指摘します。関数の値を特定の$x$の値で計算するには、その形そのままの関数に$x$の値を代入しなければなりません。この関数の場合、$x = -1, 1, 0, 2$ のすべての値は無意味な表現を導き出し、したがってこれらの値では定義されません。簡約形も同様に$x = -1, 1$については無意味になりますが、$x = 0, x = 2$の場合には値$0$を与えます。従って、再びこの二つの関数は同じものではありません。

\Item{(4)} しかし、(3)で考慮した特定の例から明らかであるように、関数が一般に標準形の有理関数として簡約されても、ある一定の$x$の値では定義されない場合があります。これらは分母が消滅する$x$の値（存在する場合）です。例えば、$(x^{2} - 7)/(x^{2} - 3x + 2)$は$x = 1$または$2$で定義されません。

\Item{(5)} 表現法として(2)と(3)で考慮したようなものについては、その簡約化過程が正当でない$x$の例外的な値を無視し、関数を標準形の有理関数に簡約することに一般に同意します。読者はこの理解に基づいて、二つの有理関数の和、積、または商もまた標準型の有理関数として簡約できることが容易に確認できます。そして一般的には「有理関数の有理関数は自身が有理関数である」と言えます：すなわち、$z = P(y)/Q(y)$において$P$と$Q$が多項式であり、$y = P_{1}(x)/Q_{1}(x)$を代入した後、簡約化により$z = P_{2}(x)/Q_{2}(x)$の形になります。

\Item{(6)} 有理関数の定義においては、係数として現れる定数が有理数である必要性はありません。有理という言葉は、変数$x$が関数内でどのように現れるかにのみ関連しています。例えば
\[
\frac{x^{2} + x + \sqrt{3}}{x\sqrt{2} - \pi}
\]
は有理関数です。

有理関数という言葉の由来は次の通りです。$P(x)/Q(x)$ のような有理関数は、$x$ に対して有限回の操作で生成できます。これらの操作には、$x$ を自身と定数との積にすること、得られた項を加算すること、そしてそれらから得られた関数同士の商を行うもののみが含まれます。変数 $x$ の観点からは、このような手続きはすべての有理数が単位元から生成される方法と非常に似ています。これは例えば以下の式で示されます：
\[
\frac{5}{3} = \frac{1 + 1 + 1 + 1 + 1}{1 + 1 + 1}.
\]

また、上記の操作を使用して $x$ から導き出された関数を各段階で使用し、それらが同じ方法で $x$ から得られた場合に、任意の関数も標準的な有理関数の型に帰着できます。このような手続きにより得られる最も一般的な種類の関数は、以下の例によって十分に示されます：
\[
\Biggl(\frac{x}{x^{2} + 1}
    + \frac{2x + 7}{x^{2} + \dfrac{11x - 3\sqrt{2}}{9x + 1}}\Biggr) \Bigg/
\left(17 + \frac{2}{x^{3}}\right),
\]
これは明らかに標準的な有理関数の型に帰着できます。

さらに、有理関数のグラフを描くことは、多項式の場合よりも微分積分学に基づいた方法を使用することで著しく容易になります。したがって、現在は非常に少ない例だけを挙げることに留めます。

以下が例示です：

1. $y = 1/x$, $y = 1/x^{2}$, $y = 1/x^{3}$，\dots のグラフを描きましょう。

（図は最初の二つの曲線のグラフを示しています。$1/0$，$1/0^{2}$，\dots は意味不明な表現であるため、これらの関数は $x = 0$ で定義されません。）

\item 2. $y = x + \frac{1}{x}$, $x - \frac{1}{x}$, $x^2 + \frac{1}{x^2}$, $x^2 - \frac{1}{x^2}$, $ax + \frac{b}{x}$ の軌跡を、$a$ と $b$ に正の値や負の値を代入して描きなさい。

\item 3. $y = \frac{x + 1}{x - 1}$, $\left(\frac{x + 1}{x - 1}\right)^2$, $\frac{1}{(x - 1)^2}$, $\frac{x^2 + 1}{x^2 - 1}$ の軌跡を描きなさい。

\item 4. $y = \frac{1}{(x - a)(x - b)}$, $y = \frac{1}{(x - a)(x - b)(x - c)}$ を $a < b < c$ として、軌跡を描きなさい。

\item 5. $m$ が大きくなるにつれて、曲線 $y = \frac{1}{x^m}$ の一般形を描きなさい。また、$m$ が奇数か偶数である場合分けをして考えなさい。

\paragraph{26.} \topic{\item C. 明示的代数関数。} 次に重要な関数のクラスは、\emph{明示的代数関数} です。これらは $x$ から有限個の操作（有理関数を生成するための操作と同様）と有限個の根引き操作によって生成できる関数です。例えば
\[
\frac{\sqrt{1 + x} - \sqrt{1 - x}}{\sqrt{1 + x} + \sqrt{1 - x}},\quad
\sqrt{x} + \sqrt{x + \sqrt{x}},\quad
\left(\frac{x^2 + x + \sqrt{3}}{x \sqrt{2} - \pi}\right)^{\frac{2}{3}}
\]
は明示的代数関数であり、$x^{m/n}$（つまり $\sqrt[n]{x^m}$）も同様です。ここで $m$ と $n$ は任意の整数です。

$y = \sqrt{x}$ のような式において符号の不確定性に注意してください。これまで $\sqrt{2}$ を $2$ の正平方根を表すものとして扱ってきましたが、$x$ が任意の正の数であるとき、$\sqrt{x}$ を正の平方根を表すものと定義すれば、$y = \sqrt{x}$ は $x$ の一対一関数となります。しかししばしば $\sqrt{x}$ を $x$ の正の平方根と負の平方根の二つの値を持つ関数として扱う方が便利です。

読者は、この方法が採用された場合、関数$\sqrt{x}$は有理関数と根本的に異なる点があることに気づくでしょう。まず第一に、有理関数は特定の少数の離散的な例外値を除いてすべての$x$の値に対して定義されています。しかし$\sqrt{x}$は、負のすべての$x$の値に対して定義されていません（つまり、全体的な範囲）。第二に、$x$が定義される値を持つとき、関数は一般に符号反対の2つの値を持ちます。

一方、関数$\sqrt{x}$は一価であり、すべての$x$の値に対して定義されています。

例 XIII.
1. $\sqrt{(x - a)(b - x)}$（$a < b$）は$a \leq x \leq b$の範囲でのみ定義されます。もし$a < x < b$であれば、2つの符号反対の値を持ちます；しかし$x = a$または$b$のときには1つの値、つまり0だけです。

2. 同様に
\begin{gather*}
\sqrt{(x - a)(x - b)(x - c)} \quad (a < b < c), \\
\sqrt{x(x^{2} - a^{2})},\quad
\sqrt{(x - a)^{2}(b - x)}\quad (a < b), \\
\frac{\sqrt{1 + x} - \sqrt{1 - x}}
     {\sqrt{1 + x} + \sqrt{1 - x}},\quad
\sqrt{(x + \sqrt{x})}.
\end{gather*}

3. 曲線$y^{2} = x$, $y^{3} = x$, $y^{2} = x^{3}$を描きなさい。

4. 関数のグラフを描きなさい
\[
y = \sqrt{a^{2} - x^{2}},\quad
y = b\sqrt{1 - (x^{2}/a^{2})}.
\]

27. 隠れ代数的関数。式
\[
y = \frac{\sqrt{1 + x} - \sqrt{1 - x}}{\sqrt{1 + x} + \sqrt{1 - x}},
\]
が満たす条件は
\[
\left(\frac{1 + y}{1 - y}\right)^6 = \frac{(1 + x)^3}{(1 - x)^2};
\]
また、式
\[
y = \sqrt{x} + \sqrt{x + \sqrt{x}},
\]
が満たす条件は
\[
y^4 - (4y^2 + 4y + 1)x = 0.
\]
これらのそれぞれの式は次の形で表すことができます：
\[
y^m + R_1 y^{m-1} + \dots + R_m = 0,
\tag{1}
\]
ここで $R_1$, $R_2$, \dots, $R_m$ は $x$ の有理関数です。そして、読者は容易に確認できるでしょう：最後の例で考慮された任意の関数 $y$ がこの形式の方程式を満たします。

自然に、これは明示的な代数的関数についても真であることが提案されます。実際、これは真であり、確かに証明は難しくありませんが、ここでは形式的な証明を書くことは控えます。例を挙げることで、このような証明の流れが読者に明らかになります。

例えば
\[
y = \frac{x + \sqrt{x} + \sqrt{x + \sqrt{x}} + \sqrt{1 + x}}
     {x - \sqrt{x} + \sqrt{x + \sqrt{x}} - \sqrt{1 + x}},
\]
が満たす条件は
\[
y = \frac{x + u + v + w}{x - u + v - w},
\]
\[
u^2 = x, \quad v^2 = x + u, \quad w^3 = 1 + x,
\]
であり、これらの式から $u$, $v$, $w$ を消去することで、求める形の方程式を得ることができます。

したがって、以下の定義を導入します：\emph{関数 $y = f(x)$ が代数的関数であるとは、その関数が有理関数 $x$ の係数を持つ $m$ 次方程式の根であることを指す。最初の係数を1とするには、明示的な一般性を失うことはありません}。

この関数のクラスには、第26節で考慮されたすべての明示的な代数的関数が含まれます。しかし、それ以外にも、明示的な代数的関数として表現できない他の関数も含まれています。一般に、式（1）のような方程式は、$m$ が 4 より大きい場合、$y$ を $x$ で明示的に解くことは不可能であり、しかし $m = 1, 2, 3,$ または $4$ のとき、そして $m$ のより高い値に対する特定の状況では常に可能であることが知られています。

代数的関数の定義は、最後の章で与えられた代数的数の定義（MiscExs I 32）と比較すべきです。

\par\medskip\noindent\textbf{【例 XIV.】}\par
\Item{1.} $m = 1$ のとき、$y$ は有理関数である。
\Item{2.} $m = 2$ のとき、方程式は $y^{2} + R_{1}y + R_{2} = 0$ であり、
\[
y = \frac{-R_{1} ± \sqrt{R_{1}^{2} - 4R_{2}}}{2}.
\]
この関数は、$R_{1}^{2} \geq 4R_{2}$ のすべての $x$ の値に対して定義される。$R_{1}^{2} > 4R_{2}$ のときには 2 値があり、$R_{1}^{2} = 4R_{2}$ のときには 1 値しかない。
\Item{3.} 方程式
\[
y^{2} - 2y - x^{2} = 0,\quad y^{2} - 2y + x^{2} = 0,\quad y^{4} - 2y^{2} + x^{2} = 0,
\]
によって定義される関数について考え、それぞれ $y$ を $x$ の明示的な関数として求め、その定義が可能となる $x$ の値を述べよ。
\Item{4.} 次の関数について、係数が $x$ で有理である代数方程式を求めよ：
\[
\sqrt{x} + \sqrt{\frac{1}{x}},\quad
\sqrt{x} + \sqrt{\frac{1}{x}},\quad
\sqrt{x + \sqrt{x}},\quad
\sqrt{x + \sqrt{x + \sqrt{x}}}.
\]
\Item{5.} 方程式 $y^{4} = x^{2}$ を考える。
[ここで $y^{2} = ±x$。$x$ が正のとき、$y = \sqrt{x}$；負のとき、$y = \sqrt{-x}$。したがって、$x = 0$ 以外のすべての $x$ の値に対して関数は 2 値を持つ。]

\item 6. $x$ の代数関数の代数関数は、$x$ の代数関数です。

[なぜなら、
\begin{alignat*}{4}
y^{m} \&+ R\\textit{{1}(z)y^{m-1} \&\&+ \dots \&\&+ R}{m}(z) \&\&= 0, \\
z^{n} \&+ S\\textit{{1}(x)z^{n-1} \&\&+ \dots \&\&+ S}{n}(x) \&\&= 0.
\end{alignat*}
を消去して $z$ を取り除くと、
\begin{equation*}
y^{p} + T_{1}(x)y^{p-1} + \dots + T_{p}(x) = 0
\end{equation*}
という形の方程式が得られます。ここですべての大文字は有理関数を表します。]

\item 7. $x$ の代数関数として明示的な代数形式で表現できない例を挙げると、方程式
\[
y^{5} - y - x = 0
\]
によって定義される関数 $y$ があります。しかし、$y$ を $x$ 用の明示的な代数式として求めることができないことを証明する方法は難しく、ここでは試みられません。

\paragraph{28.超越関数。} $x$ のすべての関数が有理または代数的でないものは、超越関数と呼ばれます。この関数のクラスは、純粋に否定的な定義により定められるため、自然とさまざまな種類の関数を無限に含みます。これらの中には現在特に興味深い2つの種類があります。

\topic{\item E. 直接逆三角関数または円関数。}これは初等三角学における正弦および余弦関数、それらの逆関数、およびそれらから派生した関数です。読者はこれらの重要な性質について既に熟悉であると仮定します。\footnote{初等三角学で与えられる円関数の定義は、任意の扇形が固有の数（その\emph{面積}）を持つことを前提としています。この仮定が正当化される方法については第VII章で明らかになります。}
ページ分割：53

\par\medskip\noindent\textbf{【例】}\par{XV.}
\item 1. $\cos x$, $\sin x$, および $a\cos x + b\sin x$ のグラフを描きなさい。
\par\medskip

$a\cos x + b\sin x = \beta\cos(x - \alpha)$ とし、$\beta = \sqrt{a^{2} + b^{2}}$、および $\alpha$ は余弦と正弦が $a/\sqrt{a^{2} + b^{2}}$ と $b/\sqrt{a^{2} + b^{2}}$ の角とするとき、これらの3つの関数のグラフは性質において類似している。

2. $\cos^{2} x$, $\sin^{2} x$, $a\cos^{2} x + b\sin^{2} x$ のグラフを描け。

3. $f(x)$ と $F(x)$ のグラフが描かれているとする。次のようなグラフ
\[ f(x)\cos^{2} x + F(x)\sin^{2} x \]
は、$y = f(x)$, $y = F(x)$ の間で振動する波状の曲線である。$f(x) = x$, $F(x) = x^{2}$ の場合のグラフを描け。

4. $\cos px + \cos qx$ のグラフは、$2\cos\frac{1}{2}(p - q)x$ と $-2\cos\frac{1}{2}(p + q)x$ の間で振動し、それぞれに接する。$(p - q)/(p + q)$ が小さい場合のグラフを素描せよ。（数学入門試験 1908年）

5. $x + \sin x$, $(1/x) + \sin x$, $x\sin x$, $(\sin x)/x$ のグラフを描け。

6. $\sin(1/x)$ のグラフを描け。
もし $y = \sin(1/x)$ なら、$x = 1/m\pi$（$m$は任意の整数）で $y = 0$ となる。同様に、$x = 1/(2m + \frac{1}{2})\pi$ で $y = 1$、$x = 1/(2m - \frac{1}{2})\pi$ で $y = -1$ となる。この曲線は $y = -1$ と $y = 1$ の間にある（図13参照）。振動は上下来し、$x$ が $0$ に近づくにつれて振動の頻度が増大する。$x = 0$ では関数は定義されていない。$x$ が大きいとき $y$ は小さい。（この文句の正確な意味については第IV章と第V章を参照のこと）負の半分の曲線は正の半分と類似している。

7. $x\sin(1/x)$ のグラフを描け。
この曲線は、最後の曲線が $y = -1$ と $y = 1$ の間にあるように（図14参照）、$y = -x$ と $y = x$ の間にある。

\item 8. \(x^{2}\sin(1/x)\), \((1/x)\sin(1/x)\)，\(\sin^{2}(1/x)\)，\(\{x\sin(1/x)\}^{2}\)，\(a\cos^{2}(1/x) + b\sin^{2}(1/x)\)，\(\sin x + \sin(1/x)\)，\(\sin x\sin(1/x)\)のグラフを描きなさい。

\item 9. \(\cos x^{2}\)，\(\sin x^{2}\)，\(a\cos x^{2} + b\sin x^{2}\)のグラフを描きなさい。

\item 10. \(\arccos x\)と\(\arcsin x\)のグラフを描きなさい。

[もし \(y = \arccos x\) なら，\(x = \cos y\) である。これにより，\(x\) を \(y\) の関数としてグラフを描くことができ，同じ曲線は \(y\) を \(x\) の関数としても表すことができる。\(y\) は \(-1 \leq x \leq 1\) の範囲で定義され，これらの値に対する \(x\) に対して無限に多価であることが明らかである。読者にはおそらく覚えているが，\(-1 < x < 1\) のとき，\(0\) よりも大きい \(\pi\) よりも小さい値の \(y\) が存在し，その値を \(\alpha\) とすると，他の値は公式 \(2n\pi ± \alpha\)（\(n\) は任意の整数）によって与えられる。]

\item 11. \(\tan x\)，\(\cot x\)，\(\sec x\)，\(\cosec x\)，\(\tan^{2} x\)，\(\cot^{2} x\)，\(\sec^{2} x\)，\(\cosec^{2} x\) のグラフを描きなさい。

\item 12. \(\arctan x\)，\(\arccot x\)，\(\arcsec x\)，\(\arccosec x\)のグラフを描きなさい。それぞれの関数のすべての値を特定の値で表す公式（例題 10 のように）を与えなさい。

\item 13. \(\tan(1/x)\)，\(\cot(1/x)\)，\(\sec(1/x)\)，\(\cosec(1/x)\)のグラフを描きなさい。

\item 14. \(\cos x\)と\(\sin x\)は \(x\) の有理関数ではないことを示しなさい。

[周期的な関数は、周期 \(a\) について \(f(x) = f(x + a)\) を満たす関数と定義される。したがって \(\cos x\) および \(\sin x\) の周期は \(2\pi\) である。周期的な関数が有理関数であることは、常に定数以外にはならないことが容易に理解できる。仮に
\[ f(x) = \frac{P(x)}{Q(x)}, \]
ここで \(P\) および \(Q\) は多項式であり、かつ \(f(x) = f(x + a)\) がすべての \(x\) の値に対して成り立つとする。\(f(0) = k\) とすると、方程式 \(P(x) - kQ(x) = 0\) は無数の \(x\) の値で満足されることがわかる。したがって、すべての \(x\) の値に対して成り立つ。従って \(f(x) = k\) であり、すなわち \(f(x)\) は定数である.]

15. もっと一般に、周期を持つ関数が \(x\) の代数的関数であることは示せ。

[代数的関数を定義する方程式を
\[ y^{m} + R_{1}y^{m-1} + \dots + R_{m} = 0 \]
とし、ここで \(R_{1}\), \(\dots\) は \(x\) の有理関数であるとする。この式は次の形に書き換えることができる：
\[ P_{0}y^{m} + P_{1}y^{m-1} + \dots + P_{m} = 0, \]
ここで \(P_{0}\), \(P_{1}\), \(\dots\) は \(x\) の多項式である。上記の議論を用いて、
\[ P_{0}k^{m} + P_{1}k^{m-1} + \dots + P_{m} = 0 \]
がすべての \(x\) の値に対して成り立つことがわかる。したがって、方程式 (1) を満たす \(y = k\) はすべての \(x\) の値に対して成り立つ。つまり、我々の代数的関数の一つの値の集合は定数に縮約される。

次に、方程式 (1) を \(y - k\) で割って再び議論を繰り返す。最終的な結論は、任意の \(x\) の値に対して、我々の代数的関数が一定の値の集合 \(k\), \(k'\), \(\dots\) から成るということである。すなわち、定数で構成されていることである.]

16. 逆正弦および逆余弦は有理関数または代数的関数ではない。[これは、任意の \(-1 < x < +1\) の値に対して \(\arcsin x\) および \(\arccos x\) が無数の値を持つことから従う。]

29. その他の超越関数のクラス。三角関数に次ぐ重要性を持つのが指数関数と対数関数で、これらは第 IX 章と第 X 章で議論される予定です。しかし、これらの関数は現在の範囲外です。また、楕円関数、ベッセル関数、 Legendre 関数、ガンマ関数など、その他の研究された超越関数の多くも本書の範囲外です。ただし、理論的には理数や代数的、または三角関数よりも重要な要素ではないが、関数関係の可能性的な多様性を示す例として特に有用ないくつかの基本的な関数があります。

例 XVI.
1. $y = [x]$ で、$[x]$ は $x$ よりも大きい整数の中で最大のものを表します。グラフは図 15a に示されています。厚線の左端点がグラフに含まれますが、右端点は含まれていません。

2. $y = x - [x]$。（図 15b）

3. $y = \sqrt{x - [x]}$。（図 15c）

4. $y = [x] + \sqrt{x - [x]}$。（図 15d）

5. $y = (x - [x])^2, [x] + (x - [x])^2$

6. $y = [\sqrt{x}], [x^2], \sqrt{x} - [\sqrt{x}], x^2 - [x^2], [1 - x^2]$。（図 15c、図 15d）

7. $y$ を $x$ の最大の素因数として定義します（例題 x.6 参照）。このとき $y$ は整数値の $x$ に対してのみ定義されます。例えば
\[
ZZZMATHTOKEN3ZZZ
\]
とします。グラフはいくつかの孤立した点から成ります。

\item{8.} $y$ を $x$ の分母とする（例7参照）。この場合、$y$ は $x$ の有理数値に対してのみ定義される。私たちはグラフ上に任意の数の点を描くことができるが、その結果は通常の意味での曲線ではなく、$x$ の無理性値に対応する点は存在しない。

$(N - 1, N)$, $(N, N)$ を結ぶ直線を引く。ここで $N$ は正の整数とする。この直線上にある locus の点の数が $N$ より小さい正の素数の数と等しいことを示せ。

\item{9.} $x$ が整数のとき $y = 0$、$x$ が整数でないとき $y = x$。この関数は直線 $y = x$ の点を $(\dots, -1, -1), (0, 0), (1, 1), (2, 2), \dots$ から取り除き、$x$ 軸上に $(-1, 0)$, $(0, 0)$, $(1, 0)$, $\dots$ の点を追加することで得られる。

読者はこの関数が不合理であると考えるかもしれません。なぜなら、彼は疑問を呈し得るでしょう：もしすべての $x$ 値（整数でない場合）に対して $y = x$ であれば、整数値の場合も同様に $y = x$ ではないでしょうか？答えは単純である：なぜそれがそうならないべきなのか？関数 $y$ は実際には関数の定義に従っている：$x$ の値が分かれば $y$ の値も確定する。私たちはこの関係を任意に選択し、いかなる制限や無意味さにも拘らずそれを採用することができる。この関数 $y$ は当然ながら、常に $y = x$ である関数とは全く異なる関数である。

\item{10.} $x$ が有理数のとき $y = 1$、$x$ が無理性数のとき $y = 0$。このグラフは $y = 1$ と $y = 0$ の直線上に点の2つの系列を配置する。目で見れば連続した直線のように見えますが、実際には各直線から無限個の点が欠けている。

Item 11. \( y = x \) であるとき \( x \) が無理数であり、\( y = \sqrt{\frac{1 + p^2}{1 + q^2}} \) であるとき \( x \) が有理分数 \( \frac{p}{q} \) である。

図表：Fig. 16。
**注釈：x = p/q の場合に描かれた式は \(\sqrt{\frac{10 + p^2}{10 + q^2}}\) であり、これは視覚的に書籍と一致するように選ばれています。詳細については末尾の抄録者注をご覧ください。**

無理数 \( x \) の値は実際には不連続な曲線をグラフに寄与しますが、直線 \( y = x \) から区別するのが困難です。

次に有理数 \( x \) の値を考えます。まず \( x \) が正の場合について考えましょう。すると \(\sqrt{\frac{1 + p^2}{1 + q^2}}\) は \( \frac{p}{q} \) と等しいことは、\( p = q \) のときのみであり、すなわち \( x = 1 \) のときです。したがって、有理数 \( x \) の値に対応するすべての点は直線から離れており、唯一の例外として点 (1, 1) があります。また、\( p < q \) の場合、\(\sqrt{\frac{1 + p^2}{1 + q^2}} > \frac{p}{q}\); \( p > q \) の場合、\(\sqrt{\frac{1 + p^2}{1 + q^2}} < \frac{p}{q}\) です。したがって、\( 0 < x < 1 \) のとき点は直線 \( y = x \) の上にあり、\( x > 1 \) のとき点は直線の下にあります。もし \( p \) と \( q \) が大きい場合、\(\sqrt{\frac{1 + p^2}{1 + q^2}}\) は \( \frac{p}{q} \) にほぼ等しくなります。任意の \( x \) の値近傍で十分大きな分子と分母を持つ有理分数を何個でも見つけることができます。したがって、グラフには直線 \( y = x \) の周りに多くの点が密集しています。正の \( x \) の値に対するグラフ全体の外観は、直線周囲に孤立した点の群れがあり、これらの点は直線に近づくにつれて密度が増すように見えます。

\( x \) の負の値に対応するグラフの部分は不連続な直線の残りと \( y \) 軸に対するすべての孤立した点の反射で構成されます。したがって、\( y \) 軸の左側では点の群れは \( y = x \) の周りではなく \( y = -x \) の周りにあり、これはグラフの一部ではありません。図 16 を参照してください。

ページ番号：58

30. 未知数が一つの方程式のグラフ的解法。多くの方程式は次の形で表すことができる：
\[
f(x) = \phi(x),
\tag{1}
\]
ここで $f(x)$ と $\phi(x)$ はグラフを容易に描くことができる関数である。そして、曲線
\[
y = f(x),\quad y = \phi(x)
\]
が点 $P$ で交わる場合、その点の横座標 $\xi$ は方程式 \Eq{(1)} の根となる。

例 XVII.
1. 二次方程式 $ax^{2} + 2bx + c = 0$. この方程式は様々な方法でグラフ的に解くことができる。例えば，
\[
y = ax + 2b,\quad y = -c/x
\]
の交点（存在する場合）が根となる。または，
\[
y = x^{2},\quad y = -(2bx + c)/a
\]
を取ることもできる。しかし、最も基本的な方法は次の円を描くことである：
\[
a(x^{2} + y^{2}) + 2bx + c = 0,
\]
その中心は $(-b/a, 0)$ であり、半径は $\{\sqrt{b^{2} - ac}\}/a$ である。この円と $x$ 軸の交点の横座標が方程式の根となる。

2. 次の方法のいずれかで解け：
\[
x^{2} + 2x - 3 = 0,\quad x^{2} - 7x + 4 = 0,\quad 3x^{2} + 2x - 2 = 0.
\]

3. 方程式 $x^{m} + ax + b = 0$. この方程式は曲線 $y = x^{m}$, $y = -ax - b$ を描くことで解ける。次に、次の表を検証せよ：
\[
\begin{alignedat}{3}
&(\ia) &&m~\text{偶数} &&\left\{
  \begin{aligned}
    &\text{$b$ 正のとき, \emph{二つまたはなし},}\\
    &\text{$b$ 負のとき, \emph{二つ};}
  \end{aligned}
\right. \\
&(\ib) &&m~\text{奇数} &&\left\{
  \begin{aligned}
    &\text{$a$ 正のとき, \emph{一つ},}\\
    &\text{$a$ 負のとき, \emph{三つまたは一つ}.}
  \end{aligned}
\right.
\end{alignedat}
\]
すべての場合を例として数値的に示す。

4. 方程式 $\tan x = ax + b$ は無限に多くの根を持つことを示せ。

\item 5. 次の方程式の解の数を求めよ：
\[
\sin x = x, \quad \sin x = \frac{1}{3}x, \quad \sin x = \frac{1}{8}x, \quad \sin x = \frac{1}{120}x.
\]

\item 6. $a$が小さな正の数であるとき（例えば$a=0.01$）、方程式
\[
x - a = \frac{1}{2}\pi \sin^2 x
\]
は3つの解を持つことを示せ。また、$a$が小さな負の数である場合も考慮し、$a$が変化するときの解の数の変化を説明せよ。

\paragraph{31. 二変数関数とそのグラフィカルな表現} 第20節（§20）では、二つの変数$x$と$y$が関係によって結びついていることを考慮しました。同様に、三つの変数（$x$, $y$, および $z$）が関係によって結びついている場合を考えることができます。この関係において、$x$と$y$の値が与えられたとき、$z$の値または値が既知である場合、$z$を$x$と$y$の二つの独立変数に対する\emph{依存関数}と呼びます。この依存関係は
\[
z = f(x, y)
\]
で表します。

第20節（§20）での議論は、このより複雑なケースに対して、適宜修正した形で全て適用できます。

二つの変数に対するこのような関数をグラフィカルに表現する方法は、一つの変数に対する関数の場合と原則的には同じです。空間の3次元において$OX$, $OY$, $OZ$という3軸を取ります。各軸が他の2軸に対して垂直であることを確認します。点$(a, b, c)$は、$YOZ$, $ZOX$, $XOY$という平面からそれぞれ$OX$, $OY$, $OZ$方向に平行に測った距離$a$, $b$, $c$を持つ点です。当然、符号を考慮する必要があります。$OX$, $OY$, $OZ$方向の長さは正とします。座標系、軸、原点の定義は以前と同じです。

今、
\[
z = f(x, y)
\]
とします。$x$ および $y$ の値が変化するとき、点 $(x, y, z)$ は空間上で動きます。この点が取るすべての位置の集合を、点 $(x, y, z)$ の \emph{軌跡} または関数 $z = f(x, y)$ の \emph{グラフ} と呼びます。$x$, $y$, $z$ の間の定義される関係が解析式で表せる場合、その式をこの点の \emph{方程式} と呼びます。例えば、方程式
\[
Ax + By + Cz + D = 0
\]
（\emph{一次方程式の一般形}）は平面を表し、任意の平面の方程式はこの形式です。また、方程式
\[
(x - \alpha)^2 + (y - \beta)^2 + (z - \gamma)^2 = \rho^2,
\]
または
\[
x^2 + y^2 + z^2 + 2Fx + 2Gy + 2Hz + C = 0,
\]
（$F^2 + G^2 + H^2 - C > 0$）は球面を表します。これらの命題の証明については、解析幾何の教科書を参照してください。

\Paragraph{32. 平面上の曲線} これまでに、関数依存性を表現するために
\[
y = f(x)
\Tag{(1)}
\]
という記号を使用してきました。これは $y$ が $x$ によって明示的に式で表される場合（例えば
\[
y = x^2, \quad \sin x, \quad a \cos^2 x + b \sin^2 x
\]）において最も適切な記法であることが明らかです。

しかし、しばしば $y$ が $x$ によって明示的に式で表すことができない関数的な関係に直面します。例えば、$y^5 - y - x = 0$ や $x^5 + y^5 - ay = 0$ のような方程式では、$y$ を $x$ の代数的関数として明示的に表すことが不可能または不便です。一方で
\[
x^2 + y^2 + 2Gx + 2Fy + C = 0
\]
の場合は、$y$ を以下のように明示的に表すことができます：
\[
y = -F + \sqrt{F^2 - x^2 - 2Gx - C};
\]
しかし、$y$ が $x$ に及ぼす関数的な依存性は、元の方程式によってより適切かつ簡潔に表現されます。

これらの二つの場合において、関数の関係は二変数$x$と$y$の関数$f(x, y)$を零に等しいという式で完全に表現されることが観察できる。すなわち、方程式
\[
f(x, y) = 0.
\Tag{(2)}
\]
によって表現される。

この方程式を関数の関係を表す標準的な方法として採用する。これは特別な形の$x$と$y$の関数である$y - f(x)$を含む方程式\Eq{(1)}の特殊な場合を包括している。したがって、点$(x, y)$の locus は$f(x, y) = 0$に制限されるもの、関数$y$によって定義されるグラフ、方程式またはlocus $f(x, y) = 0$、そしてこの曲線またはlocus の方程式について話すことができる。

また、曲線を表現する別の方法がしばしば有用である。$x$と$y$が第三の変数$t$の関数であり、$t$は主に補助的なもので、特に幾何学的な意味を持つものとして扱われるという仮定を考える。このとき、
\[
x = f(t),\quad y = F(t).
\Tag{(3)}
\]
と書くことができる。ある特定の値が$t$に割り当てられた場合、それに応じて$x$および$y$の値は分かることになる。これらの値のそれぞれのペアは点$(x, y)$を定める。異なる値の$t$に対応するすべての点をこのように構築すると、方程式\Eq{(3)}で定義されるlocus のグラフを得る。例えば
\[
x = a\cos t,\quad y = a\sin t.
\]
と仮定し、$t$が$0$から$2\pi$まで変化する場合、点$(x, y)$は中心が原点で半径が$a$の円を描くことが容易に確認できる。$t$がこれらの限界を超える場合、$(x, y)$は同じ円を繰り返し描く。この場合、$x$と$y$との直接的な関係を得るために$x$と$y$の2乗を加算するだけで$t$を消去できる。結果として得られる方程式は
\[
x^{2} + y^{2} = a^{2}.
\]

\par\medskip\noindent\textbf{【例 XVIII.】}\par
\Item{1.} 方程式が $f(x, y) = 0$, $\phi(x, y) = 0$ である2つの曲線の交点は、これらの方程式を $x$ と $y$ の連立方程式として解ける場合に求めることができます。一般には、この解は有限個の $x$ と $y$ の値のペアからなります。したがって、2つの方程式は一般的に有限個の孤立点を表します。

\Item{2.} 曲線 $(x + y)^2 = 1$, $xy = 1$, $x^2 - y^2 = 1$ を描写してください。

\Item{3.} 方程式 $f(x, y) + \lambda\phi(x, y) = 0$ は、$f = 0$ と $\phi = 0$ の交点を通る曲線を表します。

\Item{4.} 次の式が $t$ のすべての実数値を取るときの軌跡は何ですか？
\[
\Item{(\alpha)}\ x = at + b,\quad y = ct + d,\qquad
\Item{(\beta)}\ \frac{x}{a} = \frac{2t}{1 + t^2},\quad \frac{y}{a} = \frac{1 - t^2}{1 + t^2}.
\]
\]

\Paragraph{33. 三次元空間の軌跡。} 三次元空間において、基本的には2つの異なる種類の軌跡があります。最も単純な例は平面と直線です。

直線上を動く粒子には \emph{一つの自由度} しかありません。その運動方向が固定されているため、その位置は直线上の一点の距離という1つの測定値で完全に決定されます。例えば、直線を基本線 $\Lambda$ とし、その上の任意の点の位置は単一の座標 $x$ で決定されます。一方、平面内で動く粒子には \emph{二つの自由度} があります；その位置は2つの座標の決定によって完全に固定されます。

一つの方程式
\[
z = f(x, y)
\]
によって表される軌跡は、これらの2つのクラスのうちの第2のクラスに属し、平面曲線と呼ばれます。これは一般的には私たちの直感的な表面の概念を満たすかしないか（単純な明確な場合では）です。

第31節の考慮事項は、明らかに一般化され、三つの変数（または任意の数目的な変数を持つ関数）の関数$f(x, y, z)$に対する定義を提供することができます。そして、第32節で平面曲線の方程式の標準形として$f(x, y) = 0$を選択したように、現在では表面の方程式の標準形として
\[
f(x, y, z) = 0
\]
を選択することに同意します。

二つの形式 $z = f(x, y)$ や $f(x, y, z) = 0$ の方程式によって表される位置は最初の類別に属し、曲線と呼ばれます。したがって、直線は二つのタイプの方程式 $Ax + By + Cz + D = 0$ で表すことができます。空間における円は球面と平面の交差点として考えられるため、以下の形式の二つの方程式
\[
(x - \alpha)^{2} + (y - \beta)^{2} + (z - \gamma)^{2} = \rho^{2},\quad
Ax + By + Cz + D = 0.
\]
で表すことができます。

例題XIX。
1. $f(x, y, z) = 0$ のタイプの三つの方程式は何を表しますか？
2. 一般に、三次元空間における一点を表すために二つの一次方程式が必要です。例外的なケースはどのようなものですか？
3. 平面 $XOY$ 上の平面曲線 $f(x, y) = 0$ が空間上の曲線としての方程式は何ですか？ [$f(x, y) = 0$, $z = 0$.]
4. 柱体。三次元空間における一つの方程式 $f(x, y) = 0$ の意味は何か？
全ての点が表面に存在し、どの値の $z$ も $f(x, y) = 0$ を満たす場合があります。平面 $XOY$ での曲線 $f(x, y) = 0$, $z = 0$ はその位置が平面 $XOY$ で交差する曲線です。この曲線を通る全ての点に平行な $OZ$ 方向の直線を引いた表面が形成されます。このような表面は「柱体」と呼ばれます。

\Item{5.} \Topic{平面での曲面のグラフィカル表現。等高線図。}
表面を適切に平面の図で表すことが不可能であるように思われるかもしれませんが、実際にはそうではありません：しかし、表面の性質について十分な概念を得ることはしばしば可能です。表面の方程式が \(z = f(x, y)\) であるとします。

\(z\) に特定の値 \(a\) を与えると、方程式 \(f(x, y) = a\) があり、これは紙上で平面曲線を定めるものとして扱うことができます。この曲線を描き、\((a)\) とマークします。実際には、曲線 \((a)\) は平面 \(XOY\) 上の表面の断面によって投影されたものです。この作業はすべての \(a\) の値に対して（もちろん、実際的には \(a\) の選択的な値に対して）行います。このような図形が \Fig{17} に示されます。これは直ちに等高線地図を想起させます：実際には、このような地図はこの原理に基づいて作成されています。等高線 \((1000)\) は海面の平面上の投影であり、海面と平行で \(1000\) フィート上にある平面によって表面の断面によって得られます。\footnote{地球曲率の影響を無視することができます。}
\%[図版: Fig. 17.]
\Figure{17}{p063}

\Item{6.} {\Loosen 等高線の系列を描いて、表面 \(2z = 3xy\) の形状を示すことを説明してください。}

\Item{7.} \Topic{直円錐体。} 座標系の原点を円錐の頂点に取り、\(z\) 軸を円錐の軸とし、\(\alpha\) を円錐の半垂直角とする。円錐の方程式（その頂点から両方に延びる）は \(x^{2} + y^{2} - z^{2}\tan^{2} \alpha = 0\) である。

\item{8.} 一般の回転面について。Ex.7 の円錐は$ZOX$と2つの直線を交わり、その方程式は$x^{2} = z^{2}\tan^{2}\alpha$という形で組み合わされる。
つまり、曲線$y = 0$, $x^{2} = z^{2}\tan^{2}\alpha$を$z$軸周りに回転して生成された表面の方程式は、これらの2つの方程式のうちの後者から$x^{2}$を$x^{2} + y^{2}$に置き換えることで導き出される。一般に、曲線$y = 0$, $x = f(z)$を$z$軸周りに回転して生成された表面の方程式は
\[
\sqrtp{x^{2} + y^{2}} = f(z)
\]
であることを示せ。

\item{9.} 一般の円錐について。一直線を通る固定点を通る曲面を円錐と定義し、その固定点を頂点という。Ex.7 の直角円錐はその特別な例である。頂点が$O$である円錐の方程式は$f(z/x, z/y) = 0$の形であり、この形の任意の方程式は円錐を表すことを示せ。[$(x, y, z)$が円錐上にある場合、$\lambda x$, $\lambda y$, $\lambda z$もまた円錐上にある任意の値$\lambda$があるためである。]

\item{10.} 一般の直線生成曲面について。円柱と円錐は直線で構成される曲面の特別な例である。このような曲面を直線生成曲面という。

方程式
\[
x = az + b,\quad
y = cz + d,
\Tag{(1)}
\]
は2つの平面の交わりを表し、すなわち一直線を表す。ここで$a$, $b$, $c$, $d$が固定ではなく、補助変数$t$の関数であるとする。任意の特定の値$t$に対して方程式 \Eq{(1)} は一直線を与える。$t$ の値が変化するにつれてこの直線は動くと同時に曲面を生成し、その方程式は2つの方程式 \Eq{(1)} から$t$ を消去することで求めることができる。例えば、Ex.7 では円錐を生成する一直線の方程式は
\[
x = z\tan \alpha\cos t,\quad
y = z\tan \alpha\sin t,
\]
であり、ここで$t$は平面$XOZ$と一直線を通る平面および$z$軸との角度である。

別の単純な例として、次のように制御曲面を構築することができます。
円柱の軸に垂直で距離が~$l$ の二つの断面を考えます（図18a参照）。私たちは円柱の表面を、長さが~$l$ の薄い平行な剛体棒から成るものと考えることができます。例えば棒~$PQ$ があり、その端点は半径が~$a$ の二つの円棒に固定されています。

次に、同じ半径の第三の円棒を取り、最初の二つの棒のいずれかから距離が~$h$ の位置で円柱の表面に配置します（図18a参照、ここで $Pq = h$）。棒~$PQ$ の端点~$Q$ を解錨し、$P$ を中心に棒~$PQ$ を回転させて、$Q$ を第三の円棒に位置~$Q'$ で固定します。図において角度 $qOQ' = \alpha$ は次のように与えられます：
\[
l^{2} - h^{2} = qQ'^{2} = \left (2a\sin\tfrac{1}{2} \alpha\right)^{2}.
\]
円柱を構成する他のすべての棒についても同様に処理します。これにより、図18bで示されるような形状の制御曲面が得られます。この曲面は完全に直線から成り立っていますが、どこでも曲率を持ち、一般的な形としては特定のテーブルナプキンリング（図18c参照）のようなものと似ています。
\begin{figure}[hbt!]
  \begin{minipage}{0.3\textwidth}
    \centering
    \fbox{【原著図版参照】}
    \caption{図18a。}
    \label{fig:18a}
  \end{minipage}\hfill
  \begin{minipage}{0.3\textwidth}
    \centering
    \fbox{【原著図版参照】}
    \caption{図18b。}
    \label{fig:18b}
  \end{minipage}\hfill
  \begin{minipage}{0.3\textwidth}
    \centering
    \fbox{【原著図版参照】}
    \caption{図18c。}
    \label{fig:18c}
  \end{minipage}
\end{figure}

\par\medskip\noindent\textbf{【例】}\par{}
\Section{第II章の補足例題}

\begin{Examples}[]
\Item{1.} $y = f(x) = (ax + b)/(cx - a)$ のとき、$x = f(y)$ を示せ。
\par\medskip
\newpage

\Item{2.} $f(x) = f(-x)$ がすべての $x$ の値に対して成り立つとき、$f(x)$ を偶関数と呼びます。$f(x) = -f(-x)$ であるとき、奇関数と呼びます。任意の $x$ の定義域を持つ関数は、偶関数と奇関数の和として表されることが示されます。

[使用恒等式：$f(x) = \frac{1}{2}\{f(x) + f(-x)\} + \frac{1}{2}\{f(x) - f(-x)\}$]

\Item{3.} 関数のグラフを描きます：
\[ 3\sin x + 4\cos x, \quad \sin\left(\frac{\pi}{\sqrt{2}} \sin x\right). \]
[1896年数学問題]

\Item{4.} 関数のグラフを描きます：
\[ \sin x(a\cos^{2} x + b\sin^{2} x), \quad \frac{\sin x}{x}(a\cos^{2} x + b\sin^{2} x), \quad \left(\frac{\sin x}{x}\right)^{2}. \]

\Item{5.} 関数のグラフを描きます：$x[1/x]$, $[x]/x$。

\Item{6.} 関数のグラフを描きます：
\begin{align*}
\Itemp{(i)} & \arccos(2x^{2} - 1) - 2 \arccos{x}, \\
\Itemp{(ii)} & \arctan \frac{a + x}{1 - ax} - \arctan{a} - \arctan{x},
\end{align*}
ここで、$\arccos a$, $\arctan a$ は任意の $a$ の値に対して、その余弦または正接が $a$ となる最小の正な角（または零）を表します。

\Item{7.} $f\{\phi(x)\}$ のグラフを $y = x$, $f(x)$, $\phi(x)$ のグラフを使用して構築する方法を確認します：$OA = x$ を $OX$ 上に引きます。$AB$ を $OY$ に平行にして $\phi(x)$ と交わる点 $B$ を引きます。$BC$ を $OX$ に平行にして $y = x$ と交わる点 $C$ を引きます。$CD$ を $OY$ に平行にして $f(x)$ と交わる点 $D$ を引きます。$DP$ を $OX$ に平行にして $AB$ と交わる点 $P$ を引きます。すると、$P$ は求められるグラフ上の点となります。

\Item{8.} $x^{3} + px + q = 0$ の根は、放物線 $y = x^{2}$ と円
\[ x^{2} + y^{2} + (p - 1)y + qx = 0 \]
の交点（原点を除く）の横座標に等しいことが示されます。

\Item{9.} $x^{4} + nx^{3} + px^{2} + qx + r = 0$ の根は、放物線 $x^{2} = y - \frac{1}{2}nx$ と円
\[ x^{2} + y^{2}
  + (\tfrac{1}{8}n^{2} - \tfrac{1}{2}pn + \tfrac{1}{2}n + q)x
  + (p - 1 - \tfrac{1}{4}n^{2})y + r = 0.
\]
の交点の横座標に等しいことが示されます。

\item 10. 方程式
\[
x^{m} + ax^{2} + bx + c = 0
\]
を、曲線 \( y = x^{m} \), \( y = -ax^{2} - bx - c \) を用いてグラフ的に解く方法について議論せよ。異なる可能性のある根の数の表を作成せよ。

\item 11. 方程式 \(\sec\theta + \cosec\theta = 2\sqrt{2}\) の解を求める。また、方程式 \(\sec\theta + \cosec\theta = c\) が \(0 < \theta < 2\pi\) で２つの根を持つ条件は \(c^{2} < 8\) であり、４つの根を持つ条件は \(c^{2} > 8\) であることを示せ。

\item 12. 次の方程式
\[
2x = (2n + 1)\pi(1 - \cos x),
\]
で \(n\) が正の整数のとき、この方程式は \(2n + 3\) 個の根を持ち、それ以上ないことを示し、その根の位置を概略的に示せ。[1896年]

\item 13. 次の方程式
\[
\frac{2}{3}x\sin x = 1
\]
が \(-\pi < x < \pi\) の範囲で４つの根を持つことを示せ。

\item 14. 次の方程式の根の数と値について議論せよ：

(1) \(\cot x + x - \frac{3}{2}\pi = 0\),

(2) \(x^{2} + \sin^{2} x = 1\),

(3) \(\tan x = \frac{2x}{1 + x^{2}}\),

(4) \(\sin x - x + \frac{1}{6}x^{3} = 0\),

(5) \((1 - \cos x)\tan\alpha - x + \sin x = 0\).

\item 15. \(x = a, b, c\) のとき値が \(\alpha, \beta, \gamma\) となる二次多項式は
\[
\alpha\frac{(x - b)(x - c)}{(a - b)(a - c)} +
\beta \frac{(x - c)(x - a)}{(b - c)(b - a)} +
\gamma\frac{(x - a)(x - b)}{(c - a)(c - b)}
\]
である。\(x = a_{1}, a_{2}, \dots, a_{n}\) のとき値が \(\alpha_{1}, \alpha_{2}, \dots, \alpha_{n}\) となる次数が \(n-1\) の多項式の類似な公式を示せ。

\item 16. \(x = 0, 1, 2\) のとき値が \(1/c, 1/(c + 1), 1/(c + 2)\) となる二次多項式を求めよ。また、\(x = c + 2\) のときその値は \(1/(c + 1)\) であることを示せ。[1911年]

\item 17. \(x\) が \(y\) の有理関数であり、かつ \(y\) が \(x\) の有理関数のとき、\(Axy + Bx + Cy + D = 0\) であることを示せ。

\item 18. \(y\) が \(x\) の代数的関数のとき、\(x\) は \(y\) の代数的関数であることを示せ。

\item{19.} 方程式
\[
\cos\tfrac{1}{2}\pi x
  = 1 - \frac{x^{2}}{x + (x - 1)\bigsqrtp{\dfrac{2 - x}{3}}}
\]
が$0$から$1$までのすべての$x$の値に対して大まかに成り立つことを確認せよ。[各$x = 0$, $\frac{1}{6}$, $\frac{1}{3}$, $\tfrac{1}{2}$, $\frac{2}{3}$, $\frac{5}{6}$, $1$に対して表を使用して計算し、どの値で式が正確かを調べよ。]

\item{20.} 関数
\[
z = [x] + [y],\quad
z = x + y - [x] - [y]
\]
のグラフの形状はどのようなものか。

\item{21.} 関数$z = \sin x + \sin y$, $z = \sin x\sin y$, $z = \sin xy$, $z = \sin(x^{2} + y^{2})$のグラフの形状はどのようなものか。

\item{22.} 无理数的几何作图。第I章で、与えられた単位長から等しい長さ$\sqrt{2}$を得る簡単な幾何学的な作図法を一つまたは二つ示した。また、任意の二次方程式$ax^{2} + 2bx + c = 0$の根を構築する方法も示した。これは係数$a$, $b$, $c$の比が与えられた場合、つまり$a$, $b$, $c$が有理数である場合に、任意の長さを構築できることが保証されていると仮定した。これらのすべての作図法はユークリッド的なものであり、直尺と Compass だけを使用したものだった。

これらによって、どのような複雑な平方根の組み合わせでも定義される任意の無理数の長さを構築できることが明白である。例えば
\[
\bigsqrtb{\bigsqrtp{\frac{17 + 3\sqrt{11}}{17 - 3\sqrt{11}}}
           - \bigsqrtp{\frac{17 - 3\sqrt{11}}{17 + 3\sqrt{11}}}}
\]
はその一例である。この式には四次根が含まれているが、これは当然平方根の平方根である。まず$\sqrt{11}$を$1$と$11$の中間として構築し、それから$17 + 3\sqrt{11}$と$17 - 3\sqrt{11}$を構築する。あるいは、これらの複合根号は$x^{2} - 34x + 190 = 0$の解として直接構築される可能性もある。

逆に、このような無理数はユークリッドの方法でしか構築できない。単位長から始めると、任意の有理長を構築できる。したがって、$A$, $B$, $C$ の比が有理数であるとき、直線 $Ax + By + C = 0$ を構築できるし、$\alpha$, $\beta$, $\rho$ が有理数であるとき、円
\[
(x - \alpha)^{2} + (y - \beta)^{2} = \rho^{2}
\]
（または $x^{2} + y^{2} + 2gx + 2fy + c = 0$）も構築できる。この条件は、$g$, $f$, $c$ も有理数であることを意味する。

任意のユークリッドの構造において、図形に導入される新しい点は、二つの直線または円の交点、または直線と円で定められる。しかし、係数が有理数である場合、このような方程式の組
\[
Ax + By + C = 0,\quad x^{2} + y^{2} + 2gx + 2fy + c = 0
\]
を解くと、$x$ および $y$ の値は $m + n\sqrt{p}$ の形となる（$m$, $n$, $p$ が有理数）。なぜなら、$x$ を $y$ で表して第2の方程式に代入すると、$y$ の有理係数の二次方程式を得るからである。したがって、有理係数の直線や円を使って得られるすべての点の座標は、有理数と平方根で表現できる。そして、任意の二点間の距離 $\sqrt{(x_{1} - x_{2})^{2} + (y_{1} - y_{2})^{2}}$ も同様に有理数と平方根で表すことができる。

このような無理な距離を使って、さらに多くの直線や円を構築し、その係数も今度は平方根を含むものとすることができる。しかし、これらの方法を使用して構築できるすべての長さは、依然として平方根のみで表現されることが明白である。そして、私たちの無理数表現がより複雑な形を取るにせよ、これは繰り返し何回でも真実である。したがって、\emph{ユークリッドの方法は、平方根のみを含む任意の無理数表現を構築し、それ以外のものは構築できない}。

古代の有名な問題の一つが立方体倍積問題であり、これはEuclidの方法で長さ$\sqrt{2}$を構築することを意味します。$\sqrt{2}$は有理数や平方根のみによる有限な組み合わせで表すことができず、したがってこの問題は不可能であることが示されます。詳細についてはHobsonの『円の正方形』47ページ以降を参照してください；証明の最初の段階、つまり$\sqrt{2}$が有理係数を持つ二次方程式$ax^{2} + 2bx + c = 0$の根でないことを示す証明は、\okrickRef{第1章}{I}（演習問題I 24）にあります。\PageSep{68}

23. 接近的な円周平方問題。半径$R$の円の中心を$O$とする。点$A$における接線から$AP = \frac{11}{5}R$、$AQ = \frac{13}{5}R$と取る（同じ方向）。$AO$上に$AN = OP$を取って$NM$を$OQ$に平行で$AP$と交わる点$M$を通るように引く。すると
\[
AM/R = \frac{13}{25}\sqrt{146}
\]
であり、$AM$を円の周長として取り入れると$\pi$の値が小数点五桁まで正確になる。もし地球の半径$R$をもとむれば、$AM$をその周長とする誤差は約11ヤード未満である。

24. 僅用直尺从给定単位长度开始可以构造的唯一长度是有理长度。

25. 立方根2の構造。点$O$が放物線$y^{2} = 4x$の頂点で、点$S$が焦点であり、点$P$がこの放物線と放物線$x^{2} = 2y$の交点であるとする。すると$OP$は最初の放物線の焦点直線上の点$Q$を通過し、そのとき$SQ = \sqrt{2}$となる。

単位直径の円を考え、直径 \(OA\) とその接線を考える。円上の点 \(B\) および接線との交点 \(C\) を通る弦 \(OBC\) を引く。この直線上で \(OM = BC\) とする。原点を \(O\) とし、\(x\) 軸を \(OA\) として、点 \(M\) の軌跡は曲線
\[
(x^{2} + y^{2})x - y^{2} = 0
\]
（ディオクラレスの \emph{シソイド}）であることを示す。この曲線を図示せよ。\(y\) 軸方向に長さ \(OD = 2\) を取る。直線 \(AD\) が曲線と交わる点を \(P\)、直線 \(OP\) が円の接線 \(A\) との交点を \(Q\) とする。このとき \(AQ = \sqrt{2}\) であることを示す。
\par\medskip
\ページセパレーター{69}

\chapter{第3章：Chapter III: COMPLEX NUMBERS}

第三章  複素数

34. 直線と平面における移動。二つの前章で扱った「実数」~$x$は、多くの視点から見ることができます。純粋な数として考えられる一方で、少なくとも三つの異なる方法で幾何学的な意味が付け加えられます。それは、直線~$\Lambda$（第一章参照）上の長さ~$A_{0}P$の尺度として~$x$を考えてみることです。あるいは、点~$P$の標識として~$x$を考えることもできます、すなわち、~$A_{0}$から距離が~$x$である点~$P$です。また、それは直線~$\Lambda$上の位置変化または移動の尺度として~$x$を考えてみることもできます。この最後の視点に焦点を当てることが私たちの目的です。

直線~$\Lambda$上に点~$P$に置かれた小さな粒子を想像し、その後~$Q$へと移動させます。粒子から~$P$を~$Q$へ移動させるために必要な移動または位置変化を「移動~$\Seg{PQ}$」と呼びます。完全な移動を特定するには三つの要素が必要です：その「大きさ」、直線沿いの前進または後退方向の「符号」、つまり粒子の元の位置~$P$。しかし、移動によって生じる位置変化のみを考えている場合、移動の点の適用を無視し、長さと符号が同じであれば全ての移動は等価であると考えることが自然です。すると、移動は長さ~$PQ = x$と、符号が~$x$の正負によって決まる方向で完全に指定されます。したがって、曖昧性なく「移動~$[x]$」について話すことができます。そして、$\Seg{PQ} = [x]$と書くことができます。

ページ分割：70

私たちは四角括弧を使用して，位相移動$[x]$を長さまたは数$x$から区別します。\footnote{厳密には，測定する長さ$x$とその実際の長さを区別するために，似たような記号的な差異が必要です。読者はこのような区別が無駄で形式的であると考えるかもしれません。しかし，数学の経験を積むにつれて，彼は異なるものの間にはいかに密接に関連しても同じものではないことを明確に区別することが非常に重要であることが明らかになるでしょう。クリケットが数学的な科学であれば，打者Pのウィック間の移動，得点Q，およびスコアブックに記入されるマークはそれぞれ異なる概念として区別する必要があるでしょう。}
点$P$の座標が$a$であるとき，点$Q$の座標は$a + x$となります。したがって，位相移動$[x]$は点$a$から点$a + x$に粒子を移動させます。

次に平面における位相移動について考えます。\emph{位相移動$\Seg{PQ}$}の定義も同様に行います。しかし，完全に指定するためにより多くの情報が必要です。必要な情報は以下の通りです：(i) 位相移動の\emph{大きさ}，すなわち直線$PQ$の長さ；(ii) 位相移動の\emph{方向}，これは$PQ$が平面内の固定した直線と作る角度によって決定されます；(iii) 位相移動の\emph{向き}；および(iv) 位相移動の\emph{作用点}。これらの要件のうち，第四は考慮に入れずに済みます。もし二つの位相移動が大きさ，方向，向きにおいて同じであれば，これらを等価とします。すなわち，$PQ$と$RS$が等しく平行であり，$P$から$Q$への移動の向きが$R$から$S$への移動と同じであれば，位相移動$\Seg{PQ}$と$\Seg{RS}$を等価とし，以下のように表記します：
\[
\Seg{PQ} = \Seg{RS}.
\]

平面において任意の座標軸（例えば、Fig.19 の $OX$, $OY$）を取る。線分 $OA$ を引いて、$OA$ は $PQ$ に等しく平行であり、$O$ から $A$ への移動方向が $P$ から $Q$ への移動方向と一致するようにする。すると $\Seg{PQ}$ と $\Seg{OA}$ は同値な移動である。座標 $x$ および $y$ を $A$ の座標とする。このとき、$\Seg{OA}$ が完全に指定されることが明らかであり、これを「移動 $[x, y]$」と呼ぶ。そして
\[
\Seg{OA} = \Seg{PQ} = \Seg{RS} = [x, y]
\]
と書く。

Paragraph 35. 移動の同値性および数による移動の乗法。
点 $P$ の座標が $\xi$, $\eta$、点 $Q$ の座標が $\xi'$, $\eta'$ のとき、明らかに
\[
x = \xi' - \xi,\quad y = \eta' - \eta.
\]
したがって、$(\xi, \eta)$ から $(\xi', \eta')$ への移動は
\[
[\xi' - \xi, \eta' - \eta].
\]

二つの移動 $[x, y]$, $[x', y']$ が同値であることは、かつ-only-if- $x = x'$, $y = y'$ であることが明らかである。したがって
\[
[x, y] = [x', y']
\]
は
\[
x = x',\quad y = y'
\Tag{(1)}
\]
と同値である。

逆移動 $\Seg{QP}$ は $[\xi - \xi', \eta - \eta']$ であり、自然に
\begin{align*}
[\xi - \xi', \eta - \eta'] &= -[\xi' - \xi, \eta' - \eta],\\
\Seg{QP} &= -\Seg{PQ},
\end{align*}
を定義する。これらの等式は実質的に符号 $-[\xi' - \xi, \eta' - \eta]$, $-\Seg{PQ}$ の意味を定義している。したがって
\[
-[x, y] = [-x, -y],
\]
と同意し、さらに
\[
\alpha[x, y] = [\alpha x, \alpha y]
\Tag{(2)}
\]
を同意する。ここで $\alpha$ は任意の実数であり、正でも負でもよい。したがって Fig.19 の場合、$OB = -\frac{1}{2}OA$ ならば
\[
\Seg{OB} = -\tfrac{1}{2}\Seg{OA} = -\tfrac{1}{2}[x, y]
  = [-\tfrac{1}{2}x, -\tfrac{1}{2}y].
\]

等式 \Eq{(1)} および \Eq{(2)} は移動に関連する最初の二つの重要な概念を定義し、すなわち「移動の同値性」および「数による移動の乗法」である。

\Paragraph{36. 遷移の加法} まだ、次の表現に対する意味を定義する方法について述べていません：
\[
\Seg{PQ} + \Seg{P'Q'},\quad [x, y] + [x', y'].
\]
直感からすぐにわかるように、二つの遷移の和を定義するには、二つの与えられた遷移を順次適用した結果となる遷移として定義すべきです。つまり、$QQ_{1}$ を $P'Q'$ に等しく並行に引くとします。粒子が $P$ からまず $Q$ へ、その後 $Q_{1}$ へ移動するという二つの遷移 $\Seg{PQ}, \Seg{P'Q'}$ の結果として、$\Seg{PQ}$ と $\Seg{P'Q'}$ の和を $\Seg{PQ_{1}}$ とする定義が自然です。したがって、$OA$ を $PQ$ に等しく並行に引き、$OB$ を $P'Q'$ に等しく並行に引きます。平行四辺形 $OACB$ を完成すると次のようになります：
\%[図表: Fig. 20.]
\begin{figure}[h!]
    \centering
    \fbox{【図版】}
    \caption{Fig. 20}
\end{figure}

\[
\Seg{PQ} + \Seg{P'Q'} = \Seg{PQ_{1}} = \Seg{OA} + \Seg{OB} = \Seg{OC}.
\]

この定義を採用した場合の結果を考えてみましょう。$B$ の座標が $x', y'$ であるとき、$AB$ の中点の座標は $\frac{1}{2}(x + x'), \frac{1}{2} (y + y')$ であり、$C$ の座標は $x + x', y + y'$ です。したがって
\[
[x, y] + [x', y'] = [x + x', y + y'],
\Tag{(3)}
\]
これは遷移の加法をシンボリックに定義する方法として考えることができます。さらに
\begin{align*}
[x', y'] + [x, y]
  &= [x' + x, y' + y]\\
  &= [x + x', y + y'] = [x, y] + [x', y']
\end{align*}
が成り立ちます。つまり、\emph{遷移の加法は通常の代数で $a + b = b + a$ と表される交換律を満たす}ということです。この法則は、粒子がまず距離 $PQ_{2}$ を $P'Q'$ に等しく並行に移動し、その後 $PQ$ に等しく並行な距離を移動した場合でも、元の点 $Q_{1}$ に到達することを明らかにしています。
\PageSep{73}

特に
\[
[x, y] = [x, 0] + [0, y].
\Tag{(4)}
\]
ここで $[x, 0]$ は、$OX$ 方向に平行な距離 $x$ の移動を表します。これは以前、直線方向の移動のみを考えていたときに $[x]$ で表したものと同じです。$[x, 0]$ と $[0, y]$ を $[x, y]$ の \emph{成分}、そして $[x, y]$ をその \emph{結果}と呼びます。

二つの移動の加算が定義されたら、任意の数の加算にはさらに問題はありません。従って、定義により
\begin{gather*}
[x, y] + [x', y'] + [x'', y'']
  = ([x, y] + [x', y']) + [x'', y'']\\
  = [x + x', y + y'] + [x'', y'']
  = [x + x' + x'', y + y' + y''].
\end{gather*}

移動の \emph{減算} を定義するには、方程式
\[
[x, y] - [x', y'] = [x, y] + (-[x', y']),
\Tag{(5)}
\]
を用います。これは $[x, y] + [-x', -y']$ または $[x - x', y - y']$ と同じです。
特に
\[
[x, y] - [x, y] = [0, 0].
\]

移動 $[0, 0]$ は粒子の位置を変えません；これは \emph{零移動}で、$[0, 0] = 0$ と書くことにします。

\par\medskip\noindent\textbf{【例 XX.】}\par
\Item{1.} 次を証明せよ：

\SubItem{(i)} $\alpha [\beta x, \beta y] = \beta [\alpha x, \alpha y] = [\alpha \beta x, \alpha \beta y]$,

\SubItem{(ii)} $([x, y] + [x', y']) + [x'', y''] = [x, y] + ([x', y'] + [x'', y''])$,

\SubItem{(iii)} $[x, y] + [x', y'] = [x', y'] + [x, y]$,

\SubItem{(iv)} $(\alpha + \beta) [x, y] = \alpha [x, y] + \beta [x, y]$,

\SubItem{(v)} $\alpha \{[x, y] + [x', y']\} = \alpha [x, y] + \alpha [x', y']$.

[(iii) については既に証明済みです。残りの方程式は定義から等しく証明できます。各場合において、方程式の幾何学的な意味を考えてみてください。]

\Item{2.} $M$ が $PQ$ の中点であるとき、$\Seg{OM} = \frac{1}{2}(\Seg{OP} + \Seg{OQ})$ です。さらに一般に、$M$ が $PQ$ を $\mu : \lambda$ の比で分けるときには
\[
\Seg{OM}
  = \frac{\lambda}{\lambda + \mu}\, \Seg{OP}
  + \frac{\mu}{\lambda + \mu}\, \Seg{OQ}.
\]

3. $G$が等しい質量を持つ粒子$P_{1}$、$P_{2}$、…、$P_{n}$の重心であるとき、
\[
\Seg{OG} = \frac{\Seg{OP_{1}} + \Seg{OP_{2}} + \dots + \Seg{OP_{n}}}{n}.
\]

4. 平面上で共線な点$P$、$Q$、$R$が与えられたとき、実数$\alpha$、$\beta$、$\gamma$（すべてゼロではない）を用いて
\[
\alpha \cdot \Seg{OP} + \beta \cdot \Seg{OQ} + \gamma \cdot \Seg{OR} = 0;
\]
が成り立つことが可能であり、逆もまた真である。［これは実質的にEx.2の別の表現に過ぎない］。

5. 直線とは異なる方向にある二つの移動$\Seg{AB}$と$\Seg{AC}$について、
\[
\alpha \cdot \Seg{AB} + \beta \cdot \Seg{AC} = \gamma \cdot \Seg{AB} + \delta \cdot \Seg{AC},
\]
が成り立つとき、$\alpha = \gamma$ および $\beta = \delta$。

［$AB_{1} = \alpha \cdot AB$、$AC_{1} = \beta \cdot AC$と置き、平行四辺形$AB_{1}P_{1}C_{1}$を完成させる。すると$\Seg{AP_{1}} = \alpha \cdot \Seg{AB} + \beta \cdot \Seg{AC}$となる。$\Seg{AP_{1}}$はこの形で唯一の表現しかできないことが明白であり、従って定理が成り立つ］。

6. $ABCD$が平行四辺形であるとき、その内部に点$Q$をとり、$RQS$と$TQU$を$AB$、$AD$に平行な線分として引く。このとき$RU$と$TS$は$AC$上で交わることを示せ。
［図21参照］

［図2.75in: 21, p074］

点$Q$から$AB$、$AD$への比$\alpha$、$\beta$を用いて
\[
\Seg{AT} = \alpha \cdot \Seg{AB},\quad \Seg{AR} = \beta \cdot \Seg{AD},
\]
\[
\Seg{AU} = \alpha \cdot \Seg{AB} + \Seg{AD},\quad \Seg{AS} = \Seg{AB} + \beta \cdot \Seg{AD}.
\]

$RU$と$AC$の交点を$P$とする。このとき、$R$、$U$、$P$が共線であるから
\[
\Seg{AP}
  = \frac{\lambda}{\lambda + \mu} \, \Seg{AR}
  + \frac{\mu}{\lambda + \mu} \, \Seg{AU},
\]
ここで$\mu/\lambda$は$P$が$RU$を分割する比である。つまり
\[
\Seg{AP}
  = \frac{\alpha \mu}{\lambda + \mu} \, \Seg{AB}
  + \frac{\beta \lambda + \mu}{\lambda + \mu} \, \Seg{AD}.
\]

しかし、$P$が$AC$上にあるため、$\Seg{AP}$は$\Seg{AC}$の数的倍数であり、これを$\Seg{AP} = k \cdot \Seg{AC} = k \cdot (\Seg{AB} + \Seg{AD})$と表すことができる。したがって（例題5）、$\alpha\mu = \beta\lambda + \mu = (\lambda + \mu)k$から、
\[
k = \frac{\alpha\beta}{\alpha + \beta - 1}
\]
を得る。この結果の対称性は、$TS$と$AC$の交点である$P'$についても同様な議論により
\[
\Seg{AP'} = \frac{\alpha\beta}{\alpha + \beta - 1} \cdot \Seg{AC}
\]
が成り立つことを示す。したがって、$P$と$P'$は同一の点である。

7. $ABCD$は平行四辺形であり、$M$は$AB$の中点である。$DM$が$AC$を3等分し、また$AC$によって3等分されることを示せ。\footnote{前の2つの例はウィラード・ギブズの『ベクトル解析』から引用しています。}

\Paragraph{37. 移動量の乗法。}これまで、二つの移動量の積という概念にどのような意味も付与しようと試みたことはありません。考慮した唯一の種類の乗法は、移動量が数値で乗算されるものでした。式
\[
[x, y] \times [x', y']
\]
はこれまで何ら意味を持たず、これを任意に定義することができます。ただし、このような積の定義が有用であるためには、二つの移動量の積もまた移動量であることが明白です。

例えば、我々は次のように定義することができる：
\[
[x + x', y + y'];
\]
すなわち、二つの移動の積が常にその和と等しいものとする約束をすることができる。しかし、このような定義には以下の二つの重要な問題がある。まず第一に、この定義は無意味である。新しい表現方法を導入することによって何かを完全に表現できるからだ。第二に、以下のような理由により不便で誤解の原因となる可能性がある。もし$\alpha$が実数であれば、我々は既に$\alpha [x, y]$を$[\alpha x, \alpha y]$と定義している。そして、第34節で見たように、実数$\alpha$は一方から見れば移動である。すなわち、軸$OX$上の移動$[\alpha]$または後の表記法では$[\alpha, 0]$である。したがって、絶対的に必要とは言えなくとも、少なくとも望ましいのは次の定義である：
\[
[\alpha, 0] [x, y] = [\alpha x, \alpha y]
\]
であり、提案された定義はこの結果を出さない。

より合理的な定義としては次のようなものがあるかもしれない：
\[
[x, y] [x', y'] = [xx', yy']
\]
しかし、これにより
\[
[\alpha, 0] [x, y] = [\alpha x, 0]
\]
となり、この定義も第二の問題に直面する可能性がある。

実際、$[x, y] [x', y']$ の意味をどのように定義するかは明らかではありません。私たちの定義が有用であるためには、この積は$x$と$y$に依存する移動であり、つまり
\[
[x, y] [x', y'] = [X, Y],
\]
でなければならないことが明確です。ここで$X$と$Y$は$x$, $y$, $x'$, $y'$の関数です。（2）また、定義は以下の式と一致するべきであり
\[
[x, 0] [x', y'] = [xx', xy'];
\]
（3）さらに、定義は乗法の通常の交換律、分配律、結合律を遵守する必要があります。つまり
\begin{align*}
[x, y] [x', y'] &= [x', y'] [x, y],\\
([x, y] + [x', y']) [x'', y''] &= [x, y] [x'', y''] + [x', y'] [x'', y''],\\
[x, y] ([x', y'] + [x'', y'']) &= [x, y] [x', y'] + [x, y] [x'', y''],\\
\intertext{そして}
[x, y] ([x', y'] [x'', y'']) &= ([x, y] [x', y']) [x'', y''].
\end{align*}

38. その定義は以下のように提案されます。$OAB$, $OCD$が相似な三角形であるとき、順番に書かれた角度に対応する角であれば
\[
OB/OA = OD/OC,
\]
または$OB · OC = OA · OD$です。これは、移動の乗法と除法を定義することで
\[
\Seg{OB}/\Seg{OA} = \Seg{OD}/\Seg{OC},\quad
\Seg{OB} · \Seg{OC} = \Seg{OA} · \Seg{OD}
\]
となるように試みるべきであることを示唆しています。\%[図22の挿絵]

今、次のように定義する：

\[
\Seg{OB} = [x, y],\quad
\Seg{OC} = [x', y'],\quad
\Seg{OD} = [X, Y],
\]

そして点 $A$ が $(1, 0)$ のとき、$\Seg{OA} = [1, 0]$ であるとする。すると、

\[
\Seg{OA} · \Seg{OD} = [1, 0] [X, Y] = [X, Y],
\]

よって

\[
[x, y] [x', y'] = [X, Y].
\]

$\Seg{OB} · \Seg{OC}$ の積は、$OC$ 上に $OAB$ と相似な三角形を構成することで $\Seg{OD}$ として定義される。この定義が曖昧さを避けるためには、$OC$ 上で二つのような三角形 $OCD$ と $OCD'$ を描くことができる点 $D$ の選び方を考える必要がある。角度 $COD$ が $AOB$ と同じ符号と大きさを持つときのものを選ぶこととする。このとき、二つの三角形は「同じ意味での相似」であるという。

もし点 $B$ と $C$ の極座標が $(\rho, \theta)$ と $(\sigma, \phi)$ のとき、

\[
x = \rho\cos\theta,\quad
y = \rho\sin\theta,\quad
x' = \sigma\cos\phi,\quad
y' = \sigma\sin\phi,
\]

であるから、点 $D$ の極座標は明らかに $\rho\sigma$ と $\theta + \phi$ となる。したがって、

\begin{alignat*}{2}
X \&= \rho\sigma\cos(\theta + \phi) \&\&= xx' - yy',\\
Y \&= \rho\sigma\sin(\theta + \phi) \&\&= xy' + yx'.
\end{alignat*}

したがって必要な定義は、

\[
[x, y] [x', y'] = [xx' - yy', xy' + yx'].
\Tag{(6)}
\]

(1) $y = 0$ のとき、$X = xx'$, $Y = xy'$ となり、これが望んでいた通りである。(2) 右辺は $x$ と $x'$, $y$ と $y'$ を入れ替えることで変化しないので、

\[
[x, y] [x', y'] = [x', y'] [x, y];
\]

(3)

\begin{multline*}
\{[x, y] + [x', y']\} [x'', y''] = [x + x', y + y'] [x'', y'']\\
\begin{aligned}[t]
  &= [(x + x') x'' - (y + y') y'', (x + x') y'' + (y + y') x'']\\
  &= [xx'' - yy'', xy'' + yx''] + [x'x'' - y'y'', x'y'' + y'x'']\\
  &= [x, y] [x'', y''] + [x', y'] [x'', y''],
\end{aligned}
\end{multline*}

が成り立つ。

同様に、$\SecNo[§]{37}$ の末尾のすべての方程式も満たされることを確認できる。したがって定義 \Eq{(6)} は $\SecNo[§]{37}$ で述べた要求を全て満たしている。

\par\medskip\noindent\textbf{【注釈】}\par
\par\medskip\noindent\textbf{例.}\par\smallskip 以上の幾何学的定義から、移動の乗算が交換法則と分配法則を満たすことを直接示してください。[例えば交換法則について考えてみましょう。$\Seg{OB} \cdot \Seg{OC}$は$\Seg{OD}$に等しくなります（図22）、$COD$は$AOB$に相似です。$\Seg{OC} \cdot \Seg{OB}$の積を構成するには、$OB$上に$AOC$に相似な三角形$BOD_{1}$を構成しなければなりません；したがって証明したいことは$D$と$D_{1}$が一致すること、すなわち$BOD$は$AOC$に相似であること。これは初等幾何学の簡単な問題です。]
\par\medskip

\Paragraph{39. 複素数。} 一次元空間における移動$[x]$に対応する点$(x)$と実数$x$があるように、平面における移動$[x, y]$に対応する点$(x, y)$と実数のペア$x$, $y$があります。

このペアの実数$x$, $y$を次のシンボルで表すことにします。
\[
x + yi.
\]
この表記法を選んだ理由は後ほど明らかになります。現在のところ、読者は$x + yi$を単に$[x, y]$の別の書き方として考えてください。式$x + yi$は複素数と呼ばれます。

次に複素数の等価性、加算、乗算を定義します。任意の複素数に対応する移動があります。二つの複素数が等価であるとは、対応する移動が等価であることを指します。二つの複素数の和または積は、その対応する移動の和または積に対応する複素数です。したがって
\[
x + yi = x' + y'i,
\Tag{(1)}
\]
が成り立つのは$x = x'$, $y = y'$であるときだけであり、また
\begin{gather*}
(x + yi) + (x' + y'i) = (x + x') + (y + y')i;
\Tag{(2)}\\
(x + yi) (x' + y'i) = xx' - yy' + (xy' + yx')i.
\Tag{(3)}
\end{gather*}
が成り立つ。

特に、式(2)と式(3)の特殊な場合として、

\begin{gather*}
x + yi = (x + 0i) + (0 + yi),\\
(x + 0i)(x' + y'i) = xx' + xy'i;
\end{gather*}

これらの式は、複素数を扱う際に $x$ を $x + 0i$、$yi$ を $0 + yi$ として書くことで混乱が生じないことを示唆します。

複素数の正整数乗と多項式も通常の代数学で定義されます。例えば、$x = x'$, $y = y'$を式(3)に代入すると、

\[
(x + yi)^2 = (x + yi)(x + yi) = x^2 - y^2 + 2xyi.
\]

次頁（79）

読者は容易に確認できるでしょうが、複素数の加法と乗法は代数学の法則を満たし、

\begin{gather*}
(x + yi) + (x' + y'i) = (x' + y'i) + (x + yi),\\
\{(x + yi) + (x' + y'i)\} + (x'' + y''i) = (x + yi) + \{(x' + y'i) + (x'' + y''i)\},\\
(x + yi)(x' + y'i) = (x' + y'i)(x + yi),\\
(x + yi)\{(x' + y'i) + (x'' + y''i)\} = (x + yi)(x' + y'i) + (x + yi)(x'' + y''i),\\
\{(x + yi) + (x' + y'i)\}(x'' + y''i) = (x + yi)(x'' + y''i) + (x' + y'i)(x'' + y''i),\\
(x + yi)\{(x' + y'i)(x'' + y''i)\} = \{(x + yi)(x' + y'i)\}(x'' + y''i),
\end{gather*}

これらの等式的証明は、対応する移動に対する等式の証明とほとんど同じです。

複素数の引き算と除法も通常の代数学で定義されます。例えば、

\[
(x + yi) - (x' + y'i) = (x + yi) + \{- (x' + y'i)\} = x + yi + (-x' - y'i) = (x - x') + (y - y')i;
\]

または、数 $\xi + \eta i$ として定義され、

\[
(x' + y'i) + (\xi + \eta i) = x + yi,
\]

これにより同じ結果が得られます。また、

\[
(x + yi)/(x' + y'i)
\]

は複素数 $\xi + \eta i$ として定義され、

\[
(x' + y'i)(\xi + \eta i) = x + yi,
\]

または

\[
x'\xi - y'\eta + (x'\eta + y'\xi)i = x + yi,
\]

または

\[
x'\xi - y'\eta = x,\quad
x'\eta + y'\xi = y.
\tag{4}
\]

これらの方程式を$\xi$と$\eta$について解くと、以下の式を得る：
\[
\xi  = \frac{xx' + yy'}{x'^{2} + y'^{2}},\quad
\eta = \frac{yx' - xy'}{x'^{2} + y'^{2}}.
\]
この解は$x'$と$y'$がともにゼロである場合、すなわち$x' + y'i = 0$の場合には失敗する。
従って引き算は常に可能であり、除算も除数がゼロでない限り常に可能である。

---

以下が例題である：

1. 几何学的な観点から見て、移動ベクトル$\Seg{OB}$を$\Seg{OC}$で分割する問題は、三角形$COB$と$AOD$が相似であるような点$D$を見つける問題であり、これは明らかに可能（かつ解は一意的）である、ただし$C$が$0$と一致しない限り。

2. 数字$x + yi$と$x - yi$は「共役」と呼ばれる。これらが積である実数を確認するには以下の式を検証すればよい：
\[
(x + yi)(x - yi) = x^{2} + y^{2},
\]
また、共役の数字の積は実数であり、さらに
\[
\frac{x + yi}{x' + y'i}
  = \frac{(x + yi)(x' - y'i)}{(x' + y'i)(x' - y'i)}\\
  = \frac{xx' + yy' + (x'y - xy')i}{x'^{2} + y'^{2}}.
\]

---

40. 実数の最も重要な性質の一つは「因数定理」であり、これは「二つの数の積がゼロであるためには少なくとも一方がゼロである」というものである。複素数についてもこの性質が成り立つことを証明するために、前節の式（4）において$x = 0$、$y = 0$を代入する。すると以下の式を得る：
\[
x'\xi - y'\eta = 0,\quad
x'\eta + y'\xi = 0.
\]
これらの式から$\xi = 0$、$\eta = 0$であり、すなわち
\[
\xi + \eta i = 0,
\]
が得られる。ただし$x' = 0$かつ$y' = 0$、または$x' + y'i = 0$の場合はこの例外である。従って$x + yi$がゼロになるためには、$x' + y'i$または$\xi + \eta i$がゼロになる必要がある。

41. 方程式 $i^{2} = -1$。私たちは、$x + 0i$ を $x$、$0 + yi$ を $yi$ と表記することにより符号を簡略化することに同意しました。特に複素数 $1i$ は単に $i$ で表記します。これは $OY$ 方向の単位移動に対応する数です。また
\[
i^{2} = ii = (0 + 1i) (0 + 1i) = (0 \cdot 0 - 1 \cdot 1) + (0 \cdot 1 + 1 \cdot 0)i = -1.
\]
同様に $(-i)^{2} = -1$ も成立します。したがって、複素数 $i$ と $-i$ は方程式 $x^{2} = -1$ を満たします。

読者は容易に確認できるでしょうが、複素数の加法と乗法の規則により得られる結論は、実数に対する操作と同じように複素数を扱うということです。符号 $i$ は数として扱いますが、常に $ii = i^{2}$ を $-1$ に置き換えることに注意してください。例えば，
\begin{align*}
(x + yi) (x' + y'i) &= xx' + xy'i + yx'i + yy'i^{2}\\
                    &= (xx' - yy') + (xy'+ yx')i.
\end{align*}

42. $i$ による複素数の幾何学的解釈。$(x + yi)i = -y + xi$ より、$x + yi$ が $\Seg{OP}$ に対応し、$OQ$ を $\Seg{OP}$ 等しい長さで $POQ$ が正の直角となるように引くと、$(x + yi)i$ は $\Seg{OQ}$ に対応します。つまり、複素数を $i$ で乗算すると、その対応する移動が直角に回転することです。

We might have developed the whole theory of complex numbers from this point of view. Starting with the ideas of \(x\) as representing a displacement along \(OX\), and of \(i\) as a symbol of operation equivalent to turning \(x\) through a right angle, we should have been led to regard \(yi\) as a displacement of magnitude \(y\) along \(OY\). It would then have been natural to define \(x + yi\) as in §37 and §40, and \((x + yi)i\) would have represented the displacement obtained by turning \(x + yi\) through a right angle, i.e., \(-y + xi\). Finally, we should naturally have defined \((x + yi)x'\) as \(xx' + yx'i\), \((x + yi)y'i\) as \(-yy' + xy'i\), and \((x + yi)(x' + y'i)\) as the sum of these displacements, i.e., as
\[
xx' - yy' + (xy' + yx')i.
\]

43. The equations \(z^2 + 1 = 0\) and \(az^2 + 2bz + c = 0\). There is no real number \(z\) such that \(z^2 + 1 = 0\); this is expressed by saying that the equation has \emph{no real roots}. But, as we have just seen, the two complex numbers \(i\) and \(-i\) satisfy this equation. We express this by saying that the equation has \emph{the two complex roots} \(i\) and \(-i\). Since \(i\) satisfies \(z^2 = -1\), it is sometimes written in the form \(\sqrt{-1}\).

複素数はしばしば「虚数」と呼ばれる。\footnote{「実数」という用語は、「虚数」に対して反対の意味を持つものとして導入された。}
この表現は決して幸運に見えたものではないが、既に確立されており、受け入れざるを得ない。しかし読者には強調する必要があるのは、「虚数」という言葉は通常の意味での「虚幻的な」存在ではなく、「実数」よりも実在性を持つという点である。そしてそれは実数とは異なる意味での「数」でなく、前節の議論から明らかになるように\emph{数対$(x, y)$}であり、技術的な便宜のためにシンボリカルに結び付けられた形$x + yi$である。このような数対は$\frac{1}{2}$やこの紙、太陽系など何らかの通常の数と同様に「実在的」である。

例えば
\[
i = 0 + 1i
\]
は数対$(0, 1)$を表し、幾何学的には点または移動$[0, 1]$で表現できる。そして$i$が方程式$z^{2} + 1 = 0$の根であるという意味は、このような数対（または移動）を「乗算」と呼ぶ方法を定義し、$(0, 1)$を自身と組み合わせると結果として$(-1, 0)$を得ることである。

次により一般的な方程式
\[
az^{2} + 2bz + c = 0,
\]
を考える。ここで$a$、$b$、$c$は実数である。もし$b^{2} > ac$の場合は、通常の解法により二つの実数根
\[
\{-b ± \sqrt{b^{2} - ac}\}/a
\]
が得られる。もし$b^{2} < ac$の場合は方程式には実数の根がない。この方程式は次のように書き換えることができる：
\[
\{z + (b/a)\}^{2} = -(ac - b^{2})/a^{2},
\]
これは、$z + (b/a)$に$\pm i \sqrt{ac - b^{2}}/a$を代入することで満たされることが明らかである。\footnote{しばしば印刷の便宜のために$x + iy$と書くこともある。}
このことを表すために、「方程式は二つの複素根\emph{$\{-b ± i \sqrt{ac - b^{2}}\}/a$}を持つ」という表現を使う。

もし$b^2 = ac$という関係式が convenzione として成り立つとすれば（この場合、方程式はただ一つの$x$の値、すなわち$-b/a$で満たされる）、方程式は二つの等しい根を持つことになります。したがって、実数係数の二次方程式はすべての場合に二つの根を持ちます：異なる実数根、または等しい実数根、または異なる複素根です。

ページ分割{83}

複素根を認めるとするとき、二次方程式がそれ以上に多くの根を持つことは自然と疑問になります。しかし、これは不可能であることが容易に確認できます。その不可能性は、基本的な代数学で$n$次方程式が$n$個以上の実数根を持つことができないことを証明するために使われる同じ論理の流れによって証明することができます。任意の複素数$a$について、$f(z)$を$z - a$で割った余りは$f(a)$であることを順に証明します：

(1) $f(z)$が任意の実数または複素数$a$で割られたときの余りは$f(a)$である。

(2) 方程式$f(z) = 0$の根である$a$について、$f(z)$は$z - a$で割り切れる。

(3) $f(z)$が$n$次であり、$f(z) = 0$の$n$つの根が$a_1, a_2, \dots, a_n$であるとき、

\[
f(z) = A (z - a_1)(z - a_2) \dots (z - a_n),
\]

ここで$A$は実数または複素数であり、実際には$f(z)$の$x^n$の係数です。最後の結果と第40節の定理から、$f(z)$は$n$個以上の根を持つことができないことがわかります。

我々は、実数係数を持つ二次方程式が恰好2つの根を持つことを結論付けることができる。後で見るように、類似の定理は任意の次数と実数または複素数係数を持つ方程式についても真である：\emph{n次方程式は恰好nつの根を持つ}。証明の中で困難を呈する唯一の点は最初の部分であり、つまり、どんな方程式にも少なくとも1つの根があることを証明することである。この証明については現在そのままでおいておくこととする。\footnote{詳しくは\okrickRef{付録}{I}をご覧ください。}
しかし、この定理の非常に興味深い結果の一つについてすぐに注意を喚起しておくことができる。
数論においては、正整数と加法・乗法およびその逆演算である減法・除法から始めることで理論が始まる。これらの操作が常に可能であるとは限らないことが見つかる。$3 - 7$の意味を付けるには負の数を認めなければならず、$\frac{3}{7}$の意味を付けるには有理分数を認めなければならないことがわかる。
我々が平方根の抽出や方程式の解法などの算術的操作を拡張し、それらを含めると、例えば2のような完全な平方でない数の平方根の抽出のような操作は、新しい数の概念を認める必要があり、\emph{無理数}（第1章参照）を認めることで可能になることがわかる。

他の例としては、虚数√-1の抽出は、複素数をこの章で認めなければ不可能です。そして、方程式の次数が上がると、それらが複素数すらも解けないものになる可能性があると想像するのは自然でしょう。その結果、より高度な種類の、いわば超複素数への考慮が必要となるかもしれません。どのような代数方程式の根でも普通の複素数であるという事実は、これがそうではないことを示しています。普通の代数的操作を複素数に適用しても、それらの操作結果もまた必ず複素数となります。技術的な表現で言えば「複素数の域は代数的操作に対して閉じている」と言えます。

他のこと柄に移る前に、以下の点を追加します。初等代数において、足し算と掛け算の規則のみを用いて証明されたすべての定理は、その中で現れる数が実数であれ複素数であれ真です。例えば、二次方程式
\[
az^{2} + 2bz + c = 0,
\]
の根αとβに対しては
\[
\alpha + \beta = -(2b/a),\quad
\alpha\beta = (c/a).
\]

同様に、三次方程式
\[
az^{3} + 3bz^{2} + 3cz + d = 0,
\]
の根α、β、γに対しては
\[
\alpha + \beta + \gamma = -(3b/a),\quad
\beta\gamma + \gamma\alpha + \alpha\beta = (3c/a),\quad
\alpha\beta\gamma = -(d/a).
\]
このようなすべての定理は、a、b、…、α、β、…が実数であれ複素数であれ真です。

44. アルガンの図。点P（図24）を（x, y）とし、rをOPの長さ、θをXOPの角度とするので
\[
x = r\cos\theta,\quad
y = r\sin\theta,\quad
r = \sqrt{x^{2} + y^{2}},\quad
\cos\theta : \sin\theta : 1 :: x : y : r.
\]

複素数 $x + yi$ を $z$ で表し、\SecNo[§]{43} に従います。そして $z$ を \emph{複素変数} と呼びます。
\%[図版: Fig. 24.]
Figure[2.5in]{24}{p085}
点 $P$ を $z$ \emph{の点}、または $z$ に対応する点と呼び、$z$ を \emph{主座標}、$x$ を \emph{実部}、$y$ を \emph{虚部}、$r$ を \emph{絶対値}、$\theta$ を \emph{振幅} と呼びます。そして
\[
%[** TN: 行を2つに分割して元の位置に]
x = \Real(z),\quad y = \Imag(z),\quad
r = |z|,\quad \theta = \am z.
\]

$y = 0$ のとき、$z$ を \emph{実数} と呼びます。$x = 0$ のとき、$z$ を \emph{純虚数} と呼びます。虚部の符号のみが異なる複素数 $x + yi$, $x - yi$ を \emph{共役} と呼びます。共役な複素数の和 $2x$ および積 $x^{2} + y^{2}$ はともに実数であり、両者の絶対値は $\sqrtp{x^{2} + y^{2}}$ で同じであり、その積はどちらかの絶対値の平方と等しいことに注意します。例えば実数係数を持つ二次方程式の根は非実数の場合には共役です。

θ または \(\am z\) は \(x\) と \(y\) の多価関数であり、無限の値を持つ角度で異なる。これは \(2\pi\) の倍数である。

\(z\) の絶対値 \(|z|\) は点 \(P\) の極座標 \(r\) に等しく、他の極座標 \(\theta\) は \(\am z\) の一つの値に等しい。この値が必ずしも定義される主値ではないことに注意しなければならない。第22節の極座標は \(0\) から \(2\pi\) までであり、主値は \(-\pi\) から \(\pi\) までである。

\(OX\) 上に最初位置にある直線をこれらの角度のいずれかで回転させると、その直線は点 \(P\) に到達する。これらの角度のうち、\(-\pi\) から \(\pi\) までのものを \(z\) の振幅の主値と呼ぶことにしよう。この定義は、ある値が \(\pi\) の場合を除いて明確である。その場合、\(-\pi\) も値となる。この場合、主値としてどの値を考慮するか特別な規定が必要となる。

一般に、\(z\) の振幅について話すとき、逆に述べられなければ主値の振幅を指すものとする。

通常、図 \(\Fig{24}\) は Argand の図と呼ばれる。

\Paragraph{45. デ・モivre の定理。} 次の声明は加法と乗法の定義から直接導かれる。

\Item{(1)} 二つの複素数の和の実部（または虚部）は、その実部（または虚部）の和に等しい。

\Item{(2)} 二つの複素数の積の絶対値は、その絶対値の積に等しい。

\Item{(3)} 二つの複素数の積の振幅は、その振幅の和と等しくなるか、またはそれから \(2\pi\) の倍数だけ異なる。

Remarkによれば、$\am(zz')$の主値が$\am z$と$\am z'$の主値の和とは限らないことに注意すべきです。例えば、$z = z' = -1 + i$の場合、$z$と$z'$のそれぞれの主値は$\frac{3}{4}\pi$ですが、$zz' = -2i$であり、$\am(zz')$の主値は$-\frac{1}{2}\pi$でなく$\frac{3}{2}\pi$ではありません。

最後の二つの定理を以下のように表すことができます：
\[
r(\cos\theta + i\sin\theta) \times \rho(\cos\phi + i\sin\phi)
  = r\rho\{\cos(\theta + \phi) + i\sin(\theta + \phi)\}.
\]
この等式は、$\cos(\theta + \phi)$と$\sin(\theta + \phi)$の通常の三角関数公式を使用してすぐに証明できます。

より一般的には
\begin{gather*}
r_{1}(\cos\theta_{1} + i\sin\theta_{1})
  \times r_{2}(\cos\theta_{2} + i\sin\theta_{2}) \times \dots
  \times r_{n}(\cos\theta_{n} + i\sin\theta_{n})\\
  = r_{1}r_{2} \dots r_{n} \{\cos(\theta_{1} + \theta_{2} + \dots + \theta_{n})
                        + i \sin(\theta_{1} + \theta_{2} + \dots + \theta_{n})\}.
\end{gather*}

特に興味深いケースは、
\[
r_{1} = r_{2} = \dots = r_{n} = 1, \quad
\theta_{1} = \theta_{2} = \dots = \theta_{n} = \theta
\]
の場合です。

この場合、以下の等式を得ます：
\[
(\cos\theta + i\sin\theta)^{n} = \cos n\theta + i\sin n\theta,
\]
ここで$n$は任意の正整数です。これはデ・モivreの定理と呼ばれています。\footnote
  {簡略化のため、$\cos\theta + i\sin\theta$を$\Cis\theta$で表すことが時々便利です：この記法はHarkness教授とMorley教授によって提案され、デ・モivreの定理は$(\Cis\theta)^{n} = \Cis n\theta$という等式で表現されます。}

再度、もし
\[
z = r(\cos\theta + i\sin\theta)
\]
であるなら
\[
1/z = (\cos\theta - i\sin\theta)/r.
\]
したがって、$z$ の逆数の模値は $z$ の模値の逆数であり、$z$ の逆数の位相は $z$ の位相の負である。これにより、$z$ の商に関する定理を \Eq{(2)} と \Eq{(3)} に対応する形で述べることができる。

(4) 任意の複素数の比の模値はその模値の比に等しい。

(5) 任意の複素数の比の位相はその位相の差に等しく、またはそれに $2\pi$ を加えたものである。

また
\begin{align*}
(\cos\theta + i\sin\theta)^{-n} &= (\cos\theta - i\sin\theta)^{n}\\
&= \{\cos(-\theta) + i\sin(-\theta)\}^{n}\\
&= \cos(-n\theta) + i\sin(-n\theta).
\end{align*}
したがって、デ・モivre の定理はすべての整数 $n$（正または負）に対して成り立つ。

(1) から (5) までの定理に加えて、以下の定理も非常に重要な定理である。

(6) 任意の複素数の和の模値はその模値の合計よりも大きいとは限らない。

図形的に、$\Seg{OP}$, $\Seg{OP'}$, \dots がそれぞれの複素数に対応する移動を表すとする。$PQ$ を $OP'$ に平行かつ等長とし、$QR$ を $OP''$ に平行かつ等長とし、次のように連続的に線を引く。最終的には点 $U$ に達し、
\[
\Seg{OU} = \Seg{OP} + \Seg{OP'} + \Seg{OP''} + \dots.
\]
の式が成り立つ。$\Seg{OU}$ の長さは複素数の和の模値であり、その模値の合計は線分 $OPQR\dots U$ の総長であるが、これは $\Seg{OU}$ よりも等しくない。

この定理の純粋に算術的な証明は \Exs{xxi}. 1 で概説されている。

46. さらに、複素数の有理関数に関するいくつかの定理を追加する。複素数変数 $z$ の \emph{有理関数}とは、実数変数 $x$ の有理関数と同様に定義されるもので、すなわち、$z$ の多項式 quotient である。

\par\medskip\noindent\textbf{【定理】}\par
任意の有理関数 $R(z)$ は、$X + Yi$ の形に簡約されることが可能で、$X$ と $Y$ は実係数を持つ $x$ と $y$ の有理関数である。
\par\medskip

まず、任意の多項式 $P(x + yi)$ は、加法と乗法の定義により、$A + Bi$ の形に簡約されることが明らかである。ここで $A$ と $B$ は実係数を持つ $x$ と $y$ の多項式である。同様に $Q(x + yi)$ も $C + Di$ の形に簡約される。したがって
\[
R(x + yi) = \frac{P(x + yi)}{Q(x + yi)}
\]
は以下の形で表すことができる：
\begin{align*}
\frac{A + Bi}{C + Di} &= (A + Bi)(C - Di) / (C + Di)(C - Di) \\
&= \frac{AC + BD}{C^{2} + D^{2}} + \frac{BC - AD}{C^{2} + D^{2}} i,
\end{align*}
これにより定理が証明される。

\par\medskip\noindent\textbf{【定理】}\par
$R(x + yi) = X + Yi$ とし、$R$ は実係数を持つ有理関数とする。このとき $R(x - yi) = X - Yi$ である。
\par\medskip

まず、実係数を持つべき乗 $(x + yi)^n$ の場合、実際の展開により容易に確認できる。加法により、定理は実係数を持つ任意の多項式に対して成り立つことがわかる。したがって、上記の表記を使用すると
\[
R(x - yi) = \frac{A - Bi}{C - Di} = \frac{AC + BD}{C^{2} + D^{2}} - \frac{BC - AD}{C^{2} + D^{2}} i,
\]
と表すことができ、簡約は以前と同じだが $i$ の符号が変更されているだけである。この定理の結果は複素数変数を持つ関数についても同様に成り立つことが明らかである。

\par\medskip\noindent\textbf{【定理】}\par
係数が実数である方程式
\[
a_{0}z^{n} + a_{1}z^{n-1} + \dots + a_{n} = 0,
\]
の根は、自己自身が実数でない限り、共役なペアとして並べることができる。
\par\medskip

これは定理2から、$x + yi$ が根であれば $x - yi$ も根であることが従う。この定理の特別な例として、実係数を持つ二次方程式の根は実数または共役であるという結果（43節）がある。

この定理はしばしば次のように述べられる：「実数係数の方程式において複素根は共役なペアで現れる」。これは、有理数係数の方程式における無理性根が共役するペアで現れるという結果（\Exs{viii}.7）と比較されるべきである。\footnote
  {有理数 $a$, $b$ に対して $a + \sqrt{b}$, $a - \sqrt{b}$ のような数はしばしば「共役」であると言われる。}

例題 XXI.
1. 定義から直接定理（45）の(6)を証明し、幾何学的な考慮事項を使用せずに。

[まず、$|z + z'| \leq |z| + |z'|$ を証明することは $(x + x')^{2} + (y + y')^{2}
  \leq \{\sqrtp{x^{2} + y^{2}} + \sqrtp{x'^{2} + y'^{2}}\}^{2}$ を証明することと等価である。その後、一般の場合に容易に拡張できる.]

2. $|z| + |z'| + \dots = |z + z' + \dots|$ が成り立つ唯一の場合は、数 $z$, $z'$, \dots の全てが同じ振幅を持つときである。幾何学的にも解析的にこれを証明せよ。

3. 任意の複素数の和の絶対値はその実部（または虚部）の総和以上である。

4. 二つの複素数の和と積がどちらも実数である場合、その二つの数は実数か共役であるべきである。

5. $a + b\sqrt{2} + (c + d \sqrt{2})i = A + B\sqrt{2} + (C + D\sqrt{2})i$ において、$a$, $b$, $c$, $d$, $A$, $B$, $C$, $D$ が有理数であるとき、
$a = A,\quad b = B,\quad c = C,\quad d = D$ である。

6. 次の数を $A + Bi$ の形で表せ（$A$ と $B$ は実数）：
$(1 + i)^{2},\quad
\left(\frac{1 + i}{1 - i}\right)^{2},\quad
\left(\frac{1 - i}{1 + i}\right)^{2},\quad
\frac{\lambda + \mu i}{\lambda - \mu i},\quad
  \left(\frac{\lambda + \mu i}{\lambda - \mu i}\right)^{2}
- \left(\frac{\lambda - \mu i}{\lambda + \mu i}\right)^{2}$,
$\lambda$ と $\mu$ が実数である。

\Item{7.} $z = x + yi$ の関数 $z^{2}$, $z^{3}$, $z^{n}$, $1/z$, $z + (1/z)$, $(\alpha + \beta z)/(\gamma + \delta z)$ を $X + Yi$（$X$ と $Y$ は実数）の形で表せ、ここで $\alpha$, $\beta$, $\gamma$, $\delta$ は実数である。

\Item{8.} 前記の例における数値および関数の絶対値を求めよ。

\Item{9.} $z = a$, $z = b$ と $z = c$, $z = d$ の2点間を結ぶ2直線が垂直である条件は
\[
\am\left(\frac{a - b}{c - d}\right) = ±\tfrac{1}{2} \pi,
\]
すなわち $(a - b)/(c - d)$ が純虚数であるときである。これらの直線が平行である条件は何か？

\Item{10.} 三角形の3つの角点は $z = \alpha$, $z = \beta$, $z = \gamma$（$\alpha$, $\beta$, $\gamma$ は複素数）で与えられる。以下の命題を証明せよ：

\SubItem{(i)} 中心 massa（重心）は $z = \frac{1}{3}(\alpha + \beta + \gamma)$ である。

\SubItem{(ii)} 外接円の中心は $|z - \alpha| = |z - \beta| = |z - \gamma|$ を満たす点である。

\SubItem{(iii)} 各角から対辺に引いた垂線が共有点を形成し、その座標は
\[
\Re\left(\frac{z - \alpha}{\beta - \gamma}\right)
  = \Re\left(\frac{z - \beta}{\gamma - \alpha}\right)
  = \Re\left(\frac{z - \gamma}{\alpha - \beta}\right)
  = 0;
\]
を与える。

\SubItem{(iv)} 内部の点 $P$ が存在し、
\[
CBP = ACP = BAP = \omega,
\]
かつ
\[
\cot\omega = \cot A + \cot B + \cot C.
\]

三つ目の条件を証明するために、頂点 \(A\), \(B\), \(C\) と任意の点 \(z\) を考慮します。線分 \(AP\) が線分 \(BC\) に垂直である条件は（練習問題9） \((z - \alpha) / (\beta - \gamma)\) が純虚数であること、すなわち
\[
\Re(z - \alpha) \Re(\beta - \gamma) + \Im(z - \alpha) \Im(\beta - \gamma) = 0
\]
を満たすことになります。この方程式と、\(\alpha\), \(\beta\), \(\gamma\) を循環的に置き換えた2つの類似の方程式は、三つの左辺の和がゼロである事実から、同じ \(z\) の値で満たされます。

四つ目の条件を証明するために、線分 \(BC\) が \(x\) 軸の正方向に平行であるとします。このとき
\[
\gamma - \beta = a, \quad \alpha - \gamma = - b\Cis(-C), \quad \beta - \alpha = - c\Cis B.
\]
を満たします。

\(z\) と \(\omega\) を以下の方程式から決定します。
\[
\frac{(z - \alpha)(\beta_{0} - \alpha_{0})}{(z_{0} - \alpha_{0})(\beta - \alpha)}
= \frac{(z - \beta)(\gamma_{0} - \beta_{0})}{(z_{0} - \beta_{0})(\gamma - \beta)}
= \frac{(z - \gamma)(\alpha_{0} - \gamma_{0})}{(z_{0} - \gamma_{0})(\alpha - \gamma)}
= \Cis 2\omega,
\]
ここで \(z_{0}\), \(\alpha_{0}\), \(\beta_{0}\), \(\gamma_{0}\) は \(z\), \(\alpha\), \(\beta\), \(\gamma\) の共役点を表します。

これらの等しい分数の分子と分母をそれぞれ加え、方程式
\[
i\cot\omega = (1 + \Cis 2\omega)/(1 - \Cis 2\omega)
\]
を使用すると、
\[
i\cot\omega
= \frac{(\beta - \gamma)(\beta_{0} - \gamma_{0}) + (\gamma - \alpha)(\gamma_{0} - \alpha_{0}) + (\alpha - \beta)(\alpha_{0} - \beta_{0})}
{\beta\gamma_{0} - \beta_{0}\gamma + \gamma\alpha_{0} - \gamma_{0}\alpha + \alpha\beta_{0} - \alpha_{0}\beta}
\]
を得ます。この式から、\(\cot\omega\) の値は \((a^{2} + b^{2} + c^{2})/4\Delta\) であることが容易に導け、これは与えられた結果と等価です。

To determine \( z \), 我々は等しい分数の分子と分母をそれぞれ次の式で掛けてから足し合わせて新しい分数を作ります：
\[
(\gamma_{0} - \beta_{0})/(\beta - \alpha),
\]
\[
(\alpha_{0} - \gamma_{0})/(\gamma - \beta),
\]
\[
(\beta_{0} - \alpha_{0})/(\alpha - \gamma)。
\]
新しい分数は次のようになります：
\[
z = \frac{a\alpha \Cis A + b\beta \Cis B + c\gamma \Cis C}{a\Cis A + b\Cis B + c\Cis C}。
\]

11. 顶点が点 \( a \), \( b \), \( c \) と点 \( x \), \( y \), \( z \) のそれぞれの二つの三角形は相似である条件は次の行列式が0であること：
\[
\begin{vmatrix}
1   &   1   &   1\\
a   &   b   &  c \\
x   &   y   &  z
\end{vmatrix}
= 0。
\]

この条件は \(\Seg{AB}/\Seg{AC} = \Seg{XY}/\Seg{XZ}\)（大文字は対応する小文字の点を表す）であることと等価であり、これは与えられた条件と同じです。

12. 例題から導くと、点 \( x \), \( y \), \( z \) が一直線上にある場合、実数 \(\alpha\), \(\beta\), \(\gamma\) を見つけることができます。これらの数は \(\alpha + \beta + \gamma = 0\) および \(\alpha x + \beta y + \gamma z = 0\) を満たし、逆もまた真です（例題 xx.4 参照）。この場合の三角形 \( x \), \( y \), \( z \) は軸 \( OX \) 上の一定の線三角形に相似であり、最後の例題の結果を適用します。

13. 多項式方程式（複素係数）一般形。方程式 \(\alpha z + \beta = 0\) の解は \( z = -(\beta/\alpha) \) で、\(\alpha = 0\) の場合を除きます。ここで
\[
\alpha = a + Ai,
\]
\[
\beta = b + Bi,
\]
\[
z = x + yi，
\]
と置き、実部と虚部を等しくするようにします。これにより2つの方程式が得られ、これらは2つの実数 \( x \) および \( y \) を決定します。方程式が実数の根を持つ条件は \( y = 0 \) で、これは \( ax + b = 0 \), \( Ax + B = 0 \) であり、これらの方程式が一致する条件は \( aB - bA = 0 \) です。

14. 複素係数の一般形の二次方程式。この方程式は次のようになります：
\[
(a + Ai)z^{2} + 2(b + Bi)z + (c + Ci) = 0。
\]

除非 \(a\) 和 \(A\) 同時為零，我達可以を通じて除以 \(a + iA\) 来化简。因此，我達可以考虑以下形式作為方程式的标准形式：
\[
z^2 + 2(b + Bi)z + (c + Ci) = 0 \Tag{(1)}
\]
将 \(z = x + yi\) 代入并分别比较実部和虚部，我達得到关于 \(x\) 和 \(y\) 的一组联立方程式：
\[
x^2 - y^2 + 2(bx - By) + c = 0, \quad
2xy + 2(by + Bx) + C = 0.
\]

设
\[
x + b = \xi, \quad
y + B = \eta, \quad
b^2 - B^2 - c = h, \quad
2bB - C = k,
\]
上述方程式变為：
\[
\xi^2 - \eta^2 = h, \quad
2\xi\eta = k.
\]

平方并相加，我達得到：
\[
\xi^2 + \eta^2 = \sqrt{h^2 + k^2}, \quad
\xi = ±\sqrt{\frac{1}{2}\{\sqrt{h^2 + k^2} + h\}}, \quad
\eta = ±\sqrt{\frac{1}{2}\{\sqrt{h^2 + k^2} - h\}}.
\]

符号的选择应使 \(\xi\eta\) 与 \(k\) 同号：即若 \(k\) 為正，則取同号；若 \(k\) 為負，則取异号。

等根条件。两根相等僅当上述两个平方根同時消失，即 \(h = 0, k = 0\) 或 \(c = b^2 - B^2, C = 2bB\)。これ些条件等价于単个条件 \(c + Ci = (b + Bi)^2\)，明らかに表示方程式 \(\eqref{1}\) 的左辺是一个完全平方。

実根条件。若 \(x^2 + 2(b + Bi)x + (c + Ci) = 0\) 中 \(x\) 為実数，則有 \(x^2 + 2bx + c = 0, 2Bx + C = 0\)。消去 \(x\) 后，我達得到所需条件：
\[
C^2 - 4bBC + 4cB^2 = 0.
\]

純虚根条件。容易験証為：
\[
C^2 - 4bBC - 4b^2c = 0.
\]

\par\medskip\noindent\textbf{共役複素根の条件。}\par\smallskip 二つの共役複素数の和と積はともに実数であるため、$b + Bi$ および $c + Ci$ がともに実数となるためには $B = 0$, $C = 0$ でなければなりません。したがって、方程式（1）が共役複素根のペアを持つのは係数が実数である場合のみです。読者は根の明示的な表現を利用してこの結論を確認してください。さらに、$b^{2} \geq c$ の場合でも根は実数となることがあります。したがって、共役根を持つためには $B = 0$, $C = 0$, $b^{2} < c$ である必要があります。

\Item{15.} \Topic{三次方程式。} 以下の三次方程式を考えます：
\[
z^{3} + 3Hz + G = 0,
\]
ここで $G$ および $H$ は複素数であり、方程式が（a）実根、（b）純虚根、（c）共役複素根のペアを持つことが与えられています。$H = \lambda + \mu i$, $G = \rho + \sigma i$ とすると以下の結論を得ます。

\par\medskip\noindent\textbf{（a） 実根の条件。}\par\smallskip $\mu$ がゼロでない場合、実根は $-\sigma/3\mu$ であり、$\sigma^{3} + 27\lambda\mu^{2}\sigma - 27\mu^{3}\rho = 0$ である。一方、$\mu = 0$ の場合、$\sigma = 0$ も必要で、このとき方程式の係数は実数となる。この場合、三つの実根が存在する可能性がある。

\par\medskip\noindent\textbf{（b） 純虚根の条件。}\par\smallskip $\mu$ がゼロでない場合、純虚根は $(\rho/3\mu)i$ であり、$\rho^{3} - 27\lambda\mu^{2}\rho - 27\mu^{3}\sigma = 0$ である。$\mu = 0$ の場合、$\rho = 0$ も必要で、根は $yi$ となり、$y$ は方程式 $y^{3} - 3\lambda y - \sigma = 0$ の解であり、この方程式の係数は実数である。この場合、三つの純虚根が存在する可能性がある。

\par\medskip\noindent\textbf{（c） 共役複素根のペアの条件。}\par\smallskip このような根を $x + yi$ および $x - yi$ とする。三次方程式の根と係数の関係から、以下の式を得る：
\[
y^{2} - 3x^{2} = 3H,\quad
2x(x^{2} + y^{2}) = G.
\]
したがって $G$ および $H$ はともに実数である。

各場合、方程式の解を求める方法は二つあります。一つ目は根を見つけ出すことです（この場合は既知の因子で割ることにより二次方程式に簡約できます）。二つ目は実数係数を持つ三次方程式の解を求めることです。
\pagesepp{93}

16. 方程式 \(x^{3} + a_{1}x^{2} + a_{2}x + a_{3} = 0\) において、\(a_{1} = A_{1} + A_{1}'i\)、\(\dots\) の場合、解は共役複素数のペアを含みます。残りの根が \(-A_{1}'a_{3}/A_{3}'\) であることを証明し、\(A_{3}' = 0\) の場合はその例外について検討してください。

17. 方程式 \(z^{3} + 3Hz + G = 0\) が複素数の根を持つ場合、方程式
\[
8\alpha^{3} + 6\alpha H - G = 0
\]
は実数の根 \(\alpha\) を持つことを証明し、これは元の方程式の複素数の根の実部であることを示し、その実部 \(\alpha\) の符号が \(G\) の符号と一致することを示してください。

18. 任意の次数を持つ複素係数を持つ方程式は一般に実数の解や共役複素数のペアを持たないでしょう。方程式が (\ia) 実数の根、(\ib) 共役の根を持つために必要な条件は何個ありますか？

19. 交差円について。図26において、\(a\)、\(b\)、\(z\) は \(A\)、\(B\)、\(P\) の主値を用いた振幅であるとします。このとき
\[
\am\frac{z - b}{z - a} = APB,
\]
が成り立ちます。図に示す二つの円が等しい場合、\(z'\)、\(z_{1}\)、\(z_{1}'\) は \(P'\)、\(P_{1}\)、\(P_{1}'\) の振幅であり、かつ \(APB = \theta\) のとき
\[
\am\frac{z' - b}{z' - a} = \pi - \theta,\quad
\am\frac{z_{1} - b}{z_{1} - a} = -\theta,
\]
および
\[
\am\frac{z_{1}' - b}{z_{1}' - a} = -\pi + \theta.
\]
を容易に確認できます。

方程式
\[
\am\frac{z - b}{z - a} = \theta,
\]
が定める軌跡は弧 \(APB\) です。\(\theta\) を定数とし、\(\pi - \theta\)、\(-\theta\)、\(-\pi + \theta\) に置き換えることで、図に示す他の三つの弧を得ることができます。
\%[図の説明: Fig. 26.]
\figure[2.25in]{26}{p093}

方程式の系は、θがパラメータとして考えられ、-πから+πまで変化するとき得られます。これは点A、Bを通る円の系を表します。ただし、各円は異なるθの値に対応する2つの部分に分割されるべきです。

20. 今、方程式
\[
\left|\frac{z - b}{z - a}\right| = \lambda,
\Tag{(1)}
\]
を考えます。ここでλは定数です。

円ABPの点Pにおける接線がABと交わる点をKとします。すると三角形KPA、KBPは相似であり、
\[
AP/PB = PK/BK = KA/KP = \lambda.
\]
となります。したがってKA/KB = λ²，そしてKは方程式(Eq.(1))を満たすすべての点Pに対して固定された点です。またKP² = KA · KBであり、これは定数です。したがって点Pの軌跡は中心がKである円となります。

λを変化させた方程式の系は円の系を表し、この系の各円は例19の系の各円に対して直角で交わります。

例19の系は「共心点型の共軸円の系」と呼ばれ、例20の系は「極点型の共軸円の系」であり、AとBはこの系の極点です。λが非常に大きいまたは小さい場合、円はAまたはBを内部に含む非常に小さな円となります。

\Item{21.} \Topic{双線形变换。} 方程式
\[
z = Z + a,
\Tag{(1)}
\]
を考える。ここで $z = x + yi$ および $Z = X + Yi$ は複素数変数であり、これらを $xoy$ 平面および $XOY$ 平面上で表すことができる。$z$ の任意の値に対応する $Z$ の値が存在し、逆もまた真である。もし $a = \alpha + \beta i$ であれば
\[
x = X + \alpha,\quad y = Y + \beta,
\]
であり、点 $(x, y)$ は点 $(X, Y)$ に対応する。$(x, y)$ の任意の曲線が $xoy$ 平面上を描くならば、$(X, Y)$ は $XOY$ 平面上を描く。したがって、一方の平面上の任意の図形に対応するもう一方の平面上の図形がある。このような関係によって $xoy$ 平面上の図形から $XOY$ 平面上の図形へと移行することを \emph{変換} という。この特定の場合において、対応する図形間の関係は非常に簡単に定義される。$(X, Y)$ の図形は $(x, y)$ の図形と同じ大きさ、形状、および向きであり、しかしながら $\alpha$ の距離だけ左に移動し、$\beta$ の距離だけ下に移動している。このような変換を \emph{平行移動} という。

次に方程式
\[
z = \rho Z,
\Tag{(2)}
\]
を考える。ここで $\rho$ は実数である。これにより $x = \rho X$, $y = \rho Y$ となる。二つの図形は相似であり、それぞれの原点を中心に相似配置されているが、$(x, y)$ の図形のスケールは $(X, Y)$ の図形の $\rho$ 倍である。このような変換を \emph{拡大} という。

最後に方程式
\[
z = (\cos\phi + i \sin\phi)Z.
\Tag{(3)}
\]
を考える。明らかに $|z| = |Z|$ であり、$\am z$ の一つの値は $\am Z + \phi$ であり、二つの図形は $(x, y)$ の図形が原点を中心に正方向に角度 $\phi$ 回転した $(X, Y)$ の図形と異なるだけである。このような変換を \emph{回転} という。

一般線形変換
\[
z = aZ + b \Tag{(4)}
\]
は式(\(1\))、式(\(2\))、式(\(3\))の三つの変換の組合せです。なぜなら、\( |a| = \rho \) かつ \( \arg a = \phi \) のとき、式(\(4\))を以下の三つの等式に置き換えることができます
\[
z = z' + b,\quad
z' = \rho Z',\quad
Z' = (\cos\phi + i\sin\phi)Z.
\]
したがって、一般線形変換は移動、拡大、回転の組合せに相当します。

次に、変換
\[
z = 1/Z \Tag{(5)}
\]
を考慮します。もし \( |Z| = R \) かつ \( \arg Z = \Theta \) のとき、\( |z| = 1/R \) かつ \( \arg z = -\Theta \) であり、\( (x, y) \) の図形から \( (X, Y) \) の図形へ移行するには、単位半径の逆転で \( o \) を中心に前者を逆転させ、その後新しい図形の軸 \( ox \) への像（すなわち \( ox \) の反対側にある対称図形）を作ります。

最後に、変換
\[
z = \frac{aZ + b}{cZ + d} \Tag{(6)}
\]
を考慮します。これは以下の三つの変換の組合せに相当します
\[
z = (a/c) + (bc - ad)(z'/c),\quad
z' = 1/Z',\quad
Z' = cZ + d,
\]
すなわち、既に考慮した種類の変換の特定の組合せに相当します。

変換式(\(6\))は一般双線形変換と呼ばれます。\( Z \) の値に対する \( z \) の値が一対一に対応する最も一般的な種類の変換であり、逆もまた同様です。

解を求めて
\[
Z = \frac{dz - b}{cz - a}.
\]

\par\medskip\noindent\textbf{22.  一般の双線形変換は円を円に変換する。}\par\smallskip
このことは様々な方法で証明できます。純粋幾何学において、反転が円を円に変換すること（もちろん特定の場合には直線になることもある）という有名な定理を仮定することができます。または、例題19と20の結果を使用することができます。例えば、$(x, y)$円が
\[
|(z - \sigma)/(z - \rho)| = \lambda,
\]
であり、$z$を$Z$で表すと、得られるのは
\[
|(Z - \sigma')/(Z - \rho')| = \lambda',
\]
であり、ここで
\[
\sigma' = -\frac{b - \sigma d}{a - \sigma c},\quad
\rho' = -\frac{b - \rho d}{a - \rho c},\quad
\lambda' = \left|\frac{a - \rho c}{a - \sigma c}\right|\lambda.
\]

\Item{23.} $z = 1/Z$, $z = (1 + Z)/(1 - Z)$の変換を考え、$(X, Y)$曲線に対応するものを(1)中心が原点である円、(2)原点を通る直線について描きます。

\Item{24.} 円$x^{2} + y^{2} = 1$を$(X, Y)$平面上の直線に変換するための変換$z = (aZ + b)/(cZ + d)$の条件は$|a| = |c|$である。

\Item{25.} \Topic{交比。} 交比$\DPmod{(z_{1}z_{2}, z_{3}z_{4})}{(z_{1}, z_{2}; z_{3}, z_{4})}$は定義されており、
\[
\frac{(z_{1} - z_{3})(z_{2} - z_{4})}{(z_{1} - z_{4})(z_{2} - z_{3})}
\]
である。

四点$z_{1}$, $z_{2}$, $z_{3}$, $z_{4}$が同一の直線上にある場合、この定義は初等幾何学で採用されたものと一致する。$z_{1}$, $z_{2}$, $z_{3}$, $z_{4}$から形成できる交比は24個あり、これらは6つの等しい交比のグループに分けられる。一つの比が$\lambda$である場合、六つの異なる交比は$\lambda$, $1 - \lambda$, $1/\lambda$, $1/(1 - \lambda)$, $(\lambda - 1)/\lambda$, $\lambda/(\lambda - 1)$である。四点が互いに調和的または調和関係にあるとは、これらのいずれかが$-1$と等しい場合を指す。この場合、六つの比は$-1$, $2$, $-1$, $\frac{1}{2}$, $2$, $\frac{1}{2}$である。

任意の四点比が実数であるとき、すべての四点比は実数であり、その四点は円上に存在する。これは、この場合
\[
\am\frac{(z_{1} - z_{3}) (z_{2} - z_{4})}{(z_{1} - z_{4}) (z_{2} - z_{3})}
\]
が$-\pi$, $0$, $\pi$ のいずれかの値を取ることから、$\am\{(z_{1} - z_{3})/(z_{1} - z_{4})\}$ と
$\am\{(z_{2} - z_{3})/(z_{2} - z_{4})\}$ は等しいか、または$\pi$ の差を持つ（例題19 参照）。

$(z_{1}, z_{2}; z_{3}, z_{4}) = (z_{1}z_{2}, z_{3}z_{4}) / -1$ のとき、以下の二つの方程式が成り立つ：
\[
\am\frac{z_{1} - z_{3}}{z_{1} - z_{4}}
  = ±\pi + \am\frac{z_{2} - z_{3}}{z_{2} - z_{4}},\quad
\left|\frac{z_{1} - z_{3}}{z_{1} - z_{4}}\right|
  = \left|\frac{z_{2} - z_{3}}{z_{2} - z_{4}}\right|.
\]
四点 $A_{1}$, $A_{2}$, $A_{3}$, $A_{4}$ は円上にあり、$A_{1}$ と $A_{2}$ は $A_{3}$ と $A_{4}$ によって分離される。また $A_{1}A_{3}/A_{1}A_{4} = A_{2}A_{3}/A_{2}A_{4}$ である。$A_{3}A_{4}$ の中点を $O$ とする。方程式
\[
\frac{(z_{1} - z_{3}) (z_{2} - z_{4})}{(z_{1} - z_{4}) (z_{2} - z_{3})} = -1
\]
は以下の形に書き換えることができる：
\[
(z_{1} + z_{2}) (z_{3} + z_{4}) = 2(z_{1}z_{2} + z_{3}z_{4}),
\]
これはまた
\[
\{z_{1} - \tfrac{1}{2}(z_{3} + z_{4})\}
\{z_{2} - \tfrac{1}{2}(z_{3} + z_{4})\}
  = \{\tfrac{1}{2}(z_{3} - z_{4})\}^{2}.
\]
これは $\Seg{OA_{1}} · \Seg{OA_{2}} = \Seg{OA_{3}}^{2} = \Seg{OA_{4}}^{2}$ と同値である。したがって $OA_{1}$ と $OA_{2}$ は $A_{3}A_{4}$ に対して等角を成し、$OA_{1} · OA_{2} = OA_{3}^{2} = OA_{4}^{2}$ である。$A_{1}$, $A_{2}$ のペアと $A_{3}$, $A_{4}$ のペアとの間の関係は対称的であることに注意されたい。

したがって、$A_{1}A_{2}$ の中点を $O'$ とするとき、$O'A_{3}$ と $O'A_{4}$ は等角を成し、$O'A_{3} · O'A_{4} = O'A_{1}^{2} = O'A_{2}^{2}$ である。

\item[26.] $A_1, A_2$ の点が $az^2 + 2bz + c = 0$ で与えられ、$A_3, A_4$ の点が $a'z^2 + 2b'z + c' = 0$ で与えられ、$O$ は $A_3A_4$ の中点であり、かつ $ac' + a'c - 2bb' = 0$ のとき、$OA_1, OA_2$ は $A_3A_4$ に対して等角に傾き、さらに $OA_1 \cdot OA_2 = OA_3^2 = OA_4^2$ である。

\item[27.] アルガン図において交わる線分 $AB, CD$ があり、それぞれの中点を $P, Q$ とする。$AB$ が $\angle CPD$ を二等分し、かつ $PA^2 = PB^2 = PC \cdot PD$ のとき、$CD$ は $\angle AQB$ を二等分し、さらに $QC^2 = QD^2 = QA \cdot QB$ であることを証明せよ。

\item[28.] 四点が円上にある条件。十分な条件として、一つ（従って全て）の交差比が実数であること（例題25参照）がある；この条件は必要条件でもある。別の形としては、実数 $\alpha, \beta, \gamma$ を選べるとき、
\[
\begin{vmatrix}
1 & 1 & 1 \\
\alpha & \beta & \gamma \\
z_1 z_4 + z_2 z_3 & z_2 z_4 + z_3 z_1 & z_3 z_4 + z_1 z_2
\end{vmatrix}
= 0
\]
が成り立つことである。

{\Loosen[これを証明するために、$Z = 1/(z - z_4)$ の変換は点 $z_4$ に対する反転と特定の反射を伴うものと等価である（例題21参照）。もし $z_1, z_2, z_3$ が $z_4$ を通る円上にあるとき、対応する点 $Z_1 = 1/(z_1 - z_4)$, $Z_2 = 1/(z_2 - z_4)$, $Z_3 = 1/(z_3 - z_4)$ は一直線上に存在する。従って（例題12参照）、実数 $\alpha', \beta', \gamma'$ を選べるようなものが存在し、$\alpha' + \beta' + \gamma' = 0$ かつ $\alpha'/(z_1 - z_4) + \beta'/(z_2 - z_4) + \gamma'/(z_3 - z_4) = 0$ であり、これは与えられた条件と等価であることを容易に証明できる。]}

\Item{29.} デ・モivreの定理の類似形を実数について証明せよ：正角列 $\phi_{1}$, $\phi_{2}$, $\phi_{3}$, \dots が
\begin{alignat*}{2}
\tan\phi}{m+1} \&= \tan\phi\\textit{{m} \sec\phi}{1} \&\&+ \sec\phi\\textit{{m} \tan\phi}{1},\\
\intertext{であるとき}
\tan\phi\\textit{{m+n} \&= \tan\phi}{m} \sec\phi\\textit{{n} \&\&+ \sec\phi}{m} \tan\phi\\textit{{n},\\
\sec\phi}{m+n} \&= \sec\phi\\textit{{m} \sec\phi}{n} \&\&+ \tan\phi\\textit{{m} \tan\phi}{n},
\end{alignat*}
および
\[
\tan\phi_{m} + \sec\phi_{m} = (\tan\phi_{1} + \sec\phi_{1})^{m}.
\]
が成り立つことを証明せよ。[数学的帰納法を使用する]

\Item{30.} $z = Z^{m}$ の変換について．この場合，$r = R^{m}$ であり，$\theta$ と $m\Theta$ は $2\pi$ の倍数で異なる．$Z$ が原点周りの円を描くとき，$z$ は原点を中心に $m$ 回描く．
$(x, y)$ 平面全体は $(X, Y)$ 平面上の任意の $m$ 個の扇形に対応し，各扇形は角度 $2\pi/m$ を持つ．$(x, y)$ 平面上の各点は $(X, Y)$ 平面上の $m$ 点に対応する．

\Item{31.} 実変数の複素関数について．実変数 $t$ に対する二つの実関数 $f(t), \phi(t)$ が一定範囲の $t$ の値に対して定義されているとき，式
\[
z = f(t) + i\phi(t)
\Tag{(1)}
\]
を $t$ の複素関数と呼ぶ．この関数はグラフ的に曲線
\[
x = f(t),\quad y = \phi(t);
\]
によって表すことができる．この曲線の方程式はこれらの式から $t$ を消去することにより得られる．$z$ が $t$ の複素係数を持つ多項式または有理関数であるとき，この関数が表す曲線を決定するために式 \Eq{(1)} の形に表すことができる．

\SubItem{(i)} \( z = a + (b - a)t \)において、\( a \)と\( b \)は複素数である。もし \( a = \alpha + \alpha' i \), \( b = \beta + \beta' i \)ならば
\[ x = \alpha  + (\beta  - \alpha)t,\quad y = \alpha' + (\beta' - \alpha')t. \]
この曲線は複素数\( z = a \)と\( z = b \)の点を結ぶ直線である。これらの点間の区間は、\( t \)の値が0から1までの範囲に対応する。

\SubItem{(ii)} \( z = c + \rho\left(\frac{1 + ti}{1 - ti}\right) \)において、\(\rho\)は正であるとき、この曲線は中心\( c \)、半径\(\rho\)の円である。\( t \)がすべての実数を経過するとき、\( z \)は一度だけ円を描く。

\SubItem{(iii)} 一般に方程式 \( z = \frac{a + bt}{c + dt} \)は円を表す。\( x \)と\( y \)を求め、それを消去することで証明できるが、この過程はやや煩雑である。より簡便な方法は、例題22の結果を使用することである。\( z = \frac{a + bZ}{c + dZ} \), \( Z = t \)と置く。\( t \)が変化するとき、\( Z \)は\( X \)軸を描く直線を描く。したがって、\( z \)は円を描く。

\SubItem{(iv)} 方程式
\[ z = a + 2bt + ct^{2} \]
は一般に放物線を表し、もし \( b/c \)が実数であるならば直線を表す。

\SubItem{(v)} 方程式 \( z = \frac{a + 2bt + ct^{2}}{\alpha + 2\beta t + \gamma t^{2}} \), \(\alpha\)、\(\beta\)、\(\gamma\)が実数であるとき、これは楕円、双曲線、放物線のいずれかを表す。

[\( t \)から eliminare
\[ x = \frac{A  + 2Bt  + Ct^{2}}{\alpha  + 2\beta t + \gamma t^{2}},\quad y = \frac{A' + 2B't + C't^{2}}{\alpha + 2\beta t + \gamma t^{2}},
\]
ここで \( A + A'i = a \), \( B + B'i = b \), \( C + C'i = c \).]

47. 複素数の根。これまで、複素数 \(a\) に対して \( \sqrt[n]{a} \) や \( a^{m/n} \) の意味を付与していませんでした。しかし、実数の場合に与えられる初等代数学における定義を採用するのが自然です。したがって、\( n \) が正の整数であるとき、\( \sqrt[n]{a} \) や \( a^{1/n} \) を、方程式 \( z^n = a \) を満たす数 \( z \) として定義します。また、\( m \) が整数であるとき、\( a^{m/n} \) を \( (a^{1/n})^m \) と定義します。これらの定義は方程式の根が一つ以上存在するかどうかという問題を先取りしていないことを注意してください。

48. 方程式 \( z^n = a \) の解。まず
\[
a = \rho(\cos\phi + i\sin\phi),
\]
ここで、\( \rho \) は正の数で、\( \phi \) は \( -\pi < \phi \leq \pi \) を満たす角度です。次に
\[
z = r(\cos\theta + i\sin\theta)
\]
と置くと、方程式は次の形になります：
\[
r^n (\cos n\theta + i\sin n\theta) = \rho(\cos\phi + i\sin\phi);
\]
したがって
\[
r^n = \rho, \quad \cos n\theta = \cos\phi, \quad \sin n\theta = \sin\phi.
\Tag{(1)}
\]

\( r \) の唯一の可能性は \( \sqrt[n]{\rho} \) であり、これは通常の算術的な \( n \) 乗根です。最後の二つの方程式を満たすためには
\[
n\theta = \phi + 2k\pi, \quad (k \text{は整数})
\]
が必要で十分です。したがって
\[
\theta = (\phi + 2k\pi)/n.
\]
ここで \( k = pn + q \)（\( p \) と \( q \) は整数）と置くと、\( \theta \) の値は \( 2p\pi + (\phi + 2q\pi)/n \) であり、\( p \) の値は関係ありません。したがって
\[
z^n = a = \rho(\cos\phi + i\sin\phi)
\]
の解は \( n \) 個で、これらは
\[
z = r(\cos\theta + i\sin\theta), \quad r = \sqrt[n]{\rho}, \quad \theta = (\phi + 2q\pi)/n, \quad (q = 0, 1, 2, \dots, n - 1)
\]
で与えられます。

これらの \( n \) つの根が実際にはすべて異なることは、Argand の図に描くことで容易に確認できる。

特別な根
\[
\sqrt[n]{\rho}\{\cos(\phi/n) + i\sin(\phi/n)\}
\]
は、\(\sqrt[n]{a}\) の主値と呼ばれる。

\( a = 1 \), \( \rho = 1 \), \( \phi = 0 \) の場合について特に興味深い。方程式 \( x^{n} = 1 \) の \( n \) つの根は
\[
\cos(2q\pi/n) + i\sin(2q\pi/n),\quad (q = 0,\ 1,\ \dots, n - 1)
\]
である。これらの数は単位根と呼ばれ、主値は自身で単位である。\(\omega_{n} = \cos(2\pi/n) + i\sin(2\pi/n)\) とおくと、単位根は
\[
1,\quad \omega_{n},\quad \omega_{n}^{2},\ \dots,\quad \omega_{n}^{n-1}
\]
である。

例題 XXII.
1. 1 の平方根は 1, -1; 3 の立方根は 1, \( \frac{1}{2}(-1 + i\sqrt{3}) \), \( \frac{1}{2}(-1 - i\sqrt{3}) \); 4 の四乗根は 1, \( i \), -1, -\( i \); 5 の五乗根は
\[
1,\quad \tfrac{1}{4}\Bigl[\sqrt{5} - 1 + i\sqrt{10 + 2\sqrt{5}}\Bigr],\quad \tfrac{1}{4}\Bigl[-\sqrt{5} - 1 + i\sqrt{10 - 2\sqrt{5}}\Bigr],
\]
\[
\tfrac{1}{4}\Bigl[-\sqrt{5} - 1 - i\sqrt{10 - 2\sqrt{5}}\Bigr],\quad \tfrac{1}{4}\Bigl[\sqrt{5} - 1 - i\sqrt{10 + 2\sqrt{5}}\Bigr].
\]

2. \( 1 + \omega_{n} + \omega_{n}^{2} + \dots + \omega_{n}^{n-1} = 0 \) を証明せよ。

3. \( (x + y\omega_{3} + z\omega_{3}^{2})(x + y\omega_{3}^{2} + z\omega_{3}) = x^{2} + y^{2} + z^{2} - yz - zx - xy \) を証明せよ。

4. \( a \) の \( n \) 乗根は単位根と主値の \( n \) 乗根の積である。

5. \(z^{2} = \alpha + \beta i\) の根は以下のようになる：
\[
± \sqrtbr{\tfrac{1}{2} \{\sqrtp{\alpha^{2} + \beta^{2}} + \alpha\}}
± i\sqrtbr{\tfrac{1}{2} \{\sqrtp{\alpha^{2} + \beta^{2}} - \alpha\}},
\]
\(\beta\) が正のときには符号を同じにし、負のときには符号を異にする。この結果は第48節の結果と一致することを示せ。

6. \((x^{2m} - a^{2m})/(x^{2} - a^{2})\) は以下のようになる：
\[
\Bigl(x^{2} - 2ax\cos\frac{\pi}{m} + a^{2}\Bigr)
\Bigl(x^{2} - 2ax\cos\frac{2\pi}{m} + a^{2}\Bigr) \dots
\Bigl(x^{2} - 2ax\cos\frac{(m - 1)\pi}{m} + a^{2}\Bigr).
\]

\[x^{2m} - a^{2m}\] の因子は以下のようになる：
\[
(x - a),\quad (x - a\omega_{2m}),\quad (x - a\omega_{2m}^{2}), \dots, (x - a\omega_{2m}^{2m-1}).
\]
因子 \(x - a\omega_{2m}^{m}\) は \(x + a\) である。因子 \((x - a\omega_{2m}^{s})\) と \((x - a\omega_{2m}^{2m-s})\) を合わせると、因子 \(x^{2} - 2ax \cos(s\pi/m) + a^{2}\) となる。

7. \(x^{2m+1} - a^{2m+1}\), \(x^{2m} + a^{2m}\), および \(x^{2m+1} + a^{2m+1}\) を同様の方法で因子分解せよ。

8. \(x^{2n} - 2x^{n}a^{n} \cos\theta + a^{2n}\) は以下のようになる：
\begin{multline*}
\left(x^{2} - 2xa\cos\frac{\theta}{n} + a^{2}\right)
\left(x^{2} - 2xa\cos\frac{\theta + 2\pi}{n} + a^{2}\right) \dots \\
\dots\left(x^{2} - 2xa\cos\frac{\theta + 2(n - 1)\pi}{n} + a^{2}\right).
\end{multline*}

\[x^{2n} - 2x^{n}a^{n} \cos\theta + a^{2n}
  = \{x^{n} - a^{n}(\cos\theta + i\sin\theta)\}
    \{x^{n} - a^{n}(\cos\theta - i\sin\theta)\},
\]
最後の２つの式をそれぞれ \(n\) 個の因子に分解せよ。

9. 方程式 \(x^{6} - 2x^{3} + 2 = 0\) のすべての根を求めよ。

\Item{10.} 数値形式で平方根のみを含む形で正確な$\omega_{n}$の値を求めることは、単位半径の円内にユークリッド方法（直尺とコンパス）を使用して正規多角形$n$辺を内接させる幾何学的問題の代数的な同値です。この構成は可能であるのはもしこれらの長さ$\cos(2\pi/n)$および$\sin(2\pi/n)$が構築可能である場合に限ります。そして、これは可能である（第II章、補題演習II 22）のはもしこれらの数が平方根のみを含む形で表記可能である場合に限ります。

ユークリッドは$n = 3, 4, 5, 6, 8, 10, 12,$および$15$に対する構成法を与えています。明らかに、これらのいずれかの値から任意の2の冪を乗じた値に対してこの構成法は可能であることが明白です。他の特殊な$n$の値も存在し、最も興味深いものは$n = 17$です。

\par\medskip
\ページ分割{101}

\節{デ・モivre の定理の一般形}．前節の結果から、$q$が正の整数であるとき、$(\cos\theta + i\sin\theta)^{1/q}$の一つの値は
\[
\cos(\theta/q) + i\sin(\theta/q)
\]
となります。これらの表現を任意の整数（正または負）$p$で累乗すると、$(\cos\theta + i\sin\theta)^{p/q}$の一つの値が$\cos(p\theta/q) + i\sin(p\theta/q)$であることが得られ、すなわち、任意の有理数$\alpha$に対して、$(\cos\theta + i\sin\theta)^{\alpha}$の一つの値が
\[
\cos\alpha\theta + i\sin\alpha\theta
\]
であることが得られます。これはデ・モivre の定理の一般化形（第45節）です。

\章III に関する補足例題

\begin{Examples}{}
\Item{1.} 三角形$(xyz)$が正三角形である条件は
\[
x^{2} + y^{2} + z^{2} - yz - zx - xy = 0
\]
です。
\par\medskip

{\Loosen[三角形 XYZ において、移動 \Seg{ZX} は \Seg{YZ} を正または負の方向で角度 \(\frac{2}{3}\pi\) 回転させたものである。Since \(\Cis\frac{2}{3}\pi = \omega_{3}\), \(\Cis(-\frac{2}{3}\pi) = 1/\omega_{3} = \omega_{3}^{2}\)，従って \(x - z = (z - y)\omega_{3}\) や \(x - z = (z - y)\omega_{3}^{2}\) である。したがって
\(x + y\omega_{3} + z\omega_{3}^{2} = 0\) や \(x + y\omega_{3}^{2} + z\omega_{3} = 0\) である。結果は \Exs{xxii}.~3. から得られる。]}

\Item{2.} 三角形 XYZ, X'Y'Z' が二つの三角形であり、
\[
\Seg{YZ} · \Seg{Y'Z'} = \Seg{ZX} · \Seg{Z'X'} = \Seg{XY} · \Seg{X'Y'},
\]
であるとき、両方の三角形は等辺である。[方程式
\[
(y - z)(y' - z') = (z - x)(z' - x') = (x - y)(x' - y') = \kappa^{2},
\]
から \(\sum 1/(y' - z') = 0\)，または \(\sum x'^{2} - \sum y'z' = 0\) を導き、最後の例の結果を適用する。]

\Item{3.} 三角形 ABC の辺上に相似な三角形 BCX, CAY, ABZ を描く。ABC，XYZ の重心は一致することを示せ。

[私たちは \((x - c)/(b - c) = (y - a)/(c - a) = (z - b)/(a - b) = \lambda\) と置くことができる。\(\frac{1}{3}(x + y + z)\) を a, b, c で表す。]

\Item{4.} 三角形 ABC の辺上に点 X, Y, Z があり、
\[
BX/XC = CY/YA = AZ/ZB = r,
\]
であり、ABC と XYZ が相似であるとき、r = 1 または両方の三角形は等辺である。

\Item{5.} 平面上に四点 A, B, C, D があるとき、
\[
AD · BC \leq BD · CA + CD · AB.
\]
]

[複素数 \(z_{1}\), \(z_{2}\), \(z_{3}\), \(z_{4}\) がそれぞれ A, B, C, D を表すと、恒等的に
\[
(x_{1} - x_{4})(x_{2} - x_{3}) +
(x_{2} - x_{4})(x_{3} - x_{1}) +
(x_{3} - x_{4})(x_{1} - x_{2}) = 0.
\]
である。従って
\begin{align*}
|(x_{1} - x_{4})(x_{2} - x_{3})|
  &=    |(x_{2} - x_{4})(x_{3} - x_{1}) + (x_{3} - x_{4})(x_{1} - x_{2})|\\
  &\leq |(x_{2} - x_{4})(x_{3} - x_{1})| + |(x_{3} - x_{4})(x_{1} - x_{2})|.
\end{align*}]

\item{6.} 循環四角形に関するポタメイの定理を、4つの円周上にある点の交叉比が実数であることから導き出せ。［同じ恒等式を使用する］

\item{7.} $z^{2} + z'^{2} = 1$ のとき、点 $z$,~$z'$ は焦点が点 $1$,~$-1$ である楕円の共役直径の端点である。［もし $CP$,~$CD$ が楕円の共役半径であり、$S$,~$H$ がその焦点であるとき、$CD$ は $\angle SPH$ の外角の二等分線に平行であり、$SP · HP = CD^{2}$ である］

\item{8.} $|a + b|^{2} + |a - b|^{2} = 2\{|a|^{2} + |b|^{2}\}$ を証明せよ。［これは、$M$ が $PQ$ の中点であるときの幾何学的定理の解析的な相当形であり、$OP^{2} + OQ^{2} = 2OM^{2} + 2MP^{2}$ である］

\item{9.} \item{8} より
\[
|a + \sqrtp{a^{2} - b^{2}}| + |a - \sqrtp{a^{2} - b^{2}}| = |a + b| + |a - b|.
\]

を導き出せ。［もし $a + \sqrtp{a^{2} - b^{2}} = z_{1}$, $a - \sqrtp{a^{2} - b^{2}} = z_{2}$ であるとき、
\[
|z_{1}|^{2} + |z_{2}|^{2}
  = \tfrac{1}{2}|z_{1} + z_{2}|^{2}
  + \tfrac{1}{2}|z_{1} - z_{2}|^{2}
  = 2|a|^{2} + 2|a^{2} - b^{2}|,
\]
であるから
\[
(|z_{1}| + |z_{2}|)^{2}
  = 2\{|a|^{2} + |a^{2} - b^{2}| + |b|^{2}\}
  = |a + b|^{2} + |a - b|^{2} + 2|a^{2} - b^{2}|.
\]

別の言い方をすると、$z_{1}$ と $z_{2}$ が $\alpha z^{2} + 2\beta z + \gamma = 0$ の根であるとき
\[
|z_{1}| + |z_{2}| = (1/|\alpha|)
     \{(|-\beta + \sqrtp{\alpha\gamma})|
     + (|-\beta - \sqrtp{\alpha\gamma})\}.
\]

\item{10.} 方程式 $z^{2} + az + b = 0$ の両方の根が単位円周上にある必要十分条件は
\[
|a| \leq 2,\quad
|b| = 1,\quad
\am b = 2\am a.
\]

である。［振幅は必ずしも主値ではない］

\Item{11.} \(x^{4} + 4a_{1}x^{3} + 6a_{2}x^{2} + 4a_{3}x + a_{4} = 0\)が実数係数を持つ方程式で、その解が Argand 平面上で円周上に並ぶとき、
\[
a_{3}^{2} + a_{1}^{2}a_{4} + a_{2}^{3} - a_{2}a_{4} - 2a_{1}a_{2}a_{3} = 0.
\]

\Item{12.} \(a_{0}x^{4} + 4a_{1}x^{3} + 6a_{2}x^{2} + 4a_{3}x + a_{4} = 0\)の四つの解が調和関係にあるとき、
\[
a_{0}a_{3}^{2} + a_{1}^{2}a_{4} + a_{2}^{3}
  - a_{0}a_{2}a_{4} - 2a_{1}a_{2}a_{3} = 0.
\]

[式を展開し、\(Z_{23, 14} Z_{31, 24} Z_{12, 34}\) を係数の形で表す。ここで \(Z_{23, 14} = (z_{1} - z_{2}) (z_{3} - z_{4}) + (z_{1} - z_{3}) (z_{2} - z_{4})\) であり、\(z_{1}\), \(z_{2}\), \(z_{3}\), \(z_{4}\)は方程式の解である。]

\Item{13.} \Topic{虚点と直線。} \(ax + by + c = 0\)が複素係数を持つ方程式（もちろん、特殊な場合において実数となることもあります）。

\(x\)に任意の実数または複素数値を与えると、対応する \(y\) の値を求めることができます。方程式を満たす \(x\) および \(y\) の実数または複素数値のペア全体が虚直線と呼ばれ、これらの値は虚点と呼ばれ、その点は直線上にあると言います。\(x\) および \(y\) の値は点 \((x, y)\) の座標と呼ばれます。\(x\) および \(y\) が実数であるとき、その点は実点と呼ばれ、\(a\), \(b\), \(c\) がすべて実数（または共通の因数で割ることによって実数にできる）であるとき、直線は実直線と呼ばれます。点 \(x = \alpha + \beta i\), \(y = \gamma + \delta i\) および \(x = \alpha - \beta i\), \(y = \gamma - \delta i\) は共役であり、直線も同様に
\[
(A + A'i)x + (B + B'i)y + C + C'i = 0,\quad
(A - A'i)x + (B - B'i)y + C - C'i = 0.
\]

---任意の実直線上には無数の共役虚点のペアが存在する；一般に虚直線には一つだけ実点が存在し、虚直線には共役虚点のペアは存在しない：---そして、二つの与えられた虚点を結ぶ直線が実直線である条件（ia）と、二つの虚直線の交点が実点である条件（ib）を求めよ。

14. 次の恒等式を証明せよ：
\begin{gather*}
(x + y + z) (x + y\omega_{3} + z\omega_{3}^{2})
            (x + y\omega_{3}^{2} + z\omega_{3})
  = x^{3} + y^{3} + z^{3} - 3xyz,\\
(x + y + z) (x + y\omega_{5} + z\omega_{5}^{4})
            (x + y\omega_{5}^{2} + z\omega_{5}^{3})
            (x + y\omega_{5}^{3} + z\omega_{5}^{2})
            (x + y\omega_{5}^{4} + z\omega_{5})\\
= x^{5} + y^{5} + z^{5} - 5x^{3}yz + 5xy^{2}z^{2}.
\end{gather*}

15. 次の方程式を解け：
\[
x^{3} - 3ax + (a^{3} + 1) = 0,\quad
x^{5} - 5ax^{3} + 5a^{2}x + (a^{5} + 1) = 0.
\]

16. $f(x) = a_{0} + a_{1}x + \dots + a_{k}x^{k}$ のとき、
\[
\{f(x) + f(\omega x) + \dots + f(\omega^{n-1}x)\}/n
  = a_{0} + a_{n}x^{n} + a_{2n}x^{2n} + \dots + a_{\lambda n}x^{\lambda n},
\]
$\omega$ は $x^{n} = 1$ の任意の根（$x = 1$ を除く）、および $\lambda n$ は $k$ に含まれる最大の $n$ の倍数である。同様の公式を $a_{\mu} + a_{\mu+n}x^{n} + a_{\mu+2n}x^{2n} + \dots$ の形で求めよ。

17. 次の式において、$n$ は正の整数とする：
\[
(1 + x)^{n} = p_{0} + p_{1}x + p_{2}x^{2} + \dots,
\]
このとき
\[
p_{0} - p_{2} + p_{4} - \dots = 2^{\frac{1}{2} n} \cos\tfrac{1}{4}n\pi,\quad
p_{1} - p_{3} + p_{5} - \dots = 2^{\frac{1}{2} n} \sin\tfrac{1}{4}n\pi.
\]

18. 次の級数を求和せよ：
\[
\frac{x}{2! \DPchg{n - 2!}{(n - 2)!}}
  + \frac{x^{2}}{5! \DPchg{n - 5!}{(n - 5)!}}
  + \frac{x^{3}}{8! \DPchg{n - 8!}{(n - 8)!}} + \dots
  + \frac{x^{n/3}}{\DPchg{n - 1!}{(n - 1)!}},
\]
$n$ は $3$ の倍数である。

\item[19.] もし複素数$t$が$|t|=1$を満たすとき、点$x=(at+b)/(t-c)$は$t$が変化するにつれて円を描きます。ただし、$|c|=1$の場合は直線を描きます。

\item[20.] 例題と同様に$t$が変化すると、点$x=\frac{1}{2}\{at+(b/t)\}$は一般に楕円を描き、その焦点は$x^2=ab$で与えられ、軸の長さは$|a|+|b|$と$|a|-|b|$です。ただし、$|a|=|b|$の場合は有限な直線$-\sqrt{ab}$から$\sqrt{ab}$までを結ぶ点$x$が描かれます。

\item[21.] $t$が実数であり、$z=t^2-1+\sqrt{t^4-t^2}$のとき、$t^2<1$の場合、$z$は円$x^2+y^2+x=0$上の点で表されます。また、$t^2>1$の場合は$\sqrt{t^4-t^2}$が$t^4-t^2$の正の平方根を表すと仮定し、$t$が大きな正の値から大きな負の値へ減少するときの点$z$の動きについて議論してください。（1912年）

\item[22.] 転換$z=(aZ+b)/(cZ+d)$の係数は条件$ad-bc=1$を満たします。もし$c \neq 0$の場合、この転換には二つの固定点$\alpha$, $\beta$が存在し、$(a+d)^2=4$の場合は一つの固定点$\alpha$のみが存在します。これらの二つの場合では、転換は次の形で表すことができます：
\[
\frac{z-\alpha}{z-\beta}=K\frac{Z-\alpha}{Z-\beta}, \quad
\frac{1}{z-\alpha}=\frac{1}{Z-\alpha}+K.
\]

さらに$c=0$の場合、固定点$\alpha$が存在し、$a=d$の場合は一つの固定点$\alpha$のみが存在します。これらの二つの場合では、転換は次の形で表すことができます：
\[
z-\alpha=K(Z-\alpha), \quad
z=Z+K.
\]

最後に、$a$, $b$, $c$, $d$がさらに正の整数値（零を含む）に制限される場合、固定点が二つ未満である唯一の転換は$(1/z)=(1/Z)+K$, $z=Z+K$の形です。（1911年）

23. \( z = (1 + Zi)/(Z + i) \)という関係式が、\( x \)-軸上の点 \( z = 1 \)と\( z = -1 \)の部分を半円に変換することを証明せよ。その半円は点 \( Z = 1 \)と\( Z = -1 \)を通る。さらに、この変換を繰り返し適用することで得られる図形すべてを求めよ。\MathTrip{1912}。

24. \( z = 2Z + Z^{2} \)のとき、円 \( |Z| = 1 \)は\( z \)-平面上で心臓型曲線（カーディオイド）に対応する。

25. \( z = \frac{1}{2}\{Z + (1/Z)\} \)という変換について議論し、特に円 \( X^{2} + Y^{2} = \alpha^{2} \)が共役楕円（コンファコール・エルリプティクス）に変換されることを示せ。その楕円は
\[
\frac{x^{2}}{\left\{\dfrac{1}{2}\left(\alpha + \dfrac{1}{\alpha}\right)\right\}^{2}}
  + \frac{y^{2}}{\left\{\dfrac{1}{2}\left(\alpha - \dfrac{1}{\alpha}\right)\right\}^{2}}
  = 1
\]
である。

26. \( (z + 1)^{2} = 4/Z \)のとき、\( z \)-平面上の単位円は\( Z \)-平面上の放物線 \( R\cos^{2} \frac{1}{2}\Theta = 1 \)に対応し、単位円内の部分は放物線の外側に対応する。

27. \( z = \{(Z - ci)/(Z + ci)\}^{2} \)という変換により、\( z \)-平面上の上半平面は特定の半円の内部に変換されることを示せ。

28. \( z = Z^{2} - 1 \)のとき、\( z \)が円 \( |z| = \kappa \)を描くと、対応する位置の\( Z \)はそれぞれカッシニアン卵形（カッシニオイド） \( \rho_{1}\rho_{2} = \kappa \)を描く。ここで、\(\rho_{1}\), \(\rho_{2}\)は点 \( -1 \), \( 1 \)から\( Z \)までの距離である。異なる\(\kappa\)の値に対して卵形を追跡せよ。

29. 次の関係式 \( az^{2} + 2hzZ + bZ^{2} + 2gz + 2fZ + c = 0 \)について考える。\( z \)に対応する二つの値の\( Z \)が存在し、その逆もまた真であることを示せ。これらをそれぞれ \( Z \)-平面と \( z \)-平面の分岐点という。さらに、分岐点が焦点となる楕円上を描く \( z \)が同様に描かれることが示されよ。

私たちは、与えられた関係を以下のように一般性を失うことなく取ります：
\[
z^{2} + 2zZ\cos\omega + Z^{2} = 1:
\]
読者はこの形が成り立つことを確認してください。どちらの平面においても、分岐点は$\cosec\omega$と$-\cosec\omega$です。指定された形式の楕円は以下の式で与えられます：
\[
|z + \cosec\omega| + |z - \cosec\omega| = C,
\]
ここで、$C$は定数です。これは（例題9）以下に等価です：
\[
|z + \sqrt{z^{2} - \cosec^{2}\omega}| +
|z - \sqrt{z^{2} - \cosec^{2}\omega}| = C.
\]
この式を$Z$で表すと？

30. 問題。もし$z = aZ^{m} + bZ^{n}$であり、$m$、$n$は正の整数、$a$、$b$は実数であるとき、$Z$が単位円を描くにつれて、$z$は擬円または外円を描きます。

31. 次の変換を示せ：
\[
z = \frac{(a + di)Z_{0} + b}{cZ_{0} - (a - di)},
\]
ここで$a$、$b$、$c$、$d$は実数であり、$a^{2} + d^{2} + bc > 0$、そして$Z_{0}$は$Z$の共役を表します。この変換は円
\[
c(x^{2} + y^{2}) - 2ax - 2dy - b = 0
\]
に対する反転に等価です。
$a^{2} + d^{2} + bc < 0$のとき、この変換の幾何学的な解釈は何ですか？

32. 次の変換を示せ：
\[
\frac{1 - z}{1 + z} = \left(\frac{1 - Z}{1 + Z}\right)^{c},
\]
ここで$c$は有理数であり、$0 < c < 1$です。この変換により、円$|z| = 1$は角度$\pi/c$の円形月齢の境界に変換されます。
例題終了
ページ分割：106

\chapter{第4章：Chapter IV: LIMITS OF FUNCTIONS OF A POSITIVE INTEGRAL VARIABLE}

第4章 初期整数変数の関数の極限

50. 正整数変数の関数。第2章では、実数変数 \( x \) の関数の概念について議論し、そのような関数の多くの例を挙げました。そして読者には記憶に留めていただきたいが、これらの関数は重要な点において非常に異なるところがあった。一部の関数はすべての \( x \) の値で定義されており、他のものは有理数のみで定義され、さらに他のものは整数のみで定義されていたなどである。

例えば、以下の関数を考慮する：（i）\( x \)，（ii）\(\sqrt{x}\)，（iii）\( x \) の分母，（iv）\( x \) の分子と分母の積の平方根，（v）\( x \) の最大素因数，（vi）\(\sqrt{x} \) と \( x \) の最大素因数の積，（vii）\( x \) 番目の素数，（viii）ダートモア刑務所の囚人 \( x \) の身長（インチ単位）。

次に、これらの関数が定義される$x$の値の集合、すなわち関数の定義域は以下の通りである：
(i) 全ての$x$の値,
(ii) すべての正の$x$の値,
(iii) すべての有理数$x$の値,
(iv) すべての正の有理数$x$の値,
(v) すべての整数$x$の値,
(vi), (vii) すべての正の整数$x$の値,
(viii) 特定のいくつかの正の整数$x$の値、すなわち1, 2, …, N、ここでNはダートモア刑務所における特定の時間点での囚人の総数である。

最後の場合においてNは時間に依存し、xが一定の値を持つ囚人は異なる時間点で異なる個人となる。したがって、異なる時間点を考慮すると、ダートモア刑務所の設置から廃止までの特定の範囲のtの値と、特定の数の正の整数x（この数はtにより変化する）に対して定義された2つの変数y = F(x, t) の単純な例となる。

今後、上記(vii)のような関数について考えよう。この関数はすべての正の整数xに対してのみ定義され、それ以外の値では定義が失われるものとする。ページ分割点107
関数は、これまでの慣習として、実数変数xの一部の値に対してのみ定義された関数と見做すことができる。すべてのその他のxの値では定義が失われるとする。あるいは、xの値が正の整数以外である場合を考慮に入れないで、この関数を正の整数変数nの関数として考えることができる。ここでnの値は正の整数
\[
1,\ 2,\ 3,\ 4,\ \dots.
\]
とされる。

この場合、
\[
y = \phi(n)
\]
と書き、yをnのすべての値に対して定義された関数として考える。

任意の$x$に対する関数がすべての$x$の値に対して定義されているならば、その関数は$n$に対する関数として定義される。したがって、関数$y=x^{2}$から関数$y=n^{2}$を導き出すには、考慮する$x$の値を正整数以外すべて省略し、対応する$y$の値も省略すればよい。一方で、任意の$n$に対する関数から、$x$に対する何らかの関数を複数導き出すには、$x$の非正整数値に対応する$y$の値を任意に定めればよい。

以下は補足であるが、前置きや解説は含まない：

問題は、与えられた$n$に対する関数のすべての正整数値に対し、その値と一致する$x$に対する関数を求めることであり、これは上位数学において極めて重要である。この問題は、関数間の補間問題と呼ばれる。

しかし、ただ条件を満たす$x$に対する任意の関数を見つけるだけであれば、当然何の困難も生じない。先述のように、欠けている値を任意に埋めればよい。実際には与えられた$n$に対する関数の値すべてを$x$に対する関数の定義と見做し、その他の$x$の値では定義が失われるというような純粋理論的な解法でも十分であるが、通常求められるのは、与えられた値$x=1,2,\dots$においてその値を取るような、$x$に依存する式（可能な限り単純な形）である。

有些場合、特に関数が \(n\) によって定義されるときには明らかな解がある。例えば \(y = \phi(n)\) のとき、\(\phi(n)\) は \(n^2\) や \(\cos n\pi\) のような \(n\) の関数であり、\(n\) が正の整数でなくても意味を持つものであれば、自然に \(x\) の関数を \(y = \phi(x)\) と定める。

しかし、この非常に単純な場合でも、問題の他のほぼ同様に明らかな解を書くことは容易である。例えば
\[
y = \phi(x) + \sin x\pi
\]
は \(x = n\) のとき \(\phi(n)\) となる値を持つが、これは \(\sin n\pi = 0\) よりである。

他には \(\phi(n)\) が公式によって定義される場合がある。例えば \((-1)^n\) は某些の \(x\) の値に対して定義されない（ここでは分数で偶数分母を持つ \(x\) の場合や無理数の場合など）が、その公式を変形することで全ての \(x\) の値に対して定義されるようにすることができるかもしれない。この場合、例えば
\[
(-1)^n = \cos n\pi
\]
が整数 \(n\) に対して成り立つとき、補間問題は \(\cos x\pi\) の関数によって解決される。

他には \(\phi(x)\) が正の整数以外の某些の \(x\) の値で定義され、全ての \(x\) の値で定義されない場合がある。例えば \(y = n^n\) より \(y = x^x\) となる。この表現は残りの某些の \(x\) の値に対して意味を持つ。単純化のために正の \(x\) のみを考えるとき、\(x^x\) は基本的な代数で採用された分数冪の定義により全ての有理数 \(x\) に対して意味を持つ。しかし、\(x\) が無理性であるとき、\(x^x\) は（現在のところ）全く意味を持たない。したがってこの場合、補間問題はすぐに我々が定義を拡張する必要があることを示す。

その後で望む拡張がどのように実現されるかを見ることになる。

再度、次のような場合を考える：
\[
y = 1 \cdot 2 \dots n = n!.
\]
この場合において、$x=n$ のとき $n!$ に簡約される明確な $x$ の公式は存在しない。なぜなら、$x!$ は正の整数以外の値に対して意味を持たないからである。このような状況では、内插問題を解決する試みが数学において重要な進歩をもたらした。実際、数学的には、所望の性質と多くの他の興味深い重要な性質を持つ関数（ガンマ関数）を見つけることに成功した。

\Paragraph{52. 有限と無限の集合。}さらに進む前に、純粋数学において恒常的に出現する抽象的で論理的な概念についていくつか述べる必要がある。

まず第一に、読者はおそらく「クラス」という概念を理解しているだろう。ここでは、「クラス」に関する論理的な困難性については議論しないが、大まかには、ある特定の性質を持つすべての存在や物体を集めた集合または集まりを指すと説明できる。例えば、イギリス人、議員、正整数、実数などのクラスがある。

さらに、読者はおそらく「有限」または「無限」なクラスという概念についても理解しているだろう。例えば、「イギリス人」というクラスは有限なクラスであり：過去、現在、未来のすべてのイギリス人の集合は有限な数 $n$ に相当するが、もちろん現在その実際の値を知ることはできない。一方、「現在のイギリス人」のクラスには、人口調査の方法が十分であれば数え上げることで特定できる数 $n$ がある。

一方、正整数のクラスは有限ではなく無限です。これは次のようにより正確に表現できます。
$n$が任意の正整数である場合、例えば $1000$, $1,000,000$ または私たちが思い浮かべるどんな数字でも、それ以上の正整数があります。したがって、私たちが思い浮かべた数字が $1,000,000$ の場合、明らかに少なくとも $1,000,001$ の正整数があります。同様に、有理数のクラスや実数のクラスも無限です。これは便利のために、正整数や有理数や実数には \emph{無限} な個数があると言えるでしょう。しかし、読者は常にこの表現が単に該当するクラスが有限個数の要素を持たないことを意味することを忘れないように注意してください、すなわち $1000$ や $1,000,000$ のような有限な数ではありません。

53. 大きな $n$ の値における関数 $\phi(n)$ が持つ性質
現在、§50--§51 で議論していた「$n$ の関数」に戻ります。これらは、§2（第II章）で議論した $x$ の関数とは多くの違いがありますが、二つの関数のクラスには一つの基本的な特性が共有されています：\emph{定義される変数の値が無限のクラスを形成すること}。この事実は以下の考慮事項の基礎となり、次章で見るように、$x$ の関数に対しても \textit{mutatis mutandis} に適用されます。

$\phi(n)$ を $n$ の任意の関数とし、$P$ を $\phi(n)$ が持つか持たないかの性質として、例えば正整数であるか、または $1$ より大きいかのような性質を考えてみましょう。それぞれの値 $n = 1, 2, 3, \dots$ に対して $\phi(n)$ が性質 $P$ を持つか持たないかを考えます。すると以下の三つの可能性があります：---

(\ia) $\phi(n)$ はすべての $n$ の値において、または有限な数 $N$ 個の $n$ の値を除いて、性質 $P$ を持つかもしれません。

\item(\ib) $\phi(n)$は任意の$n$に対してその性質を持つことなく、または有限個の$n$，すなわち$N$個だけにおいてその性質を持つ場合があります。

\item(\ic) (\ia) と (\ib) のいずれも真でない場合があります。

(\ib)が真である場合、$\phi(n)$がその性質を持つ$n$の値は有限のクラスを形成します。(\ia)が真である場合、$\phi(n)$がその性質を持たない$n$の値は有限のクラスを形成します。第三の場合には両方のクラスが有限ではありません。いくつかの具体的な例を考えてみましょう。

\par\medskip\noindent\textbf{【注釈】}\par
\item(1) $\phi(n) = n$ とし、$P$ を正整数であるという性質とするとき、$\phi(n)$はすべての$n$に対してその性質を持つことになります。

一方で、$P$ を1000以上の正整数であるという性質とするとき、$\phi(n)$はすべての$n$に対してその性質を持つことなく、有限個の$n$，すなわち1, 2, 3, \dots, 999だけを除いてその性質を持つことになります。これらのいずれの場合も(\ia)が真です。

\item(2) $\phi(n) = n$ とし、$P$ を1000未満であるという性質とするとき、(\ib)が真です。

\item(3) $\phi(n) = n$ とし、$P$ を奇数であるという性質とするとき、(\ic)が真です。なぜなら、$\phi(n)$は$n$が奇数のときに奇数で、$n$が偶数のときに偶数であり、両方の$n$の値は無限のクラスを形成するからです。

例題．以下の各場合において、(\ia)，(\ib) または (\ic) のいずれが真かを考慮してみましょう：

\itemp(i) $\phi(n) = n$，$P$ を平方数であるという性質とするとき，

\itemp(ii) \Hang $\phi(n) = p_{n}$，ここで $p_{n}$ は$n$番目の素数を表し、$P$ を奇数であるという性質とするとき，

\itemp(iii) $\phi(n) = p_{n}$，$P$ を偶数であるという性質とするとき，

\itemp(iv) $\phi(n) = p_{n}$，$P$ を$\phi(n) > n$ という性質とするとき，

\itemp(v) $\phi(n) = 1 - (-1)^{n}(1/n)$，$P$ を$\phi(n) < 1$ という性質とするとき，

\itemp(vi) $\phi(n) = 1 - (-1)^{n}(1/n)$，$P$ を$\phi(n) < 2$ という性質とするとき，

\itemp(vii) $\phi(n) = 1000\{1 + (-1)^{n}\}/n$，$P$ を$\phi(n) < 1$ という性質とするとき。

\item{(viii)} $\phi(n) = 1/n$, $P$ は $\phi(n) < .001$ の性質を表す。

\item{(ix)} $\phi(n) = (-1)^{n}/n$, $P$ は $|\phi(n)| < .001$ の性質を表す。

\item{(x)} \Hang $\phi(n) = 10\MC000/n$, または $(-1)^{n}10\MC000/n$, $P$ は $\phi(n) < .001$ または $|\phi(n)| < .001$ のいずれかの性質を表す。

\item{(xi)} $\phi(n) = (n - 1)/(n + 1)$, $P$ は $1 - \phi(n) < .0001$ の性質を表す。

\Paragraph{54.} 現在、$\phi(n)$ と $P$ が存在し、命題 (\ia) が真であるとする。つまり、$\phi(n)$ は $P$ の性質を持つが、すべての $n$ 値に対してではなく、有限個の $N$ 個の例外的な値 $n_{1},\ n_{2},\ \dots,\ n_{N}$ を除いては持つ。これらの例外的な値を次のように表すことができる：
\[
n_{1},\ n_{2},\ \dots,\ n_{N}.
\]
当然、これらの $N$ 個の値が最初の $N$ 個の値 \(1,~2,\ \dots,\ N\) である必要はなく、前節の例からわかるように、実際にはしばしばそのような状況となる。しかし、それがそうであっても、またはそうでない場合でも、$\phi(n)$ は $n > n_{N}$ のとき $P$ の性質を持つことが知られている。したがって、第 \(n\) 番目の素数は \(n > 2\) のとき奇数であり、唯一の例外は \(n = 2\) である；また、\(1/n < .001\) は \(n > 1000\) のとき成り立つが、最初の \(1000\) 個の $n$ の値が例外であり；さらに
\[
1000\{1 + (-1)^{n}\}/n < 1
\]
は \(n > 2000\) のとき成り立つが、その例外的な値は \(2,\ 4,\ 6,\ \dots,\ 2000\) である。つまり、これらの各々の例において、ある特定の値から全ての $n$ 値に対して性質を有するということが言える。

Weしばしば，$\phi(n)$が大値，または非常に大値，または十分に大値の$n$に対してその性質を持つという表現を使用します。したがって，$\phi(n)$が性質$P$（通常は不等式を用いた関係で表される性質）を大値の$n$に対して持つと言った場合、これはある定数$n_{0}$を決定できることを意味します。そして，すべての$n_{0}$以上の$n$に対する$\phi(n)$がその性質を持つことを言います。上記の例では，この$n_{0}$は例外的な数値$n_{N}$より大きい任意の数で取ることができます：最も自然なのは$n_{N} + 1$を取ることです。

したがって，「すべての大素数は奇数である」とか「$\frac{1}{n}$は大値の$n$に対して0.001未満である」と言えるでしょう。そして，このような文脈での「大」の使用について読者は慣れなければなりません。「大」は単独で見れば，数学においても日常生活の言語においても絶対的な意味を持つものではありません。日常生活中では，一つの状況で大きな数でも別の状況では小さな数になるということが真実です：6ゴールはサッカーの試合における高得点ですが，6ランスはクリケットの試合における低得点であり，400ランスは高得点ですが400ポンドは高収入ではありません。そしてもちろん，数学においても「大」は一般的に「十分に大きい」という意味で，一つの目的では十分に大きい数でも別の目的では十分ではないかもしれません。

我々は$\phi(n)$が大値の$n$に対して性質$P$を持つという主張の意味を理解しました。この種の主張について，本章全体を通じて扱います。

55. 「n が無限大に収束する」という表現について。この問題に対する別の見方があるのは自然である。n が順次 1, 2, 3, \dots の値をとるとする。単語「順次」は時間の経過を連想させるので、もし我々が希望すれば n は各々の秒の開始時にこれらの値をとるとすることができる。すると、秒が経つにつれて n はますます大きくなり、その増加の範囲には限界がない。どんなに大きな数を考えても（例えば 2,147,483,647）、ある時刻において n はそれよりも大きい値を取る。

無尽蔵の増大を表現する短い表現が必要であるため、n が「無限大に収束する」という言い方をするか、または n \to \infty とする。この最後の記号は通常、「無限大」を略すために使用される。言葉「収束する」のように、この表現も時間における変化の概念を自然に連想させるので、n の変化を上記の方法で時間経過とともに完成させることを考えるのは便利である。しかし、これは単なる便宜上のことであり、変数 n は時間とは直接関係のない純粋な論理的-entity である。

読者は「n が無限大に収束する」という表現を強調すべきである。つまり、「n は無限大に向かって一連の値をとる」という意味である。実際には「無限大」そのものは存在しない：方程式
\[
n = \infty
\]
は現在の形で絶対的に意味不明であり、n は等しいものとして∞とはならない。なぜなら、「∞に等しい」という表現自体が意味がないからである。実際には、符号∞は「無限大に向かって収束する」の一語の中でしか意味を持たない：その意味については上記で説明した通りである。後ほど、他の符号∞を含む表現にも意味を与える方法を学ぶが、読者は常にこの点を頭に置いておくべきである。

\Item{(1)} 无限符号 $\infty$ 自身は何を意味するわけではなく、ただし含む文句によっては何かを意味することがあります，

\Item{(2)} 含む文句が何かを意味する場合でも、それは以前に特定の定義によりその文句に意味を与えたことによるものです。

今、$\phi(n)$ が大きな $n$ の値に対して性質 $P$ を持つものとし、$n$ が $\infty$ に「傾き」る場合、つまり我々が先ほど説明した意味で $n$ が十分大きい値を取ることになるので、$\phi(n)$ も性質 $P$ を持つことになります。従って、「十分大きな $n$ の値に対して $\phi(n)$ は何らかの性質を持つのか？」という問いは、「$n$ が $\infty$ に傾きるとき $\phi(n)$ の振る舞いはどうなるのか？」と表現することもできます。

\Paragraph{56. $n$ が無限大に向かうときの関数の振る舞い。} 我々は、前節で述べたことに照らし合わせて、より高度な数学において頻繁に見られる声明の意味について考察します。例えば、以下の2つの声明を考えてみましょう：(\ia) $1/n$ は大きな $n$ の値に対して小さい，(\ib) $1 - (1/n)$ は大きな $n$ の値に対してほぼ $1$ に等しい。これらは明らかにそうであるように思われますが、それらには読者にとって注意を払う価値があります。まず (\ia) を考えてみましょう、これはやや単純です。

我々はすでに「$1/n$ は大きな $n$ の値に対して $.01$ より小さい」という声明について考慮しました。これは、不等式 $1/n < .01$ がすべての $n$ の値（ある特定の値より大きい）で真であることを意味します。実際には、$1/n < .01$ は $n > 100$ のとき成立します。同様に、「$1/n$ は大きな $n$ の値に対して $.0001$ より小さい」ことも真です：実際には、$1/n < .0001$ は $n > 10,000$ のとき成立します。そして $.01$ や $.0001$ の代わりに $.000001$ や $.00000001$、または任意の正の数を取ることもできます。

しかし、この主張は複雑なもので、$\DELTA$ に特別な値（例えば $.01$）を与えることで得られる一連の主張全体を表しています。もちろん、$\DELTA$ が小さいほど $1/\DELTA$ が大きくなり、それにより「十分大きな」 $n$ の最小値も大きくなります。ある $\DELTA$ の値に対して十分大きいとされる $n$ の値は、より小さな $\DELTA$ のときには不十分です。

最後に強調された文は、$1/n$ が大なる $n$ のとき小さいという主張（式 (\ia)）を意味します。同様に、(\ib) は「もし $\phi(n) = 1 - (1/n)$ であれば、「十分大きな」 $n$ の値に対して $1 - \phi(n) < \DELTA$ が真である」という主張（任意の正の数（例えば $.01$ や $.0001$）を $\DELTA$ に与えたときも）を意味します。これは $1 - \phi(n) = 1/n$ の事実から明らかです。

\begin{align*}
&\text{$n$が大きいとき、$1/n$は小さい，} \\
&\text{$n$が無限大に近づくとき、$1/n$は0に収束する，}
\end{align*}

という主張は互いに等しく、より形式的な主張

\begin{quotation}
任意の正数$\DELTA$について、どの程度小さくても、十分大きい$n$に対して$1/n < \DELTA$である，
\end{quotation}

またはさらに形式的な主張

\begin{quotation}
任意の正数$\DELTA$について、どの程度小さくても、ある数$n_{0}$を選び、$n_{0}$以上のすべての$n$に対して$1/n < \DELTA$であるとすることができる，
\end{quotation}

に等しい。

最後の主張において現れる数$n_{0}$は当然$\DELTA$の関数であり、この事実を強調するために、$n_{0}(\DELTA)$という形で記述することもある。

以下は注記である。

読者は自分の主張が真であることを問う相手に直面すると想像してほしい。彼は小さな数列を挙げる。最初の数は$0.001$とすることができる。読者は$n > 1000$となるとき、$1/n < 0.001$であると答える。相手はこれを受け入れるが、さらに小さな数、例えば$0.000\MS000\MS1$を試みるだろう。読者は$n > 10\MC000\MC000$となるとき、$1/n < 0.000\MS000\MS1$であると答える。そしてこの単純な例では明らかに、読者が常に論争で有利であることがわかる。

次の性質を関数 \(1/n\) で表現する方法を導入します。関数 \(1/n\) の \(n\) が無限大に近づくときの `\emph{極限}は0である' という文を用います。これは記号的に以下のように表すことができます：
\[
\lim_{n\to\infty} \frac{1}{n} = 0,
\]
または単に $\lim(1/n) = 0$ と書くこともあります。また、しばしば以下のように書きます：
`$1/n \to 0$
as $n \to \infty$', これは `$1/n$ は $n$ が無限大に近づくとき $0$ に収束する' の意味です；または単に `$1/n \to 0$' と書きます。同様にして、以下の式も用います：
\[
\lim_{n\to\infty} \left(1 - \frac{1}{n}\right) = 1,\quad
\lim \left(1 - \frac{1}{n}\right) = 1,
\]
または $1 - (1/n) \to 1$。

57. 次に異なる例を考えてみましょう。関数 $\phi(n) = n^{2}$ とします。このとき `\emph{$n^{2}$ は $n$ が大きいときに大きい}' という文が成り立ちます。これはより形式的な言い方として以下のように表せます：
`任意の正の数 $\Delta$ に関わらず、十分に大きな $n$ の値に対して $n^{2} > \Delta$ である',
または
`ある数 $n_{0}(\Delta)$ を選べば、$n \geq n_{0}(\Delta)$ なるすべての $n$ について $n^{2} > \Delta$ である'。

このような場合、自然に `$n^{2}$ は $n$ が無限大に近づくとき $\infty$ に収束する' や `$n^{2}$ は $n$ とともに $\infty$ に収束する' と述べ、以下のように書きます：
\[
n^2 \to \infty.
\]

最後に関数 $\phi(n) = -n^{2}$ を考えてみましょう。このとき $n$ が大きいときに $\phi(n)$ は大きな負の値を取るため、自然に `$-n^{2}$ は $n$ が無限大に近づくとき $-\infty$ に収束する' と述べ、以下のように書きます：
\[
-n^{2} \to -\infty.
\]

このような符号 $\pm\infty$ の使用により、しばしば $n^{2} \to +\infty$ を $n^{2} \to \infty$ と書き換えることが便利であり、一般的に $\infty$ の代わりに $+\infty$ を用いることで記号の統一性を確保することが好ましい有時会。

ただし、すべてのこれらの主張において、記号 $\infty$, $+\infty$, $-\infty$ は自身で何らの意味を持たず、特定の特殊な接続の中で説明によって与えられた解釈により意味を獲得するものである。

\ Paragraph{58. 极限の定義。} 前段落の議論後、読者は一般の極限の概念を理解できるべきである。ざっくりと言えば、$\phi(n)$ が $n$ が大きいとき $l$ に近似的に等しいとすれば、$\phi(n)$ は $n$ が無限大に向かうとき $l$ に収束すると言える。しかし、この主張の意味は前段落で与えられた説明によって十分に理解できるとしても、厳密な数学的定義としては十分ではない。実際、これは「十分大きな値の $n$ に対して $\phi(n)$ が $l$ から $\DELTA$ よりも小さい差異を持つ」の一連の主張と等価である。この主張は $\DELTA = .01$ または $.0001$ または任意の正数に対して真でなければならず、そのような任意の $\DELTA$ の値に対しては、ある一定の値 $n_{0}(\DELTA)$ 以降の任意の $n$ についても真である。通常、より小さな $\DELTA$ ほどこの値 $n_{0}(\DELTA)$ が大きくなる。

従って、以下の形式的な定義を設定する：

\ Definition[I.] 
関数 $\phi(n)$ は $n$ が無限大に向かうとき極限 $l$ に収束すると、任意の正の数 $\DELTA$ に関わらず、十分大きな値の $n$ に対して $\phi(n)$ が $l$ から $\DELTA$ よりも小さい差異を持つと定義される。すなわち、任意の正の数 $\DELTA$ に対し、$\DELTA$ に応じて $n_{0}(\DELTA)$ を決定することができ、その値 $n_{0}(\DELTA)$ より大きいまたは等しいすべての $n$ に対して $\phi(n)$ が $l$ から $\DELTA$ よりも小さい差異を持つと定義される。

差分 φ(n) と l の間の正の値は |φ(n) - l| で表します。これは φ(n) - l や l - φ(n) のいずれかが正であるとき、その値を指し、現在は実数値、正または負の値のみを考えています。

この表記法では、以下のように定義を簡略化することができます：「任意に与えられた正の数 δ 无论多小，我達都能找到 n₀(δ) 使得当 n ≥ n₀(δ) 時 |φ(n) - l| < δ，則称 φ(n) 在 n 趋于无穷大時趋于极限 l，并写作
\[
\lim_{n \to \infty} \phi(n) = l.
\]」

以下は注記です：
しばしば「n → ∞」を省略することもありますし、簡略化のため、φ(n) → l と書くこともあります。

読者は、少数の簡単な例で n₀ を δ の関数として明示的に表す式を導き出す作業を通じて有益な洞察を得ることでしょう。例えば φ(x/n) = 1/n のとき l = 0 であり、条件は n ≥ n₀ に対して 1/n < δ に簡約されます。これは n₀ = 1 + [1/δ] で満たされることがわかります。（ここで以降、[x] は第2章の意味で x の最大整数を指します）

すべての δ 値に対する同じ n₀ を持つ唯一のケースがあります。
n ≥ N から一定の値 C に等しい φ(n) のとき、φ(n) - C = 0 が成り立つので、|φ(n) - C| < δ は n ≥ N およびすべての正の δ 値に対して満たされます。そして |φ(n) - l| < δ が n ≥ N およびすべての正の δ 値に対して満たされるとき、φ(n) = l が n ≥ N のときに明らかに成り立つので、そのような値のすべてに対する φ(n) は一定となります。

59. 枠限の定義は次のように幾何学的に例示することができます。関数$\phi(n)$のグラフは、$n = 1$, $2$, $3$, \dotsに対応する点から成ります。

直線$y = l$を引きます。また、その距離$\DELTA$離した平行な直線$y = l - \DELTA$, $y = l + \DELTA$も引きます。次に
\[
\lim_{n \to \infty} \phi(n) = l,
\]
\%[図表：図27。]
\ifthenelse{\boolean{Modernize}}{\%
図27（p117）
}{\%
図27（p117$orig_{notation}$）\%
}
の定義は、これらの直線を引いた後でも、どの程度近くとも、常に$x = n_{0}$という直線を引くことができ、その直線上およびそれより右側にあるグラフ上の点が、これら二つの直線の間に存在することを意味します。このような視覚的な方法は、実数変数に対して定義された関数について論じる際に特に有用です。

60. 以上は$n$に近づくとき$\phi(n)$が極限に向かう関数についての説明です。次に、$n^{2}$や$-n^{2}$のような正または負の無限大に向かう関数に対する対応する定義を述べます。この点は、既に読者には理解できるでしょう。

\par\medskip\noindent\textbf{【定義】}\par[II.]
関数$\phi(n)$が$n$とともに$+\infty$(正の無限大)に向かうとは、任意に大きな数$\Delta$を与えても、$\phi(n) > \Delta$となるような$n_{0}(\Delta)$を定めることができることを意味します。すなわち、どの程度大きな$\Delta$でも、十分大きい$n$に対して$\phi(n) > \Delta$が成り立つことを意味します。
\par\medskip

別の、より正確でない表現は「$\phi(n)$ を十分に大きい$n$によって何でも大きくできる」というものである。これは、$\phi(n)$ は $n \geq n_{0}(\Delta)$ のすべての値に対して $\Delta$ よりも大きいという基本的な点を曖昧にする可能性があるという欠点がある。しかし、この表現を使用しても問題はない。ただし、その意味が明確である限りではある。

$\phi(n)$ が $+\infty$ に収束するとき、
\[
\phi(n) \to +\infty.
\]
負の無限大に向かう関数に対する対応する定義は読者に任せることにする。

61. 定義に関するいくつかの点

次の点に注意してほしい。

(1) 明らかに、$n$ の有限個の値に対して $\phi(n)$ の値を任意の方法で変更しても、$\phi(n)$ が $n \to \infty$ に向かうときの振る舞いには影響を与えない。例えば、$1/n$ は $n \to \infty$ で $0$ に収束する。$1/n$ の有限個の値を変更することで、新しい関数を作ることができる。たとえば、$n = 1, 2, 7, 11, 101, 107, 109, 237$ に対して $\phi(n) = 3$ で、それ以外のすべての $n$ に対して $\phi(n) = 1/n$ とする関数を考えることができる。この関数についても、元の関数 $1/n$ のように $\lim \phi(n) = 0$ である。同様に、$n = 1, 2, 7, 11, 101, 107, 109, 237$ に対して $\phi(n) = 3$ で、それ以外のすべての $n$ に対して $\phi(n) = n^2$ とする関数についても、$\phi(n) \to +\infty$ であることが成り立つ。

\item{(2) 他方、通常無限個の $\phi(n)$ の値を変更することは $n$ が無限大に近づくときの $\phi(n)$ の振る舞いを根本的に影響することになる。例えば、$1/n$ を $n$ が $100$ の倍数であるとき $1$ に変更すると、$\lim \phi(n) = 0$ という命題は成り立たなくなる。有限個の値のみが影響を受ける限りでは、常に $\phi(n)$ の値が変更された最大の $n$ よりも大きい $n_0$ を定義することができる。例えば、上の例では常に $n_0 > 237$ と取ることができるし、実際には、仮想的な相手が最初の例で $\DELTA$ の値を $3$ に、または二番目の例で $\Delta$ の値を $3$ に設定した時点で、この選択を余儀なくされる。しかし、いかなる $n_0$ が大きいかに関わらず、$\phi(n)$ の値が変更されたより大きな $n$ が存在する。

\item{(3) 定義 I のテストを適用する際には、当然のことながら $|\phi(n) - l| < \DELTA$ という条件は $n = n_0$ のときだけでなく、$n \geq n_0$ のすべての $n$ に対して成り立つことが絶対的に重要である。例えば、$\phi(n)$ が最後に考慮された関数の場合、与えられた $\DELTA$ に対して $|\phi(n)| < \DELTA$ となる $n = n_0$ の $n_0$ を選ぶことができる：十分大きな値の $n$ を選び、それが $100$ の倍数でない限りである。しかし、このように $n_0$ を選んだ場合でも、$|\phi(n)| < \DELTA$ が $n \geq n_0$ のすべての $n$ に対して成り立つとは限らない：それよりも大きなすべての $100$ の倍数はこの命題に反する例外となる。}

\item{(4)} $\phi(n)$が常に$l$より大きい場合、$|\phi(n) - l|$を$\phi(n) - l$に置き換えることができる。したがって、$1/n$が$n \to \infty$のとき極限0に収束するかどうかは、$n \geq n_{0}$のとき$1/n < \DELTA$であるか否かによって決まる。しかし、$\phi(n) = (-1)^{n}/n$の場合、$l$もまた0であり、しかし$\phi(n) - l$は正でも負でもある。このような場合、条件を$|\phi(n) - l| < \DELTA$の形で表さなければならない。例えば、この特定の場合では$|\phi(n)| < \DELTA$という形で表す。

\item{(5)} \emph{$l$は$\phi(n)$の実際の値の一つである可能性がある。} たとえば、すべての$n$に対して$\phi(n) = 0$の場合、明らかに$\lim\phi(n) = 0$であることがわかる。また、(2)および(3)で述べたように、$n$が100の倍数のとき関数の値を$1$から$0$に変更した場合、$n$が100の倍数のときは$\phi(n) = 0$であり、それ以外の場合は$\phi(n) = 1/n$となる関数$\phi(n)$を得る。この関数の$n \to \infty$における極限は依然として明らかに0である。この極限値は無数の$n$の値で関数の値と一致する、すなわちすべての100の倍数においてその値が0である。

一方で\emph{極限自体は任意の$n$の値に対して関数の値とは必ずしも一致しない}ことが十分に明らかである。これは$\phi(n) = 1/n$の場合で十分に理解できる。この場合、極限値は0であるが、任意の$n$の値において関数の値は決して0にはならない。

読者にはこれらの事実を強く心に刻むべきである。極限は関数の値ではない：それはそれらの値と全く異なる何かであり、しかしその定義はそれらの関係によって与えられ、ある場合ではいくつかの値と等しいかもしれない。

関数
\[
\phi(n) = 0,\ 1,
\]
の場合、極限はすべての$\phi(n)$の値に等しい：一方で
\[
\phi(n) = \frac{1}{n},\quad
(-1)^{n}/n,\quad
1 + \frac{1}{n},\quad
1 + \left(\frac{(-1)^{n}}{n}\right)
\]
の場合、極限はどの$\phi(n)$の値とも等しくない：また
\[
\phi(n) = \frac{\sin\tfrac{1}{2}n\pi}{n},\quad
1 + \left(\frac{\sin\tfrac{1}{2}n\pi}{n}\right)
\]
の場合（これらの極限は$n$が$\infty$に近づくとき、容易に$0$と$1$であることがわかるため）、極限はすべての偶数の$n$に対して$\phi(n)$が取る値に等しいが、奇数の$n$に対する$\phi(n)$の値はすべて極限とは異なり、また互いに異なる。

(6) $n$が非常に大きいとき関数が常に非常に大きな値を取っても、それが$+\infty$または$-\infty$に向かうわけではない。これは$\phi(n) = (-1)^{n} n$によって十分な例として示される。関数は$+\infty$または$-\infty$に向かうためには、ある値の$n$を超えてから一定の符号を維持する必要がある。

例題XXIII. $n$が$\infty$に近づくとき次のような関数
\[
\phi(n) = n^{k}
\]
（$k$は正または負の整数または有理分数）の振る舞いを考える：

1. $k$が正の場合、$n^{k}$は$n$とともに$+\infty$に向かう。$k$が負の場合、$\lim n^{k} = 0$である。$k = 0$の場合、$n^{k} = 1$であり、すべての$n$に対して恒等的に成り立つ。したがって$\lim n^{k} = 1$である。

読者には、この単純な例においても、定義の条件が満たされていることを形式的な証明で確認することの有用性を体験していただきます。例えば \( k > 0 \) の場合を考えます。任意に与えられた数 \(\Delta\) に関し、\( n^{k} > \Delta \) となるような \( n_{0} \) を選ぶことを目指します。実際には、\(\sqrt[k]{\Delta}\) より大きい任意の数を \( n_{0} \) として選べばよいです。例えば \( k = 4 \) の場合、\( n^{4} > 10,000 \) となるのは \( n \geq 11 \) のときで、\( n^{4} > 100,000,000 \) となるのは \( n \geq 101 \) のときです。

2. \(\phi(n) = p_{n}\)，ここで \( p_{n} \) は第 \( n \) 番目の素数である。素数が有限個しかないという仮定下では、\(\phi(n)\) は有限個の \( n \) の値に対してのみ定義される。しかし、ユークリッドによって最初に示されたように、素数は無限に存在する。ユークリッドの証明は以下の通りです。素数が有限個しかないという仮定のもとで、それらを \( 1, 2, 3, 5, 7, 11, \dots, N \) とする。このとき、数 \( 1 + (1 · 2 · 3 · 5 · 7 · 11 \dots N) \) を考慮します。この数は明らかに \( 2, 3, 5, \dots, N \) のいずれにも割り切れないことがわかります。なぜなら、これらの数で割ったときの余りはすべて \( 1 \) だからです。したがって、この数は \( 1 \) 以外の素数で割り切れないため、その自身も素数であり、これは我々の仮定に矛盾します。

さらに明らかに、\(\phi(n) > n\)（\( n = 1, 2, 3 \) を除く）すべての \( n \) の値に対して成立します。したがって、\(\phi(n) \to +\infty\)。

3. \(\phi(n)\) を \( n \) より小さい素数の個数と定義すると、また \(\phi(n) \to +\infty\) です。

4. \(\phi(n) = [\alpha n]\)，ここで \(\alpha\) は任意の正の数である。このとき
\[
\phi(n) = 0\quad (0 \leq n < 1 / \alpha),\qquad
\phi(n) = 1\quad (1/\alpha \leq n < 2/\alpha),
\]
と続く；そして \(\phi(n) \to +\infty\)。

5. \(\phi(n) = 1,000,000/n\) の場合、\(\lim\phi(n) = 0\)：そして \(\psi(n) = n/1,000,000\) の場合、\(\psi(n) \to +\infty\)。これらの結論は最初に \(\phi(n)\) が \(\psi(n)\) よりもずっと大きいという事実によって影響を受けません。実際、\(\phi(n)\) は \( n = 1,000,000 \) まで \(\psi(n)\) より大きいです。

\item 6. $\phi(n) = 1/\{n - (-1)^{n}\}$, $n - (-1)^{n}$, $n\{1 - (-1)^{n}\}$. 最初の関数は0に収束せず、次の関数は正の無限大に収束し、最後の関数は極限値や正の無限大には収束しない。

\item 7. $\phi(n) = (\sin n\theta\pi)/n$，ここで$\theta$は任意の実数である。このとき$|\phi(n)| < 1/n$であり、$|\sin n\theta\pi| \leq 1$から$\lim\phi(n) = 0$。

\item 8. $\phi(n) = (\sin n\theta\pi)/\sqrt{n}$, $(a\cos^{2} n\theta + b\sin^{2}n\theta)/n$，ここで$a$と$b$は任意の実数である。

\item 9. $\phi(n) = \sin n\theta\pi$. もし$\theta$が整数であれば、すべての$n$に対して$\phi(n) = 0$であり、従って$\lim\phi(n) = 0$。

次に$\theta$が有理数である場合を考える。例えば、$\theta = p/q$とし、$p$と$q$は正の整数である。$n = aq + b$とおく、ここで$n$を$q$で割った商を$a$、余りを$b$とする。すると$\sin(np\pi/q) = (-1)^{ap}\sin(bp\pi/q)$となる。例えば$p$が偶数の場合、$n$が0から$q-1$まで増加するとき、$\phi(n)$は次の値を取ります：
\[ 0, \quad \sin(p\pi/q), \quad \sin(2p\pi/q), \ldots \quad \sin\{(q - 1)p\pi/q\} \]
$n$が$q$から$2q-1$まで増加するとき、これらの値は繰り返され、$n$が$2q$から$3q-1$、$3q$から$4q-1$と続くときも同様に繰り返されます。したがって$\phi(n)$の値は有限個の異なる値を無限に繰り返します。このような場合、$\phi(n)$は$n$が無限大に向かうにつれて極限値や正の無限大、負の無限大には収束しないことが明らかです。

$\theta$が無理数である場合も少し複雑で、次の例題群で議論します。

\paragraph{62. 振動関数}
 \par\medskip\noindent\textbf{【定義】}\par
$\phi(n)$が$n$が無限大に向かうとき極限値や正の無限大、負の無限大には収束しない場合、$\phi(n)$は$n$が無限大に向かうとき振動するという。
\par\medskip

関数 $\phi(n)$ は、その値が最後の例で考慮された場合のように周期的な再現を示すときには必ず振動します。しかしもちろん、この特異性を持たない場合でも振動することがあります。振動は純粋に否定的に定義されます：関数が他の特定の動作を行わないときには振動します。
\pagetotal{122}

最も単純な振動する関数の例は
\[
\phi(n) = (-1)^n,
\]
で与えられます。この関数は $n$ が偶数のときは $+1$、奇数のときは $-1$ です。この場合、値は周期的に再現されます。しかし次を考慮します
\[
\phi(n) = (-1)^n + \frac{1}{n},
\]
その値は
\[
-1 + 1,\quad
1 + \frac{1}{2},\quad
-1 + \frac{1}{3},\quad
1 + \frac{1}{4},\quad
-1 + \frac{1}{5},\ \dots.
\]
$n$ が大きいとき、すべての値は $+1$ や $-1$ に近似されます。明らかに $\phi(n)$ は極限値や $+\infty$ または $-\infty$ には収束せず、したがって振動します：しかし値は再現されません。この場合、すべての $\phi(n)$ の値は数値的に $3/2$ よりも小さいか等しいことが観察されます。同様に
\[
\phi(n) = (-1)^n \cdot 100 + \frac{1000}{n}
\]
も振動します。$n$ が大きいとき、すべての値は $100$ や $-100$ に近似されます。数値的に最大の値は $900$（$n = 1$ のとき）です。しかし $\phi(n) = (-1)^n \cdot n$ を考慮します、その値は $-1$, $2$, $-3$, $4$, $-5$,~\dots です。この関数も振動します：極限値や $+\infty$ や $-\infty$ には収束しません。そしてこの場合、項の数値的な値が一定を超えて上昇しない限界を指定することはできません。これらの二つの例の違いはさらに定義を提示します。

\par\medskip\noindent\textbf{【定義】}\par
関数 $\phi(n)$ が $n$ が無限大に近づくとき振動する場合、$\phi(n)$ は数値 $K$ を指定できるかどうかにより有限振動または無限振動と呼ばれます。すなわち、すべての $\phi(n)$ の値が数値的に $K$ よりも小さいか等しいとき、つまり $|\phi(n)| < K$ がすべての $n$ の値に対して成り立つときです。
\par\medskip

これらの定義は、§58および§60の定義も同様に、以下の例によってさらに説明される。

\#\#\# XXIV. 例

$n$が無限大に向かうとき、次の関数の振る舞いを考慮する：

1. $(-1)^{n}$, $5 + 3(-1)^{n}$, $(1\MC000\MC000/n) + (-1)^{n}$, $1\MC000\MC000(-1)^{n} + (1/n)$。

2. $(-1)^{n}n$, $1\MC000\MC000 + (-1)^{n}n$。

3. $1\MC000\MC000 - n$, $(-1)^{n}(1\MC000\MC000 - n)$。

4. $n\{1 + (-1)^{n}\}$。この場合、$\phi(n)$の値は
\[
0,\quad 4,\quad 0,\quad 8,\quad 0,\quad 12,\quad 0,\quad 16,\ \dots
\]
奇数項はすべて0であり、偶数項は無限大に収束する：$\phi(n)$は無限に振動する。

5. $n^{2} + (-1)^{n}2n$。二番目の項は無限に振動するが、$n$が大きいとき、最初の項は二番目の項よりも非常に大きくなる。実際、$\phi(n) \geq n^{2} - 2n$であり、$n^{2} - 2n = (n - 1)^{2} - 1$は$n > 1 + \sqrt{\Delta + 1}$のとき任意に与えられた値$\Delta$よりも大きい。したがって、$\phi(n) \to +\infty$である。この場合、$\phi(2k + 1)$は常に$\phi(2k)$より小さいので、関数は無限大に向かうために一連の前進と後退のステップを経て進行するが、しかし「振動」するとは定義上異なる。

6. $n^{2}\{1 + (-1)^{n}\}$, $(-1)^{n}n^{2} + n$, $n^{3} + (-1)^{n}n^{2}$。

7. $\sin n\theta\pi$。既に（例xxiii.9）示したように、$\theta$が有理数のとき$\phi(n)$は有限に振動し、$\theta$が整数でない限り$\phi(n) = 0$であり、$\phi(n) \to 0$である。

Theθが無理数の場合は少し難しい。しかし、φ(n)は有限に振動していることが直感的に理解できる。一般性を失うことなく0 < θ < 1とすることができる。まず第一に|φ(n)| < 1であるから、φ(n)は有限に振動するか極限に向かうべきである。第二の選択肢が実現可能かどうかを考える。lim sin nθπ = lと仮定する。

すると、どんなに小さなδであっても、n0を適切に選べば、すべてのn ≥ n0に対してsin nθπがl - δとl + δの間にあることがわかる。したがって、nに対応する全ての値で|sin(n + 1)θπ - sin nθπ|は2δ未満であり、|sin(1/2)θπ cos((n + 1/2)θπ)| < δである。

従って
\[
\cos \left(n + \frac{1}{2}\right)\theta\pi = \cos n\theta\pi \cos \frac{1}{2}\theta\pi - \sin n\theta\pi \sin \frac{1}{2}\theta\pi
\]
は|sin(1/2)θπ|未満である。同様に
\[
\cos \left(n - \frac{1}{2}\right)\theta\pi = \cos n\theta\pi \cos \frac{1}{2}\theta\pi + \sin n\theta\pi \sin \frac{1}{2}\theta\pi
\]
も|sin(1/2)θπ|未満であり、したがってcos nθπ cos(1/2)θπとsin nθπ sin(1/2)θπはともに|sin(1/2)θπ|未満である。すなわち、nが大きいときcos nθπ cos(1/2)θπは非常に小さく、これはcos nθπが非常に小さい場合のみ可能である。同様にsin nθπも非常に小さく、lは0であるべきである。しかし、cos nθπとsin nθπがともに非常に小さいことは不可能であり、その平方和は1であることから、sin nθπは∞に向かって振動する。

読者には、以下の議論に特に注意していただきたい。「cos nθπ cos(1/2)θπ」は非常に小さいのであり、これだけが「cos nθπ」が非常に小さい場合である」という主張を考慮してほしい。なぜ、「cos(1/2)θπ」が「非常に小さい」べきではないのかと疑問を持つかもしれない。この質問に対する答えは、この文脈における「非常に小さい」という表現の意味に求められる。大数nに対して「φ(n)は非常に小さい」と言及するとき、任意に与えられた数よりも数値的に小さくなるようなn₀を選べるという意味である。固定された数、例えばcos(1/2)θπ（これは0ではない）についてこの主張をすることは明白な誤りである。

同様にして、θが偶数整数でない限り「cos nθπ」は有限に振動することを証明せよ。

8. sin nθπ + (1/n), sin nθπ + 1, sin nθπ + n, (−1)ⁿ sin nθπ

9. a cos nθπ + b sin nθπ, sin²nθπ, a cos²nθπ + b sin²nθπ

10. a + bn + (−1)ⁿ (c + dn) + e cos nθπ + f sin nθπ

11. n sin nθπ。θが整数である場合、φ(n) = 0, φ(n) → 0。θが有理数でない（無理数）の場合、φ(n)は無限に振動する。

12. n(a cos²nθπ + b sin²nθπ)。この場合、aとbがともに正のときφ(n)は+∞に収束し、ともに負のとき−∞に収束する。a = 0, b > 0やa > 0, b = 0、またはa = 0, b = 0の場合を特別な場合として考慮せよ。aとbが符号が異なる場合、φ(n)は一般的に無限に振動する。例外的な場合も考慮せよ。

\item[13.] $\sin(n^{2}\theta\pi)$. もし$\theta$が整数であれば、$\phi(n) \to 0$となる。そうでなければ、$\phi(n)$は有限に振動し、これについては、\Exs{xxiii}.9および\Exs[]{xxiv}.7で使用した類似の手法がより複雑であるが、同様の議論によって示される。

\item[14.] $\sin(n!\theta\pi)$. $\theta$が有理数$p/q$であれば、すべての$n \geq q$に対して$n! \theta$は整数となる。従って$\phi(n) \to 0$である。$\theta$が無理性の場合には、より困難な概念を考慮する必要がある。

\item[15.] $\cos(n!\theta\pi)$, $a\cos^{2}(n!\theta\pi) + b\sin^{2}(n!\theta\pi)$, ここで$\theta$は有理数である。

\item[16.] $an - [bn]$, $(-1)^{n}(an - [bn])$。

\item[17.] $[\sqrt{n}]$, $(-1)^{n}[\sqrt{n}]$, $\sqrt{n} - [\sqrt{n}]$。

\item[18.] $n$の最小の素因数。$n$が素数であれば、$\phi(n) = n$である。$n$が偶数であれば、$\phi(n) = 2$である。従って$\phi(n)$は無限に振動する。

\item[19.] $n$の最大の素因数。

\item[20.] 平成$n$年の日数。

\par\medskip

\par\medskip\noindent\textbf{【例 XXV.】}\par
\item[1.] $\phi(n) \to +\infty$かつ$\psi(n) \geq \phi(n)$であれば、すべての$n$に対して、$\psi(n) \to +\infty$である。

\item[2.] $\phi(n) \to 0$かつ$|\psi(n)| \leq |\phi(n)|$であれば、すべての$n$に対して、$\psi(n) \to 0$である。

\item[3.] $\lim |\phi(n)| = 0$であれば、$\lim \phi(n) = 0$である。

\item[4.] $\phi(n)$が極限に収束または有限に振動し、かつ$n \geq n_{0}$に対して$|\psi(n)| \leq |\phi(n)|$であれば、$\psi(n)$も極限に収束または有限に振動する。

\item[5.] $\phi(n)$が$+\infty$または$-\infty$に収束または無限に振動し、かつ$n \geq n_{0}$に対して$|\psi(n)| \geq |\phi(n)|$であれば、$\psi(n)$も$+\infty$または$-\infty$に収束または無限に振動する。

\item6. $\phi(n)$が振動し、いかなる$n_0$においても$n > n_0$となる$n$の値に対して$\psi(n) > \phi(n)$であり、また$n > n_0$となる$n$の値に対して$\psi(n) < \phi(n)$であるなら、$\psi(n)$が振動すること。これは真か？偽であれば反例を挙げよ。

\item7. $\phi(n) \to l$ とする$n \to \infty$のとき、任意の固定整数$p$に対して$\phi(n + p) \to l$である。[これは定義からすぐに徴することができる。同様に、$\phi(n)$が$+\infty$や$-\infty$に収束したり振動するなら、$\phi(n + p)$も同じように収束または振動することを確認できる]

\item8. $p$が$n$と共に変化し、常に固定の正整数$N$未満である場合（振動の場合を除き）や、$p$が任意の方法で$n$と共に変化し、常に正である場合においても、同じ結論が成り立つ。

\item9. $n_0$の最小値を求めよ。このとき
\[
(\ia) \quad n^2 + 2n > 999 \quad (n \geq n_0), \qquad
(\ib) \quad n^2 + 2n > 1{,}000{,}000 \quad (n \geq n_0).
\]

\item10. $n_0$の最小値を求めよ。このとき
\[
(\ia) \quad n + (-1)^n > 1{,}000 \quad (n \geq n_0), \qquad
(\ib) \quad n + (-1)^n > 1{,}000{,}000 \quad (n \geq n_0).
\]

\item11. $n_0$の最小値を求めよ。このとき
\[
(\ia) \quad n^2 + 2n > \Delta \quad (n \geq n_0), \qquad
(\ib) \quad n + (-1)^n > \Delta \quad (n \geq n_0),
\]
$\Delta$が任意の正数である。

[(\ia) $n_0 = [\sqrt{\Delta + 1}]$: (\ib) $n_0 = 1 + [\Delta]$または$2 + [\Delta]$、$[\Delta]$が奇数か偶数かにより、すなわち$n_0 = 1 + [\Delta] + \frac{1}{2} \{1 + (-1)^{[\Delta]}\}$]

\item{12.} \( n_{0} \) の最小値を定めよ，使得
\[
ZZZMATHTOKEN1ZZZ
\]
が \( n \geq n_{0} \) のとき成立する。[後者の場合を取ろう。まず
\[
\frac{1}{n} + \left( \frac{(-1)^{n}}{n^{2}} \right) \leq \frac{n + 1}{n^{2}},
\]
であり，容易に \( n \geq n_{0} \) のとき \(\frac{n + 1}{n^{2}} < 0.000\MS001\) を満たす最小の \( n_{0} \) は \( 1\MC000\MC002 \) であることがわかる。しかし与えられた不等式は \( n = 1\MC000\MC001 \) のときも成立し，これが求められる \( n_{0} \) の値である。]

\paragraph{63. 极限に関するいくつかの一般定理}
\topic{\item{A.} 已知行為的两个函数之和的行為}

\par\medskip\noindent\textbf{【定理】}\par[I.] ] \(\phi(n)\) と \(\psi(n)\) が極限 \(a\)，\(b\) に収束するとき，\(\phi(n) + \psi(n)\) は極限 \(a + b\) に収束する。
\par\medskip

これはほとんど明らかである。\footnote
  {この文句には一定の曖昧さがあり，読者は注意すべきである。`ある定理はほとんど明白である'と言ったとき，それは二つの意味を持つことができる。一つは `その定理の真偽を疑うことは難しい'，`一般的な直感が本能的に受け入れるような定理である'，例えば `$2 + 2 = 4$'や`等腰三角形の底角は等しい'のような命題が真であることを指す。しかし，定理が明白であるという事実がその定理が真であることの証明を必要としない訳ではない．最も自信を持つ直感的な判断もしばしば誤りであることが判明するからである；また，定理が真であるとしても，それが同時に`明白'であるという事実はそれを証明することの理由を生じさせない．数学の目的は，ある前提から特定の結論を導くことを証明することであり，結論が前提ほど`明白'であっても，その証明の必要性と興味はしばしば損なわれることはない。}

しかし、時として（例えばここでは）「これはほとんど明白である」という表現は、その意味自体が異なる。我々は「一瞬の反芻により、読者は述べられた事実を真実だと確信し、また厳密な証明の概略も提示されるべきだ」という意味である。そしてしばしば、「この意味での明白」な主張に対しては、証明を省略してもよい。証明が欠けているからではなく、読者が容易に補完できるためであり、その証明を詳細に述べることは時間とスペースの無駄である。

読者はすぐに形成する議論の大まかな流れは次のようだ。「nが大きいとき、φ(n)はほぼaに等しく、ψ(n)はほぼbに等しいので、それらの和はほぼa+bに等しい」というものである。しかし、この議論を完全に形式化して述べる価値がある。

任意の正数$\DELTA$（例えば$0.001$, $0.0000001$, など）が与えられたとき、ある整数$n_{0}$を選び、式(1)
\[
|\phi(n) + \psi(n) - a - b| < \DELTA
\]
が$n \geq n_{0}$のときに満足することを示す必要がある。この条件は、第III章で証明された命題により、二つの数の和の絶対値はその絶対値の和以下であるというより一般的な原理から導かれる。よって
\[
|\phi(n) + \psi(n) - a - b| \leq |\phi(n) - a| + |\psi(n) - b|
\]
が成り立つ。したがって、式(2)
\[
|\phi(n) - a| + |\psi(n) - b| < \DELTA
\]
が$n \geq n_{0}$のときに満足すれば、要求される条件は満たされることがわかる。しかし、これは必ず成り立つ。なぜなら$\lim\phi(n) = a$であることから、定義により任意に小さな正数$\DELTA'$に対して$n \geq n_{1}$のとき$|\phi(n) - a| < \frac{1}{2}\DELTA$となるような整数$n_{1}$を選び得る。同様に、$\lim\psi(n) = b$であることから、任意に小さな正数$\DELTA'$に対して$n \geq n_{2}$のとき$|\psi(n) - b| < \frac{1}{2}\DELTA$となるような整数$n_{2}$を選び得る。そこで$n_{0}$を$n_{1}$と$n_{2}$のうち大きい方とする。すると、$n \geq n_{0}$のときには$|\phi(n) - a| < \frac{1}{2}\DELTA$かつ$|\psi(n) - b| < \frac{1}{2}\DELTA$となり、従って式(2)が満たされ、定理が証明される。

以下は簡潔に述べることができる。$\lim\phi(n) = a$および$\lim\psi(n) = b$であることから、任意に小さな正数$\DELTA'$に対して$n \geq n_{1}$のとき$|\phi(n) - a| < \frac{1}{2}\DELTA$かつ$n \geq n_{2}$のとき$|\psi(n) - b| < \frac{1}{2}\DELTA$となるような整数$n_{1}$と$n_{2}$を選び得る。そして、$n$が$n_{1}$または$n_{2}$以上であれば
\[
|\phi(n) + \psi(n) - a - b|
  \leq |\phi(n) - a| + |\psi(n) - b| < \DELTA;
\]
となり、従って
\[
\lim\{\phi(n) + \psi(n)\} = a + b.
\]

64. 定理Iの補助結果。読者は以下の補助結果を検証することに難儀しない筈である。

\par\medskip\noindent\textbf{【結果】}\par
\項{1.} $\phi(n)$ が極限に向かう一方、$\psi(n)$ が $+\infty$ または $-\infty$ に向かうか、有限に或は無限に振動する場合、$\phi(n) + \psi(n)$ は $\psi(n)$ の性質を示す。
\par\medskip

\par\medskip\noindent\textbf{【結果】}\par
\項{2.} $\phi(n) \to +\infty$ かつ $\psi(n) \to +\infty$ や有限に振動する場合、$\phi(n) + \psi(n) \to +\infty$。
\par\medskip

この声明において、$+\infty$ を $-\infty$ に変更しても構いません。

\par\medskip\noindent\textbf{【結果】}\par
\項{3.} $\phi(n) \to +\infty$ かつ $\psi(n) \to -\infty$ の場合、$\phi(n) + \psi(n)$ は極限に向かうか、$+\infty$ または $-\infty$ に向かうか、或は有限に或は無限に振動する可能性があります。
\par\medskip

\par\medskip\noindent\textbf{【注釈】}\par
これらの五つの可能性は順に (i) $\phi(n) = n$, $\psi(n) = -n$、(ii) $\phi(n) = n^2$, $\psi(n) = -n$、(iii) $\phi(n) = n$, $\psi(n) = -n^2$、(iv) $\phi(n) = n + (-1)^n$, $\psi(n) = -n$、(v) $\phi(n) = n^2 + (-1)^n n$, $\psi(n) = -n^2$ によって示されます。読者はそれぞれのケースについて追加の例を作成すべきです。
\par\medskip

\par\medskip\noindent\textbf{【結果】}\par
\項{4.} $\phi(n) \to +\infty$ かつ $\psi(n)$ が無限に振動する場合、$\phi(n) + \psi(n)$ は $+\infty$ に向かうか、或は有限に或は無限に振動しますが、極限に向かうことはありませんし、$-\infty$ に向かうことも、有限に振動することもありません。
\par\medskip

\par\medskip\noindent\textbf{【注釈】}\par
$\psi(n) = \{\phi(n) + \psi(n)\} - \phi(n)$；そして、$\phi(n) + \psi(n)$ が最後の三つのいずれかの方法で振る舞う場合、前段落の結果から $\psi(n) \to -\infty$ となりますが、これは事実と矛盾します。二つの可能なケースの例として (i) $\phi(n) = n^2$, $\psi(n) = (-1)^n n$、(ii) $\phi(n) = n$, $\psi(n) = (-1)^n n^2$ を考慮すべきです。また、$+\infty$ と $-\infty$ の符号は全体を通じて互換可能です。
\par\medskip

\par\medskip\noindent\textbf{【結果】}\par
\項{5.} $\phi(n)$ と $\psi(n)$ がどちらも有限に振動する場合、$\phi(n) + \psi(n)$ は極限に向かうか、或は有限に振動する必要があります。
\par\medskip

\par\medskip\noindent\textbf{【注釈】}\par
例として次のような関数を考える：
\[
\Itemp{(i)} \phi(n) = (-1)^n, \quad \psi(n) = (-1)^{n+1}, \qquad
\Itemp{(ii)} \phi(n) = \psi(n) = (-1)^n.
\]
\par\medskip

\par\medskip\noindent\textbf{【結果】}\par
\Item{6.} $\phi(n)$が有限に振動し、$\psi(n)$が無限に振動するとき、$\phi(n) + \psi(n)$は無限に振動する。
\par\medskip

\par\medskip\noindent\textbf{【注釈】}\par
$\phi(n)$の絶対値は一定の定数、例えば$K$未満である。一方で、$\psi(n)$が無限に振動することから、その絶対値は任意に大きい数（例えば$10K, 100K,\dots$）を超える値をとる。したがって、$\phi(n) + \psi(n)$も任意に大きい数（例えば$9K, 99K,\dots$）を超える値をとる。よって、$\phi(n) + \psi(n)$は無限大または無限小に収束するか、有限または無限に振動する。しかし、もし無限大に収束すると仮定すると、
\[
\psi(n) = \{\phi(n) + \psi(n)\} - \phi(n)
\]
も同様に無限大に収束することになる。したがって、$\phi(n) + \psi(n)$は無限大または無限小に収束することができず、有限または無限に振動する。
\par\medskip

\par\medskip\noindent\textbf{【結果】}\par
\Item{7.} $\phi(n)$と$\psi(n)$がともに無限に振動するとき、$\phi(n) + \psi(n)$は極限に収束するか、無限大または無限小に収束するか、有限または無限に振動する可能性がある。
\par\medskip

\par\medskip\noindent\textbf{【注釈】}\par
例えば、$\phi(n) = (-1)^n n$であり、$\psi(n)$がそれぞれ$(-1)^{n+1} n, \{1 + (-1)^{n+1}\} n, -\{1 + (-1)^n\} n, (-1)^{n+1}(n + 1), (-1)^n n$のいずれかの場合を考える。これにより、すべての可能性（極限に収束、無限大または無限小に収束、有限または無限に振動）の例が得られる。
\par\medskip

The results 1--7 cover all the cases which are really distinct.
Before passing on to consider the product of two functions, we
may point out that the result of Theorem~I may be immediately
extended to the sum of three or more functions which tend to
limits as $n\to\infty$.

65. 二つの関数の挙動が既知であるときのその積の挙動に関する定理

現在、二つの関数の積について類似した一連の定理を証明することができます。主な結果は以下の通りです。

\par\medskip\noindent\textbf{【定理】}\par[II]
$\lim \phi(n) = a$ および $\lim \psi(n) = b$ のとき，
\[
\lim \phi(n)\psi(n) = ab.
\]
\par\medskip

ここで
\[
\phi(n) = a + \phi_{1}(n),\quad
\psi(n) = b + \psi_{1}(n),
\]
とし，$\lim \phi_{1}(n) = 0$ および $\lim \psi_{1}(n) = 0$ である。すると，
\[
\phi(n)\psi(n) = ab + a\psi_{1}(n) + b\phi_{1}(n) + \phi_{1}(n)\psi_{1}(n).
\]
したがって，$\phi(n)\psi(n) - ab$ の差の数値は，$a\psi_{1}(n)$, $b\phi_{1}(n)$, $\phi_{1}(n)\psi_{1}(n)$ の数値の和よりも大きいとは限らない。これより
\[
\lim \{\phi(n)\psi(n) - ab\} = 0,
\]
が成り立つので，定理は証明された。

以下は厳密に形式的な証明である。以下の不等式が成立する。
\[
|\phi(n)\psi(n) - ab|
  \leq |a\psi_{1}(n)| + |b\phi_{1}(n)| + |\phi_{1}(n)| |\psi_{1}(n)|.
\]
$a$ や $b$ の少なくとも一方が零でない場合，$\DELTA > 0$ を選んで
\[
|\phi_{1}(n)| < \frac{1}{3}\frac{\DELTA}{|b|},\quad
|\psi_{1}(n)| < \frac{1}{3}\frac{\DELTA}{|a|}
\]
を満たす $n_{0}$ を選べる。すると，
\[
|\phi(n)\psi(n) - ab|
  < \frac{1}{3}\DELTA
  + \frac{1}{3}\DELTA
  + \frac{1}{9}\frac{\DELTA^{2}}{|a||b|}
\]
となり，$\DELTA < \frac{1}{3}|a||b|$ のとき，これは $\DELTA$ より小さい。すなわち，$n \geq n_{0}$ のとき $|\phi(n)\psi(n) - ab| < \DELTA$ となるような $n_{0}$ を選べる。したがって定理は成り立つ。読者は少なくとも一方が零である $a$ や $b$ の場合の証明を自行しなさい。

We hardly need to point out thatこの定理は、Theorem~Iと同様に任意の数目的関数の積にも即座に拡張できる。また、和に関するSecNo[§]64で述べた補助的定理と類似した一連の補助的定理がある。

$\phi(n)$が$n$に向かって無限大に近づくときには六つの異なる方法で振る舞う可能性がある。それは（1）他の限界値へと収束する、（2）0に収束する、（3a）正の無限大へと収束する、（3b）負の無限大へと収束する、（4）有限な振動を示す、（5）無限大の振動を示す。通常、規則として（3a）と（3b）は別々に考慮する必要がない。一方の場合から他方の場合への符号の変更によって結果が導き出されるためである。

補足：
これらの補助的定理を詳細に述べることは、私たちが許容できるスペースを超えてしまう。ここでは二つの例を選んで示し、それらの検証は読者に任せることにする。彼は自己でいくつかの残りの定理を形式化する経験を得られるだろう。

結果：
(i) $\phi(n)$が$+\infty$へと収束し、$\psi(n)$が有限な振動を示すとき、$\phi(n)\psi(n)$は$+\infty$や$-\infty$へと収束するか、または無限大の振動を示さなければならない。

これらの三つの可能性の例として、$\phi(n) = n$と$\psi(n)$が$2 + (-1)^{n}$、$-2 - (-1)^{n}$、$(-1)^{n}$のいずれかを選ぶことができる。

結果：
(ii) $\phi(n)$と$\psi(n)$が有限な振動を示すとき、$\phi(n)\psi(n)$は収束する限界値（これは0である可能性がある）または有限な振動を示さなければならない。

例として、(a) $\phi(n) = \psi(n) = (-1)^{n}$、(b) $\phi(n) = 1 + (-1)^{n}$、$\psi(n) = 1 - (-1)^{n}$、(c) $\phi(n) = \cos\frac{1}{3}n\pi$、$\psi(n) = \sin\tfrac{1}{3} n\pi$を選ぶことができる。

特定のケースとして重要である定理IIの特殊な場合は、$\psi(n)$が定数の場合です。この定理は単純に、$\lim k\phi(n) = ka$ が $\lim\phi(n) = a$ のとき成立することを主張します。また、$\phi(n) \to +\infty$ のとき $k\phi(n) \to +\infty$ または $k\phi(n) \to -\infty$ であり、これは $k$が正のときまたは負のときによります（ただし $k = 0$ の場合、明らかに $k\phi(n) = 0$ で $\lim k\phi(n) = 0$）。そして、$\phi(n)$が有限または無限に振動するとき、$k\phi(n)$も同様の振動を示します（ただし $k = 0$ の場合を除く）。

66．二つの関数の性質が既知であるときの差分や商の挙動について。二つの与えられた関数の差分に関する類似の定理もあります。これは前段落から明らかな推論です。商
\[
\frac{\phi(n)}{\psi(n)}
\]
を扱うために、以下の定理から始めます。

定理III．
$\lim\phi(n) = a$ であり $a \neq 0$ のとき
\[
\lim\frac{1}{\phi(n)} = \frac{1}{a}
\]
が成立します。

$\phi(n) = a + \phi_{1}(n)$ と置きますので $\lim\phi_{1}(n) = 0$。すると
\[
\left|\frac{1}{\phi(n)} - \frac{1}{a}\right| = \frac{|\phi_{1}(n)|}{|a| |a + \phi_{1}(n)|}
\]
であり、$\lim\phi_{1}(n) = 0$ より、任意の数 $\DELTA$ より小さい値を $n \geq n_{0}$ のとき選べます。

定理IIと定理IIIから、商に関する主な定理を即座に導くことができます。すなわち、

定理IV．
$\lim\phi(n) = a$ かつ $\lim\psi(n) = b$ であり $b \neq 0$ のとき
\[
\lim\frac{\phi(n)}{\psi(n)} = \frac{a}{b}
\]
が成立します。

読者は、定理IIIとIVに対応する補助的な定理について述べる、証明し、例を用いて説明することにより、有益な経験を得ることでしょう。

\Paragraph{67.}
\par\medskip\noindent\textbf{【定理】}\par[V.]
任意の有理関数 $R\{\phi(n), \psi(n), \chi(n), \dots\}$ について、$\phi(n)$, $\psi(n)$, $\chi(n)$, \dots の多項式 $P$ と $Q$ を用いて
\[
R\{\phi(n), \psi(n), \chi(n), \dots\} = P\{\phi(n), \psi(n), \chi(n), \dots\}/Q\{\phi(n), \psi(n), \chi(n), \dots\}
\]
と表すことができる任意の関数を考え、
\[
\lim_{n \to \infty} \phi(n) = a,\quad
\lim_{n \to \infty} \psi(n) = b,\quad
\lim_{n \to \infty} \chi(n) = c, \dots,
\]
かつ
\[
Q(a, b, c, \dots) \neq 0;
\]
ならば
\[
\lim_{n \to \infty} R\{\phi(n), \psi(n), \chi(n), \dots\} = R(a, b, c, \dots).
\]
\par\medskip

$P$ は有限個の項の和で表され、各項は
\[
A\{\phi(n)\}^{p} \{\psi(n)\}^{q} \dots,
\]
と表される形であり、ここで $A$ は定数、$p$, $q$, \dots は正の整数である。この項は、定理II（あるいはその明明白白な拡張）により $\lim_{n \to \infty} Aa^{p}b^{q}\dots$ に収束し、したがって $P$ も同様の拡張により $\lim_{n \to \infty} P(a, b, c, \dots)$ に収束する。同様に $Q$ も $\lim_{n \to \infty} Q(a, b, c, \dots)$ に収束し、定理IVから結論が導かれる。

\Paragraph{68.}
先の一般定理は以下の重要な特殊問題に適用できる：すなわち、$n$ の任意の有理関数 $S(n)$ の振る舞いは何であるか？（
\[
S(n) = \frac{a_{0}n^{p} + a_{1}n^{p-1} + \dots + a_{p}}
            {b_{0}n^{q} + b_{1}n^{q-1} + \dots + b_{q}},
\]
$n$ が無限大に収束するとき）
（自然には $a_{0}$ や $b_{0}$ のいずれも零でないという仮定を加える）。

$S(n)$ を変形して以下のように表すことで定理の適用が可能となる：
\[
n^{p-q}\left\{
  \biggl(a_{0} + \frac{a_{1}}{n} + \dots + \frac{a_{p}}{n^{p}}\biggr)\bigg/
  \biggl(b_{0} + \frac{b_{1}}{n} + \dots + \frac{b_{q}}{n^{q}}\biggr)
\right\}.
\]
大括弧内の関数は $R\{\phi(n)\}$ の形であり、$\phi(n) = 1/n$ であるため、$n$ が無限に近づくときには $R(0) = a_{0}/b_{0}$ に収束する。ここで $p < q$ の場合、$n^{p-q} \to 0$；$p = q$ の場合、$n^{p-q} = 1$ かつ $n^{p-q} \to 1$；$p > q$ の場合、$n^{p-q} \to +\infty$。したがって、定理IIにより以下の極限が得られる：
\begin{gather*}
\lim S(n) = 0\quad (p < q), \\
\lim S(n) = a_{0}/b_{0}\quad  (p = q), \\
S(n) \to +\infty\quad (p > q,\ a_{0}/b_{0} 正のとき), \\
S(n) \to -\infty\quad (p > q,\ a_{0}/b_{0} 負のとき).
\end{gather*}

以下は例題である：
\Item{1.} $n$ が無限大に近づくとき、以下の関数の振る舞いを調べよ：
\[
\left(\frac{n - 1}{n + 1}\right)^{2},\quad
(-1)^{n} \left(\frac{n - 1}{n + 1}\right)^{2},\quad
\frac{n^{2} + 1}{n},\quad
(-1)^{n} \frac{n^{2} + 1}{n}.
\]

\Item{2.} 以下の関数のうち、無限大に近づくとき極限値を持つものは何か：
\begin{gather*}
1/(\cos^{2}\tfrac{1}{2}n\pi + n\sin^{2}\tfrac{1}{2}n\pi),\quad
1/\{n(\cos^{2}\tfrac{1}{2}n\pi + n\sin^{2}\tfrac{1}{2}n\pi)\}, \\
  (n\cos^{2}\tfrac{1}{2}n\pi +  \sin^{2}\tfrac{1}{2}n\pi)/
\{n(\cos^{2}\tfrac{1}{2}n\pi + n\sin^{2}\tfrac{1}{2}n\pi)\}.
\end{gather*}

\Item{3.} 一般的な $n$ の有理関数を $S(n)$ とし、以下の極限を示せ：
\[
\lim\frac{S(n + 1)}{S(n)} = 1,\quad
\lim\frac{S\{n + (1/n)\}}{S(n)} = 1.
\]

69. $n$ に従って定常的に増加する関数

特別で特に重要な関数のクラスが、$n$ に向かって変化する方向が常に同じである、すなわち $n$ が増加するにつれて常に増加（または常に減少）する関数を形成している。$\phi(n)$ が常に減少すれば $-\phi(n)$ は常に増加するため、両種類の関数を別々に考慮する必要はない；一方の種類に対する証明は他方にも即座に拡張できる。

定義
関数 $\phi(n)$ が $n$ に従って定常的に増加するとするには、$\phi(n + 1) \geq \phi(n)$ がすべての $n$ の値に対して成り立つこととする。

注意すべきは、$\phi(n)$ がいくつかの $n$ の値で同じ値を取ることを排除していない点である；排除するのは可能性のある減少のみである。例えば
\[
\phi(n) = 2n + (-1)^{n}
\]
という関数があり、その $n = 0, 1, 2, 3, 4,\dots$ の値は
\[
1,\ 1,\ 5,\ 5,\ 9,\ 9,\ \dots
\]
であり、この関数は $n$ に従って定常的に増加するとされる。実際の定義では、$\phi(n)$ が某一定値以降で恒等的に一定であるような関数も含めることができる；例えば $\phi(n) = 1$ は我々の定義によれば定常的に増加する。しかし、これらの関数は極めて特殊なものであり、また $n$ に向かって無限大に近づくときの振る舞いについて疑問が生じることはないとすれば、この定義上の不整合は大きな欠点とはならない。

このクラスの関数に関する非常に重要な定理がある。

定理
$\phi(n)$ が $n$ に従って定常的に増加する場合、次のいずれかが成り立つ。
(i) $\phi(n)$ は $n$ に向かって無限大に近づくときには極限を持つ。
(ii) $\phi(n)\to +\infty$。

つまり、一般的な関数の振る舞いについて考えると通常5つの選択肢があるが、この特殊な種類の関数に関しては2つだけである。

この定理は、デデキントの定理（§17）の単純な推論である。実数ξを2つのクラスLとRに分類し、φ(n)が某nに対して≧ξとなる場合ξをクラスLに入れ、それ以外の場合ξをクラスRに入れる。

クラスLは確かに存在する；クラスRは存在するか否か未定である。もしRが存在しないならば、任意の数Δ何管いかに大きいかに関わらず、ある十分に大きいnに対してφ(n) > Δであり、従って
\[
\phi(n) \to +\infty.
\]

一方で、クラスRが存在する場合、クラスLとRは実数に対する§17の意味での分割を形成する。このとき、該当する数aを定め、任意の正数Δを考える。すると全てのnに対してφ(n) < a + Δであり、従ってΔが任意である限りφ(n) ≤ aである。一方であるnに対してφ(n) > a - Δであり、従って十分に大きいnに対して常に成り立つ。したがって
\[
a - \Delta < \phi(n) \leq a
\]
が全ての十分に大きいnに対して成立し、すなわち
\[
\phi(n) \to a.
\]

補足：
一般に任意のnに対してφ(n) < aであることに注意すべきである。なぜなら、もし某nに対してφ(n) = aならば、全てのより大きなnに対してφ(n)もまたaと等しいからである。したがってφ(n)はaと等しくなることはなく、値が最終的にすべて同じ場合を除いてLには最大要素がない。

補則：
1. φ(n)がnと共に単調に増加するならば、ある数Kが存在するか否かにより、その極限に収束するか+∞に収束するか決まる。
2. φ(n)がnと共に単調に増加し、全てのnに対してφ(n) < Kであるならば、φ(n)は極限を持つし、この極限はK以下である。

{\Loosen 限界が~$K$ と等しい場合も注意しなければならない：たとえば \eg\
$\phi(n) = 3 - (1/n)$ のとき、すべての $\phi(n)$ の値は~$3$ より小さいが、その限界は~$3$ に等しい。}

\par\medskip\noindent\textbf{【系】}\par[3.]
$n$ が増加するにつれて $\phi(n)$ が安定的に増加し、かつ限界に向かう場合、
すべての $n$ の値に対して
\[
\phi(n) \leq \lim\phi(n)
\]
である。
\par\medskip

読者は、$\phi(n)$ が $n$ が増加するにつれて減少する場合に対応する定理と補則を自分で書き出してください。

\Paragraph{70.} この定理の大きな重要性は、与えられた $n$ の関数が $n \to \infty$ のとき限界に向かうかどうかを判断できる手段を与えてくれることにある。これまでに得られていなかったものである。もし限界があるとすればそれが何であるかを予測したり推測したりする必要なく、
\[
|\phi(n) - l| < \DELTA\quad (n \geq n_{0}):
\]
というテストを使用できる。たとえば $\phi(n) = 1/n$ の場合、限界は明らかに零である。しかし、$\phi(n) = \left(1 + \frac{1}{n}\right)^{n}$ が限界に向かうかどうかを決定する必要がある場合、その限界、もしあれば、何であるかは明らかでない：そして上記のテストは $l$ を含むため、直接的に $l$ が存在するかどうかを判断するために使用することは困難である。

\par\medskip\noindent\textbf{【注釈】}\par
もちろん、このテストは間接的に使用して、$l$ が存在しないことを反証することによって使用できる。たとえば $\phi(n) = (-1)^{n}$ の場合、$l$ は~$1$ でありまた~$-1$ であることが明らかに不可能であるから、$l$ が存在しないことが示される。

\Paragraph{71.} ウェーラーストラスの定理（第19節）の別の証明。第69節の結果により、第19節で証明された重要な定理に対する別の証明を与えることができる。

如果我達把 PQ 分成两半，至少有一半包含无穷多个 S 的点。我達选择包含これ些点的那一半，或者如果两辺都包含，則选择左辺那一半；并用 \(P_1Q_1\) 表示所选的那半段（图 28）。
もし PQ を二等分すると、少なくとも一方には S の無数の点が含まれている。その一方を選び、両方にある場合は左側の半分を選ぶ。このようにして、次の列挙された区間を定義できる：
\[
PQ,\quad P_1Q_1,\quad P_2Q_2,\quad P_3Q_3,\ \dots
\]
各区間は前段の半分であり、それぞれには無数の S の点が含まれている。

点 \(P, P_1, P_2, \dots\) は左から右へと安定的に進行し、したがって \(P_n\) は極限位置 \(T\) に収束する。同様に \(Q_n\) も極限位置 \(T'\) に収束する。
しかし \(TT'\) は明らかに \(P_nQ_n\) より小さいことがわかる；そして \(P_nQ_n\) は \(PQ/2^n\) であり、0 に収束する。したがって \(T' = T\) であり、\(P_n\) と \(Q_n\) ともに \(T\) に収束する。

すると \(T\) は S の累積点である。もし \(\xi\) をその座標とするなら、任意の区間 \((\xi - \Delta, \xi + \Delta)\) または \([\xi - \Delta, \xi + \Delta]\) を考える。十分大きな \(n\) に対して、\(P_nQ_n\) はこの区間に完全に含まれる。
したがって
\[
(\xi - \Delta, \xi + \Delta)
\]
または
\[
[\xi - \Delta, \xi + \Delta]
\]
には無数の S の点が含まれている。

72. \(x^n\) の \(n\) が無限大に近づくときの極限

結果を用いて \SecNo[§]{69} で特に重要な場合、すなわち \(\phi(n) = x^n\) に対して適用します。もし \(x = 1\) であれば, \(\phi(n) = 1\), \(\lim \phi(n) = 1\) であり, \(x = 0\) の場合は \(\phi(n) = 0\), \(\lim \phi(n) = 0\) となるので,これらの特別な場合については考慮する必要はありません。
ページ分割

まず, \(x\) が正の場合を考えます。このとき, \(\phi(n + 1) = x \phi(n)\) より, \(x > 1\) の場合は \(\phi(n)\) は \(n\) とともに増加し, \(x < 1\) の場合は \(n\) が増加するにつれて減少します。

もし \(x > 1\) であれば, \(x^n\) は明らかに \(1\) より大きい極限へと収束するか, \(+\infty\) に収束する必要があります。極限を \(l\) とする場合, \(\lim \phi(n + 1) = \lim \phi(n) = l\) (例題 xxv.7) ですが,
\[
\lim \phi(n + 1) = \lim x \phi(n) = x \lim \phi(n) = xl
\]
であり, よって \(l = xl\)。そして, \(x\) と \(l\) はともに \(1\) より大きいので,これは不可能です。したがって,
\[
x^n \to +\infty \quad (x > 1)
\]

補足

例として, ニュートンの二項定理を用いて, \(\delta\) が正で \(x = 1 + \delta\) のとき, \(x^n > 1 + n\delta\) であることを示し, \(x^n \to +\infty\) を証明することができます。

一方, \(x < 1\) の場合, \(x^n\) は減少する関数であり,したがって極限へと収束するか, \(-\infty\) に収束します。\(x^n\) が正であるため,二つの選択肢のうち後者は無視できます。よって, \(\lim x^n = l\) とすると,上記と同様に \(l = xl\) であり, \(l\) はゼロでなければなりません。したがって,
\[
\lim x^n = 0 \quad (0 < x < 1)
\]

補足

例として, 前の例と同様にして, \(0 < x < 1\) のとき \((1/x)^n\) は \(+\infty\) に収束することを証明し,これより \(x^n\) は \(0\) に収束することを導きます。

\[
x
\]が負の場合を最後に考える。\[
-1 < x < 0
\]かつ\[
x = -y
\]のとき、\[
0 < y < 1
\]であるから、前節より\[
\lim y^n = 0
\]であり、従って\[
\lim x^n = 0
\]である。\[
x = -1
\]のとき、明らかに\[
x^n
\]は振動し、\[
-1, 1
\]を交互に取る。最後に\[
x < -1
\]かつ\[
x = -y
\]のとき、\[
y > 1
\]であり、\[
y^n
\]は\[
+\infty
\]に収束するから、\[
x^n
\]は正でも負でも絶対値が任意に大きい数を超える値を取る。従って\[
x^n
\]は無限に振動する。要するに：

\begin{alignat*}{2}
\&\phi(n) = x^n \to +\infty \&\&(x > 1),\\
\&\lim \phi(n) = 1            \&\&(x = 1),\\
\&\lim \phi(n) = 0            \&\&(-1 < x < 1),\\
\&\text{$\phi(n)$ \emph{有限に振動}} \&\&(x = -1),\\
\&\text{$\phi(n)$ \emph{無限に振動}}\qquad \&\&(x < -1).
\end{alignat*}

以下は例題である。

\par\medskip\noindent\textbf{【例 XXVII.】}\par
\Item{1.} \[
\phi(n)
\]が正であり、\[
\phi(n + 1) > K \phi(n)
\]、ここで\[
K > 1
\]で全ての\[
n
\]について成立するとき、\[
\phi(n) \to +\infty
\]である。
\footnotetext
  {これらの例題は特に重要であり、後述する文书中にもいくつか用いられるので、十分に研究すべきである。}
\PageSep{136}

\[
\phi(n) > K\phi(n - 1) > K^2\phi(n - 2) \dots > K^{n-1}\phi(1),
\]
から\[
K^n \to \infty
\]であることから結論がすぐに従う。

\Item{2.} 上記の条件は全ての\[
n \geq n_0
\]で成立する場合でも同じ結果が成り立つ。

\Item{3.} \[
\phi(n)
\]が正であり、\[
\phi(n + 1) < K\phi(n)
\]、ここで\[
0 < K < 1
\]であるとき、\[
\lim \phi(n) = 0
\]である。この結果も全ての\[
n \geq n_0
\]で条件が成立する場合でも成り立つ。

\Item{4.} \[
|\phi(n + 1)| < K|\phi(n)|
\]が全ての\[
n \geq n_0
\]で\[
0 < K < 1
\]であるとき、\[
\lim \phi(n) = 0
\]である。

\Item{5.} \[
\phi(n)
\]が正であり、\[
\lim \{\phi(n + 1)\}/\{\phi(n)\} = l > 1
\]であるとき、\[
\phi(n) \to +\infty
\]である。

[私たちは、\( n \geq n_{0} \)のとき\(\{\phi(n + 1)\}/\{\phi(n)\} > K > 1\)となるような\( n_{0} \)を定めることが可能です：例えば、\( K \)は1と\( l \)の中間値として取ることができます。その後、例題1を適用します.]

6. \(\lim\{\phi(n + 1)\}/\{\phi(n)\} = l\)で、\( l \)が単位より小さいとき、\(\lim\phi(n) = 0\)である。[これは例題4から導けます。例題5は例題1から導けます.]

7. \( n \to \infty \)のとき、\( r \)を任意の正整数として\(\phi(n) = n^{r}x^{n}\)の振る舞いを求めよ。

\( x = 0 \)の場合、すべての\( n \)に対して\(\phi(n) = 0\)であり、\(\phi(n) \to 0\)である。他のすべての場合において
\[
\frac{\phi(n + 1)}{\phi(n)} = \left(\frac{n + 1}{n}\right)^{r}x \to x.
\]
まず\( x \)が正のときを考える。この場合、\( x > 1 \)のときは\(\phi(n) \to +\infty\)（例題5）、\( x < 1 \)のときは\(\phi(n) \to 0\)（例題6）である。\( x = 1 \)のとき、\(\phi(n) = n^{r} \to +\infty\)である。次に\( x \)が負の場合を考える。この場合、\(|\phi(n)| = n^{r}|x|^{n}\)は\(|x| \geq 1\)のときは\(+\infty\)へ収束し、\(|x| < 1\)のときは0に収束する。したがって\( x \leq -1 \)のとき\(\phi(n)\)は無限に振動し、\(-1 < x < 0 \)のときは\(\phi(n) \to 0\)である。

8. 同様にして\( n^{-r}x^{n} \)を議論する。[結果は同じで、\( x = 1 \)または\(-1\)の場合を除き、\(\phi(n) \to 0\)である]

9. \( k \)をすべての実数値、およびすべての正負整数値として、\(\phi(n) = n^{k}x^{n}\)の\( n \to \infty \)における振る舞いを表にまとめよ。

読者は観察することができるが、\( k \)の値は特別な場合である\( x = 1 \)または\(-1\)以外では重要ではない。任意の正負でも\(\lim\{(n + 1)/n\}^{k} = 1\)であることから、比\(\phi(n + 1)/\phi(n)\)の極限は\( x \)のみに依存し、一般には因子\( x^{n} \)が主導的である。因子\( n^{k} \)は数値的に1と等しいときだけ顕著となる。]

\item 10. \( x \)が正のとき、\(\sqrt[n]{x} \to 1\) であることを証明せよ。［例えば、\( x > 1 \)とすると、\( x \)，\(\sqrt{x}\)，\(\sqrt{x}\)，…は減少する数列であり、\(\sqrt[n]{x} > 1\) がすべての \( n \) の値に対して成り立つ。したがって、\(\sqrt[n]{x} \to l\) で、\( l \geq 1 \)である。しかし、もし \( l > 1 \) であれば、十分大きな \( n \) の値に対して \(\sqrt[n]{x} > l\) または \( x > l^{n} \) となり、また \( l^{n} \to +\infty \) であるからこれは不可能である。］

\item 11. \(\sqrt[n]{n} \to 1\) であることを証明せよ。［例えば、\(\sqrt[n+1]{n + 1} < \sqrt[n]{n}\) は \((n + 1)^{n} < n^{n+1}\) や \(\left\{1 + \frac{1}{n}\right\}^{n} < n\) から成り立つ。これは \( n \geq 3 \) のとき確実に満たされる（73節参照）。したがって、\( n \) が 3 から増加するにつれて \(\sqrt[n]{n}\) は減少し、常に 1 より大きいので、その極限値は 1 より大きくまたは等しい。しかし、もし \(\sqrt[n]{n} \to l\) で \( l > 1 \) であれば、\( n > l^{n} \) となり、これは十分大きな \( n \) の値に対して明らかに成り立たない（練習問題 7, 8 参照）。］

\item 12. \(\sqrt[n]{n!} \to +\infty\) であることを証明せよ。［どんなに大きい \(\Delta\) も存在しても、十分大きな \( n \) のとき \( n! > \Delta^{n} \) である。例えば、\( u_{n} = \frac{\Delta^{n}}{n!} \) とすると、\(\frac{u_{n+1}}{u_{n}} = \frac{\Delta}{n + 1}\) であり、これは \( n \to \infty \) のときゼロに収束するから、\( u_{n} \) も同様である（練習問題 6 参照）。］

\item 13. \(-1 < x < 1\) のとき
\[
u_{n} = \frac{m(m - 1) \dots (m - n + 1)}{n!} x^{n} = \binom{m}{n} x^{n}
\]
が \( n \to \infty \) のとき零に収束することを示せ。

［もし \( m \) が正の整数であれば、\( u_{n} = 0 \) は \( n > m \) のときである。それ以外の場合、
\[
\frac{u_{n+1}}{u_{n}} = \frac{m - n}{n + 1} x \to -x,
\]
\( x = 0 \) でない限りである。］

\paragraph{73. \(\left(1 + \frac{1}{n}\right)^{n}\) の極限。} \( \phi(n) = \left(1 + \frac{1}{n}\right)^{n} \) のとき、より難しくなった問題は69節の助けを借りて解くことができる。

二項定理より、

\begin{multline*}
\left(1 + \frac{1}{n}\right)^{n} = 1 + n \cdot \frac{1}{n} + \frac{n(n - 1)}{1 \cdot 2} \cdot \frac{1}{n^{2}} + \dots + \frac{n(n - 1) \cdots (n - n + 1)}{1 \cdot 2 \cdots n} \cdot \frac{1}{n^{n}} \\
= 1 + 1 + \frac{1}{1 \cdot 2} \left(1 - \frac{1}{n}\right) + \frac{1}{1 \cdot 2 \cdot 3} \left(1 - \frac{1}{n}\right) \left(1 - \frac{2}{n}\right) + \dots \\
+ \frac{1}{1 \cdot 2 \cdots n} \left(1 - \frac{1}{n}\right) \left(1 - \frac{2}{n}\right) \cdots \left(1 - \frac{n - 1}{n}\right).
\end{multline*}

この式の $(p + 1)$ 番目の項は、

\[
\frac{1}{1 \cdot 2 \cdots p} \left(1 - \frac{1}{n}\right) \left(1 - \frac{2}{n}\right) \cdots \left(1 - \frac{p - 1}{n}\right),
\]

正であり、$n$ の増加とともに単調增加関数である。項の数も $n$ の増加とともに増加する。したがって $\left(1 + \frac{1}{n}\right)^{n}$ は $n$ の増加とともに増加し、$n \to \infty$ のときには極限に収束または $+\infty$ に向かう。

しかし、

\begin{align*}
\left(1 + \frac{1}{n}\right)^{n} &< 1 + 1 + \frac{1}{1 \cdot 2} + \frac{1}{1 \cdot 2 \cdot 3} + \dots + \frac{1}{1 \cdot 2 \cdot 3 \cdots n} \\
&< 1 + 1 + \frac{1}{2} + \frac{1}{2^{2}} + \dots + \frac{1}{2^{n-1}} < 3.
\end{align*}

したがって $\left(1 + \frac{1}{n}\right)^{n}$ は $+\infty$ に向かうことはなく、
\[
\lim_{n \to \infty} \left(1 + \frac{1}{n}\right)^{n} = e,
\]
ここで $e$ は $2 < e \leq 3$ のような数である。

74. いくつかの代数学的補題について。
この段階で、後に使用するのに便利ないくつかの基本的な不等式を証明することに便益がある。

\item (i) $\alpha > 1$ かつ $r$ が正の整数であるとき、次の不等式は明らかである：
\[
r\alpha^{r} > \alpha^{r-1} + \alpha^{r-2} + \dots + 1.
\]
両辺に $\alpha - 1$ を掛けると、
\[
r\alpha^{r}(\alpha - 1) > \alpha^{r} - 1;
\]
さらに、それぞれの辺に $r(\alpha^{r} - 1)$ を加え、両辺を $r(r + 1)$ で割ると、
\[
\frac{\alpha^{r+1} - 1}{r + 1} > \frac{\alpha^{r} - 1}{r}\quad
(\alpha > 1).
\Tag{(1)}
\]
同様に、次の不等式も証明できる：
\[
\frac{1 - \beta^{r+1}}{r + 1} < \frac{1 - \beta^{r}}{r}\quad
(0 < \beta < 1).
\Tag{(2)}
\]

したがって、$r$ と $s$ が正の整数であり、かつ $r > s$ のとき、次の不等式は成り立つ：
\[
\frac{\alpha^{r} - 1}{r} > \frac{a^{s} - 1}{s},\quad
\frac{1 - \beta^{r}}{r} < \frac{1 - \beta^{s}}{s}.
\Tag{(3)}
\]
ここで $0 < \beta < 1 < \alpha$ である。特に、$s = 1$ のときには、次の不等式が成り立つ：
\[
\alpha^{r} - 1 > r(\alpha - 1),\quad
1 - \beta^{r} < r(1 - \beta).
\Tag{(4)}
\]

\item (ii) 不等式 \Eq{(3)} と \Eq{(4)} は、$r$ と $s$ が正の整数であるという仮定のもとで証明された。しかし、$r$ と $s$ が任意の正有理数である場合でもこれらの不等式は成り立つことが容易に確認できる。例えば、不等式 \Eq{(3)} の最初のものについて考えると、$r = a/b$, $s = c/d$ とし、$a$, $b$, $c$, $d$ を正の整数とするので $ad > bc$。$\alpha = \gamma^{bd}$ と置くと、不等式は次の形になる：
\[
(\gamma^{ad} - 1)/ad > (\gamma^{bc} - 1)/bc;
\]
これは既に証明済みである。同様の議論は残りの不等式にも適用でき、$s$ が $1$ より小さい正有理数であるとき、次の不等式も成り立つことが明らかである：
\[
\alpha^{s} - 1 < s(\alpha - 1),\quad
1 - \beta^{s} > s(1 - \beta).
\Tag{(5)}
\]

\Itemp{(iii)} 以下においては、すべての文字が正の数を表し、$r$ と $s$ は有理数であり、$\alpha$ と $r$ は $1$ より大きい、$\beta$ と $s$ は $1$ より小さいとする。$\alpha$ を $1/\beta$、$\beta$ を $1/\alpha$ として \Eq{(4)} を書き換えると、

\[
\alpha^{r} - 1 < r\alpha^{r-1}(\alpha - 1),\quad
1 - \beta^{r} > r\beta^{r-1}(1 - \beta).
\Tag{(6)}
\]

同様に、\Eq{(5)} より、

\[
\alpha^{s} - 1 > s\alpha^{s-1}(\alpha - 1),\quad
1 - \beta^{s} < s\beta^{s-1}(1 - \beta).
\Tag{(7)}
\]

\Eq{(4)} と \Eq{(6)} を組み合わせると、

\[
r\alpha^{r-1}(\alpha - 1) > \alpha^{r} - 1 > r(\alpha - 1).
\Tag{(8)}
\]

\PgSep{139}

$\alpha$ を $x/y$ として、

\[
rx^{r-1} (x - y) > x^{r} - y^{r} > ry^{r-1} (x - y)
\Tag{(9)}
\]

$x > y > 0$ のときが成り立つ。同様の議論を \Eq{(5)} と \Eq{(7)} に適用すると、

\[
sx^{s-1} (x - y) < x^{s} - y^{s} < sy^{s-1} (x - y)
\Tag{(10)}
\]

が得られる。

\par\medskip

\par\medskip\noindent\textbf{【例 XXVIII.】}\par
\Item{1.} \Eq{(9)} を $r = 2, 3$、\Eq{(10)} を $s = \frac{1}{2}, \frac{1}{3}$ に対して確認せよ。

\Item{2.} $y > x > 0$ のときも \Eq{(9)} と \Eq{(10)} が成り立つことを示せ。

\Item{3.} $r < 0$ のときも \Eq{(9)} が成り立つことを示せ。[ Chrystal の Algebra, vol. ii, p. 43--45 を参照]

\Item{4.} $\phi(n) \to l$、$l > 0$ として $n \to \infty$ のとき、任意の有理数 $k$ に対して $\phi^{k} \to l^{k}$。

[定理 III（66節）により $k > 0$ と仮定できる。また、ある値の $n$ から $\frac{1}{2}l < \phi < 2l$ であることが確実である。もし $k > 1$ のとき、

\[
k\phi^{k-1}(\phi - l) > \phi^{k} - l^{k} > kl^{k-1}(\phi - l)
\]

または

\[
kl^{k-1}(l - \phi) > l^{k} - \phi^{k} > k\phi^{k-1}(l - \phi)
\]

が成り立つ。したがって、$|\phi^{k} - l^{k}|$ と $|\phi - l|$ の比は $\frac{1}{2}l)^{k-1}$ から $(2l)^{k-1}$ までに含まれる。$0 < k < 1$ の場合も同様の証明が可能である。また、$l = 0$ でも $k > 0$ のとき結果は成り立つ。]

\par\medskip

5. \Exs{xxvii}.7, 8, 9 の結果を拡張し、$r$ や $k$ が任意の有理数である場合について述べよ。

---

75. $\sqrt[n]{x} - 1$ の極限 $n(\sqrt[n]{x} - 1)$。§74 の最初の不等式（3）において $r = 1/(n - 1), s = 1/n$ を代入すると
\[
(n - 1)(\sqrt[n-1]{\alpha} - 1) > n(\sqrt[n]{\alpha} - 1)
\]
が成り立つ。したがって、$\phi(n) = n(\sqrt[n]{\alpha} - 1)$ は $n$ の増加に伴い定常的に減少し、常に正である。従って $\phi(n)$ は $n \to \infty$ のとき極限 $l$ に収束し、$l \geq 0$。

また、§74 の最初の不等式（7）において $s = 1/n$ を代入すると
\[
n(\sqrt[n]{\alpha} - 1)
  > \sqrt[n]{\alpha}\left(1 - \frac{1}{\alpha}\right)
  > 1 - \frac{1}{\alpha}
\]
が成り立つ。従って $l \geq 1 - (1/\alpha) > 0$。したがって、$\alpha > 1$ のとき
\[
\lim_{n \to \infty} n(\sqrt[n]{\alpha} - 1) = f(\alpha),
\]
$f(\alpha) > 0$。

次に $\beta < 1$ の場合を考え、$\beta = 1/\alpha$ とおく。すると $n(\sqrt[n]{\beta} - 1) = -n(\sqrt{\alpha} - 1)/\sqrt[n]{\alpha}$ となる。ここで
$n(\sqrt[n]{\alpha} - 1) \to f(\alpha)$ であり、(Ex. xxvii.10)
\[
\sqrt[n]{\alpha} \to 1.
\]
従って $\beta = 1/\alpha < 1$ のとき
\[
n(\sqrt[n]{\beta} - 1) \to -f(\alpha).
\]

最後に $x = 1$ の場合、$n(\sqrt[n]{x} - 1) = 0$ である。

---

従って、以下の結果を得る：$\sqrt[n]{x} - 1$ の極限
\[
\lim n(\sqrt[n]{x} - 1)
\]
はすべての正の $x$ に対して $x$ の関数を定義する。この関数 $f(x)$ には以下の性質がある：
\[
f(1/x) = -f(x), \quad f(1) = 0,
\]
$x > 1$ のとき正であり、$x < 1$ のとき負である。後ほど、この関数をネイピアの対数と同一視できる。

---

例題。$f(xy) = f(x) + f(y)$ を証明せよ。[以下の式を使用する：
\[
f(xy) = \lim n(\sqrt[n]{xy} - 1)
      = \lim \{n(\sqrt[n]{x} - 1)\sqrt[n]{y} + n(\sqrt[n]{y} - 1)\}.
\]

76. 無限級数。ある $n$ の任意の関数 $u(n)$ を考え，$\nu = 1, 2, \dots, n$ に対して $u(\nu)$ の値を加えると，また $n$ の関数 $s(n)$ が得られます：
\[
s(n) = u(1) + u(2) + \dots + u(n).
\]
一般的には，記号を少し変更してこの式を
\[
s_{n} = u_{1} + u_{2} + \dots + u_{n}
\]
またはより簡略に
\[
s_{n} = \sum_{\nu=1}^{n} u_{\nu}
\]
と表すことが最も便利です。

もし $s_{n}$ が $n$ が無限大に向かうとき $s$ に収束すると仮定するなら，
\[
\lim_{n \to \infty} \sum_{\nu=1}^{n} u_{\nu} = s
\]
という式が得られます。この式は通常，次のいずれかの形で表されます：
\[
\sum_{\nu=1}^{\infty} u_{\nu} = s, \quad u_{1} + u_{2} + u_{3} + \dots = s,
\]
点が $u$ の項の無限続いていることを示しています。

上記の式の意味は，ざっくりと言えば，$u$ の項を逐次加えることで，その極限 $s$ に近づいていくということです。より正確には，任意の小さな正数 $\DELTA$ を選んだとき，$n_{0}(\DELTA)$ を選び，最初の $n_{0}(\DELTA)$ 個の項の和またはそれ以上の項の和が $s - \DELTA$ と $s + \DELTA$ の間にあるようにできます。すなわち，
\[
s - \DELTA < s_{n} < s + \DELTA
\]
が成り立つとき，$n \geq n_{0}(\DELTA)$ です。このような状況では，級数
\[
u_{1} + u_{2} + \dots
\]
を収束する無限級数と呼び，その極限 $s$ を級数の和または級数のすべての項の和と呼びます。

Thus、級数 $u_{1} + u_{2} + \dots$ が収束し、その和が $s$ である、または収束して $s$ に至る、あるいは単に $s$ に収束するとは、最初の $n$ 項の和 $s_{n} = u_{1} + u_{2} + \dots + u_{n}$ が $n \to \infty$ のとき、その極限に収束することを他言する方法である。このような無限級数を考える際には、この章の前半部で既に読者に馴染みのある概念を超える新しい概念は導入されない。実際、和 $s_{n}$ は関数 $\phi(n)$ の一種であり、我々がこれまで考慮してきたような形で表現されるものである。任意の関数 $\phi(n)$ をこのような形で表すことは可能であり、
\[
\phi(n) = \phi(1) + \{\phi(2) - \phi(1)\} + \dots + \{\phi(n) - \phi(n - 1)\}
\]
と書くことで達成される。そして、$\phi(n)$ が極限 $l$ に収束する（「傾き」ではなく）という表現も時として便利である。

もし $s_{n} \to +\infty$ または $s_{n} \to -\infty$ の場合、級数 $u_{1} + u_{2} + \dots$ を収束しない、または $+\infty$ に収束する、または $-\infty$ に収束するという。これらの表現も任意の関数 $\phi(n)$ に対して適用可能である：たとえば、$\phi(n) \to +\infty$ の場合、$\phi(n)$ が $+\infty$ に収束すると述べることができる。もし $s_{n}$ が極限に収束しないか、$+\infty$ または $-\infty$ に収束しない場合、それは有限または無限の振動を示す：この場合、級数 $u_{1} + u_{2} + \dots$ は有限または無限の振動を示すと述べる。

読者は注意されたいのは、「収束」および「振動」という言葉は異なる著者が異なる意味で使用することがあることである。ここでの用語の使用はブロミッチの『無限級数』に準拠している。ホブソンの『実変数関数論』では、系列が有限の振動を示す場合のみ「振動する」という表現を使用し、無限の振動を持つ系列は「収束しない」と分類される。多くの外国の著者は、「収束しない」という言葉を単に「収束しない」ことを意味することだけを使用している。

\#\#\# 77. 無限級数に関する一般の定理

無限級数を扱う際には、以下の一般の定理を使用することが頻繁にあります。

1. \(u_{1} + u_{2} + \dots\) が収束し、その和が \(s\) のとき、\(a + u_{1} + u_{2} + \dots\) も収束し、その和は \(a + s\) です。同様に、\(a + b + c + \dots + k + u_{1} + u_{2} + \dots\) も収束し、その和は \(a + b + c + \dots + k + s\) です。

2. \(u_{1} + u_{2} + \dots\) が収束し、その和が \(s\) のとき、\(u_{m+1} + u_{m+2} + \dots\) も収束し、その和は
\[
s - u_{1} - u_{2} - \dots - u_{m}
\]
です。

3. (1) または (2) で考慮された任意の級数が発散したり振動する場合、他の級数も同様です。

4. \(u_{1} + u_{2} + \dots\) が収束し、その和が \(s\) のとき、\(ku_{1} + ku_{2} + \dots\) も収束し、その和は \(ks\) です。

5. (4) で考慮された最初の級数が発散したり振動する場合、第2の級数も同様ですが、\(k = 0\) の場合を除きます。

6. \(u_{1} + u_{2} + \dots\) と \(v_{1} + v_{2} + \dots\) がともに収束するとき、\((u_{1} + v_{1}) + (u_{2} + v_{2}) + \dots\) も収束し、その和は最初の2つの級数の和です。

これらの定理の大半はほとんど明らかであり、定義から証明することができます。または、\SecNo[§§]{63}--\SecNo{66} の結果を級数の和 \(s_{n} = u_{1} + u_{2} + \dots + u_{n}\) に適用することでも証明できます。それらの後には、多少異なる性格を持つ定理が続きます。

7. \(u_{1} + u_{2} + \dots\) が収束するとき、\(\lim u_{n} = 0\) です。

\(u_{n} = s_{n} - s_{n-1}\) であり、\(s_{n}\) と \(s_{n-1}\) は同じ極限 \(s\) をもつ。したがって \(\lim u_{n} = s - s = 0\) です。

以下は証明の翻訳です：

\par\medskip\noindent\textbf{【注釈】}\par
読者は定理の逆が真であると錯覚するかもしれません。すなわち、$\lim u_n = 0$ならば級数 $\sum u_n$ は収束すると考えられるでしょう。しかし、これは例から明らかではありません。例えば、

\[
1 + \frac{1}{2} + \frac{1}{3} + \frac{1}{4} + \dots
\]

を考えます。このとき $u_n = \frac{1}{n}$ です。最初の四項の和は

\[
1 + \frac{1}{2} + \frac{1}{3} + \frac{1}{4}
  > 1 + \frac{1}{2} + \frac{2}{4} = 1 + \frac{1}{2} + \frac{1}{2}.
\]

次四項の和は $\frac{1}{5} + \frac{1}{6} + \frac{1}{7} + \frac{1}{8} > \frac{4}{8} = \frac{1}{2}$ です。さらに八項の和は $\frac{8}{16} = \frac{1}{2}$ より大きいです。同様に、最初の

\[
4 + 4 + 8 + 16 + \dots + 2^n
\]

項の和は

\[
2^{n+1}
\]

より大きいです。そして、

\[
2 + \frac{1}{2} + \frac{1}{2} + \frac{1}{2} + \dots + \frac{1}{2}
  = \frac{1}{2}(n + 3)
\]

であり、これは $n$ とともに無限大に増加します。したがって、級数は $+\infty$ に向かって発散します。
\par\medskip

\par\medskip\noindent\textbf{【結果】}\par
(8) もし $\sum u_1 + u_2 + u_3 + \dots$ が収束するならば、その項を任意の方法で括弧内にグループ化して新しい単項を作る任意の級数も収束し、両者の和は同じです。
\par\medskip

\par\medskip\noindent\textbf{【注釈】}\par
読者はこの定理の証明を自分で補うことができるでしょう。しかし、逆は真ではありません。例えば、

\[
1 - 1 + 1 - 1 + \dots
\]

は振動しますが、

\[
(1 - 1) + (1 - 1) + \dots
\]

は $0 + 0 + 0 + \dots$ となり、$0$ に収束します。
\par\medskip

\par\medskip\noindent\textbf{【結果】}\par
(9) もしすべての項 $u_n$ が正または零であるならば、級数 $\sum u_n$ は必ず収束するか、または $+\infty$ に向かって発散します。収束した場合、その和は正でなければなりません（すべての項が零である場合を除き、その場合は当然零です）。
\par\medskip

$s_n$ は $n$ の増加に従って増加する関数であり、第69節の定義によれば、$s_n$ に対して第69節の結果を適用することができます。

\par\medskip\noindent\textbf{【結果】}\par
(10) 各項 \(u_{n}\) が正または零であるとき、級数 \(\sum u_{n}\) が収束する必要十分条件は、その任意の項の和が定数 \(K\) より小さいことが可能であることです。そして、この定数 \(K\) を見つけることができれば、級数の和は \(K\) より大きくありません。
\par\medskip

これも \SecNo[§]{69} からすぐに分かります。しかし、各項 \(u_{n}\) が正であるという条件を満たさない場合、定理は成り立ちません。例えば
\[
1 - 1 + 1 - 1 + \dots
\]
は明らかに振動し、\(s_{n}\) は偶数番目では \(1\) で奇数番目では \(0\) となります。

\par\medskip\noindent\textbf{【結果】}\par
(11) 正または零である項からなる級数 \(u_{1} + u_{2} + \dots\)，\(v_{1} + v_{2} + \dots\) があり、第二の級数が収束し、かつ \(u_{n} \leq K v_{n}\)（\(K\) は定数）が全ての \(n\) に対して成り立つとき、第一の級数も収束し、その和は第二の級数の和 \(K\) 倍以上ではありません。
\par\medskip

もし \(v_{1} + v_{2} + \dots = t\) であれば、\(v_{1} + v_{2} + \dots + v_{n} \leq t\) が全ての \(n\) に対して成り立つので、\(u_{1} + u_{2} + \dots + u_{n} \leq Kt\) となり、定理を証明します。

\par\medskip\noindent\textbf{【結果】}\par
逆に、\(\sum u_{n}\) が発散し、かつ \(v_{n} \geq Ku_{n}\) のとき、\(\sum v_{n}\) も発散します。
\par\medskip

78. 無限等比級数について。無限等比級数を考えます。その一般項は \(u_{n} = r^{n-1}\) です。この場合
\[
s_{n} = 1 + r + r^{2} + \dots + r^{n-1} = (1 - r^{n})/(1 - r),
\]
ただし、\(r = 1\) の特別な場合では
\[
s_{n} = 1 + 1 + \dots + 1 = n.
\]
最後の場合は \(s_{n} \to +\infty\) です。一般の場合、\(s_{n}\) が収束するためにはかつのみに \(r^{n}\) が収束することが必要です。72節の結果を参照すると
\par\medskip\noindent\textbf{【結果】}\par
等比級数 \(1 + r + r^{2} + \dots\) が収束し、その和は \(1/(1 - r)\) であるためにはかつのみに \(-1 < r < 1\) です。
\par\medskip

もし$r \geq 1$ならば、$s_{n} \geq n$であり、したがって$s_{n} \to +\infty$；すなわち級数は$+\infty$に収束しない。もし$r = -1$ならば、$s_{n}$は$n$が奇数のとき$1$で、偶数のとき$0$である；すなわち$s_{n}$は有限に振動する。もし$r < -1$ならば、$s_{n}$は無限大に振動する。
したがって、要するに
\par\medskip\noindent\textbf{【結果】}\par
級数$1 + r + r^{2} + \dots$は$r \geq 1$のとき$+\infty$に収束し、$-1 < r < 1$のとき$\frac{1}{1 - r}$に収束し、$r = -1$のとき有限に振動し、$r < -1$のとき無限大に振動する。
\par\medskip

例題XXIX.
1. 再帰小数点. 一般的な無限等比級数は通常の再帰小数点で表される最も一般的な例である。例えば、小数$.217\DPmod{\dot{1}\dot{3}}{\Repeat{13}}$は算術の通常の規則によれば
\[
\frac{2}{10} + \frac{1}{10^{2}} + \frac{7}{10^{3}}
  + \frac{1}{10^{4}} + \frac{3}{10^{5}} + \frac{1}{10^{6}} + \frac{3}{10^{7}}
    + \dots
  = \frac{217}{1000}
    + \frac{13}{10^{5}} \bigg/ \left(1 - \frac{1}{10^{2}}\right)
  = \frac{2687}{12375}.
\]
読者は、この簡約化においてどのような一般定理が使用されたかを考えてみてほしい。

2. $a${1}a${2}\dots a\\textit{{m} \DPmod{\dot{\alpha}}{1}\alpha\\textit{{2}\dots\dot{\alpha}}{n}}
                            {\Repeat{\alpha${1}\alpha${2}\dots \alpha\\textit{{n}}}$は一般に
\[
\frac{a_{1}a_{2}\dots a_{m}\alpha_{1}\dots \alpha_{n} - a_{1}a_{2}\dots a_{n}}
         {99\dots 900\dots 0},
\]
の形で表される。分母には$n$個の$9$と$m$個の$0$が含まれている。

3. 純再帰小数点は常に、$2$または$5$を因数に持たない真分数に等しい。

4. $m$桁の非再帰的小数と$n$桁の再帰的小数を持つ小数は、$2^{m}$または$5^{m}$で割り切れるが、それ以上のべきでは割り切れない真分数に等しい。

\Item{5.} 3，4の逆も真です。$r = p/q$とします。まず$q$が10と互いに素である場合を考えます。すべての10の冪を$q$で割ると、得られる余りは至多$q$個しか異なりません。したがって、$n_{1}$と$n_{2}$という2つの数を見つけることが可能です。ここで$\DPtypo{n_{2} > n_{1}}{n_{1} > n_{2}}$であり、$10^{n_{1}}$と$10^{n_{2}}$が同じ余りをもたらします。したがって
\[10^{n_{1}} - 10^{n_{2}} = 10^{n_{2}}(10^{n_{1}-n_{2}} - 1)\]
は$q$で割り切れるため、$10^{n} - 1$（$n = n_{1} - n_{2}$）も$q$で割り切れます。したがって$r$は$\DPtypo{P/(10^{n} - 1)}{10^{n} - 1/P}$または
\[
\frac{P}{10^{n}} + \frac{P}{10^{2n}} + \dots
\]
という形で表すことができます。つまり、$n$桁の純粋循環小数として表現できます。一方で$q = 2^{\alpha}5^{\beta}Q$であり、$Q$が10と互いに素で、$\alpha$と$\beta$の大きい方を$m$とするとき、$10^{m}r$は10と互いに素な分母を持つため、整数と純粋循環小数の和として表すことができます。しかし、どんな$\mu$の値に対しても$10^{\mu}r$はそうではなく、従って$r$の小数表現にはちょうど$m$桁の非循環部分があります。

\Item{6.} 2--5の結果に加えて、i.3の結果も考慮に入れます。最後に
\[
.\DPmod{\dot{9}}{\Repeat{9}}
  = \frac{9}{10} + \frac{9}{10^{2}} + \frac{9}{10^{3}} + \dots
  = 1
\]
から、すべての終端小数は9で構成される純循環小数として表すこともできます。例えば、$.217 = .216\DPmod{\dot{9}}{\Repeat{9}}$です。したがって、すべての真分数は循環小数として表現でき、逆もまた真です。

7. 一般の小数について。無理数を非再現的小数として表す方法。どのような小数も、再現的であれ非再現であれ、0と1の間の特定の数に対応します。小数 \( .a_{1}a_{2}a_{3}a_{4}\dots \) は級数
\[
\frac{a_{1}}{10} + \frac{a_{2}}{10^{2}} + \frac{a_{3}}{10^{3}} + \dots.
\]
を表します。すべての桁 \( a_{r} \) は正であるため、この級数の最初の n 項の和 \( s_{n} \) は n の増加とともに増大し、確かに \( .\DPmod{\dot{9}}{\Repeat{9}} \) や1未満です。したがって \( s_{n} \) は0と1の間の極限に収束します。

さらに、2つの小数は同じ数に対応することができません（特別な例外を除く）。すなわち、\( .a_{1}a_{2}a_{3}\dots \), \( .b_{1}b_{2}b_{3}\dots \) が桁 \( a_{r-1} \), \( b_{r-1} \) まで一致し、かつ \( a_{r} > b_{r} \) の場合を考えます。このとき \( a_{r} \geq b_{r} + 1 > b_{r}.b_{r+1}b_{r+2} \dots \)（\( b_{r+1} \), \( b_{r+2} \), \dots がすべて9でない限り）であり、したがって
\[
.a_{1}a_{2} \dots a_{r}a_{r+1} \dots > .b_{1}b_{2} \dots b_{r}b_{r+1} \dots.
\]
これにより、有理分数を再現的小数として表す方法（練習問題2～6）は一意的であることがわかります。また、再現しない小数は0と1の間の無理数を表します。逆に、このような任意の数もそのような小数で表現できます。それは必ずしも
\[
0,\ 1/10;\quad 1/10,\ 2/10;\ \dots;\quad 9/10,\ 1.
\]
のいずれかの区間に存在します。もし \( r/10 \) と \( (r + 1)/10 \) の間にある場合、最初の桁は \( r \) です。この区間を10分割して第2桁を決定し、そのような方法で続けていけばよいですが（練習問題3,4）、小数は再現しません。例えば、通常の方法で得られる小数 \( 1.414\dots \) は再現しません。

\item 8. 小数 $.101\MS001\MS000\MS100\MS001\MS0\dots$ および $.202\MS002\MS000\MS200\MS002\MS0\dots$ において、2つの $1$ や $2$ の間のゼロの数は各段階で 1 増えることから、これらの小数は無理数を表します。

\item 9. 小数 $.111\MS010\MS100\MS010\MS10\dots$ において、$n$ 番目の桁が $n$ が素数のとき $1$ で、それ以外は $0$ である場合、この小数は無理数を表します。[素数の数は無限であり、小数は終了しません。また、再帰的ではないことも確認できます：もし再帰的であれば、$m$ と $p$ を定めることができ、$m$, $m + p$, $m + 2p$, $m + 3p$, \dots が全て素数であるという事実は矛盾します。この系列には $m + mp$ も含まれます。]

\par\medskip

\par\medskip\noindent\textbf{【例 XXX.】}\par
\item 1. 級数 $r^{m} + r^{m+1} + \dots$ は $-1 < r < 1$ のとき収束し、その和は $1/(1 - r) - 1 - r - \dots - r^{m-1}$ である（\S77, 式(2) 参照）。

\item 2. 級数 $r^{m} + r^{m+1} + \dots$ は $-1 < r < 1$ のとき収束し、その和は $r^{m}/(1 - r)$ である（\S77, 式(4) 参照）。問題 1 と 2 の結果が一致することを確認してください。

\item 3. 級数 $1 + 2r + 2r^{2} + \dots$ は収束し、その和は $(1 + r)/(1 - r)$ であることを証明します。証明の方法は次の通りです：(α) $-1 + 2(1 + r + r^{2} + \dots)$ の形に書き換えること、(β) $1 + 2(r + r^{2} + \dots)$ の形に書き換えること、(γ) $1 + r + r^{2} + \dots$ と $r + r^{2} + \dots$ の2つの級数を加えること。それぞれの場合で使用した定理 \S77 を明記してください。
\par\medskip

\item 4. 「算術」級数
\[
a + (a + b) + (a + 2b) + \cdots
\]
は、$a$ と $b$ がともに零でない限り収束しない。もし $b$ が零でないなら、この級数は符号の $b$ によって正無限大または負無限大に収束する。また、$b = 0$ の場合、この級数は符号の $a$ によって正無限大または負無限大に収束する。

\item 5. 級数
\[
(1 - r) + (r - r^{2}) + (r^{2} - r^{3}) + \cdots
\]
が収束するときの和は何か？ この級数は $-1 < r \leq 1$ のみで収束し、その和は $r = 1$ の場合を除き $1$ である。また、$r = 1$ の場合、その和は $0$ である。

\item 6. 級数
\[
r^{2} + \frac{r^{2}}{1 + r^{2}} + \frac{r^{2}}{(1 + r^{2})^{2}} + \cdots.
\]
は常に収束し、その和は $r = 0$ の場合を除き $1 + r^{2}$ である。

\item 7. $1 + r + r^{2} + \cdots$ が収束すると仮定すれば、その和は $\frac{1}{1 - r}$ であり、これは \SecNo[§]{77}, \Eq{(1)} および \Eq{(4)} を用いて証明できる。もし $1 + r + r^{2} + \cdots = s$ なら
\[
s = 1 + r(1 + r^{2} + \cdots) = 1 + rs.
\]

\item 8. 級数
\[
r + \frac{r}{1 + r} + \frac{r}{(1 + r)^{2}} + \cdots
\]
が収束するときの和を求めよ。この級数は $-1 < \frac{1}{1 + r} < 1$、すなわち $r < -2$ や $r > 0$ の場合にのみ収束し、その和は $1 + r$ である。また、$r = 0$ の場合も収束し、その和は $0$ である。

\item 9. 級数
\[
\begin{aligned}
& r - \frac{r}{1 + r} + \frac{r}{(1 + r)^{2}} - \cdots, \\
& r + \frac{r}{1 - r} + \frac{r}{(1 - r)^{2}} + \cdots, \\
& 1 - \frac{r}{1 + r} + \left(\frac{r}{1 + r}\right)^{2} - \cdots, \\
& 1 + \frac{r}{1 - r} + \left(\frac{r}{1 - r}\right)^{2} + \cdots.
\end{aligned}
\]
の和を求めよ。

\item 10. 級数
\[
\begin{aligned}
& (1 + r) + (r^{2} + r^{3}) + \cdots, \\
& (1 + r + r^{2}) + (r^{3} + r^{4} + r^{5}) + \cdots, \\
& 1 - 2r + r^{2} + r^{3} - 2r^{4} + r^{5} + \cdots, \\
& (1 - 2r + r^{2}) + (r^{3} - 2r^{4} + r^{5}) + \cdots,
\end{aligned}
\]
の収束性を考慮し、収束するときの和を求めよ。

\item 11. \(0 \leq a_{n} \leq 1\) のとき、級数 \(a_{0} + a_{1}r + a_{2}r^{2} + \dots\) は \(0 \leq r < 1\) に対して収束し、その和は \(1/(1 - r)\) より大きくなることはない。

\item 12. さらに級数 \(a_{0} + a_{1} + a_{2} + \dots\) が収束するとき、級数 \(a_{0} + a_{1}r + a_{2}r^{2} + \dots\) は \(0 \leq r \leq 1\) に対して収束し、その和は \(a_{0} + a_{1} + a_{2} + \dots\) より小さく、かつ \(1/(1 - r)\) より小さくなる。

\item 13. 級数
\[
1 + \frac{1}{1} + \frac{1}{1·2} + \frac{1}{1·2·3} + \dots
\]
は収束する。［\(1/(1·2 \dots n) \leq 1/2^{n-1}\)］

\item 14. 級数
\[
1 + \frac{1}{1·2} + \frac{1}{1·2·3·4} + \dots,\quad
\frac{1}{1} + \frac{1}{1·2·3} + \frac{1}{1·2·3·4·5} + \dots
\]
は収束する。

\item 15. 一般の調和級数
\[
\frac{1}{a} + \frac{1}{a + b} + \frac{1}{a + 2b} + \dots,
\]
で \(a\) と \(b\) が正のとき、この級数は \(+\infty\) に発散する。

［\(u_{n} = 1/(a + nb) > 1/\{n(a + b)\}\)。これと \(1 + \frac{1}{2} + \frac{1}{3} + \dots\) を比較せよ］

\item 16. 級数
\[
(u_{0} - u_{1}) + (u_{1} - u_{2}) + (u_{2} - u_{3}) + \dots
\]
は \(u_{n}\) が無限大に近づくときのみ収束する。

\item 17. 級数 \(u_{1} + u_{2} + u_{3} + \dots\) が発散するとき、その項を任意の方法で括弧内にグループ化して形成した任意の級数も発散する。

\item 18. 正の項からなる収束する級数の部分級数は自身でも収束する。

79. 持続的な実変数の関数を極限によって表現する．前節ではしばしば以下のような極限について取り扱ってきた：
\[
\lim_{n \to \infty} \phi_{n}(x),
\]
また、以下の級数も考慮した：
\[
u_{1}(x) + u_{2}(x) + \dots
  = \lim_{n \to \infty}\{u_{1}(x) + u_{2}(x) + \dots + u_{n}(x)\},
\]
これらの極限や級数において、求めたい極限を持つ関数は$n$以外にも変数$x$を含む．このような場合、極限は当然$x$の関数となる．例えば第75節では以下の関数に遭遇した：
\[
f(x) = \lim_{n \to \infty} n(\sqrt[n]{x} - 1):
\]
また、幾何級数 $1 + x + x^{2} + \dots$ の和は$x$の関数であり、$-1 < x < 1$ のとき $1/(1 - x)$ に等しく、それ以外の$x$の値に対して未定義である．

第II章で考慮された多くの「任意」または「自然」ではない関数は、以下の例から明らかになるように、このような単純な表現が可能である．

例 XXXI.
1. $\phi_{n}(x) = x$. ここで$n$は$\phi_{n}(x)$の表記に全く現れず、$\phi(x) = \lim\phi_{n}(x) = x$ であり、すべての$x$に対する値が同じである．

2. $\phi_{n}(x) = x/n$. ここで$\phi(x) = \lim\phi_{n}(x) = 0$ であり、すべての$x$に対する値が同じである．

3. $\phi_{n}(x) = nx$. $x > 0$ のとき $\phi_{n}(x) \to +\infty$；$x < 0$ のとき $\phi_{n}(x) \to -\infty$：$x = 0$ のみで$\phi_{n}(x)$が極限を持つ（$n \to \infty$のとき $0$に収束）．したがって、$x = 0$ のとき $\phi(x) = 0$であり、それ以外の$x$に対する値は定義されていない．

4. $\phi_{n}(x) = 1/nx, nx/(nx + 1)$．

5. $\phi_{n}(x) = x^{n}$. ここで$\phi(x) = 0$ ($-1 < x < 1$)；$\phi(x) = 1$ ($x = 1$)であり、それ以外の$x$に対する値は定義されていない．

6. $\phi_{n}(x) = x^{n}(1 - x)$．ここで$\phi(x)$は例5の$\phi(x)$と異なる点として、$x = 1$ のとき $0$である．

\item 7. $\phi_{n}(x) = x^{n}/n$. この$\phi(x)$は、例6の$\phi(x)$と異なり、$x = -1$のときも$x = 1$のときも値が0となる。

\item 8. $\phi_{n}(x) = x^{n}/(x^{n} + 1)$. [$\phi(x) = 0$, ($-1 < x < 1$); $\phi(x) = \frac{1}{2}$, ($x = 1$); $\phi(x) = 1$, ($x < -1$または$x > 1$); $x = -1$のとき$\phi(x)$は定義されていない。]

\item 9. $\phi_{n}(x) = x^{n}/(x^{n} - 1)$, $(x^{n} + 1)^{-1}$, $(x^{n} - 1)^{-1}$, $(x^{n} + x^{-n})^{-1}$, $(x^{n} - x^{-n})^{-1}$。

\item 10. $\phi_{n}(x) = (x^{n} - 1)/(x^{n} + 1)$, $(nx^{n} - 1)/(nx^{n} + 1)$, $(x^{n} - n)/(x^{n} + n)$. [最初の関数では、$|x| > 1$のとき$\phi(x) = 1$, $|x| < 1$のとき$\phi(x) = -1$, $x = 1$のとき$\phi(x) = 0$であり、$x = -1$のとき$\phi(x)$は定義されていない。二番目と三番目の関数では、$x = 1$および$x = -1$のときも定義されているが、二番目の関数では値が$1$で、三番目の関数では値が$-1$である。]

\item 11. $\phi(x) = 1$, ($|x| > 1$); $\phi(x) = -1$, ($|x| < 1$); $\phi(x) = 0$, ($x = 1$および$x = -1$)の例を構築せよ。

\item 12. $\phi_{n}(x) = x\{(x^{2n} - 1)/(x^{2n} + 1)\}^{2}$, $n/(x^{n} + x^{-n} + n)$。

\item 13. $\phi_{n}(x) = \{x^{n}f(x) + g(x)\}/(x^{n} + 1)$. [この関数では、$\phi(x) = f(x)$, ($|x| > 1$); $\phi(x) = g(x)$, ($|x| < 1$); $\phi(x) = \frac{1}{2}\{f(x) + g(x)\}$, ($x = 1$); $x = -1$のとき$\phi(x)$は定義されていない。]

\item 14. $\phi_{n}(x) = (2/\pi) \arctan(nx)$. [$\phi(x) = 1$, ($x > 0$); $\phi(x) = 0$, ($x = 0$); $\phi(x) = -1$, ($x < 0$)。この関数は数論の理論において重要であり、通常$\sgn x$で表される。]

\item 15. $\phi_{n}(x) = \sin nx\pi$. [$x$が整数のとき$\phi(x) = 0$; $x$が整数以外のとき$\phi(x)$は定義されていない（例xxiv.7）。]

\item 16. $\phi_{n}(x) = \sin (n!\, x\pi)$のとき、すべての有理数$x$に対して$\phi(x) = 0$（例xxiv.14）。[無理数$x$の場合にはより大きな困難が伴う。]

\item{17.} $\phi_{n}(x) = (\cos^{2} x\pi)^{n}$. [$\phi(x) = 0$ 除く$x$が整数のとき、$\phi(x) = 1$.]

\item{18.} $N \geq 1752$ のとき、西暦 $N$ 年の日数は
\[
365 + (\cos^{2} \frac{1}{4} N\pi)^{n} - (\cos^{2} \frac{1}{100} N\pi)^{n} + (\cos^{2} \frac{1}{400} N\pi)^{n}
\]
の極限値である。

\par\medskip
\ページ分割{149}

\par\medskip\noindent\textbf{【注釈】}\par
\節番号{80. 有界集合的辺界。}任意の実数 $s$ の系または集合 $S$ について、ある数 $K$ が存在して $s \leq K$ であるとき、$S$ を \emph{上に有界}という。また、ある数 $k$ が存在して $s \geq k$ であるとき、$S$ を \emph{下に有界}という。$S$ が上に有界かつ下に有界であるとき、単に $S$ を \emph{有界}という。

まず $S$ が上に有界（必ずしも下に有界でない）とする。そのような数 $K$ の存在する無限の数が存在する；例えば、$K$ より大きい任意の数はその性質を持つ。これらの数の中で最小のものが存在することを証明する。これを $M$ と呼ぶ。この数 $M$ は $S$ のどの要素よりも大きくないが、$M$ より小さい任意の数は少なくとも $S$ の一つの要素によって超えられる。

実数 $\xi$ を二つのクラス $L$ および $R$ に分類し、$\xi$ が $S$ の要素によって超えられるとき $L$ に、超えられないとき $R$ に属すとする。すると任意の $\xi$ は $L$ や $R$ のいずれか一つに属する。これらのクラスは存在する；例えば、$S$ のどの要素よりも小さい任意の数は $L$ に属し、$K$ は $R$ に属す。さらに、$L$ の任意の要素は $S$ の少なくとも一つの要素より小さく、したがって $R$ の任意の要素より小さい。

デデキントの定理（§17）の三つの条件を満たすためには、クラス $L$ と $R$ の間に分割する数 $M$ が存在することである。

The 数字 \( M \)はその存在を証明しなければならない数字である。まず第一に、\( S \)の任意の要素 \( s \)によって \( M \)は超えられない。もし \( S \) のような要素 \( s \)が存在すると仮定すると、\( s = M + \eta \)と書けるで、\(\eta\) は正である。このとき、数 \( M + \frac{1}{2}\eta \)は \( L \)に属し、また \( R \)にも属する。これは不可能である。一方、任意の \( M \)未満の数は \( L \)に属し、したがって \( S \) の少なくとも一つの要素によって超えられる。従って \( M \)は必要な全ての性質を持つ。

この数 \( M \)を \( S \)の上界と呼び、以下の定理を述べることができる。
\par\medskip\noindent\textbf{【結果】}\par
任意の集合 \( S \)が上界で制限されているときには、その上界 \( M \)が存在する。\( S \) の任意の要素は \( M \)を超えない；しかし \( M \)未満の任意の数は少なくとも一つの \( S \) の要素によって超えられる。
\par\medskip

同様に、下界で制限されている集合 \( S \)についても定理を証明することができる（必ずしも上界で制限されていない）。
\par\medskip\noindent\textbf{【結果】}\par
任意の集合 \( S \)が下界で制限されているときには、その下界 \( m \)が存在する。\( S \) の任意の要素は \( m \)未満ではない；しかし \( m \)より大きい任意の数よりも少なくとも一つの \( S \) の要素が小さい。
\par\medskip

集合 \( S \)が上界で制限されているときには、\( M \leq K \)であり、下界で制限されているときには \( m \geq k \)である。集合 \( S \)が制限されているときには、\( k \leq m \leq M \leq K \)である。

\Paragraph{81. 制限された関数の境界} 仮定として、\(\phi(n)\)は正整数変数 \( n \) の関数とする。すべての値 \(\phi(n)\)で構成される集合を \( S \)とし、\SecNo[§]{80}における全ての議論に適用できる。もし \( S \)が上界で制限されているか、または下界で制限されているか、または両方で制限されているときには、\(\phi(n)\)は上界で制限されている、または下界で制限されている、または両方で制限されているという。もし \(\phi(n)\)が上界で制限されているとすれば、すなわちある数 \( K \)が存在して \(\phi(n) \leq K\) であるときには、ある数 \( M \)が存在して

\Itemp{(i)} \emph{\(\phi(n) \leq M\) が全ての \( n \)に対して成り立つ}。

\Itemp{(ii)} 任意の正数 $\DELTA$ に対して、$\phi(n) > M - \DELTA$ となる $n$ の少なくとも一つが存在する。
この数 $M$ を $\phi(n)$ の上限と呼びます。同様に、$\phi(n)$ が下限をもつ場合、すなわちある数 $k$ が存在して $\phi(n) \leq k$ となるすべての $n$ に対して、ある数 $m$ が存在して

\Itemp{(i)} $\phi(n) \geq m$ となるすべての $n$ に対して；
\Itemp{(ii)} 任意の正数 $\DELTA$ に対して、$\phi(n) < m + \DELTA$ となる $n$ の少なくとも一つが存在する。
この数 $m$ を $\phi(n)$ の下限と呼びます。

もし $K$ が存在すれば、$M \leq K$；もし $k$ が存在すれば、$m \geq k$；そして両方の $k$ と $K$ が存在すれば
\[
k \leq m \leq M \leq K.
\]

82. 確定限界の不確定性の制限。$\phi(n)$ を任意の有界関数とするとき、その上限を $M$、下限を $m$ とし、実数 $\xi$ を任意に取る。そして、十分大きな $n$ の値に対して $\phi(n)$ が $\xi$ と持つ不等式の関係を考える。この関係には互いに排他的な3つの可能性がある：

\Item{(1)} 充分大きな $n$ のすべてに対して $\xi \geq \phi(n)$；
\Item{(2)} 充分大きな $n$ のすべてに対して $\xi \leq \phi(n)$；
\Item{(3)} 无限の数の $n$ に対して $\xi < \phi(n)$ であり、また无限の数の $n$ に対して $\xi > \phi(n)$。

(1) の場合、$\xi$ を上界と呼び、(2) の場合、$\xi$ を下界と呼び、(3) の場合、$\xi$ を中間値と呼ぶ。明らかに、上界より小さい上界は存在しないし、下界より大きい下界も存在しない。

すべての上界を含む集合を考える。この集合は $m$ よりも小さい要素を持たないため、下限を持つ。この下限を $\Lambda$ と表す。同様に、すべての下界を含む集合にも上限があり、これを $\lambda$ で表す。

We call $\Lambda$ and $\lambda$ respectively the 上限と下限の不確定性
of $\phi(n)$ as $n$ tends to infinity; and write
\[
\Lambda = \limsup \phi(n),\quad
\lambda = \liminf \phi(n).
\]
These numbers have the following properties:

(1) \( m \leq \lambda \leq \Lambda \leq M \);

(2) $\Lambda$ and $\lambda$ are the upper and lower bounds of the aggregate of intermediate
numbers, if any such exist;

(3) If $\DELTA$ is any positive number, then $\phi(n) < \Lambda + \DELTA$ for all sufficiently large
values of $n$, and $\phi(n) > \Lambda - \DELTA$ for an infinity of values of $n$;

(4) 同様に $\phi(n) > \lambda - \DELTA$ for all sufficiently large values of $n$, and
$\phi(n) < \lambda + \DELTA$ for an infinity of values of $n$;

(5) The necessary and sufficient condition that $\phi(n)$ should tend to a limit
is that $\Lambda = \lambda$, and in this case the limit is $l$, the common value of $\lambda$ and $\Lambda$.

Of these properties, (1) is an immediate consequence of the definitions;
and we can prove (2) as follows. If $\Lambda = \lambda = l$, there can be at most one intermediate
number, viz. $l$, and there is nothing to prove. Suppose then that
$\Lambda > \lambda$. Any intermediate number $\xi$ is less than any superior and greater than
any inferior number, so that $\lambda \leq \xi \leq \Lambda$. But if $\lambda < \xi < \Lambda$ then $\xi$ must be
intermediate, since it is plainly neither superior nor inferior. Hence there are
intermediate numbers as near as we please to either $\lambda$ or $\Lambda$.

To prove (3) we observe that $\Lambda + \DELTA$ is superior and $\Lambda - \DELTA$ intermediate or
inferior. The result is then an immediate consequence of the definitions; and
the proof of (4) is substantially the same.

最後に、(5)は次のように証明できる。もし $\Lambda = \lambda = l$ ならば
\[
l - \DELTA < \phi(n) < l + \DELTA
\]
が全ての正の $\DELTA$ および十分に大きい $n$ の値に対して成り立つので、$\phi(n)\to l$ である。逆に、もし $\phi(n) \to l$ ならば、上記の不等式は十分に大きい $n$ の値に対して成り立つ。したがって $l - \DELTA$ は下界であり、$l + \DELTA$ は上界なので
\[
\lambda \geq l - \DELTA,\quad
\Lambda \leq l + \DELTA,
\]
となり、従って $\Lambda - \lambda \leq 2\DELTA$ である。$\Lambda - \lambda \geq 0$ より、これは $\Lambda = \lambda$ のみが成り立つ。

---

以下は例を示す。

1. $\Lambda$ および $\lambda$ は任意の有限数の $\phi(n)$ の値の変更に影響を受けない。

2. $\phi(n) = a$ ならば全ての $n$ に対して $m = \lambda = \Lambda = M = a$。

3. $\phi(n) = 1/n$ ならば $m = \lambda = \Lambda = 0$ であり、$M = 1$。

4. $\phi(n) = (-1)^{n}$ ならば $m = \lambda = -1$ であり、$\Lambda = M = 1$。

5. $\phi(n) = (-1)^{n}/n$ ならば $m = -1$, $\lambda = \Lambda = 0$, $M = \frac{1}{2}$。

6. $\phi(n) = (-1)^{n}\{1 + (1/n)\}$ ならば $m = -2$, $\lambda = -1$, $\Lambda = 1$, $M = \frac{3}{2}$。

7. $\phi(n) = \sin n\theta\pi$ とする。$\theta > 0$ のとき、もし $\theta$ が整数ならば $m = \lambda = \Lambda = M = 0$。もし $\theta$ が整数でない有理数であると仮定すると、様々な場合が生じる。例えば、$\theta = p/q$, $p$ および $q$ が正の奇数であり互いに素かつ $q > 1$ のとき、$\phi(n)$ は周期的な値の列
\[
\sin(p\pi/q),\quad
\sin(2p\pi/q), \dots,\quad
\sin\{(2q - 1)p\pi/q\},\quad
\sin(2qp\pi/q), \dots
\]
をとる。$\phi(n)$ の最大値および最小値は $\cos(\pi/2q)$ および $-\cos(\pi/2q)$ であることが簡単に確認できるので、
\[
m = \lambda = -\cos(\pi/2q),\quad
\Lambda = M = \cos(\pi/2q)
\]
となる。$p$ および $q$ が両方とも奇数でない場合のケースも同様に議論できる。

Theθが無理数の場合はより難しく、この場合m = λ = -1でありΛ = M = 1であることが示される。また、φ(n)の値は(-1, 1)と[-1, 1]の区間に散在しており、任意の区間内の数ξに対して、n₁, n₂, … のような数列が存在し、k → ∞のときφ(nₖ) → ξであることも示される。この結果はHardyとLittlewoodによって「Diophantine Approximationに関するいくつかの問題」（Acta Mathematica, 第37巻）で簡潔な証明が与えられている。

次に、φ(n)がnθの小数部分の場合でも、これらの結果は非常に類似している。

以下は重要な収束の一般原則について述べる。

83. 有界関数に対する収束の一般的原理。
前節の結果により、有界関数φ(n)が極限に収束するための必要十分条件を形式化できる。この条件は通常、「収束の一般的原理」と呼ばれる。

定理。
有界関数φ(n)が極限に収束するための必要十分条件は、任意の正数δに対して、n₀(δ)という数を見つけることができることである。そして、すべてのn₁, n₂に対し、n₂ > n₁ ≥ n₀(δ)のとき
\[
|\phi(n₂) - \phi(n₁)| < δ
\]
が成り立つ。

まず第一に、この条件は必要である。もしφ(n) → lであれば、n ≥ n₀のとき
\[
l - \frac{1}{2}δ < φ(n) < l + \frac{1}{2}δ
\]
であり、よって
\[
|\phi(n₂) - \phi(n₁)| < δ
\Tag{(1)}
\]
がn₁, n₂ ≥ n₀のとき成立する。

第二に、条件は十分である。これを証明するために、\(\lambda = \Lambda\) を含むことを示すだけでよい。もし \(\lambda < \Lambda\) であれば、いかなる \(\DELTA\) が小さくても無限多数の \(n\) の値に対して \(\phi(n) < \lambda + \DELTA\) であり、また無限多数の \(n\) の値に対して \(\phi(n) > \Lambda - \DELTA\)。従って、任意に与えられた数 \(n_{0}\) より大きい \(n_{1}\) と \(n_{2}\) を選べる。そして
\[
\phi(n_{2}) - \phi(n_{1}) > \Lambda - \lambda - 2\DELTA,
\]
これは十分小さい \(\DELTA\) のとき \(\frac{1}{2}(\Lambda - \lambda)\) より大きい。これにより不等式 \(\text{(1)}\) に反する。従って \(\lambda = \Lambda\) であり、したがって \(\phi(n)\) は極限に向かう。

84. 不有界関数。これまでに述べたのは有界関数に対する収束の「一般原則」であるが、有界関数と同様に不有界関数に対しても同じであり、定理1において「有界関数」という言葉は省略できる。

第一に、\(\phi(n)\) が極限 \(l\) に向かう場合、それは確かに有界である。すべての値のうち有限個以外は \(l + \DELTA\) より小さく、かつ \(l - \DELTA\) より大きい。

第二に、定理1の条件が満たされる場合、
\[
|\phi(n_{2}) - \phi(n_{1})| < \DELTA
\]
が \(n_{1} \geq n_{0}\) および \(n_{2} \geq n_{0}\) のとき成立する。ある特定の値 \(n_{1}\) を選び、これは \(n_{0}\) より大きいものとする。すると
\[
\phi(n_{1}) - \DELTA < \phi(n_{2}) < \phi(n_{1}) + \DELTA
\]
が \(n_{2} \geq n_{0}\) のとき成立する。従って \(\phi(n)\) は有界であり、したがって最後の節の証明の第二部分も適用できる。

ページ分割{153}

理論的な重要性を「収束の一般原則」として過小評価することはできない。この原理は、\SecNo[§]{69}の定理と同様に、関数$\phi(n)$が極限に向かうかどうかを決定するための方法を与えてくれるが、その前に極限が存在した場合の値を予め知らなくてもよい。また、\SecNo[§]{69}のような特殊な性質を持つ定理に比べて、その制約は避けられないものではない。しかし、基本的な作業ではこの原理を省略して、これらの特別な定理から必要なすべてを得ることが可能である。そして、原則の重要性にもかかわらず、次の章でほとんど実用的には適用されていないことが確認される。

\par\medskip\noindent\textbf{【定理】}\par[2.]
無限級数$u_{1} + u_{2} + \dots$が収束する必要十分条件は、任意の正数$\DELTA$に対して、ある$n_{0}$を選び、
\[
|u_{n_{1}+1} + u_{n_{1}+2} + \dots + u_{n_{2}}| < \DELTA
\]
がすべての$n_{1}$と$n_{2}$（$n_{2} > n_{1} \geq n_{0}$）に対して成り立つことを可能にする$n_{0}$を見つけることができることである。
\par\medskip

\Paragraph{85. 複素関数の極限と複素項の級数}
これまでの章では、$n$の実数関数およびすべての項が実数である級数についてのみ論じてきた。しかし、関数や級数の項が複素数の場合でも、アイデアや定義を拡張することは困難ではない。

$\phi(n)$が複素数であり、
\[
\rho(n) + i\sigma(n)
\]
と等しいとき、$\rho(n), \sigma(n)$は$n$の実数関数である。このとき、$n \to \infty$の極限で$\rho(n)$と$\sigma(n)$がそれぞれ極限$r$と$s$に収束するならば、$\phi(n)$は極限$l = r + is$に収束すると定義し、
\[
\lim\phi(n) = l
\]
と書く。

同様に、$u_{n}$が複素数であり、$v_{n} + iw_{n}$と等しいとき、
\[
u_{1} + u_{2} + u_{3} + \dots
\]
が収束し、その和が$l = r + is$であるとは、
\[
v_{1} + v_{2} + v_{3} + \dots,\quad w_{1} + w_{2} + w_{3} + \dots
\]
がそれぞれ$r, s$に収束するときを意味する。

$u_{1} + u_{2} + u_{3} + \dots$が収束し、その和が$l$であるとは、当然のことながら、
\[
s_{n}
  = u_{1} + u_{2} + \dots + u_{n}
  = (v_{1} + v_{2} + \dots + v_{n})
  + i(w_{1} + w_{2} + \dots + w_{n})
\]
が$n \to \infty$の極限で$l$に収束することと同義である。

実数関数や級数の場合、収束性や振動性、有限または無限について定義を与えた。しかし、複素関数や級数の場合、$\rho(n)$および$\sigma(n)$の挙動を考慮する必要があるため、多くの可能性があり、これほどまでに価値あるものとは言い難い。必要な場合にこのような区別が必要な場合は、実部または虚部が個別にどのように振舞うかにより定義することとする。

86．読者は、既に証明された実数関数や級数の定理の拡張である以下の命題を証明する上で困難を感じることはない：

（1）$\lim\phi(n) = l$ならば、任意の固定値$p$に対して$\lim\phi(n + p) = l$である。

\item{(2)} $u_{1} + u_{2} + \dots$ が収束し、その和が$l$であるとき、$a + b + c + \dots + k + u_{1} + u_{2} + \dots$ も収束し、その和は$a + b + c + \dots + k + l$であり、また $u_{p+1} + u_{p+2} + \dots$ も収束し、その和は$l - u_{1} - u_{2} - \dots - u_{p}$である。

\item{(3)} $\lim\phi(n) = l$ かつ $\lim\psi(n) = m$ のとき、
\[
\lim\{\phi(n) + \psi(n)\} = l + m.
\]

\item{(4)} $\lim\phi(n) = l$ のとき、$\lim k\phi(n) = kl$。

\item{(5)} $\lim\phi(n) = l$ かつ $\lim\psi(n) = m$ のとき、$\lim \phi(n)\psi(n) = lm$。

\item{(6)} $u_{1} + u_{2} + \dots$ が和$l$に収束し、$v_{1} + v_{2} + \dots$ が和$m$に収束するとき、$(u_{1} + v_{1}) + (u_{2}+ v_{2}) + \dots$ も和$l + m$に収束する。

\item{(7)} $u_{1} + u_{2} + \dots$ が和$l$に収束するとき、$ku_{1} + ku_{2} + \dots$ も和$kl$に収束する。

\item{(8)} $u_{1} + u_{2} + u_{3} + \dots$ が収束するとき、$\lim u_{n} = 0$。

\item{(9)} $u_{1} + u_{2} + u_{3} + \dots$ が収束するとき、その項を括弧で組み替えた任意の級数も収束し、両者の和は同じである。

\pagebreak{155}

\par\medskip\noindent\textbf{【注釈】}\par
例えば、定理(5)を証明しましょう。$\phi(n) = \rho(n) + i\sigma(n)$、$\psi(n) = \rho'(n) + i\sigma'(n)$、$l = r + is$、$m = r' + is'$ とします。

すると
\[
\rho(n) \to r,\quad
\sigma(n) \to s,\quad
\rho'(n) \to r',\quad
\sigma'(n) \to s'.
\]

また
\[
\phi(n)\psi(n) = \rho\rho' - \sigma\sigma' + i(\rho\sigma' + \rho'\sigma),
\]
そして
\[
\rho\rho' - \sigma\sigma' \to rr' - ss',\quad
\rho\sigma' + \rho'\sigma \to rs' + r's;
\]
したがって
\[
\phi(n)\psi(n) \to rr' - ss' + i(rs' + r's),
\]
すなわち
\[
\phi(n)\psi(n) \to (r + is)(r' + is') = lm.
\]
\par\medskip

次の定理は異なる性格を持つ。

\par\medskip\noindent\textbf{【結果】}\par
\item{(10)} $\phi(n) = \rho(n) + i\sigma(n)$ が$n \to \infty$ のとき零に収束するためには、かつ只有に
\[
|\phi(n)| = \sqrt{\{\rho(n)\}^{2} + \{\sigma(n)\}^{2}}
\]
が零に収束することである。
\par\medskip

\par\medskip\noindent\textbf{【注釈】}\par
もし$\rho(n)$と$\sigma(n)$がともに0に収束するならば、$\sqrt{\rho^{2} + \sigma^{2}}$も0に収束することは明らかです。逆も真で、$\rho$または$\sigma$の数値は$\sqrt{\rho^{2} + \sigma^{2}}$を超えることはないからです。
\par\medskip

\par\medskip\noindent\textbf{【結果】}\par
\Item{(11)} 一般に、$\phi(n)$が極限$l$に収束するためには、
\[
|\phi(n) - l|
\]
が0に収束することが必要かつ十分条件です。
\par\medskip

\par\medskip\noindent\textbf{【注釈】}\par
$\phi(n) - l$が0に収束することから、(10)を適用できます。
\par\medskip

\par\medskip\noindent\textbf{【結果】}\par
\Item{(12)} 第83節～第84節の定理{\upshape1}と{\upshape2}は、$\phi(n)$と$u_{n}$が複素数である場合でも依然として成り立ちます。
\par\medskip

\par\medskip\noindent\textbf{【注釈】}\par
$\phi(n)$が$l$に収束するための必要十分条件を示す必要があります。これは
\[
|\phi(n_{2}) - \phi(n_{1})| < \DELTA
\Tag{(1)}
\]
が$n_{2} > n_{1} \geq n_{0}$のとき成立することです。

もし$\phi(n) \to l$ならば、$\rho(n) \to r$と$\sigma(n) \to s$であり、よって$\DELTA$に依存する数値$n_{0}'$と$n_{0}''$を選び、
\[
|\rho(n_{2}) - \rho(n_{1})| < \tfrac{1}{2}\DELTA,\quad
|\sigma(n_{2}) - \sigma(n_{1})| < \tfrac{1}{2}\DELTA,
\]
が$n_{2} > n_{1} \geq n_{0}'$のとき最初の不等式、および$n_{2} > n_{1} \geq n_{0}''$のとき第二の不等式を満たすようにできます。したがって
\[
|\phi(n_{2}) - \phi(n_{1})|
  \leq |\rho(n_{2}) - \rho(n_{1})| + |\sigma(n_{2}) - \sigma(n_{1})|
  < \DELTA
\]
が$n_{2} > n_{1} \geq n_{0}$のとき成立します。ここで、$n_{0}$は$n_{0}'$と$n_{0}''$のうち大きい方です。したがって条件{(1)}は必要条件です。この条件が十分条件であることを証明するには、
\[
|\rho(n_{2}) - \rho(n_{1})| \leq |\phi(n_{2}) - \phi(n_{1})| < \DELTA
\]
が$n_{2} > n_{1} \geq n_{0}$のとき成立することを確認すればよいです。これにより$\rho(n)$は極限$r$に収束し、同様にして$\sigma(n)$も極限$s$に収束することが示されます。
\par\medskip
\\PageSep{156}

\paragraph{87. $z^n$ の極限について（$n \to \infty$, $z$ は任意の複素数）}
$\phi(n) = z^n$ の重要な場合を考える。この問題は、§72で実数の場合にすでに議論されている。

もし $z^n \to l$ であれば、$z^{n+1} \to l$ である（§86の式(1)）。しかし、§86の式(4)によれば
\[
z^{n+1} = zz^n \to zl,
\]
よって $l = zl$ であり、これには (\ia) $l = 0$ または (\ib) $z = 1$ のみが成り立つ。
もし $z = 1$ であれば $\lim z^n = 1$。この特殊な場合以外では、極限は存在する場合でも零に収束する。

今、$z = r(\cos\theta + i\sin\theta)$ とし、$r$ が正のとき
\[
z^n = r^n (\cos n\theta + i\sin n\theta),
\]
よって $|z^n| = r^n$。したがって $|z^n|$ は $r < 1$ の場合にのみ零に収束する。
§86の式(10)から
\[
\lim z^n = 0
\]
は $r < 1$ の場合にのみ成り立つ。他の場合では、$z^n$ は極限に収束しない（$z = 1$ で $z^n \to 1$ の場合を除く）。

\paragraph{88. 複素数 $z$ の幾何級数 $1 + z + z^2 + \dots$}
\[
s_n = 1 + z + z^2 + \dots + z^{n-1} = (1 - z^n)/(1 - z)
\]
であるから、$z = 1$ でない限り $s_n$ の値は $n$ であり、幾何級数 $1 + z + z^2 + \dots$ は $|z| = r < 1$ の場合にのみ収束し、その和は $1/(1 - z)$ である。

以下に翻訳します：

\[
z = r(\cos\theta + i\sin\theta) = r\Cis\theta
\] とし、\[
r < 1
\] のとき、

\begin{align*}
1 + z + z^{2} + \dots &= 1/(1 - r\Cis\theta),
\intertext{または}
1 + r \Cis\theta + r^{2} \Cis 2\theta + \dots
  &= 1/(1 - r\Cis\theta)\\
  &= (1 - r\cos\theta + ir\sin\theta)/(1 - 2r\cos\theta + r^{2}).
\end{align*}

実部と虚部を分離すると、

\begin{align*}
1 + r\cos\theta + r^{2}\cos 2\theta + \dots
  &= (1 - r\cos\theta)/(1 - 2r\cos\theta + r^{2}),\\
r\sin\theta + r^{2}\sin 2\theta + \dots
  &= r\sin\theta/(1 - 2r\cos\theta + r^{2}),
\end{align*}

\[
r < 1
\] のときが得られる。\[
\theta
\] を \[
\theta + \pi
\] に置き換えると、これらの結果は数値的に \[
-1 < r < 1
\] の範囲で成立することも確認できる。

従って、\[
-1 < r < 1
\] のときが成り立つ。

\par\medskip\noindent\textbf{【例 XXXIII.】}\par
\Item{1.} \[
r < 1
\] のとき \[
\phi(n) = r^{n} \cos n\theta
\] は0に収束し、\[
r = 1
\] かつ \[
\theta
\] が \[
2\pi
\] の倍数のときには1に収束することを直接証明せよ。さらに、\[
r = 1
\] かつ \[
\theta
\] が \[
2\pi
\] の倍数でないときには有限振動し、\[
r > 1
\] かつ \[
\theta
\] が \[
2\pi
\] の倍数のときには正の無限大に収束し、\[
r > 1
\] かつ \[
\theta
\] が \[
2\pi
\] の倍数でないときには無限振動することを証明せよ。

\Item{2.} 同様の結果を得るために \[
\phi(n) = r^{n} \sin n\theta
\] について証明せよ。

\Item{3.} 
\begin{gather*}
z^{m} + z^{m+1} + \dots = z^{m}/(1 - z),\\
z^{m} + 2z^{m+1} + 2z^{m+2} + \dots = z^{m}(1 + z)/(1 - z),
\end{gather*}
\[
|z| < 1
\] のときのみが成り立つことを証明せよ。第86節のどの定理を使用したかを示せ。

\Item{4.} \[
-1 < r < 1
\] のとき、

\[
1 + 2r\cos\theta + 2r^{2}\cos 2\theta + \dots
  = (1 - r^{2})/(1 - 2r\cos\theta + r^{2})
\]

が成り立つことを証明せよ。

\Item{5.} 
\[
1 + \frac{z}{1 + z} + \left(\frac{z}{1 + z}\right)^{2} + \dots
\]
\[
|z/(1 + z) | < 1
\] のとき、和は \[
1\bigg/\left(1 - \dfrac{z}{1 + z}\right) = 1 + z
\] に収束することを示せ。この条件が実部が \[
-\frac{1}{2}
\] より大きい \[
z
\] の条件と同値であることを示せ。
\par\medskip

\#\#\#  Miscellaneous Examples on Chapter IV.

1. 関数 $\phi(n)$ は $n = 0, 1, 2, \ldots$ のとき値が $1$, $0$, $0$, $0$, $1$, $0$, $0$, $0$, $1$, $\ldots$ となる。$\phi(n)$ を $n$ の式で表すと、三角関数を含まない形で $\phi(n) = \frac{1}{4}\{1 + (-1)^{n} + i^{n} + (-i)^{n}\}$ である。

2. $\phi(n)$ が $n$ に従って安定的に増加し、$\psi(n)$ が安定的に減少し、かつすべての $n$ の値に対して $\psi(n) > \phi(n)$ ならば、両方とも $\phi(n)$ と $\psi(n)$ は極限に向かうし、$\lim\phi(n) \leq \lim\psi(n)$ である。これは §69 から直接の推論である。

3. 次を証明せよ： 
\[
\phi(n) = \left(1 + \frac{1}{n}\right)^{n},\quad
\psi(n) = \left(1 - \frac{1}{n}\right)^{-n},
\]
ならば $\phi(n + 1) > \phi(n)$ かつ $\psi(n + 1) < \psi(n)$。最初の結果は §73 で既に証明されている。

4. さらに、$\psi(n) > \phi(n)$ がすべての $n$ の値に対して成り立つことを証明し、前節の例を用いて $\phi(n)$ と $\psi(n)$ は $n$ に従って極限に向かうことを導け。注：$\lim\{\psi(n) - \phi(n)\} = 0$ であることが証明されれば、各関数は極限 $e$ に向かう。 Chrystal の Algebra (第2巻、第78ページ) にはこの証明があるが、我々は §IX で異なる方法で証明する。

5. 和が $n$ であるすべての正整数の組の積の算術平均を $S_{n}$ と表す。$\lim(S_{n}/n^{2}) = \frac{1}{6}$ を示せ。
数学オリンピック 1903 年。

6. $x_{1} = \frac{1}{2}\{x + (A/x)\}$, $x_{2} = \frac{1}{2}\{x_{1} + (A/x_{1})\}$ と次々に定義するとき、$x$ と $A$ が正の値であれば $\lim x_{n} = \sqrt{A}$ であることを証明せよ。まず、$\frac{x_{n} - \sqrt{A}}{x_{n} + \sqrt{A}} = \left(\frac{x - \sqrt{A}}{x + \sqrt{A}}\right)^{2^{n}}$ を証明せよ。

\item7. \(\phi(n)\)がすべての\(n\)に対して正の整数であり、また\(n\)と共に無限大に収束するならば、\(0 < x < 1\)の場合、\(x^{\phi(n)}\)は\(0\)に収束し、\(x > 1\)の場合、\(x^{\phi(n)}\)は\(+\infty\)に収束する。また、他の\(x\)の値に対して、\(n \to \infty\)のときの\(x^{\phi(n)}\)の挙動を議論せよ。

\item8. \(a_{n}\)が\(n\)と共に単調増加または単調減少するならば、\(\frac{a_{1} + a_{2} + \dots + a_{n}}{n}\)も同様の性質を持つ。注：例題8-12はBromwichの『無限級数』から引用されたものである。

\item9. \(x_{n+1} = \sqrt{k + x_{n}}\)、かつ\(k\)と\(x_{1}\)が正ならば、\(x_{1}\), \(x_{2}\), \(x_{3}\), \(\dots\)の数列は、\(x_{1}\)が方程式\(x^{2} = x + k\)の正の根\(\alpha\)より小さいか大きいかにより単調増加または単調減少する。また、いずれの場合も、\(n \to \infty\)のとき、\(x_{n} \to \alpha\)。

\item10. \(x_{n+1} = \frac{k}{1 + x_{n}}\)、かつ\(k\)と\(x_{1}\)が正ならば、数列\(x_{1}\), \(x_{3}\), \(x_{5}\), \(\dots\)と\(x_{2}\), \(x_{4}\), \(x_{6}\), \(\dots\)は一方が単調増加で他方が単調減少であり、それぞれの数列は方程式\(x^{2} + x = k\)の正の根\(\alpha\)に収束する。

\item11. 関数\(f(x)\)はすべての\(x\)に対して増加し連続である（参照：第V章）。数列\(x_{1}\), \(x_{2}\), \(x_{3}\), \(\dots\)は方程式\(x = f(x)\)の根に収束するかどうかについて、一般的なグラフィカルな理由から議論せよ。特に、この方程式が唯一の根を持つ場合を考慮し、曲線\(y = f(x)\)が直線\(y = x\)から上から下へまたは下から上へ交差する場合に分けよ。

\item12. \(x_{1}\), \(x_{2}\)が正であり、かつ\(x_{n+1} = \frac{1}{2}(x_{n} + x_{n-1})\)ならば、数列\(x_{1}\), \(x_{3}\), \(x_{5}\), \(\dots\)と\(x_{2}\), \(x_{4}\), \(x_{6}\), \(\dots\)は一方が単調減少で他方が単調増加であり、共通の極限値\(\frac{1}{3}(x_{1} + 2x_{2})\)を持つ。

\Item{13.} 関数 \( y \) を定める方程式
\[
y = \lim_{n \to \infty} \frac{x^{2n} \sin \frac{1}{2}\pi x + x^2}{x^{2n} + 1}
\]
を用いてグラフを描きなさい。
\MathTrip{1901.}

\Item{14.} 関数
\[
y = \lim_{n \to \infty} \frac{1}{1 + n \sin^2 \pi x}
\]
は \( x \) が整数でない限り \( y = 0 \) であり、\( x \) が整数のとき \( y = 1 \) である。また
\[
y = \lim_{n \to \infty} \frac{\psi(x) + n \phi(x) \sin^2 \pi x}{1 + n \sin^2 \pi x}
\]
は \( x \) が整数でない限り \( y = \phi(x) \) であり、\( x \) が整数のとき \( y = \psi(x) \) である。

\Item{15.} 次の関数
\[
y = \lim_{n \to \infty} \frac{x^n \phi(x) + x^{-n} \psi(x)}{x^n + x^{-n}}
\]
のグラフを示せ。このグラフは \( \phi(x) \) および \( \psi(x) \) の一部からなり、通常二つの孤立点も含まれる。(\ia) \( x = 1 \), (\ib) \( x = -1 \), (\ic) \( x = 0 \) のとき \( y \) は定義されているか？

\Item{16.} \( x \) が有理数のとき \( y = 0 \)、無理性のとき \( y = 1 \) である関数 \( y \) を
\[
y = \lim_{m \to \infty} \sgn \{\sin^2 (m! \pi x)\}
\]
と表すことができる。ここで
\[
\sgn x = \lim_{n \to \infty} \left( \frac{2}{\pi} \right) \arctan(nx)
\]
である。（例題 \Ex{xxxi}. 14. 参照）[有理数 \( x \) のとき、\( \sin^2 (m! \pi x) \) および従って \( \sgn \{\sin^2 (m! \pi x)\} \) は一定の \( m \) からゼロとなる。無理性 \( x \) のとき、\( \sin^2 (m! \pi x) \) は常に正であり、従って \( \sgn \{\sin^2 (m! \pi x)\} \) も常に1である。]

また
\[
y = 1 - \lim_{m \to \infty} \left[ \lim_{n \to \infty} \{ \cos(m! \pi x) \}^{2n} \right]
\]
と表すこともできる。

\Item{17.} 次の級数を求和せよ。
\[
\sum_{1}^{\infty} \frac{1}{\nu(\nu + 1)},\quad
\sum_{1}^{\infty} \frac{1}{\nu(\nu + 1) \dots (\nu + k)}.
\]

\[
\frac{1}{\nu(\nu + 1)\dots(\nu + k)}
  = \frac{1}{k} \left\{\frac{1}{\nu(\nu + 1)\dots(\nu + k - 1)}
                     - \frac{1}{(\nu + 1)(\nu + 2)\dots(\nu + k)}\right\},
\]
より
\[
\sum_{1}^{n} \frac{1}{\nu(\nu + 1)\dots(\nu + k)}
  = \frac{1}{k} \left\{\frac{1}{1·2\dots k}
                    - \frac{1}{(n + 1)(n + 2)\dots (n + k)}\right\},
\]
ゆえに
\[
\sum_{1}^{\infty} \frac{1}{\nu(\nu + 1)\dots (\nu + k)} = \frac{1}{k(k!)},
\]

18. \( |z| < |\alpha| \)のとき
\[
\frac{L}{z - \alpha}
  = -\frac{L}{\alpha} \left(1 + \frac{z}{\alpha} + \frac{z^{2}}{\alpha^{2}} + \dots\right);
\]
また \( |z| > |\alpha| \)のとき
\[
\frac{L}{z - \alpha}
  = \frac{L}{z} \left(1 + \frac{\alpha}{z} + \frac{\alpha^{2}}{z^{2}} + \dots\right).
\]

19. \( (Az + B)/(az^{2} + 2bz + c) \)の\( z \)に関する冪級数展開について。\( az^{2} + 2bz + c = 0 \)の根を\(\alpha, \beta\)とし、\( az^{2} + 2bz + c = a(z - \alpha)(z - \beta) \)とする。\( A, B, a, b, c \)はすべて実数であり、\(\alpha\)と\(\beta\)は異なるものとする。このとき
\[
\frac{Az + B}{az^{2} + 2bz + c}
  = \frac{1}{a(\alpha - \beta)}
      \left(\frac{A\alpha + B}{z - \alpha}
          - \frac{A\beta + B}{z - \beta}\right),
\]
が成り立つ。\( b^{2} > ac \)または \( b^{2} < ac \)の2つの場合がある。

(1) \( b^{2} > ac \)のとき、根\(\alpha, \beta\)は実数で異なる。\( |z| \)が\(|\alpha|\)または\(|\beta|\)よりも小さい場合は、\( 1/(z - \alpha) \)と\( 1/(z - \beta) \)を\( z \)の昇べき級数に展開できる（例題18参照）。\( |z| \)が\(|\alpha|\)または\(|\beta|\)よりも大きい場合は、降べき級数に展開する必要がある。また、\( |z| \)が\(|\alpha|\)と\(|\beta|\)の間にある場合、1つの分数を昇べき級数、もう1つの分数を降べき級数に展開する。読者は実際に結果を書き出すことを recomendされる。\( |z| \)が\(|\alpha|\)または\(|\beta|\)と等しい場合は、このような展開は不可能である。

\Item{(2)} $b^{2} < ac$ ならば、根は共役複素数（\okrickRef{第III章}{§43}）であり、次のように書くことができる：
\[
\alpha = \rho\Cis\phi, \quad
\beta  = \rho\Cis(-\phi),
\]
ここで $\rho^{2} = \alpha\beta = c/a$、$\rho\cos\phi = \frac{1}{2}(\alpha + \beta) = - b/a$ より $\cos\phi = -\sqrt{\frac{b^{2}}{ac}}$、$\sin\phi = \sqrt{1 - \frac{b^{2}}{ac}}$。

$|z| < \rho$ ならば、各分数は $z$ の昇べき級数展開が可能である。$z^{n}$ の係数は
\[
\frac{A\rho\sin n\phi + B\sin\{(n + 1)\phi\}}{a\rho^{n+1} \sin\phi}
\]
となる。

$|z| > \rho$ の場合、同様の降べき級数展開が得られ、$|z| = \rho$ の場合、そのような展開は不可能である。

\Item{20.} $|z| < 1$ のとき
\[
1 + 2z + 3z^{2} + \dots + (n + 1)z^{n} + \dots = \frac{1}{(1 - z)^{2}}
\]
を示せ。

[$n$ 項までの和は $\frac{1 - z^{n}}{(1 - z)^{2}} - \frac{nz^{n}}{1 - z}$ である。]

\Item{21.} $|z| < |\alpha|$ のとき昇べき級数、または $|z| > |\alpha|$ のとき降べき級数として $\frac{L}{(z - \alpha)^{2}}$ を $z$ について展開せよ。

\Item{22.} $b^{2} = ac$ かつ $|az| < |b|$ のとき
\[
\frac{Az + B}{az^{2} + 2bz + c} = \sum_{0}^{\infty} p_{n}z^{n},
\]
ここで $p_{n} = \left(\frac{-a}{b}\right)^{n} \{(n + 1)aB - nbA\}$；$|az| > |b|$ のとき、降べき級数展開を求めよ。

\Item{23.} 分数 $\frac{1}{1 + z^{2}}$ の結果をEx.19で確認せよ。[$\frac{1}{1 + z^{2}} = \sum z^{n} \sin\left\{\frac{n + 1}{2}\pi\right\} = 1 - z^{2} + z^{4} - \dots$]

\Item{24.} $|z| < 1$ のとき
\[
\frac{1}{1 + z + z^{2}}
  = \frac{2}{\sqrt{3}} \sum_{0}^{\infty} z^{n} \sin\left\{\frac{2(n + 1)}{3}\pi\right\}
\]
を証明せよ。

\Item{25.} $\frac{1 + z}{1 + z^{2}}$、$\frac{1 + z^{2}}{1 + z^{3}}$、および $\frac{1 + z + z^{2}}{1 + z^{4}}$ を $z$ の昇べき級数で展開せよ。これらの結果はどのような値の $z$ に対して成り立つか？

\Item{26.} \( a/(a + bz + cz^{2}) = 1 + p_{1}z + p_{2}z^{2} + \dots \) ならば
\[
1 + p_{1}^{2}z + p_{2}^{2}z^{2} + \dots
  = \frac{a + cz}{a - cz}\, \frac{a^{2}}{a^{2} - (b^{2} - 2ac)z + c^{2}z^{2}}.
\]
\MathTrip{1900.}

\Item{27.} \( \lim_{n \to \infty} s_{n} = l \) ならば
\[
\lim_{n \to \infty} \frac{s_{1} + s_{2} + \dots + s_{n}}{n} = l.
\]

[令 \( s_{n} = l + t_{n} \). このとき、\( (t_{1} + t_{2} + \dots + t_{n})/n \) が零に収束することを証明せねばならない。
\PgSep{161}

数 \( t_{1}, t_{2}, \dots, t_{n} \) を二つの集合 \( t_{1}, t_{2}, \dots, t_{p} \) と \( t_{p+1}, t_{p+2}, \dots, t_{n} \) に分ける。ここで、\( p \) は \( n \) の関数であり、\( n \to \infty \) のとき \( p \to \infty \) だが、\( n \) よりも遅く収束するものとするので、\( p \to \infty \) かつ \( p/n \to 0 \): 例えば、\( p \) を \( \sqrt{n} \) の整数部分と仮定してもよい。

任意の正の数 \( \DELTA \) を与える。どんなに小さい \( \DELTA \) も与えても、ある \( n_{0} \) が存在して、\( n \geq n_{0} \) のとき \( t_{p+1}, t_{p+2}, \dots, t_{n} \) はすべて \( |\frac{1}{2}\DELTA| \) より小さいので
\[
|(t_{p+1} + t_{p+2} + \dots + t_{n})/n|
  < \frac{1}{2}\DELTA(n - p)/n < \frac{1}{2} \DELTA.
\]
また、\( t_{1}, t_{2}, \dots \) のすべての数の最大絶対値を \( A \) とすると
\[
|(t_{1} + t_{2} + \dots + t_{p})/n| < pA/n,
\]
これは \( n \geq n_{0} \) で十分大きい \( n_{0} \) のとき \( p/n \to 0 \) よりも小さい。したがって
\[
|(t_{1} + t_{2} + \dots + t_{n})/n|
  \leq |(t_{1} + t_{2} + \dots + t_{p})/n|
     + |(t_{p+1} + \dots + t_{n})/n| < \DELTA
\]
が \( n \geq n_{0} \) のとき成立し、定理を証明する。

読者に、極限に関する問題を扱うための専門性を得るためには、上記の議論を非常に慎重に研究すべきである。極限がゼロであることを証明する際にしばしば必要となるのは、与えられた表現を二つの部分に分割し、それぞれ異なる方法でその極限がゼロであることを証明することである。このような場合、証明は決して簡単ではない。

証明のポイントは次の通りである：$(t_{1} + t_{2} + \dots + t_{n})/n$ が $n$ が大きいとき小さいことを証明する必要がある。$t$ の値が大きいときその添字が大きい。括弧内の項を二つのグループに分ける。第一グループの項はすべて小さいわけではないが、その数は $n$ に対して小さい。第二グループの項の数は $n$ に対して小さくないが、各項はすべて小さいし、その数は少なくとも $n$ よりも少ないため、それらの和は $n$ に対して小さい。したがって、$(t_{1} + t_{2} + \dots + t_{n})/n$ を分割したそれぞれの部分は $n$ が大きいとき小さい。

28. $\phi(n) - \phi(n - 1) \to l$ であるとき $n \to \infty$ のとき $\phi(n)/n \to l$。

$\phi(n) = s_{1} + s_{2} + \dots + s_{n}$ とすると、$\phi(n) - \phi(n - 1) = s_{n}$ であり、定理は前例で証明されたものに帰着される。

29. $s_{n} = \frac{1}{2}\{1 - (-1)^{n}\}$ のとき、$n$ が奇数のとき $s_{n} = 1$、偶数のとき $s_{n} = 0$ であり、$n \to \infty$ のとき $(s_{1} + s_{2} + \dots + s_{n})/n \to \frac{1}{2}$。

この例は、Ex.27 の逆が真でないことを示す：$s_{n}$ は $n \to \infty$ のとき振動する。

30. シリーズ $\tfrac{1}{2} + \cos\theta + \cos 2\theta + \dots,\quad \sin\theta + \sin 2\theta + \dots$ の最初の $n$ 項の和を $c_{n}$、$s_{n}$ とすると、

$\lim (c}{1} + c\\textit{{2} + \dots + c}{n})/n = 0,\quad
\lim (s${1} + s${2} + \dots + s\_{n})/n = \tfrac{1}{2} \cot\tfrac{1}{2} \theta$。

例題終了

\chapter{第5章：Chapter V: LIMITS OF FUNCTIONS OF A CONTINUOUS VARIABLE. CONTINUOUS AND DISCONTINUOUS FUNCTIONS}

\章[連続変数の関数の極限]
{\V}{連続変数の関数の極限。連続関数と不连续関数}

\par\medskip\noindent\textbf{89. $x$ が無限大に向かうときの極限}\par 現在、我々は一連の関数について再び取り組みます。これらの関数は実数変数を用いて定義され、また単値関数に限定します。このような関数を $\phi(x)$ で表記します。$x$ は基底直線 $\Lambda$ 上の点に対応するすべての値を順次採用し、特定の点から始まり右方向へと進み続けます。このような状況において、$x$ が無限大に向かうという言い方をするか又は $x \to \infty$ と記します。最後の章で議論した「$n$ が無限大に向かう」ことと異なる点は、$x$ が無限大に向かうにつれてすべての値を採用することです。つまり、$x$ に対応する点 $P$ は基底直線 $\Lambda$ の初期位置より右側の全ての点に順次一致しますが、一方で $n$ は一連のジャンプによって無限大に向かいます。この差異を表現するために、$x$ が「連続的に」無限大に向かうと言えます。

最後の章の最初で説明したように、$x$ の関数と $n$ の関数との間には非常に密接な対応関係があります。任意の $n$ の関数は $x$ の関数の値から選択したものと見ることができます。最後の章では、$n$ が無限大に向かうときの関数 $\phi(n)$ の振る舞いに特徴付けられる性質について議論しました。現在は同様の問題を $x$ の関数 $\phi(x)$ に対して考慮します；そして導き出される定義と定理は、最後の章で得られたものと実質的に同じです。したがって、\SecNo[§]{58} の Def.1 に対応するものは以下の通りです：
\ページセパレーター{163}

定義．関数$\phi(x)$が$x$が無限大に近づくとき極限$l$に収束するとは、任意に小さな正の数$\DELTA$に対して、ある数$x_{0}(\DELTA)$を選び、$x \geq x_{0}(\DELTA)$となるすべての$x$について$\phi(x)$が$l$から$\DELTA$未満の差異を持つことを指します。すなわち，
\[
|\phi(x) - l| < \DELTA
\]
が成り立つとき、関数$\phi(x)$は極限$l$に収束すると定義します。

この場合、次のように記述できます：
\[
\lim_{x \to \infty} \phi(x) = l,
\]
または曖昧性がない限り単に
\[
\lim \phi(x) = l, \quad \text{または} \quad \phi(x) \to l.
\]

同様に、

定義．関数$\phi(x)$が$x$とともに無限大に収束するとは、任意に大きな数$\Delta$に対して、ある数$x_{0}(\Delta)$を選び、$x \geq x_{0}(\Delta)$となるすべての$x$について$\phi(x) > \Delta$が成り立つことを指します。

このとき、次のように記述できます：
\[
\phi(x) \to \infty.
\]

同様に、$\phi(x) \to -\infty$を定義します。
（注：しばしば$\infty$, $x \to +\infty$, $\phi(x) \to +\infty$の代わりに$+\infty$, $x \to +\infty$, $\phi(x) \to +\infty$と記述することが便利である場合があります。）

最後に、以下の定義があります：

\par\medskip\noindent\textbf{【定義】}\par[3.]
前二つの定義の条件がどちらも満たされない場合、$\phi(x)$ は $x$ が $\infty$ に近づくとき振動するという。$|\phi(x)|$ が $x \geq x_{0}$ のときに一定の定数 $K$ より小さいとすれば、$\phi(x)$ は有限振動するとし、それ以外は無限振動とする。
\par\medskip

読者は最後の章で公式 $\phi(n) \to l$, $\phi(n) \to \infty$ を表すための非形式的な表現について非常に慎重に考慮したことを思い出すだろう。現在の場合でも同様の表現法を当然用いることができる。つまり、$x$ が大きいときに $\phi(x)$ は小さいかほぼ $l$ に等しいか大きいということができ、そのような意味での「小さい」、「ほぼ等しい」、「大きい」という言葉を使用する。

\par\medskip\noindent\textbf{【例】}\par[XXXIV.]
1. 次の関数の振る舞いを $x \to \infty$ のときを考える：$\frac{1}{x}$, $1 + \frac{1}{x}$, $x^{2}$, $x^{k}$, $[x]$, $x - [x]$, $[x] + \sqrt{x - [x]}$。

最初の四つの関数は、第四章で十分に議論された関数と完全に一致している。最後の三つのグラフは第二章（例題 xvi. 1, 2, 4）で構築され、読者はすぐに $[x] \to \infty$, $x - [x]$ が有限振動し、$[x] + \sqrt{x - [x]} \to \infty$ であることがわかる。

一方、簡単な注記を加えることが可能である。関数 $\phi(x) = x - [x]$ はグラフの形から明らかに $0$ と $1$ の間で振動する。この関数は整数 $x$ のとき零となるため、導出された関数 $\phi(n)$ は常に零であり、限界値として零に収束する。同様のことは次の式においても成り立つ：
\[
\phi(x) = \sin x\pi, \quad \phi(n) = \sin n\pi = 0.
\]
$\phi(x) \to l$ や $\phi(x) \to \infty$ や $\phi(x) \to -\infty$ のような性質が存在する場合、それに対応する性質は $\phi(n)$ にも存在するが、逆は必ずしも成り立つとは限らない。

2. 同様の方法で以下の関数を考慮する：
\[
(\sin x\pi)/x, \quad x\sin x\pi, \quad (x\sin x\pi)^{2}, \quad \tan x\pi, \quad a\cos^{2} x\pi + b\sin^{2} x\pi,
\]
これらの関数のグラフを用いて注記を説明する。

3. 定義1に対する幾何学的解釈を与える。これは第IV章、§59における幾何学的解釈に類似している。

4. $\phi(x) \to l$ であり $l$ が零でない場合、$\phi(x)\cos x\pi$ および $\phi(x)\sin x\pi$ は有限に振動する。$\phi(x) \to \infty$ や $\phi(x) \to -\infty$ の場合、これらは無限に振動する。これらの関数のグラフは波状曲線であり、$y = \phi(x)$ および $y = -\phi(x)$ の間で振動する。

5. 関数
\[
y = f(x)\cos^{2} x\pi + F(x)\sin^{2} x\pi,
\]
の挙動を、$x \to \infty$ において議論する。ここで $f(x)$ および $F(x)$ は任意の単純な関数（例えば $x$ および $x^{2}$）とする。この関数 $y$ のグラフは $y = f(x)$ および $y = F(x)$ の間で振動する曲線である。

90. $x$ が $-\infty$ に近づくときの極限。読者は、"$x$ が $-\infty$ に近づく"や "$x \to -\infty$" の意味を定義するための自己での説明を作成することに難儀しないでしょう。
また
\[
\lim_{x \to -\infty} \phi(x) = l,\quad
\phi(x) \to \infty,\quad
\phi(x) \to -\infty
\]
も同様です。実際、$x = -y$ で $\phi(x) = \phi(-y) = \psi(y)$ と置くと、$y$ が $x$ が $-\infty$ に近づくとき $\infty$ に近づくことがわかります。したがって、$\phi(x)$ の挙動は $x$ が $-\infty$ に近づくときのものと同様に $\psi(y)$ の挙動を $y$ が $\infty$ に近づくときのものと一致します。

91. 第IV章 §63−§67に対応する定理。第IV章で証明された関数の和、積、商に関する定理は、連続変数 $x$ の関数に対してすべて真です（読者は自ずと補うべき明白な文言の変更を難儀することなく供給します）。これらの定理の命題だけでなく、証明も本質的に同じです。

92. 一様に増加または減少する関数。§69に対応する定義は以下の通りです：「関数 $\phi(x)$ が $x$ の増加とともに一様に増加するとされるのは、$x_{2} > x_{1}$ ならば $\phi(x_{2}) \geq \phi(x_{1})$ である」と定義します。もちろん、多くの場合、この条件は一定の値 $x_{0}$ よりも大きい $x$ の範囲でのみ満たされます。そのセクションの次の定理には変更が必要ありませんが、$n$ を $x$ に置き換える必要があります：そして証明も本質的に同じです（明白な文言の変更を除く）。

もし $\phi(x_{2}) > \phi(x_{1})$ であり、等号を除外すると、$x_{2} > x_{1}$ のすべての場合に $\phi(x)$ が「厳密な意味での一様に増加」とされることがわかります。私たちはしばしばこの区別が重要であることを発見します（§108−§109 参照）。

読者は次の関数が$x$と共に安定的に増加するかどうか（または少なくとも一定の$x$値から安定的に増加するか）を考慮すべきです：$x^{2} - x$, $x + \sin x$, $x + 2\sin x$, $x^{2} + 2\sin x$, $[x]$, $[x] + \sin x$, $[x] + \sqrtb{x - [x]}$. このすべての関数は$x \to \infty$と共に$\infty$に収束します。

---

93. $x$が0に向かうときの極限。$\phi(x)$を$\lim\limits_{x \to \infty} \phi(x) = l$となるような$x$の関数とする。そして$y = 1/x$とおく。すると
\[
\phi(x) = \phi(1/y) = \psi(y)
\]
と表すことができる。$x$が$\infty$に向かうとき、$y$は極限0に収束し、$\psi(y)$も極限$l$に収束する。

次に$x$を捨てて$\psi(y)$を単なる$y$の関数として考える。現在考慮しているのは、大きな正の$x$値に対応する$y$の値のみであり、つまり小さな正の$y$の値である。そして$\psi(y)$は、$y$を十分に小さくすれば$\psi(y)$が$l$からどれほどでも小さい差異を持つようにできるという性質を持っている。

より正確には、$\lim\phi(x) = l$という命題は、どんな小さな正の数$\DELTA$も与えられたとき、$x_{0}$を選んで$|\phi(x) - l| < \DELTA$となるような$x$の値がすべて$x_{0}$以上の値であることを意味する。これはまた、$y_{0} = 1/x_{0}$を選んで$|\psi(y) - l| < \DELTA$となるようなすべての正の$y$の値（$y_{0}$以下の値）が存在することを意味する。

したがって以下のような定義に至る：

---

A. どんな小さな正の数$\DELTA$も与えられたとき、$y_{0}(\DELTA)$を選んで
\[
|\phi(y) - l| < \DELTA
\]
となるようにすることができれば、$\phi(y)$は$y$が正の値で$0$に向かうときに極限$l$に収束すると述べ、以下のように表す：
\[
\lim_{y \to +0} \phi(y) = l.
\]

\par\medskip\noindent\textbf{【定義】}\par
\Item{B} 任意の数$\Delta$が与えられたとき、$y_{0}(\Delta)$を選んで
\[
\phi(y) > \Delta
\]
となるようにすることができれば、$\phi(y)$は$y$が正の値で$0$に近づくとき無限大に収束すると定義し、
\[
\phi(y) \to \infty
\]
と表す。
\par\medskip

同様の方法で、$y$が負の値で$0$に近づくとき$\lim\phi(y) = l$の意味を定義する。実際には、定義Aにおいて
\[
0 < y \leq y_{0}(\Delta)
\]
を
\[
-y_{0}(\Delta) \leq y < 0
\]
に変更すればよい。当然、定義Bに対応する類似の定義も存在し、$\phi(y) \to -\infty$となるような定義も同様に行うことができる。

もし$\lim\limits_{y \to +0} \phi(y) = l$かつ$\lim\limits_{y \to -0} \phi(y) = l$であるとき、単に
\[
\lim_{y \to 0} \phi(y) = l
\]
と表す。この場合は非常に重要であり、形式的な定義を付ける価値がある。

\par\medskip\noindent\textbf{【定義】}\par
任意の正数$\Delta$が与えられたとき、$y_{0}(\Delta)$を選んで、すべての$y \neq 0$で$|y| < y_{0}(\Delta)$となる$y$について
\[
|\phi(y) - l| < \Delta
\]
となるようにすることができれば、$\phi(y)$は$y$が$0$に近づくとき$l$に収束すると定義し、
\[
\lim_{y \to 0} \phi(y) = l
\]
と表す。
\par\medskip

同様に、もし$\phi(y) \to \infty$となるような$y \to +0$かつ$y \to -0$のとき、$\phi(y) \to \infty$となるように$y \to 0$とする。また、$\phi(y) \to -\infty$となるような定義も同様に行う。

最後に、もし$y \to +0$のとき$\phi(y)$が極限値や無限大、または負の無限大に収束しない場合、$\phi(y)$は$y \to +0$において振動すると述べ、有限回または無限回振動するかを考慮に入れる。同様に、$y \to -0$における振動も定義する。

これらの定義は変数$y$で表されているが、どの文字を使用しても構わないため、必要に応じて$x$と置き換えることができる。

94. \( x \) が \( a \) に近づくときの極限。もし \( y \to 0 \) のとき \( \phi(y) \to l \) であり、
\[
y = x - a, \quad \phi(y) = \phi(x - a) = \psi(x)
\]
とおくと、\( y \to 0 \) ならば \( x \to a \) で \( \psi(x) \to l \) であり、自然に
\[
\lim_{x \to a} \psi(x) = l,
\]
または単に \( \lim \psi(x) = l \) や \( \psi(x) \to l \) を書き、\( \psi(x) \) が \( x \) が \( a \) に近づくとき \( l \) に収束すると述べる。この方程式の意味は以下の定義によって形式的にとることができる：
\[
\text{与えられた} \ \DELTA \ \text{に対して、常に} \ \EPSILON(\DELTA) \ \text{を決定することができ、}
\]
\[
|\phi(x) - l| < \DELTA
\]
が \( 0 < |x - a| \leq \EPSILON(\DELTA) \) のとき成立すれば、
\[
\lim_{x \to a} \phi(x) = l.
\]

\( x \) が \( a \) より大きい値に近づく場合、すなわち \( 0 < |x - a| \leq \EPSILON(\DELTA) \) を \( a < x \leq a + \EPSILON(\DELTA) \) に置き換えることで、「\( x \) が \( a \) の右側から近づくとき \( \phi(x) \) が \( l \) に収束する」という意味を定義し、これを
\[
\lim_{x \to a+0} \phi(x) = l
\]
と表すことができる。同様にして
\[
\lim_{x \to a-0} \phi(x) = l
\]
の意味も定義することができる。

したがって、
\[
\lim_{x \to a} \phi(x) = l
\]
は
\[
\lim_{x \to a+0} \phi(x) = l, \quad \lim_{x \to a-0} \phi(x) = l
\]
の2つの主張に等しい。

\( x \to a \) のとき \( \phi(x) \to \infty \) または \( \phi(x) \to -\infty \) となる場合についても類似した定義を与えることができるが、これらは \( a = 0 \) の特別な場合における定義と全く同じであり、\( x - a = y \) とおくことで \( y \to 0 \) のときの \( \phi(x) \) の振る舞いを議論することができるため、さらに詳しく述べることはおそらく必要ない。

95. 極限的に増加または減少する関数。もしある数 \( \EPSILON \) が存在して、
\[
a - \EPSILON < x' < x'' < a + \EPSILON \quad \text{のとき} \quad \phi(x') \leq \phi(x'')
\]
ならば、\( \phi(x) \) を \( x = a \) の近傍で極限的に増加する関数と定義する。

まず $x < a$ の場合を考え、$y = 1/(a - x)$ と置きます。すると $x \to a-0$ のとき $y \to \infty$ であり、$\phi(x) = \psi(y)$ は $y$ について安定的に増加する関数で、常に $\phi(a)$ より不大です。この事実から §92 に従って $\phi(x)$ は $\phi(a)$ より不大な極限に向かうことが分かります。これを
\[
\lim_{x \to a+0} \phi(x) = \phi(a+0)
\]
と表記します。

同様の方法で $\phi(a-0)$ を定義できます；そして明らかに
\[
\phi(a-0) \leq \phi(a) \leq \phi(a+0)
\]
が成り立ちます。減少する関数についても類似の考慮事項を適用することができます。

もし $\phi(x') < \phi(x'')$ であり、等号は除かれたとき $a - \EPSILON < x' < x'' < a + \EPSILON$ の条件下で $\phi(x)$ を厳密に増加する関数と定義します。

96. 不確定の極限と収束原理。
§80 から §84 までの議論は、極限 $a$ に向かう連続変数 $x$ の関数にも適用できます。特に $\phi(x)$ が区間 $a - \EPSILON \leq x \leq a + \EPSILON$ を含む範囲で有界である場合（すなわち、$\EPSILON$, $H$, $K$ を見つけることができ、$H < \phi(x) < K$ が成り立つ）には、$x \to a$ のとき $\phi(x)$ の不確定の下限 $\lambda$ と上限 $\Lambda$ を定義し、$\phi(x) \to l$ となるための必要十分条件は $\lambda = \Lambda = l$ であることを証明できます。また、収束原理の類似形も確立できます；すなわち、$x \to a$ のとき $\phi(x)$ が極限に向かうための必要十分条件は与えられた $\DELTA$ に対して $0 < |x_{2} - a| < |x_{1} - a| \leq \EPSILON(\DELTA)$ の条件下で $|\phi(x_{2}) - \phi(x_{1})| < \DELTA$ を満たす $\EPSILON(\DELTA)$ を選べることであることを証明します。

例 XXXV．

1. \( x \to a \)のとき、\(\phi(x) \to l\)、\(\psi(x) \to l'\)ならば、\(\phi(x) + \psi(x) \to l + l'\)、\(\phi(x)\psi(x) \to ll'\)、および\(\frac{\phi(x)}{\psi(x)} \to \frac{l}{l'}\)（ただし最後の場合は \(l' = 0\) の場合を除く）。

2. 正整数 \(m\) に対して、\(x \to 0\) のとき \(x^m \to 0\)。

3. 負整数 \(m\) に対して、\(x \to +0\) のとき \(x^m \to +\infty\)、また \(x \to -0\) のとき \(x^m \to -\infty\) または \(x^m \to +\infty\)（奇数であるか偶数であるかによります）。\(m = 0\) のとき、\(x^m = 1\) であり、\(x^m \to 1\)。

4. \( x \to 0 \)のとき、\(\lim_{x \to 0} (a + bx + cx^2 + \dots + kx^m) = a\)。

5. \( x \to 0 \)のとき、\(\left\{\frac{a + bx + \dots + kx^m}{\alpha + \beta x + \dots + \kappa x^\mu}\right\} \to \frac{a}{\alpha}\)（ただし \(\alpha = 0\) の場合を除く）。\(\alpha = 0\)、\(a \neq 0\)、\(\beta \neq 0\) のとき、関数は \(x \to +0\) のとき \(+\infty\) または \(-\infty\) に収束し、符号が同じか異なるかによります。\(x \to -0\) の場合、この状況は逆転します。両方とも \(a\) と \(\alpha\) が 0 である場合については、例 XXXVI．5 を参照してください。\(a \neq 0\) であり、分母の最初のいくつかの係数が 0 であるような場合についても議論します。

\item[6.] \( \lim_{x \to a} x^{m} = a^{m} \) は、\( m \)が任意の正または負の整数である場合（\( a = 0 \)でないときおよび \( m \)が負でない場合）に成り立つ。[もし \( m > 0 \) のとき、\( x = y + a \)と置き換えて例題4を適用する。\( m < 0 \)のとき、上記例題1から結果が導かれる。これによりすぐに \( \lim P(x) = P(a) \) が成り立つことがわかる。もし \( P(x) \)が任意の多項式である場合]

\item[7.] \( \lim_{x \to a} R(x) = R(a) \) は、\( R \)が任意の有理関数を表し、かつ \( a \)がその分母の根でないとき成り立つ。

\item[8.] \( x^{m} \) の \( m \) がすべての有理数である場合（\( a = 0 \)でないおよび \( m \)が負でない場合）について、\( \lim_{x \to a} x^{m} = a^{m} \) を示す。[これは、\( a \)が正の場合、74節の不等式（9）または（10）からすぐに導かれる。つまり \( |x^{m} - a^{m}| < H|x - a| \) である。ここで \( H \)は \( mx^{m-1} \)および \( ma^{m-1} \)の絶対値の大きい方（例題xxviii.4を参照）。\( a \)が負の場合、\( x = -y \)と \( a = -b \)と置き換える。すると
\[
\lim x^{m} = \lim (-1)^{m} y^{m} = (-1)^{m} b^{m} = a^{m}.
\]
]

\paragraph{97.} 読者には、例題4, 5, 6, 7, 8 のような結果を証明する必要性がすぐに理解できるとは思われないかもしれない。彼は「なぜ \( x = 0 \)や \( x = a \)と置き換えないのか？」と疑問を持つかもしれない。もちろん、そうすると \( a \), \( a/\alpha \), \( a^{m} \), \( P(a) \), \( R(a) \)が得られる。彼は自分がどこで間違っているかを正確に理解することが非常に重要である。したがって、さらに例題を考える前にこの点を慎重に検討する必要がある。

\[
\lim_{x \to 0} \phi(x) = l
\]
これは、$x$がゼロとは異なるが僅差の値を持つとき$\phi(x)$の値についての声明です。\PageSep{170}
つまり、第93節の定義Aにおいて、$0 < y \leq y_{0}$という不等式を最初に含めることで$y = 0$の値を明示的に除外しています。

これは$x = 0$における$\phi(x)$の値についての声明ではありません。我々がこの声明を行うとき、$x$がほぼゼロと等しいとき、$\phi(x)$もほぼ$l$と等しいという主張を行います。しかし、$x$が実際にゼロと等しいときには何が起こるかについては何も主張していません。私たちは知る限りでは、$x = 0$における$\phi(x)$は定義されていないかもしれません；あるいは、$l$以外の値を取っているかもしれません。

例えば、すべての$xに対する関数\phi(x) = 0で定義される関数を考えます。明らかに
\[
\lim_{x \to 0} \phi(x) = 0
\]
が成り立ちます。\Tag{(1)}

今、$\phi(x)$と異なる点のみにおいて$\psi(x) = 1$となる関数$\psi(x)$を考慮します。つまり$x = 0$のとき$\psi(x) = 1$です。このとき
\[
\lim_{x \to 0} \psi(x) = 0
\]
が成り立ちます。\Tag{(2)}
なぜなら、$x$がほぼゼロと等しいとき、$\psi(x)$は不僅ほぼ零だが実際的にも零と等しくなるからです。しかし、$\psi(0) = 1$です。この関数のグラフは$x$軸に点$x = 0$を除き、孤立した一点$(0, 1)$で構成されます。式\Tag{(2)}は、$y$軸に向かってグラフを移動するとき、曲線の縦座標が常に零であり、限界としてゼロに収束することを表します。この事実は孤立した点$(0, 1)$の位置に関係ありません。

読者には、この例が人工的であるとして反論されるかもしれません。しかし、$x=0$ 付近でこのような振る舞いをする関数を簡単に表す式は容易に書くことができます。
一つの例は
\[
\psi(x) = [1 - x^{2}]
\]
です。ここで $[1 - x^{2}]$ は、従来通り $1 - x^{2}$ よりも小さい最大整数を表します。$x=0$ のとき $\psi(x)==1$ であり、$0 < x < 1$ または $-1 < x < 0$ のとき $0 < 1 - x^{2} < 1$ なので $\psi(x) = [1 - x^{2}] = 0$ です。

また、既に第II章 §24, 式(2)で議論した関数
\[
y = \frac{x}{x}
\]
を考えてみましょう。この関数は $x=0$ を除いてすべての $x$ の値に対して $1$ と等しいです。しかし、$x=0$ のとき $y=\frac{x}{x}$ は $1$ ではありません：実際、$x=0$ のときこの関数は定義されていません。つまり、$\phi(x)$ が $x=0$ で定義されているとは（第II章 §24, l.c. 参照）、その値を計算できるという意味です。これは $x=0$ のとき不可能です。$x=\frac{x}{x}$ を代入すると $0/0$ となり、無意味な式となります。読者は「分子と分母を $x$ で割る」と反論するかもしれませんが、$x=0$ のときこれは不可能です。したがって $y = \frac{x}{x}$ は $y = 1$ から $x=0$ で定義されていない点だけ異なる関数です。しかし
\[
\lim_{x \to 0} \frac{x}{x} = 1
\]
であり、$x \neq 0$ のとき $\frac{x}{x}$ は常に $1$ と等しいです。

同様に $\phi(x) = \frac{(x + 1)^{2} - 1}{x} = x + 2$ であり、$x \neq 0$ のときのみこの式が成り立ちますが、$x=0$ のとき未定義です。しかし
\[
\lim_{x \to 0} \phi(x) = 2
\]
です。

一方で、$\phi(x)$ の $x$ が $0$ に近づくときの極限が $\phi(0)$ と等しい場合もあります。つまり、$\phi(x) = x$ のとき $\phi(0) = 0$ であり、$\lim_{x \to 0} \phi(x) = 0$ です。実用的な観点から見れば、つまり最も頻繁に現れる応用の観点から見れば、これは通常の場合です。

例 XXXVI．

1. $\lim\limits_{x \to a} (x^{2} - a^{2})/(x - a) = 2a$.

2. $\lim\limits_{x \to a} (x^{m} - a^{m})/(x - a) = ma^{m-1}$，もし$m$が任意の整数（零を含む）であるならば。

3. $m$が有理数であり$a$が正のとき、例2の結果は依然として成り立つことを示せ。［これは直ちに第74節の不等式(Eq.(9))と(Eq.(10))から導かれる。］

4. $\lim\limits_{x \to 1} (x^{7} - 2x^{5} + 1)/(x^{3} - 3x^{2} + 2) = 1$。［注意してほしいのは、分子と分母の両方に$x-1$が因子であることに。］

5. $x$が正または負の値で$t$に近づくとき、関数
\[
\phi(x) = (a_{0}x^{m} + a_{1}x^{m+1} + \dots + a_{k}x^{m+k})
         /(b_{0}x^{n} + b_{1}x^{n+1} + \dots + b_{l}x^{n+l})
\]
の振る舞いを議論せよ。

［もし$m > n$ならば、$\lim\phi(x) = 0$。もし$m = n$ならば、$\lim\phi(x) = a_{0}/b_{0}$。もし$m < n$で$n - m$が偶数ならば、$a_{0}/b_{0} > 0$のとき$\phi(x) \to +\infty$または$a_{0}/b_{0} < 0$のとき$\phi(x) \to -\infty$。もし$m < n$で$n - m$が奇数ならば、$x \to +0$のとき$\phi(x) \to +\infty$かつ$x \to -0$のとき$\phi(x) \to -\infty$または$x \to +0$のとき$\phi(x) \to -\infty$かつ$x \to -0$のとき$\phi(x) \to +\infty$、これらは$a_{0}/b_{0} > 0$または$a_{0}/b_{0} < 0$により異なる。］

6. 小さな量に関する順序。$x$が小さいとき$x^{2}$は非常に小さく、$x^{3}$はさらに小さくなり、そのような意味で
\[
\lim_{x \to 0} (x^{2}/x) = 0,\quad
\lim_{x \to 0} (x^{3}/x^{2}) = 0,\ \dots.
\]

また、$x$が$t$に近づくとき、$x^{2}$，$x^{3}$，…もすべて$0$に近づき、しかし$x^{2}$は$x$よりも速く$0$に近づき、$x^{3}$は$x^{2}$よりも速く$0$に近づくなどという意味で述べることもできる。関数が$x$に近づくときの極限値が$0$である場合、その関数が$x$とともにどれほど速く減少するかを測る尺度が必要であり、単純な関数$x$，$x^{2}$，$x^{3}$，…を尺度として採用するのが自然である。

従って、$\phi(x)$が初等的小さな量であるとは、$x$が$0$に近づくとき$\frac{\phi(x)}{x}$が$0$以外の極限値を持つことをいう。例えば、$2x + 3x^{2} + x^{7}$は初等的小さな量であり、これは$\lim\left(\frac{2x + 3x^{2} + x^{7}}{x}\right) = 2$であるからである。

同様に、二階、三階、四階などの小さな量を定義することができる。この尺度が完全なものであるとは限らない。もし完全な尺度であれば、$x$と共に$0$に収束する任意の関数$\phi(x)$は初等的小さな量の第一、第二、またはより高い階数のいずれかであるべきである。これは明らかに成り立たない。例えば、$\phi(x) = x^{\frac{7}{5}}$は$x$よりも速く$0$に収束し、$x^{2}$よりも遅く$0$に収束する。

読者は自然と、この尺度を完璧にするために分数の階数も含めることを考えるかもしれない。例えば、$x^{\frac{7}{5}}$は$\frac{7}{5}$次の小さな量であると言えるかもしれない。しかし後で見るように、このような階数の尺度でもなお完全ではないことがわかる。そして実際、上記で定義した整数の階数の方が応用において遥かに重要であり、より正確な定義を試みる必要はほとんどない。

\Topic{大きな量の階数。} $\phi(x)$が$x$が小さいときに正または負の大きい値を持つ場合、類似の定義が提案される。$\phi(x)$が小さなとき$k$次の大きな量であるとは、$x$が$0$に近づくとき$\frac{\phi(x)}{x^{-k}} = x^{k}\phi(x)$が$0$以外の極限値を持つことをいう。

これらの定義は、$x \to 0$ の場合に適用されるものである。当然、$x \to \infty$ や $x \to a$ の場合にも対応する定義がある。

したがって、$x^{k}\phi(x)$ が $x \to \infty$ のとき零以外の極限に収束するならば、$\phi(x)$ は $x$ が大きいとき $k$ 番目の小量であるという。また、$(x - a)^{k}\phi(x)$ が $x \to a$ のとき零以外の極限に収束するならば、$\phi(x)$ は $x$ が $a$ に近いとき $k$ 番目の大量であるという。

7. $\lim\sqrt{1 + x} = \lim\sqrt{1 - x} = 1$。$1 + x = y$ や $1 - x = y$ と置き、例 xxxv．8 を用いる。

8. $\lim\{\sqrt{1 + x} - \sqrt{1 - x}\}/x = 1$。分子分母に $\sqrt{1 + x} + \sqrt{1 - x}$ を乗じる。

9. $x \to 0$ のとき、$m$ と $n$ が正の整数であるとする。$\{\sqrt{1 + x^{m}} - \sqrt{1 - x^{m}}\}/x^{n}$ の振る舞いを考慮する。

10. $\lim\{\sqrt{1 + x + x^{2}} - 1\}/x = \frac{1}{2}$。

11. $\lim\dfrac{\sqrt{1 + x} - \sqrt{1 + x^{2}}}{\sqrt{1 - x^{2}} - \sqrt{1 - x}} = 1$。

12. 関数 $y = \biggl\{\frac{1}{x - 1}
         + \frac{1}{x - \frac{1}{2}}
         + \frac{1}{x - \frac{1}{3}}
         + \frac{1}{x - \frac{1}{4}}\biggr\} \bigg/
    \biggl\{\frac{1}{x - 1}
         + \frac{1}{x - \frac{1}{2}}
         + \frac{1}{x - \frac{1}{3}}
         + \frac{1}{x - \frac{1}{4}}\biggr\}$ のグラフを描く。

$x \to 0$ のとき、極限は存在するか。ここで $y = 1$ であるが、$x = 1$, $\frac{1}{2}$, $\frac{1}{3}$, $\frac{1}{4}$ では定義されていないし、$x \to 0$ のとき $y \to 1$。

13. $\lim\dfrac{\sin x}{x} = 1$。

三角関数の比の定義から導かれるところによれば、もし$x$が正で$\frac{1}{2}\pi$より小さいならば、

\[
\sin x < x < \tan x
\]

または

\[
\cos x < \frac{\sin x}{x} < 1
\]

または

\[
0 < 1 - \frac{\sin x}{x} < 1 - \cos x = 2\sin^{2} \frac{1}{2} x.
\]

さらに$2\sin^{2} \frac{1}{2} x < 2(\frac{1}{2} x)^{2}$である。したがって、

\[
\lim_{x \to +0} \left(1 - \frac{\sin x}{x}\right) = 0
\]

および

\[
\lim_{x \to +0} \frac{\sin x}{x} = 1.
\]

$\frac{\sin x}{x}$は偶関数であるため、その結果が得られる。

14. $\lim \frac{1 - \cos x}{x^{2}} = \frac{1}{2}$。

15. $\lim \frac{\sin \alpha x}{x} = \alpha$。$\alpha = 0$のときもこれが成り立つか？

16. $\lim \frac{\arcsin x}{x} = 1$。$x = \sin y$と置く。

17. $\lim \frac{\tan \alpha x}{x} = \alpha$、$\lim \frac{\arctan \alpha x}{x} = \alpha$。

18. $\lim \frac{\cosec x - \cot x}{x} = \frac{1}{2}$。

19. $x \to 1$のとき、

\[
\lim_{x \to 1} \frac{1 + \cos \pi x}{\tan^{2}\pi x} = \frac{1}{2}.
\]

20. 関数$\sin(1/x)$、$(1/x)\sin(1/x)$、$x\sin(1/x)$は$x \to 0$のときどのように振る舞うか？最初のものは有限に振動し、第二のものは無限に振動し、第三のものは極限$0$に収束する。これらは$x = 0$のときに定義されていない。詳しくは問題集第15題6、7、8を参照のこと。

\item{21.} 次の関数
\[
y = \left( \sin \frac{1}{x} \right) / \left( \sin \frac{1}{x} \right)
\]
が $x$ が $0$ に近づくとき極限値を持つでしょうか？ [\emph{否}. この関数は $\sin(1/x) = 0$ の場合を除いて常に $1$ です；すなわち，$x = \frac{1}{\pi}$, $\frac{1}{2\pi}$, \dots, $-\frac{1}{\pi}$, $-\frac{1}{2\pi}$, \dots のとき，この式は無意味な形 $0/0$ になります。したがって，これらの値では $y$ は $x = 0$ に近い無数の値で定義されません。]

\item{22.} 整数 $m$ に対して，$[x] \to m$ および $x - [x] \to 0$ が $x \to m+0$ のとき，また $[x] \to m-1$, $x - [x] \to 1$ が $x \to m-0$ のときに証明せよ。
\par\medskip

\Paragraph{98. 実数変数の連続関数.} 読者には，連続曲線の意味についてある程度の概念があることでしょう。例えば，図29の曲線 $C$ は連続であり，曲線 $C'$ は一般に連続ですが，$x = \xi'$ および $x = \xi''$ の点では不連続です。
\%[図版: Fig. 29.]
\Figure{29}{p174}

これらのいずれの曲線も関数 $\phi(x)$ のグラフとみなすことができます。関数 $\phi(x)$ のグラフが連続曲線であるとき，その関数を \emph{連続} と呼び，それ以外は不連続と呼びます。この定義を一時的なものとして採用し，より精密な性質を区別するために試みましょう。

まず，関数 $y = \phi(x)$ の性質は各点においてある性質を持つということが明らかです。すべての $x$ に対して連続性を定義するためには，特定の値 $x$ に対する連続性を最初に定義しなければなりません。したがって，例えば点 $P$ に対応する $x = \xi$ の値について固定します。この値 $\xi$ に関連する $\phi(x)$ の特徴的な性質は何でしょうか？

まず第一に，$\phi(x)$ は $x = \xi$ で定義されていることが明らかです。これは自明な条件であり，もし $\phi(\xi)$ が定義されていなければ，曲線から点が欠けてしまいます。

第二に、\emph{$\phi(x)$は$x = \xi$近傍のすべての値に対して定義}されています；すなわち、$\xi$を含む内部を持つ区間で$\phi(x)$が定義されるものが存在します。

第三に、\emph{$x$が$\xi$から左右どちらからも接続するとき、$\phi(x)$は$\phi(\xi)$に収束}します。

これらの性質は、感覚的な視覚によって描かれる曲線の持つ性質を尽きるほど多くはありません。この曲線の視覚的な描写は、直線や円などの特定の曲線から一般化而来。しかし、これらは最も単純で基本的な性質であり：これらの性質を持つ任意の関数のグラフは、実際的に描く限りにおいて、連続曲線が何であるかの幾何学的感覚を満たすでしょう。したがって、これらの性質を選択し、数学的な連続性の概念を表します。これにより以下のような定義に至ります

\par\medskip\noindent\textbf{【定義】}\par
$x = \xi$における関数$\phi(x)$は、$x$が$\xi$から左右どちらからも接続するとき$\phi(x)$が$\phi(\xi)$に収束し、これらのそれぞれの極限が等しい場合に連続であると定義される。
\par\medskip

現在、区間全体における連続性を定義できます。関数$\phi(x)$は、その区間に属するすべての$x$について連続であれば、その区間全体において連続であると定義されます。また、任意の$x$について連続であれば、全域的に連続であると定義されます。したがって、任意の正数$\EPSILON$（$\EPSILON < \frac{1}{2}$）に対して$\DPmod{(\EPSILON, 1 - \EPSILON)}{[\EPSILON, 1 - \EPSILON]}$における$[x]$は連続であり、また$1$と$x$は全域的に連続です。\PageLabel{175}

限界の定義に戻ってみると、我々の定義は`\emph{$\phi(x)$が$x = \xi$において連続であるとは、与えられた$\DELTA$に対して$\EPSILON(\DELTA)$を選べることで$0 \leq |x - \xi| \leq \EPSILON(\DELTA)$ならば$|\phi(x) - \phi(\xi)| < \DELTA$が成り立つ}という意味と等価である'です。

Weしばしば区間$\DPmod{(a, b)}{[a, b]}$で定義される関数を考える必要がある。この場合、連続性の定義において特別な点$a$と$b$に関連する部分をわずかに変更することが便利である。それにより、$x = a$における$\phi(x)$が連続であるとは、$\phi(a + 0)$が存在し且つ$\phi(a + 0) = \phi(a)$であること、また$x = b$における$\phi(x)$が連続であるとは、$\phi(b - 0)$が存在し且つ$\phi(b - 0) = \phi(b)$であることをいう。

99節．最後の節で与えられた連続性の定義は、次のように幾何学的に説明できる。水平線$y = \phi(\xi) - \DELTA$と$y = \phi(\xi) + \DELTA$を引く。このとき、$|\phi(x) - \phi(\xi)| < \DELTA$は曲線上の$x$に対応する点がこれらの二つの直線間に存在することを表す。同様に、$|x - \xi| \leq \EPSILON$は$x$が区間$\DPmod{(\xi - \EPSILON, \xi + \EPSILON)}{[\xi - \EPSILON, \xi + \EPSILON]}$内にあることを表す。したがって、我々の定義は、どんなに近接しても二つの水平線で曲線を挟むことができる垂直帯を二本の垂直線で切り取ることができるという事実を主張する。これは、どの値の$\xi$に対しても曲線$C$（図29）において明らかである。

今後、いくつかの特殊なタイプの関数の連続性について議論する。以下の結果は（私たちはそのとき指摘したように）、第II章で暗黙に仮定されていたものである。

例XXXVII．
1．二つの点において連続な関数の和または積もまた、その点において連続である。商もまた分母がその点で零にならない限り、連続である。［これは直ちに例xxxv．1から従う］

\par\noindent\textbf{・} [2.] 任意の多項式はすべての$x$の値に対して連続である。有理分数は分母が消滅する$x$の値を除き、すべての$x$の値に対して連続である。［これは \Exs{xxxv}.\ 6,~7 の結果から導かれる。］

\par\noindent\textbf{・} [3.] $\sqrt{x}$はすべての正の$x$の値に対して連続である（\Ex{xxxv}.~8）。$x < 0$の場合には定義されていないが、$x = 0$においては終了注記 \SecNo[§]{98} の Remarkにより連続である。同様に $x^{m/n}$も$m$と$n$が任意の正整数で$n$が偶数のときにつき言える。

\par\noindent\textbf{・} [4.] $n$が奇数のときの関数 $x^{m/n}$はすべての$x$の値に対して連続である。

\par\noindent\textbf{・} [5.] $1/x$は$x = 0$において連続ではない。$x = 0$では定義されていないし、$x \to 0$のときには極限に向かうわけでもない。実際、$x \to 0$の正の値により $1/x \to +\infty$ であり、負の値により $1/x \to -\infty$。

\par\noindent\textbf{・} [6.] $m$と$n$が任意の正整数であるときの関数 $x^{-m/n}$は$x = 0$において連続性について議論せよ。

\par\noindent\textbf{・} [7.] 標準的な有理関数 $R(x) = P(x)/Q(x)$は$Q(x) = 0$の根$a$における$x = a$で不連続である。例えば $(x^{2} + 1)/(x^{2} - 3x + 2)$は$x = 1$において不連続である。有理関数の場合、不連続性は（\ia）特定の$x$の値における定義の欠如と（\ib）この値に近づくときの関数が$+\infty$または$-\infty$に向かうという特徴と常に結びついていることに注意せよ。このような特定の種類の不連続点は通常、関数の「無限大」として説明される。一般的な作業において最も頻繁に遭遇する不連続性の一種である。

\par\noindent\textbf{・} [8.] \[
\sqrtb{(x - a)(b - x)},\quad
\sqrtb{(x - a)(b - x)},\quad
\sqrtb{(x - a)/(b - x)},\quad
\sqrtb{(x - a)/(b - x)}
\] の連続性について議論せよ。

\par\noindent\textbf{・} [9.] $\sin x$と$\cos x$はすべての$x$の値に対して連続である。

［私たちは
\[
\sin(x + h) - \sin x = 2\sin \tfrac{1}{2}h \cos(x + \tfrac{1}{2}h)
\]
という関係があり、これは数値的に$h$の数値よりも小さい。］

\item10. \( x \) のどの値において \(\tan x\), \(\cot x\), \(\sec x\), および \(\cosec x\) が連続であるか，不連続であるか？

\item11. \( f(y) \) が \( y = \eta \) のとき連続であり, \( \phi(x) \) が \( x = \xi \) のときに \( \eta \) と等しい連続な関数であるならば, \( f\{\phi(x)\} \) は \( x = \xi \) のときに連続である。

\item12. \( \phi(x) \) が任意の特定の値において連続であるならば, \( \phi(x) \) の多項式, 例えば \( a\{\phi(x)\}^{m} + \dots \), も同様に連続である。

\item13. 次の関数の連続性について議論せよ：
\[
\frac{1}{a\cos^{2} x + b\sin^{2} x}, \quad
\sqrt{2 + \cos x}, \quad
\sqrt{1 + \sin x}, \quad
\frac{1}{\sqrt{1 + \sin x}}.
\]

\item14. \( \sin(1/x) \), \( x\sin(1/x) \), および \( x^{2}\sin(1/x) \) は \( x = 0 \) のときを除き連続である。

\item15. \( x \neq 0 \) のとき \( x\sin(1/x) \) であり, \( x = 0 \) のとき \( 0 \) となる関数はすべての \( x \) の値において連続である。

\item16. \( [x] \) および \( x - [x] \) はすべての整数値において不連続である。

\item17. 次の関数が不連続となる (もし存在するなら) 値 \( x \) を求めよ：\( [x^{2}] \), \( [\sqrt{x}\,] \), \( \sqrt{x - [x]} \), \( [x] + \sqrt{x - [x]} \), \( [2x] \), \( [x] + [-x] \)。

\item18. 不連続の分類。前記の例は異なる種類の不連続性の分類を示唆している。

\par\noindent\qquad\textbf{(1)}  ある定数 \( a \) に近づくとき, \( x \to a \) の値が \( a \) より小さいか大きいかにより, \( \phi(x) \) がそれぞれ \(\phi(a - 0)\) と \(\phi(a + 0)\) に収束するものとする。連続性のために必要かつ十分な条件は, \( x = a \) のとき定義されていることと, \(\phi(a - 0) = \phi(a) = \phi(a + 0)\) であることである。不連続性は様々な形で生じる可能性がある。

\par\noindent\textbf{($\alpha$)}  $\phi(a - 0)$ は $\phi(a + 0)$ と等しいかもしれませんが、$\phi(a)$ は定義されていないかもしれませんし、$\phi(a - 0)$ および $\phi(a + 0)$ とは異なる値を取るかもしれません。例えば、$\phi(x) = x \sin(1/x)$ で $a = 0$ のとき、$\phi(0 - 0) = \phi(0 + 0) = 0$ ですが、$x = 0$ での $\phi(x)$ は定義されていません。または、$\phi(x) = [1 - x^{2}]$ で $a = 0$ のとき、$\phi(0 - 0) = \phi(0 + 0) = 0$ ですが、$\phi(0) = 1$ です。

\par\noindent\textbf{($\beta$)}  $\phi(a - 0)$ と $\phi(a + 0)$ は異なる場合があります。この場合、$\phi(a)$ はどちらか一方に等しいかもしれませんし、両方とも等しくないかもしれませんし、定義されていないかもしれません。最初の例では、$\phi(x) = [x]$ のとき、$\phi(0 - 0) = -1$、$\phi(0 + 0) = \phi(0) = 0$ が示されます；二番目の例では、$\phi(x) = [x] - [-x]$ のとき、$\phi(0 - 0) = -1$、$\phi(0 + 0) = 1$、$\phi(0) = 0$ が示されます；三番目の例では、$\phi(x) = [x] + x \sin(1/x)$ のとき、$\phi(0 - 0) = -1$、$\phi(0 + 0) = 0$ であり、$\phi(0)$ は定義されていません。

これらのいずれのケースでも、$\phi(x)$ が $x = a$ において単純な不連続性を持つと言います。そして、これらのケースに加えて、$x = a$ の一方側のみで $\phi(x)$ が定義されている場合を考えます。その場合、$\phi(a - 0)$ または $\phi(a + 0)$（どちらか一方）は存在しますが、$x = a$ での $\phi(x)$ は定義されていないかもしれませんし、$\phi(a - 0)$ または $\phi(a + 0)$ とは異なる値を取るかもしれません。

$x = a$ の近傍で単調に増加または減少する関数 $\phi(x)$ は $x = a$ において最多単純な不連続性を持つことが明らかです（§95 参照）。

\SubItem{(2)} $\phi(a - 0)$ および $\phi(a + 0)$ のうちの一方（または両方）が存在しない場合もありますが、例えば $\phi(a + 0)$ が存在しない場合、$x \to a+0$ のとき $\phi(x) \to +\infty$ または $\phi(x) \to -\infty$ であり、$\phi(x)$ は $a$ に近づくにつれて極限値や $+\infty$ または $-\infty$ に向かいます。例えば $\phi(x) = 1/x$ や $\phi(x) = 1/x^{2}$ で $a = 0$ のとき、このような状況が生じます。このような場合、$x = a$ を無限大の点と呼びます（例題7参照）。また、$\phi(x)$ は一方から $x \to a$ のとき $\phi(x) \to +\infty$ または $\phi(x) \to -\infty$ であり、他方では定義されていない場合も同様の点を無限大の点と呼びます。

\SubItem{(3)} 簡単な不連続点でもないかつ無限大でもない不連続点は、振動不連続点と呼ばれます。例えば $\sin(1/x)$ および $(1/x)\sin(1/x)$ の $x = 0$ がこのような点です。

\Item{19.} $x = 0$ における関数 $(\sin x)/x$, $[x] + [-x]$, $\cosec x$, $\sqrt{1/x}$, $\sqrt{1/x}$, $\cosec(1/x)$, $\sin(1/x)/\sin(1/x)$ の不連続性の性質は何でしょうか？

\Item{20.} $x$ が有理数であるとき $1$、無理数であるとき $0$ となる関数（第II章、例題xvi.10参照）はすべての $x$ に対して不連続です。また、$x$ の有理数または無理数のみで定義される任意の関数も同様に不連続です。

\PageSep{179}

\Item{21.} $x$ が無理数であるとき $x$、$x$ が有理分数 $p/q$（$p$ および $q$ は整数）であるとき $\sqrt{(1 + p^{2})/(1 + q^{2})}$ となる関数（第II章、例題xvi.11参照）はすべての負の有理数と正の有理数において不連続ですが、正の無理数では連続です。

\Item{22.} 考慮される関数が \okrickRef{第}{IV}章, \Exs{xxxi}で定義されている点において不連続であるか、またその不連続の性質は何か？ [例えば、関数 $y = \lim x^{n}$（例題5）を考える。このとき $y$ は $-1 < x \leq 1$ の範囲で定義される：$-1 < x < 1$ のとき $y$ は $0$ に等しく、$x = 1$ のとき $y$ は $1$ に等しい。点 $x = 1$ と $x = -1$ は単純不連続点である。]

\Paragraph{100. 持続関数の基本的な性質}
持続曲線の概念を \SecNo[§]{98}で与えた分析が最も簡単または自然な方法であるとは必ずしも思われないかもしれない。持続性の概念を分析する別の方法は以下の通りである。点 $A$ と $B$ を $\phi(x)$ のグラフ上の2点とする、それぞれの座標は $x_{0}$,~$\phi(x_{0})$ および $x_{1}$,~$\phi(x_{1})$ である。任意の直線 $\lambda$ を引くが、この直線 $\lambda$ は $A$ と $B$ の間に通る。すると、直感的に言えることは、もし $\phi(x)$ のグラフが持続的であるなら、そのグラフは $\lambda$ に交わることである。

この性質を持続曲線の固有な幾何学的な性質として考えると、$\lambda$ を $x$ 軸に平行とする仮定は実質的に一般性を失うことは明らかである。この場合、点 $A$ と $B$ の縦座標は等しくない：具体的には $\phi(x_{1}) > \phi(x_{0})$ と仮定する。そして $\lambda$ を $\phi(x_{0}) < \eta < \phi(x_{1})$ の直線 $y = \eta$ とする。すると、$\phi(x)$ のグラフが $\lambda$ に交わることとは同値であり、ある $x$ の値が $x_{0}$ と $x_{1}$ の間で存在し、そのとき $\phi(x) = \eta$ であることを意味する。

したがって、持続関数 $\phi(x)$ は以下の性質を有するべきである：\emph{もし
\[
\phi(x_{0}) = y_{0},\quad
\phi(x_{1}) = y_{1},
\]
かつ $y_{0} < \eta < y_{1}$ なら、$x_{0}$ と $x_{1}$ の間で $\phi(x) = \eta$ を満たす $x$ の値が存在する}。言い換えれば、\emph{$x$ が $x_{0}$ から $x_{1}$ に変化するとき、$y$ は少なくとも一度 $y_{0}$ と $y_{1}$ の間のすべての値を取る}ことが必要である。

我達現在証明、$\phi(x)$が$x$の関数として$\S98$で定義された連続性を持つ場合、実際にはその性質を有することである。ある一定の$x$の範囲において、$x_0$の右側に存在し、$\phi(x) < \eta$となる範囲がある。これは$\phi(x_0) < \eta$であり、したがって$x = x_0$における$\phi(x)$は$\eta$未満であるから、$x$が$x_0$に十分近い場合、$\phi(x) - \phi(x_0)$が$\eta - \phi(x_0)$未満であれば成立する。また、$\phi(x)$が$x = x_0$において連続であることから、この条件は$x$が$x_0$に十分近いかぎり満たされる。同様に、ある一定の$x$の範囲において、$x_1$の左側に存在し、$\phi(x) > \eta$となる範囲がある。

$x_0$と$x_1$の間の$x$の値を二つのクラス$L$と$R$に分類する。以下のように定義する：

(1) クラス$L$には、$\phi(x) < \eta$となる$x = \xi$のすべての値$\xi$と、$x_0$から$\xi$までの間のすべての$x$の値を含める。

(2) クラス$R$にはそれ以外のすべての$x$の値を入れる。すなわち、$\phi(\xi) \geq \eta$または$x_0$から$\xi$までの間に$\phi(x) \geq \eta$となる$x$の値が存在するすべての$\xi$を含める。

これらの二つのクラスは$\S17$で課されたクラス$L$、$R$に対する条件を満たすため、実数の区間を形成する。この区間に対応する数を$\xi_0$とする。

まず、$\phi(\xi_0) > \eta$である場合を考える。つまり、$\xi_0$は上クラスに属している。そして、$\phi(\xi_0) = \eta + k$とおく。このとき、すべての$\xi' < \xi_0$に対して$\phi(\xi') < \eta$であり、
\[
\phi(\xi_0) - \phi(\xi') > k,
\]
が成り立つ。これは$x = \xi_0$における連続性の条件に矛盾する。

次に、\(\phi(\xi_{0}) = \eta - k < \eta\) とする。このとき、\(\xi'\)が任意に \(\xi_{0}\) より大きい数であれば、\(\phi(\xi') \geq \eta\) または \(\xi_{0}\) と \(\xi'\) の間にある数 \(\xi''\) について \(\phi(\xi'') \geq \eta\) を満たすものが存在する。いずれの場合も、\(\xi_{0}\) に近い数を選び、その値で \(\phi(x)\) の差が \(k\) より大きいものが見つかる。そしてこれは、\(x = \xi_{0}\) における \(\phi(x)\) の連続性の仮定と矛盾する。

したがって、\(\phi(\xi_{0}) = \eta\) であり、定理は証明される。この証明では、定理に明示的に述べられている以上のことを証明していることに注意すべきである。つまり、\(\xi_{0}\) は \(\phi(x) = \eta\) を満たす \(x\) のうち最小の値であることが証明されている。ある値をとる \(x\) の値の中で最小のものがあるとは必ずしも明らかではなく、連続関数の場合以外は一般的には真ではない。

以下は定理の逆が偽であることを示す例である。図31に表される関数 \(\phi(x)\) のグラフから、\(\phi(x_{0})\) と \(\phi(x_{1})\) の間のすべての値を少なくとも一度はとる：しかし、\(\phi(x)\) は不連続である。実際、\(\phi(x)\) が各値を一度だけとるときでも必ずしも連続とは限らないことも示すことができる。例えば、\(x = 0\) から \(x = 1\) まで \(\phi(x)\) を以下のように定義する。もし \(x = 0\) なら \(\phi(x) = 0\)；もし \(0 < x < 1\) なら \(\phi(x) = 1 - x\)；そして \(x = 1\) なら \(\phi(x) = 1\)。関数のグラフは図32に示されており、点 \(O\)，\(C\) を含みますが、点 \(A\)，\(B\) は含まれていない。明らかに、\(x\) が \(0\) より \(1\) まで変化するとき、\(\phi(x)\) は一度だけ \(\phi(0) = 0\) から \(\phi(1) = 1\) の間のすべての値をとるが、\(x = 0\) および \(x = 1\) では \(\phi(x)\) は不連続である。

実際、初等数学で通常扱われる曲線は、\(y\)が常に同一方向に変化する有限の区間から成り立っています。もし\(y = \phi(x)\)が\(x\)が\(x_0\)から\(x_1\)へ変化するとき常に同一方向に変化するとすれば、すなわち増加または減少するとすれば、連続性の二つの概念は実質的に等価であることが容易に示されます。つまり、\(\phi(x)\)が\(\phi(x_0)\)と\(\phi(x_1)\)の間のすべての値を取り、かつそのときそれが\(\SectNo[§]{98}\)で述べた意味での連続関数であるならば、それは真です。任意の\(x\)の値\(\xi\)が\(x_0\)と\(x_1\)の間にあるとする。\(x \to \xi\)として\(\xi\)より小さい値を取るとき、\(\phi(x)\)は極限\(\phi(\xi - 0)\)に収束します（\SectNo[§]{95}\Add{））。同様に、\(x \to \xi\)として\(\xi\)より大きい値を取るとき、\(\phi(x)\)は極限\(\phi(\xi + 0)\)に収束します。関数が\(x = \xi\)で連続であるためにはかつ 且つ}
\[
\phi(\xi - 0) = \phi(\xi) = \phi(\xi + 0)
\]
が成り立つ必要があります。これらの等式のいずれかが偽である場合、例えば最初の等式が偽であるとすれば、明らかに\(\phi(x)\)は\(\phi(\xi - 0)\)と\(\phi(\xi)\)の間のどの値も取らないことが明白であり、これは我々の仮定に反します。したがって\(\phi(x)\)は連続であると言えます。この節と前節の結果を総合すると、私たちの直感的な連続性の概念は実質的に正確であり、数学的言葉で精密な表現が可能であることが示されます。

\Paragraph{101.} この節および次の節では、連続関数に関するいくつかの一般的な定理を述べ、証明します。
ページ分割点：182

\par\medskip\noindent\textbf{【定理】}\par
もし\(\phi(x)\)が\(x = \xi\)で連続であり、かつ\(\phi(\xi)\)が正であるとする。このとき、正の数\(\EPSILON\)を定めることができ、その結果\(\phi(\xi)\)は区間\(\DPmod{(\xi - \EPSILON, \xi + \EPSILON)}{[\xi - \EPSILON, \xi + \EPSILON]}\)において常に正である。
\par\medskip

に対して$\DELTA = \frac{1}{2}\phi(\xi)$を基本不等式において取り入れることにより，$\EPSILON$を選べば，
\[
\DPtypo{\phi(x) - \phi(\xi)}{|\phi(x) - \phi(\xi)|} < \tfrac{1}{2}\phi(\xi)
\]
が$\DPmod{(\xi - \EPSILON, \xi + \EPSILON)}{[\xi - \EPSILON, \xi + \EPSILON]}$において成立し，次に
\[
\phi(x) \geq \phi(\xi) - |\phi(x) - \phi(\xi)| > \tfrac{1}{2}\phi(\xi) > 0,
\]
となり，$\phi(x)$は正である。明らかに負の値についても対応する定理がある。

以下は定理2である。
もし$x = \xi$において$\phi(x)$が連続であり，$\phi(x)$が$\xi$に近い任意の$x$における値で消滅するか，あるいは正でも負でもあるような値を取るならば，
$\phi(\xi) = 0$である。

これは定理1の明明白白な推論である。もし$\phi(\xi)$がゼロでないなら，それは正か負である；そしてそれが例えば正であれば，$\xi$に十分近い任意の$x$における$\phi(x)$もすべて正であり，これには定理の仮定に反する。

102. 持続関数の値域について。関数$\phi(x)$を考えるが，現在はその定義域が区間$\DPmod{(a, b)}{[a, b]}$における任意の$x$で定義されているというのみとする。

$\DPmod{(a, b)}{[a, b]}$における任意の$x$における$\phi(x)$の値は集合$S$を形成し，これは第80節において我々が第81節において関数$n$の値の集合に適用した議論を用いて扱うことができる。もしある数$K$が存在して$\phi(x) \leq K$が全ての$x$における値について成立するなら，$\phi(x)$は「上界を持つ」という。この場合，$\phi(x)$には「上界」$M$がある：$\phi(x)$のどの値も$M$を超えないが，$M$より小さい任意の数は少なくとも一つの$\phi(x)$の値によって超えられる。同様に，「下界を持つ」「下界」「有界」という用語を連続変数$x$に対する関数に適用することができる。

\par\medskip\noindent\textbf{【定理】}\par
もし$\phi(x)$が$(a, b)$および$[a, b]$において連続であれば、$\phi(x)$は$(a, b)$および$[a, b]$において有界である。
\par\medskip
PerPageSep{183}

確かに、$a$から右に延びる区間$\DPmod{(a, \xi)}{[a, \xi]}$が存在し、その区間に於いて$\phi(x)$は有界であることが決定できます。なぜなら、$\phi(x)$が$x = a$において連続であり、任意に小さな正の数$\DELTA$を与えられた場合でも、$\phi(x)$が$\phi(a) - \DELTA$と$\phi(a) + \DELTA$の間にある区間$\DPmod{(a, \xi)}{[a, \xi]}$を特定できるからです。そして明らかに、この区間に於いて$\phi(x)$は有界である。

現在、区間$\DPmod{(a, b)}{[a, b]}$の点$\xi$を二つのクラス$L$, $R$に分類する。$\phi(\xi)$が$\DPmod{(a, \xi)}{[a, \xi]}$で有界であるとき$\xi$は$L$に属し、そうでないときは$R$に属す。前文から明らかに$L$は存在する：証明したいことは$R$は存在しないことである。もし$R$が存在すると仮定すると、下クラスと上クラスが$L$と$R$のセクションに対応する数$\beta$がある。$\phi(x)$が$x = \beta$で連続であることから、どんなに小さな$\DELTA$を選んでも、以下の区間
\[
\DPmod{(\beta - \eta, \beta + \eta)}{[\beta - \eta, \beta + \eta]}
\]
（もし$\beta = b$の場合はこの区間に$\DPmod{(\beta - \eta, \beta)}{[\beta - \eta, \beta]}$を置き換える）において
\[
\phi(\beta) - \DELTA < \phi(x) < \phi(\beta) + \DELTA
\]
が成り立つ。したがって$\phi(x)$は$\DPmod{(\beta - \eta, \beta + \eta)}{[\beta - \eta, \beta + \eta]}$で有界である。今、$\beta - \eta$は$L$に属する。従って$\phi(x)$は$\DPmod{(a, \beta - \eta)}{[a, \beta - \eta]}$で有界であり、したがって全体の区間$\DPmod{(a, \beta + \eta)}{[\beta - \eta, \beta + \eta]}$でも有界である。しかし$\beta + \eta$は$R$に属するので、$\phi(x)$は$\DPmod{(a, \beta + \eta)}{[a, \beta + \eta]}$で無限大ではない。これは矛盾であり、従って$R$は存在しない。よって$\phi(x)$は全体の区間$\DPmod{(a, b)}{[a, b]}$で有界である。

以下が定理2である。
もし$\phi(x)$が$\DPmod{(a, b)}{[a, b]}$全域で連続であり、その上界と下界を$M$および$m$とするとき、$\phi(x)$は区間に少なくとも一度それぞれ$M$および$m$を取る。

任意の正数$\DELTA$について、ある$x$の値が存在して、$M - \phi(x) < \DELTA$または$\frac{1}{M - \phi(x)} > \frac{1}{\DELTA}$となる。したがって、$\frac{1}{M - \phi(x)}$は有界でない。よって、定理1によれば、連続でもない。

しかし、$M - \phi(x)$は連続関数であり、その分母が零にならない点において$\frac{1}{M - \phi(x)}$は連続である（\Ex{xxxvii}.1）。従って、分母が零になる点が存在する。この点では$\phi(x) = M$となる。

同様にして、$\phi(x) = m$となる点も存在することを示すことができる。

上記の証明は多少巧妙で間接的であり、定理の重要性を考えると、別の証明方法について説明した方が良いかもしれない。しかし、その詳細については一時的に保留する。（\SecNo[§]{104}参照）
\PageSep{184}

例題XXXVIII.
1. $\phi(x) = \frac{1}{x}$（$x = 0$のとき$\phi(x) = 0$）の場合、任意の$x = 0$を含む区間において$\phi(x)$は上界も下界も持たない。
例えば、区間$(-1, +1)$や$[-1, +1]$においてである。

2. $\phi(x) = \frac{1}{x^2}$（$x = 0$のとき$\phi(x) = 0$）の場合、区間$(-1, +1)$や$[-1, +1]$において$\phi(x)$は下界$0$を持つが上界は存在しない。

3. $\phi(x) = \sin\left(\frac{1}{x}\right)$（$x = 0$のとき$\phi(x) = 0$）の場合、$x = 0$において不連続である。任意の区間$(-\DELTA, +\DELTA)$や$[-\DELTA, +\DELTA]$において、下界は$-1$であり上界は$+1$で、これらの値は$\phi(x)$によって無数回取られる。

4. $\phi(x) = x - [x]$の場合、すべての整数$x$において不連続である。区間$(0, 1)$や$[0, 1]$において、下界は$0$であり上界は$1$で、$\phi(x)$は$x = 0$または$x = 1$のときのみ$0$となり、$1$には等しくならない。従って、$\phi(x)$は常にその上界に等しい値を取らない。

\par\noindent\textbf{5.}  \( x \)が無理数のとき\(\phi(x) = 0\)、\( x \)が有理数 fraction \( p/q \)のとき\(\phi(x) = q\)とする。このとき、任意の区間\((a, b)\)または\([a, b]\)において\(\phi(x)\)は下限 \(0\)を持つが上限は存在しない。しかし、\( x = p/q \)のとき\(\phi(x) = (-1)^{p}q\)とする場合、任意の区間において\(\phi(x)\)には上限も下限も存在しない。

\paragraph{103. 区間に於ける関数の振動} 任意に与えられた有界な関数\(\phi(x)\)について、その上界と下界をそれぞれ \(M\) および \(m\) とする。このとき、\(a \leq c \leq b\) の場合、\(M(a, b)\)は \(M(a, c)\) および \(M(c, b)\) の大きい方、また \(m(a, b)\)は \(m(a, c)\) および \(m(c, b)\) の小さい方に等しい。この数値 \(O(a, b)\)，つまり区間\((a, b)\)または\([a, b]\)における\(\phi(x)\)の上界と下界的差を、\(\phi(x)\)在区间\((a, b)\)或は\([a, b]\)的振動と呼ぶ。関数 \(M(a, b)\)，\(m(a, b)\)，\(O(a, b)\) の最も単純な性質は以下の通りである。

\par\medskip\noindent\textbf{【結果】}\par
\Item{(1)} \( a \leq c \leq b \)のとき、\( M(a, b) \)は \( M(a, c) \) および \( M(c, b) \) の大きい方に等しい。また、\( m(a, b) \)は \( m(a, c) \) および \( m(c, b) \) の小さい方に等しい。
\par\medskip

\par\medskip\noindent\textbf{【結果】}\par
\Item{(2)} \( M(a, b) \)は増加する関数であり、\( m(a, b) \)は減少する関数であり、また \( O(a, b) \)も \( b \) の増加と共に増加する。
\par\medskip

\par\medskip\noindent\textbf{【結果】}\par
\Item{(3)} \( O(a, b) \leq O(a, c) + O(c, b) \)である。
\par\medskip

最初の二つの定理は、我々の定義からほぼ即座の結果となります。$\mu$ を $M(a, c)$ と $M(c, b)$ の大きい方とするとき、任意の正数 $\DELTA$ を取ると、$\phi(x) \leq \mu$ が $\DPmod{(a, c)}{[a, c]}$ および $\DPmod{(c, b)}{[c, b]}$ において成り立ち、従って $\DPmod{(a, b)}{[a, b]}$ においても成り立つ。また、$\phi(x) > \mu - \DELTA$ が $\DPmod{(a, c)}{[a, c]}$ や $\DPmod{(c, b)}{[c, b]}$ のいずれかで満たされ、従って $\DPmod{(a, b)}{[a, b]}$ 内でも満たされる。したがって $M(a, b) = \mu$ である。$m$ に関する命題も同様に証明できる。よって（1）は証明され、（2）は明らかなる推論である。

次に、$M_{1}$ を $M(a, c)$ と $M(c, b)$ の大きい方、$M_{2}$ を小さい方とするとき、また $m_{1}$ を $m(a, c)$ と $m(c, b)$ の小さい方、$m_{2}$ を大きい方とする。このとき、$c$ は両方の区間に属するため、$\phi(c)$ は $M_{2}$ より大きくないし $m_{2}$ より小さくない。したがって $M_{2} \geq m_{2}$ であり、これらの数値が同一の区間 $\DPmod{(a, c)}{[a, c]}$ および $\DPmod{(c, b)}{[c, b]}$ に対応するかどうかに関わらず成り立つ。そして
\[
O(a, b) = M_{1} - m_{1} \leq M_{1} + M_{2} - m_{1} - m_{2}.
\]
また
\[
O(a, c) + O(c, b) = M_{1} + M_{2} - m_{1} - m_{2};
\]
であり、（3）は従って成り立つ。

\[
M(a, \beta - \eta) < M,\quad
M(a, \beta + \eta) = M
\]
が成り立つので、式(1)より
\[
M(\beta - \eta, \beta + \eta) = M
\]
である。したがって、$\phi(x)$は$\beta$に近い$x$の値に対して$\phi(x)$の値も$M$に近くなるので、連続関数$\phi(x)$より$\phi(\beta) = M$である。

もし$\beta = a$の場合には$M(a, a + \eta) = M$であり、$\beta = b$の場合には$M(a, b - \eta) < M$で、したがって$M(b - \eta, b) = M$である。これらのいずれのケースでも、以前と同じように議論を進めることができる。

また、区間$PQ$における関数$\phi(x)$の上界$M$を求めることで反復的に二分法を使用する方法でも定理は証明できる。もし区間$PQ$における$\phi(x)$の上界が$M$であり、この区間を二等分すると、そのうちの半分$P_{1}Q_{1}$においても上界が$M$であることが見つかる可能性がある。同様に、セクション71で述べたように、各区間に属する上界が常に$M$であるような一連の区間$PQ$, $P_{1}Q_{1}$, $P_{2}Q_{2}$,...を構築することができる。これらの区間はセクション71と同様に点$T$に収束し、その点における$\phi(x)$の値も容易に$M$であることが証明できる。

105. 直線上の区間の集合。ヘイネ-ボレルの定理。

いま我々は、関数の振動に関するいくつかの命題を証明することに進みます。これらの命題は抽象的ですが非常に重要なものです。特に後で見るように、積分論においても極めて重要です。これらの命題は直線上の区間に関する一般の定理に基づいています。

ページ分割点{186}

与えられた直線上の区間の集合について考えます。すなわち、各メンバーが区間$\DPmod{(\alpha, \beta)}{[\alpha, \beta]}$である集合を指します。これらの区間にどのような制限もかけません；有限個または無限個であってもよいし、重複する場合やしない場合もあります。\footnote
  {「重複」はその明確な意味で使用されます：2つの区間が共通点を持ち、それらのどちらの端点でもない場合は重複します。例えば$\DPmod{(0, \frac{2}{3})}{[0, \frac{2}{3}]}$と$\DPmod{(\frac{1}{3}, 1)}{[\frac{1}{3}, 1]}$は共通点を持ち、重複します。区間$\DPmod{(0, \frac{1}{2})}{[0, \frac{1}{2}]}$と$\DPmod{(\frac{1}{2}, 1)}{[\frac{1}{2}, 1]}$は端点を共有せず、接すると言えます。}
また、いくつかの区間が他の区間に含まれる場合もあります。

補足：
後で再び触れる機会があるため、直線上の区間の集合のいくつかの例を挙げてみましょう。\Page番号{186}

(i) 区間$\DPmod{(0, 1)}{[0, 1]}$が$n$等分された場合、その結果得られる$n$つの区間は非重複で直線全体を被覆する有限な集合を定義します。

\Itemp{(ii)} 区間 $\DPmod{(0, 1)}{[0, 1]}$ の任意の点 $\xi$ を取り、$\xi$ に区間 $\DPmod{(\xi - \EPSILON, \xi + \EPSILON)}{[\xi - \EPSILON, \xi + \EPSILON]}$ を asociate します。ただし、$\EPSILON$ は $1$ より小さい正の数であり、$0$ の場合は $\DPmod{(0, \EPSILON)}{[0, \EPSILON]}$、$1$ の場合は $\DPmod{(1 - \EPSILON, 1)}{[1 - \EPSILON, 1]}$ と asociate します。一般に、任意の区間の一部が区間 $\DPmod{(0, 1)}{[0, 1]}$ 外側に出る部分は reject します。これにより無限個の区間を定義し、多くの区間が重複していることが明らかです。

\Itemp{(iii)} 区間 $\DPmod{(0, 1)}{[0, 1]}$ 内の有理数点 $p/q$ を取り、$p/q$ に区間
\[
\DPmod{\left(\frac{p}{q} - \frac{\EPSILON}{q^{3}}, \frac{p}{q} + \frac{\EPSILON}{q^{3}}\right)}{\left[\frac{p}{q} - \frac{\EPSILON}{q^{3}}, \frac{p}{q} + \frac{\EPSILON}{q^{3}}\right]}
\]
を asociate します。ただし、$\EPSILON$ は $1$ より小さい正の数です。$0$ を $0/1$、$1$ を $1/1$ とみなしますが、これらの場合、区間の一部が区間 $\DPmod{(0, 1)}{[0, 1]}$ 外側にある部分は reject します。これにより無限個の区間を定義し、これらは明らかに重複していることがわかります。なぜなら、$p/q$ 以外の無数の有理点が $p/q$ に関連付けられた区間に存在するからです。

\par\medskip

\par\medskip\noindent\textbf{【定理 (ヘイネ・ボレルの定理)】}\par
与えられた区間 $\DPmod{(a, b)}{[a, b]}$ と、その部分に含まれる各区間が属する集合 $I$ を考える。さらに、$I$ は以下の性質を有しているとする：

\Itemp{(i)} 区間 $\DPmod{(a, b)}{[a, b]}$ の任意の点（$a$ や $b$ 除く）は少なくとも一つの区間に含まれている；

\Itemp{(ii)} $a$ は少なくとも一つの区間の左端点であり、$b$ は少なくとも一つの区間の右端点である。
\par\medskip

次に、集合$I$から有限個の区間を選んで、性質(i)と(ii)を持つ区間の集合を作ることができる。

我達知道$a$是至少一个区间$(a, a_{1})$或$[a, a_{1}]$的左端点。同時，我達也知道$a_{1}$位于至少一个区间$(a_{1}', a_{2})$或$[a_{1}', a_{2}]$内。类似地，$a_{2}$位于某个区间$(a_{2}', a_{3})$或$[a_{2}', a_{3}]$内。明らかに、除非在有限步后$a_{n}$与$b$重合，否則これ个论証可以无限重复。

如果在有限步骤后$a_{n}$确実与$b$重合，則无需进一步証明，因為我達已经得到了一个从$I$的区间中选出的有限集合，并且具有所需性质。若$a_{n}$从未与$b$重合，則点$a_{1}$、$a_{2}$、$a_{3}$等（因為每个点都位于其前驱点的右侧）将趋向于某个极限位置，但これ个极限位置目前看来可以是$(a, b)$内的任意位置。

現在假设从$a$开始执行上述过程的所有可能方式，从而得到所有类型的序列$a_{1}$、$a_{2}$、$a_{3}$等。我達可以証明，必须存在至少一个これ样的序列，在有限步骤后到达$b$。

任意の点$\xi$が$a$と$b$間に存在する場合、二つの可能性がある。すなわち（i）$\xi$は某些列の某些点$a_{n}$の左側に存在するか（ii）そうではないか。これらの点$\xi$をクラス$L$とクラス$R$に分類し、（i）が真である場合クラス$L$、（ii）が真である場合クラス$R$とする。クラス$L$は必ず存在するため、区間$\DPmod{(a, a_{1})}{[a, a_{1}]}$のすべての点がクラス$L$に属している。次にクラス$R$が存在しないことを証明し、すなわち全ての点$\xi$はクラス$L$に属することを示す。

もしクラス$R$が存在すると仮定すると、クラス$L$は完全にクラス$R$の左側に存在し、クラス$L$とクラス$R$は$a$と$b$間の実数に対する区間に該当し、その区間に対応する数$\xi_{0}$が存在する。点$\xi_{0}$は$I$の某些区間内、例えば$\DPmod{(\xi', \xi'')}{[\xi', \xi'']}$に含まれており、$\xi'$はクラス$L$に属し、したがって某些列の某些項$a_{n}$の左側にある。しかし、それにより$\DPmod{(\xi', \xi'')}{[\xi', \xi'']}$を構築された列$a_{1}, a_{2}, a_{3}, \dots$における$a_{n}$に関連する区間$\DPmod{(a_{n}', a_{n+1})}{[a_{n}', a_{n+1}]}$とみなすことができる。そして$\xi''$の左側にあるすべての点は$a_{n+1}$の左側にある。したがってクラス$L$に属する点が$\xi_{0}$の右側に存在し、これはクラス$R$の定義と矛盾する。従ってクラス$R$が存在することは不可能である。

よって全ての点$\xi$はクラス$L$に属することを示した。

$b$は$I$の某些区間の右端点であり、例えば$\DPmod{(b_{1}, b)}{[b_{1}, b]}$とみなすことができる。そして$b_{1}$はクラス$L$に属している。従ってある列$a_{1}, a_{2}, a_{3}, \dots$の某些項$a_{n}$が存在し、$a_{n} > b_{1}$であることが示される。しかしそれにより$\DPmod{(b_{1}, b)}{[b_{1}, b]}$に対応する区間$\DPmod{(a_{n}', a_{n+1})}{[a_{n}', a_{n+1}]}$を$a_{n}$とみなすことができる。そしてこれにより$n$番目の項の後ろの項が$b$と一致する列を得ることができる。したがって必要な性質を持つ有限な区間集合を得ることができ、従って定理は証明される。

\par\medskip\noindent\textbf{【注釈】}\par
この定理の観点から、第 \PageRef{p.}{186} の例を考慮することは有益である。

\Itemp{(i)} ここで定理の条件は満たされていない：点 $1/n$,~$2/n$, $3/n$,~\dots\ は任意の区間 $I$ 内に存在しない。

\Itemp{(ii)} ここで定理の条件は満たされている。点 $\EPSILON$,~$2\EPSILON$, $3\EPSILON$,~\dots, $1 - \EPSILON$ に関連する区間群
\[
\DPmod{(0, 2\EPSILON)}{[0, 2\EPSILON]}, \quad
\DPmod{(\EPSILON, 3\EPSILON)}{[\EPSILON, 3\EPSILON]}, \quad
\DPmod{(2\EPSILON, 4\EPSILON)}{[2\EPSILON, 4\EPSILON]}, \ \dots, \quad
\DPmod{(1 - 2\EPSILON, 1)}{[1 - 2\EPSILON, 1]},
\]
が要求される性質を持つ。

\Itemp{(iii)} この場合、定理を使用して証明することができる：$\EPSILON$ が十分小さいとき、区間 $\DPmod{(0, 1)}{[0, 1]}$ 内に存在する点のうち、任意の区間に含まれないものが存在すること。

もし $\DPmod{(0, 1)}{[0, 1]}$ のすべての点が $I$ の区間内にある（端点を除く）と仮定すると、ある有限数の $I$ の区間が同じ性質を持ち、したがって総長は $1$ より大きいことがわかる。しかし、$q = 1$ に対する任意の2つの区間の総長は $2\EPSILON$ であり、他の任意の値に対する $q - 1$ 個の区間の総長は $2\EPSILON(q - 1)/q^{3}$ である。したがって、任意の有限数の $I$ の区間の合計は級数
\[
1 + \frac{1}{2^{3}} + \frac{2}{3^{3}} + \frac{3}{4^{3}} + \dots,
\]
の総和 $2\EPSILON$ に比例するものより大きいことはあり得ない。この級数は第 VIII 章（\okrickRef{Ch.}{VIII}）で収束することが示される。したがって、$\EPSILON$ が十分小さいとき、すべての点が $I$ の区間に含まれるという仮定は矛盾を生むことになる。

読者には、この証明が無駄に複雑であると感じられるかもしれません。しかし、任意の区間$I$に属さない区間内の点が存在することから、これらの区間の和が1未満であることだけから直ちに導くことは遠从です。彼が援用しようとしている定理は（区間の集合が無限である場合）、明らかでなく、Heine-Borelの定理のような手法を使用して厳密に証明する必要があります。

第106節

現在、Heine-Borelの定理を連続関数$\phi(x)$の振動に関する二つの重要な定理の証明に適用します。

以下が定理Iです。
もし$\phi(x)$が区間$\DPmod{(a, b)}{[a, b]}$全体で連続であれば、この区間に有限個のサブ区間$\DPmod{(a, x_{1})}{[a, x_{1}]}$, $\DPmod{(x_{1}, x_{2})}{[x_{1}, x_{2}]}$, \dots\Add{,} $\DPmod{(x_{n}, b)}{[x_{n}, b]}$を分割することができます。これらの各サブ区間に於いて$\phi(x)$の振動が指定された正数$\DELTA$未満であることが保証されます。

任意の$a$と$b$間にある$\xi$について、$\phi(x)$が$x = \xi$で連続であることから、$\EPSILON$を定めることができます。この$\EPSILON$に対応する区間$\DPmod{(\xi - \EPSILON, \xi + \EPSILON)}{[\xi - \EPSILON, \xi + \EPSILON]}$において$\phi(x)$の振動が$\DELTA$未満であることが確認できます。実際、どの$\xi$とどの$\DELTA$に対しても無限個のこのような区間があることは明らかです。もし特定の$\EPSILON$に対する条件が満たされるなら、より小さい値の$\EPSILON$に対しては当然その条件も満たされます。適切な$\EPSILON$の値は自然に$\xi$によって異なります。一つの$\xi$に対する$\EPSILON$の値が別の$\xi$に対する$\EPSILON$の値と互換的である理由はありません。このような$\xi$に関連する区間を$\DELTA$-区間に呼びます。

もしξ = aであるならば、aとa + εの区間を定めることができ、そのような無数の区間も同様に定められます。これらはaのδ-区間に呼び、bのδ-区間も類似した方法で定義できます。

次に、区間Iを考えます。これは(a, b)内のすべての点についてのすべてのδ-区間から形成される区間の集合です。この集合がHeine-Borelの定理の条件を満たすことは明らかであり、区間に内側にある任意の点は少なくとも一つのIの区間の内側にありますし、aとbは少なくとも一つのそのような区間の端点です。したがって、有限数のIの区間から形成される集合I'を定めることができます。これはI自身と同じ性質を持ちます。

I'の構成する区間は一般に重複します（図34参照）。しかし、これらの端点は(a, b)を有限数の区間I''に分割し、それぞれがI'の区間に含まれており、各々においてφ(x)の振動がδ未満であることを明らかにします。したがって、定理Iは証明されました。

以下に定理IIを述べます。

任意の正数δに対して、ある数ηを定めることができます。このとき、区間(a, b)を長さがη未満のサブ区間にどのような方法でも分割しても、各々のサブ区間におけるφ(x)の振動はδ未満となります。

\par\medskip\noindent\textbf{【定理】}\par[II.]
任意の正数$\DELTA$に対して、ある数$\eta$を定めることができます。このとき、区間$(a, b)$を長さが$\eta$未満のサブ区間にどのような方法でも分割しても、各々のサブ区間における$\phi(x)$の振動は$\DELTA$未満となります。
\par\medskip

$\DELTA_{1} < \frac{1}{2}\DELTA$ とします。定理Iにおいて同様にして、各区間に於いて $\phi(x)$ の振動が $\DELTA_{1}$ よりも小さい有限な部分区間 $j$ を構成します。
これらの部分区間 $j$ のうち最も短いもの $j$ の長さを $\eta$ とします。今、$\DPmod{(a, b)}{[a, b]}$ を $\eta$ よりも小さい長さの部分に分割すると、どのような部分も最多二つの連続する部分区間 $j$ 内に完全に含まれます。
従って、第103節（3）により、$\eta$ よりも小さい長さの部分において $\phi(x)$ の振動は、部分区間に於いて最大の $\phi(x)$ の振動の二倍を超えてはならず、したがって $\DELTA_{1}$ 以下であり、従って $\DELTA$ 以下です。

この定理は、定積分理論において根本的に重要です（第VII章参照）。この定理又は類似の定理を使用しない限り、区間に於いて常に連続な関数がその区間に対する積分を有することを証明することは不可能です。

107節. 複数変数の連続関数。連続性と不连续性の概念は、独立した複数の変数を持つ関数に拡張することができます（第II章、31節以降参照）。このような関数に対するその適用は、私たちはこの章で考慮してきた問題よりも遥かに複雑かつ困難な問題を提起します。ここではこれらの問題について詳細に議論することは不可能ですが、後続において連続な二変数の関数が何を意味するかを知る必要がありますので、その定義を以下に示します。これは第98節最後の定義の直接的な一般化です。

\par\medskip\noindent\textbf{【定義】}\par
二変数 \( x \) および \( y \) の関数 \(\phi(x, y)\) が \( x = \xi \), \( y = \eta \) において連続であるとは、任意の正の数 \(\DELTA\) について、そのどれだけでも小さい \(\EPSILON(\DELTA)\) を選べば
\[
|\phi(x, y) - \phi(\xi, \eta)| < \DELTA
\]
が成り立つときをいう。すなわち、中心が点 \((\xi, \eta)\) で、辺が座標軸に平行であり長さ \(2\EPSILON(\DELTA)\) の正方形を描き、その内部または境界上の任意の点における \(\phi(x, y)\) の値が \(\phi(\xi, \eta)\) よりも \(\DELTA\) 未満に異なることができるときをいう。\footnote{読者はこの定義を説明する図を描いてください。}
\par\medskip

この定義は、関数 \(\phi(x, y)\) が問題の正方形上のすべての点で定義されていることを前提としており、特に点 \((\xi, \eta)\) においても定義されていることを含む。また、定義を別の方法で述べると、\(\phi(x, y)\) が \( x = \xi \), \( y = \eta \) において連続であるとは、\( x \to \xi \), \( y \to \eta \) のどのような方法であっても \(\phi(x, y) \to \phi(\xi, \eta)\) となるときをいう。この文は表面上単純に見えるが、その意味を正確に説明するためには、我々の元の文で用いた不等式が必要である。

二変数の連続関数の和、積、一般に商もまた連続であることが証明できる。二つの変数の多項式はすべての値において連続であり、日常的な解析で頻繁に現れる \( x \) および \( y \) の普通の関数も大抵連続である。すなわち、特別な関係によって結ばれた \( x \) および \( y \) の組に対してのみ不連続である。

\par\medskip

\par\medskip\noindent\textbf{【注釈】}\par
読者には、$\phi(x, y)$ の連続性を$x$ および $y$ の2変数に関して主張することは、それぞれの変数について個別に考慮したときの連続性よりも遥かに多くのことを主張していることに注意してもらいたい。明らかに、もし $\phi(x, y)$ が $x$ および $y$ に対して連続であるなら、任意の固定値を $y$（または $x$）に与えたとき、$\phi(x, y)$ は $x$ （または $y$）に関して連続であることは確かに言える。しかし、その逆は決して真ではない。例えば、

\[
\phi(x, y) = \frac{2xy}{x^{2} + y^{2}}
\]

が $x$ や $y$ が零でないときであり、$x$ または $y$ が零であるとき $\phi(x, y) = 0$ とすると、

もし $y$ に任意の固定値（零であれ非零であれ）を与えたとき、$\phi(x, y)$ は $x$ の連続関数であり、特に $x = 0$ においても連続である。その理由は、$x = 0$ のときの $\phi(x, y)$ の値が零であり、$x \to 0$ のときにはその値が極限として零に収束するからである。

同様にして、$\phi(x, y)$ は $y$ の連続関数であることも示される。しかし、$x = 0$, $y = 0$ のとき $\phi(x, y)$ は連続ではない。その値が零である一方で、$x$ および $y$ が直線 $y = ax$ 上に零に向かうとき、

\[
\phi(x, y) = \frac{2a}{1 + a^{2}},\quad
\lim\phi(x, y) = \frac{2a}{1 + a^{2}},
\]

であり、これは $-1$ から $1$ の間の任意の値を取ることができる。

\Paragraph{108. 隱函数。} 私たちは既に第II章（\okrickRef{Ch.}{II}）で、\emph{隠函数}という概念に遭遇した。例えば、

\[
y^{5} - xy - y - x = 0
\]

が $x$ と $y$ の関係を定めているとき、$y$ は $x$ の \emph{隠函数}である。

しかし、このような方程式が実際には $x$ の関数 $y$ を定めるか、あるいは複数のそのような関数を定めるかは明らかではない。第II章（\okrickRef{Ch.}{II}）では、この仮定を前提としていたが、現在ではその仮定が正当化されているかどうかを検討する立場に立っている。
\par\medskip

We shall有用な用語を用いる。点$(a, b)$の周りに，ある条件が満たされる領域として正方形が存在する場合を考える。このような正方形は$(a, b)$の近傍と呼び，その条件は$(a, b)$の近傍でまたは$(a, b)$の近くで満たされていると言える。これは単にある正方形が存在し，その内全域において条件が満たされることを意味する。変数が一つの場合には，正方形は直線上の区間に置き換わる。

以下は定理である。
(i) $f(x, y)$が点$(a, b)$の近傍で$x$と$y$の連続関数であり，
(ii) $f(a, b) = 0$であり，
(iii) $f(x, y)$がすべての$x$の値に対して，$(a)$の近傍で$y$に関して厳密な意味での単調増加関数であるとき，

(1) 方程式$f(x, y) = 0$を満たす唯一の関数$y = \phi(x)$が存在し，
(2) $\phi(x)$はすべての$x$の値に対して$(a)$の近傍で連続である。

図において、正方形は点$(a, b)$の「近傍」であり、条件(i)と(iii)が満たされています。また、$P$は点$(a, b)$です。$Q$と$R$を図のように取ると、(iii)から$f(x, y)$は$Q$で正であり、$R$で負です。このことから、かつ$f(x, y)$が$Q$および$R$で連続であるので、$QQ'$および$RR'$という直線を$OX$に平行に引くことができます。これにより$R'Q'$は$OY$に平行になり、また$f(x, y)$はすべての点において$QQ'$で正であり、すべての点において$RR'$で負です。特に、$Q'$では$f(x, y)$は正であり、$R'$では負です。したがって、(iii)および§100の性質により、$R'Q'$上の一点$P'$で一度かつしか$f(x, y)$が消えます。

同様の構築により、$RQ$と$R'Q'$間の各軸上に$f(x, y)=0$となる唯一の点を得ることができます。さらに明らかに、この構築は$RQ$の左側でも実行できます。このような点$P'$を集めた集合が求められる関数$y=\phi(x)$のグラフを形成します。

$\phi(x)$が連続であることを証明する必要があります。これは$x \to a$のとき$\phi(x)$の「未確定の極限」の概念を使用することで最も単純に達成されます（§96参照）。もし$x \to a$とするとき、$\lambda$と$\Lambda$を$\phi(x)$の$x \to a$のときの未確定の極限として取ると、明らかに点$(a, \lambda)$および$(a, \Lambda)$は直線$QR$上にある。さらに、ある$x$の値の数列が存在し、この数列を通じて$x \to a$のとき$\phi(x) \to \lambda$であり、かつ$f\{x, \phi(x)\} = 0$であることから、$f(x, y)$は$x$および$y$の連続関数であるので、

\[
f(a, \lambda) = 0
\]

が成り立ちます。これにより$\lambda = b$であり、同様に$\Lambda = b$です。したがって、$x \to a$のとき$\phi(x)$は極限$b$に向かいますので、$x=a$における$\phi(x)$は連続です。さらに明らかに、同じ方法で$x$の近傍において任意の値に対する$\phi(x)$も連続であることが示されます。

明確に述べられるように、定理の真偽は条件(iii)において「増加」を「減少」と変更しても影響を受けない。

例えば、方程式(Eq.(1))を考え、$a = 0$, $b = 0$とします。条件(i)および(ii)が満たされていることが明らかです。さらに
\[
f(x, y) - f(x, y') = (y - y')(y^4 + y^3y' + y^2y'^2 + yy'^3 + y'^4 - x - 1)
\]
は、$x$, $y$, $y'$が十分に小さいとき、$y - y'$の符号と逆の符号を持ちます。したがって条件(iii)（「増加」を「減少」とする）は満たされます。これにより、方程式(Eq.(1))を恒等的に満たし$x$とともに消滅する連続な関数$y$が一意に存在することが分かります。

同じ結論は、方程式
\[
y^2 - xy - y - x = 0
\]
でも得られます。この場合の関数は
\[
y = \frac{1}{2}\left(1 + x - \sqrt{1 + 6x + x^2}\right)
\]
であり、平方根は正です。符号を変更した別の解は$x$とともに消滅しない。

証明において注意すべき点があります。定理の仮定が「$(a, b)$の近傍で」満たされていると仮定しました、すなわち、ある一定の正方形$\xi - \varepsilon \leq x \leq \xi + \varepsilon$, $\eta - \varepsilon \leq y \leq \eta + \varepsilon$において満たされているとします。結論は「$x = a$の近傍で」成り立つ、すなわちある一定の区間$\xi - \varepsilon_1 \leq x \leq \xi + \varepsilon_1$において成り立ちます。結論の$\varepsilon_1$が仮定の$\varepsilon$と一致することを示すものはありませんし、実際これは一般的には成り立たないことが多くあります。

109. 逆関数について。特に$f(x, y)$が$F(y) - x$の形である場合を考えます。このとき以下の定理を得られます。

\par\medskip\noindent\textbf{【結果】}\par
もし $F(y)$ が $y$ の関数で、近傍において連続かつ厳密に単調増加または単調減少である（§95の意味での）とき、かつ $F(b) = a$ であれば、$x = a$ の近傍において $F(y) = x$ を満たす一意的な連続関数 $y = \phi(x)$ が存在する。
\par\medskip

このように定義された関数は、$F(y)$ の逆関数と呼ばれる。

例えば、$y^3 = x$, $a = 0$, $b = 0$ の場合を考える。このとき、定理の条件はすべて満たされる。逆関数は $x = \sqrt{y}$ である。

もし $y^2 = x$ と仮定すると、定理の条件が満たされない。なぜなら、$y^2$ は $y = 0$ を含む任意の区間において単調増加でないからである。$y$ が負のときには減少し、$y$ が正のときには増加する。この場合、定理の結論も成り立たない。なぜなら $y^2 = x$ は $x$ の2つの関数を定義し、$y = \sqrt{x}$ と $y = -\sqrt{x}$ であり、これらはともに $x = 0$ で消失し、それぞれ $x$ の正の値のみで定義されるから、方程式はしばしば2つの解を持ち、しばしば1つの解も持たない。

読者は、より一般的な方程式
\[
y^{2n} = x, \quad y^{2n+1} = x,
\]
同様に考えるべきである。また、興味深い例として、方程式
\[
y^5 - y - x = 0
\]
を考える。これは既に §xiv.7 で考慮されたものである。

同様に、方程式
\[
\sin y = x
\]
は $x$ とともに消失する1つの解を持つ。もちろん無数の解が存在し、これらは $\arcsin x$ の他の値によって与えられ（§xv.10 参照）、この条件を満たさないものも含まれる。

So 近くの x の値についてのみ考慮してきました。F(y) が (a, b) 内で正で一様に増加（または減少）していると仮定します。任意の点 ξ ∈ (a, b) に対して、ξ を含む区間 i およびその全域で一意かつ連続な逆関数 φ\\textit{i(x) を決定できます。

区間 i の集合 I からヘイネ・ボレルの定理により有限部分集合を選出し、(a, b) 全域を被覆します。選ばれた区間 i の部分集合に対応する有限個の逆関数 φ}i(x) が一意的に連続全域で定義されます。

したがって、次の定理を得ます：x = F(y) において F(y) が連続して A から B へと一様に増加すると仮定します。このとき、A から B へ x が増加するにつれて一意的な逆関数 y = φ(x) が存在し、かつその逆関数も連続して A から B へと一様に増加します。

この定理を直接的に得る方法として、より難しい §108 の定理を使わずに証明できます。A < ξ < B を仮定し、条件 (i) a < y < b および (ii) F(y) ≤ ξ を満たす y の値の集合を考えます。この集合は上界 η を持ち、明らかに F(η) ≤ ξ です。もし F(η) < ξ であれば、y > η かつ F(y) < ξ を満たす y の値が存在し、そのとき η は上界でないことが示されます。したがって F(η) = ξ であり、方程式 F(y) = ξ は一意的な解 y = η = φ(ξ) を持つことになります。また η は ξ の増加とともに一様に連続的に増加することから定理が証明されます。

第 V 章の補足例題

\Item{1.} \(a\) および \(b\) がともに零でない場合、次の式が成り立つことを示せ：
\[
ax^n + bx^{n-1} + \cdots + k = ax^n(1 + \epsilon_x),
\]
ここで、\(x\) が大きくなるとき、\(\epsilon_x\) は \(x\) の最初の位数に属する小量である。

\Item{2.} \(P(x) = ax^n + bx^{n-1} + \cdots + k\) であり、\(a\) が零でない場合、\(x\) が増加すると \(P(x)\) は最終的に \(a\) の符号を持つ。また、\(\lambda\) を任意の定数とすると、\(P(x + \lambda) - P(x)\) も同様に \(a\) の符号を持つ。

\Item{3.} 次の式が一般に成り立つことを示せ：
\[
\frac{ax^n + bx^{n-1} + \cdots + k}{Ax^n + Bx^{n-1} + \cdots + K} = \alpha + \left(\frac{\beta}{x}\right)(1 + \epsilon_x),
\]
ここで、\(\alpha = a/A\)、\(\beta = (bA - aB)/A^2\) であり、\(x\) が大きくなるとき、\(\epsilon_x\) は \(x\) の最初の位数に属する小量である。特異な場合を示す。

\Item{4.} 次の式を次の形で表せ：
\[
\frac{ax^2 + bx + c}{Ax^2 + Bx + C} = \alpha + \left(\frac{\beta}{x}\right) + \left(\frac{\gamma}{x^2}\right)(1 + \epsilon_x),
\]
ここで、\(x\) が大きくなるとき、\(\epsilon_x\) は \(x\) の最初の位数に属する小量である。

\Item{5.} 次の極限を示せ：
\[
\lim_{x \to \infty} \sqrt{x}\{\sqrt{x + a} - \sqrt{x}\} = \frac{1}{2}a.
\]

[公式を使用：\(\sqrt{x + a} - \sqrt{x} = \frac{a}{\sqrt{x + a} + \sqrt{x}}\).]

\Item{6.} 次の式が成り立つことを示せ：
\[
\sqrt{x + a} = \sqrt{x} + \frac{1}{2}\left(\frac{a}{\sqrt{x}}\right)(1 + \epsilon_x),
\]
ここで、\(x\) が大きくなるとき、\(\epsilon_x\) は \(x\) の最初の位数に属する小量である。

\Item{7.} \(\alpha\) および \(\beta\) の値を求めよ：\(\sqrt{ax^2 + 2bx + c} - \alpha x - \beta\) が \(x \to \infty\) で零となるように。さらに、\(x \to 0\) で \(\frac{x}{\sqrt{ax^2 + 2bx + c} - \alpha x - \beta}\) の極限を証明せよ。

\Item{8.} 次の極限を求めよ：
\[
\lim_{x \to \infty} x\left\{\sqrt{x^2 + \sqrt{x^4 + 1}} - x\sqrt{2}\right\}.
\]

\Item{9.} \(x \to \frac{1}{2}\pi\) のとき、\(\sec x - \tan x \to 0\) を証明せよ。

\Item{10.} \(x\) が小さいとき、\(\phi(x) = 1 - \cos(1 - \cos x)\) は四次小量であり、\(x \to 0\) で \(\frac{\phi(x)}{x^4}\) の極限を求めよ。

\par\noindent\textbf{・} {11.} \( x \)が小さいとき、\(\phi(x) = x\sin(\sin x) - \sin^{2}x\)は六階の小量であることを証明し、\(x \to 0\)のとき\(\phi(x)/x^{6}\)の極限を求めよ。

\par\noindent\textbf{・} {12.} 円の半径 \(OA\) 上に点 \(P\) があり、この半直線が円を越えて延長されているとする。点 \(P\) から円に接する接線 \(PT\) を引き、接点を \(T\) とし、\(TN\) を \(OA\) に対して垂直な線分とする。点 \(P\) が \(A\) に向かって動くとき、\(NA/AP \to 1\) であることを示せ。

\par\noindent\textbf{・} {13.} 半円弧の中央点と端点に接線を引く。この接線と弧の弦で形成される三角形の面積を \(\Delta\)、端点の接線で形成される三角形の面積を \(\Delta'\) とする。弧の長さがゼロに近づくとき、\(\Delta/\Delta' \to 4\) であることを示せ。

\par\noindent\textbf{・} {14.} \(a + \sin(1/x)\)/\(x\) の極限値を求めよ。（1）\(x \to 0\) のとき、\(\infty\) に収束する \(a > 1\) の場合、（2）\(-\infty\) に収束する \(a < -1\) の場合。また、\(-1 \leq a \leq 1\) のとき、関数は振動する。

\par\noindent\textbf{・} {15.} \(x = p/q\) のとき \(\phi(x) = 1/q\)、\(x\) が無理数のとき \(\phi(x) = 0\) とする。このとき、\(\phi(x)\) はすべての無理数に対して連続であり、すべての有理数に対して不連続である。

\par\noindent\textbf{・} {16.} \Fig{32} のグラフを描いた関数は以下のいずれかの式で表すことができる：
\[
1 - x + [x] - [1 - x],\quad
1 - x - \lim_{n\to\infty} (\cos^{2n+1}\pi x).
\]

\par\noindent\textbf{・} {17.} \(x = 0\) のとき \(\phi(x) = 0\)、\(0 < x < 1/2\) のとき \(\phi(x) = 1/2 - x\)、\(x = 1/2\) のとき \(\phi(x) = 1/2\)、\(1/2 < x < 1\) のとき \(\phi(x) = 3/2 - x\)、\(x = 1\) のとき \(\phi(x) = 1\) とする。この関数は \(0 \leq x \leq 1\) の範囲で \(0\) から \(1\) までのすべての値を一回ずつとるが、\(x = 0\)、\(x = 1/2\)、および \(x = 1\) では不連続である。また、この関数は以下の式で表すことができる：
\[
\tfrac{1}{2} - x - \tfrac{1}{2}[2x] - \tfrac{1}{2}[1 - 2x].
\]

\par\noindent\textbf{・} {18.} \( \phi(x) = x \) であるとき \( x \) が有理数であり、\( \phi(x) = 1 - x \) であるとき \( x \) が無理数。\( x \) が \( 0 \) より増加し \( 1 \) に近づくとき \( \phi(x) \) は \( 0 \) から \( 1 \) までのすべての値を一度だけとるが、任意の \( x \) の値において \( x = \frac{1}{2} \) を除き不連続である。

\par\noindent\textbf{・} {19.} \( -\frac{1}{2}\pi \) より増加し \( \frac{1}{2}\pi \) に近づくとき \( y = \sin x \) は厳密な意味で連続かつ単調に増加し、\( -1 \) から \( 1 \) までである。したがって、\( y = -1 \) より増加し \( y = 1 \) に近づくとき \( x = \arcsin y \) という連続かつ単調に増加する関数が存在することを導く。

\par\noindent\textbf{・} {20.} \( \arctan y \) の絶対値の最小値はすべての \( y \) の値に対して連続であり、\( -\frac{1}{2}\pi \) から \( \frac{1}{2}\pi \) まで \( y \) がすべての実数値を通過するとき単調に増加する。

\par\noindent\textbf{・} {21.} \( x = 0 \), \( y = 0 \) の近傍で方程式
\[
y^{2} - y - x = 0, \quad y^{4} - y^{2} - x^{2} = 0, \quad y^{4} - y^{2} + x^{2} = 0
\]
の解について、\SecNo[§§]{108}--\SecNo{109} の方法で議論する。

\par\noindent\textbf{・} {22.} \( ax^{2} + 2bxy + cy^{2} + 2dx + 2ey = 0 \) および \( \Delta = 2bde - ae^{2} - cd^{2} \) のとき、\( y \) の一つの値は
\[
y = \alpha x + \beta x^{2} + (\gamma + \epsilon_{x}) x^{3}
\]
で与えられる。ここで
\[
\alpha = -d/e, \quad \beta = \Delta/2e^{3}, \quad \gamma = (cd - be) \Delta/2e^{5},
\]
であり、\( \DPtypo{e_{x}}{\epsilon_{x}} \) は \( x \) が小さいとき一次の小量である。

[もし \( y - \alpha x = \eta \) とすると
\[
-2e\eta = ax^{2} + 2bx(\eta + \alpha x) + c(\eta + \alpha x)^{2}
= Ax^{2} + 2Bx\eta + C\eta^{2},
\]
とする。明らかに \( \eta \) は二次の小量、\( x\eta \) は三次の小量、そして \( \eta^{2} \) は四次の小量であり、また \( -2e\eta = Ax^{2} - (AB/e)x^{3} \) で誤差は四次である。]

\par\noindent\textbf{・} {23.} \( x = ay + by^2 + cy^3 \)のとき、\( y \) の一つの値は以下の式で与えられる：
\[
y = \alpha x + \beta x^2 + (\gamma + \epsilon_x) x^3,
\]
ここで、\(\alpha = 1/a\)、\(\beta = -b/a^3\)、\(\gamma = (2b^2 - ac)/a^5\) であり、\(\epsilon_x\) は \(x\) が小さいとき一次の小量である。

\par\noindent\textbf{・} {24.} \( x = ay + by^n \) のとき、\( n \) は 1 より大きい整数とする。このとき、\( y \) の一つの値は以下の式で与えられる：
\[
y = \alpha x + \beta x^n + (\gamma + \epsilon_x) x^{2n-1},
\]
ここで、\(\alpha = 1/a\)、\(\beta = -b/a^{n+1}\)、\(\gamma = nb^2/a^{2n+1}\) であり、\(\epsilon_x\) は \(x\) が小さいとき \((n-1)\) 次の小量である。

\par\noindent\textbf{・} {25.} 方程式 \( xy = \sin x \) の最小の正根は \( y \) にわたって区間 \((0, 1)\) において連続であり、\( y \) が \( 0 \) より大きく \( 1 \) に近づくときから \(\pi\) から \( 0 \) へと単調減少する。［この関数は \((\sin x)/x\) の逆関数であり、\SecNo[§]{109} を適用するとわかる］。

\par\noindent\textbf{・} {26.} 方程式 \( xy = \tan x \) の最小の正根は \( y \) にわたって区間 \((1, \infty)\) において連続であり、\( y \) が \( 1 \) より大きく無限大に向かうときから \( 0 \) から \(\frac{1}{2}\pi\) へと単調増加する。

\chapter{第6章：Chapter VI: DERIVATIVES AND INTEGRALS}

第六章 微分と積分

110. 微分または微係数。曲線の性質について考え直すことにしましょう。その中で自然に想起される最初で最も明白な性質は、前章で見たように、曲線が連続性をもつことを与える性質であり、これは我々の連続関数の定義において体現されています。

初等幾何学における一般的な曲線、すなわち直線、円、放物線などの普通の曲線には当然、多くの一般的な特性があります。これらの特性の中で最も単純で重要なのは、各点において曲線に沿った特定の「方向」を持つこと、つまりは言い換えると、曲線上の各点においてその曲線に対する「接線」を引くことができるということかもしれません。

読者にはおそらく覚えているでしょうが、初等幾何学では、点 \(P\) における曲線の接線は「点 \(Q\) が点 \(P\) に一致する極限位置である直線 \(PQ\) の極限位置」と定義されています。つまり、点 \(Q\) が点 \(P\) に向かって動くとき、直線 \(PQ\) の極限位置が存在すると仮定します。

図（図36）において、点 \(P\) は曲線上の固定点であり、点 \(Q\) は変動点です。\(PM\) と \(QN\) は \(OY\) 方向に平行で、\(PR\) は \(OX\) 方向に平行です。点 \(P\) の座標を \(x, y\)、点 \(Q\) の座標を \(x + h, y + k\) とします。\(h\) は \(N\) が \(M\) の右側にある場合は正、左側にある場合は負となります。

接線 \(PT\) が \(OX\) 方向に角度 \(\psi\) を作るという仮定があります。つまり、接線 \(PT\) が直線 \(PQ\) の極限位置であるということは、点 \(Q\) が曲線上から点 \(P\) に向かって動くとき、角 \(QPR\) の極限値が \(\psi\) であることと等価です。現在、一般の場合と例外的な場合の2つを区別しなければなりません。

The一般の場合、\(\psi\)は\(\frac{1}{2}\pi\)に等しくないため、\(PT\)は\(OY\)に並行でない。この場合、\(RPQ\)は極限\(\psi\)に収束し、
\[
RQ/PR = \tan RPQ
\]
は極限\(\tan \psi\)に収束する。ここで
\[
RQ/PR = (NQ - MP)/MN = \{\phi(x + h) - \phi(x)\}/h;
\]
であり、したがって
\[
\lim_{h \to 0} \frac{\phi(x + h) - \phi(x)}{h} = \tan\psi.
\Tag{(1)}
\]

読者はこれらの式においてすべての長さが適切な符号をもつものとして扱われていることに注意すべきであり、例えば図において\(Q\)が\(P\)の左側にあるとき\(RQ\)は負である；また極限への収束は\(h\)の符号に影響されない。

したがって、曲線\(\phi(x)\)のグラフが点\(P\)で接線を持つと仮定し、その接線が\(x\)軸に対して垂直ではない場合、\(\phi(x)\)は特定の\(x\)の値に対応して、\emph{\(\{\phi(x + h) - \phi(x)\}/h\)が\(h\)が0に収束するとき極限に収束すること}を持つ性質がある。

以下は注意点である。
これは当然、
\[
\{\phi(x + h) - \phi(x)\}/h,\quad
\{\phi(x - h) - \phi(x)\}/(-h)
\]
が\(h\)が正の値で0に収束するとき極限に収束し、両方の極限が等しいことを意味する。これらの極限が存在して等しくない場合、曲線\(y = \phi(x)\)は特定の点において角度を持つことになり、図37のように見える。

今、曲線が（円や楕円のように）長さのどの点でも接線を持ち、または特定の $x$ の変化範囲に対応する長さのどの部分でも接線を持つと仮定しましょう。さらに、この接線は $x$ 軸に対して垂直にはならないと仮定します。円の場合では、これは当然半径より長い弧を考慮することを意味します。

このような曲線について、式（1）のような方程式が $x$ の値の範囲内のすべての値で成り立ちます。各 $x$ の値に対応する $\tan\psi$ があります：$\tan\psi$ は $x$ の関数であり、考慮される $x$ の値の範囲内で定義され、元の関数 $\phi(x)$ から計算または導き出せます。この関数を $\phi(x)$ の導関数と呼び、その記号を
\[
\phi'(x)
\]
で表します。

$\phi(x)$ の導関数の別の名前は微分係数であり、$\phi(x)$ から $\phi'(x)$ を計算する操作は一般的に微分と呼ばれます。この用語は歴史的な理由により確立されています：参照 \S115。

$\psi = \frac{1}{2}\pi$ の場合についての特別な例を考慮する前に、導関数 $\phi'(x)$ を $a \leq x \leq b$ のすべての $x$ 値に対して存在することから、$\phi(x)$ がこの区間内の各点で連続であることを説明します。これは明らかに $\{\phi(x + h) - \phi(x)\}/h$ が極限に向かうためには $\lim\phi(x + h) = \phi(x)$ であることが必要であり、これが連続性を表す性質です。

\Item{(2)} 逆が真であるかどうか尋ねるのは自然です。すなわち、連続曲線は各点において一定の接線を持ち、また各点において連続な関数は任意の引数$x$に対して微分係数を持つかどうかです。\図[2インチ]{37}{p199}
ただし、曲線が$OX$に垂直な接線を有する例外的な場合（この可能性を除いて）を考慮に入れない限り、上記の質問の二つの形式は等価です。

答えは明らかに「ノー」です。角度$\alpha$と$\beta$を形成する二直線で構成される曲線を考えれば十分です。\図（Fig{37}）。
ページ分割\pageref{p.200}

読者はすぐに気づくでしょう。この場合、$\{\phi(x + h) - \phi(x)\}/h$は$h \to 0$の正の値により$\tan\beta$に、負の値により$\tan\alpha$に極限を持つ。

以下は注記です。
これは当然、曲線が点において合理的には二つの方向を持っていると述べることができる一例です。しかし、少し難しい例では、連続曲線がその一点において一つの方向も複数の方向もないと言える場合があることを示します。$x\sin(1/x)$という関数のグラフ（Fig{14}、PageRef{p.53}）を描きます。この関数は$x = 0$で定義されていないため、連続性が$x = 0$において欠けます。一方、方程式
\[
\phi(x) = x\sin(1/x)\quad (x \neq 0),\qquad
\phi(x) = 0\quad (x = 0)
\]
によって定義される関数は$x = 0$において連続であり（Exs{xxxvii}.14,15）、この関数のグラフは連続曲線です。

しかし、$\phi(x)$は$x = 0$において微分係数を持ちません。$\phi'(0)$は定義により$\lim\{\phi(h) - \phi(0)\}/h$または$\lim\sin(1/h)$であり、このような極限が存在しないからです。

関数$x$の関数が連続であっても、任意の$x$の値に対して微分係数がないことが示されています。しかし、その証明は非常に困難です。関心のある読者はBromwichの『無限級数』pp.490-1またはHobsonの『実変数函数論』pp.620-5を参照してください。

(3) 微分または微分係数の概念は、幾何学的な考慮から提示されました。しかし、その概念自体には幾何学的要素はありません。関数$\phi(x)$の導関数$\phi'(x)$は、任意の$x$の特定の値に対してこの極限が存在するか否かにより定義されます：
\[
\phi'(x) = \lim_{h \to 0} \frac{\phi(x + h) - \phi(x)}{h}
\]
曲線の幾何学は、導関数という概念が数学の多くの分野で用いられる一つの例に過ぎません。

（追記）別の重要な応用例は動力学です。粒子が一定の直線上を移動し、時間$t$における固定点からの距離が$s = \phi(t)$であるとします。そのとき、「時間$t$における粒子の速度」は定義により以下の極限に等しい：
\[
\frac{\phi(t + h) - \phi(t)}{h}
\]
$h \to 0$ のときの値です。「速度」という概念は実質的に関数の導関数という概念の一例に過ぎません。

例題XXXIX．
1. $\phi(x)$が定数であるとき、$\phi'(x) = 0$です。この結果を幾何学的に解釈してください。
2. $\phi(x) = ax + b$であるとき、$\phi'(x) = a$です。形式的定義から（i）証明し、（ii）幾何学的な観点からも証明してください。
3. $m$が正の整数であるとき、$\phi(x) = x^{m}$の場合、$\phi'(x) = mx^{m-1}$です。

[以下は省略します]
\[
ZZZMATHTOKEN0ZZZ
\]

The reader should observe that this method cannot be applied to \( x^{p/q} \), where \( p/q \) is a rational fraction, as we have no means of expressing \( (x + h)^{p/q} \) as a finite series of powers of \( h \). We shall show later on (\S118) that the result of this example holds for all rational values of \( m \). Meanwhile, the reader will find it instructive to determine \( \phi'(x) \) when \( m \) has some special fractional value (e.g., \( \frac{1}{2} \)), by means of some special device.]

4. もし \(\phi(x) = \sin x\) であるなら, \(\phi'(x) = \cos x\)；そして \(\phi(x) = \cos x\) のとき, \(\phi'(x) = -\sin x\)。

例えば, \(\phi(x) = \sin x\) の場合, 
\[
\frac{\phi(x + h) - \phi(x)}{h} = \frac{2\sin \frac{1}{2}h \cos(x + \frac{1}{2}h)}{h}
\]
であり, \( h \to 0 \) のときのこの式の極限は \(\cos x\) である。これは \(\lim \cos(x + \frac{1}{2}h) = \cos x\) (コサインが連続関数であることから) および \(\lim \{\sin \frac{1}{2}h / \frac{1}{2}h\} = 1\) (\Ex{xxxvi}. 13) の結果である。

5. 曲線 \( y = \phi(x) \) の接線と法線の方程式。曲線上の点 \((x_{0}, y_{0})\) における接線は, \( OX \) との角度 \(\psi\) を作る直線であり, \(\tan\psi = \phi'(x_{0})\) である。その方程式は
\[
y - y_{0} = (x - x_{0}) \phi'(x_{0})
\]
である。そして, 接点における法線（接線に対する垂直線）の方程式は
\[
(y - y_{0}) \phi'(x_{0}) + x - x_{0} = 0
\]
である。接線が \( y \) 軸に平行でないという仮定をした。この特殊な場合において, 接線と法線はそれぞれ \( x = x_{0} \) および \( y = y_{0} \) であることが明らかである。

6. 抛物線 \( x^{2} = 4ay \) の任意の点における接線と法線の方程式を書き下し, 点 \((x_{0}, y_{0})\) で接線が \( x = my + (a/m) \) であることを示せ。ここで \( x_{0} = 2a/m \), \( y_{0} = a/m^{2} \) である。

\paragraph{112.} 我々は、$\phi(x)$が$x$の特定の値で連続でない場合、その値での導関数を持つことはあり得ないと見た。よって、$x = 0$で定義されていない$1/x$や$x = 0$で必ず不連続となる$\sin(1/x)$のような関数は、$x = 0$での導関数を持たない。また、全ての整数値での不連続性を持つ関数$[x]$も、そのような値での導関数を持っていない。

\par\medskip\noindent\textbf{【注釈】}\par
\Paragraph{例.} $[x]$が任意の二つの整数値間で定数であるため、その導関数は存在する場合でも常に零となる。したがって、$[x]$の導関数を$[x]'$と表すとき、それは全ての$x$の値で零であり、整数値では未定義な関数となる。関数 $1 - \frac{\sin\pi x}{\sin\pi x}$ は、この性質と全く同じであることに注意するのも興味深い。
\par\medskip

私たちはまた、\Ex{xxxvii}.7 で、多項式や有理関数、三角関数などの最も単純で明白な関数において発生する最も一般的な不連続性のタイプは、$\phi(x) \to +\infty$ または $\phi(x) \to -\infty$ のような関係と関連付けられていることを確認しました。これらのすべての場合において、それらが以前考慮された場合と同じように、ある特定の $x$ の値では導関数が存在しません。
実際、第111節（§）の \Eq{(1)} で指摘したように、$\phi(x)$ のすべての不連続性は $\phi'(x)$ の不連続性でもあります。しかし、逆は成り立たないことが容易に確認できます。例えば、第110節（§）の幾何学的な視点に戻って、$OY$ に平行な $\phi(x)$ のグラフが接線を持つ特殊な場合を考えます。この場合をさらに分類すると、典型的ないくつかのケースが \Fig{38} で示されています。(\ic) や (\id) の場合では、関数は $P$ の一方側では二値であり、他方では定義されていません。このような場合、$P$ の一方側または他方で $\phi(x)$ の値の集合を、異なる関数 $\phi_{1}(x)$ と $\phi_{2}(x)$ として定義することができます。曲線の上部は $\phi_{1}(x)$ に対応します。

(\ia) の場合では
\[
\frac{\phi(x + h) - \phi(x)}{h} \to +\infty,
\]
$h \to 0$ のとき、そして (\ib) の場合では
\[
\frac{\phi(x + h) - \phi(x)}{h} \to -\infty;
\]
(\ic) の場合では
\[
\frac{\phi_{1}(x + h) - \phi_{1}(x)}{h} \to +\infty,\quad
\frac{\phi_{2}(x + h) - \phi_{2}(x)}{h} \to -\infty,
\]
そして (\id) の場合では
\[
\frac{\phi_{1}(x + h) - \phi_{1}(x)}{h} \to -\infty,\quad
\frac{\phi_{2}(x + h) - \phi_{2}(x)}{h} \to +\infty,
\]
が成り立ちます。ただし、(\ic) の場合では正の $h$ のみ、(\id) の場合では負の $h$ のみを考慮することができ、これは導関数の存在を否定します。

これらの四つの例は、方程式によって定義される関数を考えることにより得られます。
\[
\Item{(\ia)} y^{3} = x,\quad
\Item{(\ib)} y^{3} = -x,\quad
\Item{(\ic)} y^{2} = x,\quad
\Item{(\id)} y^{2} = -x,
\]
特別な$x$の値を考えるとき、その値は$x=0$です。

113. 微分に関するいくつかの一般的な規則。その後の定理において、関数$f(x)$と$F(x)$が考慮される$x$の値に対して微分$f'(x)$と$F'(x)$を持つことを仮定します。

結果
\Item{(1)} $\phi(x) = f(x) + F(x)$の場合、$\phi(x)$は微分
\[
\phi'(x) = f'(x) + F'(x).
\]
を持ちます。

結果
\Item{(2)} $\phi(x) = kf(x)$で$k$が定数の場合、$\phi(x)$は微分
\[
\phi'(x) = kf'(x).
\]
を持ちます。

これらの結果を一般の定理から導く演習として読者に残します。これはEx{xxxv}.1において述べた一般の定理から導けます。

結果
\Item{(3)} $\phi(x) = f(x)F(x)$の場合、$\phi(x)$は微分
\[
\phi'(x) = f(x)F'(x) + f'(x)F(x).
\]
を持ちます。

これは
\begin{align*}
\phi'(x)
  &= \lim\frac{f(x + h)F(x + h) - f(x)F(x)}{h}\\
  &= \lim\left\{f(x + h)\frac{F(x + h) - F(x)}{h}
              + F(x)\frac{f(x + h) - f(x)}{h}\right\}\\
  &= f(x)F'(x) + F(x)f'(x).
\end{align*}
から導けます。

結果
\Item{(4)} $\phi(x) = \dfrac{1}{f(x)}$の場合、$\phi(x)$は微分
\[
\phi'(x) = -\frac{f'(x)}{\{f(x)\}^{2}}.
\]
を持ちます。

この定理において、当然$f(x)$が考慮される特別な$x$の値に対して0でないことを仮定します。すると
\[
\phi'(x)
  = \lim \frac{1}{h} \left\{\frac{f(x) - f(x + h)}{f(x + h)f(x)}\right\}
  = -\frac{f'(x)}{\{f(x)\}^{2}}.
\]

結果
\Item{(5)} $\phi(x) = \dfrac{f(x)}{F(x)}$の場合、$\phi(x)$は微分
\[
\phi'(x) = \frac{f'(x)F(x) - f(x)F'(x)}{\{F(x)\}^{2}}.
\]
を持ちます。

これは(3)と(4)から即座に導けます。

結果
\Item{(6)} $\phi(x) = F\{f(x)\}$の場合、$\phi(x)$は微分
\[
\phi'(x) = F'\{f(x)\} f'(x).
\]
を持ちます。

For let
\[
f(x) = y,\quad
f(x + h) = y + k.
\]
すると $k \to 0$ として $h \to 0$ であり、$k/h \to f'(x)$。そして
\begin{align*}
\phi'(x)
  & = \lim \frac{F\{f(x + h)\} - F\{f(x)\}}{h}\\
  & = \lim \left\{\frac{F(y + k) - F(y)}{k}\right\}
    × \lim \left(\frac{k}{h}\right)\\
  & = F'(y)f'(x).
\end{align*}

この定理は（2）と（4）を特殊な場合として含み、$F(x) = kx$ や $F(x) = 1/x$ のときによく見ることができます。また興味深い特殊な場合としては、$f(x) = ax + b$ のときがあり、この定理は $F(ax + b)$ の導関数が $aF'(ax + b)$ であることを示します。

最後の定理には若干の前説が必要です。$x = \psi(y)$ とし、$\psi(y)$ が一定区間において連続かつ単調に増加（または減少）しているとします（厳密な意味での第95節参照）。このとき $y = \phi(x)$ と書き換えることができ、$\phi$ は $\psi$ の「逆」関数（第109節参照）です。

\par\medskip\noindent\textbf{【結果】}\par
(7) $y = \phi(x)$ が $\phi$ が $\psi$ の逆関数であり、$x = \psi(y)$ で、$\psi(y)$ の導関数 $\psi'(y)$ が零と異なるとき、$\phi(x)$ の導関数は
\[
\phi'(x) = \frac{1}{\psi'(y)}
\]
である。
\par\medskip
ページ番号分離：205

もし $\phi(x + h) = y + k$ とするとき、$k \to 0$ として $h \to 0$ であり、

\[
\phi'(x)
  = \lim_{h \to 0} \frac{\phi(x + h) - \phi(x)}{(x + h) - x}
  = \lim_{k \to 0} \frac{(y + k) - y}{\psi(y + k) - \psi(y)}
  = \frac{1}{\psi'(y)}.
\]

最後の関数は $y = \phi(x)$ の関係を用いて $x$ で表すことができ、$\phi'(x)$ は $\psi'\{\phi(x)\}$ の逆数となります。この定理により、逆関数の導関数が分かれば任意の関数を微分することができます。

114. 複素関数の導函数。これまで、$y = \phi(x)$ を$x$に関する純粋実数関数として考えていた。もし$y$が複素関数$\phi(x) + i\psi(x)$であるなら、$y$ の導函数を$\phi'(x) + i\psi'(x)$とする。読者は定理(1)～(5)が$\phi(x)$が複素数関数の場合にも依然として有効であることを難なく理解できるだろう。定理(6)と(7)も複素関数に対して類似のものがあるが、これらは「複素変数の関数」という一般的概念に依存しており、現在までには少数の特定の例しか遭遇していない。

115. 微分法の表記。我々が言う「導函数」はしばしば「微分係数」と呼ばれる。異なる名前だけでなく、異なる表記法も用いられることがある；関数$y = \phi(x)$ の導函数はしばしば以下のいずれかの表現で表される
\[
D_{x}y,\quad
\frac{dy}{dx}.
\]
これらの中で最後のもの$\frac{dy}{dx}$が最も一般的かつ便利である：しかし読者は$\frac{dy}{dx}$が「ある数$dy$を別の数$dx$で割る」という意味ではなく、「$y = \phi(x)$に対して特定の操作$D_{x}$または$d/dx$を適用した結果」であることを忘れないように注意しなければならない。この操作は$\{\phi(x + h) - \phi(x)\}/h$という商を作るもので、$h \to 0$とするものである。

\par\medskip\noindent\textbf{【注釈】}\par
当然、初見では不思議な記号が採用された理由は何かのためであり、その理由は次のように説明できます。分数$\{\phi(x + h) - \phi(x)\}/h$ の分母$h$は独立変数$x$の値$x+h$と$x$の差に相当し、分子は従属変数$y$の値$\phi(x + h)$と$\phi(x)$の差に相当します。これらの差はそれぞれ$x$および$y$の増分と呼ばれ、それぞれ$\delta x$および$\delta y$で表されます。すると分数は$\delta y/\delta x$となり、多くの目的においてその分数の極限を表すために$dy/dx$で表すことが便利です。しかし、この記号は現在の段階では純粋に象徴的なものとして扱うべきであり、$dy$および$dx$は分離できないものであり、これらが独立して存在するときには意味を持たない：特に$\lim\delta y$と$\lim\delta x$は単にゼロに等しい。読者はこの記号に慣れなければなりませんが、その記号が彼を混乱させる限りは、微分係数を$D_{x}y$の形で表すか、または先ほどの章の前半部で用いた$\phi(x)$および$\phi'(x)$の記法を使用することで回避すべきです。

ただし、第VII章では、$dx$および$dy$の符号を独立した意味を持つように定義する方法について示し、実際にはその商である導関数$dy/dx$が実際にそれらの比であることを説明します。
\par\medskip

第113節の定理は当然、この記号を使って直ちに翻訳できます。以下のように述べることができます：

\par\medskip\noindent\textbf{【結果】}\par
\Item{(1)} $y = y_{1} + y_{2}$ のとき
\[
\frac{dy}{dx} = \frac{dy_{1}}{dx} + \frac{dy_{2}}{dx};
\]

\Item{(2)} $y = ky_{1}$ のとき
\[
\frac{dy}{dx} = k\frac{dy_{1}}{dx}.
\]

\Item{(3)} $y = y_{1}y_{2}$ のとき，
\[
\frac{dy}{dx} = y_{1}\frac{dy_{2}}{dx} + y_{2}\frac{dy_{1}}{dx};
\]

\Item{(4)} $y = \dfrac{1}{y_{1}}$ のとき，
\[
\frac{dy}{dx} = -\frac{1}{y_{1}^{2}}\, \frac{dy_{1}}{dx};
\]

\Item{(5)} $y = \dfrac{y_{1}}{y_{2}}$ のとき，
\[
\frac{dy}{dx}
= \biggl(y_{2}\frac{dy_{1}}{dx} - y_{1}\frac{dy_{2}}{dx}\biggr) \bigg/ y_{2}^{2};
\]

\Item{(6)} $y$ が $x$ の関数であり、$z$ が $y$ の関数のとき，
\[
\frac{dz}{dx} = \frac{dz}{dy}\, \frac{dy}{dx};
\]
\CenterLine{\Item{(7)}}{$\dfrac{dy}{dx} = 1 \bigg/ \biggl(\dfrac{dx}{dy}\biggr)$.}
\par\medskip

\par\medskip\noindent\textbf{【例 XL.】}\par
\Item{1.} $y = y_{1}y_{2}y_{3}$ のとき，
\[
\frac{dy}{dx}
  = y_{2}y_{3}\, \frac{dy_{1}}{dx}
  + y_{3}y_{1}\, \frac{dy_{2}}{dx}
  + y_{1}y_{2}\, \frac{dy_{3}}{dx},
\]
また $y = y_{1}y_{2} \dots y_{n}$ のとき，
\[
\frac{dy}{dx}
  = \sum_{r=1}^{n} y_{1}y_{2} \dots y_{r-1}y_{r+1} \dots y_{n}\, \frac{dy_{r}}{dx}.
\]
特に $y = z^{n}$ のとき、$\dfrac{dy}{dx} = nz^{n-1}\left(\dfrac{dz}{dx}\right)$；また $y = x^{n}$ のとき，
$\dfrac{dy}{dx} = nx^{n-1}$ であり、これについては \Ex{xxxix}.3 で他に証明されている。

\Item{2.} $y = y_{1}y_{2}\dots y_{n}$ のとき，
\[
\frac{1}{y}\, \frac{dy}{dx}
  = \frac{1}{y_{1}}\, \frac{dy_{1}}{dx}
  + \frac{1}{y_{2}}\, \frac{dy_{2}}{dx} + \dots
  + \frac{1}{y_{n}}\, \frac{dy_{n}}{dx}.
\]
特に $y = z^{n}$ のとき、$\dfrac{1}{y}\, \dfrac{dy}{dx} = \dfrac{n}{z}\, \dfrac{dz}{dx}$。
\par\medskip

\Paragraph{116. 标准形式。} 今後、少数の最も基本的な関数型の導函数の形についてより系統的に調査する予定である。

多項式について\Topic{\Item{A.} Polynomials.} において、もし $\phi(x) = a_{0}x^{n} + a_{1}x^{n-1} + \dots + a_{n}$ であるとき，
\[
\phi'(x) = na_{0}x^{n-1} + (n - 1)a_{1}x^{n-2} + \dots + a_{n-1}.
\]
多項式の標準形として、次数 $n$ の $x$ に関する多項式を二項係数形式で表すことが時々便利である。すなわち，
\[
a_{0}x^{n}
  + \binom{n}{1} a_{1}x^{n-1}
  + \binom{n}{2} a_{2}x^{n-2} + \dots + a_{n}.
\]
この場合、
\[
\phi'(x) = n \left\{
    a_{0}x^{n-1}
  + \binom{n - 1}{1} a_{1}x^{n-2}
  + \binom{n - 1}{2} a_{2}x^{n-3} + \dots
  + a_{n-1}
\right\}.
\]

$\phi(x)$ の二項係数形式はしばしば符号的に $\left(a_{0}, a_{1}, \dots, a_{n} \btw x, 1\right)^{n}$ と表される；そして
\[
\phi'(x) = n\left(a_{0}, a_{1}, \dots, a_{n-1} \btw x, 1\right)^{n-1}.
\]

後で見るように、$\phi(x)$ は常に $n$ 個の因子の積として表すことができる：
\[
\phi(x) = a_{0}(x - \alpha_{1})(x - \alpha_{2}) \dots (x - \alpha_{n}),
\]
ここで $\alpha$ の値は実数または複素数である。すると
\[
\phi'(x) = a_{0}\sum (x - \alpha_{2})(x - \alpha_{3}) \dots (x - \alpha_{n}),
\]
この表記法は、$n-1$ 個の因子からすべての可能な積を形成し、それらを全て加えることを意味する。この結果の形式は、$\alpha$ のいくつかが等しい場合でも有効である；ただし当然、右辺の一部の項は重複する。読者は容易に確認できるだろう。もし
\[
\phi(x) = a_{0}(x - \alpha_{1})^{m_{1}}
               (x - \alpha_{2})^{m_{2}}\dots
               (x - \alpha_{\nu})^{m_{\nu}},
\]
であるとき、
\[
\phi'(x) = a_{0} \sum m_{1}(x - \alpha_{1})^{m_{1}-1}
    (x - \alpha_{2})^{m_{2}}\dots
    (x - \alpha_{\nu})^{m_{\nu}}.
\]

例題\par\medskip\noindent\textbf{【例 XLI.】}\par
\Item{1.} $\phi(x)$ が多項式であるとき、$\phi'(x)$ は $h$ のべき級数展開における $\phi(x + h)$ の係数であることを示せ。
\par\medskip

\item2. $\phi(x)$が$(x - \alpha)^{2}$で割り切れるとき、$\phi'(x)$は$x - \alpha$で割り切れる：一般に、$\phi(x)$が$(x - \alpha)^{m}$で割り切れるとき、$\phi'(x)$は$(x - \alpha)^{m-1}$で割り切れる。

\item3. 逆に、$\phi(x)$と$\phi'(x)$がともに$x - \alpha$で割り切れるとき、$\phi(x)$は$(x - \alpha)^{2}$で割り切れる；また、$\phi(x)$が$x - \alpha$で割り切れて$\phi'(x)$が$(x - \alpha)^{m-1}$で割り切れるとき、$\phi(x)$は$(x - \alpha)^{m}$で割り切れる。

\item4. $P(x) = 0$の多重根とその重複度をできるだけ完全に決定する方法を、基本的な代数的操作を使って示せ。

[もし$H_{1}$が$P$と$P'$の最大公約数であり、$H_{2}$が$H_{1}$と$P''$の最大公約数であり、$H_{3}$が$H_{2}$と$P'''$の最大公約数であるなどとするとき、$H_{1}H_{3}/H_{2}^{2} = 0$の根は$P = 0$の重複度2の根であり、$H_{2}H_{4}/H_{3}^{2} = 0$の根は重複度3の根であるなどとわかる。しかし、$H_{1}H_{3}/H_{2}^{2} = 0$, $H_{2}H_{4}/H_{3}^{2} = 0$などの解を完結させることが不可能な場合もある。たとえば、$P(x) = (x - 1)^{3}(x^{5} - x - 7)^{2}$の場合、$H_{1}H_{3}/H_{2}^{2} = x^{5} - x - 7$であり、$H_{2}H_{4}/H_{3}^{2} = x - 1$となり、最初の方程式を解くことができない.]

\item5. 次の式のすべての根とその重複度を求めよ。
\[
x^{4} + 3x^{3} - 3x^{2} - 11x - 6 = 0,\quad
x^{6} + 2x^{5} - 8x^{4} - 14x^{3} + 11x^{2} + 28x + 12 = 0.
\]

\item6. $ax^{2} + 2bx + c$が重複根を持つとき、すなわち$a(x - \alpha)^{2}$の形を取るとき、$2(ax + b)$は$x - \alpha$で割り切れるので、$\alpha = -b/a$となる。この値が満たすべき条件は$ax^{2} + 2bx + c = 0$であることを確認し、その結果得られる条件は$ac - b^{2} = 0$である。

\item7. 方程式 $1/(x - a) + 1/(x - b) + 1/(x - c) = 0$が等しい根を持つのは$a = b = c$のときのみである。 \MathTrip{1905.}

8. \(ax^{3} + 3bx^{2} + 3cx + d = 0\)が重根を持つことは、\(G^{2} + 4H^{3} = 0\)のときである。ここで、\(H = ac - b^{2}\), \(G = a^{2}d - 3abc + 2b^{3}\)。

\[ax + b = y\]と置くと、方程式は\(y^{3} + 3Hy + G = 0\)となる。この式は\(y^{2} + H = 0\)と共通の根を持つべきである。

9. \(\alpha\), \(\beta\), \(\gamma\), \(\delta\)が
\[ax^{4} + 4bx^{3} + 6cx^{2} + 4dx + e = 0\]
の根であるとき、その根が次の式で表される方程式は
\[\frac{1}{12}a \{
  (\alpha - \beta)(\gamma - \delta) - (\gamma - \alpha)(\beta - \delta)
\}\]
および二つの類似な表現（\(\alpha\), \(\beta\), \(\gamma\)を循環的に置き換えたもの）は
\[4\theta^{3} - g_{2}\theta - g_{3} = 0\]
である。ここで、
\[g_{2} = ae - 4bd + 3c^{2},\quad
g_{3} = ace + 2bcd - ad^{2} - eb^{2} - c^{3}.\]

もし\(\alpha\), \(\beta\), \(\gamma\), \(\delta\)のうち二つが等しいなら、この立方方程式の二つの根は等しくなる。Ex.8の結果を用いて、\(g_{2}^{3} - 27g_{3}^{2} = 0\)を得る。

10. ロールの定理（多項式の場合）：\(\phi(x)\)が任意の多項式であるとき、\(\phi(x) = 0\)の任意の二つの根の間に\(\phi'(x) = 0\)の根がある。

以下は定理の一般証明について述べています。これは多項式だけでなく、他の関数クラスにも適用されます。後で詳しく説明します。以下の証明は多項式に対してのみ有効です。$\alpha$，$\beta$ が連続する根であり、それぞれ $m$ 回と $n$ 回繰り返されていると仮定すると、

\[
\phi(x) = (x - \alpha)^{m} (x - \beta)^{n} \theta(x),
\]

ここで $\theta(x)$ は $\alpha \leq x \leq \beta$ の範囲で正の符号を持つ多項式です。すると、

{\footnotesize
\begin{align*}
\phi'(x)
  &= (x - \alpha)^{m} (x - \beta)^{n} \theta'(x)
   + \{m(x - \alpha)^{m-1} (x - \beta)^{n}
   + n(x - \alpha)^{m} (x - \beta)^{n-1}\} \theta(x)\\
  &= (x - \alpha)^{m-1} (x - \beta)^{n-1}
     [(x - \alpha) (x - \beta) \theta'(x)
      + \{m(x - \beta) + n(x - \alpha)\} \theta(x)]\\
  &= (x - \alpha)^{m-1} (x - \beta)^{n-1} F(x),
\end{align*}
}

ここで $F(x)$ とします。$F(\alpha) = m(\alpha - \beta) \theta(\alpha)$，$F(\beta) = n(\beta - \alpha) \theta(\beta)$ であり、これらは符号が異なります。したがって $F(x)$，すなわち $\phi'(x)$ は $\alpha$ と $\beta$ の間の某値で零となる。

\Paragraph{117.} \Topic{\Item{B.} 機械的関数。} 
\[
R(x) = \frac{P(x)}{Q(x)},
\]
ここで $P$ と $Q$ は多項式であるとき、\SecNo[§]{113}, (5) より、

\[
R'(x) = \frac{P'(x)Q(x) - P(x)Q'(x)}{\{Q(x)\}^{2}}
\]

を得ます。この公式は任意の機械的関数の導関数を書き下すのに使用できます。しかし、得られる形が最も単純な形であるとは限らず、$Q(x)$ と $Q'(x)$ に共通因子がない場合、つまり $Q(x)$ が重根を持たない場合、その形は最も単純です。しかし、$Q(x)$ が重根を持つ場合、得られる $R'(x)$ の表現はさらに簡略化できる可能性があります。

有理関数を微分する際に、部分分数分解を利用する場合がしばしば便利です。$Q(x)$ は第116節で述べたように、次のように表されます：
\[
a_{0}(x - \alpha_{1})^{m_{1}}(x - \alpha_{2})^{m_{2}}\dots(x - \alpha_{\nu})^{m_{\nu}}.
\]
代数学の教科書では証明されていますが、$R(x)$ は次のように表すことができます：
\begin{align*}
\Pi(x) &+ \frac{A_{1, 1}}{x - \alpha_{1}}
        + \frac{A_{1, 2}}{(x - \alpha_{1})^{2}} + \dots
        + \frac{A_{1, m_{1}}}{(x - \alpha_{1})^{m_{1}}}\\
       &+ \frac{A_{2, 1}}{x - \alpha_{2}}
        + \frac{A_{2, 2}}{(x - \alpha_{2})^{2}} + \dots
        + \frac{A_{2, m_{2}}}{(x - \alpha_{2})^{m_{2}}} + \dots,
\end{align*}
ここで $\Pi(x)$ は多項式であり、つまり多項式と
\[
\frac{A}{(x - \alpha)^{p}}
\]
の形の項の和からなります。$Q(x) = 0$ の根である $\alpha$ が与えられています。既に多項式的な部分の微分法を知っていますので、第113節の定理（4）より、または $\alpha$ が複素数の場合には第114節で示したその拡張から、最後に述べた有理関数の微分は次のようになります：
\[
-\frac{pA(x - \alpha)^{p-1}}{(x - \alpha)^{2p}}
  = -\frac{pA}{(x - \alpha)^{p+1}}.
\]

現在、一般の有理関数 $R(x)$ の微分を以下のように書き出すことができます：
\[
\Pi'(x) - \frac{A_{1, 1}}{(x - \alpha_{1})^{2}}
        - \frac{2A_{1, 2}}{(x - \alpha_{1})^{3}} - \dots
        - \frac{A_{2, 1}}{(x - \alpha_{2})^{2}}
        - \frac{2A_{2, 2}}{(x - \alpha_{2})^{3}} - \dots.
\]
さらに、$x^m$ の微分はすべての整数 $m$（正または負）に対して $mx^{m-1}$ であることを証明しました。

この節で説明した方法は、有理関数を複数回微分する際に特に有用です（第45題参照）。

\par\medskip\noindent\textbf{【例 XLII.】}\par
\Item{1.} $x$ に関する微分を計算すると次のようになります：
\[
\frac{d}{dx}\left(\frac{x}{1 + x^{2}}\right)
  = \frac{1 - x^{2}}{(1 + x^{2})^{2}},\quad
\frac{d}{dx}\left(\frac{1 - x^{2}}{1 + x^{2}}\right)
  = -\frac{4x}{(1 + x^{2})^{2}}.
\]

\Item{2.} 次の式の微分を計算すると次のようになります：
\[
\frac{d}{dx}\left(\frac{ax^{2} + 2bx + c}{Ax^{2} + 2Bx + C}\right)
  = \frac{(ax + b) (Bx + C) - (bx + c) (Ax + B)}{(Ax^{2} + 2Bx + C)^{2}}.
\]

\Item{3.} $Q$ が $(x - \alpha)^{m}$ の因数を持つ場合、$R'$（$R'$ を既約形に簡約したとき）の分母は $(x - \alpha)^{m+1}$ で割り切れますが、より高い次数の $(x - \alpha)$ のべきでは割り切れない。

\Item{4.} $R'$ の分母が単純な因数 $(x - \alpha)$ を持つことはない。したがって、分母に単純な因数を含む有理関数（例えば $1/x$）は別の有理関数の導関数ではない。
\par\medskip

\Paragraph{118.} \Topic{\Item{C.} 代数的関数。} 前節の結果と、§113 の定理 (6) を合わせて用いることで、どのような明示的な代数的関数の導関数も求めることができます。

最も重要なような関数は $x^{m}$ で、$m$ が有理数であるときです。既に §117 で述べたように、$m$ が正または負の整数の場合、この関数の導関数は $mx^{m-1}$ となります。そして、現在はすべての有理数 $m$ の値に対してこの結果が成り立つことを証明します。$y = x^{m} = x^{p/q}$ とし、ここで $p$ と $q$ は整数で $q$ は正であるとする。また $z = x^{1/q}$ とすると、$x = z^{q}$, $y = z^{p}$ となります。このとき
\[
\frac{dy}{dx}
  = \left(\frac{dy}{dz}\right) / \left(\frac{dx}{dz}\right)
  = \frac{p}{q} z^{p-q} = mx^{m-1}.
\]

This result may also be deduced as a corollary from §xxxvi.3.
もし $\phi(x) = x^{m}$ のとき、
\begin{align*}
\phi'(x)
  &= \lim_{h \to 0} \frac{(x + h)^{m} - x^{m}}{h}\\
  &= \lim_{\xi \to x} \frac{\xi^{m} - x^{m}}{\xi - x}
   = mx^{m-1}.
\end{align*}
より、より一般的な公式
\[
\frac{d}{dx} (ax + b)^{m} = ma(ax + b)^{m-1}
\]
はすべての有理数 $m$ に対して成り立つことが明らかである。

暗黙的な代数関数の微分には、理論上の困難があるが、これについては第VII章で再び扱う。
しかし、このような関数の導関数を実際に計算する際の実用的な困難は存在しない：その方法は例によって十分に説明される。
例えば $y$ が方程式
\[
x^{3} + y^{3} - 3axy = 0
\]
で与えられているとする。$x$ について微分すると、
\[
x^{2} + y^{2} \frac{dy}{dx} - a\left(y + x \frac{dy}{dx}\right) = 0
\]
となり、従って
\[
\frac{dy}{dx} = -\frac{x^{2} - ay}{y^{2} - ax}.
\]

例題 XLIII.
1. 次の関数の導関数を求めよ：
\[
\sqrt{\frac{1 + x}{1 - x}}, \quad
\sqrt{\frac{ax + b}{cx + d}}, \quad
\sqrt{\frac{ax^{2} + 2bx + c}{Ax^{2} + 2Bx + C}}, \quad
(ax + b)^{m} (cx + d)^{n}.
\]

2. 次を証明せよ：
\[
\frac{d}{dx}\left\{\frac{x}{\sqrt{a^{2} + x^{2}}}\right\}
  = \frac{a^{2}}{(a^{2} + x^{2})^{(3/2)}}, \quad
\frac{d}{dx}\left\{\frac{x}{\sqrt{a^{2} - x^{2}}}\right\}
  = \frac{a^{2}}{(a^{2} - x^{2})^{3/2}}.
\]

3. 次の $y$ の微分係数を求めよ：
\[
(i) ax^{2} + 2hxy + by^{2} + 2gx + 2fy + c = 0, \quad
(ii) x^{5} + y^{5} - 5ax^{2}y^{2} = 0.
\]

119. 超越関数．D．
既に（§xxxix.4）証明済みである：
\[
D_{x} \sin x = \cos x, \quad
D_{x} \cos x = -\sin x.
\]

由第113節の定理（4）と（5）により，読者は容易に確認できるのは，
\begin{alignat*}{2}
D\\textit{{x} \tan x \&= \sec^{2} x,          \& D}{x} \cot x \&= -\cosec^{2} x,\\
D\\textit{{x} \sec x \&= \tan x \sec x, \quad \& D}{x} \cosec x \&= -\cot x\cosec x.
\end{alignat*}
また，定理（7）を用いて通常の逆三角関数の導関数を求めることができる。読者は以下の公式を確認すべきである：
\begin{alignat*}{2}
D\\textit{{x} \arcsin x \&= ±1/\sqrtp{1 - x^{2}}, \&
D}{x} \arccos x \&= \mp 1/\sqrtp{1 - x^{2}},\\
\%
D\\textit{{x} \arctan x \&= 1/(1 + x^{2}), \&
D}{x} \arccot x \&= -1/(1 + x^{2}),\\
D\\textit{{x} \arcsec x \&= ± 1/\{x\sqrtp{x^{2} - 1}\}, \quad \&
D}{x} \arccosec x \&= \mp 1/\{x\sqrtp{x^{2} - 1}\}.
\end{alignat*}
逆正弦および逆余割の場合の不確定な符号は，$\cos(\arcsin x)$ の符号と一致し，逆余弦および逆正割の場合の不確定な符号は，$\sin(\arccos x)$ の符号と一致する。

より一般的な公式
\[
D_{x} \arcsin(x/a) = ±1/\sqrtp{a^{2} - x^{2}},\quad
D_{x} \arctan(x/a) = a/(x^{2} + a^{2}),
\]
も同様に第113節の定理（7）から容易に導くことができ，これらは重要な公式である。これらの公式の最初の式における不確定な符号は，$a\cos\{\arcsin(x/a)\}$ の符号と一致するため，
\[
a\sqrtb{1 - (x^{2}/a^{2})} = ±\sqrtp{a^{2} - x^{2}}
\]
が$a$が正または負である場合に応じて成立する。

最後に，第113節の定理（6）を用いて代数的および三角関数的な記号を含む合成関数の微分ができるようになり，任意のこのような関数の導関数を以下のような例で書き出すことができる。

\par\medskip\noindent\textbf{【例 XLIV.】}\par
\Item{1.} 次の関数の導関数を求めよ．
（これらの例において$m$は有理数であり，$a$, $b$, \dots, $\alpha$, $\beta$などの値は，それらを含む関数が実数であるようなものとする．）
\par\medskip

\begin{gather*}
\cos^{m} x, \quad \sin^{m} x, \quad
\cos x^{m}, \quad \sin x^{m}, \quad
\cos (\sin x), \quad \sin (\cos x),\\
\sqrt{a^{2}\cos^{2} x + b^{2}\sin^{2} x}, \quad
\frac{\cos x\sin x}{\sqrt{a^{2}\cos^{2} x + b^{2}\sin^{2} x}},\\
x\arcsin x + \sqrt{1 - x^{2}}, \quad
(1 + x)\arctan\sqrt{x} - \sqrt{x}.
\end{gather*}
\pageref{PageSep}{213}

\par\noindent\textbf{・} [2.] $0$ から $1$ までのすべての $x$ の値に対して $\arcsin x + \arccos x$ が定数であることを微分法で確認せよ。また、すべての正の $x$ の値に対して $\arctan x + \arccot x$ が定数であることを確認せよ。

\par\noindent\textbf{・} [3.] 次の関数の導関数を求めよ：
\[
\arcsin\sqrt{1 - x^{2}},\quad
\arcsin\{2x\sqrt{1 - x^{2}}\},\quad
\arctan \left(\frac{a + x}{1 - ax}\right).
\]
これらの結果の単純さをどのように解釈するか。

\par\noindent\textbf{・} [4.] 次の関数の導関数を求めよ：
\[
 \frac{1}{\sqrt{ac - b^{2}}} \arctan \frac{ax + b}{\sqrt{ac - b^{2}}},\quad
-\frac{1}{\sqrt{-a}} \arcsin\frac{ax + b}{\sqrt{b^{2} - ac}}.
\]

\par\noindent\textbf{・} [5.] 次の関数の導関数がそれぞれ
\[
\frac{1}{\sqrt{(\alpha - x)(x - \beta)}}
\]
であることを示せ：
\[
2\arcsin \sqrt{\frac{x - \beta}{\alpha - \beta}},\quad
2\arctan \sqrt{\frac{x - \beta}{\alpha - x}},\quad
\arcsin \frac{2\sqrt{(\alpha - x)(x - \beta)}}{\alpha - \beta}.
\]

\par\noindent\textbf{・} [6.] 次の関数の導関数を求めよ：
\[
\frac{d}{d\theta}\left\{
  \arccos \sqrt{\frac{\cos 3\theta}{\cos^{3}\theta}}
\right\}
  = \sqrt{\frac{3}{\cos\theta \cos 3\theta}}.
\]
\mathtrip{1904}.

\par\noindent\textbf{・} [7.] 次の関数の導関数を求めよ：
\[
\frac{1}{\sqrt{C(Ac - aC)}}\, \frac{d}{dx} \left[
  \arccos \sqrt{\frac{C(ax^{2} + c)}{c(Ax^{2} + C)}}
\right]
  = \frac{1}{(Ax^{2} + C) \sqrt{ax^{2} + c}}.
\]

\par\noindent\textbf{・} [8.] 次の関数の導関数がそれぞれ $1/(a + b\cos x)$ であることを示せ：
\[
\frac{1}{\sqrt{a^{2} - b^{2}}}
  \arccos \left(\frac{a\cos x + b}{a + b\cos x}\right),\quad
\frac{2}{\sqrt{a^{2} - b^{2}}}
  \arctan \left\{\sqrt{\frac{a - b}{a + b }} \tan \tfrac{1}{2}x\right\}.
\]

\item9. \(X = a + b \cos x + c \sin x\) とし、
\[
y = \frac{1}{\sqrt{a^{2} - b^{2} - c^{2}}}
  \arccos \frac{aX - a^{2} + b^{2} + c^{2}}{X \sqrt{b^{2} + c^{2}}},
\]
とすると、\( \frac{dy}{dx} = \frac{1}{X} \).

\item10. \(F[f\{\phi(x)\}\] の導関数は \(F'[f\{\phi(x)\}] f'\{\phi(x)\} \phi'(x)\) であり、さらに複雑なケースにも適用できる。

\item11. \(u\) と \(v\) を \(x\) の関数とするとき、
\[
D_{x} \arctan \left( \frac{u}{v} \right) = \frac{v D_{x} u - u D_{x} v}{u^{2} + v^{2}}.
\]

\item12. \(y = (\tan x + \sec x)^{m}\) の導関数は \(my \sec x\) である。

\item13. \(y = \cos x + i \sin x\) の導関数は \(iy\) である。

\item14. \(x \cos x\)、\(\frac{\sin x}{x}\) を微分せよ。これらの曲線 \(y = x \cos x\)、\(y = \frac{\sin x}{x}\) 上で \(x\) 軸に平行な接線を持つ \(x\) の値はそれぞれ \(\cot x = x\)、\(\tan x = x\) の根である。

\item15. \(a > 0\) における方程式 \(\sin x = ax\) が実数の解を持つかどうかについては、\(a \geq 1\) のとき \(x = 0\) 外には解がないことが容易に確認できる。また \(a < 1\) のとき、解は \(a\) が減少するにつれて有限個増加する。この方程式の根の数が変化する \(a\) の値は \(\cos \xi\) である。ここで \(\xi\) は方程式 \(\tan \xi = \xi\) の正の解である。[必要な値は \(y = ax\) と \(y = \sin x\) が接することを示す \(a\) の値である]。

\item16. \(x \neq 0\) のとき \(\phi(x) = x^{2} \sin(1/x)\)、\(x = 0\) のとき \(\phi(0) = 0\) とすると、
\[
\phi'(x) = 2x \sin(1/x) - \cos(1/x)
\]
が \(x \neq 0\) のとき成り立つ。また \(\phi'(0) = 0\) であり、\(\phi'(x)\) は \(x = 0\) で不連続である。（§111, (2) 参照）

\item17. 半径 \(a\) の円 \(x^{2} + y^{2} = a^{2}\) 上の点 \((x_{0}, y_{0})\) における接線と法線の方程式を求めよ。

[ここで \(y = \sqrt{a^{2} - x^{2}}\)、\(\frac{dy}{dx} = -\frac{x}{\sqrt{a^{2} - x^{2}}}\) であり、接線は
\[
y - y_{0} = (x - x_{0}) \left\{-\frac{x_{0}}{\sqrt{a^{2} - x_{0}^{2}}}\right\}
\]
と表される。これを整理すると \(xx_{0} + yy_{0} = a^{2}\) となる。法線は \(xy_{0} - yx_{0} = 0\) であり、これは原点を通る]。

\#\#\# 18. エッリプス $(x/a)^2 + (y/b)^2 = 1$ および双曲線 $(x/a)^2 - (y/b)^2 = 1$ の任意の点における接線と法線の方程式を求めよ。

\#\#\# 19. 参数 $t$ に対応する点における曲線 $x = \phi(t)$, $y = \psi(t)$ の接線と法線の方程式は
\[
\frac{x - \phi(t)}{\phi'(t)} = \frac{y - \psi(t)}{\psi'(t)},\quad
\{x - \phi(t)\} \phi'(t) + \{y - \psi(t)\} \psi'(t) = 0.
\]
である。

\#\#\# 120. 微分の重複。 $\phi(x)$ から $\phi'(x)$ を形成したように、$\phi'(x)$ から新しい関数 $\phi''(x)$ を形成することができる。この関数は $\phi(x)$ の二階微分または二階導数と呼ばれる。$y = \phi(x)$ の二階導数は以下のいずれかの形で表すことができる。
\[
D_x^2 y,\quad
\left(\frac{d}{dx}\right)^2 y,\quad
\frac{d^2 y}{dx^2}.
\]

同様に、$y = \phi(x)$ の $n$ 階微分または $n$ 階導数を定義することができる。これは以下のいずれかの形で表すことができる。
\[
\phi^{(n)}(x),\quad
D_x^n y,\quad
\left(\frac{d}{dx}\right)^n y,\quad
\frac{d^n y}{dx^n}.
\]
しかし、与えられた関数の $n$ 階導数を一般化した公式を容易に書き出すのは少数のケースのみである。これらのいくつかの例は以下の問題群で見ることができる。

\#\#\# XLV. の例題

1. $\phi(x) = x^m$ のとき
\[
\phi^{(n)}(x) = m(m - 1) \dots (m - n + 1)x^{m-n}.
\]
この結果により、任意の多項式の $n$ 階導数を書き出すことができる。

2. $\phi(x) = (ax + b)^m$ のとき
\[
\phi^{(n)}(x) = m(m - 1) \dots (m - n + 1)a^n(ax + b)^{m-n}.
\]
これらの例では、$m$ は任意の有理数である。もし $m$ が正の整数であり、$n > m$ のとき $\phi^{(n)}(x) = 0$.

3. 下記の公式
\[
\left(\frac{d}{dx}\right)^{n} \frac{A}{(x - \alpha)^{p}}
  = (-1)^{n} \frac{p(p + 1) \dots (p + n - 1)A}{(x - \alpha)^{p+n}}
\]
を用いることで、任意の標準形で部分分数に表された有理関数の第 \(n\) 次導関数を求めることができます。

4. \(1/(1 - x^{2})\) の第 \(n\) 次導関数は
\[
\tfrac{1}{2}(n!) \{(1 - x)^{-n-1} + (-1)^{n}(1 + x)^{-n-1}\}
\]
であることを証明します。

5. リーブニッツの定理について． \(y\) が \(uv\) の積であり，\(u\) および \(v\) の最初の \(n\) 次導関数を形成できるとき，リーブニッツの定理により \(y\) の第 \(n\) 次導関数を求めることができます．この定理は次の規則を与えます：
\[
(uv)_{n} = u_{n}v
  + \binom{n}{1}u_{n-1}v_{1}
  + \binom{n}{2}u_{n-2}v_{2} + \dots
  + \binom{n}{r}u_{n-r}v_{r} + \dots + uv_{n},
\]
ここで，下付きの文字は微分を表し，例えば \(u_{n}\) は \(u\) の第 \(n\) 次導関数を示します．定理を証明するために，まず
\begin{align*}
(uv)_{1} &= u_{1}v + uv_{1},\\
(uv)_{2} &= u_{2}v + 2u_{1}v_{1} + uv_{2},
\end{align*}
と示されます．この過程を繰り返すことで，下記の形の公式が得られます：
\[
(uv)_{n} = u_{n}v
  + a_{n, 1} u_{n-1} v_{1}
  + a_{n, 2} u_{n-2} v_{2} + \dots
  + a_{n, r} u_{n-r} v_{r} + \dots + uv_{n}.
\]

\(a_{n, r} = \binom{n}{r}\) が \(r = 1, 2, \dots, n - 1\) に対して成り立つと仮定し，この仮定が成り立つとき \(a_{n+1, r} = \binom{n + 1}{r}\) が \(r = 1, 2, \dots, n\) に対して成り立つことを示します．これにより数学的帰納法の原理から，任意の \(n\) および \(r\) に対して \(a_{n, r} = \binom{n}{r}\) が成り立つことが分かります．

\((uv)_{n+1}\) を \((uv)_{n}\) の微分により形成するとき，\(u_{n+1-r}v_{r}\) の係数は
\[
a_{n, r} + a_{n, r-1} = \binom{n}{r} + \binom{n}{r - 1} = \binom{n + 1}{r}
\]
となります．これにより定理が証明されます．

6. \(x^{m}f(x)\) の \(n\) 番目の導函数は次の式で表されます：
\[
\frac{m!}{(m - n)!} x^{m-n} f(x) + n \frac{m!}{(m - n + 1)!} x^{m-n+1} f'(x) + \frac{n(n - 1)}{1·2}\, \frac{m!}{(m - n + 2)!} x^{m-n+2} f''(x) + \dots,
\]
この級数は \(n + 1\) 項まで続けられ、または終了するまで続けることができます。

7. 次を証明せよ：\(D_{x}^{n}\cos x = \cos(x + \frac{1}{2}n\pi)\), \(D_{x}^{n}\sin x = \sin(x + \frac{1}{2}n\pi)\)。

8. \(y = A\cos mx + B\sin mx\) のとき、\(D_{x}^{2} y + m^{2} y = 0\) である。また、
\[
y = A\cos mx + B\sin mx + P_{n}(x),
\]
ここで \(P_{n}(x)\) は次数 \(n\) の多項式のとき、\(D_{x}^{n+3} y + m^{2} D_{x}^{n+1} y = 0\) である。

9. \(x^{2} D_{x}^{2}y + x D_{x} y + y = 0\) のとき、
\[
x^{2} D_{x}^{n+2} y + (2n + 1)x D_{x}^{n+1} y + (n^{2} + 1) D_{x}^{n} y = 0.
\]

[Leibniz の定理により \(n\) 回微分する。]

10. \((Lx + M)/(x^{2} - 2Bx + C)\) の \(n\) 番目の導函数を \(U_{n}\) とすると、
\[
\frac{x^{2} - 2Bx + C}{(n + 1)(n + 2)} U_{n+2}
  + \frac{2(x - B)}{n + 1} U_{n+1} + U_{n} = 0.
\]
[まず \(n = 0\) の場合の式を得てから、Leibniz の定理により \(n\) 回微分する。]

\Item{11.} \(a/(a^{2} + x^{2})\) および \(x/(a^{2} + x^{2})\) の \(n\) 番目の導関数について。以下のように表すことができる：
\[
\frac{a}{a^{2} + x^{2}}
  = \frac{1}{2i} \left(\frac{1}{x - ai} - \frac{1}{x + ai}\right), \quad
\frac{x}{a^{2} + x^{2}}
  = \frac{1}{2} \left(\frac{1}{x - ai} + \frac{1}{x + ai}\right),
\]
したがって，
\[
D_{x}^{n} \left(\frac{a}{a^{2} + x^{2}}\right)
  = \frac{(-1)^{n} n!}{2i} \left\{
      \frac{1}{(x - ai)^{n+1}} - \frac{1}{(x + ai)^{n+1}}
\right\},
\]
同様の公式が \(D_{x}^{n}\{x/(a^{2} + x^{2})\}\) にも適用される。ここで，\(\rho = \sqrt{x^{2} + a^{2}}\)，そして \(\theta\) は \(\cos \theta = x/\rho\)，\(\sin \theta = a/\rho\) の数値的に最小の角度とするとき，\(x + ai = \rho \Cis \theta\) および \(x - ai = \rho \Cis (-\theta)\) であるから
\begin{align*}
D_{x}^{n} \left(\frac{a}{a^{2} + x^{2}}\right)
  &= \left\{\frac{(-1)^{n} n!}{2i}\right\} \rho^{-n-1}
       [\Cis \{(n + 1)\theta\} - \Cis \{-(n + 1)\theta\}]\\
  &= (-1)^{n} n!\, (x^{2} + a^{2})^{-(n+1)/2} \sin \{(n + 1) \arctan(a/x)\}.
\end{align*}
同様に，
\[
D_{x}^{n} \left(\frac{x}{a^{2} + x^{2}}\right)
  = (-1)^{n} n!\, (x^{2} + a^{2})^{-(n+1)/2} \cos \{(n + 1) \arctan(a/x)\}.
\]

\Item{12.} 次の公式を証明せよ：
\begin{align*}
D_{x}^{n} \left(\frac{\cos x}{x}\right)
  &= \left\{P_{n} \cos(x + \tfrac{1}{2}n\pi) + Q_{n} \sin(x + \tfrac{1}{2}n\pi)\right\}/x^{n+1},\\
D_{x}^{n} \left(\frac{\sin x}{x}\right)
  &= \left\{P_{n} \sin(x + \tfrac{1}{2}n\pi) - Q_{n} \cos(x + \tfrac{1}{2}n\pi)\right\}/x^{n+1},
\end{align*}
ここで \(P_{n}\) および \(Q_{n}\) はそれぞれ次数 \(n\) と \(n-1\) の多項式である。

\Item{13.} 次の公式を確立せよ：
\begin{gather*}
\frac{dx}{dy} = \left(\frac{dy}{dx}\right)^{-1},\quad
\frac{d^{2} x}{dy^{2}}
  = -\frac{d^{2} y}{dx^{2}} \left(\frac{dy}{dx}\right)^{-3},\\
\frac{d^{3} x}{dy^{3}}
  = -\left\{\frac{d^{3} y}{dx^{3}} \frac{dy}{dx}
    - 3\left(\frac{d^{2} y}{dx^{2}}\right)\right\} \left(\frac{dy}{dx}\right)^{-5}.
\end{gather*}

\\PageSep{217}

\par\noindent\textbf{・} {14.} \( yz = 1 \) かつ \( y_{r} = (1/r!) D_{x}^{r}y \), \( z_{s} = (1/s!) D_{x}^{s}z \) のとき、
\[
\frac{1}{z^{3}}
\begin{vmatrix}
z    & z_{1}& z_{2}\\
z_{1}& z_{2}& z_{3}\\
z_{2}& z_{3}& z_{4}
\end{vmatrix}
= \frac{1}{y^{2}}
\begin{vmatrix}
y_{2}& y_{3}\\
y_{3}& y_{4}
\end{vmatrix}.
\]
\MathTrip{1905.}

\par\noindent\textbf{・} {15.} \( W(y, z, u) =
\begin{vmatrix}
y  & z  & u\\
y' & z' & u'\\
y''& z''& u''
\end{vmatrix},
\) 括弧内の破線は \( x \) に関する微分を表すとき、
\[
W(y, z, u) = y^{3}\, W\left(1, \frac{z}{y}, \frac{u}{y}\right).
\]

\par\noindent\textbf{・} {16.} \( ax^{2} + 2hxy + by^{2} + 2gx + 2fy + c = 0 \) のとき、
\[
\frac{dy}{dx} = -\frac{ax + hy + g}{hx + by + f}
\]
かつ
\[
\frac{d^{2}y}{dx^{2}} = \frac{abc + 2fgh - af^{2} - bg^{2} - ch^{2}}{(hx + by + f)^{3}}.
\]

\par\medskip\noindent\textbf{【例】}\par

\Paragraph{121. 微分関数に関するいくつかの一般定理}
以降では、\(\phi(x)\) を \(x\) の関数とし、任意の \(x\) の値に対して \(\phi'(x)\) が存在することを仮定する。この仮定は当然、\(\phi(x)\) の連続性を含意する。

\ParTheorem{微分符号 \(\phi'(x)\) の意味。\normalfont\textsc{定理 A}}
\(\phi'(x_{0}) > 0\) のとき、\(x < x_{0}\) かつ十分に \(x_{0}\) に近い値に対して \(\phi(x) < \phi(x_{0})\) であり、また \(x > x_{0}\) かつ十分に \(x_{0}\) に近い値に対して \(\phi(x) > \phi(x_{0})\)。

\(\{\phi(x_{0} + h) - \phi(x_{0})\}/h\) は \(h \to 0\) のとき正の極限 \(\phi'(x_{0})\) に収束する。これは \(h\) が十分小さい値に対して \(\phi(x_{0} + h) - \phi(x_{0})\) と \(h\) が同じ符号を持つことを意味し、これが定理の主張である。当然、幾何学的な視点から見れば、結果は直感的であり、不等式 \(\phi'(x) > 0\) は曲線 \(y = \phi(x)\) の接線が \(x\) 軸に対して正の鋭角を形成することを表している。読者は \(\phi'(x) < 0\) の場合に対応する定理も自己で述べるべきである。

\par\medskip\noindent\textbf{【定理】}\par[B]
$\phi(a) = 0$ かつ $\phi(b) = 0$ のとき、$a$ と $b$ の間には少なくとも一つの $x$ の値が存在し、そのとき $\phi'(x) = 0$ である。
\par\medskip

二つの可能性がある。第一に、$\phi(x)$ は区間 $(a, b)$ または $[a, b]$ 全てで零となる場合がある。この場合、$\phi'(x)$ も同様に区間に於いて零となる。他方、$\phi(x)$ が常に零でないときには、ある点における $\phi(x)$ の正または負の値を持つ $x$ の値が存在する。例えば、$\phi(x)$ が時として正であると仮定しよう。このとき、定理2（102節）により、区間内の他のどの点よりも $\phi(\xi)$ の値が等しくまたはそれ以上となる $\xi$ の値が存在し、かつ $\xi \neq a, b$ である。そして $\phi'(\xi) = 0$ であることが必要である。もし $\phi'(\xi)$ が正であれば、定理Aにより、$\xi$ より大きい $x$ の値に対して十分に近い点では $\phi(x)$ は $\phi(\xi)$ より大きくなるので、ある点での $\phi(x)$ の値が $\phi(\xi)$ より大きくなることが確定する。同様にして $\phi'(\xi) < 0$ であることも示される。

\par\medskip\noindent\textbf{【系】}\par
$\phi(a) = \phi(b) = k$ のとき、$a$ と $b$ の間には少なくとも一つの $x$ の値が存在し、そのとき $\phi'(x) = 0$ である。
\par\medskip

この命題は、$\phi(x) - k = \psi(x)$ と置き換えて定理Bを適用することで得られる。

\par\medskip\noindent\textbf{【系】}\par
ある区間において全ての $x$ の値に対して $\phi'(x) > 0$ であるとき、$\phi(x)$ は厳密な意味での $x$ の増加関数であり、その区間に於いて常に $\phi(x)$ が増加する。
\par\medskip

Let \( x_1 \)と\( x_2 \)を区間内の\( x \)の二つの値とする，并且満たす \( x_1 < x_2 \)。我々は \( \phi(x_1) < \phi(x_2) \) を示す必要がある。まず第一に、\(\phi(x_1)\) は等しくなることはできない；もし等しいとすると、定理Bにより、\( x_1 \)と\( x_2 \)の間に存在する \( x \) の値で \(\phi'(x) = 0\) となる。また、\(\phi(x_1)\) より大きいこともできない。なぜなら、\(\phi'(x_1)\) が正であるから、定理Aにより、\( x > x_1 \)かつ十分に \( x_1 \) に近いとき、\(\phi(x)\)は\(\phi(x_1)\)より大きい。従って、ある \( x = x_3 \) の値が存在し、\( x_1 < x_3 < x_2 \)で \(\phi(x_3) = \phi(x_1)\) となる；そして定理Bにより、\( x_1 \)と\( x_3 \)の間に存在する \( x \) の値で \(\phi'(x) = 0\) となる。

\par\medskip\noindent\textbf{【系】}\par[3.]
定理2の結論は、区間\((a, b)\)または\([a, b]\)に有限個の例外的な\( x \)の値が含まれる場合でも依然として成り立つ。ただし、これらの例外的な \( x \) の値に対して \(\phi(x)\) が連続であることが条件である。
\par\medskip

明らかに十分なのは、区間内に一つだけ例外的な \( x \) の値がある場合を考慮することである。この例外的な値は区間の端点に対応するものとし、それが \( a \) とする。もし \( a < x_1 < x_2 < b \) であれば、\( a + \EPSILON \) を選んで \( a + \EPSILON < x_1 \) となるようにし、\(\DPmod{(a + \EPSILON, b)}{[a + \EPSILON, b]}\) の全域で \(\phi'(x) > 0\) であるとすれば、定理2により \(\phi(x_1) < \phi(x_2)\) となる。残るは \(\phi(a) < \phi(x_1)\) を証明することである。今、\(\phi(x_1)\) は \( x_1 \) が \( a \) に減少するにつれて安定的に減少し、厳密な意味で、従って
\[
\phi(a) = \phi(a + 0) = \lim_{x_1 \to a+0} \phi(x_1) < \phi(x_1)
\]
である。

\par\medskip\noindent\textbf{【系】}\par[4.]
区間\((a, b)\)または\([a, b]\)全域で \(\phi'(x) > 0\) であり、かつ \(\phi(a) \geq 0\) のとき、\(\phi(x)\)はこの区間にわたって正である。
\par\medskip

\par\medskip\noindent\textbf{【注釈】}\par
読者には、これらの推論の第2項目を定理Aと非常に精密に比較していただきたい。もしも定理Aにおいてのように、$\phi'(x)$が単なる一点$x = x_0$で正であるという仮定だけであれば、十分近接した$x_1$と$x_2$について$\phi(x_1) < \phi(x_2)$を証明できる。これは、定理Aにより$\phi(x_1) < \phi(x_0)$かつ$\phi(x_2) > \phi(x_0)$が成り立つからである。しかし、これによって$x_0$を取り巻く任意の区間にわたって$\phi(x)$が単調に増加する関数であることを証明することはできない。なぜなら、$x_1$と$x_2$が$x_0$の両側にあるという仮定は結論にとって不可欠だからである。この点についてもう一度言及し、実際の例で説明することにする（§124）。
\par\medskip

\Paragraph{122. 最大値と最小値}
$\phi(x)$が$x = \xi$のときには$\phi(\xi)$を取る値が、$\xi$の直近の領域において他のどの値よりも大きいとき、すなわち$\xi - \EPSILON < x < \xi$および$\xi < x < \xi + \EPSILON$の範囲における$x$の値に対して$\phi(\xi) > \phi(x)$が成り立つような区間$\DPmod{(\xi - \EPSILON, \xi + \EPSILON)}{[\xi - \EPSILON, \xi + \EPSILON]}$を見つけることができるとき、その値$\phi(\xi)$を最大値と定義する。また同様に最小値も定義する。図において点Aは最大値に対応し、点Bは最小値に対応している（図39）。このように、$A_3$が最大値に対応し、$B_1$が最小値に対応しているという事実は、関数の値が$B_1$で$A_3$よりも大きいという事実とは矛盾しない。


\par\medskip\noindent\textbf{【定理】}\par[C]
$x=\xi$において$\phi(x)$の極値が存在するための必要条件は$\phi'(\xi)=0$である。
（関数が連続であっても導関数を持たない場合でも極値を取ることがあるが、ここでは導関数が存在することを仮定している）
\par\medskip
\\PageSep{220}

これは定理Aから即座に従います。この条件が十分条件ではないことは図の点Cを見れば明らかです。
例えば$y=x^3$の場合、$\phi'(x)=3x^2$であり、$x=0$で導関数は0になりますが、$x=0$では$x^3$の極値にはnessせず、グラフの形から明らかです（図10、p.45参照）。

しかし$\phi'(\xi)=0$かつ$x<\xi$で近接するすべての$x$について$\phi'(x)>0$であり、さらに$x>\xi$で近接するすべての$x$について$\phi'(x)<0$であるとき、$x=\xi$において$\phi(x)$の極大値が存在します。そしてこれらの不等式の符号を逆にすると極小値が存在します。この場合、節121の補則3により、$\DPmod{(\xi-\EPSILON, \xi)}{[\xi-\EPSILON, \xi]}$において$\phi(x)$は$x$と共に増加し、$\DPmod{(\xi, \xi+\EPSILON)}{[\xi, \xi+\EPSILON]}$において$x$と共に減少する区間を定めることができます。これにより明らかに$\phi(\xi)$が極大値となることが保証されます。

この結果はまた次のように述べることもできます。導関数$\phi'(x)$の符号が$x=\xi$で正から負に変化した場合、$x=\xi$は$\phi(x)$の極大点を生成します。そして符号が逆方向に変化した場合、$x=\xi$は$\phi(x)$の極小点を生成します。

\paragraph{123.} 最大値や最小値の条件を述べるもう一つの方法があります。これはしばしば有用です。$\phi(x)$ が二階微分 $\phi''(x)$ を持つと仮定しましょう：もちろん、$\phi'(x)$ の存在から $\phi''(x)$ の存在が導かれるわけではありませんし、逆もまた然りです。しかし、現在我々が遭遇する可能性のあるような場合においては、この条件は一般的に満たされます。

\par\medskip\noindent\textbf{【定理】}\par[D]
もし $\phi'(\xi) = 0$ かつ $\phi''(\xi) \neq 0$ であるなら、$\phi(x)$ は $x = \xi$ で最大値または最小値を持つ。$\phi''(\xi) < 0$ のときには最大値、$\phi''(\xi) > 0$ のときには最小値をとる。
\par\medskip

たとえば $\phi''(\xi) < 0$ である場合、定理 A より、$x < \xi$ で十分に近い点では $\phi'(x)$ は負であり、$x > \xi$ で十分に近い点では正です。したがって $x = \xi$ は最大値をとります。

\par\medskip\noindent\textbf{【注釈】}\par
\paragraph{124.} 前述の内容（最後の段落を除く）において、$\phi(x)$ が考慮される区間内のすべての $x$ の値に対して微分係数を持つという単なる仮定がありました。この条件が満たされない場合、これらの定理は成り立ちません。したがって、定理 B は関数
\[
y = 1 - \sqrt{x^{2}}
\]
において失敗します。ここで平方根は正の値を取ることとします。この関数のグラフは図 40 に示されています。ここでは $\phi(-1) = \phi(1) = 0$ です：しかし、図から明らかのように $x < 0$ のとき $\phi'(x)$ は 1 であり、$x > 0$ のときには -1 で、かつ常に 0 にはなりません。$x = 0$ では微分係数が存在せず、グラフ上の点 P に接線がありません。この場合 $x = 0$ は明らかに $\phi(x)$ の最大値をとりますが、$\phi'(0)$ は存在しないので 0 にはなりません。したがって最大値のテストは失敗します。

微分$\phi'(x)$の存在だけが仮定されたが、特に仮定していないものがある。
それは$\phi'(x)$自身が連続関数であるという仮定ではないことである。この点は微妙ではあるが非常に興味深い問題を提起している。
関数$\phi(x)$がすべての$x$に対して微分を持つにもかかわらず、その微分$\phi'(x)$が連続でないことは可能だろうか？つまり、曲線がすべての点において接線を持ちつつも、接線の方向が連続的に変化しないことは可能だろうか？
読者はこの質問の意味を考え、直感に基づいて答えを試みてみると、おそらく「ノー」と答えることになるだろう。しかし、この答えは間違っていることが容易に理解できる。

関数$\phi(x)$を$x \neq 0$の場合に方程式
\[
\phi(x) = x^{2}\sin(1/x)
\]
で定義し、$\phi(0) = 0$と仮定する。このとき、$\phi(x)$はすべての$x$に対して連続である。
$x \neq 0$の場合には
\[
\phi'(x) = 2x \sin(1/x) - \cos(1/x)
\]
であり、また
\[
\phi'(0) = \lim_{h \to 0} \frac{h^{2}\sin(1/h)}{h} = 0
\]
である。
したがって$\phi'(x)$はすべての$x$に対して存在する。しかし$x = 0$において$\phi'(x)$は不連続であり、$2x\sin(1/x)$は$x \to 0$で$0$に収束し、$\cos(1/x)$は不定形$-1$と$1$の間で振動するので、$\phi'(x)$も同じ範囲で振動する。

1. $\phi(x) = (x - a)^{m} (x - b)^{n}$ の場合や、$\phi(x) = (x - a)^{m} (x - b)^{n} (x - c)^{p}$ の場合を考える。ここで $m$、$n$、$p$ は正の整数であり、$a < b < c$。

最初の関数は $x = a$ および $x = b$ で零となる。そして
\[
\phi'(x) = (x - a)^{m-1} (x - b)^{n-1} \{(m + n)x - mb - na\}
\]
は $x = (mb + na)/(m + n)$ で零となる。この値は $a$ と $b$ の間にある。

2番目の場合では、二次方程式
\[
(m + n + p)x^{2} - \{m(b + c) + n(c + a) + p(a + b)\}x + mbc + nca + pab = 0
\]
が $a$ と $b$ の間および $b$ と $c$ の間に根を持つことを確認する必要がある。

2. 次の多項式は $x > 1$ のとき正であることを示す。
\[
2x^{3} + 3x^{2} - 12x + 7,\quad
3x^{4} + 8x^{3} - 6x^{2} - 24x + 19
\]

\par\noindent\textbf{3.}  \( x - \sin x \)は任意の\( x \)の値の区間において増加関数であり、また\( \tan x - x \)は\( x \)が\(-\frac{1}{2}\pi\)から\(\frac{1}{2}\pi\)まで増加するにつれて増加する。どのような値の\( a \)に対して \( ax - \sin x \)が\( x \)の単調に増加または減少する関数となるか？

\par\noindent\textbf{4.}  \( \tan x - x \)は\( x = \frac{1}{2}\pi \)から\( x = \frac{3}{2}\pi \)、また\( x = \frac{3}{2}\pi \)から\( x = \frac{5}{2}\pi \)など、それぞれの区間において増加することを示し、これより方程式 \( \tan x = x \) の解がこれらの各区間に一つだけ存在することを導く。

\par\noindent\textbf{5.}  問題3から推論すると\( x > 0 \)の場合、\( \sin x - x < 0 \)であり、これより\( \cos x - 1 + \frac{1}{2}x^{2} > 0 \)であり、さらにこれらより\( \sin x - x + \frac{1}{6}x^{3} > 0 \)である。一般に、
\begin{align*}
C_{2m} & = \cos x - 1 + \frac{x^{2}}{2!} - \dots - (-1)^{m} \frac{x^{2m}}{(2m)!},\\
S_{2m+1}& = \sin x - x + \frac{x^{3}}{3!} - \dots - (-1)^{m} \frac{x^{2m+1}}{(2m+1)!},
\end{align*}
および\( x > 0 \)のとき、\( C_{2m} \)と\( S_{2m+1} \)は\( m \)が奇数であるか偶数であるかにより正または負となる。

\par\noindent\textbf{6.}  \( f(x) \)と\( f''(x) \)が区間\((a, b)\)における任意の点で連続かつ符号が同じであれば、この区間に方程式\( f(x) = 0 \)または\( f'(x) = 0 \)の解が一つ以上存在することは不可能である。

\par\noindent\textbf{7.}  \( u \), \( v \)およびその導関数 \( u' \), \( v' \)がある\( x \)の値の区間において連続であり、また\( uv' - u'v \)がこの区間内の任意の点で消えないであれば、\( u = 0 \)の二つの解の間に一つだけ\( v = 0 \)の解があることを示す。逆も成り立つ。\( u = \cos x \), \( v = \sin x \)の場合を検証せよ。

[もし\( v \)が\( u = 0 \)の二つの解\(\alpha\)と\(\beta\)の間で消えないなら、関数\( u/v \)は区間\((\alpha, \beta)\)において連続であり端点で値を消失させる。したがって、\((u/v)' = (u'v - uv')/v^{2}\)は\(\alpha\)と\(\beta\)の間に消失しなければならないが、これは我々の仮定に矛盾する。]

\item8. $(x - 1)^{2}(x + 2)$, $x^{3} - 3x$, $2x^{3} - 3x^{2} - 36x + 10$, $4x^{3} - 18x^{2} + 27x - 7$, $3x^{4} - 4x^{3} + 1$, $x^{5} - 15x^{3} + 3$ の最大値と最小値（存在する場合）を決定せよ。それぞれの関数のグラフの形状を素描せよ。

[例えば最後の関数について考える。このとき $\phi'(x) = 5x^{2}(x^{2} - 9)$ であり、$x = -3$, $x = 0$, および $x = 3$ で消失する。$x = -3$ は最大値をもたらし、$x = 3$ は最小値をもたらす一方で、$x = 0$ ではどちらでもないことが容易に確認できる。]

\item9. $m$ および $n$ が任意の正整数であるとき、関数 $(x - a)^{m}(x - b)^{n}$ の最大値と最小値について議論せよ。$m$ および $n$ が奇数か偶数かにより異なる場合を考慮する。それぞれの関数のグラフを素描せよ。

\item10. 同様に、関数 $(x - a)(x - b)^{2}(x - c)^{3}$ の最大値と最小値について議論し、$a$, $b$, $c$ が異なる相対的な大きさを持つ場合に対応するグラフの形状を区別せよ。

\item11. $(ax + b)/(cx + d)$ はどのような $a$, $b$, $c$, $d$ の値でも最大値や最小値を持たないことを示せ。関数のグラフを描け。

\item12. 分母が複素数の根を持つとき、関数
\[
y = \frac{ax^{2} + 2bx + c}{Ax^{2} + 2Bx + C}
\]
の最大値と最小値について議論せよ。

[私たちは$a$と$A$を正とすることができます。導関数は次の方程式が成り立つとき0になります：
\[
(ax + b)(Bx + C) - (Ax + B)(bx + c) = 0.
\Tag{(1)}
\]
この方程式には実数の根が必要です。もし実数の根がないとすると、導関数は常に同じ符号を持ち続けます。これは不可能であり、$y$はすべての$x$の値に対して連続であり、$x \to +\infty$または$x \to -\infty$のとき$y \to a/A$となるためです。この曲線が直線$y = a/A$と1点で交わるという事実を容易に確認できますし、大きな正の$x$の値に対してはこの直線より上にあり、大きな負の$x$の値に対してはこの直線より下にあることが分かります。または逆の場合も同様です。つまり、方程式（1）の代数的に大きい根が$b/a > B/A$のとき最大値を、$b/a < B/A$のとき最小値を与えることになります.]

\Item{13.} 最大値と最小値自体は、$ax^{2} + 2bx + c - \lambda(Ax^{2} + 2Bx + C)$が完全平方であるような$\lambda$の値です。これは$y = \lambda$が曲線に接する条件を表します。

\Item{14.} 一般に、$R(x) = P(x)/Q(x)$の最大値と最小値は、$P(x) - \lambda Q(x) = 0$が等しい根を持つべき条件によって得られる$\lambda$の値の中から求められます。

\Item{15.} $Ax^{2} + 2Bx + C = 0$が実数の根を持つ場合、次のように進めることができます。私たちは
\[
y - (a/A) = (\lambda x + \mu)/\{A(Ax^{2} + 2Bx + C)\},
\]
ここで$\lambda = bA - aB$, $\mu = cA - aC$とします。さらに$\xi$を$\lambda x + \mu$、$\eta$を$(A/\lambda^{2})(Ay - a)$と置き換えると、次のような方程式を得ます：
\[
\eta = \xi/\{(\xi - p)(\xi - q)\}.
\]
ページ分割

この$(x, y)$から$(\xi, \eta)$への変換は、原点の移動、軸の平行移動、各軸に対するスケール変更、$\lambda < 0$の場合には$x$軸の方向の反転を含むものに過ぎません。したがって、$y$を$x$の関数として最小値を持つ場合、それは$\eta$を$\xi$の関数として最小値を持つことに相当し、逆もまた同様です。最大値についても同様です。

微分η关于ξ的导数為0，当且僅当
\[
(\xi - p)(\xi - q) - \xi(\xi - p) - \xi(\xi - q) = 0,
\]
或 \(\xi^{2} = pq\)。因此，如果p和q同号，則导数有两个根；如果异号，則没有根。后者情况下，η的图形如图41a所示。

当p和q均為正時，η的一般图形如图41b所示，并且容易看出 \(\xi = \sqrt{pq}\) 与える极大值，\(\xi = -\sqrt{pq}\) 与える极小值。最大值為 \(-1/(\sqrt{p} - \sqrt{q})^{2}\)，最小值為 \(-1/(\sqrt{p} + \sqrt{q})^{2}\)，后者较大。

特别地，当 \(p = q\) 時，函数形式為
\[
\eta = \frac{\xi}{(\xi - p)^{2}},
\]
其图形如图41c所示。

上述讨论在 \(\lambda = 0\) 即 \(a/A = b/B\) 的情况下失效。但在これ种情况下，
\[
y - \left(\frac{a}{A}\right) = \frac{\mu}{A(Ax^{2} + 2Bx + C)} = \frac{\mu}{A^{2}(x - x_{1})(x - x_{2})},
\]
设為上述形式，且 \(dy/dx = 0\) 与える単值 \(x = \frac{1}{2}(x_{1} + x_{2})\)。绘制图形后可以看出，该值与える极大值或极小值，取决于 \(\mu\) 的正負。图42所示的图形対应于前者情况。
\[
\begin{figure}[1.5in]{42}{p224d}
\]

{\Loosen[a一般関数 \( y = \frac{ax^2 + 2bx + c}{Ax^2 + 2Bx + C} \) の全般的な議論は、Chrystalの『代数学』第1巻、464−7ページに純粋代数的手法で述べられています。]}

\Item{16.} \( x - \alpha \)と\( x - \beta \)の積を\( x - \gamma \)で割った式がすべて実数を取るか否かを示せ。また、\(\gamma\)が\(\alpha\)と\(\beta\)の間にある場合、その値はすべての実数を取り、それ以外の場合には長さ \(4\sqrt{|\alpha - \gamma||\beta - \gamma|}\) の区間に含まれない実数を除いてすべての実数を取ります。

\Item{17.} \( 0 < c < 1 \)のとき、式
\[
y = \frac{x^2 + 2x + c}{x^2 + 4x + 3c}
\]
が任意の実数を取ることを示し、この場合の関数のグラフを描け。1910年。

\Item{18.} \( x = 1 \)と\( x = -1 \)で極値（最大値または最小値）がそれぞれ2と3となる、かつ\( x = 0 \)で値が2.5となる形の関数を決定せよ。1908年。

\Item{19.} \( a \)と\( b \)が正であるとき、式
\[
(x + a)(x + b)/(x - a)(x - b)
\]
の最大値と最小値は
\[
-\left(\frac{\sqrt{a} + \sqrt{b}}{\sqrt{a} - \sqrt{b}}\right)^2,\quad
-\left(\frac{\sqrt{a} - \sqrt{b}}{\sqrt{a} + \sqrt{b}}\right)^2
\]
である。

\Item{20.} 式
\[
\frac{(x - 1)^2}{(x + 1)^3}
\]
の最大値は \( \frac{2}{27} \) である。

\Item{21.} 次の関数について最大値と最小値を議論せよ。
\[
\frac{x(x - 1)}{x^2 + 3x + 3},\quad \frac{x^4}{(x - 1)(x - 3)^3},\quad
\frac{(x - 1)^2(3x^2 - 2x - 37)}{(x + 5)^2(3x^2 - 14x - 1)}
\]
[最後の関数を \( P(x)/Q(x) \) と表すと、
\[
P'Q - PQ' = 72(x - 7)(x - 3)(x - 1)(x + 1)(x + 2)(x + 5)
\]
である。]

\Item{22.} \( a \cos x + b \sin x \) の最大値と最小値を求めよ。結果を確認するために、関数を \( A \cos(x - a) \) の形で表す。

\Item{23.} 次の関数について最大値と最小値を求めよ。
\[
a^2 \cos^2 x + b^2 \sin^2 x,\quad A \cos^2 x + 2H \cos x \sin x + B \sin^2 x
\]

\Item{24.} \( \frac{\sin(x + a)}{\sin(x + b)} \) の最大値と最小値は存在しないことを示し、その関数のグラフを描け。

\par\noindent\textbf{・} [25.] 次の関数
\[
\frac{\sin^{2}x}{\sin(x + a)\sin(x + b)}\quad
(0 < a < b < \pi)
\]
は無数の最小値が \(0\) であり、最大値は
\[
-4\sin a\sin b/\sin^{2}(a - b)
\]
であることを示せ。

\par\noindent\textbf{・} [26.] $a^{2}\sec^{2}x + b^{2}\cosec^{2}x$ の最小値は $(a + b)^{2}$ である。

\par\noindent\textbf{・} [27.] $\tan 3x \cot 2x$ は $\frac{1}{9}$ と $\frac{3}{2}$ の間には存在しないことを示せ。

\par\noindent\textbf{・} [28.] 直角三角形の斜辺と他の一方の辺の和が与えられたとき、その二辺のなす角度が $60°$ のときに直角三角形の面積は最大であることを示せ。 \MathTrip{1909}。

\par\noindent\textbf{・} [29.] 固定点 $(a, b)$ を通る直線が $OX$, $OY$ 軸と交わる点を $P$, $Q$ とするとき、$PQ$, $OP + OQ$, および $OP·OQ$ の最小値はそれぞれ $(a^{2/3} + b^{2/3})^{3/2}$, $(\sqrt{a} + \sqrt{b})^{2}$, および $4ab$ であることを示せ。

\par\noindent\textbf{・} [30.] 極線が楕円と交わる点を $P$, $Q$ とするとき、$PQ$ の最小値は楕円の半長軸と半短軸の和に等しい。

\par\noindent\textbf{・} [31.] 二次曲線
\[
ax^{2} + 2hxy + by^{2} = 1
\]
の軸の長さと方向を求めよ。

[軸 $x$ 方向に角度 $\theta$ を作る半径 $r$ の長さは以下の式で与えられる。
\[
1/r^{2} = a\cos^{2} \theta + 2h\cos\theta \sin\theta + b\sin^{2} \theta.
\]
$r$ の最大値または最小値の条件は $\tan 2\theta = 2h/(a - b)$ である。この二つの式を用いて $\theta$ を消去すると
\[
\{a - (1/r^{2})\} \{b - (1/r^{2})\} = h^{2}
\]
となる。]

\par\noindent\textbf{・} [32.] $x + y = k$ である正の数 $x$ と $y$ の条件下で $x^{m}y^{n}$ の最大値は
\[
m^{m} n^{n} k^{m+n}/(m + n)^{m+n}
\]
である。

\par\noindent\textbf{・} [33.] $x^{2} + xy + y^{2} = 3\kappa^{2}$ である正の数 $x$ と $y$ の条件下で $ax + by$ の最大値は
\[
2\kappa \sqrt{a^{2} - ab + b^{2}}
\]

もし \(ax + by\) が最大値であるならば、\(a + b(dy/dx) = 0\) である。\(x\) と \(y\) の関係から \((2x + y) + (x + 2y)(dy/dx) = 0\) となる。両者の \(dy/dx\) を等しくする。

34. \(\theta\) および \(\phi\) が鋭角であり、関係式 \(a \sec\theta + b \sec\phi = c\) で結ばれており、\(a\)、\(b\)、\(c\) が正の数であるならば、\(\theta = \phi\) のとき \(a\cos\theta + b\cos\phi\) は最小値を取る。

125. 平均値の定理。次に、微分積分学において最も重要な定理の一つである平均値の定理の証明に進むことができる。

以下がその定理である。
もし \(\phi(x)\) が区間 \((a, b)\) におけるすべての \(x\) の値に対して導関数を持つならば、\(a\) より \(b\) までの区間に位置する \(x\) の値 \(\xi\) 存在して
\[
\phi(b) - \phi(a) = (b - a)\phi'(\xi)
\]
が成り立つ。

この定理の厳密な証明を与える前に、これは微分積分学において最も重要な定理であるため、その明白な幾何学的な意味を指摘しておくことが適当である。これは単に (図43 を参照) ある点 \(P\) が存在し、その接線が \(AB\) と平行であることの意味である。ここで \(\phi'(\xi)\) は点 \(P\) の接線が \(OX\) に対して作る角の正接であり、\(\{\phi(b) - \phi(a)\}/(b - a)\) は直線 \(AB\) が \(OX\) に対して作る角の正接である。

この定理を厳密な解析的手法で証明するのも容易である。関数
\[
\phi(b) - \phi(x) - \frac{b - x}{b - a}\{\phi(b) - \phi(a)\}
\]
を考える。これは \(x = a\) および \(x = b\) のとき零となる。セクション121の定理Bから、その導関数が零となる値 \(\xi\) が存在することが従う。この導関数は
\[
\frac{\phi(b) - \phi(a)}{b - a} - \phi'(x)
\]
であり、これにより定理が証明される。この証明において \(\phi'(x)\) の連続性を仮定していないことに注意しなければならない。

平均値の定理はしばしば次の形で表されます：
\[
\phi(b) = \phi(a) + (b - a) \phi'\{a + \theta(b - a)\},
\]
ここで $\theta$ は $0$ と $1$ の間の数です。もちろん $a + \theta(b - a)$ は `$a$ と $b$ の間にある数 $\xi$` を別の表現として表しています。

$b = a + h$ と置くと、
\[
\phi(a + h) = \phi(a) + h \phi'(a + \theta h)
\]
となり、これが定理の最も一般的な形です。

例題 XLVII.

1. 次を示せ：
\[
\phi(b) - \phi(x) - \frac{b - x}{b - a}\{\phi(b) - \phi(a)\}
\]
これは曲線上の点の縦座標と直線上の対応する点の縦座標との差に等しい。

2. $\phi(x) = x^2$ および $\phi(x) = x^3$ の場合で定理を検証せよ。

後者の場合、$(b^3 - a^3)/(b - a) = 3\xi^2$ を示す必要がある、ここで $a < \xi < b$；つまり、$\frac{1}{3}(b^2 + ab + a^2) = \xi^2$ のとき $\xi$ は $a$ と $b$ の間にあることを証明せよ。

3. 第 124 節末尾で述べた定理を平均値の定理を使って証明せよ。

$\phi'(0) = c$ より、小さな正数 $x$ を求めることができる。このとき $\{\phi(x) - \phi(0)\}/x$ は $c$ に近い値となるので、定理により $\xi$ の小さな正数が存在し、$\phi'(\xi)$ は $c$ に近い値となる。これは $\lim\limits_{x \to +0} \phi'(x) = a$ である場合に限って $a = c$ であり、同様に $b = c$ であることを示す。

4. 平均値の定理を使って第 113 節の定理 (6) を証明し、導出に用いる微分が連続であると仮定せよ。

[関数 \( F\{f(x)\} \) の導函数は定義により
\[
\lim_{h \to 0} \frac{F\{f(x + h)\} - F\{f(x)\}}{h}
\]
で与えられる。しかし、平均値の定理によれば \( f(x + h) = f(x) + hf'(ξ) \)，ここで \( ξ \) は \( x \) と \( x + h \) の間にある数である。そして
\[
F\{f(x) + hf'(ξ)\} = F\{f(x)\} + hf'(ξ)F'ξ_1,
\]
ここで \( ξ_1 \) は \( f(x) \) と \( f(x) + hf'(ξ) \) の間にある数である。したがって、関数 \( F\{f(x)\} \) の導函数は
\[
\lim_{h \to 0} f'(ξ)F'ξ_1 = f'(x)F'\{f(x)\}
\]
となる。これは \( h \to 0 \) のとき \( ξ \to x \) かつ \( ξ_1 \to f(x) \) であるからである。]

\ Paragraph{126.} 平均値の定理は、以下の結果を証明するのに重要な役割を果たす：もし \( \phi'(x) = 0 \) がある区間におけるすべての \( x \) の値に対して成り立つならば、\( \phi(x) \) はその区間にわたって定数である。

なぜなら、任意の \( a \) と \( b \) をこの区間に属する \( x \) の値として、
\[
\phi(b) - \phi(a) = (b - a)\phi'\{a + θ(b - a)\} = 0.
\]
が成り立つからである。これはすぐに、もし \( \phi'(x) = ψ'(x) \) がある区間におけるすべての \( x \) の値に対して成り立つならば、関数 \( \phi(x) \) と \( ψ(x) \) はその区間にわたって定数を除いて等しいことを意味する。

\ Paragraph{127.} 積分。この章では、与えられた関数 \( φ(x) \) の導函数を求めることができる場合について様々な例を見た。逆の問題を考えるのは自然である：与えられた導函数を持つ関数を決定すること。

与えられた関数 \( ψ(x) \) があるとする。そのような関数 \( φ(x) \) を求めたい。すなわち、\( φ'(x) = ψ(x) \) の関数を求めることである。少し考えれば、この問題は実質的に三つの部分に分解されることがわかる。

(1) 首に置くべきことは、このような関数 \( φ(x) \) が実際に存在するかどうかを知ることである。これは、そのような関数が存在する場合でも、それを簡単な公式で表すことができるか否かという問題とは区別して考えるべきである。

\par\noindent\textbf{・} {(2)} 一つ以上のような関数が存在する可能性を知りたい。つまり、我々の問題は一意的な解を持つ問題であるかどうかを知りたい。もし一意的でない場合、異なる解同士にはどのような単純な関係があるかを理解し、特定の解を基に他の全ての解を表現できる方法を見つけることが望まれる。

\par\noindent\textbf{・} {(3)} 解が存在する場合、その実際的な表式を見つけ出す方法を知りたい。

これらの三つの異なる質問は、関数の微分に関する三つの対応する質問と比較することで光を当てるだろう。

\par\noindent\textbf{・} {(1)} 関数$\phi(x)$はすべての$x$に対して導関数を持つことができる。例えば$m$が正の整数であるときの$x^m$や$\sin x$のように、常に導関数を持つことが多いが、必ずしもそうではない。例えば$\sqrt{x}$や$\tan x$や$\sec x$のような関数は一般的には導関数を持つが、必ずしも常にそうではない。また、導関数を持たない関数も存在する。例えば\Ex{xxxvii}.20で考慮された、どこでも連続性を欠く関数は明らかにどの$x$の値に対しても導関数を持っていない。この章では、特定の$x$の値以外で連続な関数に限定しているが、$\sqrt{x}$という例から分かるように、連続な関数でも特定の$x$の値（本例では$x=0$）において導関数を持たない可能性がある。常に導関数を持たない連続な関数や、常に接線を持つことのない連続曲線が存在するかどうかは、現在我々には未解決の問題である。常識的にはノーだが、既に\SecNo[§]{111}で述べたように、これは高次の数学が常識を誤った場合の一つである。

ただし、少なくともこの質問「\(\phi(x)\) は導関数 \(\phi'(x)\) を持つか？」に対する答えは、状況によって異なることが明確である。また、逆の問題「\(\psi(x)\) が導関数であるような関数 \(\phi(x)\) は存在するか？」に対しても異なる答えがあることも予想される。

既に述べたように、その答えが「ノー」になる場合もある：例えば \(x\) が \(0\) より小さい、等しい、または大きい場合に \(a\)、\(b\)、または \(c\) となる関数 \(\psi(x)\) の場合、答えは「ノー」である（例 xlvii.3）。

これは与えられた関数が不連続な場合である。しかし、以降では常に \(\psi(x)\) を連続とする。その場合、答えは「はい」：「\(\psi(x)\) が連続であれば、導関数 \(\phi'(x) = \psi(x)\) を持つ関数 \(\phi(x)\) は常に存在する」ということである。この証明は第 VII 章で述べられる。

次に、問題 (2) について述べる。微分の場合には導関数の直接的な定義があり、最初から唯一のものしか存在しないことが明らかである。逆問題の場合でも答えはほぼ同様に単純である：\(\phi(x)\) が問題の一つの解であれば、任意の定数 \(C\) の値に対して \(\phi(x) + C\) もまた解であり、すべての可能な解は \(\phi(x) + C\) の形で表される。これは第 126 節からすぐに導かれる。

次に、問題 (3) について述べる。ある有限な通常の関数記号による組み合わせで定義された関数 \(\psi(x)\) に対する実際的な問題は比較的単純である。逆問題は遥かに難しく、その困難性の本質は後ほど明らかになる。

\par\medskip\noindent\textbf{【定義】}\par
$\psi(x)$が$\phi(x)$の導関数であるとき、$\phi(x)$を$\psi(x)$の積分または積分関数と呼ぶ。$\phi(x)$から$\psi(x)$を得る操作を積分という。
\par\medskip

積分の実用的な問題について、この章の前半部の結果により、いくつかの一般的な関数の積分をすぐに書き出すことができる。すなわち
\[
\int x^{m}\, dx = \frac{x^{m+1}}{m + 1},\quad
\int \cos x\, dx = \sin x,\quad
\int \sin x\, dx = -\cos x.
\tag{1}
\]

これらの公式は、右辺の関数が積分記号の下にある関数の一つの積分であることを意味する。当然のことながら、最も一般的な積分は、前者に定数$C$を加えることで得られる。この定数$C$を積分の任意定数という。
\ページ分割{231}

しかしながら、最初の公式には例外があり、それが$m = -1$の場合である。この場合、公式は意味をなさない。これは予想される結果であり、既に述べたこと（§44）によれば、$1/x$は任意の多項式または有理分数の導関数ではないことを確認しているからである。

次の章で、実際には関数 \( F(x) \) が存在し、\( D_x F(x) = 1/x \) を満たすことを証明します。現時点ではその存在を仮定します。この関数 \( F(x) \) は多項式や有理関数ではないことが証明できますし、代数的関数でもないことも証明できます。実際には、\( F(x) \) はこれまでに考慮したどのクラスの関数とも独立した新しい関数であり、それらに対応する機能記号で有限な組み合わせによって表現できないことが証明できます。この証明は本書には収められませんが、その性質についてさらに議論されているのは \okrickRef{第9章}{Ch. IX} です。

まず \( x \) が正の場合を考えます。その場合、
\[
\int \frac{dx}{x} = \log x,
\Tag{(2)}
\]
と書きます。この式の右辺の関数を対数関数と呼びます：これは \( x \) の正な値に対して定義されます。

次に \( x \) が負の場合を考えます。その場合、\( -x \) は正であり、従って \( \log(-x) \) は先の定義で定義されます。また
\[
\frac{d}{dx} \log(-x) = \frac{-1}{-x} = \frac{1}{x},
\]
なので、\( x \) が負の場合には
\[
\int \frac{dx}{x} = \log(-x).
\Tag{(3)}
\]

式 \Eq{(2)} と \Eq{(3)} を統一した式は
\[
\int \frac{dx}{x} = \log(\pm x) = \log|x|,
\Tag{(4)}
\]
であり、この式の不定符号は \( \pm x \) が正になるように選ぶべきです：これらの式は \( x = 0 \) を除くすべての実数値に対して成り立ちます。

\par\medskip\noindent\textbf{【注釈】}\par
$\log x$の最も基本的な性質は、第IX章で証明されるものであり、その式は次のようになります：
\[
\log 1 = 0,\quad
\log (1/x) = -\log x,\quad
\log xy = \log x + \log y,
\]
これらの式のうち、第二式は第一式と第三式から明らかに導けます。これらの公式のいずれかを仮定する必要はありませんが、これらにより式をよりコンパクトな形で表すことが時として可能になります。

最後の式から$\log x^{2}$は$x > 0$の場合には$2\log x$であり、$x < 0$の場合には$2\log(-x)$であり、いずれの場合も$2\log |x|$と等しいことが分かります。したがって、これらの公式のいずれかは次の式に相当します：
\[
\int \frac{dx}{x} = \tfrac{1}{2}\log x^{2}.
\Tag{(5)}
\]
\par\medskip

式(\(1\))～式(\(3\))の五つの公式は積分計算における最も基本的な標準形です。これらに加えて、次の二つも追加すべきです：
\[
\int \frac{dx}{1 + x^{2}} = \arctan x,\quad
\int \frac{x}{\sqrt{1 - x^{2}}} = ±\arcsin x.\footnotemark
\Tag{(6)}
\]
\footnotetext{不明確な符号については、第119節を参照してください。}

\Paragraph{129. 多項式} 第113節のすべての一般定理は当然、積分に関する定理として述べることもできます。したがって、最初に次の公式があります：
\begin{gather*}
\int \{f(x) + F(x)\}\, dx = \int f(x) dx + \int F(x)\, dx,
\Tag{(1)}\\
\int kf(x)\, dx = k\int f(x)\, dx.
\Tag{(2)}
\end{gather*}

当然、任意定数は適切に調整されています。したがって、公式(\(1\))は$f(x)$の任意の積分と$F(x)$の任意の積和が$f(x) + F(x)$の積分であることを主張します。

これらの定理により、任意の形の関数 \(\sum A_{\nu} f_{\nu}(x)\) の積分を一括りに書き出すことができます。これは有限個の定数倍の関数の和であり、その積分が既知の場合です。特に、任意の多項式の積分も書き出せます：
\[
\int (a_{0}x^{n} + a_{1}x^{n-1} + \dots + a_{n})\, dx
  = \frac{a_{0}x^{n+1}}{n + 1} + \frac{a_{1}x^{n}}{n} + \dots + a_{n}x.
\]
ページ分割：233

130. 有理関数。多項式の積分を計算した後は、自然に有理関数へと注意を向けます。任意の有理関数 \(R(x)\) を \SecNo[§]{117} の標準形として表すと、これは多項式 \(\Pi(x)\) と \(x - \alpha\) のべき乗に比例する項の和で表されます。

多項式の積分はすぐに書き出せますが、\(p = 1\) となる項以外の他の項については、次のようになります：
\[
\int \frac{A}{(x - \alpha)^{p}}\, dx
  = -\frac{A}{p - 1}\, \frac{1}{(x - \alpha)^{p-1}},
\]
ここで \(\alpha\) は実数でも複素数でも同じです（\SecNo[§]{117}）。

\(p = 1\) となる項については、やや困難が伴います。\SecNo[§]{113} の定理 (6) よりすぐに分かります：
\[
\int F'\{f(x)\}\, f'(x)\, dx = F\{f(x)\}.
\Tag{(3)}
\]
特に、\(f(x) = ax + b\) とし、ここで \(a\) と \(b\) は実数とするとき、\(\phi(x) = F(x)\)、\(\psi(x) = F'(x)\) と定義すると \(\phi(x)\) は \(\psi(x)\) の積分であるため、次の式が得られます：
\[
\int \psi(ax + b)\, dx = \frac{1}{a}\phi(ax + b).
\Tag{(4)}
\]

例えば、
\[
\int \frac{dx}{ax + b} = \frac{1}{a} \log|ax + b|,
\]
特に、\(\alpha\) が実数のときには
\[
\int \frac{dx}{x - \alpha} = \log|x - \alpha|.
\]
したがって、\(R(x)\) の \(p = 1\) で \(\alpha\) が実数の項の積分は書き出せます。残るのは \(p = 1\) で \(\alpha\) が複素数の項です。

これらの問題に対処するため、係数がすべて実数であるという制約仮定を導入します。つまり、$Q(x) = 0$ の根 $\alpha = \gamma + \delta i$ が多重度 $m$ で存在するとすれば、その共役 $\bar{\alpha} = \gamma - \delta i$ も同様の多重度で根となります。また、部分分数展開において $(x - \alpha)^{p}$ の項 $A_{p}/(x - \alpha)^{p}$ が含まれている場合、$(x - \bar{\alpha})^{p}$ の項 $\bar{A}_{p}/(x - \bar{\alpha})^{p}$ も存在します。ここで $\bar{A}_{p}$ は $A_{p}$ の共役です。これは代数的手法による部分分数の求め方から導かれるもので、詳細は代数学の教科書に詳しく説明されています。\footnote{例えば Chrystal の \textit{Algebra} 第 1 卷 pp.151--9 を参照してください。}

したがって、$R(x)$ の部分分数展開において $(\lambda + \mu i)/(x - \gamma - \delta i)$ の項が存在するとすれば、$(\lambda - \mu i)/(x - \gamma + \delta i)$ の項も存在します。これらの 2 つの項の和は
\[
\frac{2\{\lambda(x - \gamma) - \mu\delta\}}{(x - \gamma)^{2} + \delta^{2}}
\]
となります。この分数は実質的に、$ax^{2} + 2bx + c$ の形の分母を持つ最一般の分数
\[
\frac{Ax + B}{ax^{2} + 2bx + c}
\]
で $b^{2} < ac$ の場合に等しいものです。読者は $\lambda$,~$\mu$, $\gamma$,~$\delta$ を $A$,~$B$, $a$,~$b$,~$c$ で表す公式を容易に確認できるでしょう：
\[
\lambda = A/2a,\quad
\mu = -D/(2a\sqrt{\Delta}),\quad
\gamma = -b/a,\quad
\delta = \sqrt{\Delta}/a,
\]
ここで $\Delta = ac - b^{2}$ で、$D = aB - bA$。

式 (3) の $F\{f(x)\}$ を $\log |f(x)|$ と仮定すると、次の等式が得られます：
\[
\int \frac{f'(x)}{f(x)}\, dx = \log |f(x)|;
\Tag{(5)}
\]
さらに $f(x) = (x - \lambda)^{2} + \mu^{2}$ と仮定すると、次の等式が得られます：
\[
\int \frac{2(x - \lambda)}{(x - \lambda)^{2} + \mu^{2}}\, dx
  = \log\{(x - \lambda)^{2} + \mu^{2}\}.
\]
さらに、節 §128 の式 (6) と上記の式 (4) を用いると、次の等式が得られます：
\[
\int \frac{-2\delta\mu}{(x - \lambda)^{2} + \mu^{2}}\, dx
  = -2\delta \arctan \left(\frac{x - \lambda}{\mu}\right).
\]

これらの2つの公式は、我々が$R(x)$の表現において考慮してきた2項の和を積分するのに役立ちます。これにより、分母のすべての因子が決定できるなら、任意の実数有理関数の積分を書き出すことができます。

そのような関数の積分は以下の要素から成り立っています。
\par\medskip\noindent\textbf{【結果】}\par
多項式の和、いくつかの形の有理関数
\[
-\frac{A}{p - 1}\, \frac{1}{(x - \alpha)^{p-1}},
\]
いくつかの対数関数、およびいくつかの逆正接関数。
\par\medskip

$\alpha$が複素数である場合、先に述べた有理関数は常にその共役な複素数を置き換えた別の関数と共に現れ、2つの関数の和は実数有理関数となります。

\par\medskip\noindent\textbf{【例 XLVIII.】}\par
\Item{1.} $ax^2 + 2bx + c < 0$ のとき
\[
\int \frac{Ax + B}{ax^{2} + 2bx + c}\, dx
  = \frac{A}{2a} \log |X| + \frac{D}{2a \sqrt{-\Delta}}
    \log \left|\frac{ax + b - \sqrt{-\Delta}}{ax + b + \sqrt{-\Delta}}\right|
\]
($X = ax^2 + bx + c$) であり、$\Delta > 0$ のとき
\[
\int \frac{Ax + B}{ax^{2} + 2bx + c}\, dx
  = \frac{A}{2a} \log |X| + \frac{D}{2a \sqrt{\Delta}}
     \arctan \left(\frac{ax + b}{\sqrt{\Delta}}\right)
\]
ここで、$\Delta$と$D$は第234ページで同じ意味を持つ。

\Item{2.} $ac = b^2$ の特別な場合において
\[
\int \frac{Ax + B}{ax^{2} + 2bx + c}\, dx
  = -\frac{D}{a(ax + b)} + \frac{A}{a} \log |ax + b|.
\]

\Item{3.} $Q(x) = 0$ の根がすべて実数で異なるとき、かつ$P(x)$の次数が$Q(x)$よりも低いとき
\[
\int R(x)\, dx = \sum \frac{P(\alpha)}{Q'(\alpha)} \log |x - \alpha|,
\]
ここで$\alpha$は$Q(x) = 0$のすべての根に対応します。

[$\alpha$に対する分数の形は、
\[
\frac{Q(x)}{x - \alpha} \to Q'(\alpha),\quad
(x - \alpha) R(x) \to \frac{P(\alpha)}{Q'(\alpha)},
\]
$x \to \alpha$ のときから導くことができます。]

\item 4. $Q(x)$ のすべての根が実数であり、$\alpha$ が重根で他の根は単根であるとき、$P(x)$ が $Q(x)$ より次数が低い場合、積分は $A/(x - \alpha) + A'\log |x - \alpha| + \sum B\log |x - \beta|$ となる。ここで
\[
A = -\frac{2P(\alpha)}{Q''(\alpha)},\quad
A' = \frac{2\{3P'(\alpha) Q''(\alpha) - P(\alpha) Q'''(\alpha)\}}
          {3\{Q''(\alpha)\}^{2}},\quad
B = \frac{P(\beta)}{Q'(\beta)},
\]
$\sum$ は $Q(x) = 0$ の根 $\beta$（$\beta \neq \alpha$）に対して行う。

\item 5. 次の積分を計算する。
\[
\int \frac{dx}{\{(x - 1)(x^{2} + 1)\}^{2}}.
\]

[部分分数分解は
\[
\frac{1}{4(x - 1)^{2}}
  - \frac{1}{2(x - 1)}
  - \frac{i}{8(x - i)^{2}}
  + \frac{2 - i}{8(x - i)}
  + \frac{i}{8(x + i)^{2}}
  + \frac{2 + i}{8(x + i)},
\]
で、積分は
\[
-\frac{1}{4(x - 1)} - \frac{1}{4(x^{2} + 1)}
  - \tfrac{1}{2} \log |x - 1|
  + \tfrac{1}{4} \log (x^{2} + 1)
  + \tfrac{1}{4} \arctan x
\]
となる。]

\item 6. 次の積分を計算する。
\[
\frac{x}{(x - a)(x - b)(x - c)},\quad
\frac{x}{(x - a)^{2}(x - b)},\quad
\frac{x}{(x - a)^{2} (x - b)^{2}},\quad
\frac{x}{(x - a)^{3}},
\]
\[
\frac{x}{(x^{2} + a^{2})(x^{2} + b^{2})},\quad
\frac{x^{2}}{(x^{2} + a^{2})(x^{2} + b)^{2}},\quad
\frac{x^{2} - a^{2}}{x^{2}(x^{2} + a^{2})},\quad
\frac{x^{2} - a^{2}}{x(x^{2} + a^{2})^{2}}.
\]

\Item{7.} 次の公式を証明せよ：
\begin{alignat*}{3}
\int \frac{dx}{1 + x^{4}}
  \&= \frac{1}{4\sqrt{2}} \biggl\{
    \&\&\log \biggl(\frac{1 + x\sqrt{2} + x^{2}}{1 - x\sqrt{2} + x^{2}}\biggr)
   \&\&+ 2\arctan \biggl(\frac{x\sqrt{2}}{1 - x^{2}}\biggr)\biggr\},\\
\%
\int \frac{x^{2}\, dx}{1 + x^{4}}
  \&= \frac{1}{4\sqrt{2}} \biggl\{
    \&-\&\log \biggl(\frac{1 + x\sqrt{2} + x^{2}}{1 - x\sqrt{2} + x^{2}}\biggr)
   \&\&+ 2\arctan \biggl(\frac{x\sqrt{2}}{1 - x^{2}}\biggr)\biggr\},\\
\%
\int \frac{dx}{1 + x^{2} + x^{4}}
  \&= \frac{1}{4\sqrt{3}}\biggl\{
     \&\sqrt{3}\&\log \biggl(\frac{1 + x + x^{2}}{1 - x + x^{2}}\biggr)
   \&\&+ 2\arctan \biggl(\frac{x\sqrt{3}}{1 - x^{2}}\biggr)\biggr\}.
\end{alignat*}
\par\medskip

\par\medskip\noindent\textbf{【注釈】}\par
\Paragraph{131. 理数函数の実用的な積分に関する注記}
第130節の解析から、任意の実理数函数$R(x)$の積分を求めることができる一般的方法が得られる。ただし、方程式$Q(x) = 0$を解けることが前提である。単純な場合（上記例5のように）では、この方法の適用は比較的簡単である。より複雑な場合では、作業量が時として許容不能となり、他の装置を使用せざるを得ない。本書の目的には実用的な積分問題の詳細な論述は含まれていない。より詳しい情報が必要な読者は、Goursatの『解析教程』第2版第1巻（pp. 246 et seq.）、Bertrandの『积分计算』およびBowes and Bowes社によるBromwichの小冊子『基本积分』（1911年）を参照されたい。

方程式$Q(x) = 0$が代数的に解けない場合、部分分数法は自然に失敗し、他の方法に頼らなければならない。
\footnote
  {作者の小冊子「単一变量函数的积分」（剑桥数学论文集，第2号，1915年第二版）参照。実際ではこれほど頻繁には起こらない。}
\par\medskip

\#\#\# 132. アルキメデス関数の積分

次に、アルキメデス関数の積分について問うことに自然な流れがある。$y$が$x$のアルキメデス関数であるとき、$y$を$x$で積分する問題を考える必要がある。しかし、以下のようなより一般的な積分を考えることも便利である：
\[
\int R(x, y)\, dx,
\]
ここで $R(x, y)$ は $x$ と $y$ の任意の有理関数である。この形の一般性は表面的なものであり、(\Ex{xiv}.~6) において $R(x, y)$ 自体が $x$ のアルキメデス関数であることが示される。このような形を選択した理由は実質的に利便性によるもので、例えば
\[
\frac{px + q + \sqrt{ax^{2} + 2bx + c}}
     {px + q - \sqrt{ax^{2} + 2bx + c}}
\]
のような関数は、$x$ と単純なアルキメデス関数 $\sqrt{ax^{2} + 2bx + c}$ の有理関数として扱う方が遥かに便利である。

\#\#\# 133. 置換積分と有理化による積分

\SecNo[§]{130}の式(\Eq{(3)})から、$\ds\int \psi(x)\, dx = \phi(x)$ のとき
\[
\int \psi\{f(t)\}\, f'(t)\, dt = \phi\{f(t)\}.
\]
が導かれる。

この方程式は、\(\psi(x)\) の積分が直接明らかでない多くの場合において積分を決定する方法を提供します。その規則として述べることは次の通りです：新しい変数 \(t\) についての任意の関数 \(f(t)\) を用いて \(x = f(t)\) と置き、これを \(f'(t)\) で掛けて \(\psi\{f(t)\}\, f'(t)\) の積分を求めます。その結果を \(x\) で表します。この規則によって導かれる \(t\) の関数の積分が容易に計算できることがしばしば見つかります。例えば、それが有理関数である場合や、\(x\) と \(t\) の関係を適切に選ぶことでそのような状況にすることができます。したがって、\(\sqrt{x}\) の有理関数 \(R\) の積分は、置換 \(x = t^{2}\) によって \(2tR(t^{2})\) の積分へと簡素化され、すなわち \(t\) の有理関数の積分へと簡素化されます。この積分法は「有理化による積分」と呼ばれ、非常に広範に適用できます。

その方法が現在考慮している問題に対してどのように適用されるかは明らかです。以下は結果です：もし \(x\) と \(y\) をともに \(t\) の有理関数 \(R_{1}(t)\) および \(R_{2}(t)\) として表すことができる変数 \(t\) を見つけることができれば、
\[
\int R(x, y)\, dx = \int R\{R_{1}(t), R_{2}(t)\}\, R_{1}'(t)\, dt,
\]
となります。後者の積分は \(t\) の有理関数の積分であり、第 130 節で述べた方法によって計算できます。

補助変数 \(t\) を \(x\) および \(y\) と関連付けることができるかどうかについての一般的な議論には、現在の範囲を超えるため言及しません。私たちはいくつかの単純かつ興味深い特殊なケースのみを考慮します。

\ Paragraph{134. 二項曲線に関連する積分} $x$ と $y$ が次のような方程式によって結びつけられていると仮定します：
\[
ax^{2} + 2hxy + by^{2} + 2gx + 2fy + c = 0;
\]
すなわち、$y$ を $x$ の関数として考えたグラフが二項曲線であるとします。二項曲線上の任意の点 $(\xi, \eta)$ について、$x - \xi = X$, $y - \eta = Y$ と置きます。$x$ と $y$ の関係を $X$ と $Y$ で表すと、次の形になります：
\[
aX^{2} + 2hXY + bY^{2} + 2GX + 2FY = 0,
\]
ここで $F = h\xi + b\eta + f$, $G = a\xi + h\eta + g$ です。この方程式に $Y = tX$ を代入すると、$X$ と $Y$ は $t$ の有理関数として表すことができ、従って $x$ と $y$ も同様に表すことができます。具体的には次の式が得られます：
\[
x - \xi = -\frac{2 (G + Ft)}{a + 2ht + bt^{2}},\quad
y - \eta = -\frac{2t(G + Ft)}{a + 2ht + bt^{2}}.
\]
したがって、最後の節で説明した有理化の過程は実行できます。

読者は次を確認することです：
\[
hx + by + f = -\frac{1}{2}(a + 2ht + bt^{2}) \frac{dx}{dt},
\]
従って
\[
\int \frac{dx}{hx + by + f} = -2 \int \frac{dt}{a + 2ht + bt^{2}}.
\]

$h^{2} > ab$ のとき、次の手順で進めるといくつかの面で有利です。二項曲線は漸近線が直線
\[
ax^{2} + 2hxy + by^{2} = 0,
\]
すなわち
\[
b(y - \mu x)(y - \mu' x) = 0,
\]
と平行である双曲線です。$y - \mu x = t$ と置くと、次のようになります：
\[
y - \mu x = t,\quad
y - \mu' x = -\frac{2gx + 2fy + c}{bt},
\]
従って $x$ と $y$ はこれらの方程式から $t$ の有理関数として計算できます。この過程を重要な特殊ケースへの応用で説明します。

\par\medskip\noindent\textbf{【注釈】}\par
\Paragraph{135.積分 $\ds\int \frac{dx}{\sqrt{ax^{2} + 2bx + c}}$.} 特に $y^{2} = ax^{2} + 2bx + c$ において $a > 0$ の場合、$y + x\sqrt{a} = t$ と置くと、

\[
2\frac{dx}{dt} = \frac{(t^{2} + c)\sqrt{a} + 2bt}{(t\sqrt{a} + b)^{2}},\quad
2y = \frac{(t^{2} + c)\sqrt{a} + 2bt}{t\sqrt{a} + b},
\]

よって、

\[
\int \frac{dx}{y}
  = \int \frac{dt}{t\sqrt{a} + b}
  = \frac{1}{\sqrt{a}} \log \left|x\sqrt{a} + y + \frac{b}{\sqrt{a}}\right|.
\Tag{(1)}
\]

特に $a = 1$, $b = 0$, $c = a^{2}$ の場合や $a = 1$, $b = 0$, $c = -a^{2}$ の場合、

\[
\int \frac{dx}{\sqrt{x^{2} + a^{2}}} = \log \{x + \sqrt{x^{2} + a^{2}}\},\quad
\int \frac{dx}{\sqrt{x^{2} - a^{2}}} = \log  |x + \sqrt{x^{2} - a^{2}}],
\Tag{(2)}
\]

これらの等式は直ちに微分によりその真偽を確認できる。これら公式には、$a < 0$ の一般積分の一つに対応する第3の公式

\[
\int \frac{dx}{\sqrt{a^{2} - x^{2}}} = \arcsin(x/a),
\Tag{(3)}
\]

が付随すべきである。式(3)においては $a > 0$ と仮定しているが、$a < 0$ の場合は $\arcsin(x/|a|)$ となる（参照：§119）。実際の計算では、これらの標準形の一つに帰着させる方法で一般積分を評価すべきである。

式(3)は式(2)と非常に異なる外見をしているが、その関連性を理解するには第X章を読むことが必要である。
\par\medskip

136.積分 $\ds\int \frac{\lambda x + \mu}{\sqrt{ax^{2} + 2bx + c}}\, dx$. この積分は、前節の結果を利用することで任意の場合において計算できます。以下のように進めることが最も便利です。既に
\begin{gather*}
\lambda x + \mu = \left(\frac{\lambda}{a}\right) (ax + b) + \mu - \left(\frac{\lambda b}{a}\right),\\
\int \frac{ax + b}{\sqrt{ax^{2} + 2bx + c}}\, dx = \sqrt{ax^{2} + 2bx + c},
\end{gather*}
が成り立つことから、
\[
\int \frac{\lambda x + \mu}{\sqrt{ax^{2} + 2bx + c}}\, dx
= \frac{\lambda}{a} \sqrt{ax^{2} + 2bx + c}
+ \left(\mu - \frac{\lambda b}{a}\right)
\int \frac{dx}{\sqrt{ax^{2} + 2bx + c}}
\]
を得ます。

最後の積分において $a$ は正でも負でもよいです。$a$ が正の場合、$x\sqrt{a} + \frac{b}{\sqrt{a}} = t$ と置くことで
\[
\frac{1}{\sqrt{a}} \int \frac{dt}{\sqrt{t^{2} + \kappa}},
\]
を得ます。ここで $\kappa = \frac{ac - b^{2}}{a}$ です。$a$ が負の場合、$A = -a$ と置き $x\sqrt{A} - \frac{b}{\sqrt{A}} = t$ と置くことで
\[
\frac{1}{\sqrt{-a}} \int \frac{dt}{\sqrt{-\kappa - t^{2}}}
\]
を得ます。

したがって、任意の場合においてこの積分の計算は、§135で考慮された積分の計算に依存し、その積分は以下の三つの形のいずれかに帰着されることがわかります。
\[
\int \frac{dt}{\sqrt{t^{2} + a^{2}}},\quad
\int \frac{dt}{\sqrt{t^{2} - a^{2}}},\quad
\int \frac{dt}{\sqrt{a^{2} - t^{2}}}.
\]

\par\medskip\noindent\textbf{【注釈】}\par
137.積分$\ds\int (\lambda x + \mu) \sqrtp{ax^{2} + 2bx + c}\, dx$. 同様の方法で、
\[
\int(\lambda x + \mu) \sqrtp{ax^{2} + 2bx + c}\, dx =
\left(\frac{\lambda}{3a}\right) (ax^{2} + 2bx + c)^{3/2}
+ \left(\Add{\mu} - \frac{\lambda b}{a}\right) \int \sqrtp{ax^{2} + 2bx + c}\, dx;
\]
を得る。最後の積分は、次の三つの形のいずれかに帰着される：
\[
\int \sqrtp{t^{2} + a^{2}}\, dt,\quad
\int \sqrtp{t^{2} - a^{2}}\, dt,\quad
\int \sqrtp{a^{2} - t^{2}}\, dt.
\]
これらの積分を得るためには、この点で一般の積分定理を導入するのが便利である。

138.部分積分法による積分。部分積分法の定理は、\SecNo[§]{113}で証明された積分乗法微分則の別の表現形に過ぎない。\SecNo[§]{113}の定理(3)から即座に
\[
\int f'(x)F(x)\, dx = f(x)F(x) - \int f(x)F'(x)\, dx.
\]
を得る。積分したい関数が$f'(x)F(x)$という形で表される場合、$f(x)F'(x)$が積分可能であることがあれば便利である。例えば、$\phi(x) = x\psi(x)$とし、$\psi(x)$が既知の関数$\chi(x)$の二階微分であるとする。このとき
\[
\int\phi(x)\, dx =
\int x\chi''(x)\, dx =
x\chi'(x) - \int \chi'(x)\, dx =
x\chi'(x) - \chi(x).
\]
を得る。
\par\medskip

\par\medskip\noindent\textbf{【注釈】}\par
この積分法の働きを説明するために、最後の節の積分に適用します。$f(x) = ax + b,\quad F(x) = \sqrt{ax^{2} + 2bx + c} = y$と置くと、

\[
a\int y\, dx
  = (ax + b)y - \int \frac{(ax + b)^{2}}{y}\, dx
  = (ax + b)y - a\int y\, dx + (ac - b^{2}) \int \frac{dx}{y},
\]

よって、

\[
\int y\, dx = \frac{(ax + b)y}{2a} + \frac{ac - b^{2}}{2a} \int \frac{dx}{y};
\]

最後の積分はすでに（135節）でどのように定めるかが述べられています。
\par\medskip

\par\medskip\noindent\textbf{【例 XLIX.】}\par
\Item{1.} $a > 0$ のとき、

\[
\int \sqrt{x^{2} + a^{2}}\, dx
  = \frac{1}{2}x \sqrt{x^{2} + a^{2}}
   + \frac{1}{2}a^{2} \log \left\{x + \sqrt{x^{2} + a^{2}}\right\},
\]

\[
\int \sqrt{x^{2} - a^{2}}\, dx
  = \frac{1}{2}x \sqrt{x^{2} - a^{2}}
   - \frac{1}{2}a^{2} \log |x + \sqrt{x^{2} - a^{2}}],
\]

\[
\int \sqrt{a^{2} - x^{2}}\, dx
  = \frac{1}{2}x \sqrt{a^{2} - x^{2}}
   + \frac{1}{2}a^{2} \arcsin(x/a).
\]

\PageSep{241}

\Item{2.} $x = a\sin\theta$ の置換により、

\[
\int \frac{dx}{\sqrt{a^{2} - x^{2}}},\quad
\int \sqrt{a^{2} - x^{2}}\, dx
\]

を計算し、結果が135節と例1で得られたものに一致することを確認します。

\Item{3.} $m$ が任意の有理数であるとき、

\[
\int x(x + a)^{m}\, dx
\]

を三つの方法で計算します。すなわち、(i) 部分積分法、(ii) $(x + a)^{m} = t$ の置換、(iii) $x = (x + a) - a$ とし、結果が一致することを確認します。

\Item{4.} $ax + b = \frac{1}{t}$, $x = \frac{1}{u}$ の置換により、

\[
\int \frac{dx}{y^{3}} = \frac{ax + b}{\Delta y},\quad
\int \frac{x\, dx}{y^{3}} = -\frac{bx + c}{\Delta y}
\]

を証明します。
\par\medskip

\par\noindent\textbf{5.}  \(x - a > 0\) かつ \(b - x > 0\) の条件で \(\ds\int \frac{dx}{\sqrt{(x - a)(b - x)}}\) を前節の方法、次に \((b - x)/(x - a) = t^2\) の置換法、さらに \(x = a\cos^2\theta + b\sin^2\theta\) の置換法で計算し、結果が一致することを確認せよ。

\par\noindent\textbf{6.}  \(\sqrt{(x - a)(b - x)}\) および \(\sqrt{(b - x)/(x - a)}\) を積分せよ。

\par\noindent\textbf{7.}  \(2x + a + b = \frac{1}{2}(a - b)\{t^2 + (1/t)^2\}\) の置換法または \(\sqrt{x + a} - \sqrt{x + b}\) に分数を乗じることで、\(a > b\) のとき
\[
\int \frac{dx}{\sqrt{x + a} + \sqrt{x + b}}
  = \frac{1}{2}\sqrt{a - b}\left(t + \frac{1}{3t^3}\right)
\]
が成り立つことを示せ。

\par\noindent\textbf{8.}  \(\ds\int \frac{dx}{(x + a)^{3/2} + (x - a)^{3/2}}\) を有理関数の積分に帰着させる置換法を見つけよ。 \MathTrip{1899}

\par\noindent\textbf{9.}  \(R\{x, \sqrt[n]{ax + b}\}\) の積分 \(\ds\int R\{x, \sqrt[n]{ax + b}\} dx\) を \(ax + b = y^n\) の置換法で有理関数の積分に帰着させることを示せ。

\par\noindent\textbf{10.}  次の公式を証明せよ：
\[
\int f''(x) F(x)\, dx = f'(x) F(x) - f(x) F'(x) + \int f(x) F''(x)\, dx
\]
一般化して
\begin{multline*}
\int f^{(n)}(x) F(x)\, dx \\
  = f^{(n-1)}(x) F(x) - f^{(n-2)}(x) F'(x) + \dots
  + (-1)^n \int f(x) F^{(n)}(x)\, dx.
\end{multline*}

\par\noindent\textbf{11.}  \(p\) および \(q\) が有理数であるとき、\(\ds\int (1 + x)^p x^q\, dx\) の積分は以下の 3 種類の場合で行うことができる：
(i) \(p\) が整数のとき、\(x = u^s\) と置換する（ここで \(s\) は \(q\) の分母）；
(ii) \(q\) が整数のとき、\(1 + x = t^s\) と置換する（ここで \(s\) は \(p\) の分母）；
(iii) \(p + q\) が整数のとき、\(1 + x = xt^s\) と置換する（ここで \(s\) は \(p\) の分母）。

\par\noindent\textbf{・} [12.] 不定積分 $\ds\int x^{m}(ax^{n} + b)^{q}\, dx$ は、変数置換 $ax^{n} = bt$ を用いることで、前記の不定積分に帰着される。

\par\noindent\textbf{・} [13.] 不定積分 $\ds\int R\{x, \sqrt{ax + b}, \sqrt{cx + d}\}\, dx$ は、以下の変数置換を用いて有理関数の不定積分に帰着される：
\[
4x = -\frac{b}{a} \left(t + \frac{1}{t}\right)^{2} - \frac{d}{c} \left(t - \frac{1}{t}\right)^{2}.
\]

\par\noindent\textbf{・} [14.] $y^{2}(x - y) = x$ のとき、不定積分 $\ds\int R(x, y)\, dx$ を有理関数の不定積分に帰着する。$y = tx$ と置くことで、$x = \frac{1}{t^{2}(1 - t)}$, $y = \frac{1}{t(1 - t)}$ となる。

\par\noindent\textbf{・} [15.] (\ia) $y(x - y)^{2} = x$ のとき、(\ib) $(x^{2} + y^{2})^{2} = a^{2}(x^{2} - y^{2})$ のときも同様に帰着する。(\ia) の場合、$x - y = t$ と置く；(\ib) の場合、$x^{2} + y^{2} = t(x - y)$ と置くことで、
\[
x = \frac{a^{2}t(t^{2} + a^{2})}{t^{4} + a^{4}},\quad
y = \frac{a^{2}t(t^{2} - a^{2})}{t^{4} + a^{4}}.
\]

\par\noindent\textbf{・} [16.] $y(x - y)^{2} = x$ のとき、
\[
\int \frac{dx}{x - 3y} = \frac{1}{2} \log\{(x - y)^{2} - 1\}.
\]

\par\noindent\textbf{・} [17.] $(x^{2} + y^{2})^{2} = 2c^{2}(x^{2} - y^{2})$ のとき、
\[
\int \frac{dx}{y(x^{2} + y^{2} + c^{2})}
  = - \frac{1}{c^{2}}\log\left(\frac{x^{2} + y^{2}}{x - y}\right).
\]

\par\medskip

\par\medskip\noindent\textbf{【注釈】}\par
\paragraph{139. $y^{2} = ax^{2} + 2bx + c$ のときの一般解 $\ds\int R(x, y)\, dx$.} 考察された \SecNo[§]{134} の型と、特異な楕円 $y^{2} = ax^{2} + 2bx + c$ に関連する最も一般的な不定積分は
\[
\int R(x, \sqrt{X})\, dx,
\Tag{(1)}
\]
ここで $X = y^{2} = ax^{2} + 2bx + c$ であり、$R$ は実数関数であるとする。
\par\medskip

The subject of integration is of the form \( \frac{P}{Q} \), where \( P \) and \( Q \) are polynomials in \( x \) and \( \sqrt{X} \). It may therefore be reduced to the form

\[
\frac{A + B\sqrt{X}}{C + D\sqrt{X}}
  = \frac{(A + B\sqrt{X})(C - D\sqrt{X})}{C^{2} - D^{2}X}
  = E + F\sqrt{X},
\]

where \( A \), \( B \), \(\dots\) are rational functions of \( x \). The only new problem which arises is that of the integration of a function of the form \( F\sqrt{X} \), or, what is the same thing, \( G/\sqrt{X} \), where \( G \) is a rational function of \( x \). And the integral

\[
\int \frac{G}{\sqrt{X}}\, dx
\Tag{(2)}
\]

can always be evaluated by splitting up \( G \) into partial fractions. When we do this, integrals of three different types may arise.

(i) In the first place there may be integrals of the type

\[
\int \frac{x^{m}}{\sqrt{X}}\, dx,
\Tag{(3)}
\]

where \( m \) is a positive integer. The cases in which \( m = 0 \) or \( m = 1 \) have been disposed of in Section 136. In order to calculate the integrals corresponding to larger values of \( m \), we observe that

\[
\frac{d}{dx}(x^{m-1}\sqrt{X})
  = (m - 1)x^{m-2} \sqrt{X} + \frac{(ax + b) x^{m-1}}{\sqrt{X}}
  = \frac{\alpha x^{m} + \beta x^{m-1} + \gamma x^{m-2}}{\sqrt{X}},
\]

where \( \alpha \), \( \beta \), \( \gamma \) are constants whose values may be easily calculated. It is clear that, when we integrate this equation, we obtain a relation between three successive integrals of the type Eq.(3). As we know the values of the integral for \( m = 0 \) and \( m = 1 \), we can calculate in turn its values for all other values of \( m \).

(ii) In the second place there may be integrals of the type

\[
\int \frac{dx}{(x - p)^{m}\sqrt{X}},
\Tag{(4)}
\]

where \( p \) is real. If we make the substitution \( x - p = 1/t \), then this integral is reduced to an integral in \( t \) of the type Eq.(3).

以下に、指定された原著英文を中国語無しの格調高い標準日本語で翻訳したものを出力します：

\Itemp{(iii)} 最終的に，分母 \( G \) の複素根に対応する積分が存在する可能性があります。ここでは，そのようなすべての根が単根である最も簡単な場合に限定します（参照：130節）。この場合（同様に§130を参照），\( G \) の共役複素根のペアは以下のタイプの積分を生み出します：
\[
\int \frac{Lx + M}{(Ax^{2} + 2Bx + C) \sqrt{ax^{2} + 2bx + c}}\, dx.
\Tag{(5)}
\]

この積分を評価するために，次のように置きます：
\[
x = \frac{\mu t + \nu}{t + 1},
\]
ここで，\(\mu\) と \(\nu\) を以下のように選択します：
\[
a\mu\nu + b(\mu + \nu) + c = 0,\quad
A\mu\nu + B(\mu + \nu) + C = 0;
\]
したがって，\(\mu\) と \(\nu\) は方程式の根です：
\[
(aB - bA)\xi^{2} - (cA - aC)\xi + (bC - cB) = 0.
\]
この方程式は確かに実数の根を持つため，これは例 xlvi.12 の式（1）と同一の方程式であり，したがって \(\mu\) と \(\nu\) の実数値を求めることが可能です。

積分~式(5)において代入を行うと，次の形になります。
\[
H\int \frac{t\, dt}{(\alpha t^{2} + \beta)\sqrt{\gamma t^{2} + \delta}}
  + K\int \frac{dt}{(\alpha t^{2} + \beta)\sqrt{\gamma t^{2} + \delta}}.
\]
式(6)
これらの積分のうち，二つ目の積分は次の代入によって有理化されます。
\[
\frac{t}{\sqrt{\gamma t^{2} + \delta}} = u,
\]
これにより
\[
\int \frac{dt}{(\alpha t^{2} + \beta) \sqrt{\gamma t^{2} + \delta}}
  = \int \frac{du}{\beta + (\alpha\delta - \beta\gamma) u^{2}}.
\]
ページ分割
まず，最初の積分~式(6)において \(t = 1/u\) を代入すると，二つ目のタイプの積分に変換され，先ほど説明した方法で計算できます。すなわち，\(u/\sqrt{\gamma + \delta u^{2}} = u\)，つまり \(1/\sqrt{\gamma t^{2} + \delta} = v\) と置きます。
ここで説明された積分法は，\(a/A = b/B\) の場合に失敗しますが，そのような場合は \(ax + b = t\) の代入によって簡略化できます。アルゴリズム関数の積分に関する詳細情報については，Stolz, 『微分積分学の基礎』第1巻，331ページ以降，Bromwich, 『基本積分』（Bowes and Bowes, 1911）を参照してください。Sir G. Greenhillによって提案された別の簡略化方法もあります：彼の『積分微分法の一章』，12ページ以降，および引用先の第236ページに掲載されている著者の論文を参照してください。
\par\medskip

例題
\par\medskip\noindent\textbf{【例 L.】}\par
\Item{1.} 次の積分を評価します。
\[
\int \frac{dx}{x \sqrt{x^{2} + 2x + 3}},\quad
\int \frac{dx}{(x - 1) \sqrt{x^{2} + 1}},\quad
\int \frac{dx}{(x + 1) \sqrt{1 + 2x - x^{2}}}.
\]

\Item{2.} 次の等式を証明します。
\[
\int \frac{dx}{(x - p) \sqrt{(x - p)(x - q)}}
  = \frac{2}{q - p} \sqrt{\frac{x - q}{x - p}}.
\]

\Item{3.} \(ag^{2} + ch^{2} = -\nu < 0\) のとき，次の積分が成立します。
\[
\int \frac{dx}{(hx + g) \sqrt{ax^{2} + c}}
  = -\frac{1}{\sqrt{\nu}} \arctan\left[
       \frac{\sqrt{a x^{2} + c}}{ch - agx}
  \right].
\]

4. \( x - x_{0} \)で割った積分 \(\ds\int \frac{dx}{(x - x_{0})y}\), ここで \( y^{2} = ax^{2} + 2bx + c \) は、\( ax_{0}^{2} + 2bx_{0} + c \)が正で等しい \( y_{0}^{2} \) の場合や負で等しい \( -z_{0}^{2} \) の場合に、次のいずれかの形で表すことができる：
\[
-\frac{1}{y_{0}} \log \left| \frac{axx_{0} + b(x + x_{0}) + c + yy_{0}}{x - x_{0}} \right|, \quad
\frac{1}{z_{0}} \arctan \left\{ \frac{axx_{0} + b(x + x_{0}) + c}{yz_{0}} \right\}.
\)

5. \( y = \sqrt{ax^{2} + 2bx + c}/(x - p) \) の代入により
\[
\int \frac{dx}{(x - p) \sqrt{ax^{2} + 2bx + c}}
  = \int \frac{dy}{\sqrt{\lambda y^{2} - \mu}},
\]
ここで \(\lambda = ap^{2} + 2bp + c\), \(\mu = ac - b^{2}\)。[この簡約法は優雅ではあるが、第139節で説明した方法よりも直感的ではない.]

6. \( x \sqrt{3x^{2} + 2x + 1} \) の積分
\[
\int \frac{dx}{x \sqrt{3x^{2} + 2x + 1}}
\]
は、\( x = (1 + y^{2})/(3 - y^{2}) \) の代入によって有理化される。[1911年数学トリップ]

7. 
\[
\int \frac{(x + 1)\, dx}{(x^{2} + 4) \sqrt{x^{2} + 9}}
\]
を計算せよ。

8. 
\[
\int \frac{dx}{(5x^{2} + 12x + 8) \sqrt{5x^{2} + 2x - 7}}
\]
を計算せよ。
[第139節の方法を使用。μとνが満たす方程式は \(\xi^{2} + 3\xi + 2 = 0\) であり、したがって \(\mu = -2\), \(\nu = -1\) で、適切な代入は \( x = -(2t + 1)/(t + 1) \)。これは積分を次のように簡約する：
\[
- \int \frac{dt}{(4t^{2} + 1) \sqrt{9t^{2} - 4}}
  - \int \frac{t\, dt}{(4t^{2} + 1) \sqrt{9t^{2} - 4}}.
\]
これらの積分の最初のものは \( t/\sqrt{9t^{2} - 4} = u \) と置き換え、二つ目のものは \( 1/\sqrt{9t^{2} - 4} = v \) と置き換えることで有理化される。]

9. 
\[
\int \frac{(x + 1)\, dx}{(2x^{2} - 2x + 1) \sqrt{3x^{2} - 2x + 1}}, \quad
\int \frac{(x - 1)\, dx}{(2x^{2} - 6x + 5) \sqrt{7x^{2} - 22x + 19}}
\]
を計算せよ。[1911年数学トリップ]

\Item{10.} $y^{2} = ax^{2} + 2bx + c$ のとき、$\ds\int R(x, y)\, dx$ の有理化は、放物線 $y^{2} = ax^{2} + 2bx + c$ 上の任意の点 $(p, q)$ を用いて $t = (x - p)/(y + q)$ の置換によって行われる。もちろん、$t = (x - p)/(y - q)$ の置換でも有理化される：§134 参照。
\par\medskip

\Paragraph{140.超越関数。}超越関数の異なる種類が非常に多いことから、その積分理論は有理または代数的関数の積分理論ほど系統性がない。超越関数のいくつかのクラスについて、常に積分を求めることができるものを順に考察する。

\Paragraph{141.余弦と正弦の多倍角の多項式。}任意の実数 $a$, $b$ などに対し、正整数 $m$, $m'$, $n$, $n'$ などを用いて
\[
A\cos^{m} ax \sin^{m'} ax \cos^{n} bx \sin^{n'} bx\dots,
\]
のような有限項の和である任意の関数は積分できる。このような項は、いくつかの形の有限項の和として表すことができる：
\[
\alpha\cos\{(pa + qb + \dots)x\},\quad
\beta \sin\{(pa + qb + \dots)x\}
\]
これらの項の積分はすぐに書き出すことができる。

\par\medskip\noindent\textbf{【例 LI.】}\par
\Item{1.} $\sin^{3} x \cos^{2} 2x$ を積分せよ。この場合、以下の公式を使用する：
\[
\sin^{3}  x = \tfrac{1}{4}(3\sin x - \sin 3x),\quad
\cos^{2} 2x = \tfrac{1}{2}(1 + \cos 4x).
\]
これらの式を乗じて、例えば $\sin x\cos 4x$ を $\frac{1}{2}(\sin 5x - \sin 3x)$ に置き換えることで、
\begin{multline*}
\tfrac{1}{16}\int (7\sin x - 5\sin 3x + 3\sin 5x - \sin 7x)\, dx\\
  = - \tfrac{7}{16}\cos  x + \tfrac{5}{48}\cos 3x
    - \tfrac{3}{80}\cos 5x + \tfrac{1}{112}\cos 7x.
\end{multline*}

積分は当然異なる形で異なる方法によって得られます。例えば
\[
\int \sin^{3}x \cos^{2}2x\, dx = \int(4\cos^{4}x - 4\cos^{2}x + 1) (1 - \cos^{2} x)\sin x\, dx,
\]
この積分は $\cos x = t$ の置換により
\[
\int(4t^{6} - 8t^{4} + 5t^{2} - 1)\, dt = \frac{4}{7}\cos^{7} x - \frac{8}{5}\cos^{5}x + \frac{5}{3}\cos^{3}x - \cos x.
\]
この表現と先に得られたものは定数を除いて同一であることが確認できます。

2. 任意の方法で積分せよ $\cos ax \cos bx$, $\sin ax \sin bx$, $\cos ax \sin bx$,
$\cos^{2}x$, $\sin^{3}x$, $\cos^{4}x$, $\cos x \cos 2x \cos 3x$, $\cos^{3}2x \sin^{2}3x$, $\cos^{5}x \sin^{7}x$. [このような場合、減化公式（VI 39）を使用することが時々便利である有時。]

142. 積分 $\ds\int x^{n}\cos x\, dx$, $\ds\int x^{n}\sin x\, dx$ および関連する積分。部分積分法により、前記結果を一般化できます。すなわち
\begin{alignat*}{3}
\int x^{n}\cos x\, dx \&= \& \&x^{n}\sin x \&\&- n\int x^{n-1}\sin x\, dx,\\
\int x^{n}\sin x\, dx \&= \&-\&x^{n}\cos x \&\&+ n\int x^{n-1}\cos x\, dx,
\end{alignat*}
そして明らかに、$n$ が正の整数である限りこの過程を繰り返すことでこれらの積分は完全に計算できます。したがって $n$ が正の整数である限り $\ds\int x^{n}\cos ax\, dx$ および $\ds\int x^{n}\sin ax\, dx$ を常に計算することができます；そして、前記段落と同様の過程により
\[
\int P(x, \cos ax, \sin ax, \cos bx, \sin bx, \dots)\, dx,
\]
ここで $P$ は任意の多項式である。

247ページ

例題LII.
1. 積分せよ $x\sin x$, $x^{2}\cos x$, $x^{2}\cos^{2}x$, $x^{2}\sin^{2}x \sin^{2} 2x$,
$x\sin^{2}x \cos^{4}x$, $x^{3}\sin^{3}\frac{1}{3}x$。

2. 多項式 $P$ および $Q$ を求めよ
\[
\int\{(3x - 1)\cos x + (1 - 2x)\sin x\}\, dx = P\cos x + Q\sin x.
\]

\Item{3.} \( x^{n}\cos x \)の積分 \(\ds\int x^{n}\cos x\, dx\) は \( P_{n}\cos x + Q_{n}\sin x \) で表され、
\[
P_{n} = nx^{n-1} - n(n - 1)(n - 2) x^{n-3} + \dots,\quad
Q_{n} = x^{n} - n(n - 1) x^{n-2} + \dots.
\]
となる。

\Paragraph{143. \(\cos x\)と\(\sin x\)の有理関数．} \(\cos x\)および\(\sin x\)の任意の有理関数の積分は、置換 \(\tan \frac{1}{2}x = t\) によって計算できる。実際には
\[
\cos x = \frac{1 - t^{2}}{1 + t^{2}},\quad
\sin x = \frac{2t}{1 + t^{2}},\quad
\frac{dx}{dt} = \frac{2}{1 + t^{2}},
\]
であるから、置換は積分を \(t\) の有理関数の積分に減らす。

\par\medskip\noindent\textbf{【例 LIII.】}\par
\Item{1.} 以下の等式が成り立つことを証明せよ：
\[
\int \sec x\, dx = \log |\sec x + \tan x|,\quad
\int \cosec x\, dx = \log |\tan \tfrac{1}{2}x|.
\]

[最初の積分の別の形は \(\log |\tan(\frac{1}{4}\pi + \frac{1}{2}x)|\) であり、第三の形は \(\frac{1}{2}\log |(1 + \sin x)/(1 - \sin x)|\) である。]

\Item{2.} 以下の等式が成り立つことを証明せよ：
\[
\int \tan x\, dx = -\log |\cos x|,\quad
\int \cot x\, dx = \log |\sin x|,\quad
\int \sec^{2} x\, dx = \tan x,
\]
\[
\int \cosec^{2} x\, dx = -\cot x,\quad
\int \tan x\sec x\, dx = \sec x,\quad
\int \cot x \cosec x\, dx = -\cosec x.
\]

[これらの積分は一般形に含まれているが、置換を使用する必要はない。結果は \SecNo[§]{119} および \SecNo[§]{130} の式 (5) からすぐに導き出される。]

3. \(a + b\)が正のとき、\(1/(a + b\cos x)\)の積分は以下のいずれかの形で表すことができる：
\[
\frac{2}{\sqrt{a^{2} - b^{2}}}
  \arctan \left\{t\sqrt{\frac{a - b}{a + b}}\right\},\quad
\frac{1}{\sqrt{b^{2} - a^{2}}}
  \log \left|\frac{\sqrt{b + a} + t\sqrt{b - a}}
                  {\sqrt{b + a} - t\sqrt{b - a}}\right|,
\]
ここで \(t = \tan\frac{1}{2}x\) であり、\(a^{2} > b^{2}\) のとき最初の形を用い、\(a^{2} < b^{2}\) のとき第二の形を用いる。\(a^{2} = b^{2}\) の場合、積分は \(\sec^{2}\frac{1}{2}x\) または \(\cosec^{2}\frac{1}{2}x\) の定数倍となるので、その値はすぐに求めることができる。また、\(a + b\)が負のときの積分形も導き出せ。

4. \(y\)が\(x\)を用いて方程式
\[
(a + b\cos x)(a - b\cos y) = a^{2} - b^{2},
\]
によって定義されるとき、\(x\)が0から\(\pi\)まで変化するとき、ある値の\(y\)も0から\(\pi\)まで変化する。また
\[
\sin x = \frac{\sqrt{a^{2} - b^{2}} \sin y}{a - b\cos y},\quad
\frac{\sin x}{a + b\cos x}\, \frac{dx}{dy} = \frac{\sin y}{a - b\cos y}
\]
が成り立つ。したがって、\(0 < x < \pi\)のとき
\[
\int \frac{dx}{a + b\cos x}
  = \frac{1}{\sqrt{a^{2} - b^{2}}}
    \arccos \left(\frac{a\cos x + b}{a + b\cos x}\right)
\]
が成り立つ。この結果はEx.3と一致する。

5. \(1/(a + b\cos x + c\sin x)\)の積分をどのように行うかを示せ。[\(b\cos x + c\sin x\)を\(\sqrt{b^{2} + c^{2}} \cos(x - \alpha)\)と表す]

6. \((a + b\cos x + c\sin x)/(\alpha + \beta\cos x + \gamma\sin x)\)の積分を求めよ。[適当な\(\lambda\)、\(\mu\)、\(\nu\)を定めることで
\[
a + b\cos x + c\sin x
  = \lambda + \mu(\alpha + \beta\cos x + \gamma\sin x)
            + \nu(-\beta\sin x + \gamma\cos x)
\]
と表すことができる。すると積分は
\[
\mu x + \nu \log |\alpha + \beta\cos x + \gamma\sin x|
  + \lambda \int \frac{dx}{\alpha + \beta\cos x + \gamma\sin x}
\]となる。]

\Item{7.} \(1/(a\cos^{2} x + 2b\cos x\sin x + c\sin^{2} x)\) を積分せよ。[被積分関数は、$A + B\cos 2x + C\sin 2x$ の形で表すことができるここで $A = \frac{1}{2}(a + c)$, $B = \frac{1}{2}(a - c)$, $C = b$ であるが、$\tan x = t$ と置くことでより単純に積分できる。よって
\[
\int \frac{\sec^{2} x\, dx}{a + 2b\tan x + c\tan^{2} x}
  = \int \frac{dt}{a + 2bt + ct^{2}}.
\]
]

\Paragraph{144. $\arcsin x$, $\arctan x$, および $\log x$ を含む積分。} 反正弦と反正切、対数の積分は部分積分によって容易に計算できる。例えば
\begin{align*}
\int \arcsin x\, dx
  &= x\arcsin x - \int \frac{x\, dx}{\sqrt{1 - x^{2}}}
   = x\arcsin x + \sqrt{1 - x^{2}},\\
%
\int \arctan x\, dx
  &= x\arctan x - \int \frac{x\, dx}{1 + x^{2}}
   = x\arctan x - \frac{1}{2} \log(1 + x^{2}),\\
%
\int \log x\, dx
  &= x\log x - \int dx = x(\log x - 1).
\end{align*}

$y = f(x)$ の積分が求められるならば、$\phi(y)$ の積分も常に求めることができる。なぜなら $y = f(x)$ の置換を行うと
\[
\int \phi(y)\, dy = \int xf'(x)\, dx = xf(x) - \int f(x)\, dx.
\]
読者は $\arcsin y$ および $\arctan y$ の積分も同様にして評価すべきである。

積分の形として
\[
\int P(x, \arcsin x)\, dx,\quad
\int P(x, \log x)\, dx,
\]
ここで $P$ は多項式であるとき、常に計算することができます。まず最初の形を例にとってみましょう。$\ds\int x^{m} (\arcsin x)^{n}\, dx$ のような積分の数を計算しなければなりません。変数置換 $x = \sin y$ を用いることで、$\ds\int y^{n}\sin^{m}y \cos y\, dy$ に帰着させることができます。この積分は第142節で述べた方法により求められます。第二の形の場合には $\ds\int x^{m} (\log x)^{n}\, dx$ のような積分を計算しなければなりません。部分積分を行うと
\[
\int x^{m}(\log x)^{n}\, dx
  = \frac{x^{m+1} (\log x)^{n}}{m + 1}
  - \frac{n}{m + 1} \int x^{m}(\log x)^{n-1}\, dx,
\]
を得ます。この過程を十分繰り返せば、最終的に積分の完全な値に到達します。

145. 平面曲線の面積。積分法に関する前節で説明したプロセスの最も重要な応用例は、平面曲線の面積計算です。$P_{0}PP'$（図44）が連続な曲線 $y = \phi(x)$ のグラフであり、この曲線は$x$軸の上に完全に存在し、$P$ は点 $(x, y)$、$P'$ は点 $(x + h, y + k)$ で、$h$ は正または負（図では正）であるとします。

読者には「面積」の概念は WELL 熟知されていることと思います。特に、$ONPP_{0}$ のような面積についても同様です。この概念については現在、その前提として受け入れます。実際、このような面積を定義し、数学的に分析するための最も精密な方法が既に存在しています：後ほど、$ONPP_{0}$ などの空間領域に対して「面積」をどのように付与するかを詳細に説明します。現時点では、これらの領域には明確で正の数 $(ONPP_{0})$ が関連付けられ、そのような面積は直感によく示される性質を持つと仮定します。

したがって、$ONPP_{0}$ の面積は $x$ の関数であることが明らかです。これを $\Phi(x)$ で表記します。また、$\Phi(x)$ は連続な関数です。実際には、

\[
\Phi(x + h) - \Phi(x)
= (NN'P'P)
= (NN'RP) + (PRP')
= h\phi(x) + (PRP').
\]

図が描かれた状態では、$PRP'$ の面積は $hk$ よりも小さいことが明らかです。しかし、これは必ずしも一般に成り立つわけではありません。なぜなら、弧 $PP'$ が点 $P$ から $P'$ まで一貫して上昇または下降するとは限らないからです（例えば \Fig{44a} を参照）。しかし、$PR$ からの任意の点の最大距離 $\lambda(h)$ の下では常に $PRP'$ の面積は $|h|\lambda(h)$ よりも小さいことが保証されます。さらに、$\phi(x)$ が連続な関数であるため、$h \to 0$ のとき $\lambda(h) \to 0$ です。したがって、

\[
\Phi(x + h) - \Phi(x)
= h\{\phi(x) + \mu(h)\},
\]

ここで $|\mu(h)| < \lambda(h)$ かつ $\lambda(h) \to 0$ のとき、$\Phi(x)$ は連続であることが分かります。さらに、

\[
\Phi'(x)
= \lim_{h \to 0} \frac{\Phi(x + h) - \Phi(x)}{h}
= \lim_{h \to 0} \{\phi(x) + \mu(h)\}
= \phi(x).
\]

したがって、「曲線の縦座標は面積の導関数であり、面積は縦座標の積分である」と言えます。

Weは直線以外の曲線の長さという概念を定義するのに、実際にはより難しい問題である。具体的には、図44において$P_{0}P$が一定の長さを持つことを仮定し、その長さを$S(x)$と表すことは、面積に対する同様の仮定では足りない。$S(x)$が連続であることを証明することすら不可能であり、すなわち$\lim\{S(P') - S(P)\} = 0$であることも証明できない。大きな図においても明らかに見えるが、小さな図の場合にはそうではない。実際、曲線の長さという概念自体を精密に分析せずに、何らかの程度の厳密性でさらに進むことは不可能である。

ただし、曲線が接線の方向が連続的に変化するものと仮定し、\(\phi'(x)\)が連続であるとする。このとき、曲線の長さを仮定すると以下の式を得る：
\[
\frac{S(x + h) - S(x)}{h} = \frac{PP'}{h} = \left( \frac{PP'}{h} \right) \times \left( \frac{PP'}{PP'} \right),
\]
ここで、\(PP'\)は弦が\(PP'\)である弧の長さを表す。次に
\[
PP' + \sqrt{PR^2 + RP'^2} = h \sqrt{1 + \frac{k^2}{h^2}},
\]
また
\[
k = \phi(x + h) - \phi(x) = h \phi'(\xi),
\]
ここで、\(\xi\)は\(x\)と\(x + h\)の間にある。したがって
\[
\lim \left( \frac{PP'}{h} \right) = \lim \sqrt{1 + [\phi'(\xi)]^2} = \sqrt{1 + [\phi'(x)]^2}.
\]
さらに、
\[
\lim \left( \frac{PP'}{PP'} \right) = 1,
\]
を仮定すると、以下の結果を得る：
\[
S'(x) = \lim \left( \frac{S(x + h) - S(x)}{h} \right) = \sqrt{1 + [\phi'(x)]^2}
\]
よって
\[
S(x) = \int \sqrt{1 + [\phi'(x)]^2} \, dx.
\]

以下は例題である。

\Item{1.} \(y = x^2 / 4a\)の放物線と、\(x = \xi\)の正接によって切り取られた部分の面積およびその境界となる弧の長さを計算せよ。

\Item{2.} 曲線\(ay^2 = x^3\)について同じ問題を解き、その弧の長さは
\[
\frac{8a}{27} \left\{ \left( 1 + \frac{9\xi}{4a} \right)^{3/2} - 1 \right\}
\]
であることを示せ。

\Item{3.} 式\SecNo[§§]{145}--\SecNo{146}を用いて、円\(x^2 + y^2 = a^2\)および\(x^2 + y^2 = 2ax\)の面積と弧長を計算せよ。

\Item{4.} 楕円\((x^2 / a^2) + (y^2 / b^2) = 1\)の面積は\(\pi ab\)であることを示せ。

\item 5. 曲線 \(y = \sin x\) と \(x\) 軸の区間 \(0 \leq x \leq 2\pi\) の部分で囲まれた面積を求めよ。[ここで \(\Phi(x) = -\cos x\) であり、\(x = 0\) から \(x = 2\pi\) における \(-\cos x\) の値の差はゼロである。この説明は、\(x = \pi\) から \(x = 2\pi\) の間で曲線が \(x\) 軸の下側にあるため、その部分の面積が方法を適用する際に負として計算されることが当然であることを示している。\(x = 0\) から \(x = \pi\) の区間における面積は \(-\cos \pi + \cos 0 = 2\) であり、各部分を正として計算した全体の面積はこれの二倍、すなわち \(4\) である。]

\item 6. 曲線上の任意一点の座標がパラメータ \(t\) を用いて \(x = \phi(t)\), \(y = \psi(t)\) の形で表されるものとする。ここで \(\phi\) と \(\psi\) は連続的な導関数を持つ \(t\) の関数である。\(t\) が \(t_0\) から \(t_1\) に変化するとき \(x\) が常に増加すると仮定する。このとき、対応する部分の曲線、\(x\) 軸、および \(t_0\) と \(t_1\) に対応する二つの垂直線で囲まれた領域の面積は符号を除いて \(A(t_1) - A(t_0)\) である。ここで
\[
A(t) = \int \psi(t) \phi'(t) \, dt = \int y \frac{dx}{dt} \, dt.
\]

\item 7. 単一のループから成る閉曲線 \(C\) があり、どちらの軸に対して平行な直線が二点以上で交わることはないものとする。また、任意の点 \(P\) の座標は例題6のようにパラメータ \(t\) を用いて表され、\(t\) が \(t_0\) から \(t_1\) に変化するとき \(P\) が曲線を同じ方向で一周して元の位置に戻るものとする。このとき、ループの面積は任意の一つの積分
\[
-\int y \frac{dx}{dt} \, dt,\quad
\int x \frac{dy}{dt} \, dt,\quad
\tfrac{1}{2} \int \left(x \frac{dy}{dt} - y \frac{dx}{dt}\right) dt,
\]
の初期値と最終値の差に等しい。この差は当然正として取られる。

\item[8.] Ex.7の結果を用いて、以下の曲線の面積を求めよ：
\[
ZZZMATHTOKEN17ZZZ
\]

\item[9.] 曲線 \(x^{3} + y^{3} = 3axy\) のループの面積を求めよ。［\(y = tx\)と置くと、\(x = \frac{3at}{1 + t^{3}}\), \(y = \frac{3at^{2}}{1 + t^{3}}\) となる。\(t\)が0から無限大に向かうにつれてループは一度描かれる。また
\[
\frac{1}{2} \int \left(y \frac{dx}{dt} - x \frac{dy}{dt}\right) dt = -\frac{1}{2} \int x^{2} \frac{d}{dt}\left(\frac{y}{x}\right) dt = -\frac{1}{2} \int \frac{9a^{2}t^{2}}{(1 + t^{3})^{2}} dt = \frac{3a^{2}}{2(1 + t^{3})}
\]
であり、\(t \to \infty\) のとき0に収束する。したがってループの面積は\(\frac{3}{2}a^{2}\)である。］

\item[10.] 曲線 \(x^{5} + y^{5} = 5ax^{2}y^{2}\) のループの面積を求めよ。

\item[11.] 曲線 \(x = a\sin 2t\), \(y = a\sin t\) のループの面積は \(\frac{4}{3}a^{2}\) であることを証明せよ。［1908年］

\item[12.] 点 \(t = t_{1}\) から \(t = t_{2}\) までの楕円の弧 \(x = a\cos t\), \(y = b\sin t\) の長さは \(F(t_{2}) - F(t_{1})\) であり、ここで
\[
F(t) = a \int \sqrt{1 - e^{2}\sin^{2} t} dt,
\]
\(e\) は楕円の離心率である。［しかしこの積分は現在の我々の手元にあるような関数で評価することはできない］。

\item[13.] 極座標において、関数 \(f(\theta)\) を値とする \(\theta\) の一対一関数 \(r = f(\theta)\) および半径 \(\theta = \theta_{1}\), \(\theta = \theta_{2}\) で囲まれた面積は \(F(\theta_{2}) - F(\theta_{1})\) であり、ここで
\[
F(\theta) = \frac{1}{2} \int r^{2} d\theta.
\]
また、対応する弧の長さは \(\Phi(\theta_{2}) - \Phi(\theta_{1})\) であり、
\[
\Phi(\theta) = \int \sqrt{r^{2} + \left(\frac{dr}{d\theta}\right)^{2}} d\theta.
\]

\#\#\# (i) 圏の面積と周長を決定せよ：$r = 2a\sin\theta$
(ii) $r = \frac{1}{2}l\sec^{2}\frac{1}{2}\theta$ の放物線とその焦点直線および対応する放物線の弧の長さを決定せよ：
(iii) $r = a + b\cos\theta$ のリマッソンの面積を、$a > b$, $a = b$, および $a < b$ の場合に区分して求めよ；
(iv) 次の楕円の面積を求めよ：$\frac{1}{r^{2}} = a\cos^{2}\theta + 2h\cos\theta\sin\theta + b\sin^{2}\theta$ および $\frac{l}{r} = 1 + e\cos\theta$. 最後の場合は、積分 $\int \frac{d\theta}{(1 + e\cos\theta)^{2}}$ を計算する（例題 liii.4 参照）。この積分は置換 \[
(1 + e\cos\theta)(1 - e\cos\phi) = 1 - e^{2}
\] の助けを借りて計算できる。

\#\#\#\# 14. 放物線 $2\theta = \frac{a}{r} + \frac{r}{a}$ を描き、半径ベクトル $\theta = \beta$ と点 $r = a$, $\theta = 1$ で接する二つの枝によって囲まれた領域の面積を $\frac{2}{3}a^{2}(\beta^{2} - 1)^{\frac{3}{2}}$ であることを示せ。[1900年数学オリンピック問題]

\#\#\#\# 15. 方程式 $p = f(r)$ によって与えられた曲線において、$r$ は半径ベクトル、$p$ は原点から接線への垂線の長さとする。弧と二つの半径ベクトルで囲まれた領域の面積を計算するためには、積分 $\frac{1}{2} \int \frac{pr\, dr}{\sqrt{r^{2} - p^{2}}}$ の計算が必要であることを示せ。

\#\#\# 付録：第VI章の Miscellaneous 問題

\#\#\#\# 1. 関数 $f(x)$ を次のように定める：$x \leq 0$ のとき $f(x) = 1 + x$, $0 < x < 1$ のとき $f(x) = x$, $1 \leq x \leq 2$ のとき $f(x) = 2 - x$, $x > 2$ のとき $f(x) = 3x - x^{2}$. $x = 0$, $x = 1$, $x = 2$ における $f(x)$ の連続性、および $f'(x)$ の存在と連続性について議論せよ。[1908年数学オリンピック問題]

\#\#\#\# 2. $a$, $ax + b$, $ax^{2} + 2bx + c$ 等を $u_{0}$, $u_{1}$, $u_{2}$ 等と表すとき、$u_{0}^{2} u_{3} - 3u_{0} u_{1} u_{2} + 2u_{1}^{3}$ および $u_{0} u_{4} - 4u_{1} u_{3} + 3u_{2}^{2}$ は $x$ に依存しないことを示せ。

3. \(a_0, a_1, \dots, a_{2n}\)が定数であり、\(U_r = (a_0, a_1, \dots, a_r \btw x, 1)^r\)のとき、
\[
U_0 U_{2n} - 2n U_1 U_{2n-1}
  + \frac{2n(2n - 1)}{1·2} U_2 U_{2n-2} - \dots + U_{2n} U_0
\]
は\(x\)に依存しない。

[微分し、関係式 \(U_r' = r U_{r-1}\) を用いる。]

4. \(\mu > 1\)のとき、関数 \(\arcsin(\mu \sin x) - x\) の最初の3つの導函数は、\(0 \leq x \leq \frac{1}{2} \pi\) の範囲で正である。

5. 行列式の成分が \(x\) の関数である。その微分係数を、一行だけの成分を微分し他の成分を変更しないで形成される行列式の和として示せ。

6. \(f_1, f_2, f_3, f_4\) が次数4以下の多項式であるとき、
\[
\begin{vmatrix}
f_1 & f_2 & f_3 & f_4 \\
f_1' & f_2' & f_3' & f_4' \\
f_1'' & f_2'' & f_3'' & f_4'' \\
f_1''' & f_2''' & f_3''' & f_4'''
\end{vmatrix}
\]
も次数4以下の多項式である。[5回微分し、問題5の結果を用い、消滅する行列式を排除する。]

7. \(y^3 + 3yx + 2x^3 = 0\) のとき、\(x^2(1 + x^3)y'' - \frac{3}{2}xy' + y = 0\) が成り立つ。

[1903年]

8. \(y = \phi\{\psi(y_1)\} + \phi\{x - \psi(y_1)\}\) の微分方程式、ここで \(y_1\) は \(y\) の導函数、そして \(\psi\) は \(\phi'\) の逆関数であるとき、
\(y = \phi(c) + \phi(x - c)\) や \(y = 2\phi(\frac{1}{2}x)\) が満たすことを確認せよ。

9. 同じ記号を使用して、微分方程式 \(y = \frac{x}{\psi(y_1)} \phi\{\psi(y_1)\}\) が
\(y = c \phi(x/c)\) や \(y = \beta x\) を満たすことを確認せよ。ここで \(\beta = \frac{\phi(\alpha)}{\alpha}\) であり、\(\alpha\) は方程式 \(\phi(\alpha) - \alpha \phi'(\alpha) = 0\) の任意の根である。

\Item{10.} $ax + by + c = 0$ のとき、$y_{2} = 0$（接続符は$x$に関する微分を表す）。これを述べると、すべての直線の一般微分方程式は$y_{2} = 0$であると言える。次の各曲線の一般微分方程式を求めよ：(i) $x$軸上の中心を持つすべての円、(ii) $x$軸に沿った軸を持つすべての放物線、(iii) $y$軸に平行な軸を持つすべての放物線、(iv) すべての円、(v) すべての放物線、(vi) すべての双曲線。

[方程式は以下の通りである：(i) $1 + y_{1}^{2} + yy_{2} = 0$、(ii) $y_{1}^{2} + yy_{2} = 0$、(iii) $y_{3} = 0$、(iv) $(1 + y_{1}^{2}) y_{3} = 3y_{1} y_{2}^{2}$、(v) $5y_{3}^{2} = 3y_{2} y_{4}$、(vi) $9y_{2}^{2} y_{5} - 45y_{2} y_{3} y_{4} + 40y_{3}^{3} = 0$。各場合においては、所与の曲線の方程式を書き下し、十分な数の微分方程式を得るまで微分を行うことで、任意定数を消去することができる。]

\Item{11.} すべての放物線と双曲線の一般微分方程式はそれぞれ
\[
D_{x}^{2} (y_{2}^{-2/3}) = 0,\quad
D_{x}^{3} (y_{2}^{-2/3}) = 0.
\]
を示せ。

[双曲線の方程式は以下のように表すことができる：
\[
y = ax + b ± \sqrt{px^{2} + 2qx + r}.
\]
これより
\[
y_{2} = ±(pr - q^{2})/(px^{2} + 2qx + r)^{3/2}.
\]
双曲線が放物線であるとき、$p = 0$ である。]

\Item{12.} $\dfrac{dy}{dx}$, $\dfrac{1}{2!}\, \dfrac{d^{2}y}{dx^{2}}$, $\dfrac{1}{3!}\, \dfrac{d^{3}y}{dx^{3}}$, $\dfrac{1}{4!}\, \dfrac{d^{4}y}{dx^{4}}$ を $t$, $a$, $b$, $c$ とし、$\dfrac{dx}{dy}$, $\dfrac{1}{2!}\, \dfrac{d^{2}x}{dy^{2}}$, $\dfrac{1}{3!}\, \dfrac{d^{3}x}{dy^{3}}$, $\dfrac{1}{4!}\, \dfrac{d^{4}x}{dy^{4}}$ を $\tau$, $\alpha$, $\beta$, $\gamma$ と表すとき、
\[
4ac - 5b^{2} = (4\alpha\gamma - 5\beta^{2})/\tau^{8},\quad
bt - a^{2} = - (\beta\tau - \alpha^{2})/\tau^{6}.
\]
同様の公式を関数 $a^{2}d - 3abc - 2b^{3}$, $(1 + t^{2})b - 2a^{2}t$, $2ct - 5ab$ についても確立せよ。

13. \( y_k \)が \( y = \sin(n \arcsin x) \) の \( k \) 番目の導関数であるとき、次の式が成り立つことを証明せよ：
\[
(1 - x^2)y_{k+2} - (2k + 1)xy_{k+1} + (n^2 - k^2)y_k = 0.
\]
まず \( k = 0 \) の場合を証明し、その後 Leibniz の定理を使って \( k \) 回微分すること。

14. 次の公式を証明せよ：
\[
vD_x^n u = D_x^n (uv) - nD_x^{n-1} (uD_x v) + \frac{n(n - 1)}{1 \cdot 2} D_x^{n-2} (uD_x^2 v) - \cdots
\]
ここで \( n \) は任意の正整数である。[数学的帰納法を使用する]

15. 曲線が与えられており、
\[
x = a(2 \cos t + \cos 2t), \quad y = a(2 \sin t - \sin 2t).
\]
点 \( P \) のパラメータが \( t \) のとき、接線と法線の方程式を証明せよ：
\[
x \sin \frac{1}{2} t + y \cos \frac{1}{2} t = a \sin \frac{3}{2} t, \quad x \cos \frac{1}{2} t - y \sin \frac{1}{2} t = 3a \cos \frac{3}{2} t;
\]
点 \( P \) の接線が曲線と交わる点 \( Q \), \( R \) のパラメータが \( -\frac{1}{2} t \) および \( \pi - \frac{1}{2} t \) のとき、\( QR = 4a \) を証明せよ；
接線 \( Q \) と接線 \( R \) は直角で、これらが交わる点は円 \( x^2 + y^2 = a^2 \) 上にあることを証明せよ；
点 \( P \), \( Q \), \( R \) の法線は共有点を持ち、その交点は円 \( x^2 + y^2 = 9a^2 \) 上にあることを証明せよ；
曲線の方程式を証明せよ：
\[
(x^2 + y^2 + 12ax + 9a^2)^2 = 4a(2x + 3a)^3.
\]
この曲線の形状を図示せよ。

16. 計算問題 15 の定義する曲線の方程式は、次の式で置き換えることができる：
\[
\xi/a = 2u + (1/u^2), \quad \eta/a = (2/u) + u^2,
\]
ここで \( \xi = x + yi, \eta = x - yi, u = \Cis t \) である。点 \( u \) で定義される接線と法線の方程式を証明し、計算問題 15 の (ii) から (v) の性質を導き出せ。

17. 次の方程式が等しい根を持つ条件は次の形で表すことができる：
\[
(p + q)^{2/3} - (p - q)^{2/3} = 1.
\]
この問題は 1898 年に提出された。

18. 三次方程式 \(f(x) = 0\) の根は \(\alpha\), \(\beta\), \(\gamma\) で、これらは大きさの昇順に並んでいる。もし \((\alpha, \beta)\) と \([\alpha, \beta]\) の間隔を6等分し、また \((\beta, \gamma)\) と \([\beta, \gamma]\) の間隔も同様に6等分すると、それぞれの4番目の区間に \(f'(x) = 0\) の根が存在することを示せ。\(f'(x) = 0\) の根が分割点で存在する場合、その性質はどのようなものか。

19. \(f(x)\) の極大値と極小値、および \(f(x) = 0\) の実数根を調査せよ。ここで \(f(x)\) は以下の関数のいずれかである：
\[
x - \sin x - \tan\alpha (1 - \cos x),\quad
x - \sin x - (\alpha - \sin\alpha) - \tan \tfrac{1}{2}\alpha (\cos\alpha - \cos x),
\]
ここで \(\alpha\) は \(0 < \alpha < \pi\) の角である。最初の関数の場合、重根となる条件は \(\tan\alpha - \alpha\) が \(\pi\) の倍数であること。

20. 比例 \(\lambda : \mu\) を選ぶことで、方程式 \(\lambda(ax^{2} + bx + c) + \mu(a'x^{2} + b'x + c') = 0\) の根を実数として任意の大きさに差が存在させることができる。ただし、二つの二次方程式の根がすべて実数で交差している場合を除き、その例外的な場合においても根は常に実数であり、それらの差には下限がある。

21. \(0 < x < 1\) のとき、次の不等式が成り立つことを証明せよ：
\[
\pi < \frac{\sin \pi x}{x(1 - x)} \leq 4
\]
さらに、この関数のグラフを描け。

22. 関数のグラフを描け：
\[
\pi \cot\pi x - \frac{1}{x} - \frac{1}{x - 1}.
\]

23. \(y\) の一般形のグラフを素線で描け。ただし、
\[
\frac{dy}{dx} = \frac{(6x^{2} + x - 1) (x - 1)^{2} (x + 1)^{3}}{x^{2}}
\]
である。

24. 紙を折りたたむと、一つの角が対面する端に達するようにする。紙をどのように折るかを示せ、そのときの折れ線の長さが最大となるように。

\item[25.] 椭圆 \( (x^2/a^2) + (y^2/b^2) = 1 \) が中心と共通する円によって切り取られる最大の鋭角は \( \arctan\{(a^2 - b^2)/2ab\} \) である。

\item[26.] 三角形において面積 \( \Delta \) と半周長 \( s \) が固定されているとき、一方の辺の最大値または最小値は方程式
\[ s(x - s)x^2 + 4\Delta^2 = 0 \]
の根であることを示せ。この方程式の根の実数性と、その根が最大値または最小値に対応するかを議論せよ。

[方程式 \( a + b + c = 2s \), \( s(s - a)(s - b)(s - c) = \Delta^2 \) は \( a \) と \( b \) を \( c \) の関数として定める。\( c \) について微分し、\( da/dc = 0 \) とする。すると \( b = c \), \( s - b = s - c = \frac{1}{2}a \) が得られ、これより
\[ s(a - s)a^2 + 4\Delta^2 = 0 \]
を得る。

この方程式は \( s^4 > 27\Delta^2 \) のとき3つの実数根を持つし、その他の場合1つの実数根を持つ。等辺三角形（与えられた面積に対する最小の周長を持つ三角形）において \( s^4 = 27\Delta^2 \); よって \( s^4 < 27\Delta^2 \) は不可能である。したがって方程式は3つの実数根を持つことになり、その和が正であり積が負であることから、2つの根が正で1つの根が負である。

このうちの2つの正の根は最大値に対応し、残りの1つは最小値に対応する。]

\item[27.] 3点 \( A \), \( B \), \( C \) を通る最大の等辺三角形の面積は
\[ 2\Delta + \frac{a^2 + b^2 + c^2}{2\sqrt{3}} \]
であり、\( a \), \( b \), \( c \) は辺長で \( \Delta \) は \( ABC \) の面積である。

\item[28.] 原点に頂点を置き、底角が心形線 \( r = a(1 + \cos\theta) \) 上にある2つの最大等脚三角形の面積を \( \Delta \), \( \Delta' \) とするとき、
\[ 256\Delta\Delta' = 25a^4\sqrt{5} \]
が成り立つ。

29. 点 $(x, y)$ が曲線 $x^{2}y - 4x^{2} - 4xy + y^{2} + 16x - 2y - 7 = 0$ 上で点 $(2, 3)$ に近づくとき、比 $(x^{2} - 4y + 8)/(y^{2} - 6x + 3)$ の極限値を求めよ。1903年。

もし点 $(2, 3)$ を新たな原点とすると、曲線の方程式は $\xi^{2} \eta - \xi^{2} + \eta^{2} = 0$ となり、与えられた関数は $(\xi^{2} + 4\xi - 4\eta)/(\eta^{2} + 6\eta - 6\xi)$ となる。$\eta = t\xi$ と置くと $\xi = (1 - t^{2})/t$, $\eta = 1 - t^{2}$ となる。曲線は原点を中心に分岐するループがあり、これは $t = -1$ および $t = 1$ の二つの値に対応する。与えられた関数を $t$ で表し、$t$ を $-1$ または $1$ に近づけると、極限値 $-\frac{3}{2}$, $-\frac{2}{3}$ が得られる。

30. 
$f(x) = \dfrac{1}{\sin x - \sin a} - \dfrac{1}{(x - a)\cos a}$ のとき、
\[
\frac{d}{da}\{\lim_{x \to a} f(x)\} - \lim_{x \to a}f'(x)
  = \tfrac{3}{4} \sec^{3} a - \tfrac{5}{12} \sec a.
\]
1896年。

31. $\phi(x) = 1/(1 + x^{2})$ のとき、$\phi^{n}(x) = Q_{n}(x)/(1 + x^{2})^{n+1}$ であり、$Q_{n}(x)$ は次数 $n$ の多項式であることを示せ。また、

(i) $Q_{n+1} = (1 + x^{2}) Q_{n}' - 2(n + 1) x Q_{n}$,

(ii) $Q_{n+2} + 2(n + 2) x Q_{n+1} + (n + 2)(n + 1)(1 + x^{2})Q_{n} = 0$,

(iii) $(1 + x^{2}) Q_{n}'' - 2nx Q_{n}' + n(n + 1)Q_{n} = 0$,

(iv) $Q_{n} = (-1)^{n} n!\left\{(n + 1)x^{n} - \dfrac{(n + 1)n(n - 1)}{3!} x^{n-2} + \dots\right\}$,

(v) $Q_{n} = 0$ のすべての根は実数であり、$Q_{n-1} = 0$ の根によって区切られていることを示せ。

32. $a \leq x \leq b$ の範囲で $f(x)$, $\phi(x)$, $\psi(x)$ が微分可能であるとき、$a$ と $b$ の間の値 $\xi$ が存在し、
\[
\begin{vmatrix}
f(a)   & \phi(a)   & \psi(a)\\
f(b)   & \phi(b)   & \psi(b)\\
f'(\xi)& \phi'(\xi)& \psi'(\xi)
\end{vmatrix}
=0.
\]

[第三行の要素を f(x), φ(x), ψ(x) に置き換えた関数を考える。φ(x) = x, ψ(x) = 1 の場合、この定理は平均値の定理（第125節）に帰着される.]

33. Ex. 32から公式を導く：
\[
\frac{f(b) - f(a)}{\phi(b) - \phi(a)} = \frac{f'(\xi)}{\phi'(\xi)}
\]

34. φ'(x) が x → ∞ のとき a に収束するならば、φ(x)/x は a に収束する。φ'(x) が x → ∞ のとき無限大に収束するならば、φ(x) も無限大に収束する。[公式 φ(x) - φ(x₀) = (x - x₀)φ'(\xi) を用いる。ここで x₀ < ξ < x である.]

35. φ(x) が x → ∞ のとき a に収束するならば、φ'(x) は零以外の限界に向かって収束することはできない。

36. φ(x) + φ'(x) が x → ∞ のとき a に収束するならば、φ(x) は a に収束し、φ'(x) は 0 に収束する。

φ(x) = a + ψ(x) とおくので、ψ(x) + ψ'(x) は 0 に収束する。ψ'(x) が十分大きな x の値に対して一定の符号を持つ（例えば正）ならば、ψ(x) は安定的に増加し、限界 l または無限大へと収束する。ψ(x) → ∞ の場合、ψ'(x) → -∞ となり、これは仮定に矛盾する。ψ(x) → l の場合、ψ'(x) → -l であり、これは Ex. 35 によれば l = 0 であることが不可能である。同様にして ψ'(x) が最終的に負の値を持つ場合も処理できる。ψ(x) が x がすべての限界を超越する値に対して符号を変えるならば、これらは ψ(x) の極大と極小である。x が ψ(x) の極大または極小に対応する大きな値を持つ場合、ψ(x) + ψ'(x) は小さく、ψ'(x) = 0 となり、従って ψ(x) も小さい。より強力な言葉で言えば、x が大きいときの ψ(x) の他の値も小さい。

この定理の一般化や異なる証明方法については、著者の論文「Mr Mercer の極限定理の一般化について」（Quarterly Journal of Mathematics 第43巻）を参照されたい。上記の単純な証明は Prof. E. W. Hobson による提案に基づくものである。

\#\#\# 37. x, \(\sqrt{\frac{ax + b}{mx + n}}, \sqrt{\frac{cx + d}{mx + n}}\) の積分を有理関数の積分に帰着する方法を示せ。[置換 \(mx + n = 1/t\) を用いて \Ex{xlix}.13 を参照]

\#\#\# 38. 次の積分を計算せよ：
\[
\int \frac{dx}{(1 + x^{2})^{3}},\quad
\int \sqrt{\frac{x - 1}{x + 1}} \cdot \frac{dx}{x},\quad
\int \frac{x\, dx}{\sqrt{1 + x} - \sqrt{1 + x}},\displaybreak
\]
\[
\int \sqrt{a^{2} + \sqrt{b^{2} + \frac{c}{x}}}\, dx,\quad
\int \csc^{3}x\, dx,\quad
\int \frac{5\cos x + 6}{2\cos x + \sin x + 3}\, dx,\displaybreak
\]
\[
\int \frac{dx}{(2 - \sin^{2}x) (2 + \sin x - \sin^{2} x)},\quad
\int \frac{\cos x\sin x \, dx}{\cos^{4}x + \sin^{4}x},\quad
\int \csc x \sqrt{\sec 2x}\, dx,\displaybreak
\]
\[
\int \frac{dx}{\sqrt{(1 + \sin x) (2 + \sin x)}},\quad
\int \frac{x + \sin x}{1 + \cos x}\, dx,\quad
\int \sec^{-1}x\, dx,\quad
\int (\arcsin x)^{2}\, dx,\displaybreak
\]
\[
\int x\arcsin x\, dx,\quad
\int \frac{x\arcsin x}{\sqrt{1 - x^{2}}}\, dx,\quad
\int \frac{\arcsin x}{x^{3}}\, dx,\quad
\int \frac{\arcsin x}{(1 + x)^{2}}\, dx,\displaybreak
\]
\[
\int \frac{\arctan x}{x^{2}}\, dx,\quad
\int \frac{\arctan x}{(1 + x^{2})^{3/2}}\, dx,\quad
\int \frac{\log(\alpha^{2} + \beta^{2}x^{2})}{x^{2}}\, dx,\quad
\int \frac{\log(\alpha + \beta x)}{(a + bx)^{2}}\, dx.
\]

\#\#\# 39. 約化公式。 (i) \(x^{2} + px + q\) の \(n\) 乗の逆数の積分について
\[
2(n - 1)(q - \frac{1}{4}p^{2}) \int \frac{dx}{(x^{2} + px + q)^{n}} =
\frac{x + \frac{1}{2}p}{(x^{2} + px + q)^{n-1}}
+ (2n - 3) \int \frac{dx}{(x^{2} + px + q)^{n-1}}.
\]

\[ x + \frac{1}{2}p = t, \quad q - \frac{1}{4}p^{2} = \lambda \] と置くと、次のように表されます：
\[
\int \frac{dt}{(t^{2} + \lambda)^{n}}
  = \frac{1}{\lambda} \int \frac{dt}{(t^{2} + \lambda)^{n-1}}
   - \frac{1}{\lambda} \int \frac{t^{2}\, dt}{(t^{2} + \lambda)^{n}} \\
  = \frac{1}{\lambda} \int \frac{dt}{(t^{2} + \lambda)^{n-1}}
   + \frac{1}{2\lambda(n-1)}
     \int t \frac{d}{dt} \left\{\frac{1}{(t^{2} + \lambda)^{n-1}}\right\} dt,
\]
部分積分により、結果が得られます。

このような公式は、簡約公式と呼ばれ、\( n \) が正の整数であるとき特に有用です。このとき、\(\int \frac{dx}{(x^{2} + px + q)^{n}}\) を \(\int \frac{dx}{(x^{2} + px + q)^{n-1}}\) の形で表すことができ、任意の \( n \) に対して積分を評価することができます。

(ii) \( I_{p, q} = \int x^{p}(1 + x)^{q}\, dx \) とすると
\[
(p + 1) I_{p, q} = x^{p+1}(1 + x)^{q} - qI_{p+1, q-1},
\]
また、同様の公式 \( I_{p, q} \) と \( I_{p-1, q+1} \) の関係を導きます。さらに、置換 \( x = -y/(1 + y) \) を用いて
\[
I_{p, q} = (-1)^{p+1} \int y^{p} (1 + y)^{-p-q-2}\, dy.
\]

(iii) \( X = a + bx \) とすると
\[
\int xX^{-1/3}\, dx = -\frac{3(3a - 2bx) X^{2/3}}{10b^{2}},
\]
\[
\int x^{2}X^{-1/3}\, dx = \frac{3(9a^{2} - 6abx + 5b^{2}x^{2}) X^{2/3}}{40b^{3}},
\]
\[
\int xX^{-1/4}\, dx = -\frac{4(4a - 3bx) X^{3/4}}{21b^{2}},
\]
\[
\int x^{2}X^{-1/4}\, dx = \frac{4(32a^{2} - 24abx + 21b^{2}x^{2}) X^{3/4}}{231b^{3}}.
\]

(iv) \( I_{m, n} = \int \frac{x^{m}\, dx}{(1 + x^{2})^{n}} \) とすると
\[
2(n - 1)I_{m, n} = -x^{m-1} (1 + x^{2})^{-(n-1)} + (m - 1)I_{m-2, n-1}.
\]

(v) \( I_{n} = \int x^{n} \cos\beta x\, dx \) と \( J_{n} = \int x^{n} \sin\beta x\, dx \) とすると
\[
\beta I_{n} = x^{n} \sin\beta x - nJ_{n-1}, \quad
\beta J_{n} = -x^{n} \cos\beta x + nI_{n-1}.
\]

\Itemp{(vi)} $I_{n} = \int \cos^{n} x\, dx$ および $J_{n} = \int \sin^{n} x\, dx$ のとき、
\[
nI_{n} = \sin x\cos^{n-1} x + (n - 1) I_{n-2},\quad
nJ_{n} = -\cos x\sin^{n-1} x + (n - 1) J_{n-2}.
\]

\Itemp{(vii)} $I_{n} = \int \tan^{n}x\, dx$ のとき、$(n - 1)(I_{n} + I_{n-2}) = \tan^{n-1}x$。

\Itemp{(viii)} $I_{m, n} = \int \cos^{m}x \sin^{n}x\, dx$ のとき、
\begin{alignat*}{2}
(m+n)I\\textit{{m, n}
  \&= -\&\&\cos^{m+1}x \sin^{n-1}x + (n - 1) I}{m, n-2}\\
  \&=  \&\&\cos^{m-1}x \sin^{n+1}x + (m - 1) I\_{m-2, n}.
\end{alignat*}

[これらは
\begin{align*}
(m+1)I_{m, n}
  &= -\int \sin^{n-1}x \frac{d}{dx} (\cos^{m+1}x)\, dx\\
  &= -\cos^{m+1}x \sin^{n-1}x + (n - 1)\int \cos^{m+2}x \sin^{n-2}x\, dx\\
  &= -\cos^{m+1}x \sin^{n-1}x + (n - 1)(I_{m, n-2} - I_{m, n}),
\end{align*}
から導かれる第1の簡約公式である。]

\Itemp{(ix)} $I_{m, n} = \int \sin^{m}x \sin nx\, dx$ を $I_{m-2, n}$ と結びつける。

\Itemp{(x)} $I_{m, n} = \int x^{m} \cosec^{n}x\, dx$ のとき、
\begin{multline*}
(n - 1)(n - 2)I_{m, n} = (n - 2)^{2}I_{m, n-2} + m(m - 1)I_{m-2, n-2}\\
  -x^{m-1} \cosec^{n-1}x \{m\sin x + (n - 2) x\cos x\}.
\end{multline*}

\Itemp{(xi)} $I_{n} = \int (a + b\cos x)^{-n}\, dx$ のとき、
\[
(n - 1)(a^{2} - b^{2}) I_{n}
  = -b\sin x (a + b\cos x)^{-(n-1)} + (2n - 3)aI_{n-1} - (n - 2)I_{n-2}.
\]

\Itemp{(xii)} $I_{n} = \int (a\cos^{2} x + 2h\cos x\sin x + b\sin^{2}x)^{-n}\, dx$ のとき、
\[
4n(n + 1)(ab - h^{2})I_{n+2} - 2n(2n + 1)(a + b)I_{n+1} + 4n^{2}I_{n}
  = -\frac{d^{2} I_{n}}{dx^{2}}.
\]

\Itemp{(xiii)} $I_{m, n} = \int x^{m}(\log x)^{n}\, dx$ のとき、
\[
(m + 1)I_{m, n} = x^{m+1}(\log x)^{n} - nI_{m, n-1}.
\]

\Item{40.} $n$ が正の整数であるとき、$\int x^{m}(\log x)^{n}\, dx$ の値は
\[
x^{m+1} \left\{\frac{(\log x)^{n}}{m + 1}
  - \frac{n(\log x)^{n-1}}{(m + 1)^{2}}
  + \frac{n(n - 1)(\log x)^{n-2}}{(m + 1)^{3}} - \dots
  + \frac{(-1)^{n}n!}{(m + 1)^{n+1}}\right\}.
\]

41. 任意の \( x \) の値に対して \(\phi'' + a^2\phi = 0\) を満たす最一般の関数 \(\phi(x)\) は、\(A\cos ax + B\sin ax\) や \(\rho\cos(ax + \epsilon)\) のいずれかの形で表されることが示される。ここで \( A \), \( B \), \(\rho\), \(\epsilon\) は定数である。

42. \(\omega\) を定数として、\( y' + \omega z = 0 \) および \( z' - \omega y = 0 \) を満たす最一般の関数 \( y \) と \( z \) を決定せよ。ここでダッシュは \( x \) に関する微分を表す。

43. 次の式で与えられる曲線の面積は、\(\alpha\) が正の鋭角であるとき \(\frac{1}{2}\pi(1 + \sin\alpha)^2 / \sin\alpha\) である。
\[ x = \cos\phi + \frac{\sin\alpha \sin\phi}{1 - \cos^2\alpha \sin^2\phi}, \quad y = \sin\phi - \frac{\sin\alpha \cos\phi}{1 - \cos^2\alpha \sin^2\phi} \]
（1904年）

44. 半径 \( a \) の円の弦の投影が直径に恒常的な長さ \( 2a\cos\beta \) を持つとき、その弦の中点の軌跡は二つのループを構成し、それぞれの面積は \( a^2(\beta - \cos\beta\sin\beta) \) であることを示せ。
（1903年）

45. 次の曲線 \((x/a)^{2/3} + (y/b)^{2/3} = 1\) の四分円の長さは、\((a^2 + ab + b^2) / (a + b)\) であることを示せ。

46. 半径 \( a \) の円内に点 \( A \) があり、中心から距離 \( b \) であるとき、\( A \) から引かれた接線に対する垂線の足の軌跡が囲む面積は \(\pi(a^2 + \frac{1}{2}b^2)\) であることを示せ。
（1909年）

47. 掛け算が \( (a, b, c, f, g, h \btw x, y, 1)^{2} = 0 \) である楕円の方程式があるとき、
\[
\int \frac{dx}{(lx + my + n)(hx + by + f)} = \alpha \log \frac{PT}{PT'} + \beta,
\]
ここで \( PT, \ PT' \) は点 \( P \) の垂線であり、この点が楕円上の接線の端点に引かれた chord \( lx + my + n = 0 \) 上にある。また、\( \alpha, \ \beta \) は定数である。
1902年。

48. 次の積分
\[
\int \frac{ax^{2} + 2bx + c}{(Ax^{2} + 2Bx + C)^{2}}\, dx
\]
が \( x \) の有理関数であるためにはかつ只有に、\( AC - B^{2} \) または \( aC + cA - 2bB \) のいずれかがゼロであることが必要である。
（参照：著者の引用文献を第236ページで見よ。）

49. 多項式 \( f(x) \) と \( F(x) \) を用いて、後者が重複した因子を持たないとき、
\[
\int \frac{f(x)}{\{F(x)\}^{2}}\, dx,
\]
が \( x \) の有理関数であるための必要十分条件は、\( f'F' - fF'' \) が \( F \) で割り切れることがである。
1910年。

50. 次の積分
\[
\int \frac{\alpha \cos x + \beta \sin x + \gamma}{(1 - e \cos x)^{2}}\, dx
\]
が \( \cos x \) および \( \sin x \) の有理関数であるためにはかつ只有に、\( \alpha e + \gamma = 0 \) であり、この条件を満たすときの積分を求めよ。
1910年。

\chapter{第7章：Chapter VII: ADDITIONAL THEOREMS IN THE DIFFERENTIAL AND INTEGRAL CALCULUS}

第7章
微分積分における追加の定理

147. 高次の平均値の定理。前章（§125）で証明したように、$f(x)$が区間$\DPmod{(a, b)}{[a, b]}$において導関数$f'(x)$を持つとき
\[
f(b) - f(a) = (b - a) f'(\xi),
\]
ここで$a < \xi < b$；または、$f(x)$が区間$\DPmod{(a, a + h)}{[a, a + h]}$において導関数を持つとき
\[
f(a + h) - f(a) = hf'(a + \theta_{1} h),
\Tag{(1)}
\]
ここで$0 < \theta_{1} < 1$。これは、関数$f(b) - f(x) - \frac{b - x}{b - a} \{f(b) - f(a)\}$が$x = a$および$x = b$で消えることから証明されました。

今後、$f(x)$が区間$\DPmod{(a, b)}{[a, b]}$において二階導関数$f''(x)$を持つと仮定し、一階導関数$f'(x)$の連続性を含むこの仮定を考えます。そして次の関数を考えます。
\[
f(b) - f(x) - (b - x) f'(x)
  - \left(\frac{b - x}{b - a}\right)^{2} \{f(b) - f(a) - (b - a)f'(a)\}.
\]
この関数も$x = a$および$x = b$で消え、その導関数は
\[
\frac{2(b - x)}{(b - a)^{2}}
  \{f(b) - f(a) - (b - a) f'(a) - \tfrac{1}{2}(b - a)^{2}f''(x)\}
\]
であり、これは（§121）$a$と$b$の間（$a$および$b$を除く）の某$x$で0となる必要があります。したがって、ある$x$の値$\xi$があり、$a$と$b$の間にあり、$a + \theta_{2}(b - a)$という形で表現できるため
\[
f(b) = f(a) + (b - a)f'(a) + \tfrac{1}{2}(b - a)^{2}f''(\xi).
\]

$b = a + h$と置き換えると、次の式が得られます。
\[
f(a + h) = f(a) + hf'(a) + \tfrac{1}{2}h^{2} f''(a + \theta_{2}h),
\Tag{(2)}
\]
これは二階平均値の定理と呼ばれる標準形です。

式（1）と式（2）から提示される類似性は、次の定理を導きます。

以下はTaylorの定理または一般の平均値の定理です。

もし$f(x)$が$x$に関する関数で、区間$\DPmod{(a, b)}{[a, b]}$において最初の$n$階の導関数を持つとすると、

\begin{multline*}
f(b) = f(a) + (b - a)f'(a) + \frac{(b - a)^2}{2!} f''(a) + \dots\\
  + \frac{(b - a)^{n-1}}{(n - 1)!} f^{(n-1)}(a)
  + \frac{(b - a)^n}{n!}f^{(n)}(\xi),
\end{multline*}

ここで$a < \xi < b$；そして$b = a + h$のとき、

\begin{multline*}
f(a + h) = f(a) + hf'(a) + \frac{1}{2} h^2 f''(a) + \dots\\
  + \frac{h^{n-1}}{(n - 1)!} f^{(n-1)}(a)
  + \frac{h^n}{n!} f^{(n)}(a + \theta_n h),
\end{multline*}

ここで$0 < \theta_n < 1$。

この証明は、$n = 1$や$n = 2$の特別な場合で採用された方法と同様です。関数

\[
F_n(x) - \left(\frac{b - x}{b - a}\right)^n F_n(a),
\]

を考えます。ここで、

\begin{multline*}
F_n(x) = f(b) - f(x) - (b - x)f'(x) - \frac{(b - x)^2}{2!} f''(x) - \dots\\
  - \frac{(b - x)^{n-1}}{(n - 1)!} f^{(n-1)}(x).
\end{multline*}

この関数は$x = a$と$x = b$で消え、その導関数は

\[
\frac{n(b - x)^{n-1}}{(b - a)^n}
  \left\{F_n(a) - \frac{(b - a)^n}{n!} f^{(n)}(x)\right\};
\]

であり、$a$と$b$の間には導関数が消える$x$の値があります。これにより、望む結果が得られます。

この定理の重要性を考慮し、本章末尾で異なる形を取るが本質的に同じ証明を与えることにします。

\par\medskip\noindent\textbf{【例 LV.】}\par
\Item{1.} $f(x)$が次数$r$の多項式であると仮定すると、$f^{(n)}(x)$は$n > r$のとき恒等的にゼロであり、この定理は代数的な恒等式

\[
f(a + h) = f(a) + hf'(a) + \frac{h^2}{2!} f''(a) + \dots
  + \frac{h^r}{r!} f^{(r)}(a)
\]

を導きます。
\]

\item2. \( f(x) = 1/x \)に適用し、\( x \)と\( x + h \)が正であると仮定すると、結果は以下のようになります：
\[
\frac{1}{x + h} = \frac{1}{x} - \frac{h}{x^{2}} + \frac{h^{2}}{x^{3}} - \dots
  + \frac{(-1)^{n-1} h^{n-1}}{x^{n}}
  + \frac{(-1)^{n} h^{n}}{(x + \theta_{n} h)^{n+1}}.
\]

[この結果は以下の式から確認できます：
\[
\frac{1}{x + h} = \frac{1}{x} - \frac{h}{x^{2}} + \frac{h^{2}}{x^{3}} - \dots
  + \frac{(-1)^{n-1} h^{n-1}}{x^{n}}
  + \frac{(-1)^{n} h^{n}}{x^{n}(x + h)},
\]
したがって、\( x^{n}(x + h) \)を \( (x + \theta_{n}h)^{n+1} \)の形に変形できること、または \( x^{n+1} < x^{n}(x + h) < (x + h)^{n+1} \)であること、これは明らかです。]

\item3. 以下の公式を得ます：
\begin{multline*}
\sin(x + h)
  = \sin x + h\cos x - \frac{h^{2}}{2!}\sin x - h^{3}\frac{3!}{\cos x} + \dots\\
  + (-1)^{n-1}\frac{h^{2n-1}}{(2n - 1)!}\cos x
  + (-1)^{n} h^{2n}\frac{2n!}{\sin(x + \theta_{2n} h)},
\end{multline*}
同様の公式を得ます：
\[
\cos(x + h),
\]
また、\( h \)の冪を含む類似な公式を得ます。

\item4. \( m \)が正の整数、および \( n \)が \( m \)を超えない正の整数であるとき、以下の式が成り立ちます：
\[
(x + h)^{m} = x^{m} + \binom{m}{1}x^{m-1} h + \dots
  + \binom{m}{n - 1}x^{m-n+1} h^{n-1}
  + \binom{m}{n}(x + \theta_{n} h)^{m-n} h^{n}.
\]
さらに、\( x = 0 \)を含まない区間 \( (x, x + h) \)において、この公式はすべての実数値の \( m \)とすべての正の整数値の \( n \)に対して成り立ちます。また、\( x < 0 < x + h \)または \( x + h < 0 < x \)の場合でも、\( m - n \)が正であるとき、この公式は依然として成り立ちます。

\item5. \( f(x) = 1/x \)で \( x < 0 < x + h \)のとき、公式 \( f(x + h) = f(x) + hf'(x + \theta_{1}h) \)が成り立たないことが示されます。［\( f(x + h) - f(x) > 0 \)と \( hf'(x + \theta_{1} h) = -h/(x + \theta_{1} h)^{2} < 0 \)から、平均値の定理が成り立つ条件は満たされていないことが明らかです。］

\item[6.] \( x = -a \), \( h = 2a \), \( f(x) = x^{1/3} \)のとき、方程式
\[
f(x + h) = f(x) + hf'(x + \theta_{1} h)
\]
は \(\theta_{1} = \frac{1}{2} ± \frac{1}{18}\sqrt{3} \) で満たされる。[この例は、定理の結果が証明された条件を満たさない場合でも成立する可能性があることを示している。]
\pageseppage{265}

\item[7.] 方程式 \( f(x) = 0 \) の根に近似値 \(\xi\) を用いるとき、実際の根は \(\xi + h\) であるとする。すると
\[
0 = f(\xi + h) = f(\xi) + hf'(\xi) + \frac{1}{2} h^{2}f''(\xi + \theta_{2}h),
\]
よって
\[
h = -\frac{f(\xi)}{f'(\xi)}
  - \frac{1}{2} h^{2} \frac{f''(\xi + \theta_{2}h)}{f'(\xi)}.
\]

従って、一般に \( x = \xi \) より良い近似値は
\[
x = \xi - \frac{f(\xi)}{f'(\xi)}
\]
である。根が単根であり、すなわち \( f'(\xi + h) \neq 0 \) のとき、\( h \) が十分小さい場合に、ある正の定数 \( K \) が存在し、任意の考慮される \( x \) の値に対して \( |f'(x)| > K \) であり、また \( h \) を最初の位数とみなすとき、\( f(\xi) \) も最初の位数であり、根として \(\xi - \{f(\xi)/f'(\xi)\}\) の誤差は二つ目の位数である。

\item[8.] 方程式 \( x^{2} = 2 \) に対して、初めの近似値を \(\xi = 3/2\) とする。[このとき \( h = -1/12 \), \(\xi + h = 17/12 = 1.417\dots\) となり、最初の近似値が粗いにもかかわらず、かなり良い近似値であることがわかる。今後再びこの過程を繰り返し、\(\xi = 17/12 \) を用いるとき、\(\xi + h = 577/408 = 1.414\MS215\dots\) となり、小数点五桁まで正確であることがわかる。]

\item[9.] 方程式 \( x^{2} - 1 - y = 0 \) を考え、\( y \) が小さいとき、
\[
\sqrtp{1 + y} = 1 + \frac{1}{2} y - \left\{\frac{1}{4}y^{2}/(2 + y)\right\}
\]
と近似し、誤差は四つ目の位数であることを示す。

\item{10.} $\xi - (f/f') - \frac{1}{2}(f^{2}f''/f'^{3})$ を根として取り入れた誤差は、一般に三番目の順序までであります。

\item{11.} $\sin x = \alpha x$ の方程式において $\alpha$ が小さいとき、その根は $\pi$ に近い値をとります。$(1 - \alpha)\pi$ がより良い近似であり、$(1 - \alpha + \alpha^{2})\pi$ もさらに良い近似であることを示してください。[Exs.~7--10 の方法は $f(x) = 0$ が代数方程式である必要なく、$f'$ および $f''$ が連続的である限り有効です.]

\item{12.} $h \to 0$ のときの一般の中間値定理において現れる数 $\theta_{n}$ の極限は、$f^{(n+1)}(x)$ が連続的である条件下で $1/(n + 1)$ に等しいことを示してください。

[まず $f(x + h) = f(x) + \dots + \frac{h^{n}}{n!} f^{(n)}(x + \theta_{n}h)$ および
$f(x + h) = f(x) + \dots + \frac{h^{n}}{n!} f^{(n)}(x) + \frac{h^{n+1}}{(n + 1)!} f^{(n+1)}(x + \theta_{n+1}h)$ とします。ここで $\theta_{n}$ および $\theta_{n+1}$ はともに $0$ と $1$ の間の値をとります。したがって
\[
f^{(n)}(x + \theta_{n}h) = f^{(n)}(x) + \frac{hf^{(n+1)}(x + \theta_{n+1}h)}{n + 1}.
\]
また、$f^{(n)}(x)$ の原始の中間値定理を適用し、$\theta_{n}h$ を $h$ に代入すると
\[
f^{(n)}(x + \theta_{n}h) = f^{(n)}(x) + \theta_{n}hf^{(n+1)}(x + \theta\theta_{n}h)
\]
と得られます。ここで $\theta$ も $0$ と $1$ の間の値をとります。したがって
\[
\theta_{n} f^{(n+1)}(x + \theta\theta_{n} h) = \frac{f^{(n+1)}(x + \theta_{n+1} h)}{n + 1}
\]
となり、$h \to 0$ のとき $f^{(n+1)}(x + \theta\theta_{n} h)$ および $f^{(n+1)}(x + \theta_{n+1} h)$ は同じ極限 $f^{(n+1)}(x)$ に収束するから、この結果が導けます.]

\item{13.} $f(x + 2h) - 2f(x + h) + f(x)$ の値を $h^{2}$ で割った数が $h \to 0$ のとき $f''(x)$ に収束することを、$f''(x)$ が連続的である条件下で証明してください。[方程式 (2) を用いて §147 の定理を用います.]

14. \( f(x) \) の \( x = 0 \) において連続であるとき、
\[
f(x) = a_{0} + a_{1}x + a_{2}x^{2} + \cdots + (a_{n} + \epsilon_{x}) x^{n}
\]
と表せる。ここで \( a_{r} = f^{(r)}(0)/r! \) であり、\( \epsilon_{x} \to 0 \) である。

15. 
\[
a_{0} + a_{1}x + a_{2}x^{2} + \cdots + (a_{n} + \epsilon_{x}) x^{n} = b_{0} + b_{1}x + b_{2}x^{2} + \cdots + (b_{n} + \eta_{x}) x^{n}
\]
であり、かつ \( \epsilon_{x} \) および \( \eta_{x} \) が \( x \to 0 \) のとき \( 0 \) に収束するとき、\( a_{0} = b_{0} \), \( a_{1} = b_{1} \), \(\cdots\), \( a_{n} = b_{n} \) である。\( x \to 0 \) のとき \( a_{0} = b_{0} \) となる。次に \( x \) を分母に持ち上げ、その後 \( x \to 0 \) とすると \( a_{1} = b_{1} \) が得られる。この過程は必要に応じて繰り返すことができる。したがって、\( f(x) = a_{0} + a_{1}x + a_{2}x^{2} + \cdots + (a_{n} + \epsilon_{x}) x^{n} \) であり、かつ \( f(x) \) の最初の \( n \) つの導関数が連続であるとき、\( a_{r} = f^{(r)}(0)/r! \) である。

148. テイラー級数。点 \( x = a \) を含む区間 \( (a - \eta, a + \eta) \) または \([a - \eta, a + \eta]\) 内で \( f(x) \) のすべての微分係数が連続であるとする。このとき、\( h \) が \( \eta \) より小さいとき、
\[
f(a + h) = f(a) + hf'(a) + \cdots
  + \frac{h^{n-1}}{(n - 1)!} f^{(n-1)}(a)
  + \frac{h^{n}}{n!} f^{(n)}(a + \theta_{n} h),
\]
ここで \( 0 < \theta_{n} < 1 \) であり、すべての \( n \) の値に対して成り立つ。または
\[
S_{n} = \sum_{0}^{n-1} \frac{h^{\nu}}{\nu!} f^{(\nu)}(a),\quad
R_{n} = \frac{h^{n}}{n!} f^{(n)}(a + \theta_{n} h),
\]
とおくとき、
\[
f(a + h) - S_{n} = R_{n}.
\]

さらに、\( R_{n} \to 0 \) であることを証明できる場合、
\[
f(a + h) = \lim_{n\to\infty} S_{n}
  = f(a) + hf'(a) + \frac{h^{2}}{2!} f''(a) + \cdots.
\]

This展開は$f(a + h)$に適用され、Taylorの級数と呼ばれる。
$a = 0$の場合、式は
\[
f(h) = f(0) + hf'(0) + \frac{h^{2}}{2!} f''(0) + \dots,
\]
となる。これはMaclaurinの級数と呼ばれ、残差$R_{n}$はLagrangeの形の残差と呼ばれる。

---

Remark
読者はTaylorの級数が有効であるためには$f(x)$のすべての導関数の連続性だけでは十分でないと注意すべきであり、$R_{n}$の振る舞いについて直接的な議論は常に必要である。

---

Examples LVI.
1. $f(x) = \sin x$の場合、$f(x)$のすべての導関数は$x$のすべての値に対して連続である。また、$|f^{n}(x)| \leq 1$が$x$と$n$のすべての値に対して成り立つ。したがって、この場合、$|R_{n}| \leq h^{n}/n!$であり、$n \to \infty$のとき$h$の値によらず0に収束する（例題xxvii.12）。従って
\[
\sin(x + h) = \sin x + h\cos x - \frac{h^{2}}{2!}\sin x
  - \frac{h^{3}}{3!}\cos x + \frac{h^{4}}{4!}\sin x + \dots,
\]
が$x$と$h$のすべての値に対して成り立つ。特に
\[
\sin h = h - \frac{h^{3}}{3!} + \frac{h^{5}}{5!} - \dots,
\]
が$h$のすべての値に対して成り立つ。同様に、
\[
\cos(x + h) = \cos x - h\sin x - \frac{h^{2}}{2!}\cos x
  + \frac{h^{3}}{3!} \sin x + \dots,
\]
および
\[
\cos h = 1 - \frac{h^{2}}{2!} + \frac{h^{4}}{4!} - \dots.
\]

2. バイノミアル級数。$f(x) = (1 + x)^{m}$の場合、$m$は任意の有理数であり正でも負でもよい。このとき$f^{(n)}(x) = m(m - 1) \dots (m - n + 1)(1 + x)^{m-n}$であり、Maclaurinの級数は
\[
(1 + x)^{m} = 1 + \binom{m}{1}x + \binom{m}{2}x^{2} + \dots.
\]

当$m$が正の整数であるとき、級数は終了し、通常の二項定理の公式を得る。一般の場合においては
\[
R_{n} = \frac{x^{n}}{n!} f^{(n)}(\theta_{n}x)
  = \binom{m}{n} x^{n}(1 + \theta_{n}x)^{m-n},
\]
となる。そして、$m$が正の整数でないときマクローリン級数が任意の$x$の範囲において$(1 + x)^{m}$を表すことを示すために、$R_{n} \to 0$であることが全ての$x$の値に対して成り立つことを示さなければならない。実際、$-1 < x < 1$の場合にはこれが真であり、$0 \leq x < 1$の場合には上記で与えられた$R_{n}$の表式を使用して証明できる。$(1 + \theta_{n}x)^{m-n} < 1$が$n > m$のとき成り立つことと、$\binom{m}{n} x^{n} \to 0$が$n \to \infty$のとき成り立つことを考慮すると証明できる（例題xxvii.13参照）。

しかし、$-1 < x < 0$の場合には困難が生じる。$(1 + \theta_{n}x)^{m-n} > 1$が$n > m$のとき成り立つことと、$\theta_{n}$についてのみ$0 < \theta_{n} < 1$を知っているため、$1 + \theta_{n}x$は十分に小さいか否か、$(1 + \theta_{n}x)^{m-n}$は十分に大きいか否かを保証することができない。

実際、二項定理を Tayloy の定理を使って証明するためには、$R_{n}$の異なる形式が必要であり、後で与えられるもの（162節）のような形である必要がある。

\ Paragraph{149. Taylorの定理の応用。} \ Topic{\Item{A.} 最大値と最小値。} Taylorの定理は、第VI章の \okrickRef{Ch.}{VI}, 122節--123節で述べた検定により理論的な完璧さを与えることができるが、その結果は実用的にはあまり重要ではない。$\phi(x)$が一階と二階の導関数を持つことを仮定し、$\phi(x)$の$x = \xi$における最大値または最小値が次の条件で十分であるとした：\emph{最大値の場合}では $\phi'(\xi) = 0$, $\phi''(\xi) < 0$; \emph{最小値の場合}では $\phi'(\xi) = 0$, $\phi''(\xi) > 0$。これらの検定は$\phi''(\xi)$だけでなく$\phi'(\xi)$がゼロである場合に失敗することが明らかである。

Let us suppose bahwa最初の \( n \) つの導関数
\[
\phi'(x),\quad \phi''(x),\ \dots,\quad \phi^{(n)}(x)
\]
は連続であり、これら全てが \( x = \xi \) のとき零であるものを除いて。すると、十分に小さな値の \( h \) に対して，
\[
\phi(\xi + h) - \phi(\xi) = \frac{h^{n}}{n!} \phi^{(n)} (\xi + \theta_{n} h)
\]
が成り立つ。

極大または極小があるためには、この式が十分に小さな値の \( h \) に対して常に符号を保つべきであり、正でも負でもよい。これは明らかに \( n \) が偶数であることを要求する。そして \( n \) が偶数であれば、\(\phi^{(n)}(\xi)\) が負であれば極大、正であれば極小となる。

したがって、次のテストを得る：\par\medskip\noindent\textbf{【結果】}\par極大または極小があるためには、零にならない最初の導関数は偶数であるべきであり、それが負であれば極大、正であれば極小となる。
\par\medskip

\par\medskip\noindent\textbf{【例 LVII.】}\par
\Item{1.} \(\phi(x) = (x - a)^{m}\), \( m \) が正の整数であり、\(\xi = a\) のとき結果を検証せよ。

\Item{2.} \( x = a \), \( x = b \) の点で最大値や最小値があるかどうかをテストし、曲線 \( y = (x - a)^{m} (x - b)^{n} \) の異なる可能性のある形のグラフを描け。ここで \( m \) と \( n \) は正の整数である。

\Item{3.} \( x = 0 \) での最大値や最小値があるかどうかをテストせよ：\(\sin x - x\), \(\sin x - x + \dfrac{x^{3}}{6}\), \(\sin x - x + \dfrac{x^{3}}{6} - \dfrac{x^{5}}{120}\) など、\(\cos x - 1\), \(\cos x - 1 + \dfrac{x^{2}}{2}\), \(\cos x - 1 + \dfrac{x^{2}}{2} - \dfrac{x^{4}}{24}\) など。
\par\medskip

\Paragraph{150.} \Topic{\Item{B.} 某些極限の計算。} \( f(x) \) と \(\phi(x)\) を \( x = \xi \) のとき連続な導関数 \( f'(x) \) と \(\phi'(x)\) を持つ二つの \( x \) の関数とする。そして \( f(\xi) \) と \(\phi(\xi)\) がともに零である。すると、関数
\[
\psi(x) = \frac{f(x)}{\phi(x)}
\]
は \( x = \xi \) のとき定義されない。しかし当然、\( x \to \xi \) のときこの関数は極限に向かう可能性がある。

現在
\[
f(x) = f(x) - f(\xi) = (x - \xi)f'(x_{1}),
\]
ここで \( x_{1} \) は \(\xi\) と \(x\) の間に存在します；同様に \(\phi(x) = (x - \xi)\phi'(x_{2})\) で、\( x_{2} \) も \(\xi\) と \(x\) の間に存在します。したがって
\[
\psi(x) = \frac{f'(x_{1})}{\phi'(x_{2})}.
\]
今後、四つの場合を区別する必要があります。

(1) もし \( f'(\xi) \) と \(\phi'(\xi)\) のいずれも零でないなら
\[
\frac{f(x)}{\phi(x)} \to \frac{f'(\xi)}{\phi'(\xi)}.
\]

(2) もし \( f'(\xi) = 0 \), \(\phi'(\xi) \neq 0\) の場合、
\[
\frac{f(x)}{\phi(x)} \to 0.
\]

(3) もし \( f'(\xi) \neq 0 \), \(\phi'(\xi) = 0\) の場合、\( x \to \xi \) のとき \(\frac{f(x)}{\phi(x)}\) は数値的に非常に大きくなります：しかし \(\frac{f(x)}{\phi(x)}\) が無限大に収束するか負の無限大に収束するか、または正の大きな値と負の大きな値を交互に取るかについては、\(\phi'(x) \to 0\) の方法についてさらに情報がない限りは言えません。

(4) もし \( f'(\xi) = 0 \), \(\phi'(\xi) = 0\) の場合、\( x \to \xi \) のとき \(\frac{f(x)}{\phi(x)}\) の振る舞いについて何も言えません。

しかし最後の二つの場合のいずれかで、\( f(x) \) と \(\phi(x)\) が連続的な2次の導関数を持つ可能性があります。その場合
\begin{align*}
f(x) &= f(x) - f(\xi) - (x - \xi)f'(\xi)
  = \frac{1}{2}(x - \xi)^{2} f''(x_{1}),\\
\phi(x) &= \phi(x) - \phi(\xi) - (x - \xi)\phi'(\xi)
  = \frac{1}{2}(x - \xi)^{2} \phi''(x_{2}),
\end{align*}
ここでまた \( x_{1} \) と \( x_{2} \) は \(\xi\) と \(x\) の間に存在します；したがって
\[
\psi(x)= \frac{f''(x_{1})}{\phi''(x_{2})}.
\]
今後、上記の場合に類似する様々な場合を区別することができます。特に、もし \( x = \xi \) で 2次の導関数のいずれも零でないなら
\[
\frac{f(x)}{\phi(x)} \to \frac{f''(\xi)}{\phi''(\xi)}.
\]

この議論は無限に繰り返すことができ、以下の定理を得ます：\par\medskip\noindent\textbf{【結果】}\par \( f(x) \) と \(\phi(x)\) およびそれらの導関数（必要であれば）が \( x = \xi \) で連続であると仮定します。さらに、\( f^{(p)}(x) \) と \(\phi^{(q)}(x)\) が最初に零でない \( f(x) \) と \(\phi(x)\) の導関数であると仮定します。すると

\Item{(1)} \( p = q \)のとき、\( f(x)/\phi(x) \to f^{(p)}(\xi)/\phi^{(p)}(\xi) \);

\Item{(2)} \( p > q \)のとき、\( f(x)/\phi(x) \to 0 \);

\Item{(3)} \( p < q \)かつ \( q - p \)が偶数のとき、\( f(x)/\phi(x) \to +\infty \)または\( f(x)/\phi(x) \to -\infty \)であり、その符号は\( f^{(p)}(\xi)/\phi^{(q)}(\xi) \)と同じである;

\Item{(4)} \( p < q \)かつ \( q - p \)が奇数のとき、\( x \to \xi+0 \)のとき、\( f(x)/\phi(x) \to +\infty \)または\( f(x)/\phi(x) \to -\infty \)であり、その符号は\( f^{(p)}(\xi)/\phi^{(q)}(\xi) \)と同じであるが、\( x \to \xi - 0 \)のとき符号は逆となる。

この定理は実際には以下の等式から即座に導かれるものである：
\[
f(x) = \frac{(x - \xi)^{p}}{p!} f^{(p)}(x_{1}),\quad
\phi(x) = \frac{(x - \xi)^{q}}{q!} \phi^{(q)}(x_{2}).
\]

\par\medskip\noindent\textbf{【例 LVIII.】}\par
\Item{1.} \( x - (n + 1)x^{n+1} + nx^{n+2} \)を\( (1 - x)^{2} \)で割った極限を求めよ。[ここでは、\( x = 1 \)において関数とその第1導函数が消滅し、\( f''(1) = n(n + 1) \), \( \phi''(1) = 2 \)である]

\Item{2.} \( x \to 0 \)のとき、次の極限を求めよ：
\[
\tan x - x)/(x - \sin x),\quad
(\tan nx - n\tan x)/(n\sin x - \sin nx).
\]

\Item{3.} \( x \to \infty \)のとき、\( x\{\sqrtp{x^{2} + a^{2}} - x\} \)の極限を求めよ。[ここで \( x = 1/y \)と置く]

\Item{4.} 次の極限を証明せよ：
\[
\lim_{x \to n} (x - n)\cosec x\pi = \frac{(-1)^{n}}{\pi},\quad
\lim_{x \to n} \frac{1}{x - n} \left\{
  \cosec x\pi - \frac{(-1)^{n}}{(x - n)\pi}
\right\} = \frac{(-1)^{n}\pi}{6},
\]
\( n \)は任意の整数であり、それに伴う極限を\(\cot x\pi\)を含むものも評価せよ。

\Item{5.} \( x \to 0 \)のとき、次の極限を求めよ：
\[
\frac{1}{x^{3}}\left(\cosec x - \frac{1}{x} - \frac{x}{6}\right),\quad
\frac{1}{x^{3}}\left(\cot x - \frac{1}{x} + \frac{x}{3}\right).
\]

\Item{6.} \( (\sin x\arcsin x - x^{2})/x^{6} \to \frac{1}{18} \), \( (\tan x\arctan x - x^{2})/x^{6} \to \frac{2}{9} \)である。
\par\medskip

151. 平面曲線の接点について。二つの曲線が一点で交わる（または切れる）とは、その点が二つの曲線上にあることを指します。また、二つの曲線が一点で接するとは、その点での接線が同じであることを指します。

今、$f(x)$, $\phi(x)$ が $x = \xi$ において任意の階数の導関数を連続に持つ二つの関数とし、曲線 $y = f(x)$, $y = \phi(x)$ を考える。一般には $f(\xi)$ と $\phi(\xi)$ は等しくありません。この場合、$x = \xi$ の横座標に対応する点は二つの曲線の交点ではありません。しかし、$f(\xi) = \phi(\xi)$ のとき、点 $x = \xi$, $y = f(\xi) = \phi(\xi)$ で二つの曲線が交わります。この場合、曲線がこの点で接するためには明らかに最初の導関数 $f'(x)$, $\phi'(x)$ も $x = \xi$ のとき同じ値を持つことが必要かつ十分です。

この場合の二つの曲線の接点について別の視点から見ると、図において二つの曲線は点 $P$ で接しており、$QR$ は $\phi(\xi + h) - f(\xi + h)$ に等しいか、$\phi(\xi) = f(\xi)$, $\phi'(\xi) = f'(\xi)$ より $\phi(\xi + h) - f(\xi + h) = \frac{1}{2} h^2\{\phi''(\xi + \theta h) - f''(\xi + \theta h)\}$ となる、$\theta$ は $0$ と $1$ の間の値を持つ。したがって
\[
\lim_{h \to 0} \frac{QR}{h^2} = \frac{1}{2}\{\phi''(\xi) - f''(\xi)\}
\]
となります。つまり、曲線が点 $\xi$ の横座標に対応する点で接するとき、\emph{点 $\xi + h$ の縦座標の差は $h$ が小さいときに二階の小量である}。

補足：
交わるが接しない場合、$\lim \frac{QR}{h} = \phi'(\xi) - f'(\xi)$ となり、このとき $QR$ は一階の小量に過ぎません。

弧度$QR$の小ささは、曲線の接触の密接さの一種の尺度として考えられることが明らかである。したがって、点$x = \xi$において$f$と$\phi$の最初の$n-1$階微分が等しいならば、$QR$は$n$階の小さい量となることがすぐに提案される。読者はこの事実を証明することは難しくないし、

\[
\lim_{h \to 0} \frac{QR}{h^n} = \frac{1}{n!}\{\phi^{(n)}(\xi) - f^{(n)}(\xi)\}
\]

であることも証明できる。

したがって、次の定義を導く：

\par\medskip\noindent\textbf{【定義】}\par
$n$階の接触。$f(\xi) = \phi(\xi)$, $f'(\xi) = \phi'(\xi)$, $\dots$, $f^{(n)}(\xi) = \phi^{(n)}(\xi)$であり、しかし$f^{(n+1)}(\xi) \neq \phi^{(n+1)}(\xi)$であるとき、曲線$y = f(x)$, $y = \phi(x)$は点の абсциссаが$\xi$である点において$n$階の接触があるとされる。
\par\medskip

前記の議論により、$n$階の接触の概念は軸の選択に依存し、曲線の接線が$y$軸に平行な場合には完全に失われる。この状況を扱うためには、$y$を独立変数とし、$x$を従属変数として取り扱うことができる。しかし、より良い方法は$x$と$y$をパラメータ$t$の関数として考えるということである。理論に関する優れた説明は、第266ページに引用されたFowler氏の論文やde la Vallée Poussinの『解析教程』vol. ii, pp. 396 et seq. を参照できる。

\par\medskip\noindent\textbf{【例 LIX.】}\par
1. $\phi(x) = ax + b$とするとき、$y = \phi(x)$は直線である。点$x = \xi$における接触条件は$f(\xi) = a\xi + b$, $f'(\xi) = a$である。これらの式を満たす$a$と$b$を求めることにより、$a = f'(\xi)$, $b = f(\xi) - \xi f'(\xi)$を得る。したがって、点$x = \xi$における曲線$y = f(x)$の接線の方程式は、

\[
y = xf'(\xi) + \{f(\xi) - \xi f'(\xi)\}
\]

または

\[
y - f(\xi) = (x - \xi)f'(\xi)
\]

となる。参照：例題xxxix. 5。
\par\medskip

\item{2.} 直線が曲線と単純接触を持つことは直線を完全に決定する。接線が曲線に対して二階の接触を持つためには \( f''(\xi) = \phi''(\xi) \), すなわち \( f''(\xi) = 0 \) であることが必要である。接線が曲線と二階の接触を持つ点は、不変点又は屈折点と呼ばれる。

\item{3.} 関数 \( 3x^{4} - 6x^{3} + 1 \), \( 2x/(1 + x^{2}) \), \( \sin x \), \( a\cos^{2}x + b\sin^{2}x \), \( \tan x \), \( \arctan x \) の屈折点を求めよ。

\item{4.} 二次曲線 \( ax^{2} + 2hxy + by^{2} + 2gx + 2fy + c = 0 \) は屈折点を持つことができないことを示せ。［ここで \( ax + hy + g + (hx + by + f)y_{1} = 0 \) であり
\[
a + 2hy_{1} + by_{1}^{2} + (hx + by + f)y_{2} = 0,
\]
接数は微分を表す。屈折点において
\[
a + 2hy_{1} + by_{1}^{2} = 0,
\]
または
\[
a(hx + by + f)^{2} - 2h(ax + hy + g)(hx + by + f) + b(ax + hy + g)^{2} = 0,
\]
または
\[
(ab - h^{2})\{ax^{2} + 2hxy + by^{2} + 2gx + 2fy\} + af^{2} - 2fgh + bg^{2} = 0.
\]
しかし、これは二次曲線の方程式と矛盾するため
\[
af^{2} - 2fgh + bg^{2} = c(ab - h^{2})
\]
または \( abc + 2fgh - af^{2} - bg^{2} - ch^{2} = 0 \) であることが必要であり、これは二次曲線が二つの直線に退化する条件である。］

\item{5.} 曲線 \( y = (ax^{2} + 2bx + c)/(\alpha x^{2} + 2\beta x + \gamma) \) は、\( \alpha x^{2} + 2\beta x + \gamma = 0 \) の根が実数であるか複素数であるかにより、一つまたは三つの屈折点を持つ。

[曲線の方程式は原点の変換（例 xlvi.15 参照）によって
\[
\eta = \xi/(A\xi^{2} + 2B\xi + C) = \xi/\{A(\xi - p)(\xi - q)\},
\]
と簡約される。屈折点の条件は \( \xi^{3} - 3pq\xi + pq(p + q) = 0 \) であり、\( \DPtypo{\{pq(p - q)\}}{\{pq(p - q)\}^{2}} \) が正であるか負であるにより、すなわち \( p \) と \( q \) が実数であるか複素共役であるかにより、一つまたは三つの実数根を持つ。]

\item 6. \( y = (1 - x)/(1 + x^2) \), \( y = (1 - x^2)/(1 + x^2) \), \( y = (1 + x^2)/(1 - x^2) \) の曲線について特に議論せよ。

\item 7. Ex.~5 の曲線が3つの反転点を持つとき、これらの点は一直線上に存在することを示せ。[方程式 \( \xi^3 - 3pq\xi + pq(p + q) = 0 \) を形変形して \( (\xi - p)(\xi - q)(\xi + p + q) + (p - q)^2\xi = 0 \) と書き換えることができ、したがって反転点は直線 \( \xi + A(p - q)^2\eta + p + q = 0 \) や \( A\xi - 4(AC - B^2)\eta = 2B \) 上に存在する。]

\item 8. \( y = x\sin x \), \( y = (\sin x)/x \) の各曲線は無数の反転点を持つことを示せ。

\item 9. 円と曲線の接続．曲率．
円の方程式は一般に以下の形を取る：
\[
(x - a)^2 + (y - b)^2 = r^2.
\tag{1}
\]
この方程式には3つの任意定数が含まれている。点 \( (\xi, \eta) \)，ここで \( \eta = f(\xi) \) における曲線 \( y = f(x) \) と円の接続を最高次に達するようにこれらの定数を決定しようとする。\( f'(\xi) \), \( f''(\xi) \) をそれぞれ \( \eta_1 \), \( \eta_2 \) と表す。円の方程式を2回微分すると
\begin{align}
(x - a) + (y - b)y_1 &= 0,
\tag{2}\\
1 + y_1^2 + (y - b)y_2 &= 0.
\tag{3}
\end{align}

円が曲線と接するとき、方程式 \Eq{(1)} および \Eq{(2)} は \( x = \xi \), \( y = \eta \), \( y_1 = \eta_1 \) のときに満足される。これにより \( (\xi - a)/\eta_1 = -(\eta - b) = r/\sqrt{1 + \eta_1^2} \) となる。接触が二階まで達するとき、方程式 \Eq{(3)} も \( y_2 = \eta_2 \) のときに満足される必要がある。したがって \( b = \eta + (1 + \eta_1^2)/\eta_2 \); これにより
\[
a = \xi - \frac{\eta_1(1 + \eta_1^2)}{\eta_2},\quad
b = \eta + \frac{1 + \eta_1^2}{\eta_2},\quad
r = \frac{(1 + \eta_1^2)^{3/2}}{\eta_2}.
\]

円が点 $(\xi, \eta)$ で二階接触する円は、その接曲線の \Emph{接円} と呼ばれ、その半径は \Emph{曲率半径} である。また、\Emph{曲率の測度}（単に \emph{曲率}）は半径の逆数であり、したがって曲率の測度は $f''(\xi)/\{1 + [f'(\xi)]^{2}\}^{3/2}$ または
\[
\frac{d^{2}\eta}{d\xi^{2}} \bigg/
  \biggl\{1 + \biggl(\frac{d\eta}{d\xi}\biggr)^{2}\biggr\}^{3/2}
\]
である。

10. 接円の曲率は定数であり、その値は半径の逆数に等しいことを確認し、また接円が曲率が一定の唯一の曲線であることを示せ。

11. 任意の点で放物線 $y^{2} = 4ax$ および楕円 $(x/a)^{2} + (y/b)^{2} = 1$ の \Emph{接円の中心} と \Emph{半径} を求めよ。

12. 楕円において点 $P$ での曲率半径は $CD^{3}/ab$ であり、ここで $CD$ は $CP$ に直交する半径共役線分である。

13. 一般の点 $P$ で四階接触する曲線 $y = f(x)$ 上の放物線を一意的に定めることができる事を示せ。一般的な放物線の方程式は
\[
ax^{2} + 2hxy + by^{2} + 2gx + 2fy + c = 0
\]
であり、これを $x$ について四回微分すると次の式を得る：
\begin{align*}
ax + hy + g + (hx + by + f) y_{1} &= 0,\\
a + 2hy_{1} + by_{1}^{2} + (hx + by + f) y_{2} &= 0,\\
3(h + by_{1}) y_{2} + (hx + by + f) y_{3} &= 0,\\
4(h + by_{1}) y_{3} +3by_{2}^{2} + (hx + by + f) y_{4} &= 0.
\end{align*}
四階接触を持つ放物線の場合は、これらの五つの式が $x$, $y$, $y_{1}$, $y_{2}$, $y_{3}$, $y_{4}$ を $\xi$, $\eta$, $\eta_{1}$, $\eta_{2}$, $\eta_{3}$, $\eta_{4}$ に置き換えることで満たされる。したがって、$a : b : c : f : g : h$ の比を決定するのに十分な数の式を得る。

14. 与えられた点 $P$ で三階接触する放物線は無数に存在し、その中心は一直線上にあることを示せ。

[接線と法線を座標軸とするとき、放物線の方程式は
\[2y = ax^{2} + 2hxy + by^{2}\]
の形になります。そして \(x\) が小さいとき、ある \(y\) の値は (\okrickRef{Ch.}{V}, \MiscEx{V} 22) 次のように表すことができます：
\[
y = \frac{1}{2}ax^{2} + \left(\frac{1}{2}ah + \epsilon_{x}\right)x^{3},
\]
ここで \(\epsilon_{x} \to 0\) です。しかし、この表現は次の式と等しくなければなりません：
\[
y = \frac{1}{2}f''(0)x^{2} + \left\{\frac{1}{6}f'''(0) + \epsilon'_{x}\right\}x^{3},
\]
ここで \(\epsilon'_{x} \to 0\) です。したがって \(a = f''(0)\), \(h = \frac{f'''(0)}{3f''(0)}\) となります。ただし、結果 \Ex{lv}.15 の virtue によります。しかし、中心は直線 \(ax + hy = 0\) 上にあります。]

[15. 接触が三次である放物線を決定し、楕円 \(\left(\frac{x}{a}\right)^{2} + \left(\frac{y}{b}\right)^{2} = 1\) の主軸の端点で接するようにします。]

[16. 接触が三次である放物線の中心の軌跡は、楕円 \(\left(\frac{x}{a}\right)^{2} + \left(\frac{y}{b}\right)^{2} = 1\) の点 \((a\cos\alpha, b\sin\alpha)\) で接する放物線の中心の軌跡は直径 \(x/(a\cos\alpha) = y/(b\sin\alpha)\) です。[自身がそのような放物線の一つであるため。]]

[152. 多変数関数の微分法。]
これまでに、単変数 \(x\) の関数についてのみ論じてきましたが、複数の変数 \(x\), \(y\), \dots に対する微分概念を適用することも可能です。]

次に、$f(x, y)$ を2つの実変数 $x$ および $y$ の関数とする。そして、次の極限がすべての $x$ および $y$ の値に対して存在するものとする：
\[
\lim_{h \to 0} \frac{f(x + h, y) - f(x, y)}{h}, \quad
\lim_{k \to 0} \frac{f(x, y + k) - f(x, y)}{k}
\]
これらの極限が存在する場合、$f(x, y)$ は $x$ に関して $\frac{\partial f}{\partial x}$ や $D_x f(x, y)$ と呼ばれる導関数、および $y$ に関して $\frac{\partial f}{\partial y}$ や $D_y f(x, y)$ と呼ばれる導関数を持つと言える。これらを簡略化して $f_x'$, $f_y'$, $f_x$, $f_y$ と表すことが多い。ただし、これらの新しい記号が根本的な概念の変革を意味するものではないことを忘れないでほしい。$x$ に関しての「部分微分」は通常の微分と同じ過程であり、唯一の新しさは関数 $f$ 内に独立した別の変数 $y$ の存在にある。

precedentにおいて、xとyを互いに独立な2つの実数変数として想定していました。もしxとyが関係によって結びついていたなら、状況は非常に異なることになります。この場合、私達の$f_{x}'$の定義は完全に失われてしまいます。なぜなら、xをx + hに変更することができない限り同時にyも変更されてしまうからです。しかし、その結果f(x, y)は実質的に2つの変数による関数ではなくなります。2つの変数による関数として我々が定義したものは、基本的には2つの独立な変数による関数です。もしyがxに依存しているなら、yはxの関数であり、y = φ(x)と表すことができます。そして
\[
f(x, y) = f\{x, \phi(x)\}
\]
は実質的に単一の変数xによる関数となります。もちろん、これは単一の変数yによる関数としても表現できます。または、しばしば最も便利な場合、xとyを第三の変数tによる関数として考えることができます。そして
\[
f(x, y) = f\{\phi(t), \psi(t)\}
\]
は実質的に単一の変数tによる関数となります。

例 LX．

1. x = r cos θ, y = r sin θ であるとき、r = √(x² + y²), θ = arctan(y/x) から
\[
\frac{\dd r}{\dd x} = \frac{x}{\sqrt{x^2 + y^2}}, \quad
\frac{\dd r}{\dd y} = \frac{y}{\sqrt{x^2 + y^2}},
\]
\[
\frac{\dd \theta}{\dd x} = -\frac{y}{x^2 + y^2}, \quad
\frac{\dd \theta}{\dd y} = \frac{x}{x^2 + y^2},
\]
\[
\frac{\dd x}{\dd r} = \cos \theta, \quad
\frac{\dd y}{\dd r} = \sin \theta,
\]
\[
\frac{\dd x}{\dd \theta} = -r \sin \theta, \quad
\frac{\dd y}{\dd \theta} = r \cos \theta.
\]

2. 1変数$x$の関数$y$を考えたときには$\dfrac{\mathrm{d}r}{\mathrm{d}x}\neq \left(\dfrac{\mathrm{d}x}{\mathrm{d}r}\right)^{-1}$および$\dfrac{\mathrm{d}\theta}{\mathrm{d}x}\neq \left(\dfrac{\mathrm{d}x}{\mathrm{d}\theta}\right)^{-1}$であることを説明せよ。[変数$x$の関数$y$について考えたときには、定義から$\dfrac{\mathrm{d}y}{\mathrm{d}x}$と$\dfrac{\mathrm{d}x}{\mathrm{d}y}$は互いに逆数であることが明らかだった。しかし2変数の関数を考えるときにはこの関係が成り立たない。図46の点$P$は$(x, y)$または$(r, \theta)$である。$\dd r/\dd x$を求めることを目的とするので、$x$に微小な増分$\delta x$を与えることで$P$を$P_1$へ移動させるが、$y$は一定である。このとき$OP_1$上に$OP'= OP$とし、$r$の増分$\delta r$を$P'P_1 = \delta r$とする。したがって$\dd r/\dd x = \lim\left(\frac{\delta r}{\delta x}\right)$である。一方で$\dd x/\dd r$を求めることを考えるとき、$x$と$y$は$r$および$\theta$の関数として扱われる。この場合、$r$に微小な増分$\Delta r$を与えることで$P$を$P_2$へ移動させるが、$\theta$は一定である。これにより$x$の増分$\Delta x = MM_1 = \Delta x$となる。したがって$\dd x/\dd r = \lim\left(\frac{\Delta x}{\Delta r}\right)$である。

ここで$\Delta x = \delta x$である：これは$\Delta r$の選択（ここでは$PP_2$）によるものであり、他の選択肢でも$\Delta x$および$\Delta r$はここで使用した値と比例する。しかし$\Delta r \neq \delta r$である。図から容易に確認できるように、

\[
\lim \left(\frac{\delta r}{\delta x}\right) = \lim \left(\frac{P'P_1}{PP_1}\right) = \cos\theta
\]

であり、

\[
\lim \left(\frac{\Delta r}{\Delta x}\right) = \lim \left(\frac{PP_2}{PP_1}\right) = \sec\theta
\]

であるから、

\[
\lim \left(\frac{\delta r}{\Delta r}\right) = \cos^2\theta
\]

となる。

つまり$\dd x/\dd r$および$\dd r/\dd x$は同一の仮定に基づいて形成されるわけではないという事実が明らかである。

3. $z = f(ax + by)$のとき$b(\dd z/\dd x) = a(\dd z/\dd y)$を証明せよ。

\item 4. \(X + Y = x\)、\(Y = xy\)のとき、\(\frac{\dd X}{\dd x}\), \(\frac{\dd X}{\dd y}\)，…を求めよ。\(x\), \(y\)を\(X\), \(Y\)の関数として表し、\(\frac{\dd x}{\dd X}\), \(\frac{\dd x}{\dd Y}\)，…を求めよ。

\item 5. \(X + Y + Z = x\)、\(Y + Z = xy\)、\(Z = xyz\)のとき、\(\frac{\dd X}{\dd x}\)，…を求めよ。\(x\), \(y\), \(z\)を\(X\), \(Y\), \(Z\)の関数として表し、\(\frac{\dd x}{\dd X}\)，…を求めよ。

[当然、最後の節で説明したアイデアは、任意の数の変数に対する関数にも拡張できる。しかし、読者はすべての独立変数が指定されているときのみ、複数の変数を持つ関数の偏微分の概念が確定していることを心に留めておくべきである。例えば、\(u = x + y + z\)において、\(x\), \(y\), \(z\)が独立変数であるとき、\(\frac{\dd u}{\dd x} = 1\)である。しかし、\(u\)を変数\(x\), \(x + y = \eta\), \(x + y + z = \zeta\)に対する関数と見做す場合、つまり\(u = \zeta\)であるとき、\(\frac{\dd u}{\dd x} = 0\)となる。]

\paragraph{153. 二つの関数の関数に対する微分}
一変数に対する関数の微分に関する定理として知られる「全微分係数の定理」について述べる。この定理は非常に重要で、前節で説明した二変数を持つ関数に関する概念に基づいている。この定理は
\[
f\{\phi(t), \psi(t)\}
\]
を\(t\)に関して微分する方法を与える。

まず最初に、\(f(x, y)\)が二つの変数 \(x\) と \(y\) の関数であり、その両方の変数に対するすべての値について連続な関数 \(f_{x}'\), \(f_{y}'\) であるとする（§107）。そして、\(x\) と \(y\) の変化が制限され、曲線
\[
x = \phi(t),\quad y = \psi(t),
\]
上にあるときを考える。ここで \(\phi\) と \(\psi\) は連続な微分係数を持つ関数である \(\phi'(t)\), \(\psi' (t)\) である。すると \(f(x, y)\) は単変数 \(t\) の関数 \(F(t)\) となる。問題は \(F'(t)\) を決定することである。

$t$ が $t + \tau$ に変化すると、$x$ と $y$ はそれぞれ $x + \xi$ および $y + \eta$ に変化します。このとき定義により、

\[
\frac{dF(t)}{dt} = \lim_{\tau \to 0} \frac{1}{\tau} [f(\phi(t + \tau), \psi(t + \tau)) - f(\phi(t), \psi(t))]
= \lim \frac{1}{\tau} \{f(x + \xi, y + \eta) - f(x, y)\}
= \lim \left[
    \frac{f(x + \xi, y + \eta) - f(x, y + \eta)}{\xi} \cdot \frac{\xi}{\tau}
  + \frac{f(x, y + \eta) - f(x, y)}{\eta} \cdot \frac{\eta}{\tau}
\right].
\]

平均値の定理により、

\[
\frac{f(x + \xi, y + \eta) - f(x, y + \eta)}{\xi} = f_{x}'(x + \theta\xi, y + \eta),
\quad
\frac{f(x, y + \eta) - f(x, y)}{\eta} = f_{y}'(x, y + \theta'\eta),
\]

ここで $\theta$ および $\theta'$ はそれぞれ $0$ と $1$ の間の値を取る。$\tau \to 0$ のとき、$\xi \to 0$ および $\eta \to 0$ であり、$\xi/\tau \to \phi'(t)$ および $\eta/\tau \to \psi'(t)$：また、

\[
f_{x}'(x + \theta\xi, y + \eta) \to f_{x}'(x, y), \quad
f_{y}'(x, y + \theta'\eta) \to f_{y}'(x, y).
\]

したがって、

\[
F'(t) = D_{t}f(\phi(t), \psi(t))
= f_{x}'(x, y)\phi'(t) + f_{y}'(x, y)\psi'(t),
\]

ここで $x = \phi(t)$、$y = \psi(t)$ と置く必要がある。この結果はまた、

\[
\frac{df}{dt} = \frac{\dd f}{\dd x} \cdot \frac{dx}{dt} + \frac{\dd f}{\dd y} \cdot \frac{dy}{dt}
\]

という形で表すことができる。

例 LXI．
1. $\phi(t) = (1 - t^{2})/(1 + t^{2})$、$\psi(t) = 2t/(1 + t^{2})$ のとき、$(x, y)$ の軌跡は円 $x^{2} + y^{2} = 1$ である。このとき、

\[
\phi'(t) = -4t/(1 + t^{2})^{2}, \quad \psi'(t) = 2(1 - t^{2})/(1 + t^{2})^{2},
\]

\[
F'(t) = \left\{-4t/(1 + t^{2})^{2}\right\} f_{x}' + \left\{2(1 - t^{2})/(1 + t^{2})^{2}\right\} f_{y}',
\]

ここで $x$ および $y$ は微分を実行した後、$(1 - t^{2})/(1 + t^{2})$ および $2t/(1 + t^{2})$ に置く必要がある。

特定の場合においてこの公式が成り立つことを容易に確認できます。たとえば、$f(x, y) = x^{2} + y^{2}$ とするとき、$f_{x}' = 2x$, $f_{y}' = 2y$ であり、$F'(t) = 2x\phi'(t) + 2y\psi'(t) = 0$ が容易に確認され、これは明らかに正しく、$F(t) = 1$ より成り立つことがわかります。

2. 同様の方法で次の条件において定理を確認します。(\ia) $x = t^{m}$, $y = 1 - t^{m}$, $f(x, y) = x + y$; (\ib) $x = a\cos t$, $y = a\sin t$, $f(x, y) = x^{2} + y^{2}$。

3. 最も重要な場合の一つは、$t$ が $x$ 自身であるときです。この場合、
\[ D_{x}f\{x, \psi(x)\} = D_{x}f(x, y) + D_{y}f(x, y)\psi'(x)
\]
を得ます。ここで、微分後 $\psi(x)$ に置き換える必要があります。

この場合を導入するために $\dd f/\dd x$, $\dd f/\dd y$ の記号が使用されました。$D_{x}f\{x, \psi(x)\}$ と $D_{x}f(x, y)$ のいずれかの関数に対して、$df/dx$ を使用することが自然であるように思われます。一方は微分前に $\psi(x)$ に置き換える一方は微分後に置き換えるためです。たとえば、$y = 1 - x$, $f(x, y) = x + y$ の場合、$D_{x}f(x, 1 - x) = D_{x}1 = 0$ ですが、$D_{x}f(x, y) = 1$。

二つの関数の区別は、最初のものを $df/dx$, 二つ目のものを $\dd f/\dd x$ と表記することで適切に示されます。この場合、定理は次のようになります：
\[
\frac{df}{dx} = \frac{\dd f}{\dd x} + \frac{\dd f}{\dd y}\, \frac{dy}{dx};
\]
しかし、この記号もまた一部の誤解を招く可能性があります。$f\{x, \psi(x)\}$ と $f(x, y)$ の形は $x$ の関数として全く異なるため、$df/dx$ と $\dd f/\dd x$ で同じ文字 $f$ を使用することは少し誤解を招く可能性があります。

4. $x = \phi(t)$, $y = \psi(t)$ から $t$ を消去した結果が $f(x, y) = 0$ のとき、
\[
\frac{\dd f}{\dd x}\, \frac{dx}{dt} + \frac{\dd f}{\dd y}\, \frac{dy}{dt} = 0.
\]

\Item{5.} \(x\) と \(y\) が \(t\) の関数であり、\(r\) と \(\theta\) が点 \((x, y)\) の極座標であるとき、\(r' = (xx' + yy')/r\)、\(\theta' = (xy' - yx')/r^{2}\) であり、破線は \(t\) に関する微分を表す。

\Paragraph{154. 二変数関数の平均値の定理} 最後の章で得られた多くの結果は平均値の定理に依存しており、その定理は方程式
\[
\phi(x + h) - \phi(x) = hf'(x + \theta h)
\]
または \(y = \phi(x)\) の場合、
\[
\delta y = f'(x + \theta \, \delta x) \, \delta x
\]
と表される。

今、\(z = f(x, y)\) が独立変数 \(x\) および \(y\) の関数であり、\(x\) および \(y\) にそれぞれ増分 \(h\)、\(k\) または \(\delta x\)、\(\delta y\) を与えると仮定する。そして、その対応する \(z\) の増分を表すために
\[
\delta z = f(x + h, y + k) - f(x, y)
\]
を \(h\)、\(k\) および \(x\)、\(y\) に関する \(z\) の導関数で表そうとする。

\(f(x + ht, y + kt) = F(t)\) と定義すると
\[
f(x + h, y + k) - f(x, y) = F(1) - F(0) = F'(\theta)
\]
であり、\(0 < \theta < 1\)。しかし、\SecNo[§]{153} より
\begin{align*}
F'(t) &= D_{t} f(x + ht, y + kt)\\
      &= hf_{x}'(x + ht, y + kt) + kf_{y}'(x + ht, y + kt).
\end{align*}
したがって、最終的に
\[
\delta z = f(x + h, y + k) - f(x, y)
  = hf_{x}'(x + \theta h, y + \theta k) + kf_{y}'(x + \theta h, y + \theta k)
\]
であり、これが求めた公式である。\(f_{x}'\)、\(f_{y}'\) が \(x\)、\(y\) の連続関数であると仮定すると
\begin{align*}
f_{x}'(x + \theta h, y + \theta k) &= f_{x}'(x, y) + \epsilon_{h, k}\\
f_{y}'(x + \theta h, y + \theta k) &= f_{y}'(x, y) + \eta_{h, k}
\end{align*}
であり、ここで \(\epsilon_{h, k}\) および \(\eta_{h, k}\) は \(h\)、\(k\) が零に近づくとき零に収束する。したがって、定理は
\[
\delta z = (f_{x}' + \epsilon) \, \delta x + (f_{y}' + \eta) \, \delta y
\]
と書き換えることができる。ここで \(\epsilon\) および \(\eta\) は \(\delta x\)、\(\delta y\) が零に近づくとき小さい値である。

\[
\delta z = f_{x}'\, \delta x + f_{y}'\, \delta y
\]
は近似的に成り立つと言えます；すなわち、両辺の差は$\delta x$と$\delta y$の大きい方と比較して小さいです。\footnote{あるいは$|\delta x| + |\delta y|$または$\sqrt{\delta x^{2} + \delta y^{2}}$と比較しても同じです。}
一方が他方と比べて小さい可能性があるため、私たちは「$\delta x$と$\delta y$の大きい方」を言明しなければなりません；実際には$\delta x = 0$または$\delta y = 0$であることもあります。

以下は注記ですが、任意の形で$\delta z = \lambda\, \delta x + \mu\, \delta y$がこの意味で近似的に成り立つ場合、$\lambda = f_{x}'$, $\mu = f_{y}'$であることが必要です。なぜなら
\[
\delta z - f_{x}'\, \delta x - f_{y}'\, \delta y
  = \epsilon\, \delta x + \eta\, \delta y,\quad
\delta z - \lambda\, \delta x - \mu\, \delta y
  = \epsilon'\, \delta x + \eta'\, \delta y
\]
であり、$\epsilon$, $\eta$, $\epsilon'$, $\eta'$は$\delta x$と$\delta y$が0に近づくにつれて0に収束します。したがって
\[
(\lambda - f_{x}')\, \delta x + (\mu - f_{y}')\, \delta y
  = \rho\, \delta x + \rho'\, \delta y
\]
であり、$\rho$と$\rho'$も0に収束します。したがって任意の正の数$\zeta$に対して、$\sigma$を適切に選べば
\[
|(\lambda - f_{x}')\, \delta x + (\mu - f_{y}')\, \delta y|
  < \zeta(|\delta x| + |\delta y|)
\]
が全ての数値的に$\sigma$未満の$\delta x$と$\delta y$に対して成り立ちます。$\delta y = 0$を取ると
\[
|(\lambda - f_{x}')\, \delta x|
  < \zeta|\delta x|
\]
となり、$|\lambda - f_{x}'| < \zeta$であり、$\zeta$は任意に小さくできるため、$\lambda = f_{x}'$であることが必要です。同様に$\mu = f_{y}'$も成り立ちます。

\paragraph{155. 微分}微積分の応用、特に幾何学においては、関数 \( x \), \( y \), \( z \) の増分 \( \delta x \), \( \delta y \), \( \delta z \) ではなく、それらの微分 \( dx \), \( dy \), \( dz \) を用いた方程式を扱うことが最も便利であることが多い。

単変数 \( x \) の関数 \( y = f(x) \) について考えてみよう。もし \( f'(x) \) が連続であれば
\[
\delta y = \{f'(x) + \epsilon\} \, \delta x,
\Tag{(1)}
\]
ここで \( \epsilon \to 0 \) となるとき \( \delta x \to 0 \): すなわち方程式
\[
\delta y = f'(x) \, \delta x
\Tag{(2)}
\]
は「近似して」真である。これまで私たちは単独で表されるシンボル \( dy \) の意味を任何のものとして付与してこなかったが、今では方程式
\[
dy = f'(x) \, \delta x.
\Tag{(3)}
\]
によって \( dy \) を定義することに同意する。

\( y \) に特異な関数 \( x \) を選ぶと
\[
dx = \delta x,
\Tag{(4)}
\]
となり、従って
\[
dy = f'(x) \, dx.
\Tag{(5)}
\]
両辺を \( dx \) で割ると
\[
\frac{dy}{dx} = f'(x),
\Tag{(6)}
\]
ここで \( dy/dx \) はこれまでの定義とは異なり、\( y \) の微分係数ではなく、微分 \( dy \), \( dx \) の商を表す。したがってシンボル \( dy/dx \) には二つの意味が与えられるが、これは方程式
\[
\frac{dy}{dx} = f'(x)
\Tag{(6)}
\]
がどの定義を選んでも真であるため、何らの不便も生じない。

以下は、原著英文を忠実に翻訳したものです：

\par\medskip\noindent\textbf{【注釈】}\par
方程式（5）には方程式（2）よりも二つの明確な利点があります。それは正確であるだけでなく、近似ではなく、その真偽性は \(f'(x)\) の連続性に関する仮定に依存しないことです。一方で、特定の条件下で正確な方程式（5）から近似的な方程式（2）へと移行できるという事実が前者の重要性を生み出します。「微分」記法の利点は実質的に技術的なものに過ぎません。しかし、特に複数変数関数の場合においてこれらの技術的優位性は非常に大きく、この記法を使用することはほぼ避けられません。

\(f'(x)\) が連続であるとき、
\[
\lim \frac{dy}{\delta y} = 1
\]
が \(\delta x \to 0\) のときに成り立ちます。これはしばしば \(\delta x\) が小さいときに \(dy\) は \(\delta y\) の主部であるという表現で述べ occasionally expressed by saying that \(dy\) is the *principal part* of \(\delta y\) when \(\delta x\) is small, と解釈されます。また、\(ax\) は \(x\) が小さいときに \(ax + bx^2\) の主部であるというように考えることができます。
\par\medskip
\PageSep{281}

次に、独立変数 \(x\) および \(y\) の関数 \(z\) に関連する定義について述べます。\(dz\) を方程式
\[
dz = f_{x}' \, \delta x + f_{y}' \, \delta y.
\Tag{(7)}
\]
によって定義します。

次に \(z = x\) および \(z = y\) の順に置き換えると、
\begin{align*}
dx &= \delta x,\quad
dy = \delta y,
\Tag{(8)}
\intertext{よって}
dz &= f_{x}' \, dx + f_{y}' \, dy,
\Tag{(9)}
\end{align*}
が得られます。これは近似方程式（154）節の方程式（1）に対応する正確な方程式です。また、ここでも前者は実用上の利便性のために重要であり、後者が特定の状況で導き出せるという理由から重要であることに注意しなければなりません。
\par\medskip

以下は、原著英文を忠実に翻訳したものです：

\par\medskip\noindent\textbf{【注釈】}\par
方程式（9）の性質について特別な注目が必要です。第153節で述べたように、$z = f(x, y)$ であり、$x$ と $y$ は単一の変数 $t$ の関数であるため独立ではなく、したがって $z$ も $t$ のみの関数となるとき、

\[
\frac{dz}{dt} = \frac{\partial f}{\partial x}\, \frac{dx}{dt} + \frac{\partial f}{\partial y}\, \frac{dy}{dt}
\]

が成り立ちます。この式を $dt$ で掛けて、

\[
dx = \frac{dx}{dt}\, dt,\quad
dy = \frac{dy}{dt}\, dt,\quad
dz = \frac{dz}{dt}\, dt,
\]

から

\[
dz = f_{x}'\, dx + f_{y}'\, dy
\]

を得ます。これは方程式（9）と同じ形式です。したがって、変数 $x$ と $y$ が独立であるかどうかに関わらず、$dx$ および $dy$ を用いて $dz$ を表す公式は同じです。この注目すべき点は応用において非常に重要です。

また、$z$ が二つの独立変数 $x$ と $y$ の関数であり、

\[
dz = \lambda\, dx + \mu\, dy
\]

のとき、$\lambda = f_{x}'$, $\mu = f_{y}'$ であることも注意すべきです。これは第154節の最後の段落からすぐに導けます。

前3節の定理と定義は、任意数目的変数に対する関数にも即座に拡張可能であることが明らかです。
\par\medskip

\par\medskip\noindent\textbf{【例 LXII.】}\par
\Item{1.} 楕円の面積 $A$ は $A = \pi ab$ で与えられ、$a$, $b$ は半長軸と半短軸を表します。$A$ の増分 $dA$，$a$ の増分 $da$，および $b$ の増分 $db$ の関係を証明し、

\[
\frac{dA}{A} = \frac{da}{a} + \frac{db}{b}
\]

を示すとともに、軸の増分と面積の増分との間の近似方程式を求めよ。
\par\medskip

2. 三角形 \(ABC\) の面積 \(\Delta\) を (i) \(a\), \(B\), \(C\), (ii) \(A\), \(b\), \(c\), (iii) \(a\), \(b\), \(c\) の関数として表し、以下の公式を確立します：
\[
\frac{d\Delta}{\Delta} = 2\frac{da}{a} + \frac{c\, dB}{a\sin B} + \frac{b\, dC}{a\sin C},\quad
\frac{d\Delta}{\Delta} = \cot A\, dA + \frac{db}{b} + \frac{dc}{c},
\]
\[
d\Delta = R(\cos A\, da + \cos B\, db + \cos C\, dc),
\]
ここで \(R\) は外接円の半径です。

3. 三角形の辺が変化し、面積が一定であるとき、\(a\) を \(b\) と \(c\) の関数として考えることができます。次を証明します：
\[
\frac{\dd a}{\dd b} = -\frac{\cos B}{\cos A},\quad
\frac{\dd a}{\dd c} = -\frac{\cos C}{\cos A}.
\]

これらは以下の式から導けます：
\[
da = \frac{\dd a}{\dd b}\, db + \frac{\dd a}{\dd c}\, dc,\quad
\cos A\, da + \cos B\, db + \cos C\, dc = 0.
\]

4. \(R\) が一定であるとき、\(a\), \(b\), \(c\) の変化により次の式が成り立ちます：
\[
\frac{da}{\cos A} + \frac{db}{\cos B} + \frac{dc}{\cos C} = 0,
\]
そして
\[
\frac{\dd a}{\dd b} = -\frac{\cos A}{\cos B},\quad
\frac{\dd a}{\dd c} = -\frac{\cos A}{\cos C}.
\]

これらは \(a = 2R\sin A\) 等の公式と \(R\) および \(A + B + C\) の定常性を利用することで証明できます。

5. \(z\) が \(u\) と \(v\) の関数であり、これらの \(x\) と \(y\) の関数であるとき、次の式が成り立ちます：
\[
\frac{\dd z}{\dd x} = \frac{\dd z}{\dd u}\, \frac{\dd u}{\dd x}
                    + \frac{\dd z}{\dd v}\, \frac{\dd v}{\dd x},\quad
\frac{\dd z}{\dd y} = \frac{\dd z}{\dd u}\, \frac{\dd u}{\dd y}
                    + \frac{\dd z}{\dd v}\, \frac{\dd v}{\dd y}.
\]

これらは以下の式から導けます：
\[
dz = \frac{\dd z}{\dd u}\, du + \frac{\dd z}{\dd v}\, dv,\quad
du = \frac{\dd u}{\dd x}\, dx + \frac{\dd u}{\dd y}\, dy,\quad
dv = \frac{\dd v}{\dd x}\, dx + \frac{\dd v}{\dd y}\, dy.
\]
\(du\) と \(dv\) を最初の式に代入し、結果を次の式と比較します：
\[
dz = \frac{\dd z}{\dd x}\, dx + \frac{\dd z}{\dd y}\, dy.

\item{6.} \( z \) を \( x \) および \( y \) の関数とするとき、\( X \), \( Y \), \( Z \) を以下の式で定義する：
\[
x = a_{1}X + b_{1}Y + c_{1}Z,\quad
y = a_{2}X + b_{2}Y + c_{2}Z,\quad
z = a_{3}X + b_{3}Y + c_{3}Z.
\]
このとき、\( Z \) は \( X \) および \( Y \) の関数として表すことができる。また、\(\frac{\partial Z}{\partial X}\), \(\frac{\partial Z}{\partial Y}\) を \(\frac{\partial z}{\partial x}\), \(\frac{\partial z}{\partial y}\) で表す。これらの微分係数を \( P \), \( Q \) および \( p \), \( q \) とする。このとき、
\[
dz - p\,dx - q\,dy = 0,
\]
または
\[
(c_{1}p + c_{2}q - c_{3})dZ + (a_{1}p + a_{2}q - a_{3})dX + (b_{1}p + b_{2}q - b_{3})dY = 0.
\]
この式を \( dZ - P\,dX - Q\,dY = 0 \) と比較すると、
\[
P = -\frac{a_{1}p + a_{2}q - a_{3}}{c_{1}p + c_{2}q - c_{3}},\quad
Q = -\frac{b_{1}p + b_{2}q - b_{3}}{c_{1}p + c_{2}q - c_{3}}.
\]

\item{7.} 以下の式が成り立つとき、
\[
(a_{1}x + b_{1}y + c_{1}z)p + (a_{2}x + b_{2}y + c_{2}z)q = a_{3}x + b_{3}y + c_{3}z,
\]
次に、
\[
(a_{1}X + b_{1}Y + c_{1}Z)P + (a_{2}X + b_{2}Y + c_{2}Z)Q = a_{3}X + b_{3}Y + c_{3}Z.
\]

\item{8.} \Topic{暗黙の関数の微分法。} \( f(x, y) \) およびその導函数 \( f_{y}'(x, y) \) が点 \((a, b)\) の近傍で連続であり、かつ
\[
f(a, b) = 0,\quad
f_{b}'(a, b) \neq 0,
\]
であるとき、\( f_{y}'(x, y) \) は点 \((a, b)\) の近傍で常に同一の符号を持つような近傍が存在する。例えば、\( f_{y}'(x, y) \) が点 \((a, b)\) の近傍で正であると仮定すると、任意の \( x \) の値が \( a \) に十分近いときおよび任意の \( y \) の値が \( b \) に十分近いとき、\( f(x, y) \) はより狭義の意味で \( y \) の増加関数である。このことから、\SecNo[§]{108} の定理により、点 \( x = a \) において \( y = b \) となる唯一の連続な関数 \( y \) が存在し、かつすべての \( x \) の値が \( a \) に十分近いとき、方程式 \( f(x, y) = 0 \) を満たす。

Let us now suppose \( f(x, y) \) possesses a derivative \( f_{x}'(x, y) \) which is also continuous near \( (a, b) \). If \( f(x, y) = 0 \), \( x = a + h \), \( y = b + k \), we have
\[
0 = f(x, y) - f(a, b) = (f_{a}' + \epsilon) h + (f_{b}' + \eta) k,
\]
where \( \epsilon \) and \( \eta \) tend to zero with \( h \) and \( k \). Thus
\[
\frac{k}{h} = -\frac{f_{a}' + \epsilon}{f_{b}' + \eta} \to -\frac{f_{a}'}{f_{b}'},
\]
or
\[
\frac{\mathrm{d}y}{\mathrm{d}x} = -\frac{f_{a}'}{f_{b}'}.
\]

9. The equation of the tangent to the curve \( f(x, y) = 0 \), at the point \( x_0 \), \( y_0 \), is
\[
(x - x_0) f_{x_0}'(x_0, y_0) + (y - y_0) f_{y_0}'(x_0, y_0) = 0.
\]

156. 確定積分と面積。以前に、第VI章 §145で、\( f(x) \)が \( x \)の連続関数であるとき、\( PQ \)が \( y = f(x) \)のグラフであるとき、図47に示す領域 \( PpqQ \)には、その面積とみなされる固有の数値があることを想定した。明らかに、\( Op \)を \( a \)、\( Oq \)を \( x \)とするとき、\( x \)が変化するにつれてこの面積は \( x \)の関数となり、これを \( F(x) \)と表すことができる。

この想定に基づき、§145で \( F'(x) = f(x) \)を証明し、この結果を利用して特定の曲線の面積を計算する方法を示した。しかし、まだその基本的な前提である、あるような数値としての面積 \( F(x) \)が存在することを正当化する必要がある。

私たちは確かに長方形の面積が何を意味するかを知っていますし、その面積は辺の積で測定されます。また、ユークリッドによって証明された三角形や平行四辺形、多角形の性質により、これらの図形の面積にも特定の意味を付与することができます。しかし、これまでに知っている限りでは、曲線で囲まれた図形の面積に対する直接的な定義は提供されていません。現在では、$F(x)$ の定義を与える方法を示し、その存在性を証明することを説明します。\footnote{以下の議論は、ゴルサールの『解析教程』（第2版）vol.~i, pp.~171~\textit{et~seq.} に模倣されていますが、ゴルサールの手法は格段に一般化されています。}

区間 $\DPmod{(a, b)}{[a, b]}$ 内で連続である $f(x)$ を仮定し、分割点 $x_{0}$,~$x_{1}$, $x_{2}$,~\dots, $x_{n}$ によって区間をいくつかのサブ区間に分けてみましょう。ここで
\[
a = x_{0} < x_{1} < \dots < x_{n-1} < x_{n} = b.
\]
さらに、$\delta_{\nu}$ を $\DPmod{(x_{\nu}, x_{\nu+1})}{[x_{\nu}, x_{\nu+1}]}$ とし、$\delta_{\nu}$ 内の $f(x)$ の下界を $m_{\nu}$ とします。そして
\[
s = m_{0}\delta_{0} + m_{1}\delta_{1} + \dots + m_{n}\delta_{n}
  = \sum m_{\nu}\delta_{\nu},
\]
と表記します。

明らかに、$f(x)$ の区間 $\DPmod{(a, b)}{[a, b]}$ 内の上界 $M$ が存在すれば、$s \leq M(b - a)$ となります。したがって、$s$ の値全体は第80節の言葉で言えば上界を持つ集合であり、その上界を $j$ とします。$s$ のどの値も $j$ より大きいものはありませんが、$j$ より小さい任意の数よりも大きい $s$ の値があります。

同様に、$\delta_{\nu}$ 内の $f(x)$ の上界を $M_{\nu}$ とすれば、和
\[
S = \sum M_{\nu}\delta_{\nu}
\]
を定義することができます。

{\Loosen 明らかに、$f(x)$ 在区间 $(a, b)$ 中的下界為 $m$ 時，有 $S \geq m(b - a)$. 因此，$S$ 的总和是有界的，并且具有一个下界，我達将其记作 $J$.
$S$ 的任何值都不小于 $J$，但存在一些 $S$ 的值，它達小于任意大于 $J$ 的数.}
\PageSep{285}

\par\medskip\noindent\textbf{【注釈】}\par
$s$ と $S$ の意味を明確にするために、$f(x)$ が $x = a$ から $x = b$ まで一貫して増加する単純な場合を考えると、$m_{\nu}$ は $f(x_{\nu})$ であり、$M_{\nu}$ は $f(x_{\nu+1})$. この場合、$s$ は図48に影を描いた長方形の総面積であり、$S$ は太線で囲まれた領域の面積です。一般の場合でも、$s$ と $S$ は依然として面積であり、それぞれ曲線領域を含むまたは含む長方形で構成されています。
\par\medskip

現在、任意の $s$ の和が任意の $S$ の和を超えることはないことを示します。$s$, $S$ は一つの分割方法に対応する和であり、$s'$, $S'$ は別の分割方法に対応する和です。それらについて $s \leq S'$ および $s' \leq S$ を示す必要があります.

私たちは、$s$,~$S$ または $s'$,~$S'$ のいずれかに対して区切り点となるすべての点を用いて三分割を行う第三の方法を作成できます。この第三の方法による分割に対応する和を $\mathbf{s}$,~$\mathbf{S}$ とします。すると、容易に確認できるように、

\[
\mathbf{s} \geq s,\quad \mathbf{s} \geq s',\quad
\mathbf{S} \leq S,\quad \mathbf{S} \leq S'.
\Tag{(1)}
\]

例えば、$\mathbf{s}$ は $s$ と異なる点在于少なくとも一つの区間 $\delta_{\nu}$ が $s$ の中に含まれており、この区間 $\delta_{\nu}$ はより小さな区間に分割され、

\[
\delta_{\nu, 1},\ \delta_{\nu, 2},\ \dots,\ \delta_{\nu, p},
\]

となるため、$s$ の項 $m_{\nu}\delta_{\nu}$ が $\mathbf{s}$ の中に含まれる和

\[
m_{\nu, 1}\delta_{\nu, 1} + m_{\nu, 2}\delta_{\nu, 2} + \dots + m_{\nu, p}\delta_{\nu, p},
\]

に置き換えられます。ここで $m_{\nu, 1}$, $m_{\nu, 2}$,~\dots は $\delta_{\nu, 1}$, $\delta_{\nu, 2}$,~\dots の中で $f(x)$ の下限です。しかし、明らかに $m_{\nu, 1} \geq m_{\nu}$, $m_{\nu, 2} \geq m_{\nu}$,~\dots であるため、その和は $m_{\nu}\delta_{\nu}$ よりも大きいことがわかります。したがって $\mathbf{s} \geq s;$ そして他の不等式も同様の方法で証明できます。しかし、$\mathbf{s} \leq \mathbf{S}$ であるため、

\[
s \leq \mathbf{s} \leq \mathbf{S} \leq S',
\]

が成り立ち、これが私たちが証明したいことでした。

また、$j \leq J$ も同様に導けます。なぜなら、$j$ に近い $s$ と $J$ に近い $S$ を任意に選べるため、

\[
j > J
\]

が成り立つとすると、ある $s$ と $S$ の存在を意味し、これは矛盾します。

So 迄まで、$f(x)$ の連続性を利用したことはありません。
今後、$j = J$ であることを示し、分割点$x_{\nu}$を無限に増やして各区間$\delta_{\nu}$がゼロに収束するとき、和$s$, $S$が極限$J$に収束することを示します。より正確には、任意の正数$\epsilon$に対して、ある$\delta$を選び得るような条件で
\[
0 \leq J - s < \epsilon,\quad 0 \leq S - J < \epsilon
\]
が成り立つことを示します。

定理II（§106）によれば、ある数$\delta$が存在し、
\[
M_{\nu} - m_{\nu} < \frac{\epsilon}{b - a}
\]
が全ての$\delta_{\nu}$が$\delta$未満であるとき成立します。したがって
\[
S - s = \sum (M_{\nu} - m_{\nu})\, \delta_{\nu} < \epsilon.
\]
また
\[
S - s = (S - J) + (J - j) + (j - s)
\]
であり、右辺の3つの項は全て正で、従ってそれらは全て$\epsilon$未満です。$J - j$は定数であるため、これはゼロであることになります。したがって $j = J$ であり、$0 \leq j - s < \epsilon$, $0 \leq S - J < \epsilon$ が成立します。

領域$PpqQ$の面積を$s$と$S$の共通極限、すなわち$J$として定義します。この定義はより一般的なものに拡張できます。任意の$\delta_{\nu}$内の点$x$における$f(x)$の値を$f_{\nu}$とし、和
\[
\sigma = \sum f_{\nu}\delta_{\nu}
\]
を考えます。明らかに$\sigma$は$s$と$S$の間にあるため、区間$\delta_{\nu}$がゼロに収束するとき極限$J$に向かいます。したがって面積を$\sigma$の極限として定義できます。

157. 確定積分．今後$f(x)$を連続関数とし、曲線$y = f(x)$、直線$x = a$および$x = b$、$x$軸で囲まれた領域が確定した面積を持つとする。第VI章 §145において我々は、$f(x)$の「積分関数」$F(x)$が
\[
F'(x) = f(x),\quad F(x) = \int f(x)\, dx,
\]
を満たすとき、その面積は$F(b) - F(a)$であることを証明しました。

As 常に実際的に $F(x)$ の形を決定することが困難であるため、領域 $PpqQ$ を表す式として $F(x)$ を明示的に参照しない公式を持つことは便利である。以下のように書くことにする：
\[
(PpqQ) = \int_{a}^{b} f(x)\, dx.
\]
次頁

この等式の右辺の表現は、二つの方法のいずれかで定義されるものと見做すことができる。$F(x)$ が $f(x)$ の積分関数であり、具体的な公式が知られているかどうかに関わらず、$F(b) - F(a)$ の簡略化形として見るか、もしくは直接定義された領域 $PpqQ$ の値として見るかのいずれかである。また、この領域を第156節で直接定義したものと見做すこともできる。

数
\[
\int_{a}^{b} f(x)\, dx
\]
は「確定積分」と呼ばれる。$a$ および $b$ はその「下限」および「上限」である。$f(x)$ は「積分区間の主題」または「被積関数」と呼ばれる。そして、区間 $\DPmod{(a, b)}{[a, b]}$ は「積分範囲」と呼ばれる。確定積分は $a$ および $b$ と関数 $f(x)$ の形に依存し、$x$ の関数ではない。一方で、不定積分
\[
F(x) = \int f(x)\, dx
\]
はしばしば $f(x)$ の「不定積分」と呼ばれる。

補足：
確定積分と不定積分の区別は視点の違いに過ぎない。確定積分 $\ds\int_{a}^{b} f(x)\, dx = F(b) - F(a)$ は $b$ の関数であり、特定の $f(b)$ の積分関数として見做すことができる。一方で、不定積分 $F(x)$ は常に確定積分を用いて表すことができ、
\[
F(x) = F(a) + \int_{a}^{x} f(t)\, dt.
\]

ただし、不定積分や積分関数を考えるときには、通常、ある関係性を考慮しています。その関係性により一つの関数が別の関数の導関数であることがわかります。そして、定積分を考えるときには、一般的に区間の変動は関心の対象ではありません。通常、区間は定数であり、例えば0と1です。そして
\[
\int_{0}^{1} f(x)\, dx = F(1) - F(0)
\]
は関数ではなく単なる数値です。

$\ds\int_{a}^{x} f(t)\, dt$ という積分が微分係数$f(x)$を持つことから、この積分はより一層$x$の連続な関数であることがわかります。

$x$ のすべての正の値において $1/x$ が連続であることから、前段落での調査により $\log x$ の実際の存在が証明されます。私たちは第128節で仮定的にその関数を前提としていました。

158. 圓の弧の面積。三角関数。

初等三角法の教科書で通常提示される三角関数 \(\cos x\), \(\sin x\) などの理論は、未証明の前提に依存している。角度 \(OA\), \(OP\) の配置からなる構造体であるという「幾何学的」定義は特に問題ない。しかし次に進むと、角度が数値測定可能であるとする仮定が必要となる。つまり、ある実数 \(x\) がその構造体に関連しているとされる。図47の領域 \(PpqQ\) も同様に実数に関連している。

この前提を受け入れた上で \(\cos x\) と \(\sin x\) を通常通り定義し、理論の展開において原則的な問題は発生しない。しかし全体として問題となるのは、\(\cos x\) および \(\sin x\) における \(x\) とは何かという点である。この質問に答えるためには角度の測度を定義する必要があり、現在その定義を行う立場にある。

最も自然な定義は次の通りである：点 \(O\) を中心とする半径が1の円上の弧 \(AP\) の長さを \(x\) とすると、角度の測度 \(x\) は弧 \(AP\) の長さに等しい。教科書では「円周率」の理論について述べる際にこの定義を採用している。しかし現在の目的には致命的な欠陥がある；すなわち、曲線の長さ、特に円の弧の長さが証明されていないからである。

曲線の長さという概念は面積と同様に精密な数学的分析が可能であり、その分析は前節と同様の一般的な性格を持つが、明らかに難度が高い。ここではこの主題についての一般的な扱いを提供することは不可能である。

We therefore define the measure of the angle \( AOP \) on the notion not of length but of \emph{area}. We define the measure of the angle \( AOP \) as \emph{twice the area of the sector \( AOP \) of the unit circle}.

Suppose, in particular, that \( OA \) is \( y = 0 \) and that \( OP \) is \( y = mx \), where \( m > 0 \). The area is a function of \( m \), which we may denote by \( \phi(m) \). If we write \( \mu \) for \( (1 + m^{2})^{-\frac{1}{2}} \), then \( P \) is the point \( (\mu, m\mu) \), and
\[
\phi(m) = \tfrac{1}{2} m\mu^{2} + \int_{\mu}^{1} \sqrt{1 - x^{2}}\, dx.
\]
Differentiating with respect to \( m \), we find
\[
\phi'(m) = \frac{1}{2(1 + m^{2})},\quad
\phi(m) = \tfrac{1}{2} \int_{0}^{m} \frac{dt}{1 + t^{2}}.
\]
Thus the analytical equivalent of our definition would be to define \( \arctan m \) by the equation
\[
\arctan m = \int_{0}^{m} \frac{dt}{1 + t^{2}};
\]
and the whole theory of the circular functions could be worked out from this starting point, just as the theory of the logarithm is worked out from a similar definition in Chapter IX. See Appendix III.

\par\medskip\noindent\textbf{【例 LXIII.】}\par \Topic{Calculation of the definite integral from the indefinite integral.}
\Item{1.} Show that
\[
\int_{a}^{b} x^{n}\, dx = \frac{b^{n+1} - a^{n+1}}{n + 1},
\]
and in particular that
\[
\int_{0}^{1} x^{n}\, dx = \frac{1}{n + 1}.
\]

\Item{2.} \(\ds\int_{a}^{b} \cos mx\, dx = \frac{\sin mb - \sin ma}{m}\),
\(\ds\int_{a}^{b} \sin mx\, dx = \frac{\cos ma - \cos mb}{m}\).

\Item{3.} \(\ds\int_{a}^{b}\frac{dx}{1 + x^{2}} = \arctan b - \arctan a\),
\(\ds\int_{0}^{1}\frac{dx}{1 + x^{2}} = \tfrac{1}{4}\pi\).

以下は翻訳です：

[ここに明らかな困難があるのは、\(\arctan x\)が多価関数であるという事実によるものです。この困難を避けるためには、以下の等式において
\[
\int_{0}^{x} \frac{dt}{1 + t^{2}} = \arctan x,
\]
\(\arctan x\)は\(-\frac{1}{2}\pi\)と\(\frac{1}{2}\pi\)の間にある角であることを観察する必要があります。積分は\(x = 0\)のとき消失し、\(x\)が増加するにつれて定常かつ連続に増加します。したがって、同じことが\(\arctan x\)にも言え、従って\(x \to \infty\)のとき\(\arctan x \to \frac{1}{2}\pi\)となります。同様にして、\(x \to -\infty\)のとき\(\arctan x \to -\frac{1}{2}\pi\)を示すことができます。

また、以下の等式においても同様です：
\[
\int_{0}^{x} \frac{dt}{\sqrt{1 - t^{2}}} = \arcsin x,
\]
ここで\(-1 < x < 1\)、\(\arcsin x\)は\(-\frac{1}{2}\pi\)と\(\frac{1}{2}\pi\)の間にある角を表します。したがって、\(a\)と\(b\)がともに単位より小さいとき、
\[
\int_{a}^{b} \frac{dx}{\sqrt{1 - x^{2}}} = \arcsin b - \arcsin a.
\]

4. \(\ds\int_{0}^{1} \frac{dx}{1 - x + x^{2}} = \frac{2\pi}{3\sqrt{3}}\)、\(\ds\int_{0}^{1} \frac{dx}{1 + x + x^{2}} = \frac{\pi}{3\sqrt{3}}\).

5. \(\ds\int_{0}^{1} \frac{dx}{1 + 2x\cos\alpha + x^{2}} = \frac{\alpha}{2\sin\alpha}\) は\(-\pi < \alpha < \pi\)において成り立ち、ただし\(\alpha = 0\)のとき積分値は\(\frac{1}{2}\)であり、これは\(\alpha \to 0\)のとき\(\frac{1}{2}\alpha\cosec\alpha\)の極限です。

6. \(\ds\int_{0}^{1} \sqrt{1 - x^{2}}\, dx = \frac{1}{4}\pi\)、\(\ds\int_{0}^{a} \sqrt{a^{2} - x^{2}}\, dx = \frac{1}{4}\pi a^{2}\) は\(a > 0\)のとき成り立つ。

7. \(\ds\int_{0}^{\pi} \frac{dx}{a + b\cos x} = \frac{\pi}{\sqrt{a^{2} - b^{2}}}\) は \(a > |b|\) のとき成り立ちます。[不定積分の形については、\Exs{liii}.3,4を参照してください。また、\(|a| < |b|\)の場合、被積分関数には\(0\)と\(\pi\)の間にある無限大があります。\(a\)が負で\(-a > |b|\)であるときの積分値はどのようなものでしょうか？]

\item 8. \( \ds\int_{0}^{\frac{1}{2}\pi} \frac{dx}{a^{2}\cos^{2}x + b^{2}\sin^{2}x} = \frac{\pi}{2ab} \), もし \( a \) と \( b \) が正の数であるとき。\( a \) と \( b \) の符号が異なるか、または両方とも負の数であるとき、この積分の値は何か？

\item 9. フーリエ積分について。もし \( m \) と \( n \) が正の整数であるなら
\[
\int_{0}^{2\pi} \cos mx \sin nx\, dx
\]
は常にゼロであり、
\[
\int_{0}^{2\pi} \cos mx \cos nx\, dx,\quad
\int_{0}^{2\pi} \sin mx \sin nx\, dx
\]
は \( m = n \) でない限りゼロであり、\( m = n \) のとき各々は \( \pi \) に等しい。

\item 10. \( \ds\int_{0}^{\pi} \cos mx \cos nx\, dx \) と \( \ds\int_{0}^{\pi} \sin mx \sin nx\, dx \) はそれぞれ \( m = n \) でない限りゼロであり、\( m = n \) のとき各々は \( \frac{1}{2}\pi \) に等しい。また
\[
\int_{0}^{\pi} \cos mx \sin nx\, dx = \frac{2n}{n^{2} - m^{2}},\quad
\int_{0}^{\pi} \cos mx \sin nx\, dx = 0,
\]
は \( n - m \) が奇数か偶数である場合に応じて成立する。

\paragraph{159. 定積分の計算から定義への推算。}少数の場合、定積分を直接計算して、\SecNo[§§]{156} と \SecNo{157} の定義から始めることができます。通常は不定積分を使用するのがずっと簡単ですが、読者はいくつかの例を通じて作業を行うことで学ぶことが有益です。

\par\medskip\noindent\textbf{【例 LXIV.】}\par
\item 1. \( \ds\int_{a}^{b} x\, dx \) を計算するために、区間 \([a, b]\) を点 \( a = x_{0} \), \( x_{1} \), \( x_{2} \), \dots, \( x_{n} = b \) で等分し、分割数 \( n \to \infty \) のときの極限
\[
(x_{1} - x_{0})f(x_{0}) + (x_{2} - x_{1})f(x_{1}) + \dots + (x_{n} - x_{n-1})f(x_{n-1})
\]
を計算する。

この和は次のようになります。
\[
\frac{b - a}{n}\left[
  a + \left(a + \frac{b - a}{n}\right) + \left(a + 2\frac{b - a}{n}\right)
    + \dots + \left\{a + (n - 1)\frac{b - a}{n}\right\}
  \right]
  = \frac{b - a}{n}\left[
    na + \frac{b - a}{n} \{1 + 2 + \dots + (n - 1)\}
  \right]
  = (b - a)\left\{a + (b - a)\frac{n(n - 1)}{2n^{2}}\right\},
\]
これは$n \to \infty$のとき$\frac{1}{2} (b^{2} - a^{2})$に収束します。グラフィカルな理由で結果を確認してください。

2. $x^2$の積分$\ds\int_{a}^{b} x^{2}\, dx$も同様にして計算します。

3. $0 < a < b$の場合、$x$の積分$\ds\int_{a}^{b} x\, dx$を$a$, $ar$, $ar^2$, $\dots$, $ar^{n-1}$, $ar^n$で区切った$a$から$b$までの$n$分割によって計算します。$r^n = b/a$です。同様の方法で、より一般的な積分$\ds\int_{a}^{b} x^{m}\, dx$も計算します。

4. $\cos mx$と$\sin mx$の積分$\ds\int_{a}^{b}\cos mx\, dx$および$\ds\int_{a}^{b}\sin mx\, dx$は例題1の方法で計算します。

5. $n\sum\limits_{r=0}^{n-1} \dfrac{1}{n^{2} + r^{2}} \to \tfrac{1}{4}\pi$が$n \to \infty$のとき成立することを証明します。
これは次の事実から導けます。
\[
\frac{n}{n^{2}} + \frac{n}{n^{2} + 1^{2}} + \dots + \frac{n}{n^{2} + (n - 1)^{2}}
  = \sum_{r=0}^{n-1} \frac{(1/n)}{1 + (r/n)^{2}},
\]
これは$n \to \infty$のとき$\ds\int_{0}^{1} \frac{dx}{1 + x^{2}}$に収束します。積分の直接的な定義によります。

6. $\dfrac{1}{n^{2}} \sum\limits_{r=0}^{n-1} \sqrtp{n^{2} - r^{2}} \to \tfrac{1}{4}\pi$が成立することを証明します。
この極限は$\ds\int_{0}^{1} \sqrtp{1 - x^{2}}\, dx$です。

定数積分の一般的な性質。定数積分には、以下の式で表される重要な性質があります。\footnote
  {これらの式に登場するすべての関数は当然連続であり、定数積分は連続な関数に対してのみ定義されています。}
\[
(1)\quad \ds\int_{a}^{b} f(x)\, dx = -\int_{b}^{a} f(x)\, dx.
\]

このことは、積分関数 \( F(x) \) を用いて積分の定義から即座に導けます。すなわち \( F(b) - F(a) = -\{F(a) - F(b)\} \) です。直接的な定義においては下限が上限より大きいことを前提としていましたので、この方法では \( a < b \) のときの積分 \(\ds\int_{b}^{a} f(x)\, dx\) には適用できません。この定義を基本とすれば、そのような場合にもこれを拡張して等式 (1) を右辺の定義として採用する必要があります。

\centerline{\item{(2)}} \(\ds\int_{a}^{a} f(x)\, dx = 0\)。
\centerline{\item{(3)}} \(\ds\int_{a}^{b}f(x)\, dx + \int_{b}^{c}f(x)\, dx = \int_{a}^{c}f(x)\, dx\)。
\centerline{\item{(4)}} \(\ds\int_{a}^{b}kf(x)\, dx = k \int_{a}^{b}f(x)\, dx\)。
\centerline{\item{(5)}} \(\ds\int_{a}^{b}\{f(x) + \phi(x)\}\, dx
= \int_{a}^{b}f(x)\, dx + \int_{a}^{b}\phi(x)\, dx\)。

注記
読者は、これらの性質の形式的な証明を書き出すことを試みることで有益な練習となるでしょう。それぞれの場合においては、(α) 積分関数による定義と (β) 直接的定義から始める証明を与えます。

次の定理も重要なものです。

結果
\centerline{\item{(6)}} \( f(x) \geq 0 \) のとき \( a \leq x \leq b \)，則 \(\ds\int_{a}^{b}f(x)\, dx \geq 0\)。

注記
和 \( s \)（156節参照）が負になることはないことを確認すればよいです。後で示される（VII 節41問）では、積分の値が零であるためには \( f(x) \) が常に零と等しいことが必要であり、これは121節の第二補則からも導けます。

結果
\centerline{\item{(7)}} \( H \leq f(x) \leq K \) のとき \( a \leq x \leq b \)，則
\[ H(b - a) \leq \int_{a}^{b}f(x)\, dx \leq K(b - a)。 \]

注記
これは、(6) を \( f(x) - H \) および \( K - f(x) \) に適用することで即座に導けます。

結果
\centerline{\item{(8)}} \(\ds\int_{a}^{b}f(x)\, dx = (b-a)f(\xi)\)，ここで \(\xi\) は \( a \) と \( b \) の間にある。

\par\medskip\noindent\textbf{【注釈】}\par
これは（7）から導けます。なぜなら、$H$ を $\DPmod{(a, b)}{[a, b]}$ 内で $f(x)$ の最小値とし、$K$ を最大値とすると、積分は $\eta(b - a)$ に等しくなります。ここで $\eta$ は $H$ と $K$ の間の値です。また、$f(x)$ が連続であるからには、ある $\xi$ の値で $f(\xi) = \eta$ となるものが存在します（100節）。

積分関数 $F(x)$ を用いると、（8）の結果は以下の形で表すことができます：
\[
F(b) - F(a) = (b - a)F'(\xi),
\]
これにより、（8）は125節の平均値定理の別の表現に過ぎないことがわかります。私たちは（8）を「積分の最初の平均値定理」と呼ぶことにします。
\par\medskip
\PgSep{293}

\par\medskip\noindent\textbf{【結果】}\par
\Item{(9)} \Topic{積分に対する一般化された平均値定理}
$\phi(x)$ が正であるとき、（7）で定義した $H$ と $K$ を用いて以下の不等式が成り立ちます：
\[
H\int_{a}^{b} \phi(x)\, dx
  \leq \int_{a}^{b} f(x)\phi(x)\, dx
  \leq K\int_{a}^{b} \phi(x)\, dx;
\]
さらに、
\[
\int_{a}^{b} f(x)\phi(x)\, dx = f(\xi) \int_{a}^{b} \phi(x)\, dx,
\]
ここで $\xi$ は（8）で定義された値です。
\par\medskip

\par\medskip\noindent\textbf{【注釈】}\par
これは、（6）の定理を以下の積分に適用することで即座に導けます：
\[
\int_{a}^{b} \{f(x) - H\}\phi(x)\, dx,\quad
\int_{a}^{b} \{K - f(x)\}\phi(x)\, dx.
\]
$\phi(x)$ が常に負である場合にも同様の結果が成り立つことを読者は自分で確認すべきです。
\par\medskip

\par\medskip\noindent\textbf{【結果】}\par
\Itemp{(10)} \Topic{積分計算の基本定理}
関数
\[
F(x) = \int_{a}^{x} f(t)\, dt
\]
は $f(x)$ に等しい導函数を持ちます。
\par\medskip

これは既に145節で証明されていますが、ここでは形式的な定理として再掲することに便益があります。また、157節で指摘したように、これより派生する結果として $F(x)$ は $x$ の連続関数であることがわかります。

例 LXV．

1. 直接の定義による定積分と上記の方程式（1）～（5）により示せ：\[
\int_{-a}^{a} \phi(x^{2})\, dx = 2\int_{0}^{a} \phi(x^{2})\, dx,\quad
\int_{-a}^{a} x\phi(x^{2})\, dx = 0;
\]
\[
\int_{0}^{\frac{1}{2}\pi} \phi(\cos x)\, dx
= \int_{0}^{\frac{1}{2} \pi} \phi(\sin x)\, dx
= \tfrac{1}{2} \int_{0}^{\pi} \phi(\sin x)\, dx;
\]
\[
\int_{0}^{m\pi} \phi(\cos^{2} x)\, dx = m\int_{0}^{\pi} \phi(\cos^{2} x)\, dx,\quad
m は整数．[これらの等式的真偽は、積分記号内の関数のグラフを描くことで幾何学的に直感的に理解できる]

2. $\ds\int_{0}^{\pi} \frac{\sin nx}{\sin x}\, dx$ は $n$ が奇数であるとき $\pi$、偶数であるとき $0$ に等しいことを証明せよ．[最後の項が $1$ または $2\cos x$ となる公式 $(\sin nx)/(\sin x) = 2\cos \{(n - 1)x\} + 2\cos \{(n - 3)x\} + \dots$ を用いる]

3. $\ds\int_{0}^{\pi} \sin nx \cot x\, dx$ は $n$ が奇数であるとき $0$、偶数であるとき $\pi$ に等しいことを証明せよ．

4. $\phi(x) = a_{0} + a_{1}\cos x + b_{1}\sin x + a_{2}\cos 2x + \dots + a_{n}\cos nx + b_{n}\sin nx$ であり、$k$ は $n$ 以下の正の整数であるとき
\[
\int_{0}^{2\pi} \phi(x)\, dx = 2\pi a_{0},\quad
\int_{0}^{2\pi} \cos kx \phi(x)\, dx = \pi a_{k},\quad
\int_{0}^{2\pi} \sin kx \phi(x)\, dx = \pi b_{k}.
\]
$k > n$ のとき、最後の二つの積分の値は $0$ である．[ lxiii.9 を用いる]

5. $\phi(x) = a_{0} + a_{1} \cos x + a_{2}\cos 2x + \dots + a_{n}\cos nx$ であり、$k$ は $n$ 以下の正の整数であるとき
\[
\int_{0}^{\pi} \phi(x)\, dx = \pi a_{0},\quad
\int_{0}^{\pi} \cos kx \phi(x)\, dx = \tfrac{1}{2}\pi a_{k}.
\]
$k > n$ のとき、最後の積分の値は $0$ である．[ lxiii.10 を用いる]

6. \(a\) と \(b\) が正のとき、次の等式を証明せよ：
\[
\int_{0}^{2\pi} \frac{dx}{a^{2}\cos^{2} x + b^{2}\sin^{2} x} = \frac{2\pi}{ab}.
\]

7. \(a \leq x \leq b\) において \(f(x) \leq \phi(x)\) ならば、
\[
\int_{a}^{b} f\, dx \leq \int_{a}^{b}\phi\, dx.
\]

8. 次の不等式を証明せよ：
\[
0 < \int_{0}^{\frac{1}{2}\pi} \sin^{n+1}x\, dx < \int_{0}^{\frac{1}{2}\pi} \sin^{n}x\, dx,
\]
\[
0 < \int_{0}^{\frac{1}{4}\pi} \tan^{n+1}x\, dx < \int_{0}^{\frac{1}{4}\pi} \tan^{n}x\, dx.
\]

9. \(n > 1\) のとき、次の不等式を証明せよ：
\[
0.5 < \int_{0}^{\frac{1}{2}} \frac{dx}{\sqrt{1 - x^{2n}}} < 0.524.
\]

10. 次の不等式を証明せよ：
\[
0.5 < \int_{0}^{1} \frac{dx}{\sqrt{4 - x^{2} + x^{3}}} < \frac{\pi}{6}.
\]

11. \(0 < x < 1\) のとき、次の不等式を証明せよ：
\[
\frac{3x + 8}{16} < \frac{1}{\sqrt{4 - 3x + x^{3}}} < \frac{1}{\sqrt{4 - 3x}},
\]
これより
\[
0.573 < \int_{0}^{1} \frac{dx}{\sqrt{4 - 3x + x^{3}}} < \frac{2}{3}.
\]

12. 次の不等式を証明せよ：
\[
0.573 < \int_{1}^{2} \frac{dx}{\sqrt{4 - 3x + x^{3}}} < 0.595.
\]
\(x = 1 + u\) と置き換えて、
\[
2 + 3u^{2} + u^{3}
\]
を
\[
2 + 4u^{2}, \quad 2 + 3u^{2}
\]
に近似する。

13. 正の鋭角 \(\alpha\) と \(\phi\) に対して、次の不等式が成り立つ：
\[
\phi < \int_{0}^{\phi} \frac{dx}{\sqrt{1 - \sin^{2}\alpha \sin^{2} x}} < \frac{\phi}{\sqrt{1 - \sin^{2}\alpha \sin^{2}\phi}}.
\]
\(\alpha = \phi = \frac{1}{6}\pi\) のとき、積分値は \(0.523\) と \(0.541\) の間にある。

14. \( f(x) \) の定積分の絶対値が \( f(x) \) の絶対値の定積分で上から制限されることを証明せよ：
\[
\left|\int_{a}^{b} f(x)\, dx\right| \leq \int_{a}^{b}|f(x)|\, dx.
\]

もし \(\sigma\) が第156節末尾で考慮される和であり、\(\sigma'\) が関数 \( |f(x)| \) から形成される対応する和であるならば、\( |\sigma| \leq \sigma' \) である。

15. \( |f(x)| \leq M \) のとき、
\[
\left|\int_{a}^{b} f(x)\phi(x)\, dx\right| \leq M\int_{a}^{b}|\phi(x)|\, dx.
\]

161. 部分積分と置換積分。第138節から以下の公式が導かれる：
\[
\int_{a}^{b} f(x)\phi'(x)\, dx
  = f(b)\phi(b) - f(a)\phi(a) - \int_{a}^{b} f'(x)\phi(x)\, dx.
\]
この公式は、定積分の部分積分の公式と呼ばれる。

また、第133節から知っているように、\( F(t) \) が \( f(t) \) の積分関数であるとき、
\[
\int f\{\phi(x)\}\phi'(x)\, dx = F\{\phi(x)\}.
\]
したがって、\(\phi(a) = c\)、\(\phi(b) = d\) のとき、
\[
\int_{c}^{d} f(t)\, dt
  = F(d) - F(c)
  = F\{\phi(b)\} - F\{\phi(a)\}
  = \int_{a}^{b} f\{\phi(x)\}\phi'(x)\, dx;
\]
これは、置換によって定積分を変形するための公式である。

積分の部分積分法や変換公式はしばしば、被積分関数の不定積分を求めることなく定積分の値を求めるのに役立ちます。また、不定積分を求めることができない場合でも、そのような場合もあります。以下にいくつかの例を見ていきましょう。定積分の値が不定積分の知識なしで求められることがあるのは当然であり、関数 \(F(x)\) の一般形を決定することができないという事実は、その特定の二つの値の差 \(F(b) - F(a)\) を決定できる可能性を阻害しないからです。しかし、通常これは現在の方法では達成できないことが多いです。

例 LXVI.

1. 以下を証明せよ：
\[
\int_{a}^{b} x f''(x)\, dx = \{bf'(b) - f(b)\} - \{af'(a) - f(a)\}.
\]

2. より一般に、以下の等式が成り立つことを示せ：
\[
\int_{a}^{b} x^{m} f^{(m+1)}(x)\, dx = F(b) - F(a),
\]
ここで
\begin{multline*}
F(x) = x^{m} f^{(m)}(x)
     - mx^{m-1} f^{(m-1)}(x) \\
     + m(m - 1)x^{m-2} f^{(m-2)}(x) - \dots
     + (-1)^{m} m!\, f(x).
\end{multline*}

3. 以下の等式を証明せよ：
\[
\int_{0}^{1} \arcsin x\, dx = \frac{\pi}{2} - 1,\quad
\int_{0}^{1}x\arctan x\, dx = \frac{\pi}{4} - \frac{1}{2}.
\]

4. \(a\) と \(b\) が正のとき、以下の等式を証明せよ：
\[
\int_{0}^{\frac{1}{2}\pi}
  \frac{x\cos x\sin x\, dx}{(a^{2}\cos^{2}x + b^{2}\sin^{2}x)^{2}}
  = \frac{\pi}{4ab^{2}(a + b)}.
\]

部分積分法を用いて、例 LXIII.8 を使用せよ。

5. \(f_{1}(x) = \int_{0}^{x}f(t)\, dt,\quad f_{2}(x) = \int_{0}^{x}f_{1}(t)\, dt,\ \dots,\quad f_{k}(x) = \int_{0}^{x} f_{k-1}(t)\, dt\) のとき、以下の等式が成り立つことを示せ：
\[
f_{k}(x) = \frac{1}{(k - 1)!} \int_{0}^{x} f(t)(x - t)^{k-1}\, dt.
\]

部分積分法を繰り返し使用せよ。

6. \( u_{m, n} = \int_{0}^{1} x^{m} (1 - x)^{n}\, dx \) において、\( m \) と \( n \) が正の整数であるとき、積分部分積分法を用いて \( (m + n + 1) u_{m, n} = nu_{m, n-1} \) を証明し、さらに
\[
u_{m, n} = \frac{m! \, n!}{(m + n + 1)!}
\]
を得よ。

7. \( u_{n} = \int_{0}^{\frac{1}{4}\pi} \tan^{n} x\, dx \) のとき、\( u_{n} + u_{n-2} = \frac{1}{n - 1} \) を証明し、これよりすべての正の整数 \( n \) における積分を評価せよ。

[置換 \(\tan^{n} x = \tan^{n-2} x (\sec^{2} x - 1)\) して部分積分を行う。]

8. 前例から \( u_{n} \) は \(\frac{1}{2(n - 1)}\) と \(\frac{1}{2(n + 1)}\) の間にあることを導け。

9. \( u_{n} = \int_{0}^{\frac{1}{2}\pi} \sin^{n} x\, dx \) のとき、\( u_{n} = \frac{n - 1}{n} u_{n-2} \) を証明せよ。[置換 \(\sin^{n-1} x \sin x\) として部分積分を行う。]

10. 前問から \( n \) が奇数または偶数のとき、
\[
u_{n} = \frac{2 \cdot 4 \cdot 6 \dots (n - 1)}{3 \cdot 5 \cdot 7 \dots n}, \quad \frac{1}{2}\pi \frac{1 \cdot 3 \cdot 5 \dots (n - 1)}{2 \cdot 4 \cdot 6 \dots n}
\]
を得よ。

11. 第二の平均値定理。関数 \( f(x) \) が \( x = a \) から \( x = b \) のすべての \( x \) の値に対して符号が一定な微分係数を持つとき、\( a \) と \( b \) の間に数 \( \xi \) が存在し、
\[
\int_{a}^{b} f(x)\phi(x)\, dx = f(a) \int_{a}^{\xi} \phi(x)\, dx + f(b) \int_{\xi}^{b} \phi(x)\, dx
\]
を満たす。

[置換 \(\int_{a}^{x}\phi(t)\, dt = \Phi(x)\) とすると、
\begin{align*}
\int_{a}^{b} f(x)\phi(x)\, dx &= \int_{a}^{b} f(x)\Phi'(x)\, dx \\
&= f(b)\Phi(b) - \int_{a}^{b} f'(x)\Phi(x)\, dx \\
&= f(b)\Phi(b) - \Phi(\xi) \int_{a}^{b}f'(x)\, dx,
\end{align*}
一般化された平均値定理（§160）により、
\[
\int_{a}^{b} f(x)\phi(x)\, dx = f(b)\Phi(b) + (f(a) - f(b))\Phi(\xi)
\]
となり、これが要求される結果と等価である。]

\item{12.} ボネッットの第2中值定理の形式。$f'(x)$が恒等的に符号を持つ場合、かつ$f(b)$と$f(a) - f(b)$が同じ符号を持つとき、
\[
\int_{a}^{b} f(x)\phi(x)\, dx = f(a) \int_{a}^{X} \phi(x)\, dx,
\]
ここで$X$は$a$と$b$の間にある。[定数$\mu$が$\Phi(\xi)$と$\Phi(b)$の間にあり、かつ$\Phi(x)$の$x$の値に対応する値としてこの式の値であるとき、
$f(b)\Phi(b) + \{f(a) - f(b)\}\Phi(\xi) = \mu f(a)$が成り立つ。重要な場合としては、$0 \leq f(b) \leq f(x) \leq f(a)$の場合がある。]

同様に証明すると、$f(a)$と$f(b) - f(a)$が同じ符号を持つとき、
\[
\int_{a}^{b} f(x)\phi(x)\, dx = f(b) \int_{X}^{b} \phi(x)\, dx,
\]
ここで$X$は$a$と$b$の間にある。[関数$\Psi(\xi) = \int_{\xi}^{b} \phi(x)\, dx$を使用する。この積分は次の形で表すことができる：
\[
f(a)\Psi(a) + \{f(b) - f(a)\}\Psi(\xi).
\]
重要な場合としては、$0 \leq f(a) \leq f(x) \leq f(b)$の場合がある。]

\item{13.} 次の不等式を証明せよ：
\[
\left|\int_{X}^{X'} \frac{\sin x}{x}\, dx\right| < \frac{2}{X}
\]
$X' > X > 0$のとき。[第1項の公式を使用し、任意の区間に対する$\sin x$の積分が数値的に2未満であることに注意せよ。]

\item{14.} \Ex{lxv}.1の結果を置換則を利用して証明せよ。
[(i) の場合、積分範囲を$\DPmod{(-a, 0)}{[-a, 0]}$と$\DPmod{(0, a)}{[0, a]}$に分割し、最初の部分では$x = -y$とする。 (ii) の場合、置換則$x = \frac{1}{2}\pi - y$を使用して最初の方程式を得る。二つ目の方程式を得るためには、積分範囲$\DPmod{(0, \pi)}{[0, \pi]}$を二等分し、$x = \frac{1}{2}\pi + y$の置換を利用する。 (iii) の場合、積分範囲を$m$等分割し、$x = \pi + y$, $x = 2\pi + y$,~\dotsの置換を利用する。]

\item{15.} 次の等式を証明せよ：
\[
\int_{a}^{b} F(x)\, dx = \int_{a}^{b} F(a + b - x)\, dx.
\]

\item{16.} 次を証明せよ：
\[
\int_{0}^{\frac{1}{2}\pi} \cos^{m} x \sin^{m} x \, dx = 2^{-m} \int_{0}^{\frac{1}{2}\pi} \cos^{m} x \, dx.
\]

\item{17.} 次を証明せよ：
\[
\int_{0}^{\pi} x \phi(\sin x) \, dx = \frac{1}{2}\pi \int_{0}^{\pi} \phi(\sin x) \, dx.
\]

\[x = \pi - y\]と置く。

\item{18.} 次を証明せよ：
\[
\int_{0}^{\pi} \frac{x \sin x}{1 + \cos^{2} x} \, dx = \frac{1}{4}\pi^{2}.
\]

\item{19.} $x = a \cos^{2} \theta + b \sin^{2} \theta$ の変換により次を証明せよ：
\[
\int_{a}^{b} \sqrt{(x - a)(b - x)} \, dx = \frac{1}{8}\pi (b - a)^{2}.
\]

\item{20.} $(a + b \cos x) (a - b \cos y) = a^{2} - b^{2}$ の置換により次を証明せよ：
\[
\int_{0}^{\pi} (a + b \cos x)^{-n} \, dx = (a^{2} - b^{2})^{-(n - \frac{1}{2})} \int_{0}^{\pi} (a - b \cos y)^{n-1} \, dy,
\]
$n$ が正の整数であり $a > |b|$ のとき、また $n = 1$, $2$, $3$ の場合を評価せよ。

\item{21.} $m$ と $n$ が正の整数であるとき次を証明せよ：
\[
\int_{a}^{b} (x - a)^{m} (b - x)^{n} \, dx = (b - a)^{m+n+1} \frac{m! \, n!}{(m + n + 1)!}.
\]

\[x = a + (b - a) y\]と置き、例題6を用いる。

\paragraph{162. テイラーの定理の積分法による証明} 我々は \SecNo[§]{147} で言及した別の証明形式を与える。

$f(x)$ を第 $n$ 次導関数が連続な関数とするとき、
\[
F_{n}(x) = f(b) - f(x) - (b - x)f'(x) - \dots - \frac{(b - x)^{n-1}}{(n - 1)!} f^{(n-1)}(x).
\]

すると
\[
F_{n}'(x) = -\frac{(b - x)^{n-1}}{(n - 1)!} f^{(n)}(x),
\]
よって
\[
F_{n}(a) = F_{n}(b) - \int_{a}^{b} F_{n}'(x) \, dx = \frac{1}{(n - 1)!} \int_{a}^{b} (b - x)^{n-1} f^{(n)}(x) \, dx.
\]
$a + h$ を $b$ と置き、積分を $x = a + th$ で変換すると
\[
f(a + h) = f(a) + hf'(a) + \dots + \frac{h^{n-1}}{(n - 1)!} f^{(n-1)}(a) + R_{n},
\]
ここで
\[
R_{n} = \frac{h^{n}}{(n - 1)!} \int_{0}^{1} (1 - t)^{n-1} f^{(n)}(a + th) \, dt.
\]

現在、任意の正整数$p$が$n$より小さいとき、定理(9)（160節）により、

\begin{align*}
\int_{0}^{1} (1 - t)^{n-1} f^{(n)}(a + th)\, dt
  &= \int_{0}^{1}(1 - t)^{n-p} (1 - t)^{p-1} f^{(n)}(a + th)\, dt \\
  &= (1 - \theta)^{n-p} f^{(n)}(a + \theta h) \int_{0}^{1} (1 - t)^{p-1}\, dt,
\end{align*}

ここで $0 < \theta < 1$。従って、

\[
R_{n} = \frac{(1 - \theta)^{n-p} f^{(n)}(a + \theta h)h^{n}}{p(n - 1)!}.
\Tag{(3)}
\]

$p = n$とすると、ラグランジュの形式が得られます（148節）。一方で $p = 1$とすると、カuchyの形式が得られ、

\[
R_{n} = \frac{(1 - \theta)^{n-1} f^{(n)}(a + \theta h) h^{n}}{(n - 1)!}.\footnotemark
\Tag{(4)}
\]
\footnotetext{147節で使用された方法も、これらの余項の異なる形式を導くために修正することができます。}

カuchyの形式を二項級数に適用すると、$f(x) = (1 + x)^{m}$（$m$が正整数でないとき）、余項は

\[
R_{n} = \frac{m(m - 1)\dots (m - n + 1)}{1·2\dots (n - 1)}\,
        \frac{(1 - \theta )^{n-1} x^{n}}{(1 + \theta x)^{n-m}}
\]

となります。

$(1 - \theta)/(1 + \theta x)$は$-1 < x < 1$の範囲で常に1未満であり、$x$が正でも負でも同じです。また、$(1 + \theta x)^{m-1}$はすべての$n$に対して一定の定数$K$より小さい値を取るため、実際には$m > 1$の場合$(1 + |x|)^{m-1}$よりも小さく、$m < 1$の場合$(1 - |x|)^{m-1}$よりも小さいです。従って、

\[
|R_{n}| < K |m| \left|\binom{m - 1}{n - 1}\right| |x^{n}| = \rho_{n}
\]

となります。ただし$\rho_{n} \to 0$（$n \to \infty$）であり、これにより$R_{n} \to 0$が成り立ちます。二項定理の真偽は、すべての有理数$m$と$-1 < x < 1$の範囲での$x$に対して証明されます。

ラグランジュの形式では、\Ex{xxvii}.2において負の$x$値に関連する難題が生じていました。

164. 実変数の複素関数の積分
これまで、定積分の対象が常に実数であると仮定してきました。実変数 $x$ の複素関数 $f(x) = \phi(x) + i\psi(x)$ に対する定積分を、以下の等式により定義します：
\[
\int_{a}^{b} f(x)\, dx
  = \int_{a}^{b} \{\phi(x) + i\psi(x)\}\, dx
  = \int_{a}^{b} \phi(x)\, dx + i \int_{a}^{b} \psi(x)\, dx;
\]
このような積分の性質は、既に考慮した実積分の性質から導き出されます。

これらの性質のうち、後で用いるものがあります。これは不等式によって表現され、以下のようになります：
\[
\left|\int_{a}^{b} f(x)\, dx\right| \leq \int_{a}^{b} |f(x)|\, dx.\footnotemark
\Tag{(1)}
\]
\footnotetext{同様の実積分に対する不等式は、§lxv. 14で証明されています。}

この不等式は、§156と§157での定義から容易に導けます。$\delta_{\nu}$ の意味が §156 と同じであり、$\phi_{\nu}$ および $\psi_{\nu}$ は $\delta_{\nu}$ 上の点における $\phi$ および $\psi$ の値で、$f_{\nu} = \phi_{\nu} + i\psi_{\nu}$ とします。このとき、以下のようになります：
\begin{align*}
\int_{a}^{b} f\, dx
   &= \int_{a}^{b} \phi\, dx + i \int_{a}^{b} \psi\, dx \\
  &= \lim \sum \phi_{\nu}\, \delta_{\nu} + i \lim \sum \psi_{\nu}\, \delta_{\nu} \\
  &= \lim \sum (\phi_{\nu} +  i\psi_{\nu})\, \delta_{\nu}
   = \lim \sum f_{\nu}\, \delta_{\nu},
\end{align*}
したがって
\[
\int_{a}^{b} f\, dx
  = |\lim \sum f_{\nu}\, \delta_{\nu}|
  = \lim | \sum f_{\nu}\, \delta_{\nu}|;
\]
一方で
\[
\int_{a}^{b} |f|\, dx = \lim \sum |f_{\nu}|\, \delta_{\nu}.
\]

不等式 $|\sum f_{\nu}\, \delta_{\nu}| \leq \sum |f_{\nu}|\, \delta_{\nu}$ から、結果はすぐに導けます。

また、§162の公式 (1) および (2) は、$f$ が複素関数 $\phi + i\psi$ の場合でも真です。

第7章の補足例
\Section{第7章の補足例}

1. 次のTaylor級数の項を確認せよ：
\begin{alignat*}{2}
\& (1) \& \tan x \&= x + \frac{1}{3}x^{3} + \frac{2}{15}x^{5} + \cdots, \\
\& (2) \& \sec x \&= 1 + \frac{1}{2}x^{2} + \frac{5}{24}x^{4} + \cdots, \\
\& (3) \& x\cosec x \&= 1 + \frac{1}{6}x^{2} + \frac{7}{360}x^{4} + \cdots, \\
\& (4) \& x\cot x \&= 1 - \frac{1}{3}x^{2} - \frac{1}{45}x^{4} - \cdots.
\end{alignat*}

2. $f(x)$ およびその最初の $n + 2$ 階微分が連続であり、$f^{(n+1)}(0) \neq 0$ かつ Taylor級数の $n$ 項後のラグランジュ形式の余項で現れる $\theta_{n}$ の値が $\theta$ の値であるとき、
\[
\theta_{n} = \frac{1}{n + 1}
  + \frac{n}{2(n + 1)^{2}(n + 2)}
    \left\{\frac{f^{(n+2)}(0)}{f^{(n+1)}(0)} + \epsilon_{x}\right\}x,
\]
ここで $\epsilon_{x} \to 0$ であるとき、$x \to 0$。[lv.12の方法を用いて証明せよ]。

3. $f(x) = \frac{1}{1 + x}$ の場合、最後の結果を確認せよ。[ここで $(1 + \theta_{n}x)^{n+1} = 1 + x$]

4. $f(x)$ が最初の三階微分を持つとき、
\[
f(b) = f(a) + \frac{1}{2}(b - a)\{f'(a) + f'(b)\}
  - \frac{1}{12}(b - a)^{3} f'''(\alpha),
\]
ここで $a < \alpha < b$。[147節の議論に類似した方法を用いて証明せよ]。

5. 同じ条件下で、
\[
f(b) = f(a) + (b - a)f'\left\{\frac{1}{2}(a + b)\right\}
  + \frac{1}{24}(b - a)^{3}f'''(\alpha).
\]

6. $f(x)$ が最初の五階微分を持つとき、
\[
f(b) = f(a) + \frac{1}{6}(b - a)\left[f'(a) + f'(b) + 4f'\left\{\frac{1}{2}(a + b)\right\}\right]
  - \frac{1}{2880}(b - a)^{5} f^{(5)}(\alpha).
\]

7. 同じ条件下において、
\[
f(b) = f(a) + \frac{1}{2}(b - a)\{f'(a) + f'(b)\}
  - \frac{1}{12}(b - a)^{2}\{f''(b) - f''(a)\}
  + \frac{1}{720}(b - a)^{5} f^{(5)}(\alpha).
\]

8. 次の公式を証明せよ：
\[
\begin{vmatrix}
  f(a) &  f(b)\\
  g(a) &  g(b)
\end{vmatrix}
= (b - a)
\begin{vmatrix}
  f(a)  &  f'(\beta)\\
  g(a)  &  g'(\beta)
\end{vmatrix},
\]
ここで $\beta$ は $a$ と $b$ の間にある。また、
\[
\begin{vmatrix}
  f(a) & f(b) & f(c)\\
  g(a) & g(b) & g(c)\\
  h(a) & h(b) & h(c)
\end{vmatrix}
= \frac{1}{2} (b - c)(c - a)(a - b)
\begin{vmatrix}
  f(a) & f'(\beta) & f''(\gamma)\\
  g(a) & g'(\beta) & g''(\gamma)\\
  h(a) & h'(\beta) & h''(\gamma)
\end{vmatrix},
\]
ここで $\beta$ と $\gamma$ は $a$, $b$, $c$ の最小値と最大値の間にある。[公式(ii) を証明するには、関数
\[
\phi(x) =
\begin{vmatrix}
  f(a) & f(b) & f(x)\\
  g(a) & g(b) & g(x)\\
  h(a) & h(b) & h(x)
\end{vmatrix}
- \frac{(x - a)(x - b)}{(c - a)(c - b)}
\begin{vmatrix}
  f(a) & f(b) & f(c)\\
  g(a) & g(b) & g(c)\\
  h(a) & h(b) & h(c)
\end{vmatrix},
\]
を考える。この関数は $x = a$, $x = b$, $x = c$ のとき零となる。その第1導函数は、定理 B（§121）により、$a$, $b$, $c$ の最小値と最大値の間で異なる 2 個の値において零となる。したがって、その第2導函数も同様の条件を満たす $\gamma$ の値において零となる。これにより、次の公式を得る：
\[
\begin{vmatrix}
  f(a) & f(b) & f(c)\\
  g(a) & g(b) & g(c)\\
  h(a) & h(b) & h(c)
\end{vmatrix}
= \frac{1}{2}(c - a)(c - b)
\begin{vmatrix}
  f(a) & f(b) & f''(\gamma)\\
  g(a) & g(b) & g''(\gamma)\\
  h(a) & h(b) & h''(\gamma)
\end{vmatrix}.
\]
読者はこの証明を難なく完成させるであろう。

\item9. \(F(x)\)が\(x = 0\)で第\(n-1\)階の導数がすべて零であり、かつ第\(n\)階の導数について\(A \leq F^{(n)}(x) \leq B\)が\(0 \leq x \leq h\)において成り立つとき、\(0 \leq x \leq h\)において\(A(x^n/n!) \leq F(x) \leq B(x^n/n!)\)である。

この結果を適用し
\[
f(x) - f(0) - xf'(0) - \dots - \frac{x^{n-1}}{(n-1)!} f^{(n-1)}(0)
\]
について考え、Taylorの定理を得る。

\item10. \(\Delta_h \phi(x) = \phi(x) - \phi(x + h)\), \(\Delta_h^2 \phi(x) = \Delta_h \{\Delta_h \phi(x)\}\)，およびその他の差分を定義し、\(\phi(x)\)が第\(n\)階の導数を持つとき
\[
\Delta_h^n \phi(x)
  = \sum_{r=0}^n (-1)^r \binom{n}{r} \phi(x + rh)
  = (-h)^n \phi^{(n)}(\xi),
\]
ここで\(\xi\)は\(x\)と\(x + nh\)の間にある。したがって、\(\phi^{(n)}(x)\)が連続であるとき
\[
\frac{\Delta_h^n \phi(x)}{h^n} \to (-1)^n \phi^{(n)}(x)
\]
が\(h \to 0\)のときに成り立つ。この結果は既に\(n = 2\)のとき、\Ex{lv}.13において述べられている。

\item11. \Ex{10}から推論し、\(x^n - m \Delta_h^n x^m \to m(m-1) \dots (m-n+1)h^n\)が\(x \to \infty\)のときに成り立つことを示す。ここで\(m\)は任意の有理数であり、\(n\)は任意の正整数である。特に
\[
x\sqrt{x} \{\sqrt{x} - 2\sqrt{x+1} + \sqrt{x+2}\} \to -\frac{1}{4}
\]
を証明する。

\item12. \(y = \phi(x)\)が\(x\)の関数であり、少なくとも第4階までの導数が連続であるとし、\(\phi(0) = 0\), \(\phi'(0) = 1\)とする。すると
\[
y = \phi(x) = x + a_2 x^2 + a_3 x^3 + (a_4 + \epsilon_x)x^4,
\]
ここで\(\epsilon_x \to 0\)が\(x \to 0\)のときに成り立つ。この関数について
\[
x = \psi(y)
  = y - a_2 y^2 + (2a_2^2 - a_3)y^3
      - (5a_2^3 - 5a_2 a_3 + a_4 + \epsilon_y)y^4,
\]
ここで\(\epsilon_y \to 0\)が\(y \to 0\)のときに成り立つ。を満たす\(x\)の値についての公式を確立し、
\[
\frac{\phi(x)\psi(x) - x^2}{x^4} \to a_2^2
\]
が\(x \to 0\)のときに成り立つことを証明する。

\item{13.} 曲線 \( x = f(t), y = F(t) \) の点 \((x, y)\) における曲率中心の座標 \((\xi, \eta)\) は以下の式で与えられる：
\[
- \frac{\xi - x}{y'} = \frac{\eta - y}{x'} = \frac{x'^2 + y'^2}{x'y'' - x''y'};
\]
また、曲線の曲率半径は
\[
\left( x'^2 + y'^2 \right)^{3/2} / (x'y'' - x''y')
\]
であり、破線は \( t \) に関する微分を表す。

\item{14.} 曲線 \( 27ay^2 = 4x^3 \) の点 \((x, y)\) における曲率中心の座標 \((\xi, \eta)\) は以下の式で与えられる：
\[
3a(\xi + x) + 2x^2 = 0, \quad
\eta = 4y + \frac{9ay}{x}.
\]
\MathTrip{1899.}

\item{15.} 点 \((x, y)\) における接点において \( (1 + y_1^2)y_3 = 3y_1y_2^2 \) を満たすとき、曲線に対する接円はその点で３次接触を有することを証明せよ。また、この性質を持つ唯一の曲線であることを証明し、さらに、双曲線におけるこの性質を持つ唯一の点は軸の端点であることを見よ。[参照：\okrickRef{第VI章}{VI}, \MiscEx{VI}\ 10(iv)]。

\item{16.} 曲線 \( y = ax^2 + bx^3 + cx^4 + \cdots + kx^n \) の原点における最接近楕円は \( a^3y = a^4x^2 + a^2bxy + (ac - b^2)y^2 \) である。したがって、曲線 \( y = f(x) \) 上の点 \((\xi, \eta)\) における最接近楕円は以下の式で与えられる：
\[
18\eta_2^3T
= 9\eta_2^4(x - \xi)^2
+ 6\eta_2^2\eta_3(x - \xi)T
+ (3\eta_2\eta_4 - 4\eta_3^2)T^2,
\]
ここで \( T = (y - \eta) - \eta_1(x - \xi) \)。
\MathTrip{1907.}

\Item{17.} \Topic{同次関数。\footnote{この例と次の例では、出現在するすべての導函数が連続であることを読者には仮定してほしい。}}
$u = x^{n} f(y/x, z/x, \dots)$ とするとき、$x$, $y$, $z$, \dots が $\lambda : 1$ の比で増加したときに $u$ は元の形を保ち、因子 $\lambda^{n}$ を除いて変化しない。このような状況では、$u$ を $x$, $y$, $z$, \dots の同次関数と呼ぶ。$u$ が同次で次数 $n$ のとき、
\[
x\frac{\partial u}{\partial x} + y\frac{\partial u}{\partial y} + z\frac{\partial u}{\partial z} + \dots = nu.
\]
この結果は、同次関数に関する \Emph{オイラーの定理} と呼ばれる。

\Item{18.} $u$ が同次で次数 $n$ のとき、$\frac{\partial u}{\partial x}$, $\frac{\partial u}{\partial y}$, \dots はそれぞれ次数 $n - 1$ の同次関数である。

\Item{19.} 方程式 $f(x, y) = 0$（例：$x^{n} + y^{n} - x = 0$）を $z$ を 1 に置き換えることで同次化した形として $F(x, y, z) = 0$ とする。曲線 $f(x, y) = 0$ 上の点 $(\xi, \eta)$ の接線の方程式は
\[
xF_{\xi} + yF_{\eta} + zF_{\zeta} = 0,
\]
ここで、$F_{\xi}$, $F_{\eta}$, $F_{\zeta}$ は $x = \xi$, $y = \eta$, $z = \zeta = 1$ のときの $F_{x}$, $F_{y}$, $F_{z}$ の値を表す。

\Item{20.} \Topic{依存関数と独立関数。ヤコビアンや関数行列式。} $u$ と $v$ を $x$ と $y$ の関数として、同一の関係
\[
\phi(u, v) = 0.
\Tag{(1)}
\]

異なるiating \Eq{(1)} について $x$ と $y$ で微分すると、次を得る：
\[
\frac{\dd \phi}{\dd u}\, \frac{\dd u}{\dd x} +
\frac{\dd \phi}{\dd v}\, \frac{\dd v}{\dd x} = 0,\quad
\frac{\dd \phi}{\dd u}\, \frac{\dd u}{\dd y} +
\frac{\dd \phi}{\dd v}\, \frac{\dd v}{\dd y} = 0,
\Tag{(2)}
\]
そして、$\phi$ の微分を消去すると、次を得る：
\[
J =
\begin{vmatrix}
  u_{x} & u_{y}\\
  v_{x} & v_{y}
\end{vmatrix}
= u_{x}v_{y} - u_{y}v_{x} = 0,
\Tag{(3)}
\]
ここで $u_{x}$, $u_{y}$, $v_{x}$, $v_{y}$ は、$u$ と $v$ について $x$ と $y$ で微分した値である。
この条件は、次のような関係が存在するための必要条件である。この条件が十分条件でもあることは証明できる；これについては Goursat の \textit{Cours d' Analyse} 第1巻 pp. 125 以降を参照せよ。

二つの関数 $u$ と $v$ は、次のような関係 $\phi(u, v) = 0$ によって結びつけられているか否かにより、依存的または独立的であるという。通常、$J$ を $u$ と $v$ について $x$ と $y$ に対する Jacobi 階行列または関数行列と呼び、次のように表す：
\[
J = \frac{\dd(u, v)}{\dd(x, y)}.
\]

任意の数目的な変数を持つ関数についても類似の結果が成り立つ。したがって、三つの関数 $u$, $v$, $w$ は、三つの変数 $x$, $y$, $z$ について $\phi(u, v, w) = 0$ のような関係によって結びつけられているか否かにより、次のように表される：
\[
J =
\begin{vmatrix}
  u_{x} & u_{y} & u_{z}\\
  v_{x} & v_{y} & v_{z}\\
  w_{x} & w_{y} & w_{z}
\end{vmatrix}
= \frac{\dd(u, v, w)}{\dd(x, y, z)}
\]
が全ての $x$, $y$, $z$ の値に対して消えているか否かにより、独立的または依存的である。

21. $ax^{2} + 2hxy + by_{2}$ と $Ax^{2} + 2Hxy + By^{2}$ が独立であるためには、$a/A = h/H = b/B$ でなければならないことを示せ。

22. $ax^{2} + by^{2} + cz^{2} + 2fyz + 2gzx + 2hxy$ を $x$, $y$, $z$ の一次関数の積として表すことができるためには、かつのみに
\[
abc + 2fgh - af^{2} - bg^{2} - ch^{2} = 0.
\]

[px + qy + rzとp'x + q'y + r'zが与えられた関数によって関数関係にある条件を書き下す.]

\Item{23.} u および v は ξ および η の関数であり、そのうちの ξ および η は x および y の関数であるとき
\[
\frac{\partial(u, v)}{\partial(x, y)}
  = \frac{\partial(u, v)}{\partial(\xi, \eta)}\, \frac{\partial(\xi, \eta)}{\partial(x, y)}
\]
が成り立つ。この結果を任意の数の変数に拡張する。

\Item{24.} x の関数 f(x) があり、その導函数は 1/x であり、x = 1 のとき f(x) は 0 となる。u = f(x) + f(y), v = xy のとき
\[
u_{x}v_{y} - u_{y}v_{x} = 0
\]
が成り立つから、u および v は関数関係によって結びつけられている。y = 1 を代入すると、この関係式は f(x) + f(y) = f(xy) であることが示される。同様の方法で証明すると、f(x) の導函数が 1/(1 + x^2) であり、f(0) = 0 のとき
\[
f(x) + f(y) = f\left(\frac{x + y}{1 - xy}\right)
\]
が成り立つことが示される。

\Item{25.} \[
f(x) = \int_{0}^{x} \frac{dt}{\sqrt{1 - t^{4}}}
\] のとき
\[
f(x) + f(y) = f\left\{
  \frac{x\sqrt{1 - y^{4}} + y\sqrt{1 - x^{4}}}{1 + x^{2}y^{2}}
\right\}.
\]

\Item{26.} u = f(x) + f(y) + f(z), v = f(y)f(z) + f(z)f(x) + f(x)f(y), w = f(x)f(y)f(z) の関数関係が存在するとき、f は定数である。関数関係の条件は
\[
f'(x)f'(y)f'(z) \{f(y) - f(z)\} \{f(z) - f(x)\} \{f(x) - f(y)\} = 0
\]
となる。

\Item{27.} f(y, z), f(z, x), f(x, y) が関数関係によって結びつけられているとき、f(x, x) は x に依存しない。

\Item{28.} u = 0, v = 0, w = 0 が Ex.19 で homogeneous 化された3つの円の方程式であるとき、
\[
\frac{\partial(u, v, w)}{\partial(x, y, z)} = 0
\]
はこれらを正交的に切り取る円の方程式を表す。

\item29. \(A\), \(B\), \(C\) が \(x\) の関数であり、
\[
\begin{vmatrix}
  A & A' & A''\\
  B & B' & B''\\
  C & C' & C''
\end{vmatrix}
\]
が恒等的に消滅するとき、定数 \(\lambda\), \(\mu\), \(\nu\) を用いて \(\lambda A + \mu B + \nu C\) が恒等的に消滅することを示すことができる。逆も成り立つ。［逆はほとんど明らかである。直接の定理を証明するには、\(\alpha = BC' - B'C\)，…とする。すると \(\alpha' = BC'' - B''C\)，…となり、行列式が消滅することから \(\beta\gamma' - \beta'\gamma = 0\)，…であり、したがって比 \(\alpha : \beta : \gamma\) は定数である。しかし \(\alpha A + \beta B + \gamma C = 0\) である。］

\item30. 三つの変数 \(x\), \(y\), \(z\) が関係によって結びつけられており、(i) \(z\) は \(x\) と \(y\) の関数であり、導関数 \(z_{x}\), \(z_{y}\)，および (ii) \(x\) は \(y\) と \(z\) の関数であり、導関数 \(x_{y}\), \(x_{z}\) を持つ。このとき
\[
x_{y} = - \frac{z_{y}}{z_{x}},\quad x_{z} = \frac{1}{z_{x}}
\]
を証明せよ。
\pagetsep{305}

［以下に示す通り、
\[
dz = z_{x}\, dx + z_{y}\, dy,\quad dx = x_{y}\, dy + x_{z}\, dz.
\]
最初の式において \(dx\) を置き換えると
\[
dz = (z_{x} x_{y} + z_{y})\, dy + z_{x} x_{z}\, dz,
\]
これが成り立つには \(z_{x} x_{y} + z_{y} = 0\)，\(z_{x} x_{z} = 1\) である必要がある。］

\item31. 四つの変数 \(x\), \(y\), \(z\), \(u\) が二つの関係によって結びつけられており、どの二つも他のものとしての関数で表すことができる。次のように示せ：
\[
y_{z}^{u}z_{x}^{u}x_{y}^{u} = - y_{z}^{x}z_{x}^{y}x_{y}^{z} = 1,\quad x_{z}^{u}z_{x}^{y} + y_{z}^{u}z_{y}^{x} = 1,
\]
ここで \(y_{z}^{u}\) は \(y\) を \(z\) と \(u\) の関数として \(z\) に関する導関数を表す。［1897年］

\item 32. \(A\), \(B\), \(C\), \(\lambda\) を求めよ。これらの定数は、次の式の最初の4つの導関数が \(x = 0\) のとき零となるようにする：
\[
\int_{a}^{a+x} f(t)\, dt - x[Af(a) + Bf(a + \lambda x) + Cf(a + x)]
\]
また、\(A\), \(B\), \(C\), \(D\), \(\lambda\), \(\mu\) を求めよ。これらの定数は、次の式の最初の6つの導関数が \(x = 0\) のとき零となるようにする：
\[
\int_{a}^{a+x} f(t)\, dt
  - x[Af(a) + Bf(a + \lambda x) + Cf(a + \mu x) + Df(a + x)]
\]

\item 33. \(a > 0\)、\(ac - b^2 > 0\)、および \(x_1 > x_0\) のとき、次の積分が成立する：
\[
\int_{x_0}^{x_1} \frac{dx}{ax^2 + 2bx + c}
  = \frac{1}{\sqrt{ac - b^2}} \arctan\left\{
      \frac{(x_1 - x_0) \sqrt{ac - b^2}}
           {ax_1 x_0 + b(x_1 + x_0) + c}
  \right\},
\]
逆接弦は \(0\) から \(\pi\) の範囲にある。\footnote
  {Exs.~33--35, 38, および 40 に関連して、Dr~Bromwich の論文（\textit{Messenger of Mathematics} 第 xxxv 巻）を参照。}

\item 34. 次の積分を評価せよ：
\[
\int_{-1}^{1} \frac{\sin\alpha\, dx}{1 - 2x\cos\alpha + x^2}.
\]
どのような値の \(\alpha\) に対して、この積分が \(\alpha\) の不連続な関数となるか？

\MathTrip{1904.}

\Loosen[この積分の値は \(2n\pi < \alpha < (2n + 1)\pi\) のとき \(\frac{1}{2}\pi\)、および \((2n - 1)\pi < \alpha < 2n\pi\) のとき \(-\frac{1}{2}\pi\) であり、\(n\) は任意の整数。また、\(\alpha\) が \(\pi\) の倍数であるとき \(0\) となる。]

\item 35. \(ax^2 + 2bx + c > 0\) が \(x_0 \leq x \leq x_1\) の範囲で成立し、\(f(x) = \sqrt{ax^2 + 2bx + c}\)、および
\[
y = f(x),\quad y_0 = f(x_0),\quad y_1 = f(x_1),\quad X = (x_1 - x_0)/(y_1 + y_0),
\]
のとき、次の積分は \(a\) が正または負の場合でそれぞれ次のように表される：
\[
\int_{x_0}^{x_1} \frac{dx}{y}
  = \frac{1}{\sqrt{a}} \log \frac{1 + X\sqrt{a}}{1 - X\sqrt{a}},\quad
\frac{-2}{\sqrt{-a}} \arctan\{X\sqrt{-a}\}.
\]
後者の場合、逆接弦は \(0\) から \(\frac{1}{2}\pi\) の範囲にある。[この積分を \(t = \frac{x - x_0}{y + y_0}\) の置換により \(\ds 2\int_{0}^{X} \frac{dt}{1 - at^2}\) の形に変換することができる。]

36. \(a\) に対して
\[
\int_{0}^{a} \frac{dx}{x + \sqrt{a^{2} - x^{2}}} = \frac{\pi}{4}.
\]
数学トリップ 1913.

37. \(a > 1\) のとき
\[
\int_{-1}^{1} \frac{\sqrt{1 - x^{2}}}{a - x}\, dx = \pi \{a - \sqrt{a^{2} - 1}\}.
\]
ページ分割 306.

38. \(p > 1\), \(0 < q < 1\) のとき
\[
\int_{0}^{1} \frac{dx}{\sqrt{\{(1 + (p^{2} - 1)x)(1 - (1 - q^{2})x)\}}} = \frac{2\omega}{(p + q)\sin\omega},
\]
ここで \(\omega\) は余弦が \(\frac{1 + pq}{p + q}\) の正の鋭角である角度である。

39. \(a > b > 0\) のとき
\[
\int_{0}^{2\pi} \frac{\sin^{2}\theta\, d\theta}{a - b\cos\theta}
  = \frac{2\pi}{b^{2}} \{a - \sqrt{a^{2} - b^{2}}\}.
\]
数学トリップ 1904.

40. \(a > \sqrt{b^{2} + c^{2}}\) のとき
\[
\int_{0}^{\pi} \frac{d\theta}{a + b\cos\theta + c\sin\theta}
  = \frac{2}{\sqrt{a^{2} - b^{2} - c^{2}}} \arctan \left\{\frac{\sqrt{a^{2} - b^{2} - c^{2}}}{c}\right\},
\]
逆接線は \(0\) から \(\pi\) の間である。

41. \(f(x)\) が連続で常に非負であり，\(\int_{a}^{b} f(x)\, dx = 0\) のとき
\(f(x) = 0\) である。［もし \(f(x)\) が \(\xi\) における正の数 \(k\) に等しいとすると，連続性により \((\xi - \delta, \xi + \delta)\) または \([\xi - \delta, \xi + \delta]\) の区間において \(f(x) > \frac{1}{2}k\) となる．そして積分の値は \(\delta k\) より大きい．］

42. ＜シュワルツの不等式＞を証明せよ．
\[
\left(\int_{a}^{b} \phi\psi\, dx\right)^{2}
  \leq \int_{a}^{b} \phi^{2}\, dx \int_{a}^{b} \psi^{2}\, dx.
\]

［定義を用いて §156 と §157 の定義を利用する．また不等式
\[
\left(\sum \phi_{\nu}\psi_{\nu}\delta_{\nu}\right)^{2}
  \leq \sum \phi_{\nu}^{2}\delta_{\nu} \sum \psi_{\nu}^{2}\delta_{\nu}
\]
（第1章，問題集 I 10 参照）を利用する．］

43. \(P_n(x) = \frac{1}{(\beta - \alpha)^n n!} \left( \frac{d}{dx} \right)^n \{(x - \alpha)(\beta - x)\}^n\) と定義すると、\(P_n(x)\)は次数 \(n\) の多項式であり、任意の次数 \(n-1\) より小さい多項式 \(\theta(x)\) に対して
\[
\int_{\alpha}^{\beta} P_n(x) \theta(x) \, dx = 0
\]
を満たす。

44. \(m \neq n\) のとき、
\[
\int_{\alpha}^{\beta} P_m(x) P_n(x) \, dx = 0
\]
であり、\(m = n\) のとき、積分の値は \(\frac{\beta - \alpha}{2n + 1}\) である。

45. \(Q_n(x)\)が次数 \(n\) の多項式で、任意の次数 \(n-1\) より小さい多項式 \(\theta(x)\) に対して
\[
\int_{\alpha}^{\beta} Q_n(x) \theta(x) \, dx = 0
\]
を満たすとき、\(Q_n(x)\)は \(P_n(x)\) の定数倍である。

\(Q_n(x) - \kappa P_n(x)\) を次数 \(n-1\) の多項式と選べば、
\[
\int_{\alpha}^{\beta} Q_n(x) (Q_n(x) - \kappa P_n(x)) \, dx = 0,\quad
\int_{\alpha}^{\beta} P_n(x) (Q_n(x) - \kappa P_n(x)) \, dx = 0,
\]
であり、これより
\[
\int_{\alpha}^{\beta} (Q_n(x) - \kappa P_n(x))^2 \, dx = 0.
\]
Ex.41 を適用する。

46. 定積分の近似値について。被積分関数を$\phi(x)$とし、定積分$\ds\int_{a}^{b} \phi(x)\, dx$を$\tfrac{1}{2}(b - a) \{\phi(a) + \phi(b)\}$で表すときの誤差は$\tfrac{1}{12}M(b - a)^{3}$未満であり、ここで$M$は区間$(a, b)$における$\phi''(x)$の最大値である。また、$\ds\int_{a}^{b} \phi(x)\, dx$を$(b - a)\phi\{\tfrac{1}{2}(a + b)\}$で表すときの誤差は$\tfrac{1}{24}M(b - a)^{3}$未満である。[Exs.\ 4~および5において$f'(x)= \phi(x)$とおく]さらに、$\ds\int_{a}^{b} \phi(x)\, dx$を
\[
\tfrac{1}{6}(b - a)[\phi(a) + \phi(b) + 4\phi\{\tfrac{1}{2}(a + b)\}]
\]
で表すときの誤差は$\tfrac{1}{2880}M(b - a)^{5}$未満であり、ここで$M$は$\phi^{(4)}(x)$の最大値である。[Ex.~6を用いる。この規則は非常に良い近似を与えるもので、Simpsonの法則と呼ばれる。これは先に述べた近似値の一つ thirds と二つ目の三分之二 の合計である]

Simpsonの法則によって与えられる近似値が、直線$x = a$、$x = b$、$y = 0$および曲線$y = \phi(x)$上の点$a$、$\tfrac{1}{2}(a + b)$、$b$を通る平行于$OY$軸の放物線によって囲まれた面積であることを示す。

もし$\phi(x)$が三次多項式であれば、$\phi^{(4)}(x) = 0$であり、Simpsonの法則は正確となる。つまり、点$a$、$\tfrac{1}{2}(a + b)$、$b$を通る三次曲線を無数に描くことができるが、これらすべての曲線は同じ面積を与える。

47. $\phi(x)$が五次多項式であるとき
\[
\int_{0}^{1} \phi(x)\, dx = \tfrac{1}{18}\{5\phi(\alpha) + 8\phi(\tfrac{1}{2}) + 5\phi(\beta)\}
\]
であり、$\alpha$と$\beta$は方程式$x^{2} - x + \frac{1}{10} = 0$の根である。

\item[48.] シンプソンの則を次の公式から円周率$\pi$の計算に適用する：\[
\frac{1}{4}\pi = \int_{0}^{1} \frac{dx}{1 + x^{2}}.
\] [結果は$.7833\dots$である。積分区間を$0$から$\tfrac{1}{2}$と$\tfrac{1}{2}$から$1$に分割し、それぞれの積分に対してシンプソンの則を適用すると、結果は$.785\MS391\MS6\dots$となる。正確な値は$.785\MS398\MS1\dots$である。]

\item[49.] 次の不等式が成り立つことを示す：\[
8.9 < \int_{3}^{5} \sqrt{4 + x^{2}}\, dx < 9.
\] [1903年]

\item[50.] 次の積分を小数点以下二桁まで計算する：\[
\int_{0}^{1} \frac{dx}{1 + x},\quad
\int_{0}^{1} \frac{dx}{\sqrt{1 + x^{4}}},\quad
\int_{0}^{\pi} \sqrt{\sin x}\, dx,\quad
\int_{0}^{\pi} \frac{\sin x}{x}\, dx.
\] [最後の積分において$x = 0$では被積分関数が定義されていないが、$x = 0$における値を1と定めることにより、積分区間全体で連続となる。]

\chapter{第8章：Chapter VIII: THE CONVERGENCE OF INFINITE SERIES AND INFINITE INTEGRALS}

第8章 不定期級数および不定積分の収束

165．第四章で、無限級数が収束する、発散する、または振動するとは何を意味するかを説明し、幾何級数
\[
1 + x + x^{2} + \dots
\]
およびそれに密接に関連した他の級数からいくつかの簡単な例でその定義を示しました。この章では、より系統的な方法で主題を追跡し、解析において頻繁に出現する最も単純な級数が収束するかどうかを判定できる多くの定理を証明します。

しばしば以下の表記を使用します：
\[
u_{m} + u_{m+1} + \dots + u_{n} = \sum_{m}^{n} \phi(\nu),
\]
無限級数 $u_{0} + u_{1} + u_{2} + \dots$ は $\sum\limits_{0}^{\infty} u_{n}$、または単に $\sum u_{n}$ と書きます。\footnote
  {当然、我々の系列を $u_{1} + u_{2} + \dots$（第四章では如此）または $u_{0} + u_{1} + \dots$（ここでは如此）で表すことに何らかの差異はありません。この章の後半では、型 $a_{0} + a_{1}x + a_{2}x^{2} + \dots$ の系列に関心を持ちます：これらのためには後者の表記が明らかに便利です。したがって、これを使用する標準的な表記として採用します。しかし、これを系統的に使用することはありませんし、$u_{1}$ が最初の項であることがより便利な場合を仮定します。例えば、級数 $1 + \frac{1}{2} + \frac{1}{3} + \dots$ を扱う際には、$u_{n} = 1/n$ とし、系列は $u_{1}$ から始まる方がより便利です。一方で、$u_{n} = 1/(n + 1)$ とし、系列は $u_{0}$ から始まる方がより便利な場合もあります。この注意点は、例 lxviii.4 にも適用されます。}

166. 正項級数の収束論は、考慮されるすべての項が正である場合に比較的単純です。\footnote{ここおよびその後の文脈において、「正」は0を含むものとみなされます。}このような級数について最初に考える理由は、それらが最も扱いやすいだけでなく、負または複素な項を含む級数の収束性の議論がしばしば単純な正項級数の議論に依存するからです。

級数の収束性や発散性について議論しているときには、有限個の項は無視できます。したがって、級数に有限個の負または複素な項のみを含む場合、それらを省略し、残りの項に対して以下の定理を適用することができます。

167. 以下は、§77で証明された基本的な定理です。

\par\medskip\noindent\textbf{【結果】}\par
\Item{A.} 正項級数は収束するか発散するかでなければならず、振動することはできません。
\par\medskip

\par\medskip\noindent\textbf{【結果】}\par
\Item{B.} \sum u\\textit{{n} が収束するための必要十分条件は、ある定数 K が存在して
\[
u_{0} + u_{1} + \dots + u_{n} < K
\]
がすべての n に対して成り立つことです。
\par\medskip

\par\medskip\noindent\textbf{【結果】}\par
\Topic{\Item{C.} 比較定理。} \sum u}{n} が収束し、v${n} \leq u${n} がすべての n に対して成り立つ場合、\sum v\\textit{{n} は収束し、\sum v}{n} \leq \sum u\\textit{{n} です。より一般的には、v}{n} \leq Ku\\textit{{n} で K は定数である場合、\sum v}{n} は収束し、\sum v\\textit{{n} \leq K \sum u}{n} です。さらに、\sum u\\textit{{n} が発散し、v}{n} \geq Ku\\textit{{n} がすべての n に対して成り立つ場合、\sum v}{n} は発散します。\footnote{この定理の最後の部分は §77 で明示的に述べられていませんが、読者は証明を容易に補完することができます。}
\par\medskip

さらに、これらのテストを使用して級数 $\sum v_{n}$ の収束性または発散性を推定するには、そのテストが十分に大きい $n$ に対して満足することだけを知る必要がある。すなわち、すべての $n > n_{0}$（一定の値 $n_{0}$ より大きい $n$）に対してテストが満足すればよい。しかし、当然のこととして、$\sum v_{n} \leq K \sum u_{n}$ という結論はこの場合必ず成り立つわけではなく、$K$ の値によります。

この定理の特に有用なケースは、

\par\medskip\noindent\textbf{【結果】}\par
\Item{D.} $\sum u_{n}$ が収束（発散）であり、かつ $u_{n}/v_{n}$ が $n \to \infty$ のとき一定の値（ゼロ以外）に収束するならば、$\sum v_{n}$ も収束（発散）である。
\par\medskip

\Paragraph{168. このテストの最初の応用。} 目前、特定の級数クラスの収束性に関する重要な事実は、$r < 1$ のとき $\sum r^{n}$ が収束し、$r \geq 1$ のとき発散することである。\footnote
  {この章では $r$ を常に正またはゼロの数で表すことにする。}
したがって、定理 C を適用して $u_{n} = r^{n}$ と仮定すると自然である。これによりすぐに、

\par\medskip\noindent\textbf{【結果】}\par
\Item{1.} $v_{n} \leq Kr^{n}$（$r < 1$）が十分に大きい $n$ に対して常に成り立つとき、級数 $\sum v_{n}$ は収束する。
\par\medskip

$K = 1$ の場合、この条件は $v_{n}^{1/n} \leq r$ の形で表すことができる。したがって、

\par\medskip\noindent\textbf{【結果】}\par
\Item{2.} $v_{n}^{1/n} \leq r$（$r < 1$）が十分に大きい $n$ に対して常に成り立つとき、級数 $\sum v_{n}$ は収束する。
\par\medskip

発散性のテストも存在し、

\par\medskip\noindent\textbf{【結果】}\par
\Item{2\ia.} $v_{n}^{1/n} \geq 1$ が無限個の $n$ の値に対して成り立つとき、級数 $\sum v_{n}$ は発散する。
\par\medskip

この定理はほとんど証明不要で、$v_{n}^{1/n} \geq 1$ は $v_{n} \geq 1$ を含意している。定理 2 と 2\ia は非常に広範な応用が可能であるが、一部の目的では収束性の異なるテストを使用することが便利である。

\par\medskip\noindent\textbf{【結果】}\par
\Item{3.} $v_{n+1}/v_{n} \leq r$, $r < 1$ となるすべて十分大きな$n$に対して級数$\sum v_{n}$が収束する。
\par\medskip

これを証明するために、$n \geq n_{0}$ のとき $v_{n+1}/v_{n} \leq r$ であると仮定すると
\[
v_{n} = \frac{v_{n}}{v_{n-1}} \cdot \frac{v_{n-1}}{v_{n-2}} \dots \frac{v_{n_{0}+1}}{v_{n_{0}}} v_{n_{0}}
  \leq \frac{v_{n_{0}}}{r^{n_{0}}} r^{n};
\]
となり、収束級数$\sum r^{n}$との比較により結果が得られる。このテストはダランベールのテストと呼ばれ、後で見るように、理論的にはカuchyのテストよりも一般性が小さく、ダランベールのテストが適用できる場合にはカuchyのテストも適用可能であり、時には後者では適用できない場合もある。また収束しない級数に対するテストは、収束する級数に対するダランベールのテストよりもずっと一般的ではない。読者は容易に証明できるように、$v_{n+1}/v_{n} \geq r \geq 1$ がすべての$n$に対してまたは十分大きな$n$に対して成り立つ場合、$\sum v_{n}$は収束しないことが真である。しかし、この条件が無限個の$n$に対して成り立つ場合でも、それが常にそうなるとは限らない（例 lxvii. 9 参照）。一方で、定理2\ia\ のテストはそのような無限個の$n$に対してのみ満たされる必要があった。それでもダランベールのテストは実用上非常に有用であり、$v_{n}$が複雑な関数である場合でも、$v_{n+1}/v_{n}$はしばしばより単純で扱いやすいものとなる。

解析において最も単純なケースではしばしば $v_{n+1}/v_{n}$ または $v_{n}^{1/n}$ が$n \to \infty$ のとき極限に収束することがある。\footnote{第九章（\okrickRef{Ch.}{IX}、例 lxxxvii. 36）で証明されるように、$v_{n+1}/v_{n} \to l$ である場合、$v_{n}^{1/n} \to l$ であることが示される。逆は真ではないことを、$n$が奇数のとき$v_{n} = 1$、$n$が偶数のとき$v_{n} = 2$と仮定することで確認できる。}
この極限が$1$未満の場合、定理2または3の条件が満たされることが明らかである。したがって

\par\medskip\noindent\textbf{【結果】}\par
\Item{4.} $v_{n}^{1/n}$ または $v_{n+1}/v_{n}$ が $n \to \infty$ のとき極限をとるが、その極限が1未満であるならば、級数 $\sum v_{n}$ は収束します。
\par\medskip

もしいずれの関数も極限をとり、その極限が1より大きいならば、$\sum v_{n}$ は発散するとほぼ明らかです。この形式的な証明は読者に任せるものとします。しかし $v_{n}^{1/n}$ または $v_{n+1}/v_{n}$ が1に収束するとき、これらのテストは完全に失敗し、また常に1未満であるが無数の値において1に近づくような極限をとるときも同様です。そして、$v_{n+1}/v_{n}$ の関連するテストはその比がしばしば1より小さいときおよびしばしば1より大きいときがある場合にも失敗します。$v_{n}^{1/n}$ がこのような行動をするとき、定理2\ia\ は級数の発散を証明するために十分です。しかし明らかに、いくつかのより精巧なテストが必要となる多くのケースがあります。

\par\medskip\noindent\textbf{【例 LXVII.】}\par
\Item{1.} $k$ が正の有理数であるとき、級数 $\sum n^{k} r^{n}$ に対してカuchy のテストと d'Alembert のテスト（4で特定）を適用します。

[ここで $v_{n+1}/v_{n} = \{(n + 1)/n\}^{k} r \to r$ なので、d'Alembert のテストにより、$r < 1$ のとき級数は収束し、$r > 1$ のとき発散することがすぐにわかります。テストが失敗する場合 $r = 1$: しかしそのとき明らかに発散します。$\lim n^{1/n} = 1$（\Ex{xxvii}. 11）なので、カuchy のテストも同様の結論を導きます.]

\Item{2.} 級数 $\sum(An^{k} + Bn^{k-1} + \dots + K) r^{n}$ を考慮します。$A$ を正と仮定できます。$r^{n}$ の係数を $P(n)$ とすると、$P(n)/n^{k} \to A$ であり、\SecNo[§]{167} D の結果により、この級数は $\sum n^{k} r^{n}$ に似た行動をします。

\Item{3.} 
\[
\sum \frac{An^{k} + Bn^{k-1} + \dots + K}
          {\alpha n^{l} + \beta n^{l-1} + \dots + \kappa} r^{n}\quad
(A > 0,\ \alpha > 0).
\]

[級数は$\sum n^{k-l} r^{n}$のように振舞います。$r = 1, k < l$の場合にはさらに考慮が必要です。]

\Item{4.} 我々は(\okrickRef{第}{IV}章、\MiscEx{IV}問題17)で、級数
\[
\sum \frac{1}{n(n + 1)},\quad
\sum \frac{1}{n(n + 1)\dots (n + p)}
\]
が収束することを見ました。これらに対してカuchyの試験法とd'Alembertの試験法は失敗しますことを示してください。[$u_{n}^{1/n}$ の極限、$(u_{n+1}/u_{n})$ の極限がともに 1 であるためです。]

\Item{5.} $p$が整数で $p \geq 2$ のとき、級数 $\sum n^{-p}$ は収束することを示してください。[$n(n + 1)\dots (n + p - 1)/n^{p}$ の極限が 1 であるため、これは問題4で考慮した級数の収束性から導けます。$p = 1$ のときにはすでに \SecNo[§]{77}、式(7)において収束しないことが示されており、$p \leq 0$ のときも明らかに収束しません。]

\Item{6.} $l > k + 1$ のとき収束し、$l \leq k + 1$ のとき発散する級数
\[
\sum \frac{An^{k} + Bn^{k-1} + \dots + K}
          {\alpha n^{l} + \beta n^{l-1} + \dots + \kappa}
\]
を示してください。

\Item{7.} $m_{n}$が正の整数であり、$m_{n+1} > m_{n}$ のとき、級数 $\sum r^{m_{n}}$ は$r < 1$ のとき収束し、$r \geq 1$ のとき発散します。例えば、級数 $1 + r + r^{4} + r^{9} + \dots$ は$r < 1$ のとき収束し、$r \geq 1$ のとき発散します。

\Item{8.} $r = .1$ のとき24桁、$r = .9$ のとき2桁まで級数 $1 + 2r + 2r^{4} + \dots$ を計算してください。[$r = .1$ のとき、最初の5項で和は$1.200\MS200\MS002\MS000\MS000\MS2$となり、誤差は
\[
2r^{25} + 2r^{36} + \dots
  < 2r^{25} + 2r^{36} + 2r^{47} + \dots
  = 2r^{25}/(1 - r^{11})
  < 3/10^{25}
\]
となります。$r = .9$ のとき、最初の8項で和は$5.458\dots$となり、誤差は$2r^{64}/(1 - r^{17}) < .003$未満です。]

9. \(0 < a < b < 1\) のとき、級数 \(a + b + a^{2} + b^{2} + a^{3} + \dots\) は収束する。Cauchyの試験法はこの級数に対して適用できるが、d'Alembertの試験法は失敗する。［例えば
\[
v_{2n+1}/v_{2n} = (b/a)^{n+1} \to \infty,\quad v_{2n+2}/v_{2n+1} = b(a/b)^{n+2} \to 0.
\]

10. \(r\) をすべての正の値に対して、級数 \(1 + r + \frac{r^{2}}{2!} + \frac{r^{3}}{3!} + \dots\) および \(1 + r + \frac{r^{2}}{2^{2}} + \frac{r^{3}}{3^{3}} + \dots\) は収束する。

11. \(\sum u_{n}\) が収束すれば、\(\sum u_{n}^{2}\) および \(\sum \frac{u_{n}}{1 + u_{n}}\) も収束する。

12. \(\sum u_{n}^{2}\) が収束すれば、\(\sum \frac{u_{n}}{n}\) も収束する。［例えば \(2u_{n}/n \leq u_{n}^{2} + (1/n^{2})\) であり、\(\sum (1/n^{2})\) は収束する。］

13. 次の等式を示せ：
\[
1 + \frac{1}{3^{2}} + \frac{1}{5^{2}} + \dots
= \frac{3}{4}\left(1 + \frac{1}{2^{2}} + \frac{1}{3^{2}} + \dots \right)
\]
および
\[
1 + \frac{1}{2^{2}} + \frac{1}{3^{2}}
+ \frac{1}{5^{2}} + \frac{1}{6^{2}}
+ \frac{1}{7^{2}} + \frac{1}{9^{2}} + \dots
= \frac{15}{16} \left(1 + \frac{1}{2^{2}} + \frac{1}{3^{2}} + \dots\right).
\]
ページ分割 313

最初の結果を証明するために、次の定理 \Eq{(8)} および \Eq{(6)} の §77 を用いて
\[
1 + \frac{1}{2^{2}} + \frac{1}{3^{2}} + \dots
= \left(1 + \frac{1}{2^{2}}\right)
+ \left(\frac{1}{3^{2}} + \frac{1}{4^{2}}\right) + \dots
= 1 + \frac{1}{3^{2}} + \frac{1}{5^{2}} + \dots
+ \frac{1}{2^{2}} \left(1 + \frac{1}{2^{2}} + \frac{1}{3^{2}} + \dots\right)
\]
を得る。

\item 14. $\sum (1/n)$ の収束性を否定する反証法で証明せよ。もし級数が収束すると仮定すると、例13で用いた議論により、

\[
1 + \frac{1}{2} + \frac{1}{3} + \cdots
  = (1 + \frac{1}{3} + \frac{1}{5} + \cdots)
  + \frac{1}{2} (1 + \frac{1}{2} + \frac{1}{3} + \cdots),
\]

または

\[
\frac{1}{2} + \frac{1}{4} + \frac{1}{6} + \cdots
  = 1 + \frac{1}{3} + \frac{1}{5} + \cdots
\]

となり、これは明らかに矛盾する。最初の級数の各項は二番目の級数の対応する項よりも小さいからである。

\Paragraph{169.} 級数の収束性や発散性を更に調査する前に、正項の級数に関する重要な一般定理を証明しよう。

\par\medskip\noindent\textbf{【定理 (ディリクレの定理)】}\par
正項の級数の和は項の順序に関係なく同じである。
\par\medskip

この定理は、正項の収束する級数 $u_{0} + u_{1} + u_{2} + \cdots$ があり、それから異なる級数
\[
v_{0} + v_{1} + v_{2} + \cdots
\]

を新たに順序で項を取り出して作ると、その第二の級数も収束し、最初の級数と同じ和を持つことを主張する。もちろん、どの項も省略してはならない：各 $u$ はどこかの $v$ の中に含まれる必要があり、逆もまた然りである。

証明は非常に簡単である。$s$ を $u$ の項の和とする。それならば、$u$ の中から選んだ任意の数の項の和は $s$ よりも大きくなることはない。しかし、すべての $v$ は $u$ であり、したがって $v$ の中から選んだ任意の数の項の和も $s$ よりも大きくなることはない。従って $\sum v_{n}$ は収束し、その和 $t$ は $s$ よりも大きくなることはない。しかし、同じ方法で示すことができるように、$s \leq t$ であることもできる。したがって $s = t$ である。

170. 正の項を持つ級数の積について。Dirichletの定理からすぐに得られる結果は以下の通りである：もし $u_0 + u_1 + u_2 + \dots$ と $v_0 + v_1 + v_2 + \dots$ が二つの収束する正項級数であり、それぞれの和を $s$ と $t$ とするならば、級数
\[
u_0 v_0 + (u_1 v_0 + u_0 v_1)
  + (u_2 v_0 + u_1 v_1 + u_0 v_2) + \dots
\]
は収束し、その和は $st$ である。

すべての可能なペア $u_m v_n$ の積を無限に並べた双方向配列として次のように表すことができる：
\[
ZZZMATHTOKEN0ZZZ
\]
これらの項を様々な方法で単純無限級数の形に再配置することができる。その中には以下のようなものがある。

(1) $m + n = 0$ のときの単項 $u_0 v_0$ から始める；次に $m + n = 1$ の二項 $u_1 v_0, u_0 v_1$ を取り、さらに $m + n = 2$ の三項 $u_2 v_0, u_1 v_1, u_0 v_2$ を加える。そして次々とこのパターンを繰り返すことで、定理の級数
\[
u_0 v_0 + (u_1 v_0 + u_0 v_1)
  + (u_2 v_0 + u_1 v_1 + u_0 v_2) + \dots
\]
を得ることができる。

\Item{(2)} まず、両方の接尾辞がゼロである単項式 $u_{0}v_{0}$ から始めます。次に、接尾辞が1であるがより大きな接尾辞は含まれない項 $u_{1}v_{0}, u_{1}v_{1}, u_{0}v_{1}$ を取り上げます。その後、接尾辞が2でありより大きな接尾辞は含まれない項 $u_{2}v_{0}, u_{2}v_{1}, u_{2}v_{2}, u_{1}v_{2}, u_{0}v_{2}$ を考慮します。このようにして、これらのグループの項の合計はそれぞれ次のようになります：
\begin{multline*}
u_{0}v_{0}, \quad
(u_{0} + u_{1})(v_{0} + v_{1}) - u_{0}v_{0},\\
(u_{0} + u_{1} + u_{2})(v_{0} + v_{1} + v_{2}) - (u_{0} + u_{1})(v_{0} + v_{1}), \quad \dots
\end{multline*}
そして、最初の $n+1$ つのグループの合計は次のようになります：
\[
(u_{0} + u_{1} + \cdots + u_{n})(v_{0} + v_{1} + \cdots + v_{n}),
\]
これは $n \to \infty$ のとき $st$ に収束します。この方法で級数の合計が形成されるとき、最初の一組、二組、三組、\dots の合計はそれぞれ最初の一矩形、二矩形、三矩形、\dots を示す図においてすべての項を含みます。

二番目の方法で形成された級数の合計は $st$ です。しかし、第一の級数（括弧が取り除かれた場合）は第二の級数の再配置であり、したがって、Dirichlet の定理により、それは合計 $st$ に収束します。従って、この定理は証明されました。
\PageSep{315}

\par\medskip\noindent\textbf{【例 LXVIII.】}\par
\Item{1} $r < 1$ のとき、次のことが確認されます：
\[
1 + r^{2} + r + r^{4} + r^{6} + r^{3} + \cdots
= 1 + r + r^{3} + r^{2} + r^{5} + r^{7} + \cdots
= 1/(1 - r).
\]

\Item{2.} $u_{0} + u_{1} + \cdots$ や $v_{0} + v_{1} + \cdots$ のいずれかの級数が発散している場合、それらの項 $(u_{1}v_{0} + u_{0}v_{1}) + (u_{2}v_{0} + u_{1}v_{1} + u_{0}v_{2}) + \cdots$ も発散します。ただし、一方の級数のすべての項がゼロである trivial な場合を除きます。

\item 3. シリーズ $u_{0} + u_{1} + \dots$, $v_{0} + v_{1} + \dots$, $w_{0} + w_{1} + \dots$ が和 $r$, $s$, $t$ に収束するとき、$\lambda_{k} = \sum u_{m}v_{n}w_{p}$ の系列 $\sum \lambda_{k}$ がすべての $m$, $n$, $p$ に対して $m + n + p = k$ となる値を含めて総和を取ることで収束し、その和は $rst$ である。

\item 4. シリーズ $\sum u_{n}$, $\sum v_{n}$ が和 $s$, $t$ に収束するとき、$w_{n} = \sum u_{l}v_{m}$ の系列 $\sum w_{n}$ がすべてのペア $l$, $m$ に対して $lm = n$ となる値を含めて総和を取ることで収束し、その和は $st$ である。

\par\medskip

\Paragraph{171. 收敛与发散的进一步测试。}
第311頁から第313頁の例により、一般的なテスト \SecNo[§]{168} では扱えない単純かつ興味深い正項級数の種類があることが示される。実際、$u_{n+1}/u_{n}$ が $n \to \infty$ のとき極限に収束する最も簡単な系列の場合、一般的なテスト \SecNo[§]{168} はこの極限が $1$ のとき一般に失敗する。したがって、例 lxvii.5 ではこれらのテストが失敗し、特別な装置を用いることになった。その本質は、幾何級数ではなく、例 lxvii.4 の系列を比較系列として使用することだった。

\par\medskip\noindent\textbf{【注釈】}\par
実際、比較によって得られた \SecNo[§]{168} のテストの対象となった幾何級数は収束するだけでなく、非常に速やかに収束し、収束を確保するために必要なほど速く収束するよりも遥かに速い。それゆえ、これらの比較によるテストは自然に粗雑なものであり、より精密なテストが必要となることが多い。
\par\medskip

私たちは\Ex{xxvii}.7で示したように、$r < 1$ の条件下では$n^{k}r^{n} \to 0$ であり、任意の$k$の値に関わらず$n \to \infty$ のときも成り立つ。また、\Ex{lxvii}.1でさらに示したように、級数$\sum n^{k}r^{n}$は収束する。従って、$r < 1$の条件のもとでの数列$r$, $r^{2}$, $r^{3}$, \dots, $r^{n}$, \dotsは、$1^{-k}$, $2^{-k}$, $3^{-k}$, \dots, $n^{-k}$, \dotsよりも急速に減少する。最初の目には矛盾のように思えるが、$r$が単位より大きくないほど小さい値でかつ$k$が大きいときである。

例えば以下の二つの数列を考える：
\[
\dfrac{2}{3}, \quad \dfrac{4}{9}, \quad \dfrac{8}{27}, \dots; \qquad
1, \quad \dfrac{1}{4096}, \quad \dfrac{1}{531441}, \dots
\]
これらはそれぞれ一般項が$(\frac{2}{3})^{n}$と$n^{-12}$である。最初の目には、第二の数列の方が著しく減少しているように思えるが、実際にはそうではない。十分遠くまで数列を調べれば、第一数列の項が非常に小さくなることがわかる。

例えば：
\[
(\frac{2}{3})^{4} = \frac{16}{81} < \frac{1}{5}, \quad
(\frac{2}{3})^{12} < (\frac{1}{5})^{3} < (\frac{1}{10})^{2}, \quad
(\frac{2}{3})^{1000} < (\frac{1}{10})^{166},
\]
そして
\[
1000^{-12} = 10^{-36};
\]
したがって、第一数列の1000番目の項は第二数列の対応する項の10^130分の1以下である。従って、級数$\sum (\frac{2}{3})^{n}$は$\sum n^{-12}$よりも著しく収束し、さらに$\sum n^{-12}$も$\sum n^{-2}$よりも著しく収束する。

このような性質を示す多くの数値結果は、作者の論文「無限の階序」（Cambridge Math. Tracts, No. 12）の付録III（Mr J. Jacksonによって作成された）に見ることができる。この級数$\sum n^{-12}$の和を小数点7桁まで正確に求めることは5項で十分だが、同様の精度を得るには$\sum n^{-2}$の和は約1億万項が必要である。

\Paragraph{172.} 我々は、正項級数の収束性または発散性を確認するためのさらに２つのテスト、マク Lau やカuchy の積分テストとカuchy の凝縮テストを確立することに進みます。これらのテストは非常に完全ではないものの、本書において必要な範囲では十分に一般性があります。

どちらのテストも適用する際には、関数 \( u_{n} \) の性質についてさらなる仮定を行います。これまでに、\( u_{n} \) が正であるという仮定のみをしていたのです。ここでは、\( u_{n} \) は \( n \) とともに減少することを仮定します：すなわち、すべての \( n \) の値に対して \( u_{n+1} \leq u_{n} \) または少なくとも十分に大きな値に対してこの不等式が成り立つことを仮定します。

\par\medskip\noindent\textbf{【注釈】}\par
この条件は最も重要な多くの場合で満たされます。一方から見れば、正項級数の場合には制約としては何ら問題ないかもしれません：ディリクレの定理により、項を再配置しても収束性または発散性に影響を与えることはありませんので、項を大小順に並び替えて得られる減少する項の級数に対してテストを行うことができます。
\par\medskip

しかし、これらの２つのテストを述べる前に、単純で重要な定理を述べて証明します。これをアベルの定理と呼びます。\footnote{この定理はアベルによって発見されましたが忘れられ、プリングシェイムによって再発見されました。}
これは一方的な定理であり、収束性のための十分なテストを提供するのみで、収束性については提供しませんが、上記の２つの定理よりもより基本的な性格を持っています。

\par\medskip\noindent\textbf{【注釈】}\par
\Paragraph{173. アベル（またはプリングシェイム）の定理。} \par\medskip\noindent\textbf{【結果】}\par正項かつ減少する項で収束する級数 \(\sum u_{n}\) の場合、\(\lim nu_{n} = 0\) である。
\par\medskip

\%[** TN: Keeping notation \delta]
$n_{n}$ がゼロに収束しない場合、ある正の数 $\delta$ が存在し、無限個の $n$ の値に対して $n_{n} \geq \delta$ となる。このとき、$n_{1}$ を最初のそのような $n$ の値とする；$n_{2}$ を次に $n_{1}$ より二倍以上大きいそのような $n$ の値とする；$n_{3}$ を次に $n_{2}$ より二倍以上大きいそのような $n$ の値とする；このようにして、$n_{2} > 2n_{1}$, $n_{3} > 2n_{2}$, と続く数列 $n_{1}, n_{2}, n_{3}, \dots$ を得る。すると、$n_{2} - n_{1} > \frac{1}{2}n_{2}$, $n_{3} - n_{2} > \frac{1}{2}n_{3}$, と続く関係が成り立つとともに、$n_{1}u_{n_{1}} \geq \delta$, $n_{2}u_{n_{2}} \geq \delta$, と続く不等式も成立する。また、$u_{n}$ は $n$ の増加に伴って減少することから、

\begin{gather*}
u_{0} + u_{1} + \dots + u_{n_{1} - 1} \geq n_{1}u_{n_{1}} \geq \delta,\\
u_{n_{1}} + \dots + u_{n_{2} - 1} \geq (n_{2} - n_{1})u_{n_{2}} > \frac{1}{2} n_{2}u_{n_{2}} \geq \frac{1}{2} \delta,\\
u_{n_{2}} + \dots + u_{n_{3} - 1} \geq (n_{3} - n_{2})u_{n_{3}} > \frac{1}{2} n_{3}u_{n_{3}} \geq \frac{1}{2} \delta,
\end{gather*}

と続く。したがって、級数 $\sum u_{n}$ の項を括り合わせて新たな級数を得ることができ、その項は発散する級数

\[
\delta + \frac{1}{2} \delta + \frac{1}{2} \delta + \dots
\]

の各項より大きい。従って $\sum u_{n}$ は発散する。

\par\medskip

\par\medskip\noindent\textbf{【例 LXIX.】}\par
1. アベルの定理を用いて、$\sum (1/n)$ と $\sum \{1/(an + b)\}$ が発散することを示せ。[ここで $nu_{n} \to 1$ または $nu_{n} \to 1/a$.]
\par\medskip

2. \( u_n \)が\( n \)に増加するときアベルの定理は成り立たないことを示せ。[級数
\[
1 + \frac{1}{2^2} + \frac{1}{3^2}
  + \frac{1}{4}
  + \frac{1}{5^2} + \frac{1}{6^2} + \frac{1}{7^2} + \frac{1}{8^2}
  + \frac{1}{9}
  + \frac{1}{10^2} + \dots,
\]
において \( u_n = 1/n \) または \( 1/n^2 \)，\( n \)が平方数であるか否かにより定まるとき，この級数は収束する。これは，次のように再配置できるからである：
\[
\frac{1}{2^2} + \frac{1}{3^2}
  + \frac{1}{5^2} + \frac{1}{6^2} + \frac{1}{7^2} + \frac{1}{8^2}
  + \frac{1}{10^2}
  + \dots + \left(1 + \frac{1}{4} + \frac{1}{9} + \dots\right),
\]
これらの級数はすべて収束する。しかし，\( u_n \)が平方数であるとき \( nu_n = 1 \)，従って明らかに \( nu_n \to 0 \) とはならない。

3. アベルの定理の逆は成り立たない。すなわち，\( u_n \)が\( n \)に減少し，\(\lim nu_n = 0\)であるとき，\(\sum u_n\)は収束しないという命題は真ではない。

級数 $\sum(1/n)$ を取り、最初の項に$1$を乗じ、次の項に$\frac{1}{2}$を乗じ、その後の二つの項に$\frac{1}{3}$を乗じ、その後の四つの項に$\frac{1}{4}$を乗じ、その後の八つの項に$\frac{1}{5}$を乗じ、そのような順序で項を乗じます。新しい級数から得られる項を括弧でグループ化すると、

\[
1 + \tfrac{1}{2} · \tfrac{1}{2}
  + \tfrac{1}{3} \left(\tfrac{1}{3} + \tfrac{1}{4}\right)
  + \tfrac{1}{4} \left(\tfrac{1}{5} + \tfrac{1}{6} + \tfrac{1}{7} + \tfrac{1}{8}\right) + \dots;
\]

この級数は収束しないことがわかります。なぜなら、その項は次のように定義される級数よりも大きいからです。

\[
1 + \tfrac{1}{2} · \tfrac{1}{2}
  + \tfrac{1}{3} · \tfrac{1}{2}
  + \tfrac{1}{4} · \tfrac{1}{2} + \dots,
\]

これは収束しないことが明らかです。しかし、級数

\[
1 + \tfrac{1}{2} · \tfrac{1}{2}
  + \tfrac{1}{3} · \tfrac{1}{3}
  + \tfrac{1}{3} · \tfrac{1}{4}
  + \tfrac{1}{4} · \tfrac{1}{5}
  + \tfrac{1}{4} · \tfrac{1}{6} + \dots
\]

の項は$n u_n \to 0$という条件を満たします。実際、$2^{\nu-2} < n \leq 2^{\nu-1}$ のとき $n u_n = 1/\nu$であり、$n \to \infty$ のとき $\nu \to \infty$ です。

174. Maclaurinの（またはCauchyの）積分試験。
$n$が増加するにつれて定義された$u_{n}$が減少すると、$u_{n} = \phi(n)$と書くことができ、かつ$\phi(n)$は連続で減少する関数$\phi(x)$の$x = n$における値であると仮定できる。このとき、任意の正整数$\nu$に対して
\[
\phi(\nu - 1) \geq \phi(x) \geq \phi(\nu)
\]
が$n - 1 \leq x \leq \nu$において成り立つ。次に
\[
v_{\nu}
  = \phi(\nu - 1) - \int_{\nu-1}^{\nu} \phi(x)\, dx
  = \int_{\nu-1}^{\nu} \{\phi(\nu - 1) - \phi(x)\}\, dx,
\]
と定義すると、
\[
0 \leq v_{\nu} \leq \phi(\nu - 1) - \phi(\nu).
\]
となる。このとき$\sum v_{\nu}$は正項級数であり、
\[
v_{2} + v_{3} + \dots + v_{n} \leq \phi(1) - \phi(n) \leq \phi(1)
\]
が成り立つ。したがって$\sum v_{\nu}$は収束し、これより
\[
v_{2} + v_{3} + \dots + v_{n}
\]
または
\[
\sum_{1}^{n-1} \phi(\nu) - \int_{1}^{n} \phi(x)\, dx
\]
は$n \to \infty$のとき正の極限に収束する。

\[
\Phi(\xi) = \int_{1}^{\xi} \phi(x)\, dx
\]とおくと、$\Phi(\xi)$は$\xi$の連続かつ単調増加な関数である。すると
\[
u_{1} + u_{2} + \dots + u_{n-1} - \Phi(n)
\]
は$n \to \infty$のとき正の極限に収束し、その値は$\phi(1)$を超えない。従って$\sum u_{\nu}$は$\Phi(n)$が有限の極限に向かうとき収束し、無限大に向かうとき発散する。そして$\Phi(n)$は単調増加であることから、$\xi \to \infty$のとき$\Phi(\xi)$が有限の極限に向かうときまたは無限大に向かうとき、$\sum u_{\nu}$も同様に収束または発散する。したがって
\par\medskip\noindent\textbf{【結果】}\par
$x > 1$におけるすべての値に対して正で連続な関数$\phi(x)$があり、かつ$x$が増加するにつれて単調減少するとき、級数
\[
\phi(1) + \phi(2) + \dots
\]
は$\xi \to \infty$のとき$\Phi(\xi) = \int_{1}^{\xi} \phi(x)\, dx$が有限の極限$l$に向かうときまたは無限大に向かうとき、収束または発散する。また最初の場合において、級数の和は$\phi(1) + l$を超えない。
\par\medskip

\par\medskip\noindent\textbf{【注釈】}\par
実際にはその和は$\phi(1) + l$より小さい。なぜなら、§160の式（6）と第VII章、Miscellaneous Examples VII 41から$v_{\nu} < \phi(\nu - 1) - \phi(\nu)$であり、$\phi(x) = \phi(\nu)$が区間$(\nu - 1, \nu)$または$[\nu - 1, \nu]$において常に成り立つことは不可能であるからである。
\par\medskip

\par\medskip\noindent\textbf{【例 LXX.】}\par
\Item{1.}級数
\[
\sum_{1}^{\infty} \frac{1}{n^{2} + 1}
\]
は$\tfrac{1}{2} + \tfrac{1}{4}\pi$未満である。

\Item{2.}級数
\[
-\tfrac{1}{2} \pi < \sum_{1}^{\infty} \frac{a}{a^{2} + n^{2}} < \tfrac{1}{2} \pi
\]
が成り立つ。

\Item{3.}もし$m > 0$ならば
\[
\frac{1}{m^{2}} + \frac{1}{(m + 1)^{2}} + \frac{1}{(m + 2)^{2}} + \dots < \frac{m + 1}{m}
\]
である。
\par\medskip

175. 級数 $\sum n^{-s}$ について。積分テストの最も重要な応用は、級数
\[
1^{-s} + 2^{-s} + 3^{-s} + \dots + n^{-s} + \dots,
\]
において $s$ が任意の有理数である場合にあります。既に述べた通り（第77節および練習 lxvii.14, lxix.1）は、$s = 1$ のときこの級数は収束しません。

$s \leq 0$ のとき、この級数は明らかに収束しません。$s > 0$ のとき $u_{n}$ は $n$ が増加するにつれて減少し、テストを適用できます。ここで
\[
\Phi(\xi) = \int_{1}^{\xi} \frac{dx}{x^{s}} = \frac{\xi^{1-s} - 1}{1 - s},
\]
ただし $s \neq 1$ の場合があります。$s > 1$ のとき $\xi^{1-s} \to 0$ であり、
\[
\Phi(\xi) \to \frac{1}{(s - 1)} = l,
\]
とします。また $s < 1$ のとき $\xi^{1-s} \to \infty$ であり、したがって $\Phi(\xi) \to \infty$。

したがって、級数 $\sum n^{-s}$ は $s > 1$ のとき収束し、$s \leq 1$ のとき発散し、最初の場合はその和は $s/(s - 1)$ より小さいという結果を得ます。

以下に注記があります。$s < 1$ の場合の発散については、$\sum (1/n)$ との比較からすぐにこの結果を導くこともできました。しかし、前節の分析が失敗した場合にも積分テストをどのように適用するかを見るのは興味深いです。この場合
\[
\Phi(\xi) = \int_{1}^{\xi} \frac{dx}{x},
\]
であり、$\Phi(\xi) \to \infty$ であることが容易にわかります。もし $\xi > 2^{n}$ のとき
\[
\Phi(\xi) > \int_{1}^{2^{n}} \frac{dx}{x}
  = \int_{1}^{2} \frac{dx}{x}
  + \int_{2}^{4} \frac{dx}{x} + \dots
  + \int_{2^{n-1}}^{2^{n}} \frac{dx}{x}.
\]
とします。ここで $x = 2^{r}u$ と置き換えると
\[
\int_{2^{r}}^{2^{r+1}} \frac{dx}{x} = \int_{1}^{2} \frac{du}{u},
\]
であり、したがって $\ds\Phi(\xi) > n\int_{1}^{2} \frac{du}{u}$ となり、これは $\Phi(\xi) \to \infty$ を示します。

例題 LXXI．

1. $\ds \Phi(\xi) = \int_{1}^{\xi} \frac{dx}{x^{s}}$ において $s < 1$ のとき、$\xi$ と共に無限大に収束することを類似の議論により積分を使わず証明せよ。

2. 級数 $\sum n^{-2}$, $\sum n^{-3/2}$, $\sum n^{-11/10}$ は収束し、それぞれの和は $2$, $3$, $11$ より小さい。級数 $\sum n^{-1/2}$, $\sum n^{-10/11}$ は発散する。

3. $a > 0$ のとき、級数 $\sum \frac{n^{s}}{n^{t} + a}$ は $t > 1 + s$ のとき収束し、$t \leq 1 + s$ のとき発散する。[比較対象：$\sum n^{s-t}$]

4. 次の級数について収束性または発散性を議論せよ：
\[
\tsum \frac{a_{1}n^{s_{1}} + a_{2}n^{s_{2}} + \dots + a_{k}n^{s_{k}}}{b_{1}n^{t_{1}} + b_{2}n^{t_{2}} + \dots + b_{l}n^{t_{l}}},
\]
ここで全ての文字は正の数を表し、$s$'s と $t$'s は有理数であり、大小順に並べられている。

5. 次の不等式を証明せよ：
\begin{gather*}
2\sqrt{n} - 2
  < \frac{1}{\sqrt{1}} + \frac{1}{\sqrt{2}} + \dots + \frac{1}{\sqrt{n}}
  < 2\sqrt{n} - 1, \\
\tfrac{1}{2}\pi
  < \frac{1}{2\sqrt{1}} + \frac{1}{3\sqrt{2}} + \frac{1}{4\sqrt{3}} + \dots
  < \tfrac{1}{2}(\pi + 1).
\end{gather*}
1911年数学トリップ

6. $\phi(n) \to l > 1$ のとき、級数 $\sum n^{-\phi(n)}$ は収束し、$\phi(n) \to l < 1$ のとき発散する。

176．カuchyの凝縮試験法。第2節[§]172で述べた二つの試験法のうちの一つは以下の通りである：もし $u_{n} = \phi(n)$ が $n$ の減少関数であれば、級数 $\sum \phi(n)$ は $\sum 2^{n}\phi(2^{n})$ が収束するか発散するかにより収束するか発散するかを定める。

We can prove this by an argument which we have used already in the special case of the series \(\sum (1/n)\). In the first place,
\begin{gather*}
\phi(3) + \phi(4) \geq 2\phi(4), \\
\phi(5) + \phi(6) + \dots + \phi(8) \geq 4\phi(8), \\
\phi(2^{n} + 1) + \phi(2^{n} + 2) + \dots + \phi(2^{n+1}) \geq 2^{n}\phi(2^{n+1}).
\end{gather*}
If \(\sum 2^{n}\phi(2^{n})\) diverges, then so do \(\sum 2^{n+1}\phi(2^{n+1})\) and \(\sum 2^{n}\phi(2^{n+1})\), and then the inequalities just obtained show that \(\sum \phi(n)\) diverges.
\PageSep{321}

On the other hand,
\[
\phi(2) + \phi(3) \leq 2\phi(2),\quad
\phi(4) + \phi(5) + \dots + \phi(7) \leq 4\phi(4),
\]
and so on. And from this set of inequalities it follows that if \(\sum 2^{n}\phi(2^{n})\) converges, then so does \(\sum \phi(n)\). Thus the theorem is established.

For our present purposes the field of application of this test is practically the same as that of the Integral Test. It enables us to discuss the series \(\sum n^{-s}\) with equal ease. For \(\sum n^{-s}\) will converge or diverge according as \(\sum 2^{n}2^{-ns}\) converges or diverges, i.e., according as \(s > 1\) or \(s \leq 1\).

\par\medskip\noindent\textbf{【例 LXXII.】}\par
\Item{1.} Show that if \(a\) is any positive integer greater than 1 then \(\sum \phi(n)\) is convergent or divergent according as \(\sum a^{n}\phi(a^{n})\) is convergent or divergent. [Use the same arguments as above, taking groups of \(a\), \(a^{2}\), \(a^{3}\), \dots terms.]

\Item{2.} If \(\sum 2^{n}\phi(2^{n})\) converges then it is obvious that \(\lim 2^{n}\phi(2^{n}) = 0\). Hence deduce Abel's Theorem of §173.
\par\medskip

177. 不定積分。第174節の積分試験により、$\phi(x)$が$x$について正で減少する関数であるとき、級数$\sum \phi(n)$は積分関数$\Phi(x) = \int_{1}^{x} \phi(t)\, dt$が無限大に近づき一定の値へ収束するか否かにより収束または発散します。$\Phi(x)$が一定の値へ収束すると仮定し、
\[
\lim_{x \to \infty} \int_{1}^{x} \phi(t)\, dt = l
\]
が成り立つとする。このとき、不定積分
\[
\int_{1}^{\infty} \phi(t)\, dt
\]
を収束し、その値は$l$であると定義します。また、このような積分を無限積分と呼びます。

これまで$\phi(t)$が正で減少する関数であるという仮定をしてきましたが、他の場合にもこの定義を拡張することは自然です。下限が1であることに特別な意味はなく、$t \geq a$のとき連続な関数$\phi(t)$について
\[
\lim_{x \to \infty} \int_{a}^{x} \phi(t)\, dt = l
\]
が成り立つ場合、無限積分
\[
\int_{a}^{\infty}\phi(t)\, dt \quad \Tag{(1)}
\]
を収束し、その値は$l$であると定義します。

上限$a$と下限$A$間の通常の積分（第VII章で定義）は、対比的に有限積分と呼びます。

一方、
\[
\int_{a}^{x}\phi(t)\, dt \to \infty
\]
の場合には、この積分を無限大に発散すると定義し、同様の定義を与えることができます。さらに、これらの選択肢がいずれも成り立たない場合、この積分は$x \to \infty$で有限または無限に振動すると定義します。

これらの定義から以下のような注意点を抽出できます。

\par\medskip\noindent\textbf{【注釈】}\par
\Itemp{(i)} もし
\[
\int_{a}^{x}\phi(t)\, dt = \Phi(x)
\]
と書くなら、積分は$\Phi(x)$が極限に収束するか、無限大（または負の無限大）に発散するか、あるいは振動するかによって収束、発散、または振動します。$x \to \infty$ のとき、$\Phi(x)$ が極限に収束すれば、その値は$\Phi(\infty)$となります。より一般的には、$\Phi(x)$が$\phi(x)$の任意の積分関数であれば、積分の値は$\Phi(\infty) - \Phi(a)$となります。

\Itemp{(ii)} 特殊な場合において、$\phi(t)$ が常に正であるなら、$\Phi(x)$ は$x$ の単調増加関数であり、従って収束または無限大への発散のいずれかのみが可能な選択肢となります。

\Itemp{(iii)} 積分式 \Eq{(1)} は$a$に依存しますが、$t$には独立し、任意の他の文字で$t$を置き換えることで積分の値は変化しません（§157参照）。

\Itemp{(iv)} 読者には「無限積分」という用語が何か明確な値を持つもの（例えば2または$\frac{1}{2}\pi$）を表すために使用されていることに混乱はないと理解していただければと思います。無限積分と有限積分の区別は、無限級数と有限級数の区別に似ています：誰もが無限級数が必ず発散するとは思わないように。

\Itemp{(v)} 積分式 $\ds\int_{a}^{x} \phi(t)\, dt$ は§156および§157で単純な極限、すなわちある有限和の極限として定義されました。したがって無限積分は「極限の極限」、つまり重複した極限となります。無限積分の概念は実質的に有限積分よりもより複雑であり、それがその発展形であるという点で異なります。

\Itemp{(vi)} §174 の積分テストを以下のように述べることができます：もし$\phi(x)$が$x$が増加するにつれて正で安定して減少すれば、無限級数 $\sum\phi(n)$ と無限積分 $\ds\int_{1}^{\infty} \phi(x)\, dx$ は共に収束または発散します。

\Itemp{(vii)} 読者には、無限積分に対する式（1）から式（6）と類似した定理を述べて証明することに困難はありません。すなわち、式（2）に対応する結果は以下の通りです：もし $\ds\int_{a}^{\infty} \phi(x)\, dx$ が収束し、$b > a$ のとき、$\ds\int_{b}^{\infty} \phi(x)\, dx$ も収束し、

\[
\int_{a}^{\infty} \phi(x)\, dx
  = \int_{a}^{b} \phi(x)\, dx + \int_{b}^{\infty}\phi(x)\, dx.
\]

\par\medskip
\par\medskip

\Paragraph{178. $\phi(x)$ が正の場合} シリーズに対する定理 A--D と類似した、無限積分式（1）（第 177 節）の収束性または発散性に関する一般定理を考慮すること自然です。第 177 節で既に述べたように、A は列に対しても積分に対しても真であることが示されています（式（ii）参照）。B に対応する定理は以下の通りです：無限積分式（1）の収束性の必要十分条件は、ある定数 $K$ を用いて、

\[
\int_{a}^{x} \phi(t)\, dt < K
\]

がすべての $x > a$ に対して成り立つことである。

同様に、C に対応する定理は以下の通りです：もし $\ds\int_{a}^{\infty} \phi(x)\, dx$ が収束し、$\psi(x) \leq K\phi(x)$ がすべての $x > a$ に対して成り立つとき、$\ds\int_{a}^{\infty} \psi(x)\, dx$ も収束し、

\[
\int_{a}^{\infty} \psi(x)\, dx \leq K\int_{a}^{\infty} \phi(x)\, dx.
\]

発散性に対する対応するテストは読者に任せます。

We may observe that d'Alembert's test (\S168), depending as it does on the notion of successive terms, has no analogue for integrals; and that the analogue of Cauchy's test is not of much importance, and in any case could only be formulated when we have investigated in greater detail the theory of the function
\[
\phi(x) = r^{x},
\]
as we shall do in Chapter IX. The most important special tests are obtained by comparison with the integral
\[
\int_{a}^{\infty} \frac{dx}{x^{s}} \quad (a > 0),
\]
whose convergence or divergence we have investigated in \S175, and are as follows: 
\par\medskip\noindent\textbf{【結果】}\par
if $\phi(x) < Kx^{-s}$, where $s > 1$, when $x \geq a$, then
\[
\int_{a}^{\infty} \phi(x)\, dx
\]
is convergent; and if $\phi(x) > Kx^{-s}$, where $s \leq 1$, when
$x \geq a$, then the integral is divergent; and in particular, if
$\lim x^{s}\phi(x) = l$, where $l > 0$, then the integral is convergent or
divergent according as $s > 1$ or $s \leq 1$.
\par\medskip

\par\medskip\noindent\textbf{【注釈】}\par
There is one fundamental property of a convergent infinite series in regard to which the analogy between infinite series and infinite integrals breaks down. If $\sum \phi(n)$ is convergent then $\phi(n) \to 0$; but it is not always true, even when $\phi(x)$ is always positive, that if $\int_{a}^{\infty} \phi(x)\, dx$ is convergent then $\phi(x) \to 0$.
\par\medskip

\[
\int_{0}^{\xi} \phi(x)\, dx < \sum_{0}^{\infty} \frac{1}{(n + 1)^{2}}
\]
が成り立つため、$\ds\int_{0}^{\infty} \phi(x)\, dx$は収束する。しかし、$\phi(x) \to 0$とは限らない。

例LXXIII．
1. $\alpha x^{r} + \beta x^{r-1} + \dots + \lambda$を分子、$Ax^{s} + Bx^{s-1} + \dots + L$を分母とする積分
\[
\int_{a}^{\infty} \frac{\alpha x^{r} + \beta x^{r-1} + \dots + \lambda}
                      {Ax^{s} + Bx^{s-1} + \dots + L}\, dx,
\]
において、$\alpha$と$A$が正であり、かつ$a$が分母の最大根より大きいとき、$s > r + 1$であれば収束し、それ以外は発散する。

2. 次の積分
\[
\int_{a}^{\infty} \frac{dx}{\sqrt{x}},\quad
\int_{a}^{\infty} \frac{dx}{x^{4/3}},\quad
\int_{a}^{\infty} \frac{dx}{c^{2} + x^{2}},\quad
\int_{a}^{\infty} \frac{x\, dx}{c^{2} + x^{2}},\quad
\int_{a}^{\infty} \frac{x^{2}\, dx}{c^{2} + x^{2}},\quad
\int_{a}^{\infty} \frac{x^{2}\, dx}{\alpha + 2\beta x^{2} + \gamma x^{4}}
\]
は収束するか否か。最初の二つの積分において$a > 0$とし、最後のものは分母の最大根（存在すれば）より大きい$a$とする。

3. 次の積分
\[
\int_{a}^{\xi} \cos x\, dx,\quad
\int_{a}^{\xi} \sin x\, dx,\quad
\int_{a}^{\xi} \cos(\alpha x + \beta)\, dx
\]
は$\xi \to \infty$で有限に振動する。

4. 次の積分
\[
\int_{a}^{\xi} x\cos x\, dx,\quad
\int_{a}^{\xi} x^{2}\sin x\, dx,\quad
\int_{a}^{\xi} x^{n} \cos(\alpha x + \beta)\, dx,
\]
において$n$は任意の正整数であれば$\xi \to \infty$で無限に振動する。

5. 不尽積分至負无穷。もし$\ds\int_{\xi}^{a} \phi(x)\, dx$が$\xi \to -\infty$のとき極限$l$に収束するならば、我々は$\ds\int_{-\infty}^{a} \phi(x)\, dx$を収束し、その値が$l$であると定義する。このような積分は前節で議論した積分と同じく、すべての性質において類似している：読者はそれらを形式化することに困難を感じないだろう。

6. 不尽積分从負无穷至正无穷。もし積分
\[
\int_{-\infty}^{a} \phi(x)\, dx,\quad
\int_{a}^{\infty} \phi(x)\, dx
\]
が両方とも収束し、それぞれの値が$k$と$l$であるならば、我々は
\[
\int_{-\infty}^{\infty} \phi(x)\, dx
\]
を収束し、その値が$k + l$であると定義する。

7. 次を証明せよ：
\[
\int_{-\infty}^{0} \frac{dx}{1 + x^{2}}
  = \int_{0}^{\infty} \frac{dx}{1 + x^{2}}
  = \tfrac{1}{2} \int_{-\infty}^{\infty} \frac{dx}{1 + x^{2}}
  = \tfrac{1}{2}\pi.
\]

8. 次を一般に証明せよ：
\[
\int_{-\infty}^{\infty} \phi(x^{2})\, dx = 2\int_{0}^{\infty} \phi(x^{2})\, dx,
\]
ただし積分$\ds\int_{0}^{\infty} \phi(x^{2})\, dx$が収束するとき。

9. 次を証明せよ：
もし$\ds\int_{0}^{\infty} x\phi(x^{2})\, dx$が収束するならば、$\ds\int_{-\infty}^{\infty} x\phi(x^{2})\, dx = 0$。

10. アベルの定理の類似（第173節参照）。もし$\phi(x)$が正で一様に減少し、かつ$\ds\int_{a}^{\infty} \phi(x)\, dx$が収束するならば、$x\phi(x) \to 0$である。この事実を（ア）アベルの定理と積分試験を使用して、（イ）直接的に、第173節の議論に類似した手法で証明せよ。

11. $a = x_{0} < x_{1} < x_{2} < \dots$ であり、$x_{n} \to \infty$、そして$\ds u_{n}= \int_{x_{n}}^{x_{n+1}} \phi(x)\, dx$とするとき、$\ds\int_{a}^{\infty} \phi(x)\, dx$の収束は$\sum u_{n}$の収束を含む。もし$\phi(x)$が常に正であるならば、逆も真である。[一般には逆が真でないことを、$\phi(x) = \cos x$, $x_{n} = n\pi$という例によって示すことができる]。

179. 代入と部分積分による無限積分への応用。§161で議論した定積分の変換規則は、無限積分にも拡張される可能性がある。

(1) 代入による変換。\(\int_{a}^{\infty} \phi(x)\, dx\)が収束する場合を考える。さらに、任意の\(\xi > a\)に対して、§161で述べたように
\[
\int_{a}^{\xi} \phi(x)\, dx = \int_{b}^{\tau} \phi\{f(t)\}f'(t)\, dt,
\]
が成り立つとする。ここで \(a = f(b)\), \(\xi = f(\tau)\) である。さらに、関数関係 \(x = f(t)\) は \(t \to \infty\) のとき \(x \to \infty\) を満たすとする。このとき、\(\tau\) および従って \(\xi\) を無限大に近づけると
\[
\int_{b}^{\infty} \phi\{f(t)\}f'(t)\, dt
\]
が収束し、\(\int_{a}^{\infty} \phi(x)\, dx\) と等しいことがわかる。

一方で、\(\xi \to \infty\) が \(\tau \to -\infty\) のときまたは \(\tau \to c\) のときに成り立つ場合もある。最初の場合には
\[
\int_{a}^{\infty} \phi(x)\, dx = -\lim_{\tau\to-\infty} \int_{\tau}^{b} \phi\{f(t)\}f'(t)\, dt = -\int_{-\infty}^{b} \phi\{f(t)\}f'(t)\, dt
\]
が得られる。次の場合には
\[
\int_{a}^{\infty} \phi(x)\, dx = \lim_{\tau\to c} \int_{b}^{\tau} \phi\{f(t)\}f'(t)\, dt.
\]
この式は§181で再び扱われる。

当然、\(\int_{-\infty}^{a} \phi(x)\, dx\) および \(\int_{-\infty}^{\infty} \phi(x)\, dx\) の類似の結果も存在する。詳細を述べる価値はなく、読者は自分自身でこれらを形式化できるだろう。

\par\medskip\noindent\textbf{【例 LXXIV.】}\par
\Item{1.} $x = t^{\alpha}$ の置換により示すと、$s > 1$ および $\alpha > 0$ のとき
\[
\int_{1}^{\infty} x^{-s}\, dx = \alpha\int_{1}^{\infty} t^{\alpha(1-s) - 1}\, dt;
\]
この結果を直接積分により確認する。

\Item{2.} $\ds\int_{a}^{\infty} \phi(x)\, dx$ が収束すると、$\alpha$ が正のとき
\[
 \alpha\int_{(a-\beta)/\alpha}^{\infty} \phi(\alpha t + \beta)\, dt,
\]
負のとき
\[
 -\alpha\int_{-\infty}^{(a-\beta)/\alpha} \phi(\alpha t + \beta)\, dt
\]
と等しい。

\Item{3.} $\phi(x)$ が $x$ の正の定常減少関数であり、$\alpha$ および $\beta$ が任意の正の数であるとき、級数 $\sum \phi(n)$ の収束は級数 $\sum \phi(\alpha n + \beta)$ の収束を含み、逆もまた然り。

\Item{4.} 次のように示す。
\[
%[** TN: In-line in the original]
\int_{1}^{\infty} \frac{dx}{(1 + x)\sqrt{x}} = \tfrac{1}{2} \pi.
\]

\%[** TN: Added paragraph break]
$x = t^{2}$ と置換する。

\Item{5.} 次のように示す。
\[
\int_{0}^{\infty} \frac{\sqrt{x}}{(1 + x)^{2}}\, dx = \tfrac{1}{2}\pi.
\]

$x = t^{2}$ と置換し、部分積分により計算する。

\Item{6.} $\phi(x) \to h$ となる $x \to \infty$ のとき、$\phi(x) \to k$ となる $x \to -\infty$ のとき
\[
\int_{-\infty}^{\infty} \{\phi(x - a) - \phi(x - b)\}\, dx = -(a - b)(h - k).
\]

\begin{alignat*}{2}
\%[** TN: Re-broken]
\int\\textit{{-\xi'}^{\xi} \{\phi(x - a) - \phi(x - b)\}\, dx
  \&= \int}{-\xi'}^{\xi} \phi(x - a)\, dx \&\&- \int\\textit{{-\xi'}^{\xi} \phi(x - b)\, dx\\
  \&= \int}{-\xi'-a}^{\xi-a} \phi(t)\, dt \&\&- \int\\textit{{-\xi'-b}^{\xi-b} \phi(t)\, dt\\
  \&= \int}{-\xi'-a}^{-\xi'-b} \phi(t)\, dt \&\&- \int\\textit{{\xi-a}^{\xi-b} \phi(t)\, dt.
\end{alignat*}
\PageSep{328}
これらの積分の最初の項は、次の形で表すことができる：
\[
(a - b) k + \int_{-\xi'-a}^{-\xi'-b} \rho\, dt,
\]
ここで $\rho \to 0$ である $\xi' \to \infty$ のとき，最後の積分の絶対値は $|a - b| \kappa$ 以下であり、$\kappa$ は区間 $\DPmod{(-\xi' - a, -\xi' - b)}{[-\xi' - a, -\xi' - b]}$ 全ての点における $\rho$ の最大値である。したがって
\[
\int_{-\xi'-a}^{-\xi'-b} \phi(t)\, dt \to (a - b) k.
\]
二番目の積分は同様に議論することができる。

(2) 部分積分による積分. 部分積分の公式（§161）は
\[
\int_{a}^{\xi} f(x)\phi'(x)\, dx
  = f(\xi)\phi(\xi) - f(a)\phi(a) - \int_{a}^{\xi} f'(x)\phi(x)\, dx.
\]

今、$\xi \to \infty$ とする。このとき、上記の式において $\xi$ を含む任意の二つの項が極限に収束するならば、その他の項も極限に収束し、結果として
\[
\int_{a}^{\infty} f(x)\phi'(x)\, dx
  = \lim_{\xi\to\infty} f(\xi)\phi(\xi) - f(a)\phi(a)
    - \int_{a}^{\infty} f'(x)\phi(x)\, dx.
\]
当然、$-\infty$ までの積分や $-\infty$ から $\infty$ までの積分についても類似の結果が得られる。

\par\medskip\noindent\textbf{【例 LXXV.】}\par
\Item{1.} 次を示せ：
\[
%[** TN: In-line in the original]
\int_{0}^{\infty} \frac{x}{(1 + x)^{3}}\, dx
  = \tfrac{1}{2} \int_{0}^{\infty} \frac{dx}{(1 + x)^{2}}
  = \tfrac{1}{2}.
\]

\Item{2.} $\ds\int_{0}^{\infty} \frac{x^{2}}{(1 + x)^{4}}\, dx = \tfrac{2}{3} \int_{0}^{\infty} \frac{x}{(1 + x)^{3}}\, dx = \tfrac{1}{3}$.
\]

\Item{3.} \( m \) と \( n \) が正の整数であるとき、
\[
I_{m, n} = \int_{0}^{\infty} \frac{x^{m}\, dx}{(1 + x)^{m+n}},
\]
ならば
\[
I_{m, n} = \left\{\frac{m}{m + n - 1}\right\} I_{m-1, n}.
\]
従って、\( I_{m, n} = m! (n - 2)! / (m + n - 1)! \) を証明せよ。

\Item{4.} 同様に、
\[
I_{m, n} = \int_{0}^{\infty} \frac{x^{2m+1}\, dx}{(1 + x^{2})^{m+n}}
\]
ならば
\[
I_{m, n} = \left\{\frac{m}{m + n - 1}\right\} I_{m-1, n},\quad
2I_{m, n} = m! (n - 2)! / (m + n - 1)!.
\]
結果を例題3の結果に \( x = t^{2} \) の置換を適用して確認せよ。

\Paragraph{180. 他の無限積分の種類。} 普通または有限積分の定義において、（1）積分区間が有限であり（2）被積分関数が連続であるという条件が満たされていると仮定した。

しかし、これらの条件を満たさない多くの場合でも、「定積分」の概念を拡張して適用できる可能性がある。例えば、先ほどの節で議論した無限積分は、（1）の条件である積分区間が無限大である点で、第VII章の積分とは異なる。今後、これらの条件（1），（2）のうちの一つが満たされていない場合を考える。特に、被積分関数 \(\phi(x)\) が積分区間に有限個の値 \( x = \xi_{1} \)，\( \xi_{2} \)，\(\dots\) で連続性を失い、これらの例外的な値に \( x \) が一方から近づくと \(\phi(x) \to \infty\) や \(\phi(x) \to -\infty\) となる場合について定義を考える自然さがある。

明らかなように、$\DPmod{(a, A)}{[a, A]}$ 内で $\DPmod{(a, A)}{[a, A]}$ に対して $\emph{一つ}$ のような点 $\xi$ を含む場合のみを考えればよい。このような点が複数存在する場合は、$\DPmod{(a, A)}{[a, A]}$ を有限個のサブ区間に分割し、それぞれのサブ区間内に一つだけ含まれるようにできる；そして、これらのサブ区間ごとの積分値が定義されている場合、全体の区間に対する積分を各サブ区間での積分値の和として定義できる。さらに、$\DPmod{(a, A)}{[a, A]}$ 内の一つの点 $\xi$ は $a$ や $A$ のいずれかの端点に存在すると仮定できる。もし $\xi$ が $a$ と $A$ の間にある場合、$\ds\int_{a}^{A} \phi(x)\, dx$ を
\[
\int_{a}^{\xi} \phi(x)\, dx + \int_{\xi}^{A} \phi(x)\, dx,
\]
として定義できる。これらの積分が適切に定義されていると仮定すればよい。そして、$\xi = a$ とする。明らかに、$\xi = A$ の場合にも、導き出される定義は微小な変更を伴って適用可能であることがわかる。

次に、$\phi(x)$ を $\DPmod{(a, A)}{[a, A]}$ 内の $x = a$ 以外で連続でありながら、$x \to a$ のとき $\phi(x) \to \infty$ となるような関数として考える。このような関数の典型的な例は
\[
\phi(x) = (x - a)^{-s},
\]
$s > 0$ のときである；特に $a = 0$ の場合、$\phi(x) = x^{-s}$ となる。したがって、$s > 0$ のときの
\[
\int_{0}^{A} \frac{dx}{x^{s}},
\Tag{(1)}
\]
をどのように定義するかを考える。

積分 $\ds\int_{1/A}^{\infty} y^{s-2}\, dy$ は $s < 1$ のとき収束し、$\lim\limits_{\eta\to\infty} \ds\int_{1/A}^{\eta} y^{s-2}\, dy$ と意味する（§175）。しかし、変数の置換 $y = 1/x$ をすると、

\[
\int_{1/A}^{\eta} y^{s-2}\, dy = \int_{1/\eta}^{A} x^{-s}\, dx
\]

を得る。したがって、$\lim\limits_{\eta\to\infty} \ds\int_{1/\eta}^{A} x^{-s}\, dx$ または同値に、

\[
% [** TN: Keeping notation \epsilon]
\lim_{\epsilon\to +0} \int_{\epsilon}^{A} x^{-s}\, dx
\]

が $s < 1$ のとき存在する。そして、この極限の値を積分 $\Eq{(1)}$ の値として定義するのが自然である。同様の考慮から、

\[
\int_{a}^{A} (x - a)^{-s}\, dx = \lim_{\epsilon\to +0} \int_{a+\epsilon}^{A} (x - a)^{-s}\, dx
\]

という等式により $\ds\int_{a}^{A} (x - a)^{-s}\, dx$ を定義する。

したがって、以下の一般的な定義を得る：\par\medskip\noindent\textbf{【定義】}\par積分 \[
\int_{a+\epsilon}^{A} \phi(x)\, dx
\]
が $\epsilon \to +0$ のとき極限 $l$ に収束するならば、積分 \[
\int_{a}^{A} \phi(x)\, dx
\]
は収束し、その値を $l$ とする。
\par\medskip

また、$\phi(x) \to \infty$ のとき $x \to A$ のとき、

\[
\int_{a}^{A} \phi(x)\, dx = \lim_{\epsilon \to +0} \int_{a}^{A-\epsilon} \phi(x)\, dx
\]

と定義する。そして、上記の説明により、積分区間 $\DPmod{(a, A)}{[a, A]}$ に $\phi(x)$ の有限個の無限大が含まれる場合にもこの定義を拡張できる。

被積分関数が $x$ が積分区間に含まれる値や値群に向かって $\infty$ または $-\infty$ に収束する積分は、第二種の無限積分と呼ばれる。第一種の無限積分は §177 以降で議論されている。§177 の終わりで述べたほとんどの注記（i）～（vii）は第一種の無限積分だけでなく第二種の無限積分にも適用される。

\par\medskip\noindent\textbf{【注釈】}\par
\Paragraph{181.} いまや、第179節の式(4)を次の形で表すことができる：
\[
\int_{a}^{\infty} \phi(x)\, dx = \int_{b}^{c} \phi\{f(t)\}f'(t)\, dt.
\Tag{(1)}
\]
右辺の積分は、$\tau \to c$ の極限として、範囲 $\DPmod{(b, \tau)}{[b, \tau]}$ 上の対応する積分で定義される、すなわち二種類目の無限積分である。そして、$\phi\{f(t)\}f'(t)$ が $t = c$ において無限大を持つとき、この積分は主に二種類目の無限積分となる。たとえば、$\phi(x) = (1 + x)^{-m}$ で $1 < m < 2$、$a = 0$、および $f(t) = t/(1 - t)$ の場合を考えると、$b = 0$、$c = 1$ であり、式(1)は次のようになる：
\[
\int_{0}^{\infty} \frac{dx}{(1 + x)^{m}} = \int_{0}^{1} (1 - t)^{m-2}\, dt;
\Tag{(2)}
\]
右辺の積分は二種類目の無限積分である。

一方で、$\phi\{f(t)\}f'(t)$ が $t = c$ において連続である場合もある。この場合、
\[
\int_{b}^{c} \phi\{f(t)\}f'(t)\, dt
\]
は有限積分であり、
\[
\lim_{\tau \to c} \int_{b}^{\tau} \phi\{f(t)\}f'(t)\, dt = \int_{b}^{c} \phi\{f(t)\}f'(t)\, dt,
\]
は第160節の定理(10)の推論により成り立つ。この場合、変数置換 $x = f(t)$ によって無限積分が有限積分に変換される。この状況は先ほど考えた例において $m \geq 2$ のときによく見られる。
\par\medskip

\par\medskip\noindent\textbf{【例 LXXVI.】}\par
\Item{1.} $\phi(x)$ が $x = a$ を除き連続であり、かつ $x \to a$ のとき $\phi(x) \to \infty$ である場合、$\ds\int_{a}^{A} \phi(x)\, dx$ が収束するための必要十分条件は、任意に小さな $\epsilon$ に対して
\[
\int_{a+\epsilon}^{A} \phi(x)\, dx < K
\]
となる定数 $K$ を見つけることができることである（第178節参照）。
\par\medskip

数は明瞭に、$a$ と $A$ の間で $A'$ を選べるようである。このとき $\phi(x)$ は $\DPmod{(a, A')}{[a, A']}$ 全ての範囲で正である。もし $\phi(x)$ が $\DPmod{(a, A)}{[a, A]}$ の全体の区間で正であれば、当然 $A'$ と $A$ を同一視できる。次に
\[
\int_{a-\epsilon}^{A} \phi(x)\, dx = \int_{a-\epsilon}^{A'} \phi(x)\, dx + \int_{A'}^{A} \phi(x)\, dx.
\]
上記の式の右辺の最初の積分は $\epsilon$ が減少するにつれて増加し、したがって無限大に収束または無限大へと傾き、結果の真実性は明らかになる。

条件が満たされない場合、$\ds\int_{a-\epsilon}^{A} \phi(x)\, dx \to \infty$ である。そのとき、積分 $\ds\int_{a}^{A} \phi(x)\, dx$ は無限大へと収束するといい、これを発散とする。もし $x \to a + 0$ のときに $\phi(x) \to \infty$ であれば、収束と発散は積分に対する唯一の選択肢であることが明瞭である。同様に、$\phi(x) \to -\infty$ の場合についても議論できる。

2. 次に証明せよ：
\[
\int_{a}^{A} (x - a)^{-s}\, dx = \frac{(A - a)^{1-s}}{1 - s}
\]
が $s < 1$ のとき成り立つ。また、$s \geq 1$ のとき積分は発散する。

3. もし $x \to a + 0$ のときに $\phi(x) \to \infty$ であり、かつ $\phi(x) < K(x - a)^{-s}$ で $s < 1$ であれば、$\ds\int_{a}^{A} \phi(x)\, dx$ は収束する。また、もし $\phi(x) > K(x - a)^{-s}$ で $s \geq 1$ であれば、積分は発散する。これは §178 において述べられた一般の比較定理の特殊な場合である。

4. 次の積分は収束するか発散するか：
\begin{gather*}
\int_{a}^{A} \frac{dx}{\sqrtb{(x - a)(A - x)}},\quad
\int_{a}^{A} \frac{dx}{(A - x)\sqrtp{x - a}},\quad
\int_{a}^{A} \frac{dx}{(A - x)\sqrtp{A - x}},\\
\int_{a}^{A} \frac{dx}{\sqrtp{x^{2} - a^{2}}},\quad
\int_{a}^{A} \frac{dx}{\sqrtp{A^{3} - x^{3}}},\quad
\int_{a}^{A} \frac{dx}{x^{2} - a^{2}},\quad
\int_{a}^{A} \frac{dx}{A^{3} - x^{3}}
\end{gather*}

5. 次の積分は収束し、それぞれの値はゼロである。
\[
\int_{-1}^{1}\frac{dx}{\sqrt{x}},\quad
\int_{a-1}^{a+1} \frac{dx}{\sqrt{x - a}}
\]

6. 積分は収束する。
\[
\int_{0}^{\pi} \frac{dx}{\sqrt{\sin x}}
\]
被積分関数は$x$がそれぞれの極限に向かうとき無限大に傾き、従って収束する。

7. 積分
\[
\int_{0}^{\pi} \frac{dx}{(\sin x)^{s}}
\]
$s < 1$ のときのみ収束する。

8. 積分
\[
\int_{0}^{\frac{1}{2}\pi} \frac{x^{s}}{(\sin x)^{t}}\, dx
\]
$t < s + 1$ のとき収束する。

9. $h > 0$ のとき、積分
\[
\int_{0}^{h} \frac{\sin x}{x^{p}}\, dx
\]
$p < 2$ のとき収束することを示す。また、$0 < p < 2$ のとき、次の積分
\[
\int_{0}^{\pi} \frac{\sin x}{x^{p}} dx,\quad
\int_{\pi}^{2\pi} \frac{\sin x}{x^{p}}\, dx,\quad
\int_{2\pi}^{3\pi} \frac{\sin x}{x^{p}}\, dx,\ \dots
\]
符号を交互に変えて絶対値は減少することを示す。被積分関数の極限が$k\pi$と$(k + 1)\pi$である積分を$x = k\pi + y$という置換で変形する。

10. $0 < p < 2$ のとき、積分
\[
\int_{0}^{h} \frac{\sin x}{x^{p}}\, dx
\]
が最大値を取るのは$h = \pi$のときである。数学的旅行1911年。

11. 積分
\[
\int_{0}^{\frac{1}{2} \pi}(\cos x)^{l}(\sin x)^{m}\, dx
\]
$l > -1$, $m > -1$ のときのみ収束する。

12. $s < 1$ のとき、次の積分は以前の定義のいずれにも直接該当しない。
\[
\int_{0}^{\infty} \frac{x^{s-1}\, dx}{1 + x}
\]
被積分関数は$x \to +0$のときに無限大に傾き、積分区間が無限である。従って、この積分を次のように定義する自然な方法がある。
\[
\int_{0}^{1} \frac{x^{s-1}\, dx}{1 + x}
  + \int_{1}^{\infty} \frac{x^{s-1}\, dx}{1 + x}
\]
これらの積分が収束するときのみ定義される。

{\Loosen 第一の積分は、$0 < s < 1$ のとき第二種の収束無限積分です。第二の積分は、$s < 1$ のとき第一種の収束無限積分です。$s > 1$ のとき、第一の積分は通常の有限積分になりますが、その場合第二の積分は発散します。したがって、$0$ から $\infty$ までの積分は $0 < s < 1$ のときにのみ収束します。}

\Item{13.} 次を証明せよ：
\[
\int_{0}^{\infty} \frac{x^{s-1}}{1 + x^{t}}\, dx
\]
は $0 < s < t$ のときにのみ収束する。

\Item{14.} 次の積分を証明せよ：
\[
\int_{0}^{\infty} \frac{x^{s-1} - x^{t-1}}{1 - x}\, dx
\]
は $0 < s < 1$、$0 < t < 1$ のときにのみ収束する。[被積分関数は $x = 1$ のとき定義が不明瞭ですが、$(x^{s-1} - x^{t-1})/(1 - x) \to t - s$ として $x \to 1$ の左右から極限を取ることにより、被積分関数は $x = 1$ のとき $t - s$ の値を持つ連続な関数となります。]

被積分関数が特定の点で不連続であり、その不連続性は定義上の欠如によるもので、その点での特定の値を付与することで解消される場合があります。この場合、通常は被積分関数の定義をその点で補完します。例えば
\[
\int_{0}^{\frac{1}{2} \pi} \frac{\sin mx}{x}\, dx,\quad
\int_{0}^{\frac{1}{2} \pi} \frac{\sin mx}{\sin x}\, dx
\]
は $x = 0$ のとき被積分関数を値 $m$ として定義すると通常の有限積分となります。]

\Item{15.} 置換と部分積分について。置換による変形公式および部分積分の公式は、第一種だけでなく第二種の無限積分に対しても当然拡張できます。読者は自己で一般の定理を導出すべきです。これは第 179 節の線に従って行います。

\Item{16.} 部分積分により証明せよ：$s > 0$, $t > 1$ のとき
\[
\int_{0}^{1} x^{s-1}(1 - x)^{t-1}\, dx
  = \frac{t - 1}{s} \int_{0}^{1} x^{s} (1 - x)^{t-2}\, dx.
\]

17. \( s > 0 \)のとき、
\[
\int_{0}^{1} \frac{x^{s-1}\, dx}{1 + x}
  = \int_{1}^{\infty} \frac{t^{-s}\, dt}{1 + t}.
\]

18. \( 0 < s < 1 \)のとき、
\[
\int_{0}^{1} \frac{x^{s-1} + x^{-s}}{1 + x}\, dx
  = \int_{0}^{\infty} \frac{t^{-s}\, dt}{1 + t}
  = \int_{0}^{\infty} \frac{t^{s-1}\, dt}{1 + t}.
\]

19. \( a + b > 0 \)のとき、
\[
\int_{b}^{\infty} \frac{dx}{(x + a)\sqrt{x - b}} = \frac{\pi}{\sqrt{a + b}}.
\]
（1909年）

20. \( x = t/(1 - t) \)という置換により示すと、\( l \)と\( m \)がともに正のとき、
\[
\int_{0}^{\infty} \frac{x^{l-1}}{(1 + x)^{l+m}}\, dx
  = \int_{0}^{1} t^{l-1} (1 - t)^{m-1}\, dt.
\]

21. \( x = pt/(p + 1 - t) \)という置換により示すと、\( l \), \( m \), および \( p \)がすべて正のとき、
\[
\int_{0}^{1} x^{l-1} (1 - x)^{m-1}\, \frac{dx}{(x + p)^{l + m}}
  = \frac{1}{(1 + p)^{l} p^{m}} \int_{0}^{1} t^{l-1} (1 - t)^{m-1}\, dt.
\]

22. \( a \)から\( b \)までの積分
\[
\int_{a}^{b} \frac{dx}{\sqrt{(x - a)(b - x)}} = \pi\quad\text{および}\quad
\int_{a}^{b} \frac{x\, dx}{\sqrt{(x - a)(b - x)}} = \tfrac{1}{2} \pi (a + b),
\]
(i) \( x = a + (b - a)t^{2} \)という置換により、(ii) \( (b - x)/(x - a) = t \)という置換により、および(iii) \( x = a\cos^{2} t + b\sin^{2} t \)という置換により証明する。

23. \( s > -1 \)のとき、
\[
\int_{0}^{\frac{1}{2} \pi} (\sin\theta)^{s}\, d\theta
  = \int_{0}^{1} \frac{x^{s}\, dx}{\sqrt{1 - x^{2}}}
  = \tfrac{1}{2} \int_{0}^{1} \frac{x^{\frac{1}{2}(s-1)}\, dx}{\sqrt{1 - x}}
  = \tfrac{1}{2} \int_{0}^{1} (1 - x)^{\frac{1}{2}(s-1)} \frac{dx}{\sqrt{x}}.
\]

\Item{24.} 以下の方程式を確立せよ：
\begin{align*}
&\int_{0}^{1} \frac{f(x)\, dx}{\sqrt{1 - x^{2}}}
  = \int_{0}^{\frac{1}{2}\pi} f(\sin\theta)\, d\theta,\\
%
&\int_{a}^{b} \frac{f(x)\, dx}{\sqrt{(x - a)(b - x)}}
  = 2\int_{0}^{\frac{1}{2}\pi} f(a\cos^{2}\theta + b\sin^{2}\theta)\, d\theta,\\
%
&\int_{-a}^{a} f\left\{\sqrt{\frac{a - x}{a + x}}\right\} dx
  = 4a\int_{0}^{\frac{1}{2}\pi} f(\tan\theta) \cos\theta \sin\theta\, d\theta\Add{.}
\end{align*}

\Item{25.} 次の積分を証明せよ：
\[
\int_{0}^{1} \frac{dx}{(1 + x)(2 + x) \sqrt{x(1 - x)}}
  = \pi\left(\frac{1}{\sqrt{2}} - \frac{1}{\sqrt{6}}\right)\Add{.}
\]

[置換によって変換する規則を適用する際には時として注意が必要である。以下の例はこれの良い説明を与える。]

次に
\[
J = \int_{1}^{7} (x^{2} - 6x + 13)\, dx.
\]
直接積分により $J = 48$ が得られる。今、置換
\[
y = x^{2} - 6x + 13,
\]
により $x = 3 ± \sqrt{y - 4}$ となる。$x = 1$ のとき $y = 8$ であり、$x = 7$ のとき $y = 20$ であるから、
\[
J = \int_{8}^{20} y\frac{dx}{dy}\, dy
  = ±\tfrac{1}{2}\int_{8}^{20} \frac{y\, dy}{\sqrt{y - 4}}
\]
という結果に至る。不定積分は
\[
\tfrac{1}{3}(y - 4)^{3/2} + 4(y - 4)^{1/2}
\]
であり、これにより $±\frac{80}{3}$ の値が得られ、どの符号を選んでも明らかに誤りである。

The 解説は \(x\) と \(y\) の関係をより精密に考察することで得られる。函数 \(x^{2} - 6x + 13\) は \(x = 3\) で最小値を取るとき、\(y = 4\) である。\(x\) が 1 から 3 へ増加するとき、\(y\) は 8 から 4 へ減少し、\(\frac{dx}{dy}\) は負であり
\[
\frac{dx}{dy} = -\frac{1}{2\sqrtp{y - 4}}
\]
となる。\(x\) が 3 から 7 へ増加するとき、\(y\) は 4 から 20 へ増加し、他の符号を選ばなければならない。したがって
\[
J = \int_{1}^{7} y\, dx
  = \int_{8}^{4} \left\{-\frac{y}{2\sqrtp{y - 4}}\right\} dy
  + \int_{4}^{20} \frac{y}{2\sqrtp{y - 4}}\, dy,
\]
という公式が得られ、これは正しい結果を導く。

同様に、被積分関数 \(x = \arcsin y\) の代入により \(\ds\int_{0}^{\pi} dx = \pi\) を変換するとき、\(dx/dy = 1/\sqrtp{1 - y^{2}}\) もしくは \(dx/dy = -1/\sqrtp{1 - y^{2}}\) であり、これは \(0 \leq x < \frac{1}{2}\pi\) のときもしくは \(\frac{1}{2}\pi < x \leq \pi\) のときである。

例題。以下の積分をそれぞれの代入により確認せよ。
\[
\int_{0}^{1} (4x^{2} - x + \tfrac{1}{16})\, dx,\quad
\int_{0}^{\pi} \cos^{2}x\, dx
\]
被積分関数 \(4x^{2} - x + \frac{1}{16} = y\)、\(x = \arcsin y\) により。

183. 正の項と負の項を持つ級数。無限級数の和や無限積分（第1種または第2種）の定義は、各項や関数の値が正でも負でもよい場合にも適用される。しかし、この章で確立した収束性や発散性の特別なテストと、それらを例示するために用いた例題は、すべてのこれらの値が正である場合にほぼ完全に限定されている。もちろん、すべての値が負である場合も基本的には異なるものではなく、\(u_{n}\) を \(-u_{n\) または \(\phi(x)\) を \(-\phi(x)\) に変更することで前記の場合に帰着される。

級数の場合、$u_{n}$に課された条件が有限項に対してのみ破られることが明示的または暗黙的に許されていると常に想定されてきた。つまり、ある一定の項から全ての項がその条件（例えばすべての項が正である）を満たすことが必要であるだけである。同様に無限積分の場合、条件は特定の値以上のすべての$x$に対してまたは特定の区間$\DPmod{(a, a + \delta)}{[a, a + \delta]}$（$a$近傍で積分の対象が無限大に収束する）において満たされるものと想定されてきた。したがって、当のテストは次の級数
\[
\sum \frac{n^{2} - 10}{n^{4}}
\]
や積分
\[
\int_{1}^{\infty} \frac{3x - 7}{(x + 1)^{3}}\, dx,\quad
\int_{0}^{1} \frac{1 - 2x}{\sqrt{x}}\, dx
\]
に対して適用される。これらはそれぞれ$n \geq 4$で$n^{2} - 10 > 0$、および$x > \frac{7}{3}$で$3x - 7 > 0$、また$0 < x < \frac{1}{2}$で$1 - 2x > 0$である。

しかし、$u_{n}$の符号の変化が系列全体を通じて継続する場合、すなわち正と負の項の数が無限にあるような系列の場合（例えば$1 - \frac{1}{2} + \frac{1}{3} - \frac{1}{4} + \dots$）、または$\phi(x)$が$x \to \infty$で符号を継続的に変化する場合（例えば積分
\[
\int_{1}^{\infty} \frac{\sin x}{x^{s}}\, dx
\]
）、または$x \to a$で$\phi(x)$の不連続点$a$近傍で符号を継続的に変化する場合（例えば積分
\[
\int_{a}^{A} \sin\left(\frac{1}{x - a}\right) \frac{dx}{x - a}
\]
）、収束や発散の問題はより難しくなる。なぜなら、ここでは振動の可能性も考慮しなければならないからである。

この巻においては積分に対するより一般的な問題には触れることはないが、次章で無限に正と負の項を含むいくつかの単純な級数の例について考える必要がある。

184.絶対収束級数について．ある項が正でも負でもよい級数$\sum u_{n}$を考えます．このとき，$|u_{n}| = \alpha_{n}$と定義します．つまり，$u_{n}$が正のとき$\alpha_{n} = u_{n}$であり，$u_{n}$が負のとき$\alpha_{n} = -u_{n}$です．さらに，$v_{n} = u_{n}$または$v_{n} = 0$と定義します．これは$u_{n}$が正のとき$v_{n} = u_{n}$であり，$u_{n}$が負のとき$v_{n} = 0$です．また，$w_{n} = -u_{n}$または$w_{n} = 0$と定義します．これは$u_{n}$が負のとき$w_{n} = -u_{n}$であり，$u_{n}$が正のとき$w_{n} = 0$です；あるいは，$v_{n}$または$w_{n}$は$u_{n}$が正のとき$\alpha_{n}$に等しく，負のとき零となります．すると，明らかに$v_{n}$と$w_{n}$は常に正であり，
\[
u_{n} = v_{n} - w_{n}, \quad
\alpha_{n} = v_{n} + w_{n}.
\]
ページ分割：337

以下は注記です．例えば，級数が$1 - (1/2)^{2} + (1/3)^{2} - \dots$であるとき，$u_{n} = (-1)^{n-1}/n^{2}$であり$\alpha_{n} = 1/n^{2}$です．そして$v_{n} = 1/n^{2}$または$v_{n} = 0$は$n$が奇数のときと偶数のときで，$w_{n} = 1/n^{2}$または$w_{n} = 0$は$n$が偶数のときと奇数のときで定義されます．

現在，二つの場合を区別することができます．

A. 級数$\sum \alpha_{n}$が収束する場合を考えます．例えば上の例では，$\sum \alpha_{n}$は
\[
1 + (1/2)^{2} + (1/3)^{2} + \dots.
\]
となります．このとき，$\sum v_{n}$と$\sum w_{n}$も収束します：これは（Ex. xxx. 18）任意の項からなる級数が正の項の収束級数から選ばれたものであれば，その級数は収束するからです．そして，定理（Eq. (6) of §77）により，$\sum u_{n}$または$\sum(v_{n} - w_{n})$も収束し，$\sum v_{n} - \sum w_{n}$と等しいです．

したがって，以下の定義を導きます．

以下は定義です．

\par\medskip\noindent\textbf{【定義】}\par
級数$\sum \alpha_{n}$または$\sum |u_{n}|$が収束するとき，級数$\sum u_{n}$は絶対収束と呼ばれます．
\par\medskip

以下が訳文です：

そして、上記で証明したことは以下の結果に相当します：\par\medskip\noindent\textbf{【結果】}\parもし$\sum u_{n}$が絶対収束するなら、それは収束します。その正の項と負の項を別々に取り出した系列も収束し、これらの系列の和は正の項の和に加えて負の項の和で表されます。
\par\medskip

以下が注意です：読者は「絶対収束する系列は収束する」という主張が単なる同義語であると誤解しないように注意してください。$\sum u_{n}$が「絶対収束する」と言っているとき、直接$\sum u_{n}$が収束することを主張しているわけではありません。別の系列$\sum |u_{n}|$の収束性を主張していますし、これは事前に$\sum u_{n}$が振動しないことを明らかにすることは決してありません。

以下は例です：LXXVII.

1. 一般収束原理（§84）を利用して、絶対収束する系列は収束することを証明します。$\sum |u_{n}|$が収束しているので、任意の正数$\DELTA$に対して$n_{0}$を選べます。このとき
\[
|u_{n_{1}+1}| + |u_{n_{1}+2}| + \dots + |u_{n_{2}}| < \DELTA
\]
が$n_{2} > n_{1} \geq n_{0}$のときに成り立ちます。従って
\[
|u_{n_{1}+1} + u_{n_{1}+2} + \dots + u_{n_{2}}| < \DELTA
\]
であり、したがって$\sum u_{n}$は収束します。

2. $\sum a_{n}$が正の項からなる収束系列で、$|b_{n}| \leq K a_{n}$であるなら、$\sum b_{n}$は絶対収束します。

3. $\sum a_{n}$が正の項からなる収束系列であるなら、$-1 \leq x \leq 1$のとき$\sum a_{n}x^{n}$は絶対収束します。

4. $\sum a_{n}$が正の項からなる収束系列であるなら、すべての$\theta$に対して$\sum a_{n}\cos n\theta$、$\sum a_{n}\sin n\theta$は絶対収束します。例としては、§88の系列$\sum r^{n}\cos n\theta$、$\sum r^{n}\sin n\theta$が挙げられます。

\Item{5.} 任意の絶対収束級数から選んだ項で構成される級数は絶対収束級数である。[その級数の項の絶対値の級数は、元の級数の項の絶対値の級数から選ばれたものであるためである.]

\Item{6.} $\sum |u_{n}|$が収束するとき、
\[
|\tsum u_{n}| \leq \tsum |u_{n}|
\]
であり、等号が成立する唯一の状況は各項が同一の符号を持つ場合である。

\Paragraph{185. ディリクレの定理への絶対収束級数の拡張} ディリクレの定理（§169）によれば、正項級数の項は任意に並び替えられてもその和には影響を与えない。これにより、任意の絶対収束級数も同様の性質を持つことが容易に理解できる。元の級数$\sum u_{n}$を並び替えて$\sum u'_{n}$とするとき、$\alpha'_{n}$, $v'_{n}$, $w'_{n}$は$u'_{n}$から$\alpha_{n}$, $v_{n}$, $w_{n}$と同じように形成される。すると、$\sum \alpha'_{n}$は$\sum \alpha_{n}$の並び替えであり収束し、同様に$\sum v'_{n}$, $\sum w'_{n}$もそれぞれ$\sum v_{n}$, $\sum w_{n}$の並び替えであるため収束する。また、ディリクレの定理により、$\sum v'_{n} = \sum v_{n}$, $\sum w'_{n} = \sum w_{n}$であり、従って
\[
\tsum u'_{n}
  = \tsum v'_{n} - \tsum w'_{n}
  = \tsum v_{n} - \tsum w_{n}
  = \tsum u_{n}.
\]

\Paragraph{186. 条件付き収束級数} \Item{B.} 現在、項の絶対値の級数$\sum \alpha_{n}$が発散する場合（つまり$\infty$に収束しない場合）を考慮しなければならない。

\par\medskip\noindent\textbf{【定義】}\par
元の級数$\sum u_{n}$が収束し、しかし$\sum |u_{n}|$が発散するとき、その級数は「条件付き収束」と呼ばれる。

まず最初に、\(\sum u_n\) が条件付き収束する場合、\(\SecNo[§]{184}\) の \(\sum v_n\)，\(\sum w_n\) の両方の級数は \(\infty\) に向かって発散することに注意します。これらは明らかに両方とも収束することはできません、なぜならこれには \(\sum (v_n + w_n)\) や \(\sum \alpha_n\) の収束が含まれるからです。そして、もし一つの級数、例えば \(\sum w_n\) が収束し、\(\sum v_n\) が発散する場合、
\[
\sum_{n=0}^{N} u_n = \sum_{n=0}^{N} v_n - \sum_{n=0}^{N} w_n,
\Tag{(1)}
\]
であり、したがって \(N\) と共に \(\infty\) に収束することになります。これは \(\sum u_n\) の収束性と矛盾します。

したがって、\(\sum v_n\)，\(\sum w_n\) は両方とも発散します。式 \Eq{(1)} から明らかであるように、条件付き収束する級数の和は、それぞれ \(n\) と共に \(\infty\) に収束する2つの関数の差の極限です。また、\(\sum u_n\) は正項級数の収束性を持つ性質（\Ex{xxx}.18）や絶対収束する級数（\Ex{lxxvii}.5）が持つ性質（任意の項を選択しても収束する級数を形成すること）を持たなくなります。そして、ディリクレの定理によって規定された性質は条件付き収束する級数にはおそらく存在しないと思われます；少なくとも \SecNo[§]{185} の証明は完全に失敗し、これは \(\sum v_n\) と \(\sum w_n\) のそれぞれの収束性に依存していました。すぐに見るようにこの予想は正当で、現在考慮しているような級数に対してこの定理が成り立たないことがわかります。

\Paragraph{187. 条件付き収束級数の収束テスト} 条件付き収束級数に対するテストは、§167 以降で述べたような単純かつ一般的なものを見つけることは期待できません。収束性が方程式（1）によって示されるように、正項と負項の打ち消しが基本となる級数に対して収束性を定式化する問題は自然に難しくなります。最初の一歩として、条件付き収束級数に対する比較テストは存在しません。

たとえば、$\sum v_{n}$ の収敛性から $\sum u_{n}$ の収敛性を推論したいとする場合があります。その際には
\[
v_{0} + v_{1} + \dots + v_{n},\quad
u_{0} + u_{1} + \dots + u_{n}
\]
を比較しなければなりません。もしすべての $u$ と $v$ が正で、さらに各 $v$ が対応する $u$ よりも小さい場合、すぐに
\[
v_{0} + v_{1} + \dots + v_{n} < u_{0} + \dots + u_{n}
\]
を推論し、したがって $\sum v_{n}$ は収敛すると結論付けることができます。また、$u$ のみが正で、各 $v$ が対応する $u$ よりも数値的に小さい場合、
\[
|v_{0}| + |v_{1}| + \dots + |v_{n}| < u_{0} + \dots + u_{n}
\]
を推論し、したがって $\sum v_{n}$ は絶対収敛すると結論付けることができます。しかし、$u$ の項と $v$ の項が符号に制限がない一般の場合には、
\[
|v_{0}| + |v_{1}| + \dots + |v_{n}| < |u_{0}| + \dots + |u_{n}|
\]
を推論するしかありません。これにより $\sum v_{n}$ の絶対収敛性が $\sum u_{n}$ の絶対収敛性から推論できるが、$\sum u_{n}$ が条件付き収敛である場合、何らの推論もできません。

\Paragraph{例} 略して見ると、級数 $1 - \frac{1}{2} + \frac{1}{3} - \frac{1}{4} + \dots$ は収敛します。しかし、各項が前者の級数の対応する項よりも数値的に小さいにもかかわらず、級数 $\frac{1}{2} + \frac{1}{3} + \frac{1}{4} + \frac{1}{5} + \dots$ は収散します。

従って、このようなテストは、第一章の前半部で与えられたものよりも著しく特殊な性質を持つべきである。

\#\#\# 188. 離散列
最も単純かつ一般的な条件付き収束列は、項が交互に正と負となる「交替列」と呼ばれる列である。この種の最も重要な列の収束性は以下の定理によって証明される。

\#\#\# 定理
$\phi(n)$ が $n$ の正の関数であり、$n \to \infty$ のとき安定的に0に収束する場合、次のような級数
\[
\phi(0) - \phi(1) + \phi(2) - \dots
\]
は収束し、その和は $\phi(0)$ と $\phi(0) - \phi(1)$ の間にある。

ここで $\phi_{0}$, $\phi_{1}$, \dots は $\phi(0)$, $\phi(1)$, \dots であり、
\[
s_{n} = \phi_{0} - \phi_{1} + \phi_{2} - \dots + (-1)^{n}\phi_{n}
\]
とおく。すると
\[
s_{2n+1} - s_{2n-1} = \phi_{2n} - \phi_{2n+1} \geq 0,\quad
s_{2n} - s_{2n-2} = -(\phi_{2n-1} - \phi_{2n}) \leq 0.
\]
したがって $s_{0}$, $s_{2}$, $s_{4}$, \dots, $s_{2n}$, \dots の列は減少し、従って有限の極限に収束するか無限に発散する一方で、$s_{1}$, $s_{3}$, $s_{5}$, \dots, $s_{2n+1}$, \dots の列は増加し、同様に有限の極限に収束するか無限に発散する。しかし
\[
\lim (s_{2n+1} - s_{2n}) = \lim (-1)^{2n+1} \phi_{2n+1} = 0
\]
から、両方の列は有限の極限に収束し、その極限値は等しい。つまり、列 $s_{0}$, $s_{1}$, \dots, $s_{n}$, \dots の極限が存在することを意味する。$s_{0} = \phi_{0}$, $s_{1} = \phi_{0} - \phi_{1}$ から明らかに、この極限値は $\phi_{0}$ と $\phi_{0} - \phi_{1}$ の間にある。

例 LXXVIII．

1. 級数
\[
1 - \frac{1}{2} + \frac{1}{3} - \frac{1}{4} + \dots,\quad
1 - \frac{1}{\sqrt{2}} + \frac{1}{\sqrt{3}} - \frac{1}{\sqrt{4}} + \dots,\\
\sum \frac{(-1)^{n}}{(n + a)},\quad
\sum \frac{(-1)^{n}}{\sqrt{n + a}},\quad
\sum \frac{(-1)^{n}}{(\sqrt{n} + \sqrt{a})},\quad
\sum \frac{(-1)^{n}}{(\sqrt{n} + \sqrt{a})^{2}},
\]
$a > 0$ のとき、これらは条件付き収束する。

2. 級数 $\sum (-1)^{n}(n + a)^{-s}$, $a > 0$ のとき、$s > 1$ のとき絶対収束し、$0 < s \leq 1$ のとき条件付き収束し、$s \leq 0$ のとき振動する。

3. 第 §188 の級数の和はすべての $n$ の値に対して $s_{n}$ と $s_{n+1}$ の間にある；最初の $n$ 項の和を全体の和に代える誤差は、$(n + 1)$ 番目の項の絶対値よりも小さくなる。

4. 級数
\[
\sum \frac{(-1)^{n}}{\sqrt{n} + (-1)^{n}},
\]
を $n = 2$ の項から始めるものとし、最初のいくつかの項の定義上の問題を避けるために。この級数は以下の形で書くことができる：
\[
\sum \left[\left\{
    \frac{(-1)^{n}}{\sqrt{n} + (-1)^{n}}
  - \frac{(-1)^{n}}{\sqrt{n}}\right\}
  + \frac{(-1)^{n}}{\sqrt{n}}\right]
\]
または
\[
\sum \left\{\frac{(-1)^{n}}{\sqrt{n}} - \frac{1}{n + (-1)^{n}\sqrt{n}}\right\}
  = \sum (\psi_{n} - \chi_{n}),
\]
とおく。級数 $\sum \psi_{n}$ は収束するが、$\sum \chi_{n}$ は発散し、すべての項が正であり $\lim n\chi_{n} = 1$ であるから。したがって元の級数は発散するが、これは $\phi_{2} - \phi_{3} + \phi_{4} - \dots$ の形であり $\phi_{n} \to 0$ であるにもかかわらず。この例は定理において $\phi_{n}$ が安定的に $0$ に収束する必要性を示している。読者は容易に確認できるだろう $\sqrt{2n + 1} - 1 < \sqrt{2n} + 1$ であるから、この条件は満たされていない。

5. 188節の条件が満たされ、$\phi_{n}$は定常的に正の極限$l$に収束する場合、級数$\sum (-1)^{n}\phi_{n}$は有限な振動を示します。

6. 確率収束級数の項の再配置による和の変更。級数$1 - \frac{1}{2} + \frac{1}{3} - \frac{1}{4} + \dots$の和を$s$、その最初の$2n$項の和を$s_{2n}$としますから、$\lim s_{2n} = s$。

次に、以下の級数を考えます：
\[
1 + \frac{1}{3} - \frac{1}{2} + \frac{1}{5} + \frac{1}{7} - \frac{1}{4} + \dots
\Tag{(1)}
\]
この級数では、二つの正の項が一つの負の項を続けており、最初の$3n$項の和を$t_{3n}$とします。すると
\begin{align*}
t_{3n}
  &= 1 + \frac{1}{3} + \dots + \frac{1}{4n-1}
       - \frac{1}{2} - \frac{1}{4} - \dots - \frac{1}{2n}\\
  &= s_{2n} + \frac{1}{2n + 1} + \frac{1}{2n + 3} + \dots + \frac{1}{4n - 1}.
\end{align*}

ここで
\[
\lim \left[\frac{1}{2n + 1} - \frac{1}{2n + 2} + \frac{1}{2n + 3} - \dots
         + \frac{1}{4n - 1} - \frac{1}{4n}\right] = 0,
\]
{\Loosen これは括弧内の項の和が明らかに$\frac{n}{(2n + 1)(2n + 2)}$未満であるから；そして}
\[
\lim \left(\frac{1}{2n + 2} + \frac{1}{2n + 4} + \dots + \frac{1}{4n}\right)
  = \frac{1}{2} \lim \frac{1}{n} \sum_{r=1}^{n} \frac{1}{1 + (r/n)}
  = \frac{1}{2} \int_{1}^{2} \frac{dx}{x},
\]
が成立し、\SecNo[§§]{156}と\SecNo{158}により。したがって
\[
\lim t_{3n} = s + \frac{1}{2} \int_{1}^{2} \frac{dx}{x},
\]
\PageSep{342}
従って、級数~\Eq{(1)}の和は$s$ではなく、最後の方程式の右辺となります。後述するところでは、この二つの級数の和の実際的な値を示します：\SecNo[§]{213}と\okrickRef{第9章}{Ch. IX}、\MiscEx{IX} 19。

確率収束級数は常に任意の和に収束させたり、または$\infty$や$-\infty$に向かって発散させることが可能です。証明についてはBromwichの著書『無限級数』第68ページを参照してください。

\[
1 + \frac{1}{\sqrt{3}} - \frac{1}{\sqrt{2}}
  + \frac{1}{\sqrt{5}} + \frac{1}{\sqrt{7}} - \frac{1}{\sqrt{4}} + \dots
\]
ここで、
\[
t_{3n} = s_{2n} + \frac{1}{\sqrt{2n + 1}} + \frac{1}{\sqrt{2n + 3}} + \dots
  + \frac{1}{\sqrt{4n - 1}}
  > s_{2n} + \frac{n}{\sqrt{4n - 1}},
\]
であり、$s_{2n} = 1 - \frac{1}{\sqrt{2}} + \dots - \frac{1}{\sqrt{2n}}$ で、$n \to \infty$ のとき $s_{2n}$ は収束する。

以下に補足説明がある：
\Paragraph{189. アベルの試験とディリクレの試験}
より一般的な試験法があり、これは第188節の試験を特別な場合として含む。それは以下の通りである。

\par\medskip\noindent\textbf{【定理 (ディリクレの試験)】}\par
$\phi_{n}$ が第188節で述べた条件を満足し、$\sum a_{n}$ が収束または有限回振動する任意の数列であるとき、数列
\[
a_{0}\phi_{0} + a_{1}\phi_{1} + a_{2}\phi_{2} + \dots
\]
は収束する。
\par\medskip

読者は容易に確認できる identities が次の式で表されます：
\[
a_{0}\phi_{0} + a_{1}\phi_{1} + \dots + a_{n}\phi_{n}
  = s_{0}(\phi_{0} - \phi_{1})
  + s_{1}(\phi_{1} - \phi_{2}) + \dots
  + s_{n-1}(\phi_{n-1} - \phi_{n}) + s_{n}\phi_{n},
\]
ここで $s_{n} = a_{0} + a_{1} + \dots + a_{n}$ です。現在、級数 $(\phi_{0} - \phi_{1}) + (\phi_{1} - \phi_{2}) + \dots$ は収束します。なぜなら $n$ 項までの和は $\phi_{0} - \phi_{n}$ であり、$\lim \phi_{n} = 0$ であるからです；またすべての項は正です。さらに $\sum a_{n}$ が実際には収束しないとしても有限に振動するので、定数 $K$ を定めることができます。すなわち $|s_{\nu}| < K$ が全ての $\nu$ の値に対して成り立ちます。したがって級数
\[
\tsum s_{\nu}(\phi_{\nu} - \phi_{\nu+1})
\]
は絶対収束し、よって
\[
s_{0}(\phi_{0} - \phi_{1})
  + s_{1}(\phi_{1} - \phi_{2}) + \dots
  + s_{n-1}(\phi_{n-1} - \phi_{n})
\]
は $n \to \infty$ のとき極限に収束します。また $\phi_{n}$、従って $s_{n}\phi_{n}$ も零に収束します。したがって
\[
a_{0}\phi_{0} + a_{1}\phi_{1} + \dots + a_{n}\phi_{n}
\]
は極限に収束し、すなわち級数 $\sum a_{\nu}\phi_{\nu}$ は収束します。

アベルのテストについて。ディリクレのテストよりも頻繁には用いられませんが、時として有用となる別のテストがあります。

$\phi_{n}$ が $n$ の正で減少する関数であり、その極限が $n \to \infty$ のとき必ずしも零であるとは限らないとします。すなわち $\phi_{n}$ に関する仮定は少ないですが、代わりに $\sum a_{n}$ についての仮定を増やします。それは収束することです。すると次の定理が成り立ちます：もし $\phi_{n}$ が $n$ の正で減少する関数であり、$\sum a_{n}$ が収束するとき、$\sum a_{n}\phi_{n}$ も収束します。

なぜなら $\phi_{n}$ は $n \to \infty$ のときに極限 $l$ を持つとします：そして $\lim (\phi_{n} - l) = 0$ です。したがって、ディリクレのテストにより $\sum a_{n}(\phi_{n} - l)$ は収束し、$\sum{a_{n}}$ も収束するので $\sum a_{n}\phi_{n}$ も収束します。

ページ番号：343

\par\medskip\noindent\textbf{【結果】}\par収束する級数は、その項を任意の正で減少する因子で乗じたとしても収束性が保たれる。
\par\medskip
\par\medskip

\par\medskip\noindent\textbf{【例 LXXIX.】}\par
\Item{1.} ディリクレの試験とアベルの試験も、収束性の一般原理（§84）を用いて証明できる。例えば、アベルの試験の条件が満たされているとする。次のように恒等的に表すことができる：
{\setlength{\multlinegap}{0pt}
\begin{multline*}
a_{m}\phi_{m} + a_{m+1}\phi_{m+1} + \dots + a_{n}\phi_{n}
  = s_{m, m}(\phi_{m} - \phi_{m+1}) + s_{m, m+1}(\phi_{m+1} - \phi_{m+2})\\
  + \dots + s_{m, n-1}(\phi_{n-1} - \phi_{n}) + s_{m, n}\phi_{n}\dots,
\Tag{(1)}
\end{multline*}}\%
ここで
\[
s_{m, \nu} = a_{m} + a_{m+1} + \dots + a_{\nu}.
\]

式(1)の左辺は、$h\phi_{m}$ と $H\phi_{m}$ の間にある。ここで $h$ と $H$ は $s_{m, m}$, $s_{m, m+1}$, \dots, $s_{m, n}$ 中で代数的に最小値と最大値である。任意の正の数 $\DELTA$ に対して、$m \geq m_{0}$ のとき $|s_{m, \nu}| < \DELTA$ となるような $m_{0}$ を選べるから、
\[
|a_{m}\phi_{m} + a_{m+1}\phi_{m+1} + \dots + a_{n}\phi_{n}|
  < \DELTA \phi_{m} \leq \DELTA \phi_{1}
\]
が $n > m \geq m_{0}$ のとき成立する。したがって、級数 $\sum a_{n}\phi_{n}$ は収束する。

\Item{2.} $\theta$ が $\pi$ の倍数でないとき、級数 $\sum \cos n\theta$ と $\sum \sin n\theta$ は有限に振動する。例えば、二つの級数の最初の $n$ 項の和を $s_{n}$ と $t_{n}$ とするとき、$z = \Cis\theta$ とし $|z| = 1$, $z \neq 1$ とする。次のように表すことができる：
\[
|s_{n} + it_{n}|
  = \left|\frac{1 - z^{n}}{1 - z}\right|
  \leq \frac{1 + |z^{n}|}{|1 - z|}
  \leq \frac{2}{|1 - z|};
\]
したがって、$|s_{n}|$ と $|t_{n}|$ も $2/|1 - z|$ を超えない。これらの級数が実際には収束していないことは、その $n$ 番目の項がゼロに収束しないこと（例題 xxiv. 7, 8）から従う。
\par\medskip

The sine series converges to zero if \(\theta\) is a multiple of \(\pi\). The cosine series oscillates finitely if \(\theta\) is an odd multiple of \(\pi\) and diverges if \(\theta\) is an even multiple of \(\pi\).

従って、\(\theta_{n}\)が\(n\)の正の関数であり、\(n \to \infty\)のとき定常的に零に収束する場合、級数
\[
\sum \phi_{n} \cos n\theta,\quad
\sum \phi_{n} \sin n\theta
\]
は収束し、ただし最初の級数が\(\theta\)が\(2\pi\)の倍数である場合を除いて。この場合、最初の級数は\(\sum \phi_{n}\)に帰着され、これは収束するかしないか定かでない：第二の級数は恒等的に零となる。もし\(\sum \phi_{n}\)が収束するならば、両方の級数はすべての\(\theta\)に対して絶対収束し（\Ex{lxxvii}.4）、その結果の全趣旨は\(\sum \phi_{n}\)が発散する場合に適用される点にある。この場合、上記の級数は条件付きでかつ絶対的に収束しないことが\Ex{lxxix}.6で証明される。

\(\theta = \pi\)を余弦級数に代入すると、\SecNo[§]{188}の結果に戻るため、\(\cos n\pi = (-1)^{n}\)である。

3. 級数 \(\sum n^{-s} \cos n\theta, \sum n^{-s} \sin n\theta\)は\(s > 0\)の場合に収束し、ただし最初の級数が\(2\pi\)の倍数であり \(0 < s \leq 1\)である場合を除く。

4. 線3の級数は一般に\(s > 1\)の場合絶対収束し、\(0 < s \leq 1\)の場合条件付き収束し、\(s \leq 0\)の場合振動する（\(s = 0\)では有限振動し、\(s < 0\)では無限振動）。例外的な場合を挙げよ。

5. \(\sum a_{n}n^{-s}\)が収束または有限振動するならば、\(t > s\)のとき\(\sum a_{n}n^{-t}\)は収束する。

6. $\phi_{n}$が$n$の正の関数で、$n \to \infty$のとき定常的に0に収束し、かつ$\sum \phi_{n}$が発散する場合、級数$\sum \phi_{n} \cos n\theta$, $\sum \phi_{n} \sin n\theta$は絶対収束しません。ただし、$\theta$が$\pi$の倍数であるときには正弦級数のみ例外です。[例えば、$\sum \phi_{n} |\cos n\theta|$が収束すると仮定します。$\cos^{2} n\theta \leq |\cos n\theta|$より、
\[
\sum \phi_{n} \cos^{2} n\theta \quad \text{または} \quad \tfrac{1}{2} \sum \phi_{n} (1 + \cos 2n\theta)
\]
が収束すると分かります。しかし、$\sum \phi_{n}$が発散し、かつDirichletの試験により$\sum \phi_{n} \cos 2n\theta$が収束する（$\theta$が$\pi$の倍数でない限り）からこれは不可能です。したがって、$\theta$が$\pi$の倍数であるとき、$\sum \phi_{n} |\cos n\theta|$は発散します。読者は正弦級数について同様の議論を書き出しましょう。$\theta$が$\pi$の倍数である場合に失敗するところを注意してください。]

190. 複素項を持つ級数。これまで、すべての項が実数である級数だけを考え来ました。今後は
\[
\tsum u_{n} = \tsum (v_{n} + iw_{n}),
\]
ここで$v_{n}$と$w_{n}$は実数であるような級数について考えます。このような級数の考慮は、当然のことながら何か新しい要素を導入することはありません。級数が収束するためでなく、かつのみ$\sum v_{n}$, $\sum w_{n}$がそれぞれ収束することが必要です。ただし、この種の級数の一つのクラスは非常に重要であり、特別な扱いが必要です。従って、次のように定義します。

定義
$u_{n} = v_{n} + iw_{n}$である級数$\sum u_{n}$が絶対収束するとは、$\sum v_{n}$と$\sum w_{n}$がそれぞれ絶対収束することを意味します。

定理
$u_{n}$の絶対収束の必要十分条件は$\sum |u_{n}|$または$\sum \sqrt{v_{n}^{2} + w_{n}^{2}}$が収束することです。

もし絶対収束する級数 $\sum u_{n}$ が与えられた場合、両方の級数 $\sum |v_{n}|$, $\sum |w_{n}|$ は収束し、したがって $\sum \{|v_{n}| + |w_{n}|\}$ も収束します：しかし
\[
|u_{n}| = \sqrt{v_{n}^{2} + w_{n}^{2}} \leq |v_{n}| + |w_{n}|,
\]
PerPageSep{345}
従って $\sum |u_{n}|$ も収束します。一方で
\[
|v_{n}| \leq \sqrt{v_{n}^{2} + w_{n}^{2}},\quad
|w_{n}| \leq \sqrt{v_{n}^{2} + w_{n}^{2}},
\]
したがって $\sum |u_{n}|$ が収束するとき、$\sum |v_{n}|$ および $\sum |w_{n}|$ も収束します。

絶対収束する級数は明らかに収束します。実部と虚部がそれぞれ収束するからです。また、ディリクレの定理（§169, §185）は、$\sum v_{n}$ および $\sum w_{n}$ の個別の級数に対して適用することで、絶対収束する複素級数にも即座に拡張できます。

以下は注記です。
絶対収束する級数の収束性は、一般の収束原理（例題 lxxvii. 1 参照）から直接導くことも可能です。これは読者に課題として残します。

191. 平方級数。通常の初等解析で扱われる三角関数や対数、指数などの普通の関数の理論において最も重要な部分は、そのような関数を $x^{n}$ の形の級数 $\sum a_{n}x^{n}$ に展開することに関連しています。このような級数は $x$ の平方級数と呼ばれます。私たちは既に、テイラー級数やマクローリン級数（§148）との関連でこの種の展開に遭遇してきました。しかし当時は実変数 $x$ についてのみ考慮していました。今後は複素変数 $z$ の平方級数 $\sum a_{n}z^{n}$ のいくつかの一般的な性質を検討します。

以下は結果です。
A. 平方級数 $\sum a_{n}z^{n}$ はすべての $z$ に対して収束するか、ある領域の値に対して収束するか、または $z = 0$ のみで収束するかもしれません。

それぞれの可能性を例示するために十分なだけの例を挙げます。

\par\medskip\noindent\textbf{【注釈】}\par
\Item{1.} \emph{級数 $\sum \frac{z^n}{n!}$ はすべての $z$ の値に対して収束する。} ここで $u_n = \frac{z^n}{n!}$ とすると
\[
\left|\frac{u_{n+1}}{u_n}\right| = \left|\frac{z}{n+1}\right| \to 0
\]
が$n \to \infty$ のとき任意の $z$ の値に対して成り立つ。したがってダランベールの試験により $\sum |u_n|$ はすべての $z$ の値に対して収束し、元の級数はすべての $z$ の値に対して絶対収束する。後で見るように収束する冪級数は一般に絶対収束する。

\Item{2.} \emph{$\sum n! z^n$ は $z = 0$ 外のすべての $z$ の値に対して収束しない。} ここで $u_n = n! z^n$ とすると
\[
\left|\frac{u_{n+1}}{u_n}\right| = (n + 1)|z|
\]
が$n \to \infty$ のとき無限大に収束し、$z = 0$ 以外の場合である。したがって（参照：§27の例1, 2, 5）第 $n$ 項の絶対値は$n \to \infty$ のときに無限大に収束する；よって級数は除く $z = 0$ の場合を含め収束しない。明らかに $z = 0$ のとき任意の冪級数は収束する。

\Item{3.} \emph{$|z| < 1$ のとき $\sum z^n$ は常に収束し、$|z| \geq 1$ のとき収束しない。} このことは§88で証明された。したがって各可能性の実例を得た。

\par\medskip

\Paragraph{192.} \par\medskip\noindent\textbf{【結果】}\par\Item{B.} 冪級数 $\sum a_n z^n$ が特定の $z$ の値、例えば $z_1 = r_1 (\cos \theta_1 + i \sin \theta_1)$ に対して収束するとき、すべての $|z| < r_1$ の $z$ の値に対して絶対収束する。

\par\medskip

$\lim a_n z_1^n = 0$ であるから $\sum a_n z_1^n$ が収束することから、定数 $K$ を選べるような存在を確認できる。すなわち $|a_n z_1^n| < K$ がすべての $n$ の値に対して成り立つ。ここで $|z| = r < r_1$ のとき
\[
|a_n z^n| = |a_n z_1^n| \left(\frac{r}{r_1}\right)^n < K \left(\frac{r}{r_1}\right)^n,
\]
となり、収束する幾何級数 $\sum (r/r_1)^n$ と比較すると結果がすぐに得られる。

换句话说、級数が点 \( P \) で収束するならば、原点から \( P \) よりも近いすべての点でも絶対収束します。

---

以下に補足例を示す。\( z = z_1 \) のとき級数が有限回振動している場合でも結果は成り立ちます。[もし \( s_n = a_0 + a_1z_1 + \dots + a_nz_1^n \) であるならば、ある定数 \( K \) を見つけることができますので、すべての \( n \) の値に対して \( |s_n| < K \) です。また、\( |a_nz_1^n| = |s_n - s_{n-1}| \leq |s_{n-1}| + |s_n| < 2K \) であり、以前の議論と同様に証明を完了できます。]

---

193. 幾何級数の収束領域。収束円。\( z = r \) を正実軸上の任意の点とする。もし \( z = r \) のとき幾何級数が収束するならば、半径 \( r \) の円 \( |z| = r \) 内のすべての点でも絶対収束します。特に、\( r \) よりも小さいすべての実数 \( z \) に対して収束します。

次に、正実軸上の点 \( r \) を二つのクラスに分けることにしましょう。一つは級数が収束するクラス、もう一つは収束しないクラスです。最初のクラスには少なくとも \( z = 0 \) の一点があります。一方、第二のクラスは存在せず、級数がすべての \( z \) の値に対して収束することもあります。しかし、もし第二のクラスが存在する場合、そして最初のクラスに \( z = 0 \) 以外の点も含まれる場合、明らかに第一のクラスのどの点も第二のクラスのどの点よりも左側にあります。したがって、ある点 \( z = R \) が二つのクラスを分けることがわかります。この点はどちらのクラスにも属する可能性があります。

すると、\( |z| = R \) の円内のすべての点でも絶対収束します。

---

任意の点 \( P \) を考えます。原点 \( O \) から半径 \( R \) よりも小さい円を描き、この円は \( P \) を含むようにします。この円が \( OA \) と交わる点を \( Q \) とするとき、級数は \( Q \) で収束し、従って定理 B より \( P \) でも絶対収束します。

一方、級数
\[
\sum_{n=0}^{\infty} a_n z^n
\]
は円外の任意の点 \(P'\) において収束しません。なぜなら、もし \(P'\) で収束するとすれば、\(O\) より \(P\) に近いすべての点で絶対収束することになりますが、これは矛盾します。実際、この級数は \(A\) と \(Q'\) の間の任意の点でも収束しません（図51参照）。

これまでの場合では、以下の2つの状況を除いて考察してきました：
(1) 正実軸上において収束しない点が \(z = 0\) のみである場合、または
(2) 正実軸上のすべての点で収束する場合。

明らかに、状況 (1) では \(z = 0\) 外では収束せず、状況 (2) では任意の \(z\) について絶対収束します。したがって、以下の結果を得ます：

\par\medskip\noindent\textbf{【結果】}\par
冪級数は次のいずれかの形で収束します：
\Item{(1)} \(z = 0\) とその他の値において収束する；
\Item{(2)} 任意の \(z\) の値について絶対収束する；または
\Item{(3)} 半径 \(R\) の円内のすべての \(z\) の値について絶対収束し、円外の任意の \(z\) の値では収束しない。
\par\medskip

状況 (3) においては、その円が収束円と呼ばれ、その半径が冪級数の収束半径となります。

この一般の結果は、収束円上の級数の振る舞いについて絶対的な情報を与えません。以下の例をみると、実際には非常に多様な可能性があることがわかります。

\par\medskip\noindent\textbf{【例 LXXX.】}\par
\Item{1.} 級数 \(1 + az + a^2 z^2 + \dots\) において \(a > 0\) のとき、収束円の半径は \(1/a\) です。収束円上では任意の点でも収束せず、\(z = 1/a\) では発散し、円上の他のすべての点では有限に振動します。

\Item{2.} 級数 \(\frac{z}{1^2} + \frac{z^2}{2^2} + \frac{z^3}{3^2} + \dots\) の収束円の半径は \(1\) です；収束円上のすべての点で絶対収束します。
\par\medskip

\item 3. さらに、一般に \( |a_{n+1}| / |a_n| \to \lambda \) または \( |a_n|^{1/n} \to \lambda \) となるとき、級数 \( a_0 + a_1 z + a_2 z^2 + \cdots \) の収束半径は \( 1 / \lambda \) である。最初の場合において
\[
\lim |a_{n+1} z^{n+1}| / |a_n z^n| = \lambda |z|,
\]
これは \( |z| < 1 / \lambda \) のとき小さく、\( |z| > 1 / \lambda \) のとき大きいため、この場合ではダランベールの試験法（§168, 3）を用いることができる。二つ目の場合においても同様にカuchyの試験法（§168, 2）を用いることができる。

\item 4. \Topic{対数級数} 級数
\[
z - \frac{1}{2} z^2 + \frac{1}{3} z^3 - \cdots
\]
は（後で明らかになる理由から）「対数的」級数と呼ばれる。この級数の収束半径は Ex. 3 より \( 1 \) である。

収束円周上の \( z \) のとき、\( z = \cos\theta + i\sin\theta \) と表記できるので、級数は次の形に変形される：
\[
\cos\theta - \frac{1}{2} \cos 2\theta + \frac{1}{3} \cos 3\theta - \cdots
+ i(\sin\theta - \frac{1}{2} \sin 2\theta + \frac{1}{3} \sin 3\theta - \cdots).
\]

実部と虚部は両方とも収束するが、絶対収束しない（\( \theta \) が \( \pi \) の奇数倍でない限り）。\( \theta \) が \( \pi \) の奇数倍のとき、\( z = -1 \) となり、級数は \( -1 - \frac{1}{2} - \frac{1}{3} - \cdots \) となり発散し、\(-\infty\) に収束する。したがって、対数的級数はその収束円周上のすべての点で収束し、\( z = -1 \) の点を除く。

5. 二項級数について考える。級数は次のようである：
\[
1 + mz + \frac{m(m - 1)}{2!} z^{2} + \frac{m(m - 1)(m - 2)}{3!} z^{3} + \dots
\]
$m$が正の整数の場合、級数は終了する。一般に
\[
\frac{|a_{n+1}|}{|a_{n}|} = \frac{|m - n|}{n + 1} \to 1,
\]
であるから収束半径は1である。ここでは円周上の収束について議論しないが、それには少し難しい点がある。

194. 平方級数の唯一性について。$z = 0$以外の某些値においても収束する平方級数$\sum a_{n} z^{n}$があり、その和を$f(z)$とするとき、$f(z)$は容易に以下のように表すことができる：
\[
a_{0} + a_{1}z + a_{2}z^{2} + \dots + (a_{n} + \epsilon_{z})z^{n},
\]
ここで$\epsilon_{z} \to 0$である。もし収束半径を$\mu$とし、$|z| < \mu$のとき、$|a_{n}| \mu^{n} < K$（$K$は定数）であり、従って
\begin{align*}
\left|f(z) - \sum_{0}^{n} a_{\nu}z^{\nu}\right|
  &\leq |a_{n+1}| |z^{n+1}| + |a_{n+2}| |z^{n+2}| + \dots\\
  &< K \left(\frac{|z|}{\mu}\right)^{n+1}
       \left(1 + \frac{|z|}{\mu} + \frac{|z|^{2}}{\mu^{2}} + \dots\right)
   = \frac{K |z|^{n+1}}{\mu^{n} (\mu - |z|)},
\end{align*}
ここで$K$は$z$によらない定数である。$\Ex{lv}.~15$から、ある数$\mu$の絶対値より小さい全ての$z$について
\[
\sum a_{n}z^{n} = \sum b_{n}z^{n}
\]
が成り立つとき、$a_{n} = b_{n}$である。この結果は著しく一般化され、現在ではその詳細には触れない。これは同じ関数$f(z)$を異なる平方級数で表すことはできないことを示している。

\Paragraph{195. 級数の乗法} 170節で見たように、正項級数 $\sum u_{n}$ および $\sum v_{n}$ が収束するとき、$\sum u_{n} × \sum v_{n} = \sum w_{n}$ であり、ここで
\[
w_{n} = u_{0}v_{n} + u_{1}v_{n-1} + \dots + u_{n}v_{0}.
\]
今では、$\sum u_{n}$ および $\sum v_{n}$ が絶対収束するすべての場合にこの結果を拡張できます。なぜなら、我々の証明は Dirichlet の定理の単純な応用に過ぎなかったからです。その定理はすでに絶対収束する級数に対して拡張されています。

\par\medskip\noindent\textbf{【例 LXXXI.】}\par
\Item{1.} $|z|$ が列 $\sum a_{n}z^{n}$ や $\sum b_{n}z^{n}$ の収束半径よりも小さいとき、二つの級数の積は $\sum c_{n}z^{n}$ であり、ここで $c_{n} = a_{0}b_{n} + a_{1}b_{n-1} + \dots + a_{n}b_{0}$。

\Item{2.} もし列 $\sum a_{n}z^{n}$ の収束半径が $R$ であり、$f(z)$ は $|z| < R$ のときの級数の和で、かつ $|z|$ が $R$ や 1 よりも小さいとき、$f(z)/(1 - z) = \sum s_{n}z^{n}$ であり、ここで $s_{n} = a_{0} + a_{1} + \dots + a_{n}$。

\Item{3.} $1/(1 - z)$ の級数を二乗することで、$|z| < 1$ のとき $1/(1 - z)^{2} = 1 + 2z + 3z^{2} + \dots$ を証明せよ。

\Item{4.} 同様に、$1/(1 - z)^{3} = 1 + 3z + 6z^{2} + \dots$ を証明し、一般項は $\frac{1}{2}(n + 1)(n + 2)z^{n}$ であることを示せ。

\Item{5.} 負の整数指数に対する二項定理。もし $|z| < 1$ であり、$m$ が正の整数であるとき
\[
\frac{1}{(1 - z)^{m}}
  = 1 + mz + \frac{m(m + 1)}{1·2} z^{2} + \dots
      + \frac{m(m + 1) \dots (m + n - 1)}{1·2 \dots n} z^{n} + \dots.
\]

[定理の真偽を $m$ までのすべての指標に対して仮定し、例題 2 を用いて $\frac{1}{(1 - z)^{m+1}} = \sum s_{n}z^{n}$ であり、ここで
\begin{align*}
s_{n}
  &= 1 + m + \frac{m(m + 1)}{1·2} + \dots
      + \frac{m(m + 1) \dots (m + n - 1)}{1·2 \dots n} \\
  &= \frac{(m + 1)(m + 2) \dots (m + n)}{1·2 \dots n},
\end{align*}
これは数学的帰納法で容易に証明できる。]

6. シリーズの乗算により、次式が成り立つことを証明せよ：
\[
f(m, z) = 1 + \binom{m}{1} z + \binom{m}{2} z^{2} + \dots,
\]
かつ $|z| < 1$ のとき、$f(m, z)f(m', z) = f(m + m', z)$。この式はオイラーによる二項定理の証明の基礎となる。積の級数における $z^{n}$ の係数は以下のようになる：
\[
\binom{m'}{n}
  + \binom{m}{1} \binom{m'}{n - 1}
  + \binom{m}{2} \binom{m'}{n - 2} + \dots
  + \binom{m}{n - 1} \binom{m'}{1}
  + \binom{m}{n}.
\]
これは $m$ と $m'$ の多項式であるが、$m$ と $m'$ が正の整数のとき、二項定理によりこの多項式は $\dbinom{m + m'}{k}$ に簡約される。そして、このような多項式が全ての正の整数値に対して等しいならば、これらは恒等的に等しい。

7. $f(z) = 1 + z + \dfrac{z^{2}}{2!} + \dots$ のとき、$f(z)f(z') = f(z + z')$。級数の $f(z)$ は全ての $z$ の値に対して絶対収束する：そして $u_{n} = \dfrac{z^{n}}{n!}$, $v_{n} = \dfrac{z'^{n}}{n!}$ のとき、$w_{n} = \dfrac{(z + z')^{n}}{n!}$ であることが容易に確認できる。

8. 次の式が成り立つ：
\[
C(z) = 1 - \frac{z^{2}}{2!} + \frac{z^{4}}{4!} - \dots,\quad
S(z) = z - \frac{z^{3}}{3!} + \frac{z^{5}}{5!} - \dots,
\]
すると：
\[
C(z + z') = C(z)C(z') - S(z)S(z'),\quad
S(z + z') = S(z)C(z') + C(z)S(z'),
\]
かつ
\[
\{C(z)\}^{2} + \{S(z)\}^{2} = 1.
\]

9. 多項定理の失敗。$\sum u_{n}$ および $\sum v_{n}$ が絶対収束していないとき、多項定理は常に成り立たないことが示される。$u_{n} = v_{n} = \frac{(-1)^{n}}{\sqrt{n + 1}}$ の場合を考慮すると：
\[
w_{n} = (-1)^{n} \sum_{r=0}^{n} \frac{1}{\sqrt{(r + 1)(n + 1 - r)}}.
\]
しかし、$\sqrt{(r + 1)(n + 1 - r)} \leq \frac{1}{2}(n + 2)$ であり、よって $|w_{n}| > (2n + 2)/(n + 2)$ で、これは $2$ に収束する。したがって $\sum w_{n}$ は収束しないことが確認できる。

 Miscellaneous Examples on Chapter VIII.

\#\#\# 例題

1. 級数 $\sum n^{k}\left(\sqrt{n + 1} - 2\sqrt{n} + \sqrt{n - 1}\right)$ の収束性について議論せよ。ここで $k$ は実数である。

2. 
\[
\sum n^{r} \Delta^{k}(n^{s}),
\]
において
\[
\Delta u_{n} = u_{n} - u_{n+1},\quad
\Delta^{2} u_{n} = \Delta(\Delta u_{n}),
\]
と定義される $\Delta^{k}(n^{s})$ が一般に $n^{s-k}$ の位数を持つことから、級数は $k > r + s + 1$ のときかつ $s$ が $k$ より小さい正の整数でない限り収束する。

3. 
\[
\sum_{1}^{\infty} \frac{n^{2} + 9n + 5}{(n + 1)(2n + 3)(2n + 5)(n + 4)}
  = \frac{5}{36}.
\]
を示せ。

一般項を部分分数分解するとよい。

4. $R(n)$ を任意の $n$ の有理関数とするとき、多項式 $P(n)$ と定数 $A$ を求め、級数 $\sum \{R(n) - P(n) - (A/n)\}$ が収束することを示せ。特に $R(n)$ が $1/(an + b)$, $(an^{2} + 2bn + c)/(\alpha n^{2} + 2\beta n + \gamma)$ のいずれかの場合についても考慮せよ。

5. 
\[
1 - \frac{1}{1 + z}
  + \frac{1}{2} - \frac{1}{2 + z}
  + \frac{1}{3} - \frac{1}{3 + z} + \dots
\]
は $z$ が負の整数でない限り収束する。

6. 
\begin{gather*}
%[** TN: Set on one line in the original]
\sum \sin\frac{a}{n},\quad
\sum \frac{1}{n} \sin\frac{a}{n},\quad
\sum (-1)^{n} \sin\frac{a}{n},\\
\sum \left(1 - \cos\frac{a}{n}\right),\quad
\sum (-1)^{n} n\left(1 - \cos\frac{a}{n}\right),
\end{gather*}
において $a$ は実数である。これらの級数の収束性または発散性を調査せよ。

7. 
\[
\sum_{1}^{\infty} \left(1 + \frac{1}{2} + \frac{1}{3} + \dots + \frac{1}{n}\right)
  \frac{\sin(n\theta + \alpha)}{n},
\]
において $\theta$ と $\alpha$ は実数である。この級数の収束性について議論せよ。

\MathTrip{1989.}

8. シリーズ
\[
1 - \frac{1}{2} - \frac{1}{3} + \frac{1}{4} + \frac{1}{5} + \frac{1}{6}
  - \frac{1}{7} - \frac{1}{8} - \frac{1}{9} - \frac{1}{10} + \cdots,
\]
において、符号が同じ項はそれぞれ1項、2項、3項、4項などと形成されるグループで構成されているが、このシリーズは収束する。しかし、それに対応する系列は有限の振動を示し、そのグループにはそれぞれ1項、2項、4項、8項などと形成されるグループが含まれる。

9. $u_{1}$, $u_{2}$, $u_{3}$, \dots は正で減少する数列であり、その極限値はゼロであるとする。このとき、次の系列は収束する。
\[
u_{1} - \frac{1}{2}(u_{1} + u_{2}) + \frac{1}{3}(u_{1} + u_{2} + u_{3}) - \cdots,
\quad
u_{1} - \frac{1}{3}(u_{1} + u_{3}) + \frac{1}{5}(u_{1} + u_{3} + u_{5}) - \cdots.
\]
ここで、$(u_{1} + u_{2} + \dots + u_{n})/n = v_{n}$ とすると、$v_{1}$, $v_{2}$, $v_{3}$, \dots も正で減少する数列であり、その極限値はゼロである（IV章、IV節の8、27問参照）。これにより最初の系列が収束することを示すことができる；第二の系列については読者に任せる。特に次の系列は収束する。
\[
1 - \frac{1}{2}\left(1 + \frac{1}{2}\right)
  + \frac{1}{3}\left(1 + \frac{1}{2} + \frac{1}{3}\right) - \cdots,\quad
1 - \frac{1}{3}\left(1 + \frac{1}{5}\right)
  + \frac{1}{5}\left(1 + \frac{1}{3} + \frac{1}{5}\right) - \cdots.
\]

10. $u_{0} + u_{1} + u_{2} + \dots$ は正で減少する項から成る発散する系列であり、
\[
\frac{u_{0} + u_{2} + \dots + u_{2n}}{u_{1} + u_{3} + \dots + u_{2n+1}} \to 1.
\]

11. $\alpha > 0$ のとき、$\lim\limits_{p\to\infty} \sum\limits_{n=0}^{\infty} (p + n)^{-1-\alpha} = 0$。

12. $\lim\limits_{\alpha\to 0+} \alpha \sum\limits_{1}^{\infty} n^{-1-\alpha} = 1$。これは§174から導かれる。
\[
0 < 1^{-1-\alpha} + 2^{-1-\alpha} + \cdots + (n - 1)^{-1-\alpha}
    - \int_{1}^{n} x^{-1-\alpha}\, dx \leq 1,
\]
そして、$\sum n^{-1-\alpha}$ の値は $1/\alpha$ と $(1/\alpha) + 1$ の間にあることが容易に導かれる。

13. 無限級数 $\sum\limits_{1}^{\infty} u_{n}$ の和を求めよ、ここで
\[
u_{n} = \frac{x^{n} - x^{-n-1}}{(x^{n} + x^{-n})(x^{n+1} + x^{-n-1}) }
  = \frac{1}{x - 1}
      \left(\frac{1}{x^{n} + x^{-n}} - \frac{1}{x^{n+1} + x^{-n-1}}\right),
\]
は級数が収束するすべての実数 $x$ に対して成り立つ。[もし $|x|$ が1でないなら、級数の和は $\frac{x}{(x - 1)(x^{2} + 1)}$ である。もし $x = 1$ の場合、$u_{n} = 0$ であり、そのときの和も $0$ となる。もし $x = -1$ の場合、$u_{n} = \frac{1}{2}(-1)^{n+1}$ であり、級数は有限に振動する。]

14. 次の無限級数の和を求めよ：
\[
\frac{z}{1 + z} + \frac{2z^{2}}{1 + z^{2}} + \frac{4z^{4}}{1 + z^{4}} + \dots,\quad
\frac{z}{1 - z^{2}} + \frac{z^{2}}{1 - z^{4}} + \frac{z^{4}}{1 - z^{8}} + \dots
\]
（すべての指数が $2$ の冪である）は、それらが収束する限り求めよ。

[最初の級数は $|z| < 1$ のときのみ収束し、その和は $\frac{z}{1 - z}$ となる；二つ目の級数は $|z| < 1$ のときには $\frac{z}{1 - z}$ に収束し、$|z| > 1$ のときには $\frac{1}{1 - z}$ に収束する。]

15. $|a_{n}| \leq 1$ がすべての $n$ に対して成り立つなら、方程式
\[
0 = 1 + a_{1}z + a_{2}z^{2} + \dots
\]
はその絶対値が $\frac{1}{2}$ よりも小さい根を持つことができず、その絶対値が $\frac{1}{2}$ である根を持つ唯一の状況は $a_{n} = -\Cis(n\theta)$ のときであり、この場合 $z = \frac{1}{2} \Cis(-\theta)$ が根となる。

16. 再現級数について。冪級数 $\sum a_{n}z^{n}$ が再現級数と呼ばれるのは、その係数が次の関係式を満たすときです：
\[
a_{n} + p_{1}a_{n-1} + p_{2}a_{n-2} + \dots + p_{k}a_{n-k} = 0,
\Tag{(1)}
\]
ここで $n \geq k$ であり、$p_{1}$, $p_{2}$, \dots, $p_{k}$ は $n$ に依存していません。任意の再現級数は $z$ の有理関数の展開と等価です。これを証明するために、まず $z$ の絶対値が十分小さいときには級数は収束することを観察します。すなわち、$G$ を次の2つの数のうち大きい方として
\[
1,\quad
|p_{1}| + |p_{2}| + \dots + |p_{k}|
\]
と定めます。式 (1) から $|a_{n}| \leq G\alpha_{n}$ であり、ここで $\alpha_{n}$ は前回の係数の中で絶対値が最大のものであるので、$|a_{n}| < KG^{n}$ となります。ここで $K$ は $n$ に依存していません。したがって、$z$ の絶対値が $1/G$ より小さいときには再現級数は収束します。

しかし、級数 $f(z) = \sum a_{n}z^{n}$ を $p_{1}z$, $p_{2}z^{2}$, \dots, $p_{k}z^{k}$ でそれぞれ乗じて加えると、$(k-1)$ 番目の係数以降が全て消え、新たな級数を得ます。この関係式 (1) より
\[
(1 + p_{1}z + p_{2}z^{2} + \dots + p_{k}z^{k})f(z)
  = P_{0} + P_{1}z + \dots + P_{k-1}z^{k-1}
\]
となります。ここで $P_{0}$, $P_{1}$, \dots, $P_{k-1}$ は定数です。多項式 $1 + p_{1}z + p_{2}z^{2} + \dots + p_{k}z^{k}$ は級数の「関係尺度」と呼ばれます。

逆に、既知の結果から任意の有理関数を多項式と特定の部分分数の和として表すことは可能であり、その部分分数は型 $A/(z - a)^{p}$ であることが示されます。また、負の整数指数に対する二項定理から、分母が $z$ で割り切れない任意の有理関数は、$|z|$ の絶対値が十分小さいとき収束する冪級数展開を持つことが分かります。実際、$|z| < \rho$ において $\rho$ は分母の根の絶対値の中で最小のものである（同書IV章，演習IV 18問以降参照）。また、先に述べた議論を逆にすると級数が再帰的であることが明らかになります。したがって、

結果：冪級数が再帰的である必要十分条件は、その展開が $z$ のような有理関数であることである。

17. 差分方程式の解法。式（1）のような形をした関係式 $a_{n}$ に対するものは、定数係数を持つ線形差分方程式と呼ばれる。このような方程式は、例によって十分に説明できる方法で解くことができる。例えば、方程式が
\[
a_{n} - a_{n-1} - 8a_{n-2} + 12a_{n-3} = 0
\]
であるとする。再帰的な冪級数 $\sum a_{n}z^{n}$ を考慮する。式（16）の例と同様に、その和は
\[
\frac{a_{0} + (a_{1} - a_{0}) z + (a_{2} - a_{1} - 8a_{0}) z^{2}}
     {1 - z - 8z^{2} + 12z^{3}}
  = \frac{A_{1}}{1 - 2z} + \frac{A_{2}}{(1 - 2z)^{2}} + \frac{B}{1 + 3z}
\]
と表されることが分かり、ここで $A_{1}$, $A_{2}$, および $B$ は $a_{0}$, $a_{1}$, $a_{2}$ を用いて容易に表現できる数である。各分数を個別に展開すると、$z^{n}$ の係数は
\[
a_{n} = 2^{n}\{A_{1} + (n + 1) A_{2}\} + (-3)^{n} B
\]
と表されることが分かる。$A_{1}$, $A_{2}$, $B$ の値は最初の三つの係数 $a_{0}$, $a_{1}$, $a_{2}$ に依存し、これらは当然任意に選べる。

\item{18.} 差分方程式 \( u_{n} - 2\cos\theta u_{n-1} + u_{n-2} = 0 \) の解は \( u_{n} = A\cos n\theta + B\sin n\theta \) であり、\( A \) と \( B \) は任意の定数である。

\item{19.} \( u_{n} \) が \( n \) の多項式で次数が \( k \) のとき、\( \sum u_{n} z^{n} \) は関係式の周期が \( (1 - z)^{k+1} \) の再発する級数である。

\item{20.} \( \frac{9}{(z - 1)(z + 2)^{2}} \) を \( z \) の昇べきに展開せよ。

\item{21.} \( f(n) \) が \( z/(1 + z + z^{2}) \) の \( z \) のべき級数展開における \( z^{n} \) の係数であるとき、
\[
(1)\quad f(n) + f(n - 1) + f(n - 2) = 0,\quad
(2)\quad f(n) = \frac{\omega_{3}^{n} - \omega_{3}^{2n}}{\omega_{3} - \omega_{3}^{2}},
\]
ここで \( \omega_{3} \) は単位円上の複素立方根である。したがって、\( n \) が \( 3k \), \( 3k + 1 \), \( 3k + 2 \) の形のとき \( f(n) \) は \( 0 \), \( 1 \), \( -1 \) となることを示し、この結果を \( z/(1 + z + z^{2}) = z(1 - z)/(1 - z^{3}) \) の恒等式を利用して確認せよ。

\item{22.} 硬貨を投げて点数を得るプレイヤーが表が出たときは 1 点、裏が出たときは 2 点とし、合計 \( n \) 点以上になるまで続けます。彼の合計 \( n \) 点に達する確率は \( \frac{1}{3} \{ 2 + (-\frac{1}{2})^{n} \} \) であることを証明せよ。

[もし \( p_{n} \) がその確率であるとき、\( p_{n} = \frac{1}{2} (p_{n-1} + p_{n-2}) \) であり、また \( p_{0} = 1 \), \( p_{1} = \frac{1}{2} \) である。]

\item{23.} 
\[
\frac{1}{a + 1} + \frac{1}{a + 2} + \dots + \frac{1}{a + n}
= \binom{n}{1}\frac{1}{a + 1}
- \binom{n}{2}\frac{1!}{(a + 1)(a + 2)} + \dots
\]
が \( n \) が正の整数、かつ \( a \) が \( -1, -2, \ldots, -n \) のいずれでもないとき成り立つことを証明せよ。

各項を部分分数に分解することで得られます。$a > -1$ のとき、式
\[
\int_{0}^{1} x^{a}\frac{1 - x^{n}}{1 - x}\, dx
  = \int_{0}^{1} (1 - x)^{a}\{1 - (1 - x)^{n}\}\frac{dx}{x}
\]
から、$x$ のべき級数展開して各項を別々に積分することで簡単に導けます。この結果は代数的な恒等式であるため、$a = -1, -2, \dots, -n$ 以外のすべての値に対して成り立ちます。

24. シリーズの乗算により証明せよ：
\[
\sum_{0}^{\infty} \frac{z^{n}}{n!} \sum_{1}^{\infty} \frac{(-1)^{n-1}z^{n}}{n·n!}
  = \sum_{1}^{\infty} \left(1 + \frac{1}{2} + \frac{1}{3} + \dots + \frac{1}{n}\right) \frac{z^{n}}{n!}.
\]

$z^{n}$ の係数は
\[
\frac{1}{n!}\left\{
  \binom{n}{1} - \frac{1}{2}\binom{n}{2} + \frac{1}{3}\binom{n}{3} - \dots
\right\}
\]
となる。Ex.~23 を用いて $a = 0$ の場合を取ればよい。

25. $n \to \infty$ のとき $A_{n} \to A, B_{n} \to B$ ならば
\[
(A_{1}B_{n} + A_{2}B_{n-1} + \dots + A_{n}B_{1})/n \to AB.
\]

$A_{n} = A + \epsilon_{n}$ と置く。与えられた式は
\[
A \frac{B_{1} + B_{2} + \dots + B_{n}}{n}
  + \frac{\epsilon_{1}B_{n} + \epsilon_{2}B_{n-1} + \dots + \epsilon_{n}B_{1}}{n}
\]
と等しい。

最初の項は $AB$ に収束する（第IV章、MiscEx IV 27）。第二項の絶対値は $\beta\{|\epsilon_{1}| + |\epsilon_{2}| + \dots + |\epsilon_{n}|\}/n$ 以下であり、$\beta$ は $|B_{\nu}|$ の最大値より大きい任意の数である。この式はゼロに収束する。

26. $c_{n} = a_{1}b_{n} + a_{2}b_{n-1} + \dots + a_{n}b_{1}$ とし
\[
A_{n} = a_{1} + a_{2} + \dots + a_{n},\quad
B_{n} = b_{1} + b_{2} + \dots + b_{n},\quad
C_{n} = c_{1} + c_{2} + \dots + c_{n},
\]
とおくとき、
\[
C_{n} = a_{1}B_{n} + a_{2}B_{n-1} + \dots + a_{n}B_{1}
      = b_{1}A_{n} + b_{2}A_{n-1} + \dots + b_{n}A_{1}
\]
および
\[
C_{1} + C_{2} + \dots + C_{n} = A_{1}B_{n} + A_{2}B_{n-1} + \dots + A_{n}B_{1}.
\]

したがって、級数 $\sum a_n$, $\sum b_n$ が収束し、その和を $A$, $B$ とするとき、$A_n \to A$, $B_n \to B$ であれば、

\[
\frac{C_1 + C_2 + \dots + C_n}{n} \to AB.
\]

従って、$\sum c_n$ が収束すればその和は $AB$ である。この結果は「Abelの級数乗法の定理」と呼ばれる。

\Item{27.} 次を証明せよ：

\begin{align*}
\tfrac{1}{2} \left(1 - \tfrac{1}{2} + \tfrac{1}{3} - \dots\right)^2
  &= \tfrac{1}{2} - \tfrac{1}{3} \left(1 + \tfrac{1}{2}\right)
       + \tfrac{1}{4} \left(1 + \tfrac{1}{2} + \tfrac{1}{3}\right) - \dots,\\
%
\tfrac{1}{2} \left(1 - \tfrac{1}{3} + \tfrac{1}{5} - \dots\right)^2
  &= \tfrac{1}{2} - \tfrac{1}{4} \left(1 + \tfrac{1}{3}\right)
       + \tfrac{1}{6} \left(1 + \tfrac{1}{3} + \tfrac{1}{5}\right) - \dots.
\end{align*}

[Ex. 9 を用いて級数の収束性を証明せよ。]

\Item{28.} $m$ と $n$ のどのような値に対して、積分 $\ds\int_{0}^{\pi} \sin^{m} x (1 - \cos x)^{n}\, dx$ が収束するか？ [もし $m + 1$ および $m + 2n + 1$ が正であれば。]

\Item{29.} $a > 1$ のとき、

\[
\int_{-1}^{1} \frac{dx}{(a - x) \sqrt{1 - x^2}}
  = \frac{\pi}{\sqrt{a^2 - 1}}
\]

を証明せよ。

\Item{30.} 次の公式を立証せよ：

\begin{alignat*}{2}
\int}{0}^{\infty} F\left(\sqrt{x^2 + 1} + x\right)\, dx
  \&= \tfrac{1}{2}\int\\textit{{1}^{\infty} \&\&\left(1 + \frac{1}{y^2}\right) F(y) \, dy,\\
\%
\int}{0}^{\infty} F\left(\sqrt{x^2 + 1} - x\right)\, dx
  \&= \tfrac{1}{2}\int\\textit{{0}^{1} \&\&\left(1 + \frac{1}{y^2}\right) F(y)\, dy.
\end{alignat*}

特に、$n > 1$ のとき、

\[
\int_{0}^{\infty} \frac{dx}{\left(\sqrt{x^2 + 1} + x\right)^n}
  = \int_{0}^{\infty} \left(\sqrt{x^2 + 1} - x\right)^n\, dx
  = \frac{n}{n^2 - 1}.
\]

[これらの例においては、任意の関数が考慮される積分が第177節以降で定義された意味を持つものであることが当然仮定されている]

31. \(2y = ax - \frac{b}{x}\) において \(a\) と \(b\) が正のとき、\(x\) が \(0\) から \(\infty\) に増加するにつれて \(y\) は \(-\infty\) から \(\infty\) へと安定的に増加することを示せ。これより、
\[
\int_{0}^{\infty} f\left\{\frac{1}{2}\left(ax + \frac{b}{x}\right)\right\} dx = \frac{1}{a} \int_{-\infty}^{\infty} f\{\sqrt{y^{2} + ab}\} \left\{1 + \frac{y}{\sqrt{y^{2} + ab}}\right\} dy
= \frac{2}{a} \int_{0}^{\infty} f\{\sqrt{y^{2} + ab}\} dy.
\]

32. \(2y = ax + \frac{b}{x}\) において \(a\) と \(b\) が正のとき、任意の \(\sqrtp{ab}\) より大きい値の \(y\) に対して二つの \(x\) の値に対応することを示せ。これらのうち大きい方を \(x_{1}\)、小さい方を \(x_{2}\) とおくと、\(y\) が \(\sqrtp{ab}\) から \(\infty\) へと増加するにつれて \(x_{1}\) は \(\sqrtp{\frac{b}{a}}\) から \(\infty\) へと増加し、\(x_{2}\) は \(\sqrtp{\frac{b}{a}}\) から \(0\) へと減少することを示せ。これより
\[
\int_{\sqrtp{\frac{b}{a}}}^{\infty} f(y)\, dx_{1} = \frac{1}{a} \int_{\sqrtp{ab}}^{\infty} f(y) \left\{\frac{y}{\sqrt{y^{2} - ab}} + 1\right\} dy,
\]
\[
\int_{0}^{\sqrtp{\frac{b}{a}}} f(y)\, dx_{2} = \frac{1}{a} \int_{\sqrtp{ab}}^{\infty} f(y) \left\{\frac{y}{\sqrt{y^{2} - ab}} - 1\right\} dy,
\]
および
\[
\int_{0}^{\infty} f\left\{\frac{1}{2}\left(ax + \frac{b}{x}\right)\right\} dx = \frac{2}{a} \int_{\sqrtp{ab}}^{\infty} \frac{yf(y)}{\sqrt{y^{2} - ab}} dy
= \frac{2}{a} \int_{0}^{\infty} f\{\sqrt{z^{2} + ab}\} dz.
\]

33. 
\[
\int_{0}^{\pi} f(\sec\tfrac{1}{2}x + \tan\tfrac{1}{2}x) \frac{dx}{\sqrt{\sin x}} = \int_{0}^{\pi} f(\cosec x) \frac{dx}{\sqrt{\sin x}}
\] を証明せよ。

34. \(a\) と \(b\) が正のとき、次の等式が成り立つ：
\[
\int_{0}^{\infty} \frac{dx}{(x^{2} + a^{2})(x^{2} + b^{2})} = \frac{\pi}{2ab(a + b)}, \quad
\int_{0}^{\infty} \frac{x^{2}\, dx}{(x^{2} + a^{2})(x^{2} + b^{2})} = \frac{\pi}{2(a + b)}.
\]
次に、\(\alpha\)、\(\beta\)、\(\gamma\) が正のとき、\(\beta^{2} \geq \alpha\gamma\) の条件下で次の等式が成り立つ：
\[
\int_{0}^{\infty} \frac{dx}{\alpha x^{4} + 2\beta x^{2} + \gamma} = \frac{\pi}{2\sqrtp{2\gamma A}}, \quad
\int_{0}^{\infty} \frac{x^{2}\, dx}{\alpha x^{4} + 2\beta x^{2} + \gamma} = \frac{\pi}{2\sqrtp{2\alpha A}},
\]
ここで \(A = \beta + \sqrtp{\alpha\gamma}\)。最後の結果は、例題31を用いて \(f(y) = 1/(c^{2} + y^{2})\) の置換により導くことができる。これらの最終的な結果は、\(\beta^{2} < \alpha\gamma\) の条件下でも成り立つが、その証明はそれほど簡単ではない。

35. \(b\) が正のとき、次の等式が成り立つ：
\[
\int_{0}^{\infty} \frac{x^{2}\, dx}{(x^{2} - a^{2})^{2} + b^{2}x^{2}} = \frac{\pi}{2b}, \quad
\int_{0}^{\infty} \frac{x^{4}\, dx}{\{(x^{2} - a^{2})^{2} + b^{2}x^{2}\}^{2}} = \frac{\pi}{4b^{3}}.
\]

36. シュワルツの不等式（第7章、補題7-42）を無限積分に拡張せよ。

37. \(φ(x)\) を178節末尾で考慮した関数とするとき、次の等式が成り立つ：
\[
\int_{0}^{\infty} φ(x)\, dx = \sum_{n=0}^{\infty} \frac{1}{(n + 1)^{2}}.
\]

38. 次の等式が成り立つ：
\begin{alignat*}{2}
\int}{1}^{\infty} dx \left(\int\\textit{{1}^{\infty} \frac{x - y}{(x + y)^{3}}\, dy\right) \&= -1, \& \quad
\int}{1}^{\infty} dy \left(\int\\textit{{1}^{\infty} \frac{x - y}{(x + y)^{3}}\, dx\right) \&= 1;\\
\%
\int}{1}^{\infty} dx \left(\int\\textit{{1}^{\infty} \frac{x^{2} - y^{2}}{(x^{2} + y^{2})^{2}}\, dy\right) \&= -\frac{1}{4}\pi, \& \quad
\int}{1}^{\infty} dy \left(\int\_{1}^{\infty} \frac{x^{2} - y^{2}}{(x^{2} + y^{2})^{2}}\, dx\right) \&= \frac{1}{4}\pi.
\end{alignat*}

これらの積分の限界を0と1に置き換えた類似の結果も確立せよ。
1913年の数学トリップ。

ページ分割線：357

\chapter{第9章：Chapter IX: THE LOGARITHMIC AND EXPONENTIAL FUNCTIONS OF A REAL VARIABLE}

第九章

196．これまで前幾章で扱った本質的に異なる種類の関数は多くありません。その中でも通常の目的では最も重要なものは多項式、有理函数、代數函数（明示的または暗黙的）、三角函数（直接または逆）です。

しかし、数学において重要な関数のリストはまだ尽きていません。数学的知識の範囲が徐々に拡大されるにつれて、一つ又は一つの新しい関数のクラスが分析に導入されてきました。これらの新しい関数は、既知の関数で問題を解決できないことが見受けられたため、しばしば導入されました。この過程は無理数と複素数が最初に導入されたときのものと比較して公正に評価できます。それは、既存の数で解けない代数方程式が見つかったときに初めて無理数と複素数が導入されたようにです。

新しい関数の最も豊富な源は積分問題でした。ある関数$f(x)$を既知の関数を使って積分しようと試みましたが、これらの試みは失敗に終わりました。一定数の失敗後には、問題が解けないと考えられるようになりました。しばしばこれは証明されましたが、通常はその厳密な証明は後にしか得られませんでした。一般に、数学的には既存の不可能性を確信したと一時的に受け入れた後、新しい関数$F(x)$をその定義によって導入し、必要な性質を持つことを定義しました。すなわち、$F'(x) = f(x)$です。この定義から、$F(x)$の性質を調査し始めると、$F(x)$は既知の関数の有限な組み合わせでは絶対に持つことができない性質を持つことが明らかになり、元の問題が解けないと仮定した正しさが確認されました。そのような例は前頁で見られました。第六章（\okrickRef{Ch.}{VI}）において、方程式
\[
\log x = \int \frac{dx}{x}
\]
を使って関数$\log x$を定義しました。

\par\medskip\noindent\textbf{【注釈】}\par
$\log x$が真新しい関数であると推測する根拠は何でしょうか。既に述べた通り（\Ex{xlii}.4）$\log x$は有理関数でないことが分かります。有理関数の導函数は、分母に重複因子のみを含む有理関数となるからです。しかし、代数的または三角関数であるかどうかについては、より難しい問題となります。しかしながら、少量の実験により微分が代数的な無理性を取り除くことはないことが容易に確認できます。例えば、$\sqrt{1 + x}$を何回でも微分すると、その結果は常に$\sqrt{1 + x}$と有理関数の積となります。読者には、多くの例でこの主張の正確性を検証してみてください。同様に、$\sin x$や$\cos x$が含まれる関数を微分すると、これらのいずれかは結果において残ります。

したがって、$\log x$が新しい関数であるという厳密な証明は得られません（そのような証明については、236ページに引用された著者の論文をご覧ください）。しかし、合理的な推測としては、それが新しい関数であることが分かります。したがって、それを新しい関数として扱います。そして、検討するとその性質はこれまで遭遇したどの関数とも異なることが分かります。
\par\medskip

\Paragraph{197. $\log x$の定義}
$\log x$、すなわち$x$の対数を以下の等式により定義します：
\[
\log x = \int_{1}^{x} \frac{dt}{t}.
\]
$x$は正であることが必要で、（\Ex{lxxvi}.2）積分の範囲に$x = 0$が含まれる場合、この積分には意味がないからです。下限を1以外に選んでもよいですが、1が最も便利です。この定義により$\log 1 = 0$となります。

We shall now consider how $\log x$ behaves as $x$ varies from $0$ towards $\infty$. It follows at once from the definition that $\log x$ is a continuous function of $x$ which increases steadily with $x$ and has a derivative
\[
D_{x} \log x = 1/x;
\]
and it follows from §175 that $\log x$ tends to $\infty$ as $x \to \infty$.

If $x$ is positive but less than $1$, then $\log x$ is negative. For
\[
\log x = \int_{1}^{x} \frac{dt}{t} = -\int_{x}^{1} \frac{dt}{t} < 0.
\]
Moreover, if we make the substitution $t = 1/u$ in the integral, we obtain
\[
\log x = \int_{1}^{x} \frac{dt}{t}
       = -\int_{1}^{1/x} \frac{du}{u}
       = -\log(1/x).
\]
Thus $\log x$ tends steadily to $-\infty$ as $x$ decreases from $1$ to $0$.

The general form of the graph of the logarithmic function is shown in Fig. 52. Since the derivative of $\log x$ is $1/x$, the slope of the curve is very gentle when $x$ is very large, and very steep when $x$ is very small.

\par\medskip\noindent\textbf{【例 LXXXII.】}\par
\Item{1.} Prove from the definition that if $u > 0$ then
\[
\frac{u}{1 + u} < \log(1 + u) < u.
\]

[For $\log(1 + u) = \int_{0}^{u} \frac{dt}{1 + t}$, and the subject of integration lies between $1$ and $1/(1 + u)$.]

\Item{2.} Prove that $\log(1 + u)$ lies between $u - \frac{u^2}{2}$ and $u - \frac{u^2}{2(1 + u)}$ when $u$ is positive. [Use the fact that $\log(1 + u) = u - \int_{0}^{u} \frac{t\, dt}{1 + t}$.]

\Item{3.} If $0 < u < 1$ then $u < -\log(1 - u) < \frac{u}{1 - u}$.

\Item{4.} Prove that
\[
\lim_{x \to 1} \frac{\log x}{x - 1} = \lim_{t \to 0} \frac{\log (1 + t)}{t} = 1.
\]

[Use Ex.~1.]
\par\medskip
\PageSep{360}

\Paragraph{198. $\log x$ が満足する関数方程式} 関数 $\log x$ は次の関数方程式を満たす：
\[ f(xy) = f(x) + f(y). \tag{1} \]
証明は、変数 $t = yu$ の置換により得られる。したがって
\begin{align*}
\log xy &= \int_{1}^{xy} \frac{dt}{t} \\
        &= \int_{1/y}^{x} \frac{du}{u} \\
        &= \int_{1}^{x} \frac{du}{u} - \int_{1}^{1/y} \frac{du}{u} \\
        &= \log x - \log(1/y) = \log x + \log y,
\end{align*}
これにより定理が証明される。

\par\medskip\noindent\textbf{【例 LXXXIII.】}\par
\Item{1.} 方程式 (1) を満たし、$\log x$ から根本的に異なる解は存在しないことを示すことができる。方程式を $x$ と $y$ についてそれぞれ微分すると、次の2つの方程式を得る：
\[ yf'(xy) = f'(x), \quad xf'(xy) = f'(y); \]
これらから $f'(xy)$ を消去すると、$xf'(x) = yf'(y)$ となる。この関係が任意の $x$ と $y$ の値に対して成り立つならば、$xf'(x) = C$ であり、ここで $C$ は定数である。したがって
\[ f(x) = \int \frac{C}{x} \, dx + C' = C\log x + C', \]
そして $C' = 0$ であることから明らかである。従って、$\log x$ から根本的に異なる解は存在せず、$f(x) = 0$（$C = 0$ のとき）が唯一の平凡な解である。

\Item{2.} 同様にして、方程式
\[ f(x) + f(y) = f\left(\frac{x + y}{1 - xy}\right) \]
を満たし、$\arctan x$ から根本的に異なる解は存在しないことを示す。
\par\medskip

\Paragraph{199. $\log x$ の $x$ とともに無限大に収束する方法} 我々は xxxvi\.6 で、ある関数が $x$ とともに無限大に収束するときのいくつかの異なる方法を定義したことを思い出す。特に、$x$ が大きくなると、その関数が最初の、第二の、第三の \dots 種類の greatness のいずれかであることを区別した。関数 $f(x)$ は $k$ 番目の greatness のとき、$f(x)/x^k$ が $x$ が無限大に向かうにつれて非ゼロの極限に収束することを言う。

以下は関数の級数を定義する方法で、$x$ とともに無限大に収束するが、その大きさの順序が最初のものよりも小さいものがある。例えば $\sqrt{x}$, $\sqrt{x}$, $\sqrt{x}$, \dots はそのような関数である。一般に、$\alpha$ が任意の正有理数であれば、$x$ が大きいとき $x^{\alpha}$ は $\alpha$ 番目の順序を持つ大きさ関数となる。そして $\alpha$ をどれほど小さくしてもよい。ページ分割点で361ページに移行し、$\alpha$ の値を任意に小さくできる。そして $\alpha$ のすべての値を与えることで、可能であるすべての「無限大の順序」を尽きるとは考えられなかったかもしれない。少なくとも $x$ とともに無限大に収束する関数 $f(x)$ について、どれほどゆっくりであっても、$\alpha$ の適切な値を見つけることができると考えられたし、逆に、どれほど速くであっても、$\alpha$ の適切な値を見つけることができると考えられた。

関数 $\log x$ の最も興味深い特徴は、$x$ が無限大に向かうときの挙動にある。これは上記で述べた前提が自然であるように思われても、誤りであることを示している。すなわち、$\log x$ は $x$ とともに無限大に収束するが、任意の正の冪（整数または分数）よりも遅い速度で無限大に向かう。言い換えれば $\log x \to \infty$ だが
\[
\frac{\log x}{x^{\alpha}} \to 0
\]
がすべての正の $\alpha$ の値に対して成り立つ。この事実は、しばしば「$\log x$ の無限大の順序は無限に小さい」というように非公式に表現されることがあるが、読者はこの段階でそのような表現を避ける必要はないだろう。

\paragraph{200. $(\log x)/x^{\alpha}$ が $x \to \infty$ のとき 0 に収束することの証明。}
任意の正数 $\beta$ を取る。このとき $t > 1$ のとき $1/t < 1/t^{1-\beta}$ であるから
\[
\log x = \int_{1}^{x} \frac{dt}{t} < \int_{1}^{x} \frac{dt}{t^{1-\beta}},
\]
または
\[
\log x < \frac{x^{\beta} - 1}{\beta} < \frac{x^{\beta}}{\beta},
\]
$x > 1$ のときが成り立つ。今、任意の正数 $\alpha$ を取る。このときより小さな正数 $\beta$ を選べる。そして
\[
0 < \frac{\log x}{x^{\alpha}} < \frac{x^{\beta-\alpha}}{\beta} \quad (x > 1)
\]
が成り立つ。しかし、$\alpha > \beta$ より $x^{\beta-\alpha}/\beta \to 0$ であるから
\[
\frac{\log x}{x^{\alpha}} \to 0.
\]

\paragraph{201. $x \to +0$ のとき $\log x$ の挙動。}
$x = 1/y$ のとき $(\log x)/x^{\alpha} = -y^{\alpha} \log y$ であるから、上記の定理より
\[
\lim_{y \to +0} y^{\alpha} \log y = -\lim_{x \to +\infty} \frac{\log x}{x^{\alpha}} = 0.
\]
したがって $\log x$ は $-\infty$ に収束し、$\log(1/x) = -\log x$ は正の値で $x$ が $0$ に近づくとき $\infty$ に収束する。しかし、$\log(1/x)$ は任意の正数の $1/x$ のべき乗（整数または分数）よりも遅い速度で $\infty$ に収束する。

以下に翻訳した文章を出力します：

\par\medskip\noindent\textbf{【注釈】}\par
\Paragraph{202. 意味の尺度。対数尺度。} 再び次のような関数列を考えましょう：
\[
x,\quad \sqrt{x},\quad \sqrt{x},\ \dots,\quad \sqrt[n]{x},\ \dots,
\]
この列は、任意の二つの関数 $f(x)$ と $\phi(x)$ を含む場合、$x \to \infty$ のとき、$f(x)$ および $\phi(x)$ ともに無限大に収束し、また $f(x)/\phi(x)$ は $f(x)$ が $\phi(x)$ より右にある場合は $0$ に、左にある場合は無限大に収束することを有する性質を持っています。この列のすべての既存項の右側に新たな項を挿入することで、この列を延長することができます。まず $\log x$ を挿入します。これは任意の古い項よりも遅い速度で無限大に収束します。次に $\sqrt{\log x}$ は $\log x$ より遅い速度で無限大に収束し、$\sqrt{\log x}$ は $\sqrt{\log x}$ より遅い速度で無限大に収束します。このようにして、次のような列を得ます：
\[
x,\quad \sqrt{x},\quad \sqrt{x},\ \dots,\quad \sqrt[n]{x},\ \dots\quad
\log x,\quad \sqrt{\log x},\quad
\sqrt{\log x},\ \dots\quad
\sqrt[n]{\log x},\ \dots
\]
これは二つの単純無限列が順に並べられたものです。しかしこれだけではありません。関数 $\log\log x$ を考慮します。$\log x/x^{\alpha} \to 0$ であるすべての正の $\alpha$ の値に対して、$x = \log y$ と置き換えると
\[
(\log\log y)/(\log y)^{\alpha} = (\log x)/x^{\alpha} \to 0.
\]
したがって $\log\log y$ は $y$ とともに無限大に収束しますが、任意の $\log y$ の冪よりも遅い速度で収束します。従って、次のような列を延長することができます：
\begin{gather*}
x,\quad \sqrt{x},\quad \sqrt{x},\ \dots\qquad
\log x,\quad \sqrt{\log x},\quad \sqrt{\log x},\ \dots\\
\log\log x,\quad \sqrt{\log\log x},\ \dots\quad \sqrt[n]{\log\log x},\ \dots;
\end{gather*}
さらに、$\log\log\log x$，$\log\log\log\log x$，\dots を導入することで、この列を任意の長さまで延長することができます。$x = 1/y$ と置き換えることで、$y \to 0$ のとき無限大に収束する $y$ の関数に対する類似の意味の尺度を得ることができます。\footnote{「規模の秩序」に関するより詳しい情報については、著者の小論文「無限の順序」（Camb.\ Math.\ Tracts, No.~12）を参照してください。}
\par\medskip

\par\medskip\noindent\textbf{【例 LXXXIV.】}\par
\Item{1.} 任意の項 $f(x),~F(x)$ の級数において、新しい項 $\phi(x)$ を挿入することができる。$\phi(x)$ は $f(x)$ よりも速く無限大に収束し、$F(x)$ よりも遅く無限大に収束する。[例えば、$\sqrt{x}$ と $\sqrt{x}$ の間に $x^{5/12}$ を挿入できる：$\sqrtp{\log x}$ と $\sqrtp{\log x}$ の間に $(\log x)^{5/12}$ を挿入できる。そして一般に、$\phi(x) = \sqrtb{f(x) F(x)}$ は述べた条件を満たす.]

\Item{2.} $\sqrt{x}$ よりも速く無限大に収束し、$x^{\alpha}$ よりも遅く無限大に収束する関数を求めよ。ここで $\alpha$ は $1/2$ より小さい任意の有理数である。[$\sqrt{x}/(\log x)$ はそのような関数であり；または $\beta$ を任意の正の有理数として $\sqrt{x}/(\log x)^{\beta}$ も同様である。]

\Item{3.} $\sqrt{x}$ よりも速く無限大に収束し、$\sqrt{x}/(\log x)^{\alpha}$ よりも遅く無限大に収束する関数を求めよ。ここで $\alpha$ は任意の有理数である。[関数 $\sqrt{x}/(\log\log x)$ がそのような関数である。これらの例から、無限の大尺度において \emph{不完全性} は固有の特性であることが推測される。]
\par\medskip

4. 関数
\[
f(x) = \frac{x^{\alpha} (\log x)^{\alpha'} (\log\log x)^{\alpha''}}{x^{\beta}  (\log x)^{\beta'}  (\log\log x)^{\beta''}}
\]
が $x$ が無限大に近づくときどう振る舞うか？ $\alpha \neq \beta$ の場合、
\[
f(x) = x^{\alpha-\beta} (\log x)^{\alpha'-\beta'} (\log\log x)^{\alpha''-\beta''}
\]
の振舞いは $x^{\alpha-\beta}$ によって支配される。$\alpha = \beta$ の場合、$x$ のべき乗が消え、
$(\log x)^{\alpha'-\beta'}$ によって支配される（$\alpha' = \beta'$ の場合を除き）、そのとき $(\log\log x)^{\alpha''-\beta''}$ によって支配される。したがって $\alpha > \beta$ のとき $f(x) \to \infty$、または $\alpha = \beta$, $\alpha' > \beta'$ のとき、または $\alpha = \beta$, $\alpha' = \beta'$, $\alpha'' > \beta''$ のとき $f(x) \to \infty$ であり、$\alpha < \beta$ のとき、または $\alpha = \beta$, $\alpha' < \beta'$ のとき、または $\alpha = \beta$, $\alpha' = \beta'$, $\alpha'' < \beta''$ のとき $f(x) \to 0$。

5. 関数 $x/\sqrt{\log x}$, $x\sqrt{\log x}/\log\log x$, $x\log\log x/\sqrt{\log x}$, $(x\log\log\log x)/\sqrt{\log\log x}$ を、$x \to \infty$ で無限大に収束する速度に従って並べよ。

6. 関数
\[
\frac{\log\log x}{x\log x},\quad
\frac{\log x}{x},\quad
\frac{x\log\log x}{\sqrt{x^{2} + 1}},\quad
\frac{\sqrt{x+1}}{x(\log x)^{2}}
\]
を、$x \to \infty$ で0に収束する速度に従って並べよ。

7. 関数
\[
x\log\log(1/x),\quad
\sqrt{x}/\{\log(1/x)\},\quad
\sqrt{x\sin x\log(1/x)},\quad
(1 - \cos x)\log(1/x)
\]
を、$x \to +0$ で0に収束する速度に従って並べよ。

8. 次の関係を示せ：
\[
D_{x}\log\log x = \frac{1}{x\log x},\quad
D_{x}\log\log\log x = \frac{1}{x\log x\log\log x},
\]
と続く。

9. 次の関係を示せ：
\[
D_{x}(\log x)^{\alpha} = \frac{\alpha}{x(\log x)^{1-\alpha}},\quad
D_{x}(\log\log x)^{\alpha} = \frac{\alpha}{x\log x(\log\log x)^{1-\alpha}},
\]
と続く。

第203節. 数字 \( e \)．今後、通常 \( e \) を表す数字を導入します。この数字は、より高度な数学において極めて重要なものです。また、\(\pi\) とともに解析学の基本的な定数の一つです。

\( e \) を \emph{対数が1である数字}として定義します。つまり
\[
1 = \int_{1}^{e} \frac{dt}{t}
\]
という方程式によって \( e \) は定義されます。対数関数 \(\log x\) は、第95節で厳密に述べた意味での \( x \) の増加函数であり、値1をただ一度通るだけです。したがって、我々の定義は実質的に一つの特定の数字を定義します。

また
\[
\log xy = \log x + \log y
\]
から
\[
\log x^{2} = 2\log x, \quad \log x^{3} = 3\log x, \ldots, \quad \log x^{n} = n\log x,
\]
ここで \( n \) は任意の正整数です。したがって
\[
\log e^{n} = n\log e = n.
\]

また、任意の正整数 \( p \) と \( q \) に対して、\( e^{p/q} \) を \( e^{p} \) の正の \( q \) 次根とするとき
\[
p = \log e^{p} = \log(e^{p/q})^{q} = q\log e^{p/q},
\]
から \(\log e^{p/q} = p/q\) となります。したがって、任意の正有理数 \( y \) のとき、\( e^{y} \) を \( e \) の正の \( y \) 次冪とするとき
\[
\log e^{y} = y,
\]
式（1）が成り立ちます。また \(\log e^{-y} = -\log e^{y} = -y\) から、式（1）はすべての有理数 \( y \) の正または負の値に対して真です。つまり
\[
y = \log x, \quad x = e^{y}
\]
が互いに導く方程式（2）は、\( y \) が有理数で \( e^{y} \) がその正の値を持つ限り成立します。

現在、指数が無理性であるような冪 \( e^{y} \) の定義を提供していませんので、式 \( e^{y} \) は \( y \) が有理数のときのみ定義されます。

例．\( 2 < e < 3 \) を証明せよ。まず
\[
\int_{1}^{2} \frac{dt}{t} < 1,
\]
から \( 2 < e \) であることが明らかです。また
\[
\int_{1}^{3} \frac{dt}{t} = \int_{1}^{2} \frac{dt}{t} + \int_{2}^{3} \frac{dt}{t}
= \int_{0}^{1} \frac{du}{2 - u} + \int_{0}^{1} \frac{du}{2 + u}
= 4\int_{0}^{1} \frac{du}{4 - u^{2}} > 1,
\]
から \( e < 3 \) であることが分かります。

\Paragraph{204. 指数関数。} 現在，実数 $y$ に対する \emph{指数関数} $e^{y}$ を対数関数の逆関数として定義します。つまり，
\[ x = e^{y}
\]
を満たすとき、$y = \log x$ です。

私たちは，$x$ が $0$ から $\infty$ に向かって変化するにつれて，$y$ はより厳密な意味で $-\infty$ から $\infty$ に安定的に増加することを見ました。したがって，$x$ の一つの値に対応する $y$ の一つの値があり，逆もまた然りです。また，$y$ は $x$ の連続関数であり，\SecNo[§]{109}から $x$ も同様に $y$ の連続関数であることがわかります。

\par\medskip\noindent\textbf{【注釈】}\par
指数関数の連続性を直接証明することも容易です。もし $x = e^{y}$ で，$x + \xi = e^{y+\eta}$ なら，
\[ \eta = \int_{x}^{x+\xi} \frac{dt}{t}.
\]
したがって，$\xi > 0$ のとき $|\eta|$ は $\xi/(x + \xi)$ より大きいし，$\xi < 0$ のとき $|\eta|$ は $|\xi|/x$ より大きいです。そして，$\eta$ が非常に小さいなら $\xi$ も非常に小さいことになります。
\par\medskip
\PageSep{365}

したがって，$e^{y}$ は $y$ の正の連続関数であり，$y$ が $-\infty$ から $\infty$ に向かうにつれて安定的に $0$ から $\infty$ に増加します。さらに，$e^{y}$ は数 $e$ の正の $y$ 乗であり，基本的な定義に従います。特に，$y = 0$ のとき $e^{y} = 1$ です。指数関数 $e^{y}$ のグラフの一般形は \Fig{53} に示されています。
\%[図版: Fig. 53.]
\Figure[3インチ]{53}{p365}

\Paragraph{205. 指数関数の主要な性質。}
\Item{(1)} $x = e^{y}$ であるとき，$y = \log x$ であり，$\frac{dy}{dx} = \frac{1}{x}$ および
\[ \frac{dx}{dy} = x = e^{y}.
\]
したがって，\emph{指数関数の導関数はその自身と等しい}です。より一般に，$x = e^{ay}$ のとき $\frac{dx}{dy} = ae^{ay}$ です。

\Item{(2)} \emph{指数関数は以下の関数方程式を満たす}
\[ f(y + z) = f(y)f(z).
\]

Thisは、$y$と$z$が有理数のとき指数法則から導かれる。もし$y$または$z$、あるいは両方が無理性であるなら、有理数の列$y_{1}$, $y_{2}$, \dots, $y_{n}$, \dots\ と$z_{1}$, $z_{2}$, \dots, $z_{n}$, \dots\ を選ぶことができる使得$\lim y_{n} = y$, $\lim z_{n} = z$。指数関数が連続であるから、
\[
e^{y} × e^{z} = \lim e^{y_{n}} × \lim e^{z_{n}} = \lim e^{y_{n}+z_{n}} = e^{y+z}.
\]
特に$e^{y} × e^{-y} = e^{0} = 1$、または$e^{-y} = 1/e^{y}$。

また$\log x$が満たす関数方程式から$e^{y}$が満たす関数方程式を導くこともできる。もし$y_{1} = \log x_{1}$, $y_{2} = \log x_{2}$とすれば、$x_{1} = e^{y_{1}}$, $x_{2} = e^{y_{2}}$であるから、$y_{1} + y_{2} = \log x_{1} + \log x_{2} = \log x_{1}x_{2}$であり
\[
e^{y_{1}+y_{2}} = e^{\log x_{1}x_{2}} = x_{1}x_{2} = e^{y_{1}} × e^{y_{2}}。
\]
\\PageSep{366}

\par\medskip\noindent\textbf{【例 LXXXV.】}\par
\Item{1.} $dx/dy = ax$のとき、$x = Ke^{ay}$、$K$は定数。

\Item{2.} 方程式$f(y + z) = f(y)f(z)$に根本的に異なる解は指数関数から導けない。[もし$f(y)$が微分可能であると仮定する。方程式を$y$と$z$についてそれぞれ微分すると
\[
f'(y + z) = f'(y)f(z),\quad
f'(y + z) = f(y)f'(z)
\]
となり、これより$f'(y)/f(y) = f'(z)/f(z)$であり、従って各々は定数である。したがって$x = f(y)$のとき、$dx/dy = ax$（例題1）となる。]

\Item{3.} $(e^{ay} - 1)/y \to a$として$y \to 0$を証明せよ。[平均値の定理により、$e^{ay} - 1 = aye^{a\eta}$を得るここで$0 < |\eta| < |y|$である。]
\par\medskip

\Paragraph{206.} \par\medskip\noindent\textbf{【結果】}\par\Item{(3)} 関数$e^{y}$は$y$と共に無限大に発散し、任意の$y$の冪よりも速く発散する。すなわち
\[
\lim y^{\alpha}/e^{y} = \lim e^{-y}y^{\alpha} = 0
\]
が$y \to \infty$で全ての$\alpha$に対して成り立つ。
\par\medskip

We saw that $\frac{\log x}{x^{\beta}} \to 0$ as $x \to \infty$, for any positive value of~$\beta$ however small. Writing $\alpha$ for~$\frac{1}{\beta}$, we see that $\frac{(\log x)^{\alpha}}{x} \to 0$ for any value of~$\alpha$ however large. The result follows on putting $x = e^{y}$. It is also clear that $e^{\gamma y}$ tends to~$\infty$ if $\gamma > 0$, and to~$0$ if $\gamma < 0$, and in each case more rapidly than any power of~$y$.

\par\medskip\noindent\textbf{【注釈】}\par
From this result, we can construct a `scale of infinity' similar to that constructed in \SecNo[§]{202}, but extending in the opposite direction; i.e., a scale of functions which tend to~$\infty$ more and more rapidly as $x \to \infty$. The exponential function was introduced by inverting the equation $y = \log x$ into $x = e^{y}$; and we have accordingly, up to the present, used $y$ as the independent and $x$ as the dependent variable in discussing its properties. We shall now revert to the more natural plan of taking~$x$ as the independent variable, except when it is necessary to consider a pair of equations of the type $y = \log x$, $x = e^{y}$ simultaneously, or when there is some other special reason to the contrary. The scale is
\[
x,\quad x^{2},\quad x^{3},\ \dots\quad
e^{x},\quad e^{2x},\ \dots\quad
e^{x^{2}},\ \dots,\quad e^{x^{3}},\ \dots,\quad e^{e^{x}},\ \dots,
\]
where of course $e^{x^{2}}$,~\dots, $e^{e^{x}}$,~\dots denote $e^{(x^{2})}$,~\dots, $e^{(e^{x})}$,~\dots.

The reader should try to apply the remarks about the logarithmic scale, made in \SecNo[§]{202} and \Exs{lxxxiv}, to this `exponential scale' also. The two scales may of course (if the order of one is reversed) be combined into one scale
\[
\dots\ \log\log x,\ \dots\quad \log x,\ \dots\quad
x,\ \dots\quad e^{x},\ \dots\quad e^{e^{x}},\ \dots.
\]
\par\medskip

207. 一般のべき乗 \(a^x\)。関数 \(a^x\) は、\(x\) が有理数である場合にのみ定義されており、特に \(a = e\) の場合を除いてです。今後、\(a\) が任意の正の数である場合を考えます。もし \(x\) が正の有理数 \(p/q\) であるとするならば、べき乗 \(a^{p/q}\) の正の値 \(y\) は \(y^q = a^p\) で与えられ、これより
\[
q \log y = p \log a, \quad \log y = \left(\frac{p}{q}\right) \log a = x \log a,
\]
よって
\[
y = e^{x \log a}.
\]

私たちは、\(x\) が無理数であるとき、この定義を \(a^x\) の定義とします。したがって \(10^{\sqrt{2}} = e^{\sqrt{2} \log 10}\) です。注意すべきは、\(x\) が無理数であるとき、\(a^x\) は正の \(a\) のみで定義され、またその値も必ず正であること、および \(\log a^x = x \log a\) であることです。関数 \(a^x\) の最も重要な性質は以下の通りです。

(1) どんな値を取っても \(a^x \times a^y = a^{x+y}\) で、\((a^x)^y = a^{xy}\) です。つまり、指数法則は有理数の指数と同様に無理数の指数に対して成り立ちます。まず第一に
\[
a^x \times a^y = e^{x \log a} \times e^{y \log a} = e^{(x+y) \log a} = a^{x+y};
\]
そして第二に
\[
(a^x)^y = e^{y \log a^x} = e^{xy \log a} = a^{xy}.
\]

(2) \(a > 1\) のとき、\(a^x = e^{x \log a} = e^{\alpha x}\)，ここで \(\alpha\) は正の数です。この場合、関数 \(a^x\) のグラフは \(e^x\) のグラフと相似であり、また \(a^x \to \infty\) で \(x \to \infty\) となり、任意の \(x\) のべきより速く増大します。

\(a < 1\) のとき、\(a^x = e^{x \log a} = e^{-\beta x}\)，ここで \(\beta\) は正の数です。この場合、関数 \(a^x\) のグラフは \(e^x\) の形状と相似ですが左右逆転しており、また \(a^x \to 0\) で \(x \to \infty\) となり、任意の \(1/x\) のべきより速く減少します。

(3) \(a^x\) は \(x\) の連続関数であり
\[
D_x a^x = D_x e^{x \log a} = e^{x \log a} \log a = a^x \log a.
\]

(4) \(a^x\) はまた \(a\) の連続関数であり
\[
D_a a^x = D_a e^{x \log a} = e^{x \log a} (x/a) = x a^{x-1}.
\]

\Item{(5)} $x \to 0$ のとき $(a^{x} - 1)/x \to \log a$ である。これは明らかに $D_{x}a^{x} = a^{x}\log a$ という事実から派生するものだが、その結果の特定の形はしばしば有用であり、当然ながら $\alpha x \to 0$ のとき $(e^{\alpha x} - 1)/x \to \alpha$ であるという結果（\Ex{lxxxv}.3）と等価である。

\par\medskip\noindent\textbf{【注釈】}\par
前章で多くの結果が $x$ が有理数であるという制限付きで述べられていた。本節で与えられた定義と定理により、この制限を除去することができる。
\par\medskip
\PageSep{368}

\Paragraph{208. $e^{x}$ の極限表現} 第IV章 \SecNo[§]{73} では $\{1 + (1/n)\}^{n}$ が $n \to \infty$ のときある極限に収束することを証明し、その極限を仮定的に $e$ と表していた。現在はこの極限を前節で述べた数 $e$ と同一視する。しかしより一般的な結果を確立できる。すなわち
\[
\lim_{n\to\infty} \left(1 + \frac{x}{n}\right)^{n}
  = \lim_{n\to\infty} \left(1 - \frac{x}{n}\right)^{-n}
  = e^{x}.
\Tag{(1)}
\]
この結果は非常に重要であるため、異なる証明法を示す。

\Item{(1)} 
\[
\frac{d}{dt} \log(1 + xt) = \frac{x}{1 + xt}
\] 
から
\[
\lim_{h\to 0} \frac{\log(1 + xh)}{h} = x
\]
が導かれる。$h = 1/\xi$ を置くと
\[
\lim \xi \log\left(1 + \frac{x}{\xi}\right) = x
\]
が $\xi \to \infty$ や $\xi \to -\infty$ のとき成立する。指数関数が連続であるから
\[
\left(1 + \frac{x}{\xi}\right)^{\xi} = e^{\xi\log\{1+(x/\xi)\}} \to e^{x}
\]
が $\xi \to \infty$ や $\xi \to -\infty$ のとき成立し、すなわち
\[
\lim_{\xi\to\infty} \left(1 + \frac{x}{\xi}\right)^{\xi}
  = \lim_{\xi\to -\infty} \left(1 + \frac{x}{\xi}\right)^{\xi}
  = e^{x}.
\Tag{(2)}
\]

$\xi \to \infty$ や $\xi \to -\infty$ のとき整数値を取る限りでは、結果は式 (1) で表される。

\par\medskip\noindent\textbf{【注釈】}\par
(2) 任意の正整数 \( n \)，いかなる大きさであっても，および \( x > 1 \) のとき，
\[
\int_{1}^{x} \frac{dt}{t^{1+(1/n)}}
  < \int_{1}^{x} \frac{dt}{t}
  < \int_{1}^{x} \frac{dt}{t^{1-(1/n)}},
\]
または
\[
n(1 - x^{-1/n}) < \log x < n(x^{1/n} - 1).
\Tag{(3)}
\]
ここで \( y = \log x \) とおくと，\( y \) は正であり，\( x = e^y \) である．その後，簡単な変形により，
\[
\left(1 + \frac{y}{n}\right)^{n} < x < \left(1 - \frac{y}{n}\right)^{-n}.
\Tag{(4)}
\]
次に，
\[
1 + \frac{y}{n} = \eta_{1},\quad
1 - \frac{y}{n} = \frac{1}{\eta_{2}}.
\]
ここで \( 0 < \eta_{1} < \eta_{2} \) であり，特に十分大きな \( n \) の値に対しては；さらに，
\[
\eta_{2}^{n} - \eta_{1}^{n}
  < n\eta_{2}^{n-1} (\eta_{2} - \eta_{1})
  = y^{2}\eta_{2}^{n}/n,
\]
これは \( n \to \infty \) のとき明らかに0に収束する．以上の不等式（4）から結果が導かれる．より一般的な結果（2）も同様の方法で証明できる．ただし，\( 1/n \) を連続変数 \( h \) に置き換える．

\Paragraph{209. \(\log x\) の極限表現} 私たちはまた（75節参照），
\[
\lim n(1 - x^{-1/n}) = \lim n(x^{1/n} - 1) = \log x.
\]

これは，
\[
n(x^{1/n} - 1) - n(1 - x^{-1/n}) = n(x^{1/n} - 1)(1 - x^{-1/n}),
\]
が \( n \to \infty \) のとき0に収束することから導かれる．ここで，\( n(x^{1/n} - 1) \) は極限に収束し（75節参照），かつ \( x^{-1/n} \) は1に収束する（27例題10）．以上の不等式（3）（208節参照）から結果が導かれる．

\par\medskip

\par\medskip\noindent\textbf{【例 LXXXVI.】}\par
1. 不等式（4）（208節参照）において \( y = 1 \) および \( n = 6 \) を取ることにより，\( 2.5 < e < 2.9 \) を証明せよ．

2. \( t > 1 \) のとき，
\[
\frac{t^{1/n} - t^{-1/n}}{t - t^{-1}} < \frac{1}{n},
\]
および \( x > 1 \) のとき，
\[
\int_{1}^{x} \frac{dt}{t^{1-(1/n)}} - \int_{1}^{x} \frac{dt}{t^{1+(1/n)}}
  < \frac{1}{n} \int_{1}^{x} \left(t - \frac{1}{t}\right) \frac{dt}{t}
  = \frac{1}{n} \left(x + \frac{1}{x} - 2\right).
\]
これより，209節の結果を導け．

\item{3.} $\xi_{n}$が$n$の関数であり、かつ$n\xi_{n} \to l$となるとき、$(1 + \xi_{n})^{n} \to e^{l}$である。

この式を$n\log(1 + \xi_{n})$の形に書き直すと、

\[
l \left(\frac{n\xi_{n}}{l}\right) \frac{\log(1 + \xi_{n})}{\xi_{n}}
\]

となり、\Ex{lxxxii}.4を用いると$n\log(1 + \xi_{n}) \to l$であることがわかる。

\item{4.} $n\xi_{n} \to \infty$のとき、$(1 + \xi_{n})^{n} \to \infty$であり、また$1 + \xi_{n} > 0$かつ$n\xi_{n} \to -\infty$のとき、

\[
(1 + \xi_{n})^{n} \to 0
\]

である。

\item{5.} \SecNo[§]{208}の\Eq{(1)}から、$e^{y}$は任意の$y$の冪よりも速く無限大に収束することを導き出すことができる。

\Paragraph{210. 共通対数。} 読者には対数の概念とその計算における応用が既に知られているべきである。初等代数では、底$a$に対する$x$の対数$\log_{a} x$は方程式

\[
x = a^{y}, \quad
y = \log_{a} x
\]

によって定義される。

この定義は当然、$y$が有理数である場合にのみ適用可能であるが、しばしば無言でこの点を省略される。

ページ番号分離：370

したがって我々の対数は底$e$に対する対数である。実際的な計算では底$10$の対数を使用する。もし、

\[
y = \log x = \log_{e} x, \quad
z = \log_{10} x,
\]

ならば$x = e^{y}$であり、また$x = 10^{z} = e^{z\log 10}$であるから、

\[
\log_{10} x = (\log_{e} x)/(\log_{e} 10)
\]

となる。したがって一度$\log_{e} 10$を計算すれば、両対数系間の変換は容易である。

この本において我々の目的は対数の実用的な応用に関する詳細な説明ではない。読者がそれについて不慣れであれば、初等代数や三角法の教科書を参照すべきである。\footnote
  {例えばChrystalの『代数学』第1巻、第xxi章を参照されたい。$\log_{e} 10$の値は$2.302\dots$であり、その逆数は$.434\dots$である。}

\par\medskip\noindent\textbf{【例】}\par{LXXXVII.}
\Item{1.} \(D_{x} e^{ax}\cos bx = re^{ax} \cos(bx + \theta),\quad D_{x} e^{ax}\sin bx = re^{ax} \sin(bx + \theta)\) を示せ。ここで \(r = \sqrt{a^{2} + b^{2}}\), \(\cos\theta = a/r\), \(\sin\theta = b/r\) である。さらに、これらの関数の \(n\) 番目の導関数を求めよ。特に、\(D_{x}^{n} e^{ax} = a^{n} e^{ax}\) を示せ。

\Item{2.} 曲線 \(y = e^{-ax}\sin bx\) の形状を追跡せよ。ここで \(a\) と \(b\) は正である。\(y\) は無限の極大値を持つことが示され、その値は等比数列を形成し、これらの極大値は曲線 \(y = \frac{b}{\sqrt{a^{2} + b^{2}}} e^{-ax}\) 上に存在することが証明される。

\Item{3.} 指数関数を含む積分。次の等式を証明せよ。
\[
\int e^{ax}\cos bx\, dx = \frac{a\cos bx + b\sin bx}{a^{2} + b^{2}}\, e^{ax},\quad
\int e^{ax}\sin bx\, dx = \frac{a\sin bx - b\cos bx}{a^{2} + b^{2}}\, e^{ax}.
\]
ここで \(I\) と \(J\) をそれぞれの積分とし、部分積分により次の式を得る。
\[
aI = e^{ax}\cos bx + bJ,\quad aJ = e^{ax}\sin bx - bI.
\]
これらの方程式を解いて \(I\) と \(J\) を求めよ。

\Item{4.} 例題2の曲線と \(x\) の正半軸で囲まれる連続した面積は等比数列を形成し、その和は
\[
\frac{b}{a^{2} + b^{2}} \cdot \frac{1 + e^{-a\pi/b}}{1 - e^{-a\pi/b}}
\]
であることを証明せよ。

\Item{5.} \(a > 0\) のとき、次の積分を証明せよ。
\[
\int_{0}^{\infty} e^{-ax}\cos bx\, dx = \frac{a}{a^{2} + b^{2}},\quad
\int_{0}^{\infty} e^{-ax}\sin bx\, dx = \frac{b}{a^{2} + b^{2}}.
\]

\Item{6.} \(I_{n} = \int e^{ax}x^{n}\, dx\) とすると、\(aI_{n} = e^{ax}x^{n} - nI_{n-1}\) である。部分積分により、すべての正整数 \(n\) の値に対して \(I_{n}\) を計算することが可能である。
\(\PageSep{371}\)

\item7. 正の整数 \( n \) に対して、次の等式が成り立つことを証明せよ：
\[
\int_{0}^{\xi} e^{-x} x^{n} \, dx = n! \, e^{-\xi} \left(
e^{\xi} - 1 - \xi - \frac{\xi^{2}}{2!} - \cdots - \frac{\xi^{n}}{n!}
\right)
\]
および
\[
\int_{0}^{\infty} e^{-x} x^{n} \, dx = n!.
\]

\item8. 任意の \( e^{x} \) の有理関数の積分法を示せ。[\( x = \log u \) と置くと、\( e^{x} = u \), \( dx/du = 1/u \) となり、積分は \( u \) の有理関数の積分に変換される]

\item9. 次の積分を計算せよ：
\[
\frac{e^{2x}}{(c^{2} e^{x} + a^{2} e^{-x})(c^{2} e^{x} + b^{2} e^{-x})}
\]
\( a \) が \( b \) と等しい場合と異なる場合を区別せよ。

\item10. 形 \( P(x, e^{ax}, e^{bx}, \dots) \) の任意の関数の積分法を証明せよ。ここで \( P \) は多項式を表す。[これは \( P \) が \( Ax^{m} e^{kx} \) のような型の項の和に表されることが知られているから]

\item11. 形 \( P(x, e^{ax}, e^{bx}, \dots, \cos lx, \cos mx, \dots, \sin lx, \sin mx, \dots) \) の任意の関数の積分法を示せ。

\item12. \( \lambda > 0 \) かつ \( a \) が \( R(x) \) の分母の最大の根より大きいとき、次の積分
\[
\int_{a}^{\infty} e^{-\lambda x} R(x) \, dx
\]
は収束することを証明せよ。[これは \( e^{\lambda x} \) が任意の \( x \) の冪よりも速く発散するという事実から]

\item13. \( \lambda > 0 \) に対して、次の積分
\[
\int_{-\infty}^{\infty} e^{-\lambda x^{2} + \mu x} \, dx
\]
は任意の \( \mu \) の値に対して収束し、また任意の正整数 \( n \) に対して次の積分も同様に収束する：
\[
\int_{-\infty}^{\infty} e^{-\lambda x^{2} + \mu x} x^{n} \, dx
\]

\item14. \( e^{x^{2}} \), \( e^{-x^{2}} \), \( xe^{x} \), \( xe^{-x} \), \( xe^{x^{2}} \), \( xe^{-x^{2}} \), および \( x \log x \) のグラフを描き、各関数の極大値・極小値とグラフ上の屈折点を求めよ。

\item 15. 方程式 \(e^{ax} = bx\) が正の定数 \(a\) と \(b\) の条件下で、二つの実数解を持つことを示せ。また、\(b > ae\) のときには一つの実数解もなしであることも示せ。[曲線 \(y = e^{ax}\) 上の点 \((\xi, e^{a\xi})\) における接線は
\[
y - e^{a\xi} = ae^{a\xi}(x - \xi)
\]
であり、この直線が原点を通るときには \(a\xi = 1\) である。したがって、直線 \(y = aex\) は点 \((1/a, e)\) で曲線に接している。これらの結果を図示することで明らかになる。読者は \(a\) や \(b\) のいずれかまたは両方が負の場合についても議論すべきである。]
\pagesepp{372}

\item 16. 方程式 \(e^{x} = 1 + x\) は \(x = 0\) を除き実数解を持たず、また方程式 \(e^{x} = 1 + x + \frac{1}{2}x^{2}\) は三つの実数解を持つことを示せ。

\item 17. 関数のグラフを描け：
\[
\log(x + \sqrt{x^{2} + 1}),\quad
\log\left(\frac{1 + x}{1 - x}\right),\quad
e^{-ax}\cos^{2}bx,\\
e^{-(1/x)^{2}},\quad
e^{-(1/x)^{2}}\sqrt{1/x},\quad
e^{-\cot x},\quad
e^{-\cot^{2} x}.
\]

\item 18. 方程式
\[
\log(x + \sqrt{x^{2} + 1}) = \frac{x}{100},\quad
e^{x} - \frac{2 + x}{2 - x} = \frac{1}{10\MC000},\quad
e^{x}\sin x = 7,\quad
e^{x^{2}}\sin x = 10\MC000
\]
の実数解の大まかな位置を決定せよ。

\item 19. \topic{双曲線関数。} 双曲線関数 \(\cosh x\)、\(\sinh x\)、\dots は方程式によって定義される：
\[
\cosh x = \frac{1}{2}(e^{x} + e^{-x}),\quad
\sinh x = \frac{1}{2}(e^{x} - e^{-x}),
\]
\[
\tanh x = \frac{\sinh x}{\cosh x},\quad
\coth x = \frac{\cosh x}{\sinh x},
\]
\[
\sech x = \frac{1}{\cosh x},\quad
\cosech x = \frac{1}{\sinh x}.
\]
これらの関数のグラフを描け。

\Item{20.} 以下の方程式を確立せよ：
\begin{gather*}
\cosh(-x) = \cosh x,\quad
\sinh(-x) = -\sinh x,\quad
\tanh(-x) = -\tanh x, \displaybreak\\
%
\cosh^{2} x - \sinh^{2} x = 1,\quad
\sech^{2} x + \tanh^{2} x = 1,\quad
\coth^{2} x - \cosech^{2} x = 1, \displaybreak\\
%
\cosh 2x = \cosh^{2} x + \sinh^{2} x,\quad
\sinh 2x = 2\sinh x\cosh x, \displaybreak\\
%
\begin{alignedat}{2}
\cosh(x + y) &= \cosh x\cosh y &&+ \sinh x\sinh y,\\
\sinh(x + y) &= \sinh x\cosh y &&+ \cosh x\sinh y.
\end{alignedat}
\end{gather*}

\Item{21.} これらの公式は、$\cos x$ および $\sin x$ の対応する公式から導かれるものと同様に、$\cosh x$ を $\cos x$ として $i\sinh x$ を $\sin x$ として書き換えることで得られる。

[これにより、すべての $\cos nx$ および $\sin nx$ の公式も同様に得られることになる。この類似性の理由は、第X章で明らかになります。]

\Item{22.} $\cosh x$ および $\sinh x$ を（a）$\cosh 2x$ （b）$\sinh 2x$ の形で表す。符号に関するあいまいさについて議論せよ。1908年数学問題。

\Item{23.} 以下を証明せよ：
\begin{gather*}
D_{x}\cosh x = \sinh x,\quad
D_{x}\sinh x = \cosh x,\\
%
D_{x}\tanh x = \sech^{2}x,\quad
D_{x}\coth x = -\cosech^{2}x,\\
%
D_{x}\sech x = -\sech x\tanh x,\quad
D_{x}\cosech x = -\cosech x\coth x,\\
%
D_{x}\log \cosh x = \tanh x,\quad
D_{x}\log|\sinh x| = \coth x,\\
%
D_{x}\arctan e^{x} = \tfrac{1}{2}\sech x,\quad
D_{x}\log |\tanh \tfrac{1}{2} x| = \cosech x.
\end{gather*}

[これらのすべての公式は、もちろん積分の公式に変換できる。]
373ページ

\Item{24.} $\cosh x > 1$ および $-1 < \tanh x < 1$ を証明せよ。

\Item{25.} $y = \cosh x$ のとき $x = \log\{y ± \sqrt{y^{2} - 1}\}$, $y = \sinh x$ のとき
$x = \log\{y + \sqrt{y^{2} + 1}\}$, $y = \tanh x$ のとき $x = \frac{1}{2}\log\{(1 + y)/(1 - y)\}$ であることを証明せよ。最初のケースにおける符号に関するあいまいさについて説明せよ。

26. 我々は、$\cosh x$, $\sinh x$, $\tanh x$ の逆関数をそれぞれ $\argcosh x$, $\argsinh x$, $\argtanh x$ で表すことにする。$\argcosh x$ は $x \geq 1$ のみで定義され、一般には二価であることを示せ。一方、$\argsinh x$ はすべての実数値 $x$ に対して定義され、また $\argtanh x$ は $-1 < x < 1$ の範囲で定義され、この2つの関数はいずれも一価であることを示せ。これらの関数のグラフを素描せよ。

27. $-\frac{1}{2}\pi < x < \frac{1}{2}\pi$ かつ $y$ が正のとき、$\cos x\cosh y = 1$ の条件下で、
\[
y = \log(\sec x + \tan x),\quad
D_{x} y = \sec x,\quad
D_{y} x = \sech y.
\]
を示せ。

28. $a > 0$ のとき、
\[
\int \frac{dx}{\sqrt{x^{2} + a^{2}}} = \argsinh(x/a)
\]
であり、また
\[
\int \frac{dx}{\sqrt{x^{2} - a^{2}}}
\]
は $x > 0$ のとき $\argcosh(x/a)$ で、$x < 0$ のとき $-\argcosh(-x/a)$ であることを証明せよ。

29. $a > 0$ のとき、
\[
\int \frac{dx}{x^{2} - a^{2}}
\]
は $|x|$ が $a$ より小さいとき $-(1/a)\argtanh(x/a)$ で、$|x|$ が $a$ より大きいとき $-(1/a)\argcoth(x/a)$ であることを証明せよ。[28問と29問の結果は第VI章の多くの公式を異なる方法で表すのに用いることができる.]

30. 
\begin{alignat*}{3}
\int \frac{dx}{\sqrt{(x - a)(x - b)}}
  \&=  \&\&2\log\{\sqrt{x - a} + \sqrt{x - b}\}      \&\&(a < b < x),\\
\int \frac{dx}{\sqrt{(a - x)(b - x)}}
  \&= -\&\&2\log\{\sqrt{a - x} + \sqrt{b - x}\}\quad \&\&(x < a < b),\\
\int \frac{dx}{\sqrt{(x - a)(b - x)}}
  \&=  \&\&2\arctan\left(\sqrt{\frac{x - a}{b - x}}\right)      \&\&(a < x < b).
\end{alignat*}

31. 
\[
\int_{0}^{1} x \log\left(1 + \frac{1}{2}x\right)\, dx
  = \frac{3}{4} - \frac{3}{2}\log\frac{3}{2}
  < \frac{1}{2} \int_{0}^{1} x^{2}\, dx
  = \frac{1}{6}.
\]
1913年数学 Tripos 問題

32. 方程式 \(a\cosh x + b\sinh x = c\) を解き、\(c > 0\) の条件下で、もし \(b^{2} + c^{2} - a^{2} < 0\) のとき、実数の根は存在しないことを示し、一方で \(b^{2} + c^{2} - a^{2} > 0\) のとき、正負の符号により二つの根、一つの根、または実数の根が存在しないことを証明せよ。また、\(b^{2} + c^{2} - a^{2} = 0\) の場合についても議論せよ。

33. 同時方程式 \(\cosh x\cosh y = a, \sinh x\sinh y = b\) を解け。

34. \(x^{1/x} \to 1\) である。\(x \to \infty\) のとき、\(x^{1/x} = e^{(\log x)/x}\) であり、\((\log x)/x \to 0\) であるから証明せよ。また、関数 \(x^{1/x}\) は \(x = e\) のときに最大値を取ることを示し、\(x > 0\) のときの関数のグラフを描け。

35. \(x^{x} \to 1\) である。\(x \to +0\) のとき、

36. \(\{f(n + 1)\}/\{f(n)\} \to l\) であり、\(l > 0\) のとき、\(n \to \infty\) のとき \(\sqrt[n]{f(n)} \to l\) である。証明は \(\log f(n + 1) - \log f(n) \to \log l\) より、\((1/n)\log f(n) \to \log l\) (\okrickRef{Ch.}{IV}, \MiscEx{IV}~27) を用いる。

37. \(\sqrt[n]{n!}/n \to 1/e\) である。\(n \to \infty\) のとき、\(f(n) = n^{-n} n!\) とすると、\(\{f(n + 1)\}/\{f(n)\} = \{1 + (1/n)\}^{-n} \to 1/e\) である。この結果を用いて例題36の方法で証明せよ。

38. \(\sqrt[n]{(2n)!/(n!)^{2}} \to 4\) である。\(n \to \infty\) のとき、

39. 方程式 \(e^{x} = x^{1\MC000\MC000}\) の近似解について議論せよ。

方程式は、一般のグラフィカルな考慮によれば、二つの正の根と一つの負の根を持つことが容易にわかります。正の根のうちの一つは約1より少し大きいもので、もう一つは非常に大きなものであり、負の根は約-1より少し大きいものです。大きな正の根の大きさを粗略に決定するためには、次のようにすることができます。もし $e^{x} = x^{1\MC000\MC000}$ であるとすると、

\[
x = 10^{6} \log x,\quad
\log x = 13.82 + \log\log x,\quad
\log\log x = 2.63 + \log \left(1 + \frac{\log\log x}{13.82}\right),
\]

と大まかに近似できます。ここで $13.82$ および $2.63$ はそれぞれ $\log 10^{6}$ および $\log\log 10^{6}$ の近似値です。これらの方程式から、比 $\log x : 13.82$ と $\log\log x : 2.63$ が大体1に近いことがわかります。したがって、

\[
x = 10^{6}(13.82 + \log\log x) = 10^{6}(13.82 + 2.63) = 16\MC450\MC000
\]

という値は根の粗略な近似を与えるものであり、誤差は $10^{6}(\log\log x - 2.63)$ や $(10^{6} \log\log x)/13.82$ や $(10^{6} × 2.63)/13.82$ で測定され、これらは約 $200,000$ 未満です。これらの近似は当然粗略ですが、根の尺度を把握するのに十分なものです。

40．同様に方程式
\[
e^{x} = 1\MC000\MC000 x^{1\MC000\MC000},\quad
e^{x^{2}} = x^{1\MC000\MC000\MC000}
\]
を議論します。

211．級数および積分の収束性に対する対数的テスト。第VIII章（§175以降）で示したように、

\[
\sum_{1}^{\infty} \frac{1}{n^{s}},\quad
\int_{a}^{\infty} \frac{dx}{x^{s}}\qquad (a > 0)
\]

は $s > 1$ のとき収束し、$s \leq 1$ のとき発散します。したがって $\sum (1/n)$ は発散し、しかしすべての正の値 $\alpha$ に対して $\sum n^{-1-\alpha}$ は収束します。

私たちは第200節で、対数を用いることで、$n \to \infty$ のとき $1/n$ より速くゼロに収束しながら $n^{-1-\alpha}$ より遅く収束する関数を作ることができるという事実を見た。ただし $\alpha$ は任意の正の数であってもよい。例えば、$1/(n \log n)$ はそのような関数であり、級数
\[
\sum \frac{1}{n \log n}
\]
が収束するか発散するかについては、どのような型の級数 $\sum n^{-s}$ と比較しても決定できない。

同様に、以下の級数も同様である：
\[
\sum \frac{1}{n (\log n)^2},\quad
\sum \frac{\log \log n}{n \sqrt{\log n}}.
\]
これらの級数が収束するか発散するかを決定するためのテストを見つけることは興味深い問題であり、第174節の積分テストから容易に導くことができる。

なぜなら
\[
D_x (\log x)^{1-s} = \frac{1 - s}{x (\log x)^s},\quad
D_x \log \log x = \frac{1}{x \log x},
\]
であるから、
\[
\int_a^\xi \frac{dx}{x (\log x)^s}
  = \frac{(\log \xi)^{1-s} - (\log a)^{1-s}}{1 - s},\quad
\int_{a}^{\xi} \frac{dx}{x \log x} = \log \log \xi - \log \log a,
\]
である。ただし $a > 1$。最初の積分は、$s > 1$ のとき $\xi \to \infty$ のとき $-(\log a)^{1-s}/(1 - s)$ に収束し、$s < 1$ のとき $\infty$ に収束する。二つ目の積分は $\infty$ に収束する。したがって、以下の級数と積分
\[
\sum_{n_0}^\infty \frac{1}{n (\log n)^s},\quad
\int_a^\infty \frac{dx}{x (\log x)^s},
\]
は $n_0$ および $a$ が 1 より大きいとき、$s > 1$ のとき収束し、$s \leq 1$ のとき発散する。

したがって、$\phi(n)$ が正で全ての $n$ に対してある定数 $K$ より小さいとき $\sum \phi(n)$ は収束し、$\phi(n)$ が正で全ての $n$ に対してある定数 $K$ より大きいとき $\sum \phi(n)$ は発散する。同様の定理が積分についても成り立つことは読者に任せることにする。

例 LXXXVIII．

1. 級数
\[
\sum \frac{1}{n(\log n)^{2}},\quad
\sum \frac{(\log n)^{100}}{n^{101/100}},\quad
\sum \frac{n^{2} - 1}{n^{2} + 1}\, \frac{1}{n(\log n)^{7/6}}
\]
は収束する。初めの級数の収束性は前節の定理の直接的な結果である。二つ目の級数の収束性は十分に大きい$n$に対して$(\log n)^{100}$は任意に小さな正の$\beta$を用いても$n^{\beta}$より小さいという事実から導かれる。従って、$\beta = 1/200$とすると、十分に大きい$n$に対して$(\log n)^{100} n^{-101/100}$は$n^{-201/200}$よりも小さい。三つ目の級数の収束性は前節末尾の比較試験から導かれる。

2. 級数
\[
\sum \frac{1}{n(\log n)^{6/7}},\quad
\sum \frac{1}{n^{100/101}(\log n)^{100}},\quad
\sum \frac{n\log n}{(n\log n)^{2} + 1}
\]
は発散する。

3. 級数
\[
\sum \frac{(\log n)^{p}}{n^{1+s}},\quad
\sum \frac{(\log n)^{p} (\log\log n)^{q}}{n^{1+s}},\quad
\sum \frac{(\log\log n)^{p}}{n(\log n)^{1+s}},
\]
$s > 0$のとき、すべての$p$と$q$に対して収束する。同様に、級数
\[
\sum \frac{1}{n^{1-s}(\log n)^{p}},\quad
\sum \frac{1}{n^{1-s}(\log n)^{p}(\log\log n)^{q}},\quad
\sum \frac{1}{n(\log n)^{1-s}(\log\log n)^{p}}
\]
は発散する。

\item 4. 以下のような級数の収束性または発散性についての問題
\[
\sum \frac{1}{n \log n \log \log n},\quad
\sum \frac{\log \log \log n}{n \log n (\log \log n)^{1/2}}
\]
は、定理 \PageRef{p.}{375} では解決できません。各々の場合において、級数の項は $1/(n \log n)$ より速くゼロに収束しますが、$n^{-1} (\log n)^{-1-\alpha}$ より遅い速度でゼロに収束します（$\alpha$ はどんなに小さな正数でもよい）。このような級数に対してはより精密なテストが必要です。読者は、以下の方程式から証明できるべきです：
\begin{align*}
D_{x} (\log_{k} x)^{1-s}
  &= \frac{1 - s}{x \log x \log_{2} x \dots \log_{k-1} x (\log_{k} x)^{s}},\\
D_{x} \log_{k+1} x
  &= \frac{1}{x \log x \log_{2} x \dots \log_{k-1} x \log_{k} x},
\end{align*}
ここで $\log_{2} x = \log \log x$, $\log_{3} x = \log \log \log x$ などから、以下の定理を証明すべきです：\emph{級数と積分
\[
\sum_{n_{0}}^{\infty} \frac{1}{n \log n \log_{2} n \dots \log_{k-1} n (\log_{k} n)^{s}},\quad
\int_{a}^{\infty} \frac{dx}{x \log x \log_{2} x \dots \log_{k-1} x (\log_{k} x)^{s}}
\]
は $s > 1$ のとき収束し、$s \leq 1$ のとき発散する}（$n_{0}$ と $a$ は $\log_{k} n$ および $\log_{k} x$ が $n \geq n_{0}$ や $x \geq a$ のとき正となるような十分に大きな数である）。これらの $n_{0}$ と $a$ の値は $k$ が増加するにつれて非常に急速に増大します：したがって $\log x > 0$ には $x > 1$, $\log_{2} x > 0$ には $x > e$, $\log_{3} x = \log \log \log x > 0$ には $x > e^{e}$ が必要であり、そのような値が容易に確認できます：$e^{e} > 10$, $e^{e^{e}} > e^{10} > 20,000$, $e^{e^{e^{e}}} > e^{20,000} > 10^{8000}$。

読者には、高い指数関数如きもの、例えば \(e^{e^x}\) および \(e^{e^{e^x}}\) のような関数が \(x\) を増加させるに従って極めて速く増大することを観察していただきたい。同じように、任意の値を持つ \(a > 1\) に対して \(a^{a^x}\) および \(a^{a^{a^x}}\) のような関数についても同様の注目すべき事実が言える。計算により \(9^{9^9}\) は 369,693,100桁であり、一方で \(10^{10^{10}}\) は当然 10,000,000,000桁であることが確認されている。逆に、高い対数関数の増大速度は非常に遅い。例えば \(\log\log\log\log x > 1\) を満たすには \(x\) が 8000桁以上の数である必要がある。

5. 次の積分 \(\ds\int_{0}^{a} \frac{1}{x} \left\{\log \left(\frac{1}{x}\right)\right\}^{s} dx\) (0 < a < 1) が \(s < -1\) のとき収束し、\(s \geq -1\) のとき発散することを証明せよ。[ε → +0 のとき次の積分の振る舞いを考えよ：
\[
\int_{\epsilon}^{a} \frac{1}{x} \left\{\log \left(\frac{1}{x}\right)\right\}^{s} dx
\]
この結果は、より高い対数因子を導入することでさらに精密化される可能性がある。]

6. 任意の \(s\) の値に対して \(\ds\int_{0}^{1} \frac{1}{x} \left\{\log \left(\frac{1}{x}\right)\right\}^{s} dx\) は意味を持たないことを証明せよ。[最後の例から、収束のためには \(s < -1\) が必要条件であることが示されている。しかし、\(s\) が負のとき \(\{\log(1/x)\}^s\) は \(x → 1 - 0\) のとき \((1 - x)^s\) に同値であり、したがって \(s < -1\) のとき上限において積分は発散する。]

7. \(\ds\int_{0}^{1} x^{a-1} \left\{\log \left(\frac{1}{x}\right)\right\}^{s} dx\) が収束するための必要十分条件は \(a > 0, s > -1\) である。

\par\medskip\noindent\textbf{【例 LXXXIX.】}\par
\Item{1.} エルミの極限について。\( n \)が無限大に近づくとき、
\[
\phi(n) = 1 + \frac{1}{2} + \frac{1}{3} + \dots + \frac{1}{n - 1} - \log n
\]
は極限\(\gamma\)に収束し、\(0 < \gamma \leq 1\)であることを示せ。[これは\SecNo[§]{174}からすぐに従う。実際の値は\(\gamma = .577\dots\)であり、通常\(\gamma\)はエルミの定数と呼ばれる.]

\Item{2.} \(a\)と\(b\)が正のとき、
\[
\frac{1}{a} + \frac{1}{a + b} + \frac{1}{a + 2b} + \dots + \frac{1}{a + (n - 1) b} - \frac{1}{b}\log(a + nb)
\]
は\(n\)が無限大に近づくとき、極限に収束する。

\Item{3.} \(0 < s < 1\)のとき、
\[
\phi(n) = 1 + 2^{-s} + 3^{-s} + \dots + (n - 1)^{-s} - \frac{n^{1-s}}{1 - s}
\]
は\(n\)が無限大に近づくとき、極限に収束する。

\Item{4.} 次の級数
\[
\frac{1}{1}
  + \frac{1}{2(1 + \frac{1}{2})}
  + \frac{1}{3(1 + \frac{1}{2} + \frac{1}{3})} + \dots
\]
は収束しないことを示せ。[級数の一般項を\(1/(n\log n)\)と比較する。]また、同様に上記の級数が\(\sum (1/n)\)から導かれるように、\(\sum n^{-s}\)から導かれた級数は\(s > 1\)のとき収束し、それ以外の場合は収束しないことを示せ。

\Item{5.} \(\sum u_{n}\)が正項の級数であり、
\[
s_{n} = u_{1} + u_{2} + \dots + u_{n},
\]
であるとき、\(\sum (u_{n}/s_{n-1})\)は\(\sum u_{n}\)が収束するか収束しないかに応じて収束または発散することを証明せよ。[\(\sum u_{n}\)が収束するとき、\(s_{n-1}\)は正の極限\(l\)に収束し、したがって\(\sum (u_{n}/s_{n-1})\)は収束する。\(\sum u_{n}\)が発散するとき、\(s_{n-1} \to \infty\)であり
\[
u_{n}/s_{n-1} > \log\{1 + (u_{n}/s_{n-1})\} = \log (s_{n}/s_{n-1})
\]
(\Ex{lxxxii}. 1)；これは明らかに
\[
\log(s_{2}/s_{1}) + \log(s_{3}/s_{2}) + \dots + \log(s_{n}/s_{n-1})
  = \log(s_{n}/s_{1})
\]
は\(n\)が無限大に近づくとき、無限大に収束することである。]

\item{6.} 同じ結果が級数 $\sum (u_{n}/s_{n})$ にも成り立つことを証明せよ。収束する場合、同様の証明が使用される。$\sum u_{n}$ が発散し、ある一定以上の $n$ の値から $u_{n} < s_{n-1}$ であるならば、$s_{n} < 2s_{n-1}$ であり、$\sum (u_{n}/s_{n-1})$ の発散性から $\sum (u_{n}/s_{n})$ も発散する。一方で、$u_{n} \geq s_{n-1}$ が無数の $n$ の値において成り立つ場合（急速に発散する級数の場合など）、すべてのこれらの $n$ の値に対して $u_{n}/s_{n} \geq \frac{1}{2}$ である。

\item{7.} 級数 $1 - \frac{1}{2} + \frac{1}{3} - \dots$ を求和せよ。次の式が成り立つ：
\[
1 + \frac{1}{2} + \dots + \frac{1}{2n} = \log(2n + 1) + \gamma + \epsilon_{n},
\quad
2\left(\frac{1}{2} + \frac{1}{4} + \dots + \frac{1}{2n}\right)
  = \log(n + 1) + \gamma + \epsilon_{n}',
\]
ここで $\gamma$ はオイラーの定数であり、$\epsilon_{n}$,~$\epsilon_{n}'$ は $n \to \infty$ のとき零に収束する数である。これらの式を引いて $n \to \infty$ とすると与えられた級数の和は $\log 2$ であることがわかる。詳しくは §213 を参照せよ。

\item{8.} 級数
\[
\sum_{0}^{\infty} (-1)^{n}\left(1 + \frac{1}{2} + \dots + \frac{1}{n + 1} - \log n - C\right)
\]
が有限振動し、$C = \gamma$ のとき収束する。

\Paragraph{212. 指数関数と対数関数に関連した級数。テイラーの定理による $e^{x}$ の展開。}
指数関数のすべての導函数はその自身に等しいため、次の式が成り立つ：
\[
e^{x} = 1 + x + \frac{x^{2}}{2!} + \dots + \frac{x^{n-1}}{(n - 1)!}
  + \frac{x^{n}}{n!} e^{\theta x}
\]
ここで $0 < \theta < 1$。また、$x^{n}/n! \to 0$ であり、$n \to \infty$ のとき任意の $x$ の値においても成り立つ（例 xxvii. 12）。さらに $e^{\theta x} < e^{x}$。したがって、$n$ を $\infty$ に近づけると次の式を得る：
\[
e^{x} = 1 + x + \frac{x^{2}}{2!} + \dots + \frac{x^{n}}{n!} + \dots.
\Tag{(1)}
\]

この方程式の右辺に記された級数は、\emph{指数級数}と呼ばれる。特に
\[
e = 1 + 1 + \frac{1}{2!} + \dots + \frac{1}{n!} + \dots;
\tag{(2)}
\]
よって
\[
\left(1 + 1 + \frac{1}{2!} + \dots + \frac{1}{n!} + \dots\right)^x = 1 + x + \frac{x^2}{2!} + \dots + \frac{x^n}{n!} + \dots,
\tag{(3)}
\]
という結果があり、これは\emph{指数定理}と呼ばれる。また
\[
a^x = e^{x \log a} = 1 + (x \log a) + \frac{(x \log a)^2}{2!} + \dots
\tag{(4)}
\]
がすべての正の$a$に対して成り立つ。

補足：
読者は、指数級数が各項を微分したときに自身を再現する性質を持つことに気づくだろう。他の$x$の冪の級数にはこの性質は存在しない：関連する追加の注記については、\okrickRef{付録}{II}を参照されたい。

冪級数 \(e^{x}\) は非常に重要であり、 Tayler の定理を使わない別の方法で調べる価値があります。
\[
E_{n}(x) = 1 + x + \frac{x^{2}}{2!} + \dots + \frac{x^{n}}{n!},
\]
を定義し、\(x > 0\) とする。すると
\[
\left(1 + \frac{x}{n}\right)^{n}
  = 1 + n\left(\frac{x}{n}\right)
      + \frac{n(n - 1)}{1·2} \left(\frac{x}{n}\right)^{2} + \dots
      + \frac{n(n - 1)\dots 1}{1·2\dots n} \left(\frac{x}{n}\right)^{n}
\]
であり、これは \(E_{n}(x)\) より小さい。また、\(n > x\) の条件下では、負の整数指数の二項定理により
\[
\left(1 - \frac{x}{n}\right)^{-n}
  = 1 + n\left(\frac{x}{n}\right)
      + \frac{n(n + 1)}{1·2} \left(\frac{x}{n}\right)^{2} + \dots
  > E_{n}(x).
\]
したがって
\[
\left(1 + \frac{x}{n}\right)^{n} < E_{n}(x) < \left(1 - \frac{x}{n}\right)^{-n}.
\]
しかし (\S208) には最初と最後の関数は \(n \to \infty\) のとき極限 \(e^{x}\) に収束し、したがって \(E_{n}(x)\) も同じである。このことから、\(x > 0\) のとき方程式 (1) が成り立つことが分かり、\(x < 0\) の場合も指数級数の性質により同様に成り立つ。

例題 XC.
1. \(\cosh x = 1 + \frac{x^{2}}{2!} + \frac{x^{4}}{4!} + \dots,\quad
\sinh x = x + \frac{x^{3}}{3!} + \frac{x^{5}}{5!} + \dots.\)

2. \(x > 0\) のとき指数級数の最大項は \([x] + 1\) 番目であり、\(x\) が整数でない場合、前項も等しい。

3. \(n! > (n/e)^{n}\) を示せ。［\(n^{n}/n!\) は \(e^{n}\) の級数の一つの項である］

4. \(e^{n} = (n^{n}/n!)(2 + S_{1} + S_{2})\) を証明し、ここで
\[
S_{1} = \frac{1}{1 + \nu} + \frac{1}{(1 + \nu)(1 + 2\nu)} + \dots,\quad
S_{2} = (1 - \nu) + (1 - \nu)(1 - 2\nu) + \dots,
\]
および \(\nu = 1/n\)；さらに \(n!\) は \(2(n/e)^{n}\) と \(2(n + 1)(n/e)^{n}\) の間にあることを導け。

\item5. 指数級数を利用すると、$e^x$は任意の$x$の冪より速く無限大に収束することを証明できる。[不等式 $e^x > \frac{x^n}{n!}$ を用いる。]

\item6. $e$は有理数ではないことを示す。[$e = p/q$ とし、$p$ および $q$ が整数であるとする。このとき
\[
\frac{p}{q} = 1 + \frac{1}{2!} + \frac{1}{3!} + \dots + \frac{1}{q!} + \dots
\]
または、両辺を $q!$ で掛けると
\[
q! \left(\frac{p}{q} - 1 - 1 - \frac{1}{2!} - \dots - \frac{1}{q!}\right) = \frac{1}{q + 1} + \frac{1}{(q + 1)(q + 2)} + \dots
\]
となり、左辺は整数であり、右辺は $\frac{1}{q + 1} + \left(\frac{1}{q + 1}\right)^2 + \dots = \frac{1}{q}$ より小さい。これは矛盾する。]

\item7. 級数 $\sum_{0}^{\infty} P_r(n) \frac{x^n}{n!}$ の和を求める。[$P_r(n)$ を
\[
A_0 + A_1 n + A_2 n(n - 1) + \dots + A_r n(n - 1) \dots (n - r + 1)
\]
と表すとき、
\begin{align*}
\sum_{0}^{\infty} P_r(n) \frac{x^n}{n!} &= A_0 \sum_{0}^{\infty} \frac{x^n}{n!} + A_1 \sum_{1}^{\infty} \frac{x^n}{(n - 1)!} + \dots + A_r \sum_{r}^{\infty} \frac{x^n}{(n - r)!}\\
&= (A_0 + A_1 x + A_2 x^2 + \dots + A_r x^r) e^x.
\end{align*}
]

\item8. 次の等式を証明する。
\[
\sum_{1}^{\infty} \frac{n^3}{n!} x^n = (x + 3x^2 + x^3) e^x, \quad
\sum_{1}^{\infty} \frac{n^4}{n!} x^n = (x + 7x^2 + 6x^3 + x^4) e^x;
\]
また、$S_n = 1^3 + 2^3 + \dots + n^3$ とするとき
\[
\sum_{1}^{\infty} S_n \frac{x^n}{n!} = \frac{1}{4}(4x + 14x^2 + 8x^3 + x^4) e^x.
\]
特に、最後の級数は $x = -2$ のとき零となる。

\item9. 次の等式を証明する。
\[
\sum \frac{n}{n!} = e, \quad
\sum \frac{n^2}{n!} = 2e, \quad
\sum \frac{n^3}{n!} = 5e;
\]
また、$k$が任意の正整数であるとき $\sum \frac{n^k}{n!}$ は $e$ の正の整数倍となる。

\item10. 次の等式を証明する。
\[
\sum_{1}^{\infty} \frac{(n - 1)x^n}{(n + 2)n!} = \left\{ (x^2 - 3x + 3)e^x + \frac{1}{2} x^2 - 3 \right\} / x^2.
\]

[M numeratorとdenominatorにn + 1を乗じて、Ex.~7と同様の手順で進めよ.]

\Item{11.} a, b, c を定め、(x + a)e^x + (bx + c) / x^3 が x → 0 のとき極限を持つようにし、その極限を評価し、関数 e^x + (bx + c) / (x + a) のグラフを描け。

\Item{12.} 1 + x, 1 + x + 1/2 x^2, 1 + x + 1/2 x^2 + 1/6 x^3 のグラフを描き、それらを e^x のグラフと比較せよ。

\Item{13.} e^{-x} - 1 + x - x^n / 2! + ... - (-1)^n x^n / n! が奇数または偶数であるとき正または負であることを証明し、指数定理を導け。PageSep{381}

\Item{14.} X$0 = e^x, X$1 = e^x - 1, X\\textit{2 = e^x - 1 - x, X}3 = e^x - 1 - x - (x^2 / 2!) ,... のとき dX$ν / dx = X$ν-1 である。従って、t > 0 のとき
X$1(t) = ∫$0^t X\\textit{0 dx < te^t, X}2(t) = ∫$0^t X$1 dx < ∫\\textit{0^t xe^x dx < e^t ∫}0^t x dx = t^2 / 2! e^t
であり、一般に X\_ν(t) < t^ν / ν! e^t であることを証明せよ。指数定理を導け。

\Item{15.} x^(2+p) = a^2 の正の根の p のべき級数展開は以下の項から始まる：
a {1 - 1/2 p log a + 1/8 p^2 log a (2 + log a)}。
\MathTrip{1909.}
\par\medskip

\paragraph{213. ログarithmic series.} 又、$x$ の冪級数展開において非常に重要なのは $\log(1 + x)$ に対する展開です。$\log(1 + x) = \int_{0}^{x} \frac{dt}{1 + t}$ であり、もし $t$ が 1 より小さいなら $1/(1 + t) = 1 - t + t^2 - \dots$ となるため、$\log(1 + x)$ は $-1 < x < 1$ のとき、級数 $1 - t + t^2 - \dots$ の各項を $t = 0$ から $t = x$ まで積分した級数 $x - \frac{1}{2} x^2 + \frac{1}{3} x^3 - \dots$ と等しいだろうと考えられます。実際、これが正しいです。

\[
1/(1 + t) = 1 - t + t^2 - \dots + (-1)^{m-1} t^{m-1} + \frac{(-1)^m t^m}{1 + t},
\]
よって $x > -1$ のとき
\[
\log(1 + x) = \int_{0}^{x} \frac{dt}{1 + t} = x - \frac{x^2}{2} + \dots + (-1)^{m-1} \frac{x^m}{m} + (-1)^m R_m,
\]
ここで
\[
R_m = \int_{0}^{x} \frac{t^m \, dt}{1 + t}.
\]

$R_m$ の $m \to \infty$ における極限がゼロであることを示す必要があります。これは $0 < x \leq 1$ のときほとんど明らかです；その場合、$R_m$ は正で
\[
\int_{0}^{x} t^m \, dt = \frac{x^{m+1}}{m + 1}
\]
より小さいため、従って $1/(m + 1)$ よりも小さいです。一方、$-1 < x < 0$ のとき、$t = -u$、$x = -\xi$ と置き換えることで
\[
R_m = (-1)^m \int_{0}^{\xi} \frac{u^m \, du}{1 - u}
\]
となり、$R_m$ の符号が $(-1)^m$ であることがわかります。また、積分範囲における $1/(1 - u)$ の最大値は $1/(1 - \xi)$ であり、
\[
0 < |R_m| < \frac{1}{1 - \xi} \int_{0}^{\xi} u^m \, du = \frac{\xi^m}{(m + 1)(1 - \xi)} < \frac{1}{(m + 1)(1 - \xi)}
\]
であるため、$R_m \to 0$。

したがって
\[
\log(1 + x) = x - \frac{1}{2} x^2 + \frac{1}{3} x^3 - \dots,
\]
が $-1 < x \leq 1$ のとき成立します。$x$ がこれらの範囲外にある場合、級数は収束しません。$x = 1$ のとき
\[
\log 2 = 1 - \frac{1}{2} + \frac{1}{3} - \dots,
\]
という結果を得ます。これは既に他の方法で証明済みです（\Ex{lxxxix}.7）。

\paragraph{214. 逆正接関数の級数}
同様な方法で証明すると、次のようになります：
\begin{align*}
\arctan x = \int_{0}^{x} \frac{dt}{1 + t^{2}}
  &= \int_{0}^{x}(1 - t^{2} + t^{4} - \dots)\, dt\\
  &= x - \tfrac{1}{3} x^{3} + \tfrac{1}{5} x^{5} - \dots,
\end{align*}
ただし $-1 \leq x \leq 1$ のときです。唯一の違いは、証明が少し簡単である点にあります；すなわち、$\arctan x$ は $x$ の奇関数であるため、正の値のみを考慮すればよいことになります。また、この級数は $x = -1$ のときも $x = 1$ のときも収束します。詳細な議論は読者に任せます。この級数によって表される $\arctan x$ の値は、$-1 \leq x \leq 1$ のとき $-\frac{1}{4}\pi$ と $\frac{1}{4}\pi$ の間のものであり、第 VII 章（\Ex{lxiii}.3）で見たように、この値は積分によっても表されます。$x = 1$ のときには、次の公式を得ます：
\[
\tfrac{1}{4}\pi = 1 - \tfrac{1}{3} + \tfrac{1}{5} - \dots.
\]

\par\medskip\noindent\textbf{【例 XCI.】}\par
\Item{1.} $-x \leq x < 1$ のとき、$\log \left(\dfrac{1}{1 - x}\right) = x + \frac{1}{2} x^{2} + \frac{1}{3} x^{3} + \dots$。

\Item{2.} $-1 < x < 1$ のとき、$\argtanh x = \frac{1}{2} \log\left(\dfrac{1 + x}{1 - x}\right) = x + \frac{1}{3} x^{3} + \frac{1}{5} x^{5} + \dots$。

\Item{3.} $x$ が正のとき、次のことを証明せよ：
\[
\log(1 + x) = \frac{x}{1 + x}
  + \tfrac{1}{2} \left(\frac{x}{1 + x}\right)^{2}
  + \tfrac{1}{3} \left(\frac{x}{1 + x}\right)^{3} + \dots。
\]
1911 年式。

\Item{4.} テイラーの定理を利用して $\log(1 + x)$ と $\arctan x$ の級数を得よ。

[最初の級数における $x$ が負の場合、ラグランジュの形式 $R_{n} = (-1)^{n-1} \frac{x^{n}}{n(1 + \theta x)^{n}}$ を用いるときの余項についての議論では困難が生じる。その場合、カuchy の形式、すなわち
\[
R_{n} = (-1)^{n-1} (1 - \theta)^{n-1} \frac{x^{n}}{(1 + \theta x)^{n}}
\]
を使用すべきである（二項級数の対応する議論を参照、\Ex{lvi}.2 および §163 参照）。
]

第二級数の場合には
\begin{align*}
D_{x}^{n} \arctan x
  &= D_{x}^{n-1} \left\{\frac{1}{1 + x^{2}}\right\}\\
  &= (-1)^{n-1} (n - 1)! (x^{2} + 1)^{-n/2} \sin \left\{n \arctan \left(\frac{1}{x}\right)\right\}
\end{align*}
（Ex. xlv. 11）であり、余項については問題ありません。その絶対値は明らかに$1/n$未満です。

公式 $D_{x}^{n} \arctan x$ は$x = 0$のとき失敗します。$\arctan(1/x)$が定義されないからです。$\arctan(1/x)$は$\frac{1}{2}\pi$と解釈されることが容易にわかります（Ex. xlv. 11 参照）。

5. $y > 0$のとき
\[
\log y = 2 \left\{\frac{y - 1}{y + 1} + \frac{1}{3} \left(\frac{y - 1}{y + 1}\right)^{3} + \frac{1}{5} \left(\frac{y - 1}{y + 1}\right)^{5} + \dots\right\}.
\]

この恒等式 $y = \left(1 + \frac{y - 1}{y + 1}\right) / \left(1 - \frac{y - 1}{y + 1}\right)$ を用います。この級数は$\log 2$を計算するのに使用できますが、収束が遅いため、級数 $1 - \frac{1}{2} + \frac{1}{3} - \dots$ より実用的には無意味です。$y = 2$と置き、小数点以下3桁で$\log 2$を求めます。

6. $\log 10$を小数点以下3桁で求めます。
\[
\log 10 = 3 \log 2 + \log \left(1 + \frac{1}{4}\right).
\]

7. $x > 0$のとき
\[
\log \left(\frac{x + 1}{x}\right) = 2 \left\{\frac{1}{2x + 1} + \frac{1}{3(2x + 1)^{3}} + \frac{1}{5(2x + 1)^{5}} + \dots\right\}
\]
であり、$x > 2$のとき
\[
\log \left(\frac{(x - 1)^{2}(x + 2)}{(x + 1)^{2}(x - 2)}\right) = 2 \left\{\frac{2}{x^{3} - 3x} + \frac{1}{3} \left(\frac{2}{x^{3} - 3x}\right)^{3} + \frac{1}{5} \left(\frac{2}{x^{3} - 3x}\right)^{5} + \dots\right\}.
\]
$\log 2 = .693\MS147\MS1\dots$、$\log 3 = 1.098\MS612\MS3\dots$と与えられたとき、$x = 10$を第二の公式に代入して$\log 11 = 2.397\MS895\dots$を示します。
数学トリップ 1912。

\item 8. $\log 2$, $\log 5$, および $\log 11$ が分かっている場合、以下の公式により $\log 13$ を計算することができます：
\[
\log 13 = 3 \log 11 + \log 5 - 9 \log 2
\]
この公式を使用すると、$\log 13$ の誤差はほぼ $0.000\MS15$ となります。数学トリップ 1910。

\item 9. 次の式が成り立ちます：
\[
\frac{1}{2} \log 2 = 7a + 5b + 3c, \quad
\frac{1}{2} \log 3 = 11a + 8b + 5c, \quad
\frac{1}{2} \log 5 = 16a + 12b + 7c,
\]
ここで $a = \argtanh(1/31)$, $b = \argtanh(1/49)$, $c = \argtanh(1/161)$。

これらの公式は、$\log 2$, $\log 3$, および $\log 5$ の高速かつ任意の精度での計算を可能にします。

\item 10. 次の式が成り立ちます：
\[
\frac{1}{4}\pi = \arctan(1/2) + \arctan(1/3) = 4\arctan(1/5) - \arctan(1/239),
\]
そして $\pi$ を小数点以下6桁で計算します。

\item 11. $(1 + x)^{1+x}$ の $x$ のべき乗展開は、次の項から始まります：
\[
1 + x + x^{2} + \frac{1}{2} x^{3}.
\]
数学トリップ 1910。

\item 12. 大きな値の $x$ に対して以下の公式が成り立ちます：
\[
\log_{10} e - \sqrtb{x(x + 1)} \log_{10}\left(\frac{1 + x}{x}\right)
  = \frac{\log_{10} e}{24x^{2}},
\]
この公式を $x = 10$ の場合に適用し、$\log_{10} e$ の近似値を求めます。その結果の精度を評価します。
数学トリップ 1910。

\item 13. $-1 < x < 1$ の範囲で以下の展開が成り立ちます：
\[
\frac{1}{1 - x} \log\left(\frac{1}{1 - x}\right)
  = x + \left(1 + \frac{1}{2}\right)x^{2}
      + \left(1 + \frac{1}{2} + \frac{1}{3}\right)x^{3} + \dots,
\]
[Ex. lxxxi. 2 を用います。]

\item 14. 指数級数と $\log_{10} 2.3758 = .375\MS809\MS9\dots$ および $\log_{10} e = .4343\dots$ の事実を用いて、方程式 $x = 100 \log_{10} x$ の近似解は $237.581\MS21$ となります。
数学トリップ 1910。

\item 15. $\log \cos x$ および $\log(\sin x/x)$ を $x^4$ まで $x$ のべき乗展開し、次の式が成り立つことを確認します：
\[
\log \sin x
  = \log x - \frac{1}{45} \log \cos x + \frac{64}{45} \log \cos \frac{1}{2}x.
\]
数学トリップ 1908。

\Item{16.} \( x \) が \(-1 \leq x \leq 1\) の範囲内であるとき、
\[
\int_{0}^{x} \frac{dt}{1 + t^{4}} = x - \tfrac{1}{5}x^{5} + \tfrac{1}{9}x^{9} - \dots
\]
を示せ。さらに、
\[
1 - \tfrac{1}{5} + \tfrac{1}{9} - \dots = \{\pi + 2\log(\sqrt{2} + 1)\}/4\sqrt{2}
\]
を得よ。

\MathTrip{1896.}

[同様に \SecNo[§]{214} の方法を用いて、\Ex{xlviii}.7 の結果を利用する。]

\Item{17.} 同様にして証明せよ：
\[
\tfrac{1}{3} - \tfrac{1}{7} + \tfrac{1}{11} - \dots = \int_{0}^{1} \frac{t^{2}\, dt}{1 + t^{4}} = \{\pi - 2\log(\sqrt{2} + 1)\}/4\sqrt{2}.
\]

\Item{18.} \( a \) と \( b \) を正の整数とするとき、
\[
\frac{1}{a} - \frac{1}{a + b} + \frac{1}{a + 2b} - \dots = \int_{0}^{1} \frac{t^{a-1}\, dt}{1 + t^{b}}
\]
を証明せよ。さらに、級数の和を求めよ。この方法で \( 1 - \frac{1}{4} + \frac{1}{7} - \dots \) と \( \frac{1}{2} - \frac{1}{5} + \frac{1}{8} - \dots \) の和を計算せよ。

\par\medskip

\Paragraph{215. バイノミアル級数。} 我々は既に (\SecNo[§]{163}) で、
\[
(1 + x)^{m} = 1 + \binom{m}{1}x + \binom{m}{2}x^{2} + \dots,
\]
を \( -1 < x < 1 \) および \( m \) が有理数であると仮定して調査した。\( m \) が無理数のとき、
\begin{gather*}
(1 + x)^{m} = e^{m\log(1+ x)},\\
D_{x}(1 + x)^{m} = \{m/(1 + x)\} e^{m\log(1 + x)} = m(1 + x)^{m-1},
\end{gather*}
となるので、\( (1 + x)^{m} \) の微分法則は同じであり、\SecNo[§]{163} で与えられた定理の証明もその有効性を失わない。\( x = 1 \) や \( x = -1 \) のときの級数収束に関する問題については議論しない。

\par\medskip\noindent\textbf{【例 XCII.】}\par
\Item{1.} \( -1 < x < 1 \) のとき、
\[
\frac{1}{\sqrt{1 + x^{2}}} = 1 - \frac{1}{2}x^{2} + \frac{1·3}{2·4}x^{4} - \dots,\quad
\frac{1}{\sqrt{1 - x^{2}}} = 1 + \frac{1}{2}x^{2} + \frac{1·3}{2·4}x^{4} + \dots.
\]

2. 二次根号および他の根号の近似

$\sqrt{M}$ を数値を求めたい二次根号とする。$N^{2}$ が $M$ の最も近い平方数であり、$M = N^{2} + x$ や $M = N^{2} - x$ と表される。ここで $x$ は正である。$x$ は $N$ よりも大きくなることはないため、$x/N^{2}$ は比較的小さい値であり、根号 $\sqrt{M} = N\sqrtb{1 ± (x/N^{2})}$ を級数で表すことができる。

\[
= N\left\{
  1 ± \frac{1}{2}\left(\frac{x}{N^{2}}\right)
    - \frac{1·1}{2·4}\left(\frac{x}{N^{2}}\right)^{2} ± \dots
\right\},
\]

この級数は少なくとも速やかに収束し、非常に速やかに収束する可能性がある。例えば、

\[
\sqrt{67} = \sqrtp{64 + 3}
  = 8\left\{
    1 + \frac{1}{2}\left(\frac{3}{64}\right)
      - \frac{1·1}{2·4}\left(\frac{3}{64}\right)^{2} + \dots
\right\}.
\]

$8\frac{3}{16}$（最初の二項による近似値）を用いて誤差が犯される。第二項以降では項は符号を交互に変えて減少するため、誤差は上回りであり、その大きさは $3^{2}/64^{2}$ 未満で $.003$ よりも小さい。

3. $x$ が $N^{2}$ に対して小さい場合、

\[
\sqrtp{N^{2} + x} = N + \frac{x}{4N} + \frac{Nx}{2(2N^{2} + x)},
\]

誤差は $x^{4}/N^{7}$ の量级である。この過程を $\sqrt{907}$ に適用する。

[二項定理により展開すると、

\[
\sqrtp{N^{2} + x}
  = N + \frac{x}{2N} - \frac{x^{2}}{8N^{3}} + \frac{x^{3}}{16N^{5}},
\]

誤差は次の項の数値未満である。すなわち $5x^{4}/128N^{7}$ 未満である。また、

\[
\frac{Nx}{2(2N^{2} + x)}
  = \frac{x}{4N} \left(1 + \frac{x}{2N^{2}}\right)^{-1}
  = \frac{x}{4N} - \frac{x^{2}}{8N^{3}} + \frac{x^{3}}{16N^{5}},
\]

誤差は $x^{4}/32N^{7}$ 未満である。結果はこれより得られる。同じ方法は二次根号以外の根号にも適用できる。例えば $\sqrt{1031}$ も同様に扱える。]
\pagebreak{386}

4. \( M \)が\( N^3 \)から1％未満の差異がある場合、\(\sqrt{M}\)は\(\frac{2}{3}N + \frac{1}{3}(M/N^2)\)から\( N/90 \MC000 \)未満の差異がある。

5. \( M = N^4 + x \)で、かつ\( x \)が\( N \)に対して小さい場合、\(\sqrt{M}\)の良い近似値は
\[
\frac{51}{56}N + \frac{5}{56} \frac{M}{N^3} + \frac{27Nx}{14(7M + 5N^4)}
\]
である。\( N = 10 \), \( x = 1 \)のとき、この近似値は小数点16桁まで正確であることを示せ。

6. 次の級数を求和せよ：
\[
\sum_{n=0}^{\infty} P_r(n) \binom{m}{n} x^n,
\]
ここで\( P_r(n) \)は\( n \)に関する次数\( r \)の多項式である。

[多項式 \( P_r(n) \)を形 \( A_0 + A_1 n + A_2 n(n-1) + \cdots \) として表すことを\Ex{xc}.7で示す。]

7. 次の級数を求和せよ：
\[
\sum_{n=0}^{\infty} n \binom{m}{n} x^n, \quad \sum_{n=0}^{\infty} n^2 \binom{m}{n} x^n
\]
さらに、
\[
\sum_{n=0}^{\infty} n^3 \binom{m}{n} x^n = \{ m^3 x^3 + m(3m - 1) x^2 + mx \}(1 + x)^{m-3}.
\]

216. 指数関数と対数関数の理論の別の展開法。次に、指数関数 \( e^x \) および対数関数 \( \log x \) の性質を調査するための方法について概略を与える。この方法は前頁で採用した論理的順序とは異なるものである。この方法は指数級数 \( 1 + x + \frac{x^2}{2!} + \cdots \) から始まる。この級数はすべての\( x \)に対する値に対して収束することが知られているため、関数 \( \exp x \) を方程式
\[
\exp x = 1 + x + \frac{x^2}{2!} + \cdots.
\Tag{(1)}
\]
によって定義することができる。

次に、\Ex{lxxxi}.7で示すように証明する：
\[
\exp x × \exp y = \exp(x + y).
\Tag{(2)}
\]

再度
\[
\frac{\exp h - 1}{h}
  = 1 + \frac{h}{2!} + \frac{h^{2}}{3!} + \dots
  = 1 + \rho(h),
\]
ここで $\rho(h)$は数値的に
\[
\left|\frac{1}{2}h\right| + \left|\frac{1}{2}h\right|^{2} + \left|\frac{1}{2}h\right|^{3} + \dots
  = \frac{\left|\frac{1}{2}h\right|}{1 - \left|\frac{1}{2}h\right|}
\]
よりも小さいので、$\rho(h) \to 0$ となります。従って
\[
\frac{\exp(x + h) - \exp x}{h}
  = \exp x \left(\frac{\exp h - 1}{h}\right) \to \exp x
\]
となり、$h \to 0$ のとき、または
\[
D_{x} \exp x = \exp x.
\Tag{(3)}
\]
さらに $\exp x$は連続関数であることを証明した。

今後、手順を選択する余地がある。$y = \exp x$とし、$\exp 0 = 1$を観察すると
\[
\frac{dy}{dx} = y,\quad
x = \int_{1}^{y} \frac{dt}{t},
\]
となる。そして、指数関数の逆関数として対数関数を定義すれば、この章の初期視点に戻る。

しかし異なる方法で進めることも可能である。式(2)から、正の整数$n$に対して
\[
(\exp x)^{n} = \exp nx,\quad
(\exp 1)^{n} = \exp n.
\]
$x$が正の有理数$m/n$の場合、
\[
\left\{\exp\left(\frac{m}{n}\right)\right\}^{n} = \exp m = (\exp 1)^{m},
\]
より $\exp(m/n)$は $(\exp 1)^{m/n}$の正の値と等しい。この結果を負の有理数$x$に対して
\[
\exp x \exp(-x) = 1;
\]
により拡張し、
\[
\exp x = (\exp 1)^{x} = e^{x},
\]
とすることができる。ここで
\[
e = \exp 1 = 1 + 1 + \frac{1}{2!} + \frac{1}{3!} + \dots,
\]
であり、すべての有理数$x$に対して成立する。最後に、無理数$x$に対して $e^{x}$を $\exp x$と定義する。対数関数は$\exp x$または$e^{x}$の逆関数として定義される。

例．二項級数理論
\[
1 + \binom{m}{1} x + \binom{m}{2} x^{2} + \dots = f(m, x),
\]
ここで $-1 < x < 1$から類似的に展開し、方程式
\[
f(m, x) f(m', x) = f(m + m'\Add{,} x)
\]
(\Ex{lxxxi}.~6)から始める。

\footnotetext{これらの例の多くはBromwichの『無限級数』から引用されています。}

\begin{Examples}{}
\Item{1.} $\log_{10} e = .4343$とし、$2^{10}$および$3^{21}$が$10$の近似べきであることを与えれば、$\log_{10}2$および$\log_{10}3$を小数点以下四位まで計算せよ。[1905年]

\Item{2.} $(\frac{1}{2}e)^{\sqrt{3}}$と$(\sqrt{2})^{\frac{1}{2}\pi}$のうち、大きい方を求めよ。[対数を取って$\sqrt{3}/(\sqrt{3} + \frac{1}{4}\pi) < \frac{2}{5} \sqrt{3} < .6929 < \log 2$と観察せよ。]

\Item{3.} $\log_{10}n$が有理数でないことを示せ。ただし$n$は$10$のべきでない正の整数とする。[もし$n$が$10$で割り切れないとき、$\log_{10}n = p/q$とすれば$10^{p} = n^{q}$となりこれは不可能である。なぜなら$10^{p}$は$0$で終わるのに対し、$n^{q}$はそうではないからである。また$n = 10^{a}N$とし、$N$が$10$で割り切れないとき$\log_{10}N$および]
\[
\log_{10}n = a + \log_{10}N
\]
は有理数でない。

\Item{4.} どのような$x$の値に対して、関数$\log x$, $\log\log x$, $\log\log\log x$,~\dots\ (a)~$0$に等しい (b)~$1$に等しい (c)~定義されていないか？また同様の問いを関数$lx$, $llx$, $lllx$,~\dots\ に対して考える。ここで$lx = \log |x|$とする。

\Item{5.} 
\[
\log x - \binom{n}{1} \log(x + 1) + \binom{n}{2} \log(x + 2) - \dots
  + (-1)^{n} \log(x + n)
\]
は負であり、$x$が$0$から$\infty$に向かって増加するにつれて定常的に$0$に近づく。

[関数の導関数は
\[
\sum_{0}^{n} (-1)^{r} \binom{n}{r} \frac{1}{x + r}
  = \frac{n!}{x(x + 1) \dots (x + n)}
\]
であり、これは容易に部分分数分解により確認できる。この式は正であり、$x \to \infty$のとき関数自体も$0$に収束するためである。なぜなら
\[
\log(x + r) = \log x + \epsilon_{x},
\]
ここで$\epsilon_{x} \to 0$であり、$1 - \dbinom{n}{1} + \dbinom{n}{2} - \dots = 0$であるからである。]

6. \( x \)に関する微分作用素の\( n \)乗を用いて、次の式を証明せよ：
\[
\left(\frac{d}{dx}\right)^n \frac{\log x}{x} = \frac{(-1)^n n!}{x^{n+1}} \left( \log x - 1 - \frac{1}{2} - \dots - \frac{1}{n} \right).
\]
（1909）

7. \( x > -1 \)のとき、\( x^2 > (1 + x) \{\log(1 + x)\}^2 \)である。 （1906）

[置換 \( 1 + x = e^\xi \) し、\( \xi > 0 \) のとき \( \sinh \xi > \xi \) という事実を利用する.]

8. \( x \)が \( 0 \)から無限大へ増加するにつれて、\(\frac{\log(1 + x)}{x}\)と\(\frac{x}{(1 + x)\log(1 + x)}\)は両方とも単調に減少することを示せ。

9. \( -1 \)から無限大へ増加するにつれて、関数 \((1 + x)^{-1/x}\) は \( 0 \)と\( 1 \)の間で一度だけすべての値を取り出すことを示せ。 （1910）

10. \( x \to 0 \) のとき、\(\frac{1}{\log(1 + x)} - \frac{1}{x} \to \frac{1}{2}\)であることを示せ。

11. \( x \to 0 \) のとき、\( -1 \)から無限大へ増加するにつれて、関数 \(\frac{1}{\log(1 + x)} - \frac{1}{x}\)は単調に減少し、値が \( 1 \)から \( 0 \)へと減少することを示せ。 [\( x = 0 \)のとき関数は定義されていないが、\( x = 0 \)のとき値を\(\frac{1}{2}\)として与えることで、\( x = 0 \)での連続性を得る。この結果を示すために、Ex.7を利用すること.]

12. \( \xi \)が正であるとき、関数 \(\frac{\log \xi - \log x}{\xi - x}\)は \( x \)が \( 0 \)から \(\xi\)へ増加するにつれて単調に減少し、\( x \to \xi \)のときの極限を求めよ。

13. \( M \)と \( N \)が大きな正数であるとき、\( x \)が \( 2\log M \)と \( 16N^2 \)の大きい方より大きいとき、次の不等式が成り立つことを示せ：
\[ e^x > Mx^N. \]
[容易に証明できるのは、\(\log x < 2\sqrt{x}\)である。したがって、与えられた不等式は
\[
x > \log M + 2N\sqrt{x}
\]
を満たすとき、そして従って
\[
\frac{1}{2}x > \log M, \quad \frac{1}{2}x > 2N\sqrt{x}
\]
を満たすときも成り立つ。]

14. \( f(x) \)と\( \phi(x) \)が \( x \to \infty \)のとき無限大に収束し、かつ \( \frac{f'(x)}{\phi'(x)} \to \infty \)であるとき、\( \frac{f(x)}{\phi(x)} \to \infty \)であることを示せ。 [Ch. VI, Ex. 33 の結果を利用する。] \( f(x) = x^\alpha \), \( \phi(x) = \log x \)と置き換えることで、すべての正の\(\alpha\)に対して \(\frac{\log x}{x^\alpha} \to 0\)であることを証明せよ。

\item[15.] \( p \) と \( q \) が正の整数であるとき、
\[
\frac{1}{pn + 1} + \frac{1}{pn + 2} + \cdots + \frac{1}{qn}
  \to \log\left(\frac{q}{p}\right)
\]
であり、これは \( n \to \infty \) のときである。[参照：\Ex{lxxviii}.~6.]

\item[16.] 正の数 \( x \) に対して、
\[
n\log\left\{\frac{1}{2}(1 + x^{1/n})\right\} \to -\frac{1}{2}\log x
\]
であり、これは \( n \to \infty \) のときである。[式は
\[
n\log\left\{\frac{1}{2}(1 + x^{1/n})\right\}
  = n\log\left\{1 - \frac{1}{2}(1 - x^{1/n})\right\}
  = \frac{1}{2}n(1 - x^{1/n}) \frac{\log(1 - u)}{u}
\]
ここで \( u = \frac{1}{2}(1 - x^{1/n}) \)。次に、\SecNo[§]{209} および \Ex{lxxxii}.~4. を用いる.]

\item[17.] 正の数 \( a \) と \( b \) に対して、
\[
\left\{\frac{1}{2}(a^{1/n} + b^{1/n})\right\}^n \to \sqrt{ab}.
\]

\%[** TN: No paragraph break in the original]
対数を取って \Ex{16.} を用いる。

\item[18.] 次の式が成り立つことを示す。
\[
1 + \frac{1}{3} + \frac{1}{5} + \cdots + \frac{1}{2n - 1}
  = \frac{1}{2}\log n + \log 2 + \frac{1}{2} \gamma + \epsilon_n,
\]
ここで \( \gamma \) はオイラーの定数（\Ex{lxxxix}.~1）であり、\( \epsilon_n \to 0 \) である。

\item[19.] 次の式が成り立つことを示す。
\[
1 + \frac{1}{3} - \frac{1}{2} + \frac{1}{5}
  + \frac{1}{7} - \frac{1}{4} + \frac{1}{9} + \cdots
  = \frac{3}{2} \log 2,
\]
この級数は、正の項を2つと負の項を1つ交互に取り、元の級数 \( 1 - \frac{1}{2} + \frac{1}{3} - \cdots \) を形成したものである。最初の \( 3n \) 項の和は
\begin{multline*}
1 + \frac{1}{3} + \frac{1}{5} + \cdots + \frac{1}{4n - 1}
  - \frac{1}{2} \left(1 + \frac{1}{2} + \cdots + \frac{1}{n}\right)\\
  = \frac{1}{2}\log 2n + \log 2 + \frac{1}{2}\gamma + \epsilon_n
  - \frac{1}{2}(\log n + \gamma + \epsilon_n'),
\end{multline*}
ここで \( \epsilon_n \) と \( \epsilon'_n \) は \( n \to \infty \) のとき \( 0 \) に収束する。[参照：\Ex{lxxviii}.~6]。

\item[20.] 次の式が成り立つことを示す。
\[
1 - \frac{1}{2} - \frac{1}{4} + \frac{1}{3} - \frac{1}{6} - \frac{1}{8} + \frac{1}{5} - \frac{1}{10} - \cdots = \frac{1}{2}\log 2.
\]

21. \( n \)までの項の和を証明する。
\[
\sum_{1}^{n} \frac{1}{\nu(36\nu^{2} - 1)} = -3 + 3\Sigma_{3n+1} - \Sigma_{n} - S_{n}
\]
ここで \( S_{n} = 1 + \dfrac{1}{2} + \dots + \dfrac{1}{n} \), \(\Sigma_{n} = 1 + \dfrac{1}{3} + \dots + \dfrac{1}{2n - 1} \)。したがって、無限級数の和は
\[
-3 + \tfrac{3}{2}\log 3 + 2\log 2
\]
であることを証明する。

22. 次の級数の和を示す。
\[
\sum_{1}^{\infty} \frac{1}{n(4n^{2} - 1)} = 2\log 2 - 1, \quad
\sum_{1}^{\infty} \frac{1}{n(9n^{2} - 1)} = \tfrac{3}{2}(\log 3 - 1).
\]

23. 次の4つの級数の和を証明する。
\[
\sum_{1}^{\infty} \frac{1}{4n^{2} - 1}, \quad
\sum_{1}^{\infty} \frac{(-1)^{n-1}}{4n^{2} - 1}, \quad
\sum_{1}^{\infty} \frac{1}{(2n + 1)^{2} - 1}, \quad
\sum_{1}^{\infty} \frac{(-1)^{n-1}}{(2n + 1)^{2} - 1}
\]
それぞれの和は \(\frac{1}{2}\), \(\frac{1}{4}\pi - \frac{1}{2}\), \(\frac{1}{4}\), \(\frac{1}{2}\log 2 - \frac{1}{4} \)である。

24. \( n! (a/n)^n \)が \( a < e \)のとき0に収束し、\( a > e \)のとき無限大に発散することを証明する。
[もし \( u_n = n! (a/n)^n \)とすると, \( u_{n+1}/u_n = a\{1 + (1/n)\}^{-n} \to a/e \)である。\( a = e \)のとき関数が無限大に発散することを示すためには、この章の定理を超える範囲の証明が必要で、Bromwich の『無限級数』pp. 461 et seq. を参照せよ。]

25. \( x \to \infty \)のとき次の式の極限を求めよ。
\[
\left(\frac{a_{0} + a_{1} x + \dots + a_{r} x^{r}}
           {b_{0} + b_{1} x + \dots + b_{r} x^{r}}\right)^{\lambda_{0}+\lambda_{1}x}
\]
異なる場合を区別して証明する。 \MathTrip{1886.}

26. \( x > 0 \)のとき次の級数の和が無限大に発散することを示す。
\[
\sum \log \left(1 + \frac{x}{n}\right)
\]
従って、\( x \)が正のとき
\[
(1 + x)(2 + x) \dots (n + x)/n! \to \infty
\]
であることを導く。[関数の対数は \(\sum_{1}^{n} \log \left(1 + \dfrac{x}{\nu}\right)\) である。]

\item27. \( x > -1 \)のとき、次の等式を証明せよ：
\[
\frac{1}{(x + 1)^2}
= \frac{1}{(x + 1)(x + 2)}
+ \frac{1!}{(x + 1)(x + 2)(x + 3)}
+ \frac{2!}{(x + 1)(x + 2)(x + 3)(x + 4)} + \cdots.
\]
\mathtrip{1908.}

\(\frac{1}{(x + 1)^2}\)と級数の最初の\( n \)項の和との差は
\[
\frac{1}{(x + 1)^2} \cdot \frac{n!}{(x + 2)(x + 3) \cdots (x + n + 1)}.
\]

\item28. \( A e^{\alpha x} + B e^{\beta x} + \cdots = 0 \)という形の方程式は、\( A, B, \cdots \)が多項式であり、\(\alpha, \beta, \cdots\)が異なる実数であるとき、すべての\( x \)に対して成り立つことはない。[\(\alpha\)が\(\alpha, \beta, \cdots\)の中で代数的に最大のものであるとき、\( x \to \infty \)のとき項\( Ae^{\alpha x} \)は他のすべてを上回る.]

\item29. 次の数列
\[
a_1 = e,\quad a_2 = e^{e^2},\quad a_3 = e^{e^{e^3}},\ \cdots
\]
が指数尺度の任意の項よりも速く無限大に収束することを示せ。
[Let \( e_1(x) = e^x \), \( e_2(x) = e^{e_1(x)} \)，そして次のように定義する。もし\( e_k(x) \)が指数尺度の任意の項であるとき、\( n > k \)のとき\( a_n > e_k(n) \).]

\item30. 次の等式を証明せよ：
\[
\frac{d}{dx} \{\phi(x)\}^{\psi(x)}
= \frac{d}{dx} \{\phi(x)\}^{\alpha} + \frac{d}{dx} \{\beta^{\psi(x)}\}
\]
ここで、微分後に\(\alpha = \psi(x) \)，\(\beta = \phi(x) \)とする。同様の規則を\(\phi(x)^{[\{\psi(x)\}^{\chi(x)}]} \)の微分にも確立せよ。

\item{31.} \(D_{x}^{n} e^{-x^{2}} = e^{-x^{2}} \phi_{n}(x)\)であることを証明すると、(i) \(\phi_{n}(x)\)は次数\(n\)の多項式である、(ii) \(\phi_{n+1} = -2x\phi_{n} + \phi_{n}'\)であり、(iii) \(\phi_{n} = 0\)のすべての根は実数で互いに異なるものであり、かつ\(\phi_{n-1} = 0\)の根と区切られている。[結果が\(\DPtypo{n}{\kappa} = 1, 2, \dots, \DPtypo{\kappa}{n}\)に対して真であると仮定し、\(\DPtypo{\phi_{\kappa+1}}{\phi_{n+1}}\)の符号を\(\DPtypo{\phi_{\kappa}}{\phi_{n}} = 0\)となる\(x\)の\(n\)つの値および大規模な（正または負）\(x\)の値に対して考慮する。]

\item{32.} \(f(xy) = f(x)f(y)\)の一般解は\(x^{a}\)であり、\(a\)は定数である：また\(\cosh ax\)または\(\cos ax\)は\(f''(0)\)が正か負かにより、\[
f(x + y) + f(x - y) = 2f(x)f(y)
\]の解である。[第二結果を証明するには、\(f\)が一階、二階、三階微分を持つと仮定する。次に
\[
2f(x) + y^{2}\{f''(x) + \epsilon_{y}\}
  = 2f(x)[f(0) + yf'(0) + \tfrac{1}{2} y^{2}\{f''(0) + \epsilon_{y}'\}],
\]
ここで、\(\epsilon_{y}\)と\(\epsilon_{y}'\)は\(y\)とともに零に収束する。従って、\(f(0) = 1, f'(0) = 0, f''(x) = f''(0)f(x)\)であり、したがって \(a = \sqrtb{f''(0)}\)または\(a = \sqrtb{-f''(0)}\)である。]

\item{33.} \(x \to +0\)のとき、関数 \(x^{\sin(1/x)}, x^{\sin^{2}(1/x)}, x^{\cosec(1/x)}\) の振る舞いはいかなるものか？

\item{34.} 曲線 \(y = \tan x e^{\tan x}, y = \sin x \log \tan \frac{1}{2}x\) を追跡せよ。

\item{35.} 方程式 \(e^{x} = ax + b\)は\(a < 0\)または\(a = 0, b > 0\)のとき一つの実数解を持つ。また、\(a > 0\)のとき、\(a \log a > b - a\)の場合に二つの実数解を持つか、それともない。

\item{36.} \(e^{x} = ax^{2} + 2bx + c\)の方程式は\(a > 0\)のとき一つ、二つ、または三つの実数解を持つか、\(a < 0\)のとき一つもしくは二つの実数解を持たないかをグラフィカルな考慮により示し、異なる場合を区別する方法を示せ。

\item{37.} 曲線 \( y = \dfrac{1}{x} \log\left(\dfrac{e^{x} - 1}{x}\right) \) を追跡し、点 \( (0, \frac{1}{2}) \) が対称中心であることを示し、また \( x \) の値がすべての実数を経過するにつれて \( y \) は定常的に \( 0 \) より大きく \( 1 \) に近づくことを証明せよ。この結果から方程式
\[
\dfrac{1}{x} \log\left(\dfrac{e^{x} - 1}{x}\right) = \alpha
\]
が実数解を持たない場合、\( 0 < \alpha < 1 \) であることが必要であり、そのとき一方の実数解があり、その符号は \( \alpha - \frac{1}{2} \) の符号と一致することを導き出せ。まず
\[
y - \tfrac{1}{2}
  = \dfrac{1}{x} \left\{\log\left(\dfrac{e^{x} - 1}{x}\right) - \log e^{\frac{1}{2} x}\right\}
  = \dfrac{1}{x} \log\left(\dfrac{\sinh \frac{1}{2}x}{\frac{1}{2}x}\right)
\]
は明らかに \( x \) の奇関数である。また
\[
\dfrac{dy}{dx}
  = \dfrac{1}{x^{2}} \left\{\tfrac{1}{2} x \coth \tfrac{1}{2}x - 1
      - \log\left(\dfrac{\sinh \frac{1}{2}x}{\frac{1}{2}x}\right)\right\}.
\]
大括弧内の関数は \( x \to 0 \) のとき零に収束し、その導函数は
\[
\dfrac{1}{x} \left\{1 - \left(\dfrac{\frac{1}{2}x}{\sinh \frac{1}{2}x}\right)^2\right\}
\]
であり、これは \( x \) の符号と一致する。したがって \( dy/dx > 0 \) である。

\item{38.} 曲線 \( y = e^{1/x} \sqrt{x^{2} + 2x} \) を追跡し、方程式
\[
e^{1/x} \sqrt{x^{2} + 2x} = \alpha
\]
が負の \(\alpha\) のとき実数解を持たないことを示し、\( 0 < \alpha < a = e^{1/\sqrt{2}} \sqrt{2 + 2\sqrt{2}} \) のとき一方の負の実数解を持つことを証明せよ。さらに \( \alpha > a \) のとき二つの正の実数解と一方の負の実数解を持つことを示せ。

\item{39.} 方程式 \( f_{n}(x) = 1 + x + \dfrac{x^{2}}{2!} + \dots + \dfrac{x^{n}}{n!} = 0 \) が奇数である \( n \) のとき一方の実数解を持ち、偶数である \( n \) のとき実数解を持たないことを示せ。

[$n = 1$, $2$, … $2k$についてこの命題が成り立つと仮定すると、$f_{2k+1}(x) = 0$は次数が奇数であるため少なくとも一つの実根を持つ。また、その根がより多くなることはない、なぜならそれであれば$f'_{2k+1}(x)$または$f_{2k}(x)$は少なくとも一度消失しなければならないからである。したがって、$f_{2k+1}(x) = 0$はただ一つの実根を持つことになり、これにより$f_{2k+2}(x) = 0$は最多二つの実根を持つことができない。もし二つあるとすると、$\alpha$および$\beta$とするとき、$f'_{2k+2}(x)$または$f_{2k+1}(x)$は$\alpha$および$\beta$の間で少なくとも一度消失しなければならない。そして
\[
f_{2k+2}(\gamma) = f_{2k+1}(\gamma) + \frac{\gamma^{2k+2}}{(2k + 2)!} > 0.
\]
また、$x$が大きいとき（正の方向でも負の方向でも）$f_{2k+2}(x)$も正であり、図をみるとこれらの結果は矛盾していることがわかる。したがって、$f_{2k+2}(x) = 0$は実根を持たない。

40. $a$および$b$が正でほぼ等しいとき
\[
\log \frac{a}{b} = \frac{1}{2}(a - b) \left(\frac{1}{a} + \frac{1}{b}\right),
\]
近似して、誤差は$\frac{1}{6}\{(a - b)/a\}^{3}$程度であることを証明せよ。対数級数を利用せよ。この公式は歴史的に興味深いものであり、ナピアが対数の計算に使用した。

41. $-1 < x < 1$の場合、級数乗法を用いて
\begin{align*}
\tfrac{1}{2}\{\log(1 + x)\}^{2}
  &= \tfrac{1}{2} x^{2}
   - \tfrac{1}{3}(1 + \tfrac{1}{2})x^{3}
   + \tfrac{1}{4}(1 + \tfrac{1}{2} + \tfrac{1}{3})x^{4} - \dots,\\
\tfrac{1}{2}(\arctan x)^{2}
  &= \tfrac{1}{2} x^{2}
   - \tfrac{1}{4}(1 + \tfrac{1}{3})x^{4}
   + \tfrac{1}{6}(1 + \tfrac{1}{3} + \tfrac{1}{5})x^{6} - \dots.
\end{align*}

42. 以下を証明せよ：
\[
(1 + \alpha x)^{1/x}
  = e^{\alpha}\{1 - \tfrac{1}{2} a^{2}x
                 + \tfrac{1}{24}(8 + 3a)a^{3}x^{2}(1 + \epsilon_{x})\},
\]
ここで$\epsilon_{x} \to 0$である。

43. 1 + x + \frac{x^{2}}{2!} + \dots + \frac{x^{n}}{n!} の自然対数の x のべき級数展開の最初の n + 2 項は
\[
x - \frac{x^{n+1}}{n!}
  \left\{\frac{1}{n + 1}
       - \frac{x}{1!\, (n + 2)}
       + \frac{x^{2}}{2!\, (n + 3)} - \dots
       + (-1)^{n} \frac{x^{n}}{n!\, (2n + 1)}
  \right\}.
\]

44. 
\[
\exp \left(-x - \frac{x^{2}}{2} - \dots - \frac{x^{n}}{n}\right)
\]
の x のべき級数展開は最初に以下の項から始まる：
\[
1 - x + \frac{x^{n+1}}{n + 1}
  - \sum_{s=1}^{n} \frac{x^{n+s+1}}{(n + s)(n + s + 1)}.
\]

45. -1 < x < 1 のとき、以下の式が成り立つ：
\begin{align*}
\frac{1}{3}x + \frac{1·4}{3·6}2^{2}x^{2} + \frac{1·4·7}{3·6·9}3^{2}x^{3} + \dots
  &= \frac{x(x + 3)}{9(1 - x)^{7/3}},\\
\frac{1}{3}x + \frac{1·4}{3·6}2^{3}x^{2} + \frac{1·4·7}{3·6·9}3^{3}x^{3} + \dots
  &= \frac{x(x^{2} + 18x + 9)}{27(1 - x)^{10/3}}.
\end{align*}

[これらの結果は、例 xcii.6 の方法を使用して証明できます。ただし、無限級数の微分に関する問題は私たちの範囲外です。より簡単に得られる方法として、微分法がありますが、その詳細は省略します。]

\item 46. \(a\) と \(b\) が正のとき、次の積分を証明せよ：
\[
\int_{0}^{\infty} \frac{dx}{(x + a)(x + b)} = \frac{1}{a - b} \log\left(\frac{a}{b}\right),
\]
\[
\int_{0}^{\infty} \frac{dx}{(x + a)(x + b)^2} = \frac{1}{(a - b)^2 b} \left\{a - b - b \log\left(\frac{a}{b}\right)\right\},
\]
\[
\int_{0}^{\infty} \frac{x \, dx}{(x + a)(x + b)^2} = \frac{1}{(a - b)^2} \left\{a \log\left(\frac{a}{b}\right) - a + b\right\},
\]
\[
\int_{0}^{\infty} \frac{dx}{(x + a)(x^2 + b^2)} = \frac{1}{(a^2 + b^2) b} \left\{\frac{1}{2}\pi a - b \log\left(\frac{a}{b}\right)\right\},
\]
\[
\int_{0}^{\infty} \frac{x \, dx}{(x + a)(x^2 + b^2)} = \frac{1}{a^2 + b^2} \left\{\frac{1}{2}\pi b + a \log\left(\frac{a}{b}\right)\right\}.
\]
これらの積分の値を推定し、独立に検証せよ。各関数
\[
a - 1 - \log a, \quad a \log a - a + 1, \quad \frac{1}{2} \pi a - \log a, \quad \frac{1}{2} \pi + a \log a
\]
がすべての正の \(a\) の値に対して正であることを確認せよ。

\item 47. \(\alpha\), \(\beta\), \(\gamma\) が全て正であり、\(\beta^2 > \alpha \gamma\) のとき、次の積分を証明せよ：
\[
\int_{0}^{\infty} \frac{dx}{\alpha x^2 + 2 \beta x + \gamma} = \frac{1}{\sqrt{\beta^2 - \alpha \gamma}} \log \left\{\frac{\beta + \sqrt{\beta^2 - \alpha \gamma}}{\sqrt{\alpha \gamma}}\right\}.
\]
また、\(\alpha\) が正であり、\(\alpha \gamma > \beta^2\) のとき、積分の値は
\[
\int_{0}^{\infty} \frac{dx}{\alpha x^2 + 2 \beta x + \gamma} = \frac{1}{\sqrt{\alpha \gamma - \beta^2}} \arctan \left\{\frac{\sqrt{\alpha \gamma - \beta^2}}{\beta}\right\}
\]
であり、逆接続の値は \(0\) と \(\pi\) の間を選ぶこと。他の本当に異なる場合で収束する積分があるか否かを確認せよ。

48. \(a > -1\) のとき、以下の積分が成り立つことを証明せよ：
\[
\int_{1}^{\infty} \frac{dx}{(x + a)\sqrt{x^{2} - 1}}
  = \int_{0}^{\infty} \frac{dt}{\cosh t + a}
  = 2\int_{1}^{\infty}\frac{du}{u^{2} + 2au + 1};
\]
さらに、積分の値は以下のようになる：
\[
-1 < a < 1 のとき
\quad \frac{2}{\sqrt{1 - a^{2}}} \arctan \sqrt{\frac{1 - a}{1 + a}}
\]
また、\(a > 1\) のとき
\[
\frac{1}{\sqrt{a^{2} - 1}} \log \frac{\sqrt{a + 1} + \sqrt{a - 1}}{\sqrt{a + 1} - \sqrt{a - 1}}
  = \frac{2}{\sqrt{a^{2} - 1}} \operatorname{arctanh} \sqrt{\frac{a - 1}{a + 1}};
\]
また、\(a = 1\) の場合についても議論せよ。

49. \(a > 0\) のとき、以下の積分を同様に変形し、その値は次のようになることを示せ：
\[
\int_{0}^{\infty} \frac{dx}{(x + a) \sqrt{x^{2} + 1}}
  = \frac{1}{\sqrt{a^{2} + 1}} \log \frac{a + 1 + \sqrt{a^{2} + 1}}{a + 1 - \sqrt{a^{2} + 1}}
  = \frac{2}{\sqrt{a^{2} + 1}} \operatorname{arctanh} \frac{\sqrt{a^{2} + 1}}{a + 1}.
\]

50. 次の積分の値は次のようになることを証明せよ：
\[
\int_{0}^{1} \arctan x\, dx = \frac{1}{4}\pi - \frac{1}{2}\log 2.
\]

51. \(0 < \alpha < 1\)、\(0 < \beta < 1\) のとき、以下の積分の値は次のようになることを証明せよ：
\[
\int_{-1}^{1} \frac{dx}{\sqrt{(1 - 2\alpha x + \alpha^{2})(1 - 2\beta x + \beta^{2})}}
  = \frac{1}{\sqrt{\alpha\beta}} \log \frac{1 + \sqrt{\alpha\beta}}{1 - \sqrt{\alpha\beta}}.
\]

52. \(a > b > 0\) のとき、以下の積分の値は次のようになることを証明せよ：
\[
\int_{-\infty}^{\infty} \frac{d\theta}{a\cosh \theta + b\sinh \theta}
  = \frac{\pi}{\sqrt{a^{2} - b^{2}}}.
\]

53. 次の積分の値は次のようになることを証明せよ：
\[
\int_{0}^{1} \frac{\log x}{1 + x^{2}}\, dx
  = -\int_{1}^{\infty} \frac{\log x}{1 + x^{2}}\, dx,\quad
\int_{0}^{\infty} \frac{\log x}{1 + x^{2}}\, dx = 0;
\]
さらに、\(a > 0\) のとき、以下の積分の値は次のようになることを導き出せ：
\[
\int_{0}^{\infty} \frac{\log x}{a^{2} + x^{2}}\, dx = \frac{\pi}{2a}\log a.
\]
証明には \(x = 1/t\)、\(x = au\) の置換を利用する。

54. \(a > 0\) のとき、以下の積分の値は次のようになることを証明せよ：
\[
\int_{0}^{\infty} \log \left(1 + \frac{a^{2}}{x^{2}}\right) dx = \pi a.
\]
[部分積分を利用する。]

\chapter{第10章：Chapter X: THE GENERAL THEORY OF THE LOGARITHMIC, EXPONENTIAL, AND CIRCULAR FUNCTIONS}

\par\medskip\noindent\textbf{【対数関数、指数関数および三角関数の一般理論】}\par

\par\medskip\noindent\textbf{217. 複素変数の函数}\par第III章において，複素変数
\[
z = x + iy
\]
を定義した。この章では，$x + iy$ の代わりに $x + yi$ を書くことが一般的に便利であるため，しばしば後者を使用する。

いくつかのクラスの $z$ を含む式について，多項式 $P(z)$ などの単純な性質を考慮した。このような式を $z$ の関数と呼ぶのは自然であり，実際には，多項式 $P(z)$ および $Q(z)$ の比 $P(z)/Q(z)$ を有理関数と呼んだ。

しかし，$z$ の関数の一般的な定義は与えていない。

Itは$z$についての関数を定義する際に、実変数$x$についての関数を定義したときと同じ方法で定義するのが自然に思われ Oversight: The translation should maintain the original meaning and structure without unnecessary oversight.

以下のように翻訳します：

$z$についての関数を定義する際に、実変数$x$についての関数を定義したときと同じ方法で定義するのが自然に思われるかもしれません。すなわち、$z$と$Z$間にどのような関係が存在すれば、あるまたはすべての$z$の値に対応する$Z$の値や値群があるという関係が存することを意味して$Z$は$z$の関数であると言えるでしょう。しかし、精密に検討すると、この定義からは何らの利益も得られないことがわかります。なぜなら、$z$が与えられれば$x$と$y$も与えられることになり、逆もまた真です。$z$の値を指定することは、$x$と$y$のペアの値を指定することと同じことです。したがって、提案された定義によれば、「$z$についての関数」は実変数$x$と$y$の複素関数
\[
f(x, y) + ig(x, y)
\]
に他なりません。例えば
\[
x - iy,\quad xy,\quad |z| = \sqrt{x^{2} + y^{2}},\quad \arg z = \arctan(y/x)
\]
は「$z$についての関数」です。この定義は完全に正当ですが、無意味であり、新しい概念を真正に定義していないからです。したがって、「複素変数$z$についての関数」という表現を使う際には、より制限された意味で使用することが便利です。すなわち、実変数$x$と$y$の一般的な複素関数のクラスの中から特別なクラスを選出し、その表現が限定されるようにします。しかし、この選択方法を説明し、選ばれた特別なクラスの関数の特性を説明しようとすると、本書の範囲を超えるところまで行かざるを得ません。したがって、一般的な定義は与えず、直接定義された特殊な関数に限定して扱います。

218. 既に \( z \) の多項式（§39）、\( z \) の有理函数（§46）および \( z \) の根（§47）を定義しました。実変数 \( x \) の場合で与えた（§26--§27）明示的および暗黙的な代数函数の定義は、複素変数へ拡張するには難しくありません。これらのすべての場合において、我々は複素数 \( z \) を、点 \( z \) の座標（§44）である \( z \) の引数、関数 \( f(z) \) の引数と呼びます。この章で扱う問題は、対数函数、指数函数および三角函数や円関数の定義と主要な性質を定める問題です。これらの函数は当然、実数値 \( z \) に対してのみ定義されており、対数函数は正の値に対してのみ定義されています。

対数函数について最初に扱います。その定義を以下の拡張により試みます：
\[
\log x = \int_{1}^{x} \frac{dt}{t}\quad (x > 0);
\]
これを行うために、積分概念のいくつかの拡張について簡単に考察する必要があります。

219. 実数および複素曲線積分。弧 \( AB \) は方程式
\[
x = \phi(t),\quad y = \psi(t),
\]
で定義される曲線 \( C \) の一部であり、ここで \(\phi\) と \(\psi\) は連続微分係数 \(\phi'\) および \(\psi'\) を持つ関数です。そして、\( t \) が \( t_{0} \) から \( t_{1} \) に変化するとき、点 \( (x, y) \) は同じ方向で曲線を動いて \( A \) から \( B \) へ移動します。

次に、曲線積分
\[
\int_{C} \{g(x, y)\, dx + h(x, y)\, dy\}
\]
を定義する。ここで $g$ および $h$ は $x$ と $y$ の連続な関数であり、これは形式的な置換 $x = \phi(t)$, $y = \psi(t)$ を通じて得られる通常の積分に等しいものとする。
すなわち
\[
\int_{t_{0}}^{t_{1}} \{g(\phi, \psi) \phi' + h(\phi, \psi) \psi'\}\, dt
\]
である。このとき、$C$ を積分路と呼ぶ。

次に、
\[
z = x + iy = \phi(t) + i\psi(t)
\]
とする。これにより $t$ が変化するとき Argand の図において $z$ は曲線 $C$ 上を描く。
さらに、
\[
f(z) = u + iv
\]
は $z$ の多項式または有理関数であるとする。

このとき、
\[
\int_{C} f(z)\, dz
\]
を定義する。これは
\[
\int_{C} (u + iv) (dx + i\, dy)
\]
と定義され、さらにこれは
\[
\int_{C} (u\, dx - v\, dy) + i\int_{C} (v\, dx + u\, dy)
\]
と定義される。

220. $\Log \zeta$ の定義。任意の複素数 $\zeta = \xi + i\eta$ に対して，一般対数 $\Log \zeta$ を以下の式で定義する：
\[
\Log \zeta = \int_{C} \frac{dz}{z},
\]
ここで $C$ は $1$ から $\zeta$ まで始まり，原点を通らない曲線である。したがって（図54参照），道 (ia), (ib), (ic) は定義に含まれるような道の例である。$\Log z$ の値は特定の積分路を選んだときに定義される。しかし，現在では，その定義から導かれる $\Log z$ の値がどの程度積分路によって依存するかは明確でない。例えば，$\zeta$ が実数かつ正であるとする，すなわち $\xi$ に等しいとする。このときの積分路の一つは $1$ から $\xi$ までの直線であり，これを $x = t$, $y = 0$ の方程式で定義できる。この場合，特定の積分路を選んだときに
\[
\Log \xi = \int_{1}^{\xi} \frac{dt}{t}
\]
となり，$\Log \xi$ は $\log \xi$ に等しい。すなわち，$\xi$ が実数かつ正であるときの $\Log \xi$ の一つの値は $\log \xi$ である。しかし，この場合も一般の場合も，積分路は無限多様な方法で選べる。すべての $\Log \xi$ の値が必ずしも $\log \xi$ に等しいとは限らないことが示されることはなく，実際にはそのようなことは起こらないことがわかる。これが $\Log \zeta$, $\Log \xi$ の記号を $\log \zeta$, $\log \xi$ から採用した理由である。$\Log \xi$ は（おそらく）多価関数であり，$\log \xi$ はその一つの値に過ぎない。そして，現在までで見られる限り，一般の場合には三つの可能性が等しく存在する：

(1) どの積分路を選んでも $1$ から $\zeta$ までの $\Log \zeta$ の値が常に同じである。

\Item{(2)} 異なる経路に対応する異なる値が得られる場合がある;

\Item{(3)} いくつかの異なる値が得られ、それぞれは一連の経路に対応している:

これらのいずれかの選択肢の真偽は、私たちの定義によって暗示されるものではない。
(PageSep{399})

\Paragraph{221. $\Log \zeta$ の値} 点 $z = \zeta$ の極座標を $\rho$ と $\phi$ とするので、
\[
\zeta = \rho(\cos\phi + i\sin\phi).
\]
現在、$-\pi < \phi < \pi$ であり、$\rho$ は任意の正の値を持つ。したがって、$\zeta$ は零または実数の負の値以外の任意の値を持つ。

経路 $C$ 上の任意の点 $(x, y)$ の座標は $t$ の関数であり、その極座標 $(r, \theta)$ も同様である。また
\begin{align*}
\Log \zeta
  &= \int_{C} \frac{dz}{z}
   = \int_{C} \frac{dx + i\, dy}{x + iy} \\
  &= \int_{t_{0}}^{t_{1}} \frac{1}{x + iy} \left(\frac{dx}{dt} + i\frac{dy}{dt}\right) dt,
\end{align*}
は \SecNo[§]{219} の定義による。しかし $x = r\cos\theta$, $y = r\sin\theta$ であり、
\begin{align*}
\frac{dx}{dt} + i\frac{dy}{dt}
  &= \left(\cos\theta\, \frac{dr}{dt} - r\sin\theta\, \frac{d\theta}{dt}\right)
  + i\left(\sin\theta\, \frac{dr}{dt} + r\cos\theta\, \frac{d\theta}{dt}\right) \\
  &= (\cos\theta + i\sin\theta) \left(\frac{dr}{dt} + ir\frac{d\theta}{dt}\right);
\end{align*}
したがって
\[
\Log \zeta
  = \int_{t_{0}}^{t_{1}} \frac{1}{r}\, \frac{dr}{dt}\, dt
 + i\int_{t_{0}}^{t_{1}} \frac{d\theta}{dt}\, dt
  = [\log r] + i[\theta],
\]
ここで $[\log r]$ は $t = t_{1}$ と $t = t_{0}$ の点に対応する $\log r$ の値の差を表し、$[\theta]$ も同様の意味を持つ。

明確には
\[
[\log r] = \log \rho - \log 1 = \log \rho;
\]
ですが、$[\theta]$ の値は少し考慮が必要です。まず、積分路が $1$ から $\zeta$ までの直線である場合を考えましょう。$\theta$ の初期値は $1$ の振幅の一つ、すなわち $2k\pi$（$k$ は任意の整数）です。仮に初期の $\theta = 2k\pi$ とします。図から明らかに、$t$ が直線を動くにつれて $\theta$ は $2k\pi$ から $2k\pi + \phi$ へと増加します。したがって
\[
[\theta] = (2k\pi + \phi) - 2k\pi = \phi,
\]
そして、積分路が直線であるとき、$\Log \zeta = \log \rho + i\phi$ です。
ページ分割点{400}

この特定の値を $\Log \zeta$ の主値と呼びます。$\zeta$ が実数で正の場合、$\zeta = \rho$ かつ $\phi = 0$ なので、主値は通常の対数 $\log \zeta$ です。従って、一般に主値を表すために $\Log \zeta$ を $\log \zeta$ で表すことが便利です。したがって
\[
\log \zeta = \log \rho + i\phi,
\]
そして主値はその虚部が $-\pi$ から $\pi$ の間にあることを特徴づけます。

次に、図56で示すような任意の経路を考える。この経路と点1からζまでの直線が囲む領域（または領域）が原点を含まない場合がある。このような経路上では、θは依然としてφに等しいことが容易に理解できる。例えば、図で実線で示したような曲線上では、θは初期値2kπから最初に
\[
2k\pi - XOP
\]
まで減少し、その後再び増加してQ点で2kπとなり、最後には2kπ + φとなる。点線の経路は、直線と曲線が原点を含まない2つの領域を囲む相似だがやや複雑な例を示している。したがって、

\[
\Log \zeta = \log \zeta = \log \rho + i\phi
\]
となる。

一方で、θがφに等しくないような積分経路も容易に構成できる。例えば、図57で実線で示したような曲線を考える。θが初期値2kπである場合、P点まで到達するときには2π増加し、Q点まで到達するときには4π増加することになる；そして最終的な値は
\[
2k\pi + 4\pi + \phi
\]
となり、θは4π + φであり
\[
\Log \zeta = \log \rho + i(4\pi + \phi)
\]
となる。

この場合、積分路は正の方向で原点を2周します。もし$k$回原点を周り、そのような経路を選んだとすれば、同様の方法で
\[
[\theta] = 2k\pi + \phi
\]
および
\[
\Log \zeta = \log \rho + i(2k\pi + \phi)
\]
が得られます。ここで$k$は正の整数です。経路を原点に対して反対方向に回すと（図57の虚線経路のように）、$k$は負の整数になります。$\zeta$の絶対値$|\zeta| = \rho$であり、異なる角度$2k\pi + \phi$が$\am \zeta$の異なる値であることに注意すると、すべての$\log |\zeta| + i\am \zeta$の値は$\Log \zeta$の値と一致することが分かります。前段落での議論から明らかに、すべての$\Log \zeta$の値はこのような形で表されます。

結論をまとめると以下の通りです：\par\medskip\noindent\textbf{【結果】}\par$\Log \zeta$の一般値は
\[
\log |\zeta| + i\am \zeta = \log \rho + i(2k\pi + \phi)
\]
であり、ここで$k$は任意の正または負の整数です。絶対値が決定される経路によって$k$の値が定まります。もし経路が一直線であれば$k = 0$で
\[
\Log \zeta = \log \zeta = \log \rho + i\phi
\]
となります。\par\medskip

これまで、$\zeta$を関数$\Log \zeta$の引数とし、$(\xi, \eta)$または$(\rho, \phi)$を$\zeta$の座標とし、$z$, $(x, y)$, $(r, \theta)$を積分経路上の任意の点およびその座標と表記してきましたが、現在では関数$\Log z$の引数として$z$を使用することも可能です。次に示す例でこの表記法に戻ります。

例 XCIII．

1. 以上で \(-\pi < \theta < \pi\) と仮定し、\(z\) が実数かつ負の値である場合を除外しました。この場合、点 \(1\) から点 \(z\) への直線は \(0\) を通りますので、積分経路として適さないものとなります。両端 \(\pi\) および \(-\pi\) は \(\am z\) の値であり、\(\theta\) はこれらいずれかの値に等しいことになります。また \(r = -z\) です。任意の値の \(\Log z\) の実部と虚部は \(x\) および \(y\) に対して連続な関数ですが、\(x = 0\), \(y = 0\) の場合を除きます。

2. \(\Log z\) の任意の値の実部と虚部は \(x\) および \(y\) に対して連続な関数です。

3. \(\Log z\) が満足する機能方程式。関数 \(\Log z\) は以下の等式を満たします：
\[
\Log z_{1}z_{2} = \Log z_{1} + \Log z_{2},
\]
これは両辺の任意の値が他方のいずれかの値と一致することを意味します。これは、
\[
z_{1} = r_{1}(\cos\theta_{1} + i\sin\theta_{1}),\quad
z_{2} = r_{2}(\cos\theta_{2} + i\sin\theta_{2}),
\]
と置き、第 401 頁の公式を適用することですぐに分かります。しかし、すべての場合においては以下の等式が成り立つわけではありません：
\[
\log z_{1}z_{2} = \log z_{1} + \log z_{2}.
\]
例えば、
\[
z_{1} = z_{2} = \frac{1}{2}(-1 + i\sqrt{3})
  = \cos \frac{2}{3}\pi + i \sin \frac{2}{3}\pi,
\]
とすると、\(\log z_{1} = \log z_{2} = \frac{2}{3}\pi i\) であり、\(\log z_{1} + \log z_{2} = \frac{4}{3}\pi i\) ですが、これは \(\Log z_{1}z_{2}\) の主値ではありません。実際には、\(\log z_{1}z_{2} = -\frac{2}{3}\pi i\) です。

方程式（1）のようなものにおいて、一方の値が他方の値と一致するとき、その方程式を完全な方程式または完全に真なる方程式と呼びます。

4. 方程式 $\Log z^{m} = m\Log z$ は完全には真ではありません。右辺のすべての値は左辺の値と一致しますが、逆は成り立ちません。

5. 方程式 $\Log (1/z) = -\Log z$ は完全に真です。また、$\log (1/z) = -\log z$ も真ですが、実数で負のときを除きます。

6. 方程式
\[
\log \left(\frac{z - a}{z - b}\right) = \log (z - a) - \log (z - b)
\]
は、点 $z = a$ と $z = b$ を結ぶ直線およびこれらの点を通る $OX$ 方向に無限に延びる直線で囲まれた領域の外側にあるときには真です。

7. 方程式
\[
\log \left(\frac{a - z}{b - z}\right)
  = \log \left(1 - \frac{a}{z}\right) - \log \left(1 - \frac{b}{z}\right)
\]
は、点 $O$、$a$、$b$ を頂点とする三角形の外側にあるときには真です。

8. 実数変数 $x$ の関数 $\Imag(\Log x)$ のグラフを描きます。[このグラフは $y = 2k\pi$ の正半分と $y = (2k + 1)\pi$ の負半分で構成されます。]

9. 実数変数 $x$ の定義された関数 $f(x)$ は
\[
\pi f(x) = p\pi + (q - p)\Imag(\log x)
\]
であり、$x > 0$ のとき $p$ に等しく、$x < 0$ のとき $q$ に等しい。

10. 実数変数 $x$ の定義された関数 $f(x)$ は
\[
\pi f(x) = p\pi + (q - p)\Imag\{\log(x - 1)\} + (r - q)\Imag(\log x)
\]
であり、$x > 1$ のとき $p$ に等しい、$0 < x < 1$ のとき $q$ に等しい、$x < 0$ のとき $r$ に等しい。

11. 方程式 $\log z$ (i) および任意の値 $\Log z$ (ii) が実数または虚数単位純粋な形となるときの $z$ の値は何か？（a）実数か （b）純粋虚数か？

\Item{12.} $z = x + iy$ のとき、$\Log\Log z = \log R + i(\Theta + 2k'\pi)$ であり、
\[
R^{2} = (\log r)^{2} + (\theta + 2k\pi)^{2}
\]
で、$\Theta$ は方程式群
\[
\cos\Theta : \sin\Theta : 1 ::
\log r : \theta + 2k\pi: \sqrtb{(\log r)^{2} + (\theta + 2k\pi)^{2}}
\]
によって決定される最小の正角である。
$1 + i\sqrt{3}$ の $\Log\Log$ の双無限集合の値を概略的に描き、そのうち哪些が $\log\Log(1 + i \sqrt{3})$ の値であり、哪些が $\Log\log(1 + i\sqrt{3})$ の値であるか示せ。

\Paragraph{222. 指数関数。} 第 IX 章で、実変数 $y$ に対する関数 $e^{y}$ を、関数 $y = \log x$ の逆関数として定義した。複素変数 $z$ に対する関数を $\Log z$ の逆関数として定義すべき自然な提案がなされる。

\par\medskip\noindent\textbf{【定義】}\par
任意の値 $\zeta$ に対して $\Log z = \zeta$ であるとき、$z$ を $\zeta$ の指数と呼び、
\[
z = \exp \zeta.
\]
\par\medskip

したがって $z = \exp \zeta$ は $\zeta = \Log z$ のときである。任意の与えられた値 $z$ に対して無数の異なる値 $\zeta$ に対応することが確実である。逆に、任意の与えられた値 $\zeta$ に対して無数の値 $z$ が存在すると考えるのは自然であるかのように思われるが、これは定理によって証明される。

\par\medskip\noindent\textbf{【定理】}\par
指数関数 $\exp \zeta$ は $\zeta$ の単值関数である。
\par\medskip

\[
z_{1} = r_{1}(\cos\theta_{1} + i\sin\theta_{1}),\quad z_{2} = r_{2}(\cos\theta_{2} + i\sin\theta_{2})
\]
がともに$\exp \zeta$の値であると仮定すると、
\[
\zeta = \Log z_{1} = \Log z_{2},
\]
であり、これより
\[
\log r_{1} + i(\theta_{1} + 2m\pi) = \log r_{2} + i(\theta_{2} + 2n\pi),
\]
ここで$m$と$n$は整数である。この式から
\[
\log r_{1} = \log r_{2},\quad \theta_{1} + 2m\pi = \theta_{2} + 2n\pi,
\]
を得る。したがって$r_1 = r_2$であり、$\theta_{1}$と$\theta_{2}$は$2\pi$の倍数で異なる。従って
\[
z_{1} = z_{2}.
\]

\Paragraph{コロラリー}
もし$\zeta$が実数であるなら、$\exp \zeta = e^{\zeta}$であり、これは$\zeta$の実数指数関数である。

もし$z = e^{\zeta}$とすると、$\log z = \zeta$すなわち$\Log z$の値の一つが$\zeta$である。従って
\[
z = \exp \zeta.
\]

\Paragraph{223. $\exp \zeta$の値}
$\zeta = \xi + i\eta$とし、
\[
z = \exp \zeta = r(\cos\theta + i\sin\theta).
\]
すると
\[
\xi + i\eta = \Log z = \log r + i(\theta + 2m\pi),
\]
ここで$m$は整数である。従って$\xi = \log r$, $\eta = \theta + 2m\pi$であり、または
\[
r = e^{\xi},\quad \theta = \eta - 2m\pi;
\]
よって
\[
\exp (\xi + i\eta) = e^{\xi}(\cos\eta + i\sin\eta).
\]

もし$\eta = 0$とすると、$\exp \xi = e^{\xi}$であり、これは既に第222節で推測した通りである。明らかに$\exp (\xi + i\eta)$の実部と虚部は全ての$\xi$と$\eta$に対して$\xi$と$\eta$の連続関数である。

224. $\exp \zeta$ が満足する関数方程式

$\zeta_{1} = \xi_{1} + i\eta_{1}$, $\zeta_{2} = \xi_{2} + i\eta_{2}$ とすると、

\begin{align*}
\exp \zeta_{1} × \exp \zeta_{2}
  &= e^{\xi_{1}} (\cos\eta_{1} + i\sin\eta_{1})
   × e^{\xi_{2}} (\cos\eta_{2} + i\sin\eta_{2}) \\
  &= e^{\xi_{1}+\xi_{2}} \{\cos(\eta_{1} + \eta_{2}) + i\sin(\eta_{1} + \eta_{2})\} \\
  &= \exp(\zeta_{1} + \zeta_{2}).
\end{align*}

したがって、指数関数は実数 $\zeta_{1}$, $\zeta_{2}$ のときを含め、

$f(\zeta_{1} + \zeta_{2}) = f(\zeta_{1}) f(\zeta_{2})$

という関数方程式を満足することが証明されている（第205節）。

225. 一般のべき乗 $a^{\zeta}$

実数 $\zeta$ のとき $\exp \zeta = e^{\zeta}$ であることに鑑み、複素数 $\zeta$ のときも同じ表記を採用し、$\exp \zeta$ の表記を完全に廃止すべきかと思われたかもしれない。しかし、私たちは $e^{\zeta}$ の意味をより一般的な定義を与える必要があり、

このとき $e^{\zeta}$ は $\exp \zeta$ のみならず無数の値を持つ関数を表すことが明らかになる。

既に実数正または負の $a$ および有理数 $\zeta$ の場合、または偶数でない分母を持つ有理数 $\zeta$ の場合における符号 $a^{\zeta}$ の意味を定義している。初等代数学では、

$a^{\zeta}$ は最大でも2つの値を持つ。

第3章では任意の実数または複素数 $a$ および任意の有理数 $\zeta = p/q$ の場合にこの定義を拡張し、

第9章では新たな定義、すなわち

$a^{\zeta} = e^{\zeta \log a}$

という方程式を与え、これは実数 $\zeta$ および実数正の $a$ のときには常に成り立つ。

次のように、様々な表現に意味を付与しました：
\[
3^{1/2},\quad (-1)^{1/3},\quad (\sqrt{3} + \frac{1}{2}i)^{-1/2},\quad (3.5)^{1+\sqrt{2}};
\]
しかし、次のような表現には意味を付与する定義がまだ提供されていません：
\[
(1 + i)^{\sqrt{2}},\quad 2^{i},\quad (3 + 2i)^{2+3i}.
\]
現在、$a$ と $\zeta$ が実数または複素数である限り、$a$ が零でない限り全ての値に対して一般的な定義を与えます。

定義
指数関数 $a^{\zeta}$ を以下の式によって定義します：
\[
a^{\zeta} = \exp (\zeta \Log a)
\]
ここで $\Log a$ は $a$ の対数の任意の値です。

この定義が以前の定義と矛盾せず、それらを全て含むことを確認する必要があります。

(1) $a$ が正で $\zeta$ が実数の場合、$\zeta \Log a$ の一つの値は実数であり、すなわち $\zeta \log a$ です。そして $\exp (\zeta \log a) = e^{\zeta \log a}$ となり、これは第 IX 章で採用した定義と一致します。当時、第 IX 章の定義は初等代数での定義と一致していましたから、新しい定義も同様です。

(2) $a = e^{\tau} (\cos \psi + i \sin \psi)$ の場合、
\begin{gather*}
\Log a = \tau + i(\psi + 2m\pi), \\
\exp \{(p/q)\Log a\} = e^{p\tau/q} \Cis \{(p/q)(\psi + 2m\pi)\},
\end{gather*}
ここで $m$ は任意の整数値を持てます。$m$ が全ての可能な整数値を取ると、この式は $q$ 個だけ異なる値を取り、これらは第 48 節で見つけた $a^{p/q}$ の値と一致します。したがって、新しい定義も第 III 章の定義と一致しています。

226. 一般の $a^{\zeta}$ の値について
\[
\zeta = \xi + i\eta,\quad a = \sigma(\cos\psi + i\sin\psi)
\]
ここで $\psi$ は $-\pi < \psi \leq \pi$ を満たすようにします。この表記法では、$\sigma = e^{\tau}$ または $\tau = \log \sigma$ です。

次に
\[
\zeta \Log a = (\xi + i\eta)\{\log \sigma + i(\psi + 2m\pi)\} = L + iM,
\]
で定義されるようにします。ここで
\[
L = \xi \log \sigma - \eta(\psi + 2m\pi),\quad
M = \eta\log \sigma + \xi (\psi + 2m\pi);
\]
そして
\[
a^{\zeta} = \exp(\zeta\Log a) = e^{L}(\cos M + i\sin M).
\]
したがって、$a^{\zeta}$ の一般値は
\[
e^{\xi\log \sigma - \eta(\psi+2m\pi)}
    [\cos\{\eta\log \sigma + \xi(\psi + 2m\pi)\}
  + i\sin\{\eta\log \sigma + \xi(\psi + 2m\pi)\}]
\]
となります。

一般に $a^{\zeta}$ は無数の値を持つ関数です。実際には
\[
|a^{\zeta}| = e^{\xi\log \sigma - \eta(\psi+2m\pi)}
\]
が、異なる $m$ の値に対して異なる値を取るためです。ただし $\eta = 0$ の場合、すべての異なる値の絶対値は同じになります。しかし、どの二つの値も、その振幅が同じであるか、または $2\pi$ の倍数だけ異なる場合にのみ等しくなります。これは $\xi(\psi + 2m\pi)$ と $\xi(\psi + 2n\pi)$（$m$ と $n$ は異なる整数）が、もし異なるなら $2\pi$ の倍数だけ異なることを意味します。したがって
\[
\xi(\psi + 2m\pi) - \xi(\psi + 2n\pi) = 2k\pi,
\]
が成り立つとき、$\xi = k/(m - n)$ は有理数です。結論として、$a^{\zeta}$ は $\zeta$ が実数かつ有理数でない限り無数の値を持つことになります。一方で、$\zeta$ が実数かつ有理数であるとき、$a^{\zeta}$ は有限個の値を持つことが既に示されています。

補足 Remark: $a^{\zeta} = \exp (\zeta\Log a)$ の主値は、$\Log a$ の主値を与えることにより得られます。すなわち、一般式において $m = 0$ と仮定することで得られます。したがって、$a^{\zeta}$ の主値は
\[
e^{\xi\log \sigma - \eta\psi}
  \{\cos(\eta\log \sigma + \xi\psi) + i\sin(\eta\log \sigma + \xi\psi)\}
\]
となります。

二つの特別な場合が特に興味深い。もし \(a\) が実数で正であり、\(\zeta\) が実数であるとき、\(\sigma = a\), \(\psi = 0\), \(\xi = \zeta\), \(\eta = 0\) であり、主値の \(a^{\zeta}\) は \(e^{\zeta \log a}\) で、これは最後の章で定義された値である。もし \(|a| = 1\) で \(\zeta\) が実数であるとき、\(\sigma = 1\), \(\xi = \zeta\), \(\eta = 0\) であり、主値の \((\cos \psi + i \sin \psi)^{\zeta}\) は \(\cos \zeta \psi + i \sin \zeta \psi\) である。これはデ・モivre の定理のさらに一般的な形を表している（§45, §49）。

---

例題 XCIV.

1. \(i^{i}\) のすべての値を求めよ。[定義により
\[
i^{i} = \exp (i \log i).
\]
しかし
\[
i = \cos \frac{1}{2}\pi + i \sin \frac{1}{2}\pi,\quad
\log i = (2k + \frac{1}{2})\pi i,
\]
ここで \(k\) は任意の整数である。したがって
\[
i^{i} = \exp \left\{ -(2k + \frac{1}{2})\pi \right\} = e^{-(2k + \frac{1}{2})\pi}.
\]
\(i^{i}\) のすべての値は実数で正である。]

2. \( (1 + i)^{i}\), \(i^{1+i}\), \((1 + i)^{1+i}\) のすべての値を求めよ。

3. \(a^{\zeta}\) の値を Argand 幾何図に描くと、等角多角形の頂点となり、その多角形は等角螺線に内接し、角度は \(a\) によらず一定である。[1899年]

[もし \(a^{\zeta} = r(\cos \theta + i \sin \theta)\) であれば
\[
r = e^{\xi \log \sigma - \eta (\psi + 2m\pi)},\quad
\theta = \eta \log \sigma + \xi (\psi + 2m\pi);
\]
そしてすべての点は螺旋 \(r = \sigma^{(\xi^{2} + \eta^{2})/\xi} e^{-\eta \theta/\xi}\) 上に存在する。]

\item{4.} \Topic{関数 $e^{\zeta}$ について。} 一般公式において $a$ を $e$ と書き換えると、$\log \sigma = 1$, $\psi = 0$ となるので
\[
e^{\zeta} = e^{\xi - 2m\pi\eta} \{ \cos(\eta + 2m\pi\xi) + i\sin(\eta + 2m\pi\xi) \}
\]
を得る。関数 $e^{\zeta}$ の主値は $e^{\xi} (\cos\eta + i\sin\eta)$ であり、これは $\exp \zeta$ （§223）と等しい。
特に、$\zeta$ が実数であるとき、すなわち $\eta = 0$ の場合、
\[
e^{\zeta} (\cos 2m\pi\zeta + i\sin 2m\pi\zeta)
\]
を得る。これは一般値であり、主値は $e^{\zeta}$ と等しい。ここで $e^{\zeta}$ は指数関数の正の値を表す（第9章参照）。

\item{5.} $\Log e^{\zeta} = (1 + 2m\pi i)\zeta + 2n\pi i$ を示し、ここで $m$ と $n$ は任意の整数である。また、一般には $\Log a^{\zeta}$ は双無限個の値を持つ。

\item{6.} 方程式 $1/a^{\zeta} = a^{-\zeta}$ は完全に成り立つ（例 xciii.3）。主値においても成り立つ。

\item{7.} 方程式 $a^{\zeta} \times b^{\zeta} = (ab)^{\zeta}$ は完全に成り立つが、主値においては必ずしも成り立たない。

\item{8.} 方程式 $a^{\zeta} \times a^{\zeta'} = a^{\zeta + \zeta'}$ は完全には成り立たないが、主値においては成り立つ。右辺のすべての値は左辺の値であるが、一般値 $a^{\zeta} \times a^{\zeta'}$ のうち
\[
\exp \{ \zeta (\log a + 2m\pi i) + \zeta' (\log a + 2n\pi i) \}
\]
は通常、$a^{\zeta + \zeta'}$ の値ではない（$m = n$ の場合を除く）。

\item{9.} 方程式
\[
\Log a^{\zeta} = \zeta \Log a, \quad (a^{\zeta})^{\zeta'} = (a^{\zeta'})^{\zeta} = a^{\zeta \zeta'}
\]
についての対応する結果は何か？

\item{10.} どのような値の $\zeta$ に対して (\ia) 任意の値、(\ib) 主値であるかを次のように分類せよ。 (i) 実数 (ii) 純虚数 (iii) 単位長度？

\item 11. $a^{\zeta}$ のすべての値が実数である必要十分条件は、$2\xi$ と $\{\eta \log |a| + \xi \am a\} / \pi$ がともに整数であることである。これらの値がすべて単位長さを持つための対応する条件は何であるか？

\item 12. $x > 0$ のとき、$|x^{i} + x^{-i}|$ の一般値は
\[
e^{-(m-n)\pi} \sqrt{2\{\cosh 2(m + n)\pi + \cos(2\log x)\}}
\]
である。

\item 13. 次の議論における誤りを説明せよ：$e^{2m\pi i} = e^{2n\pi i} = 1$ であり、$m$ と $n$ は任意の整数であるから、両辺を $i$ 乗すると $e^{-2m\pi} = e^{-2n\pi}$ となる。

\item 14. $x$ が実数であるとき、$x^{x}$ のどの値が自身で実数であるか？ [もし $x > 0$ なら
\[
x^{x} = \exp (x \Log x) = \exp (x \log x) \Cis 2m\pi x,
\]
最初の因子は実数である。主値、つまり $m = 0$ のとき、常に実数である。

もし $x$ が有理分数 $p/(2q + 1)$ または無理性であるなら、他の実数値はありません。しかし、$x$ が形 $p/2q$ のとき、一つの他の実数値があり、それが $m = q$ で与えられる $-\exp (x \log x)$ である。

もし $x = -\xi < 0$ なら
\[
x^{x} = \exp \{-\xi \Log (-\xi)\}
      = \exp (-\xi \log \xi) \Cis \{-(2m + 1)\pi \xi\}.
\]
どの値も実数である唯一のケースは $\xi = p/(2q + 1)$ のときであり、$m = q$ で与えられる実数値は
\[
\exp (-\xi \log \xi) \Cis (-p\pi) = (-1)^{p} \xi^{-\xi}
\]
である。

実数性の例としては次のものが示される：
\[
(\tfrac{1}{3})^{1/3} = \sqrt{\tfrac{1}{3}},\quad
(\tfrac{1}{2})^{\frac{1}{2}} = ±\sqrt{\tfrac{1}{2}},\quad
(-\tfrac{2}{3})^{-\frac{2}{3}} = \sqrt{\tfrac{9}{4}},\quad
(-\tfrac{1}{3})^{-\frac{1}{3}} = -\sqrt{3}.
\]

\Item{15.} \Topic{対数の任意底について。} $\zeta = \Log_{a} z$ は二つの異なる方法で定義できます。一つには（i）$\zeta = \Log_{a} z$ が主値である $a^{\zeta}$ の値が $z$ と等しいとき；二つには（ii）$\zeta = \Log_{a} z$ が $\zeta$ の任意の値における $a^{\zeta}$ の値が $z$ と等しいときとして定義できます。

例えば、$a = e$ の場合、最初の定義によれば $\zeta = \Log_{e} z$ は主値である $e^{\zeta}$ の値が $z$ と等しいときであり、または $\exp \zeta = z$ のときです。したがって $\Log_{e} z$ は $\Log z$ に一致します。しかし、二番目の定義によれば $\zeta = \Log_{e} z$ は
\[
e^{\zeta} = \exp (\zeta \Log e) = z, \quad
\zeta \Log e = \Log z,
\]
または $\zeta = (\Log z)/(\Log e)$ の任意の値が取れるときです。したがって
\[
\zeta = \Log_{e} z = \frac{\log |z| + (\am z + 2m\pi)i}{1 + 2n\pi i},
\]
となり、$\zeta$ は $z$ の二重無限多価関数となります。そして一般に、この定義によれば $\Log_{a} z = (\Log z)/(\Log a)$ です。

\Item{16.} $\Log_{e} 1 = 2m\pi i/(1 + 2n\pi i)$, $\Log_{e}(-1) = (2m + 1)\pi i/(1 + 2n\pi i)$ で、$m$ と $n$ は任意の整数です。

\Paragraph{227. 指数形式の正弦と余弦。}
公式
\[
\exp (\xi + i\eta) = \exp \xi(\cos\eta + i\sin\eta)
\]
から、多くの極めて重要な補助公式を導くことができます。$\xi = 0$ の場合、$\exp (i\eta) = \cos\eta + i\sin\eta$ を得ます；そして $\eta$ の符号を変えると、$\exp (-i\eta) = \cos\eta - i\sin\eta$ です。したがって
\begin{alignat*}{3}
\cos\eta \&= \&\&\tfrac{1}{2} \&\&\{\exp (i\eta) + \exp (-i\eta)\},\\
\sin\eta \&= -\&\&\tfrac{1}{2}i\&\&\{\exp (i\eta) - \exp (-i\eta)\}.
\end{alignat*}
もちろん、$\exp (i\eta)$ を用いて $\eta$ の任意の三角関数を導くことができます。

228. $\sin\zeta$および$\cos\zeta$の定義（全ての$\zeta$に対する）
前節で見たように、$\zeta$が実数であるとき、以下の式が成り立ちます。
\begin{alignat*}{3}
\cos\zeta \&= \&\&\tfrac{1}{2} \&\&\{\exp (i\zeta) + \exp (-i\zeta)\},\Tag{(1a)}\\
\sin\zeta \&= -\&\&\tfrac{1}{2}i\&\&\{\exp (i\zeta) - \exp (-i\zeta)\}.
\Tag{(1b)}
\end{alignat*}

これらの式の左辺は、初等三角法で採用された通常の幾何学的定義により、$\zeta$が実数である場合にのみ定義されています。一方、右辺は全ての$\zeta$（実数または複素数）に対して定義されています。したがって、これらの公式を$\cos \zeta$および$\sin \zeta$の全ての$\zeta$に対する定義として採用するのが自然です。これらの定義は、$\zeta$が実数である場合の初等的な定義と一致します。

$\cos \zeta$および$\sin \zeta$を定義した後で、他の三角関数を以下の式により定義します。
\[
\tan \zeta = \frac{\sin \zeta}{\cos \zeta},\quad
\cot \zeta = \frac{\cos \zeta}{\sin \zeta},\quad
\sec \zeta = \frac{1}{\cos \zeta},\quad
\cosec \zeta = \frac{1}{\sin \zeta}.
\Tag{(2)}
\]
$\cos \zeta$および$\sec \zeta$は偶関数であり、$\sin \zeta$、$\tan \zeta$、$\cot \zeta$、および$\cosec \zeta$は奇関数です。また、$\exp (i\zeta) = t$とすると、以下の式が成り立ちます。
\begin{gather*}
\cos \zeta =  \tfrac{1}{2}  \{t + (1/t)\},\quad
\sin \zeta = -\tfrac{1}{2}i \{t - (1/t)\},\\
\cos^{2} \zeta + \sin^{2} \zeta
  = \tfrac{1}{4}[\{t + (1/t)\}^{2} - \{t - (1/t)\}^{2}] = 1.
\Tag{(3)}
\end{gather*}

We さらに三角関数 $\zeta + \zeta'$ を $\zeta$ および $\zeta'$ の三角関数で表すことは、初等三角法における公式と全く同じ形で可能である。実際、$\exp (i\zeta) = t$, $\exp (i\zeta') = t'$ とするとき、

\begin{align*}
%[** TN: Set on two lines in the original, not aligned]
\cos (\zeta + \zeta')
  &= \tfrac{1}{2} \left(tt' + \frac{1}{tt'}\right) \\
  &= \tfrac{1}{4} \left\{
      \left(t + \frac{1}{t}\right) \left(t' + \frac{1}{t'}\right)
    + \left(t - \frac{1}{t}\right) \left(t' - \frac{1}{t'}\right)\right\}\\
  &= \cos\zeta \cos\zeta' - \sin\zeta \sin\zeta';
\Tag{(4)}
\end{align*}

と同様に、

\[
\sin (\zeta + \zeta') = \sin\zeta \cos\zeta' + \cos\zeta \sin\zeta'.
\Tag{(5)}
\]

特に、

\[
\cos(\zeta + \tfrac{1}{2}\pi) = -\sin\zeta,\quad
\sin(\zeta + \tfrac{1}{2}\pi) =  \cos\zeta.
\Tag{(6)}
\]

これらの普通の三角関数の公式は、式~\Eq{(2)}--\Eq{(6)} の代数的推論であり、したがって、この節で定義された一般化された三角関数についても同様の関係が成り立つ。

\par\medskip\noindent\textbf{【注釈】}\par
\Paragraph{229. 一般化された双曲線関数。} \Ex{lxxxvii}.19 では、実数 $\zeta$ の値に対して $\cosh \zeta$ および $\sinh \zeta$ を式

\[
\cosh\zeta = \tfrac{1}{2} \{\exp \zeta + \exp (-\zeta)\},\quad
\sinh\zeta = \tfrac{1}{2} \{\exp \zeta - \exp (-\zeta)\}.
\Tag{(1)}
\]

で定義した。

現在、変数の複素値に対してこの定義を拡張することができる；すなわち、式~\Eq{(1)} はすべての実数または複素数 $\zeta$ の値に対して $\cosh \zeta$ および $\sinh \zeta$ を定めるものとする。読者は以下の関係を容易に確認できる：

\[
\cos i\zeta =  \cosh \zeta,\quad
\sin i\zeta = i\sinh \zeta,\quad
\cosh i\zeta =  \cos \zeta,\quad
\sinh i\zeta = i\sin \zeta.
\]

任意の基本的な三角関数公式、例えば
\[
\cos 2\zeta = \cos^{2}\zeta - \sin^{2}\zeta
\]
は、$\zeta$が複素数値を取ることも許されるときでも真です。したがって、$\cos i\zeta$ を $\cos \zeta$, $\sin i\zeta$ を $\sin \zeta$, そして $\cos 2i\zeta$ を $\cos 2\zeta$ と書き換えることも可能であり、これは $\cosh \zeta$ を $\cos \zeta$, $i\sinh \zeta$ を $\sin \zeta$, そして $\cosh 2\zeta$ を $\cos 2\zeta$ と書き換えることに相当します。従って
\[
\cosh 2\zeta = \cosh^{2}\zeta + \sinh^{2}\zeta.
\]
このような変換は任意の三角関数恒等式にも適用できます。これは、双曲線関数と通常の三角関数の公式との間の対応が \Ex{lxxxvii}. 21 で述べられていることに説明されます。

230. $\cos(\xi + i\eta)$, $\sin(\xi + i\eta)$ 等の公式
加法定理から、以下の公式を得ます。
\[
\begin{alignat*}{4}
\cos (\xi + i\eta)
  &= \cos\xi \cos i\eta &&- \sin\xi \sin i\eta
  &&= \cos\xi \cosh \eta &&- i\sin\xi \sinh \eta,\\
\sin (\xi + i\eta)
  &= \sin\xi \cos i\eta &&+ \cos\xi \sin i\eta
  &&= \sin\xi \cosh \eta &&+ i\cos\xi \sinh \eta.
\end{alignat*}
\]
これらの公式はすべての $\xi$ と $\eta$ の値に対して成り立ちます。興味深いのは、$\xi$ と $\eta$ が実数である場合です。このとき、複素数の余弦および正弦の実部と虚部を表す式を得ることができます。

1. $\cos\zeta$ と $\sin\zeta$ の値を求めよ。
   (i) 実数
   (ii) 純粹な虚数。例えば、$\cos\zeta$ は $\eta = 0$ や $\xi$ が $\pi$ の何倍でも実数となる。

2.
\[
\begin{alignat*}{2}
|\cos (\xi + i\eta)|
  &= \sqrt{\cos^{2}\xi + \sinh^{2}\eta}
  &&= \sqrt{\frac{1}{2} (\cosh 2\eta + \cos 2\xi)}, \\
|\sin (\xi + i\eta)|
  &= \sqrt{\sin^{2}\xi + \sinh^{2}\eta}
  &&= \sqrt{\frac{1}{2} (\cosh 2\eta - \cos 2\xi)}.
\end{alignat*}
\]

[使用する式は、$|\cos(\xi + i\eta)| = \sqrt{\cos(\xi + i\eta) \cos(\xi - i\eta)}$ である。]

3.
$\tan (\xi + i \eta)
  = \dfrac{\sin 2\xi + i\sinh 2\eta}{\cosh 2\eta + \cos 2\xi}$,\quad
$\cot (\xi + i \eta)
  = \dfrac{\sin 2\xi - i\sinh 2\eta}{\cosh 2\eta - \cos 2\xi}$.

[例えば、
\[
\tan (\xi + i\eta)
  = \frac{\sin (\xi + i\eta) \cos (\xi - i\eta)}
         {\cos (\xi + i\eta) \cos (\xi - i\eta)}
  = \frac{\sin 2\xi + \sin 2i\eta}{\cos 2\xi + \cos 2i\eta},
\]
これにより、与えられた結果が得られる。]

4.
\begin{align*}
\sec (\xi + i \eta)
  &= \frac{\cos\xi \cosh\eta + i\sin\xi \sinh\eta}
          {\frac{1}{2} (\cosh 2\eta + \cos 2\xi)}, \\
\cosec (\xi + i \eta)
  &= \frac{\sin\xi \cosh\eta - i\cos\xi \sinh\eta}
          {\frac{1}{2} (\cosh 2\eta - \cos 2\xi)}.
\end{align*}

5. $|\cos (\xi + i\eta)| = 1$ のとき、$\sin^{2} \xi = \sinh^{2} \eta$ であり、$|\sin (\xi + i\eta)| = 1$ のとき、$\cos^{2} \xi = \sinh^{2} \eta$。

6. $|\cos (\xi + i\eta)| = 1$ のとき、
\[
\sin \{\am \cos (\xi + i\eta)\} = ±\sin^{2} \xi =  ±\sinh^{2} \eta.
\]

7. $\Log \cos (\xi + i\eta) = A + iB$ を証明し、ここで
\[
A = \tfrac{1}{2} \log \{\tfrac{1}{2} (\cosh 2\eta + \cos 2\xi)\}
\]
であり、$B$ は次の条件を満たす任意の角である：
\[
\frac{\cos B}{\cos\xi \cosh\eta}
  = -\frac{\sin B}{\sin\xi \sinh\eta}
  = \frac{1}{\sqrt{\frac{1}{2} (\cosh 2\eta + \cos 2\xi)}}.
\]
$\Log \sin (\xi + i\eta)$ の類似公式も求めよ。

\Item{8.} \Topic{方程式 $\cos\zeta = a$ の解、ここで $a$ は実数。} $\zeta = \xi + i\eta$ と置き、実部と虚部を比較すると
\[
\cos\xi \cosh\eta = a,\quad
\sin\xi \sinh\eta = 0.
\]
したがって、$\eta = 0$ または $\xi$ は $\pi$ の倍数である。もし (i) $\eta = 0$ ならば $\cos\xi = a$ となり、これは $-1 \leq a \leq 1$ のときのみ可能である。この仮定から解は
\[
\zeta = 2k\pi ± \arccos a,
\]
ここで $\arccos a$ は $0$ と $\frac{1}{2}\pi$ の間にある。もし (ii) $\xi = m\pi$ ならば $\cosh\eta = (-1)^m a$ となり、これは $a \geq 1$ で $m$ が偶数、または $a \leq -1$ で $m$ が奇数である。もし $a = ±1$ ならば $\eta = 0$ となり、最初のケースに戻る。もし $|a| > 1$ ならば $\cosh\eta = |a|$ となり、解は
\begin{alignat*}{4}
\zeta \&=\& 2k      \&\pi ± i\log \{ \&\&a + \sqrt{a^{2} - 1}\}\quad \&\&(a > 1), \\
\zeta \&=\&(2k + 1) \&\pi ± i\log \{-\&\&a + \sqrt{a^{2} - 1}\}\quad \&\&(a < -1).
\end{alignat*}
例えば、$\cos\zeta = -\frac{5}{3}$ の一般解は $\zeta = (2k + 1)\pi ± i\log 3$ である。
\PageSep{412}

\Item{9.} $\sin\zeta = \alpha$ を解く、ここで $\alpha$ は実数。

\Item{10.} \Topic{$\cos\zeta = \alpha + i\beta$ の解、ここで $\beta \neq 0$。} $\beta > 0$ と仮定してもよい、$\beta < 0$ の結果は $i$ の符号を変えることで導くことができる。この場合
\[
\cos\xi \cosh\eta = \alpha,\quad
\sin\xi \sinh\eta = -\beta,
\Tag{(1)}
\]
および
\[
(\alpha/\cosh\eta)^{2} + (\beta/\sinh\eta)^{2} = 1.
\]

如果我達令 $\cosh^{2} \eta = x$，則有
\[
x^{2} - (1 + \alpha^{2} + \beta^{2})x + \alpha^{2} = 0
\]
または $x = (A_{1} ± A_{2})^{2}$, ここで
\[
A_{1} = \frac{1}{2}\sqrt{(\alpha + 1)^{2} + \beta^{2}},\quad
A_{2} = \frac{1}{2}\sqrt{(\alpha - 1)^{2} + \beta^{2}}.
\]
$\alpha > 0$ の仮定のもとで、$A_{1} > A_{2} > 0$ であり、$\cosh\eta = A_{1} ± A_{2}$。また
\[
\cos\xi = \frac{\alpha}{\cosh\eta} = A_{1} \mp A_{2},
\]
そして $\cosh\eta > \cos\xi$ より、$\cosh\eta = A_{1} + A_{2}$ および $\cos\xi = A_{1} - A_{2}$ を取る。これらの方程式の一般解は
\[
\xi = 2k\pi ± \arccos M,\quad
\eta = ±\log \{L + \sqrt{L^{2} - 1}\}
\]
であり、ここで $L = A_{1} + A_{2}$, $M = A_{1} - A_{2}$ で $\arccos M$ は $0$ と $\frac{1}{2}\pi$ の間にある。

しかしながら、上記で求めた $\eta$ および $\xi$ の値は、方程式
\[
\cos\xi \cosh\eta = \alpha,\quad
\sin\xi \sinh\eta = \beta
\]
の解を含むだけでなく、方程式~(1) の解も含んでいる。これは後者の方程式の平方根を用いたためである。二つの解の区別のために、$\sin\xi$ の符号は方程式~(2) の最初の方程式の不確定な符号と一致し、$\sinh\eta$ の符号は方程式~(2) の第二方程式の不確定な符号と一致する。そして $\beta > 0$ より、これらの二つの符号は異なる。したがって必要な一般解は
\[
\zeta = 2k\pi ± [\arccos M - i\log \{L + \sqrt{L^{2} - 1}\}]
\]
である。

11. $\alpha < 0$ 及び $\alpha = 0$ の場合について同様に計算せよ。

12. $\beta = 0$ のとき、$L = \frac{1}{2}|\alpha + 1| + \frac{1}{2}|\alpha - 1|$ および $M = \frac{1}{2}|\alpha + 1| - \frac{1}{2}|\alpha - 1|$。得られた結果が Ex.~8 の結果と一致することを確認せよ。

13. 正の数αとβに対して、方程式sinζ = α + iβの一般解は以下の式で表される：
\[
ζ = kπ + (-1)^k [\arcsin M + i \log \{L + \sqrt{L^2 - 1}\}],
\]
ここで0から1/2πまでの範囲に収まるarcsin Mがある。他の可能な場合の解も求めよ。

14. tanζ = αを解く（αは実数）。すべての根は実数である。

15. β ≠ 0の場合、方程式tan ζ = α + iβの一般解は以下の式で表される：
\[
ζ = kπ + \frac{1}{2}θ + \frac{1}{4}i \log \left\{
  \frac{α^2 + (1 + β)^2}
       {α^2 + (1 - β)^2}
  \right\},
\]
ここでθは以下の比が最も小さい数値の角度である：
\[
\cos θ : \sin θ : 1 ::
1 - α^2 - β^2 : 2α :
  \sqrt{(1 - α^2 - β^2)^2 + 4α^2}.
\]

16. z = ξexp(1/4πi)（ξは実数）かつcも実数の場合、cos2πz - cos2πcの絶対値は以下の式で表される：
\[
\surd \left[\frac{1}{2} \{1 + \cos 4πc + \cos(2πξ√2) + \cosh(2πξ√2) - 4\cos 2πc \cos(πξ√2) \cosh(πξ√2)\}\right].
\]

17. 次の式を証明せよ：
\[
|\exp \exp(\xi + iη)| = \exp(\exp\xi \cos η),
\]
また
\[
\begin{aligned}
\Re \{\cos \cos(\xi + iη)\} &= \cos(\cos\xi \coshη) \cosh(\sin\xi \sinhη),\\
\Im \{\sin \sin(\xi + iη)\} &= \cos(\sin\xi \coshη) \sinh(\cos\xi \sinhη).
\end{aligned}
\]

18. ζが原点を通る直線沿いに移動し、その直線とOX軸との角度が1/2π未満の場合、|expζ|は無限大へと収束し、角度が1/2πより大きい場合、|expζ|は0へと収束する。

19. ゾテ（ζ）が原点を通る任意の直線を動かして無限大に向かうとき、|\cos\zeta| と |\sin\zeta| は無限大に収束する。

20. ゾテ（ζ）が19問の例で述べた直線を動かして無限大に向かうとき、その直線が実軸より上にある場合は \tan\zeta は -i に収束し、下にある場合は i に収束する。

231. ログ関数と逆三角関数の関連性について。第VI章で述べたように、定数α, β, …を含む有理または代数的関数φ(x, α, β, …) の積分は、α, β, …の値によって異なる形を取りうる。時としてそれは対数で表され、時として逆三角関数で表される。例えば、

\[
\int \frac{dx}{x^{2} + \alpha}
  = \frac{1}{\sqrt{\alpha}} \arctan \frac{x}{\sqrt{\alpha}}
\Tag{(1)}
\]

がα > 0 のときであり、また

\[
\int \frac{dx}{x^{2} + \alpha}
  = \frac{1}{2\sqrtp{-\alpha}}
      \log \left|\frac{x - \sqrtp{-\alpha}}{x + \sqrtp{-\alpha}}\right|
\Tag{(2)}
\]

がα < 0 のときである。これらの事実は、対数関数と逆三角関数との間の機能的な関連性の存在を示唆する。そのような関連性があることはまた、円関数を \exp i\zeta で表すことができ、対数が指数関数の逆であるという事実からも推測できる。

方程式
\[
\int \frac{dx}{x^{2} - \alpha^{2}}
  = \frac{1}{2\alpha} \log \left(\frac{x - \alpha}{x + \alpha}\right)
\]
は、$\alpha$が実数かつ$(x - \alpha)/(x + \alpha)$が正のとき成り立ちます。この方程式において$\alpha$を$i\alpha$に書き換えることができれば、以下の式を得ます。
\[
\arctan \left(\frac{x}{\alpha}\right)
  = \frac{1}{2i} \log\left(\frac{x - i\alpha}{x + i\alpha}\right) + C
\]
ここで$C$は定数です。複素数の対数を定義した今、この方程式が実際に成り立つかどうかが問われます。

現在（節番号221）
\[
\Log(x ± i\alpha) = \frac{1}{2} \log(x^{2} + \alpha^{2}) ± i(\phi + 2k\pi)
\]
ここで$k$は整数、$\phi$は$\cos\phi = x/\sqrt{x^{2} + \alpha^{2}}$かつ$\sin\phi = \alpha/\sqrt{x^{2} + \alpha^{2}}$を満たす数値的に最小の角です。したがって
\[
\frac{1}{2i} \Log\left(\frac{x - i\alpha}{x + i\alpha}\right) = -\phi - l\pi
\]
ここで$l$は整数であり、これは実際には$\arctan(x/\alpha)$の任意の値から定数を異なっています。

対数関数と逆三角関数との間の標準的な公式は以下の通りです。
\[
\arctan x = \frac{1}{2i} \Log\left(\frac{1 + ix}{1 - ix}\right)
\]
ここで$x$は実数です。この公式を$\tan y = x$と置き換えることで容易に検証できます。右辺は
\[
\frac{1}{2i} \Log\left(\frac{\cos y + i\sin y}{\cos y - i\sin y}\right)
  = \frac{1}{2i} \Log(\exp 2iy) = y + k\pi
\]
となり、$k$は任意の整数です。したがって方程式（4）は「完全に」成り立ちます（例 xciii.3）。読者は以下の公式も検証すべきです。
\[
\arccos x = -i \Log\{x ± i\sqrt{1 - x^{2}}\},\quad
\arcsin x = -i \Log\{ix ± \sqrt{1 - x^{2}}\}
\]
ここで$-1 \leq x \leq 1$：これらの公式も「完全に」成り立ちます。

例として、方程式
\[
\cos u = x = \frac{1}{2}\{y + (1/y)\},
\]
ここで $y = \exp(iu)$ に対して解くと、$y = x ± i\sqrt{1 - x^{2}}$ を得る。従って：
\[
u = -i \log y = -i \log\{x ± i\sqrt{1 - x^{2}}\},
\]
これは方程式 (5) の最初の式と等価である。残りの方程式 (4) および (5) を類似した論理で得ることができる。

232. 指数関数 $\exp z$ の冪級数．
§212 において、$z$ が実数のとき
\[
\exp z = 1 + z + \frac{z^{2}}{2!} + \dots.
\]
Tag (1)
と見た。さらに §191 においては、右辺の級数が $z$ が複素数である場合でも収束（実際には絶対収束）することも確認した。従って方程式 (1) も同様に成り立つことが自然に推測され、この事実を証明する。

級数 (1) の和を $F(z)$ と表す。級数が絶対収束していることから、直接乗算（例 lxxxi.7 参照）により $F(z)$ は関数方程式
\[
F(z) F(h) = F(z + h).
\]
Tag (2)
を満たすことがわかる。ここで $z = iy$ とし、$F(z) = f(y)$ とする。すると
\[
f(y) f(k) = f(y + k);
\]
であり、従って
\[
\frac{f(y + k) - f(y)}{k} = f(y) \left\{\frac{f(k) - 1}{k}\right\}.
\]

ここで
\[
\frac{f(k) - 1}{k} = i \left\{1 + \frac{ik}{2!} + \frac{(ik)^{2}}{3!} + \dots\right\};
\]
であり、従って $|k| < 1$ のとき
\[
\left|\frac{f(k) - 1}{k} - i\right| < \left(\frac{1}{2!} + \frac{1}{3!} + \dots\right)|k| < (e - 2)|k|.
\]
したがって $\frac{f(k) - 1}{k} \to i$ であり、従って
\[
f'(y) = \lim_{k \to 0} \frac{f(y + k) - f(y)}{k} = if(y).
\]
Tag (3)
となる。

現在
\[
f(y) = F(iy) = 1 + (iy) + \frac{(iy)^2}{2!} + \cdots
     = \phi(y) + i\psi(y),
\]
ここで $\phi(y)$ は偶関数、$\psi(y)$ は奇関数であり、従って
\[
|f(y)| = \sqrt{\{\phi(y)\}^2 + \{\psi(y)\}^2}
       = \sqrt{\{\phi(y) + i\psi(y)\}\{\phi(y) - i\psi(y)\}}
       = \sqrt{F(iy) F(-iy)} = \sqrt{F(0)} = 1;
\]
したがって
\[
f(y) = \cos Y + i \sin Y,
\]
ここで $Y$ は $-\pi < Y \leq \pi$ を満たす $y$ の関数である。また $f(y)$ は微分可能であり、実部と虚部 $\cos Y$ および $\sin Y$ も微分可能で、従って連続関数である。したがって $Y$ も $y$ の連続関数である。もし $y$ から $y + k$ に変化すると $Y$ が $Y + K$ に変化するとき、$K$ は $k$ とともに零に収束し
\[
\frac{K}{k}
  = \left\{\frac{\cos(Y + K) - \cos Y}{k}\right\} / \left\{\frac{\cos(Y + K) - \cos Y}{K}\right\}.
\]
右辺の2つの商のうち、最初のものは $k \to 0$ のとき零に収束し、$\cos Y$ が $y$ に関する微分可能であることから、第二項は $\lim -\sin Y$ に収束する。したがって $\frac{K}{k}$ も零に収束し、従って $Y$ は $y$ に関する微分可能である。

さらに
\[
f'(y) = (-\sin Y + i \cos Y) \frac{dY}{dy}.
\]
かつすでに確認したように
\[
f'(y) = if(y) = -\sin Y + i \cos Y.
\]
従って
\[
\frac{dY}{dy} = 1,\quad
Y = y + C,
\]
ここで $C$ は定数であり、
\[
f(y) = \cos(y + C) + i \sin(y + C).
\]

しかし $y = 0$ のとき $f(0) = 1$ より $C$ は $2\pi$ の倍数であり、従って $f(y) = \cos y + i \sin y$ である。したがってすべての実数 $y$ に対して
\[
F(iy) = \cos y + i \sin y.
\]
またもし $x$ も実数であるとき
\[
F(x + iy) = F(x) F(iy) = \exp x(\cos y + i \sin y) = \exp(x + iy),
\]
または
\[
\exp z = 1 + z + \frac{z^2}{2!} + \cdots,
\]
すべての $z$ の値に対して成立する。

233. コサインｚとサインｚの冪級数

前節の結果と第228節の式（1）から即座に知ることができるが、

\[
\cos z = 1 - \frac{z^{2}}{2!} + \frac{z^{4}}{4!} - \dots,\quad
\sin z = z - \frac{z^{3}}{3!} + \frac{z^{5}}{5!} - \dots
\]

がすべてのｚに対して成り立つ。これらの結果は第228節の例lvi.1で実数ｚに対する場合を証明している。

例XCVI.

1. コサインｉとサインｉを小数点以下二桁まで冪級数を使用して計算せよ。

2. \(|\cos z| \leq \cosh|z|\) および \(|\sin z| \leq \sinh|z|\) を証明せよ。

3. \(|z| < 1\) のとき、\(|\cos z| < 2\) および \(|\sin z| < \frac{6}{5}|z|\) を証明せよ。

4. \(\sin 2z = 2\sin z \cos z\) より、

\[
(2z) - \frac{(2z)^{3}}{3!} + \frac{(2z)^{5}}{5!} - \dots
  = 2\left(z - \frac{z^{3}}{3!} + \dots\right)
     \left(1 - \frac{z^{2}}{2!} + \dots\right).
\]

右辺の二つの級数を乗算（第195節）し、係数を比較（第194節）することにより、

\[
\binom{2n + 1}{1} + \binom{2n + 1}{3} + \dots + \binom{2n + 1}{2n + 1} = 2^{2n}.
\]

を証明せよ。二項定理を使用して結果を検証せよ。同様の恒等式は、

\[
\cos^{2}z + \sin^{2}z = 1,\quad
\cos2z = 2\cos^{2}z - 1 = 1 - 2\sin^{2}z.
\]

から導き出せ。

5. 次を示せ。

\[
\exp\{(1 + i)z\}
  = \sum_{0}^{\infty} 2^{\frac{1}{2}n} \exp(\tfrac{1}{4}n\pi i) \frac{z^{n}}{n!}.
\]

\item{6.} $\cos z \cosh z$ を $z$ のべき級数に展開する。[私たちは以下のようになります：
\begin{align*}
\cos z \cosh z + i \sin z \sinh z 
  &= \cos\{(1 - i)z\} 
   = \frac{1}{2} [\exp\{(1 + i)z\} + \exp\{-(1 + i)z\}] \\
  &= \frac{1}{2} \sum_{n=0}^{\infty} 2^{\frac{1}{2}n}
       \{1 + (-1)^{n}\} \exp(\frac{1}{4}n\pi i) \frac{z^{n}}{n!},
\end{align*}
同様に
\[
\cos z \cosh z - i \sin z \sinh z = \cos (1 + i)z 
  = \frac{1}{2} \sum_{n=0}^{\infty} 2^{\frac{1}{2}n}
      \{1 + (-1)^{n}\} \exp(-\frac{1}{4}n\pi i) \frac{z^{n}}{n!}.
\]
したがって
\[
\cos z \cosh z 
  = \frac{1}{2} \sum_{n=0}^{\infty} 2^{\frac{1}{2}n}\{1 + (-1)^{n}\} \cos \frac{1}{4}n\pi \frac{z^{n}}{n!}
  = 1 - \frac{2^{2}z^{4}}{4!} + \frac{2^{4}z^{8}}{8!} - \dots.
\]

\item{7.} $\sin z \sinh z$, $\cos z \sinh z$, および $\sin z \cosh z$ を $z$ のべき級数に展開する。

\item{8.} $\sin^{2} z$ および $\sin^{3} z$ を $z$ のべき級数に展開する。[以下の公式を使用します：
\[
\sin^{2} z = \frac{1}{2} (1 - \cos 2z),\quad
\sin^{3} z = \frac{1}{4} (3\sin z - \sin 3z),\ \dots.
\]
明らかに、同じ方法で $\cos^{n} z$ および $\sin^{n} z$ を展開することができます（$n$ は任意の整数）.]

\item{9.} シリーズを求める：
\[
C = 1 + \frac{\cos z}{1!} + \frac{\cos 2z}{2!} + \frac{\cos 3z}{3!} + \dots,\quad
S = \frac{\sin z}{1!} + \frac{\sin 2z}{2!} + \frac{\sin 3z}{3!} + \dots.
\]

[ここで
\begin{align*}
C + iS &= 1 + \frac{\exp(iz)}{1!} + \frac{\exp(2iz)}{2!} + \dots 
   = \exp\{\exp(iz)\} \\
  &= \exp(\cos z) \{\cos(\sin z) + i\sin(\sin z)\},
\end{align*}
同様に
\[
C - iS = \exp\{\exp(-iz)\} = \exp(\cos z)\{\cos(\sin z) - i\sin(\sin z)\}.
\]
したがって
\[
C = \exp(\cos z)\cos(\sin z),\quad
S = \exp(\cos z)\sin(\sin z).
\]

\item{10.} 次のものを求める：
\[
1 + \frac{a\cos z}{1!} + \frac{a^{2}\cos 2z}{2!} + \dots,\quad
\frac{a\sin z}{1!} + \frac{a^{2}\sin 2z}{2!} + \dots.
\]

\item{11.} 次のものを求める：
\[
1 - \frac{\cos 2z}{2!} + \frac{\cos 4z}{4!} - \dots,\quad
\frac{\cos z}{1!} - \frac{\cos 3z}{3!} + \dots
\]
およびそれに対応する正弦関数を含む級数。

\item{12.} 1+\frac{\cos 4z}{4!} + \frac{\cos 8z}{8!} + \dots = \frac{1}{2}\{\cos(\cos z) \cosh(\sin z) + \cos(\sin z) \cosh(\cos z)\} を示せ。

\item{13.} \cos(x+h) 及び \sin(x+h) の h に関する冪級数展開（\Ex{lvi}.1）は、x 及び h が実数または複素数であるすべての値に対して有効であることを示せ。

\Paragraph{234. ログarithmic series.} \SecNo[§]{213} で得た結果より
\[
\log(1 + z) = z - \frac{1}{2}z^{2} + \frac{1}{3}z^{3} - \dots
\Tag{(1)}
\]
が、z が実数かつ絶対値が 1 より小さいときには成り立つ。右辺の級数は z の絶対値が 1 より小さい任意の複素数に対して収束し、さらに絶対収束する。このような複素数値の z に対する方程式 \Eq{(1)} が真であることが自然に推測される。この真実性を証明するために、\SecNo[§]{213} の議論を一部修正して示すことができる。実際には、|z| ≤ 1 であり -1 以外のすべての z 値に対して \Eq{(1)} が成り立つことを証明する。

\log(1 + z) は \Log(1 + z) の主値であり、
\[
\log(1 + z) = \int_{C} \frac{du}{u}
\]
ここで C は複素数変数 u 平面上の点 1 と 1+z を結ぶ直線である。z が実数でないことを仮定してもよい、既に実数値 z の場合について式 \Eq{(1)} の証明がなされているからである。

もし
\[
z = r(\cos\theta + i\sin\theta) = \zeta r,
\]
とし、$|r| \leq 1$ であり、また
\[
u = 1 + \zeta t,
\]
とすると、$t$ が $0$ から $r$ に増加するとき、$u$ は $C$ を描きます。そして
\begin{align*}
\int_{C} \frac{du}{u}
  &= \int_{0}^{r} \frac{\zeta\, dt}{1 + \zeta t} \\
  &= \int_{0}^{r} \left\{\zeta - \zeta^{2} t + \zeta^{3} t^{2} - \dots
     + (-1)^{m-1} \zeta^{m} t^{m-1}
     + \frac{(-1)^{m} \zeta^{m+1} t^{m}}{1 + \zeta t}\right\} dt \\
  &= \zeta r - \frac{(\zeta r)^{2}}{2} + \frac{(\zeta r)^{3}}{3} - \dots
     + (-1)^{m-1} \frac{(\zeta r)^{m}}{m} + R_{m} \\
  &= z - \frac{z^{2}}{2} + \frac{z^{3}}{3} - \dots
     + (-1)^{m-1} \frac{z^{m}}{m} + R_{m},
\Tag{(2)}
\end{align*}
ここで
\[
R_{m} = (-1)^{m} \zeta^{m+1} \int_{0}^{r} \frac{t^{m}\, dt}{1 + \zeta t}.
\Tag{(3)}
\]

$\SecNo[§]{164}$ の $\Eq{(1)}$ から、
\[
|R_{m}| \leq \int_{0}^{r} \frac{t^{m}\, dt}{|1 + \zeta t|}.
\Tag{(4)}
\]
今、$|1 + \zeta t|$ や $|u|$ は直線 $C$ 上の点から原点 $O$ への垂線の長さ $\varpi$ よりも決して小さくない。\footnote
  {実数でない $z$ のため、$C$ は延長しても $O$ を通らない。読者はこの議論を図で示すことを推奨する。}
したがって
\[
|R_{m}| \leq \frac{1}{\varpi} \int_{0}^{r} t^{m}\, dt
  = \frac{r^{m+1}}{(m + 1) \varpi}
  \leq \frac{1}{(m + 1) \varpi},
\]
\PageSep{419}
よって $R_{m} \to 0$ が $m \to \infty$ のとき成立します。$\Eq{(2)}$ から
\[
\log(1 + z) = z - \tfrac{1}{2} z^{2} + \tfrac{1}{3} z^{3} - \dots.
\Tag{(5)}
\]

もちろん、証明の過程で級数が収束することを示しました：これは既に $\Ex{lxxx}.~4$ で証明されています。実際、$\DPtypo{z|}{|z|} < 1$ のときには絶対収束し、$|z| = 1$ のときには条件付き収束します。

$z$ を $-z$ に置き換えると
\[
\log \left(\frac{1}{1 - z}\right)
  = -\log(1 - z)
  = z + \tfrac{1}{2} z^{2} + \tfrac{1}{3} z^{3} + \dots.
\Tag{(6)}
\]

\paragraph{235.} 現に
\begin{align*}
\log(1 + z)
  &= \log\{(1 + r \cos\theta) + ir\sin\theta\} \\
  &= \frac{1}{2} \log(1 + 2r\cos\theta + r^{2})
     + i\arctan \left(\frac{r\sin\theta}{1 + r\cos\theta}\right).
\end{align*}
この逆接弦関数の値は、$-\frac{1}{2}\pi$ と $\frac{1}{2}\pi$ の間に取らなければなりません。なぜなら、$1 + z$ は $-1$ から $z$ までのベクトルを表すため、$z$ が半径 $|z| = 1$ の円内にあるとき、主値 $\am(1 + z)$ はこれらの限界の間にあるからです。\footnote{前段落の脚注を参照してください。}

また、$z^{m} = r^{m}(\cos m\theta + i\sin m\theta)$ より、式 (5) の \SecNo[§]{234} における実部と虚部を比較すると
\begin{align*}
\frac{1}{2} \log(1 + 2r\cos\theta + r^{2})
  &= r\cos\theta - \frac{1}{2}r^{2} \cos 2\theta
                 + \frac{1}{3}r^{3} \cos 3\theta - \dots, \\
\arctan \left(\frac{r\sin\theta}{1 + r\cos\theta}\right)
  &= r\sin\theta - \frac{1}{2}r^{2} \sin 2\theta
                 + \frac{1}{3}r^{3} \sin 3\theta - \dots.
\end{align*}
これらの式は $0 \leq r \leq 1$ のとき、そしてすべての $\theta$ の値に対して成り立ちますが、$r = 1$ の場合、$\theta$ は $\pi$ の奇数倍には等しくならない必要があります。また、$-1 \leq r \leq 0$ のときも同様で、$r = -1$ の場合、$\theta$ は $\pi$ の偶数倍には等しくならない必要があります。

特に興味深いケースは \( r = 1 \) の場合である。この場合、次のようになる：
\[
\log(1 + z) = \log(1 + \Cis\theta)
  = \frac{1}{2} \log(2 + 2\cos\theta)
       + i\arctan\left(\frac{\sin\theta}{1 + \cos\theta}\right),
\]
ここで \( -\pi < \theta < \pi \) のとき、さらに
\[
\begin{alignat*}{4}
\cos\theta &- \frac{1}{2} \cos 2\theta &&+ \frac{1}{3} \cos 3\theta &&- \dots
  &&= \frac{1}{2} \log(4\cos^{2} \frac{1}{2}\theta), \\
\sin\theta &- \frac{1}{2} \sin 2\theta &&+ \frac{1}{3} \sin 3\theta &&- \dots
  &&= \frac{1}{2} \theta.
\end{alignat*}
\]
他の値の \( \theta \) に対する級数の和は、周期関数であることを考慮すると簡単に求められる。したがって、余弦級数の和は、すべての \( \theta \) の値に対して（奇数倍の \( \pi \) を除き）\(\frac{1}{2} \log(4\cos^{2} \frac{1}{2}\theta)\) であり、正弦級数の和は、\((2k - 1)\pi < \theta < (2k + 1)\pi\) のとき \(\frac{1}{2} (\theta - 2k\pi)\) で、\(\theta\) が奇数倍の \( \pi \) のとき零である。正弦級数によって表される関数のグラフは図58に示されている。この関数は \(\theta = (2k + 1)\pi\) において不連続である。
\begin{center}
\fbox{【図58】}
\end{center}

\par\medskip\noindent\textbf{【注釈】}\par
\( z \) を \( iz \) および \( -iz \) に書き換えて、引くと次のようになる：
\[
\frac{1}{2i} \log\left(\frac{1 + iz}{1 - iz}\right)
  = z - \frac{1}{3}z^{3} + \frac{1}{5}z^{5} - \dots.
\]
\( z \) が実数で、絶対値が1未満のとき、第231節の結果から次の公式を得る：
\[
\arctan z = z - \frac{1}{3}z^{3} + \frac{1}{5}z^{5} - \dots,
\]
これはすでに第214節で異なる方法で証明されている。
\par\medskip

\par\medskip\noindent\textbf{【例 XCVII.】}\par
\Item{1.} 任意の三角形において \( a > b \) のとき、次の式が成り立つことを証明せよ：
\[
\log c = \log a - \frac{b}{a} \cos C - \frac{b^{2}}{2a^{2}} \cos 2C - \dots.
\]

\[ \log c = \frac{1}{2} \log(a^{2} + b^{2} - 2ab\cos C ). \]

2. \( -1 < r < 1 \) および \( -\frac{1}{2}\pi < \theta < \frac{1}{2}\pi \) のとき、
\[ r\sin 2\theta - \frac{1}{2}r^{2} \sin 4\theta + \frac{1}{3}r^{3} \sin 6\theta - \cdots = \theta - \arctan \left\{\left(\frac{1 - r}{1 + r}\right) \tan\theta\right\}, \]
逆接線は \( -\frac{1}{2}\pi \) から \( \frac{1}{2}\pi \) の間にある。他のすべての \( \theta \) の値に対する級数の和を決定せよ。

3. \( \log(1 + iz) \) および \( \log(1 - iz) \) を \( z \) のべき乗展開することにより、\( -1 < r < 1 \) のとき、
\[
\begin{aligned}
r\sin\theta &+ \frac{1}{2}r^{2} \cos 2\theta - \frac{1}{3}r^{3} \sin 3\theta - \frac{1}{4}r^{4} \cos 4\theta + \cdots = \frac{1}{2} \log(1 + 2r \sin\theta + r^{2}), \\
r\cos\theta &+ \frac{1}{2}r^{2} \sin 2\theta - \frac{1}{3}r^{3} \cos 3\theta - \frac{1}{4}r^{4} \sin 4\theta + \cdots = \arctan \left(\frac{r\cos\theta}{1 - r\sin\theta}\right),
\end{aligned}
\]
逆接線は \( -\frac{1}{2}\pi \) から \( \frac{1}{2}\pi \) の間にある。また、
\[
\begin{aligned}
r\sin\theta &- \frac{1}{3}r^{3} \sin 3\theta + \cdots = \frac{1}{4} \log\left(\frac{1 + 2r \sin\theta + r^{2}}{1 - 2r \sin\theta + r^{2}}\right), \\
r\cos\theta &- \frac{1}{3}r^{3} \cos 3\theta + \cdots = \frac{1}{2} \arctan \left(\frac{2r\cos\theta}{1 - r^{2}}\right),
\end{aligned}
\]
逆接線は \( -\frac{1}{2}\pi \) から \( \frac{1}{2}\pi \) の間にある。

4. 以下の級数を証明せよ：
\begin{alignat*}{3}
\cos\theta \cos\theta \&- \tfrac{1}{2} \cos 2\theta \cos^{2}\theta
  \&\&+ \tfrac{1}{3} \cos 3\theta \cos^{3} \theta - \dots
  \&\&= \tfrac{1}{2} \log(1 + 3\cos^{2} \theta),\\
\sin\theta \sin\theta \&- \tfrac{1}{2} \sin 2\theta \sin^{2}\theta
  \&\&+ \tfrac{1}{3} \sin 3\theta \sin^{3} \theta - \dots
  \&\&= \arccot (1 + \cot\theta + \cot^{2}\theta),
\end{alignat*}
逆正接は $-\frac{1}{2}\pi$ から $\frac{1}{2}\pi$ の間である；類似の表現を級数
\[
\cos\theta \sin\theta - \tfrac{1}{2} \cos 2\theta \sin^{2}\theta + \dots,\quad
\sin\theta \cos\theta - \tfrac{1}{2} \sin 2\theta \cos^{2}\theta + \dots
\]
の和にも見出す。

236. ログ級数の応用。指数極限。任意の複素数 $z$ と十分小さい実数 $h$ を考える。このとき
\[
\log(1 + hz) = hz - \tfrac{1}{2}(hz)^{2} + \tfrac{1}{3}(hz)^{3} - \dots,
\]
よって
\[
\frac{\log(1 + hz)}{h} = z + \phi(h, z),
\]
ここで
\begin{gather*}
\phi(h, z) = -\tfrac{1}{2}hz^{2}
  + \tfrac{1}{3}h^{2}z^{3}
  - \tfrac{1}{4}h^{3}z^{4} + \dots,\\
|\phi(h, z)| < |hz^{2}| (1 + |hz| + |h^{2}z^{2}| + \dots)
  = \frac{|hz^{2}|}{1 - |hz|},
\end{gather*}
したがって $\phi(h, z) \to 0$ である。これより
\[
\lim_{h\to 0} \frac{\log(1 + hz)}{h} = z.
\Tag{(1)}
\]

特に $h = 1/n$ のとき、ここで $n$ は正の整数とする。このとき
\[
\lim_{n\to \infty} n\log \left(1 + \frac{z}{n}\right) = z,
\]
よって
\[
\lim_{n\to \infty} \left(1 + \frac{z}{n}\right)^{n}
  = \lim_{n\to \infty} \exp\left\{n\log\left(1 + \frac{z}{n}\right)\right\}
  = \exp z.
\Tag{(2)}
\]
これは実数 $z$ の場合で証明された結果の一般化である。

从式（1）我達可以推导出一些在下一節中需要用到的结果。もし$t$と$h$が実数であり、$h$が十分に小さいとき、

\[
\frac{\log(1 + tz + hz) - \log(1 + tz)}{h} = \frac{1}{h}\log\left(1 + \frac{hz}{1 + tz}\right)
\]

これは$h \to 0$の極限として$z/(1 + tz)$に収束します。従って、

\[
\frac{d}{dt} \{\log(1 + tz)\} = \frac{z}{1 + tz}.
\Tag{(3)}
\]

次節ではさらに、任意の実数または複素数$m$に対して$t$に関する$(1 + tz)^m$の微分公式が必要です。まず、複素関数$\phi(t) = \psi(t) + i\chi(t)$について考えます。実部と虚部が導関数を持つとき、

\[
\frac{d}{dt}(\exp\phi) = \frac{d}{dt}\{(\cos\chi + i\sin\chi) \exp\psi\}
\]

となり、

\[
= (\cos\chi + i\sin\chi) \psi' + (-\sin\chi + i\cos\chi)\chi'\right) \exp\psi
\]

\[
= (\psi' + i\chi')(\cos\chi + i\sin\chi) \exp\psi = (\psi' + i\chi') \exp(\psi + i\chi) = \phi' \exp\phi.
\]

したがって、$\exp\phi$の微分法則は実数の場合と同様です。これにより、

\[
\frac{d}{dt}(1 + tz)^m = \frac{d}{dt} \exp\{m\log(1 + tz)\}
\]

\[
= \frac{mz}{1 + tz} \exp\{m\log(1 + tz)\} = mz(1 + tz)^{m-1}.
\Tag{(4)}
\]

ここで、$(1 + tz)^m$と$(1 + tz)^{m-1}$は主値を持つ。

237. ビノミアル定理の一般形

既に（§215）証明済みであるが、級数
\[
1 + \binom{m}{1} z + \binom{m}{2} z^{2} + \dots
\]
の和はすべての実数$m$と$-1 < z < 1$のすべての実数に対する$(1 + z)^{m} = \exp\{m\log(1 + z)\}$である。$a_{n}$が$z^{n}$の係数であるとき
\[
\left|\frac{a_{n+1}}{a_{n}}\right| = \left|\frac{m - n}{n + 1}\right| \to 1,
\]
$m$が実数であれ複素数であれ、常に成り立つ。したがって（例題lxxx.3）この級数は$z$の絶対値が1未満であるとき常に収束し、その和は$\exp\{m\log(1 + z)\}$であり、すなわち$(1 + z)^{m}$の主値と等しいことを証明する。

§236から知るとおり、$t$が実数であるとき
\[
\frac{d}{dt}(1 + tz)^{m} = mz(1 + tz)^{m-1}
\]
であり、$z$と$m$はそれぞれ任意の実数または複素数で、両辺は主値を持つ。したがって、$\phi(t) = (1 + tz)^{m}$の場合には
\[
\phi^{(n)}(t) = m(m - 1) \dots (m - n + 1)z^{n} (1 + tz)^{m-n}
\]
である。この公式は$t = 0$でも成り立つので、
\[
\frac{\phi^{(n)}(0)}{n!} = \binom{m}{n} z^{n}.
\]

現在、§164の末尾で述べたことに従って
\[
\phi(1) = \phi(0) + \phi'(0) + \frac{\phi''(0)}{2!} + \dots + \frac{\phi^{(n-1)}(0)}{(n - 1)!} + R_{n},
\]
ここで
\[
R_{n} = \frac{1}{(n - 1)!}\int_{0}^{1} (1 - t)^{n-1} \phi^{(n)}(t)\, dt.
\]
また、もし $z = r(\cos\theta + i\sin\theta)$ であるなら
\[
|1 + tz| = \sqrt{1 + 2tr\cos\theta + t^{2}r^{2}} \geq 1 - tr,
\]
従って
\begin{align*}
|R_{n}|
  &< \frac{|m(m - 1) \dots (m - n + 1)|}{(n - 1)!}\, r^{n} \int_{0}^{1} \frac{(1 - t)^{n-1}}{(1 - tr)^{n-m}}\, dt\\
  &< \frac{|m(m - 1) \dots (m - n + 1)|}{(n - 1)!}\,
       \frac{(1 - \theta)^{n-1} r^{n}}{(1 - \theta r)^{n-m}},
\end{align*}
ここで $0 < \theta < 1$；したがって（§163を参照）
\[
|R_{n}| < K\frac{|m(m - 1) \dots (m - n + 1)|}{(n - 1)!}\, r^{n} = \rho_{n},
\]
とおく。しかし
\[
\frac{\rho_{n+1}}{\rho_{n}} = \frac{|m - n|}{n}r \to r,
\]
従って（Ex.xxvii.6）$\rho_{n} \to 0$ であり、したがって $R_{n} \to 0$ である。よって $n \to \infty$ のとき
\[
|R_{n}| \to 0.
\]
これにより以下の定理を得る。

定理
二項級数の和は
\[
1 + \binom{m}{1} z + \binom{m}{2} z^{2} + \dots
\]
が $\exp\{m\log(1 + z)\}$ であり、ここで対数は主値を取ることで、すべての実数または複素数 $m$ および $|z| < 1$ のすべての $z$ に対して成立する。

二項級数についてより完全な議論は、Bromwichの『無限級数』225ページ以降を参照されたい。

次頁

\par\medskip\noindent\textbf{【例 XCVIII.】}\par
\Item{1.} $m$が実数であるとき、次の式から
\[
\log(1 + z) = \frac{1}{2} \log(1 + 2r\cos\theta + r^{2}) + i\arctan\left(\frac{r\sin\theta}{1 + r\cos\theta}\right),
\]
を得る。したがって
\begin{align*}
\sum_{0}^{\infty} \binom{m}{n} z^{n}
  &= \exp\left\{\frac{1}{2}m \log(1 + 2r\cos\theta + r^{2})\right\}
     \Cis \left\{m\arctan\left(\frac{r\sin\theta}{1 + r\cos\theta}\right)\right\} \\
  &= (1 + 2r\cos\theta + r^{2})^{\frac{1}{2}m}
     \Cis \left\{m\arctan\left(\frac{r\sin\theta}{1 + r\cos\theta}\right)\right\},
\end{align*}
すべての逆正接が$-\frac{1}{2}\pi$から$\frac{1}{2}\pi$までである。特に、$\theta = \frac{1}{2}\pi$, $z = ir$と仮定し、実部と虚部を比較すると
\begin{align*}
1 - \binom{m}{2} r^{2} + \binom{m}{4} r^{4} - \dots
  &= (1 + r^{2})^{\frac{1}{2}m} \cos(m\arctan r), \\
\binom{m}{1} r - \binom{m}{3} r^{3} + \binom{m}{5} r^{5} - \dots
  &= (1 + r^{2})^{\frac{1}{2}m} \sin(m\arctan r).
\end{align*}

\Item{2.} $m = 1$, $2$, $3$のとき、公式1を検証せよ。当然、$m$が正の整数である場合、級数は有限である。

\Item{3.} $0 \leq r < 1$のとき
\begin{align*}
1 - \frac{1·3}{2·4} r^{2} + \frac{1·3·5·7}{2·4·6·8} r^{4} - \dots
  &= \sqrt{\frac{\sqrt{1 + r^{2}} + 1}{2(1 + r^{2})}}, \\
\frac{1}{2} r - \frac{1·3·5}{2·4·6} r^{3} + \frac{1·3·5·7·9}{2·4·6·8·10} r^{5} - \dots
  &= \sqrt{\frac{\sqrt{1 + r^{2}} - 1}{2(1 + r^{2})}}
\end{align*}
を証明せよ。

[最後の2つの公式1において$m = -\frac{1}{2}$と置く。]

4. 任意の実数 \( m \) に対して、\( -\frac{1}{4}\pi < \theta < \frac{1}{4}\pi \) のとき、

\[
\cos m\theta = \cos^{m} \theta \left\{1 - \binom{m}{2} \tan^{2} \theta
      + \binom{m}{4} \tan^{4} \theta - \dots\right\},
\]

\[
\sin m\theta = \cos^{m} \theta \left\{\binom{m}{1} \tan\theta
      - \binom{m}{3} \tan^{3} \theta + \dots\right\}.
\]

これらの結果は、

\[
\cos m\theta + i\sin m\theta = (\cos\theta + i\sin\theta )^{m}
  = \cos^{m} \theta(1 + i\tan\theta)^{m}
\]

から即座に導かれる。

5. \( f(m, z) = \sum\dbinom{m}{n} z^{n} \)（\( |z| < 1 \)）は、直接列級数の乗算により (\Ex{lxxxi}.~6) で証明され、関数方程式

\[
f(m, z) f(m', z) = f(m + m', z)
\]

を満たす。一般結果 (\PageRef{p.}{423}) を仮定せずに、\SecNo[§]{216} に類似した議論により導くと、

\[
f(m, z) = \exp\{m\log(1 + z)\}
\]

が成り立つことを示す。

6. \( z \) および \( \mu \) が実数で、\( -1 < z < 1 \) のとき、

\[
\sum \binom{i\mu}{n} z^{n} = \cos\{\mu\log(1 + z)\} + i\sin\{\mu\log(1 + z)\}.
\]

7. 1. \( i^{\log(1+i)} \) の実部は、

\[
e^{(4k+1)\pi^{2}/8 } \cos \left\{\frac{1}{4}(4k + 1)\pi\log 2\right\}
\]

である。ここで \( k \) は任意の整数。

2. \( a\cos\theta + b\sin\theta + c = 0 \)（\( a,~b,~c \) が実数で、\( c^{2} > a^{2} + b^{2} \)）のとき、

\[
\theta = m\pi + \alpha
  ± i\log \frac{|c| + \sqrt{c^{2} - a^{2} - b^{2}}}{\sqrt{a^{2} + b^{2}}},
\]

ここで \( m \) は任意の奇数または偶数、\( c > 0 \) のとき奇数、\( c < 0 \) のとき偶数であり、

\[
\cos\alpha = \frac{a}{\sqrt{a^{2} + b^{2}}}, \quad \sin\alpha = \frac{b}{\sqrt{a^{2} + b^{2}}}.
\]

3. \( \theta \) が実数で、\( \sin\theta \sin\phi = 1 \) のとき、

\[
\phi = (k + \tfrac{1}{2})\pi ± i\log \cot \tfrac{1}{2}(k\pi + \theta),
\]

ここで \( k \) は任意の偶数または奇数、\( \sin\theta > 0 \) のとき偶数、\( \sin\theta < 0 \) のとき奇数。

4. \(x\)が実数のとき、以下の関係を示せ：
\[
\frac{d}{dx} \exp\{(a + ib)x\} = (a + ib) \exp\{(a + ib) x\},
\]
\[
\int \exp \{(a + ib)x\}\, dx = \frac{\exp{(a + ib)x}}{a + ib}.
\]
これらを用いて、問題487の結果を導け。

5. \(a > 0\)のとき、以下の積分を示せ：
\[
\int_{0}^{\infty} \exp\{-(a + ib)x\}\, dx = \frac{1}{a + ib},
\]
これらを用いて、問題487の結果を導け。

6. 極座標方程式 \((x/a)^2 + (y/b)^2 = 1\)が楕円を表すとき、他の代数曲線の方程式において最高次項を表す \(f(x, y)\) を用いて、楕円とその曲線の交点の偏角の和は、
\[
-i\{\log f(a, ib) - \log f(a, -ib)\}
\]
から2πの倍数だけ異なることを示せ。

[これらの偏角は、\(f(a\cos\alpha, b\sin\alpha) + \dots = 0\)または
\[
f\left\{\frac{1}{2} a \left(u + \frac{1}{u}\right), -\frac{1}{2} ib \left(u - \frac{1}{u}\right)\right\} + \dots = 0,
\]
ここで \(u = \exp i\alpha\)とし、\(\sum\alpha\)はこの方程式の根の積 \(P\) の \(-i\log P\) の値の一つに等しい]

7. 方程式 \(\tan z = az\) の解の数と近似的位置を決定せよ。ここで \(a\) は実数である。

既に知っている（\Ex{xvii}.4）ことでは、方程式は無数の実数解を持つ。今、$z = x + iy$ と置き、実部と虚部を等しくする。次のように得られる：
\[
\frac{\sin 2x}{\cos 2x + \cosh 2y} = ax,\quad
\frac{\sinh 2y}{\cos 2x + \cosh 2y} = ay,
\]
従って、$x$ や $y$ が零でない場合、次が成り立つ：
\[
\frac{\sin 2x}{2x} = \frac{\sinh 2y}{2y}.
\]
このことは不可能であり、左辺は絶対値で言えば1未満であり、右辺は絶対値で言えば1より大きい。したがって $x = 0$ や $y = 0$ である。$y = 0$ の場合、方程式の実数解に戻る。$x = 0$ の場合、$\tanh y = ay$ となる。$a \leq 0$ や $a \geq 1$ のとき、この方程式は零以外の実数解を持たず、$0 < a < 1$ のときには二つのような実数解を持つ。したがって、$0 < a < 1$ のときには二つの純虚数解があり、それ以外の場合にはすべての解が実数である。

8. 方程式 $\tan z = az + b$（$a$ と $b$ は実数で、$b \neq 0$）は $a \leq 0$ のとき複素数解を持たない。$a > 0$ の場合、すべての複素数解の実部は絶対値で言えば $|b/2a|$ より大きい。

9. 方程式 $\tan z = a/z$（$a$ は実数）は複素数解を持たないが、$a < 0$ のとき二つの純虚数解を持つ。

10. 方程式 $\tan z = a\tanh cz$（$a$ と $c$ は実数）は無数の実数解と純虚数解を持つが、複素数解を持たない。

11. $x$ が実数であるとき、次のように表される：
\[
e^{ax} \cos bx = \sum_{0}^{\infty} \frac{x^{n}}{n!} \left\{
  a^{n} - \binom{n}{2} a^{n-2} b^{2} + \binom{n}{4} a^{n-4} b^{4} - \dots
\right\},
\]
大括弧内の項は $\frac{1}{2}(n + 1)$ または $\frac{1}{2}(n + 2)$ 個である。同様の級数を $e^{ax} \sin bx$ についても求めよ。

12. $n\phi(z, n) \to z$ として $n \to \infty$ のとき、$\{1 + \phi(z, n)\}^{n} \to \exp z$。

13. 実数変数 $t$ の複素関数 $\phi(t)$ に対しては、次が成り立つ：
\[
\frac{d}{dt} \log \phi(t) = \frac{\phi'(t)}{\phi(t)}.
\]

\%[** TN: Paragraph break added]
[次の公式を使用する：
\[
\phi = \psi + i\chi,\quad
\log \phi = \frac{1}{2}\log(\psi^{2} + \chi^{2}) + i\arctan\left(\frac{\chi}{\psi}\right).]
\]

14. 变換。第III章（§xxi．21 以降，および§MiscExs III 22 以降）では，二変数 $z$，$Z$ の平面間に関係 $z = f(Z)$ を持つ図形の幾何学的関係についていくつかの単純な例を検討しました。現在は，対数関数，指数関数，または三角関数を含む関係の場合について考察します。

まず，次の関係を考えます：
\[
z = \exp\left(\frac{\pi Z}{a}\right),\quad
Z = \frac{a}{\pi} \Log z
\]
ここで $a$ は正の定数です。一つの $Z$ の値に対応する $z$ の値が一つですが，一方の $z$ の値には無数の $Z$ の値が対応します。座標 $(x, y)$ で $z$ を表し，$(X, Y)$ で $Z$ を表すと，次の関係があります：
\begin{alignat*}{2}
x \&= e^{\frac{\pi X}{a}} \cos\left(\frac{\pi Y}{a}\right), \& y \&= e^{\frac{\pi X}{a}} \sin\left(\frac{\pi Y}{a}\right),\\
X \&= \frac{a}{\pi} \log r, \& Y \&= \frac{a\theta}{\pi} + 2ka,
\end{alignat*}
ここで $k$ は任意の整数です。$\theta < -\pi$ または $\theta > \pi$ の場合，および主値 $\Log z = \log z$ を仮定するとき，$k = 0$ であり，$Z$ は $OX$ 軸に平行な平面の帯状領域に制限されます。この帯状領域の一つの点が全体の $z$-平面上の一つの点に対応し，逆もまた然りです。主値以外の $\Log z$ の値を仮定すると，$z$-平面と $Z$-平面内の幅 $2a$ の別の帯状領域との間にも類似の関係が得られます。

$Z$-平面上の$X$と$Y$が一定である直線に対応するものは、$z$-平面上の$r$と$\theta$が一定である円や半径ベクトル全体です。一方で、$r$が一定な円に対応するのは$OY$に平行な直線の一部であり、その長さは$2a$です。$Z$が後者の直線全体を描くためには、$z$は円周上を無限に回る必要があります。

15. $Z$-平面上の直線に対応するものは、$z$-平面上の等角螺旋であることを示せ。

16. 同様にして、変換$z = c\cosh(\pi Z/a)$について議論し、特に、$z$-平面上の全てが無数の$OX$に平行な幅$2a$の帯状領域に対応することを示せ。また、$X = X_{0}$の直線に対応する楕円は
\[
\left\{\frac{x}{c\cosh(\pi X_{0}/a)}\right\}^{2} +
\left\{\frac{y}{c\sinh(\pi X_{0}/a)}\right\}^{2} = 1,
\]
であり、異なる$X_{0}$の値に対してこれらの楕円は共軸系を形成し、$Y = Y_{0}$の直線に対応する共役共軸双曲線系も形成します。$Z$が$X = X_{0}$または$Y = Y_{0}$の直線全体を描くとき、$z$の変化を追跡せよ。また、$z$が共軸系の焦点間のセグメントと$x$軸の残りの部分で形成される退化した楕円や双曲線を描くとき、$Z$はどのように変化するか。

17. 16問の結果が14問および第III章の25問の結果と一致することを確認せよ。変換$z = c\cosh(\pi Z/a)$は以下の変換から成り立つものとみなすことができる：
\[
z = cz_{1},\quad
z_{1} = \tfrac{1}{2}\{z_{2} + (1/z_{2})\},\quad
z_{2} = \exp(\pi Z/a).
\]

\item{18.} 同様に，変換 \( z = c \tanh(\pi Z/a) \) を議論し，直線 \( X = X_{0} \) に対応する共軸円系が
\[
\left\{ x - c \coth\left(2\pi X_{0}/a\right) \right\}^{2} + y^{2} = c^{2} \cosech^{2}\left(2\pi X_{0}/a\right)
\]
であり，直線 \( Y = Y_{0} \) に対応する直交系の共軸円系であることを示す。

\item{19.} \Topic{立体投影とメルカトル投影。} 単位球の中心が原点である点を南極（座標 \(0, 0, -1\)）から北極の接平面に射影する．球面上の点の座標は \(\xi\), \(\eta\), \(\zeta\) であり，接平面上の直交座標系 \(OX\), \(OY\) を \(\xi\) 軸と \(\eta\) 軸と並行に取る．このとき，点の射影座標は
\[
x = \frac{2\xi}{1 + \zeta}, \quad y = \frac{2\eta}{1 + \zeta}
\]
であり，さらに \( x + iy = 2 \tan \frac{1}{2}\theta \Cis \phi \) と表される．ここで，\(\phi\) は経度（平面 \(\eta = 0\) から測定）で，\(\theta\) は球面上の点の北極距離である．

この射影は，球面を接平面上に投影するものであり，一般的には立体投影と呼ばれる．次に新しい複素変数
\[
Z = X + iY = -i \log \frac{1}{2}z = -i \log \frac{1}{2}(x + iy)
\]
を導入し，\(X = \phi\), \(Y = \log \cot \frac{1}{2}\theta\) とおくことで，平面 \(Z\) 上の別の投影を得る．この投影では，緯度と経度の平行線はそれぞれ \(X\) 軸と \(Y\) 軸に並行な直線で表される．

\item{20.} 方程式
\[
z = \Log \left( \frac{Z - a}{Z - b} \right)
\]
によって与えられる変換について議論し，\(x\) と \(y\) が定数となる直線は \(Z\) 平面上の2つの直交系の共軸円系に対応することを示す．

\Item{21.} 複素数の変換
\[
z = \Log \left\{\frac{\sqrt{Z - a} + \sqrt{Z - b}}{\sqrt{b - a}}\right\}
\]
について議論し、$x$ と $y$ が一定である直線は共役楕円と双曲線の族に対応し、その焦点は点 $Z = a$ および $Z = b$ であることを示す。

[私たちは
\begin{alignat*}{2}
\sqrt{Z - a} + \sqrt{Z - b} \&= \sqrt{b - a}\, \exp(\& \&x + iy), \\
\sqrt{Z - a} - \sqrt{Z - b} \&= \sqrt{b - a}\, \exp(\&-\&x - iy);
\end{alignat*}
とし、
\[
|Z - a| + |Z - b| = |b - a|\cosh 2x,\quad
|Z - a| - |Z - b| = |b - a|\cos 2y.
\]
]

\Item{22.} \Topic{変換 $z = Z^{i}$。} $z = Z^{i}$ において、虚数の冪が主値を持つとき
\[
\exp(\log r + i\theta) = z = \exp(i\log Z) = \exp(i\log R - \Theta),
\]
であるから $\log r = -\Theta$, $\theta = \log R + 2k\pi$ であり、ここで $k$ は整数。すべての $k$ の値が同じ点 $z$ を与えるため、$k = 0$ と仮定し
\[
\log r = -\Theta,\quad
\theta = \log R.
\Tag{(1)}
\]

$Z$ 平面上の全体は、$R$ がすべての正の値をとり、$\Theta$ が $-\pi$ から $\pi$ まで動くときに覆われる。これにより $r$ の範囲は $\exp(-\pi)$ から $\exp(\pi)$ までとなり、$\theta$ はすべての実数値を通る。したがって $Z$ 平面は半径 $\exp(-\pi)$ と $\exp(\pi)$ で囲まれたリングに対応し、このリングは無限に覆われる。ただし、$\theta$ が $-\pi$ から $\pi$ の範囲だけ動くとき、リングは一度しか覆われないため、$R$ は $\exp(-\pi)$ から $\exp(\pi)$ の範囲で動く。これにより $Z$ の変化は、$z$ が動くリングと同様の性質を持つリングに制限される。さらに、各リングには負の実軸を通り得ない境界があり、これは $z$（または $Z$）がその振幅が $-\pi$ から $\pi$ の範囲を超えないことを意味する。

We これで，2つの環間の対応を次のように得る．その対応は以下の式によって与えられる：
\[
z = Z^{i}, \quad Z = z^{-i},
\]
ここで各べき乗は主値を持つ。一方の平面における円が原点を中心とする円に対応するものは，他方の平面における原点を通る直線である。

23. $Z$ が点 $\exp \pi$ から正方向に大きい円を一周して点 $-\exp \pi$ まで移動し，障害物を通過して小さい円を逆方向に一周し，再び障害物を通過して大きい円の残り部分を一周して元の位置に戻るとき，$z$ の変化を追跡せよ。

24. 各平面を，半径が
\[
\dots, \quad e^{-(2n+1)\pi}, \ \dots, \quad e^{-\pi}, \quad e^{\pi}, \quad e^{3\pi}, \ \dots, \quad e^{(2n+1)\pi}, \ \dots
\]
の円で無限に分割すると，一方の平面における任意の環に対応する環を他方の平面における任意の環と対応させる方法として，式 $z = Z^{i}$, $Z = z^{-i}$ の適切なべき乗値を選ぶことで示せ。

25. $z = Z^{i}$ において，$Z$ のどのべき乗値も選んでもよいが，$Z$ がその平面における極が原点である等角螺旋を移動するとき，$z$ もその平面における極が原点である等角螺旋を移動する。

26. $Z = z^{ai}$, $a$ が実数のとき，$z$ が実軸を通って原点に近づくとき，$Z$ は原点を中心とする円周上を無限に回り続け，$Z$ の実部と虚部は有限に振動する。

27. $Z = z^{a+bi}$, $a$ および $b$ が任意の実数のとき，同じ問題について議論せよ。

28. 実数 $a$ に対する型 $\sum\limits_{-\infty}^{\infty} a_{n}z^{nai}$ の級数の収束領域は角度であることを示せ，すなわち，不等式 $\theta_{0} < \am z < \theta_{1}$ によって定義される領域である．（角度は線に縮小するか，全体の平面を覆うことができる。）

\item[29.] \topic{等高線．}複素変数$z$の関数$f(z)$に対して、$|f(z)|$が一定値となる曲線を$f(z)$の等高線と呼ぶ．以下の関数の等高線の形状を図示せよ：
\begin{alignat*}{2}
z - a \quad\& (\text{同心円}), \qquad\&
(z - a)(z - b) \quad\& (\text{カートシアン卵形}), \\
\frac{z - a}{z - b} \quad\& (\text{共軸円}), \qquad\&
\exp z \quad\& (\text{直線}).
\end{alignat*}

\item[30.] $(z - a)(z - b)(z - c)$, $\frac{1 + z\sqrt{3} + z^{2}}{z}$の等高線の形状を図示せよ．[後者の関数の一部の等高線は図59に描かれている．曲線が標記された \textsc{i}--\textsc{vii} に対応する値は
\[
.10,\quad 2 - \sqrt{3} = .27,\quad
.40,\quad 1.00,\quad 2.00,\quad
2 + \sqrt{3} = 3.73,\quad 4.53
\]
である．読者は、与えられた任意の有理関数の等高線の形状について一般的な概念を形成するのにほとんど困難を感じないだろうが、詳細に掘り下げると複素変数の関数の一般理論に到達することになる．]

\item[31.] (i) $z \exp z$, (ii) $\sin z$の等高線の形状を図示せよ．[図60は$\sin z$の等高線を表している．曲線が標記された \textsc{i}--\textsc{viii} に対応する値は $k = .35$, $.50$, $.71$, $1.00$, $1.41$, $2.00$, $2.83$,~$4.00$ である．]

\item[32.] 水平面曲線の形を描き，$c$が実定数であるとき，$\exp z - c$ の等位曲線を概略的に示す．[\Fig{61}は$|\exp z - 1|$の等位曲線を示し，\textsc{i}--\textsc{vii}は$k = -1.00, -0.20, -0.05, 0.00, 0.05, 0.20, 1.00$に対応する曲線である．]

\item[33.] $c$が正の定数であるとき，$\sin z - c$ の等位曲線は，\Fig{62}，\Fig{63}に示されている．これらの曲線の性質は$c < 1$または$c > 1$により異なる．\Fig{62}では$c = 0.5$を採用し，\textsc{i}--\textsc{viii}は$k = 0.29, 0.37, 0.50, 0.87, 1.50, 2.60, 4.50, 7.79$に対応する．\Fig{63}では$c = 2$を採用し，\textsc{i}--\textsc{vii}は$k = 0.58, 1.00, 1.73, 3.00, 5.20, 9.00, 15.59$に対応する．$c = 1$のとき，これらの曲線は\Fig{60}と同様であり，ただし原点とスケールが異なる．

\item[34.] $0 < \theta < \pi$ のとき，
\begin{alignat*}{3}
\cos\theta \&+ \frac{1}{3} \cos 3\theta \&\&+ \frac{1}{5} \cos 5\theta \&\&+ \cdots
  = \frac{1}{4} \log \cot^{2}\frac{1}{2}\theta,\\
\sin\theta \&+ \frac{1}{3} \sin 3\theta \&\&+ \frac{1}{5} \sin 5\theta \&\&+ \cdots
  = \frac{1}{4}\pi,
\end{alignat*}
が成り立つことを証明し，それらの級数が収束するすべての他の値における和を決定せよ．[方程式
\[
z + \frac{1}{3}z^{3} + \frac{1}{5}z^{5} + \cdots
  = \frac{1}{2} \log \left(\frac{1 + z}{1 - z}\right)
\]
を用いる．$\theta$が$\pi$に増加するとき，各級数の和は単純に符号が変化する．したがって，最初の公式はすべての$\theta$の値（$\pi$の倍数では級数が発散する）で成り立つことがわかる．一方，第二級数の和は$2k\pi < \theta < (2k + 1)\pi$のとき$\frac{1}{4}\pi$，$(2k + 1)\pi < \theta < (2k + 2)\pi$のとき$-\frac{1}{4}\pi$，および$\theta$が$\pi$の倍数のとき$0$である．]

\Item{35.} \(0 < \theta < \frac{1}{2}\pi\) のとき、次の等式が成り立つことを証明せよ：
\begin{alignat*}{3}
\cos\theta \&- \tfrac{1}{3} \cos 3\theta \&\&+ \tfrac{1}{5} \cos 5\theta \&\&- \dots
  = \tfrac{1}{4}\pi,\\
\sin\theta \&- \tfrac{1}{3} \sin 3\theta \&\&+ \tfrac{1}{5} \sin 5\theta \&\&- \dots
  = \tfrac{1}{4} \log (\sec\theta + \tan\theta)^{2};
\end{alignat*}
これらの級数の和を、それらが収束するすべての他の値の \(\theta\) についても定める。

\Item{36.} \(0 < \theta - \alpha, \theta + \alpha < 2\pi\) のとき以外に、次の等式が成り立つことを証明せよ：
\[
\cos\theta \cos\alpha
  + \tfrac{1}{2} \cos 2\theta \cos 2\alpha
  + \tfrac{1}{3} \cos 3\theta \cos 3\alpha + \dots
  = -\tfrac{1}{4} \log \{4(\cos\theta - \cos\alpha)^{2}\}.
\]

\Item{37.} \(a, b\) のいずれも実数でないとき、次の積分が成り立つことを証明せよ：
\[
\int_{0}^{\infty} \frac{dx}{(x - a)(x - b)}
  = -\frac{\log(-a) - \log(-b)}{a - b},
\]
各対数は主値を持つ。\(a = ci, b = -ci\)（\(c > 0\)）のとき、この結果を検証せよ。また、\(a\) や \(b\) が実数で負である場合についても議論せよ。

\Item{38.} \(\alpha, \beta\) が実数であり、\(\beta > 0\) のとき、次の積分が成り立つことを証明せよ：
\[
\int_{0}^{\infty} \frac{d}{x^{2} - (\alpha + i\beta)^{2}}
  = \frac{\pi i}{2(\alpha + i\beta)}.
\]
\(\beta < 0\) のとき、この積分の値を求めよ。

\Item{39.} \(Ax^{2} + 2Bx + C = 0\) の根が虚部符号異なるとき、次の積分が成り立つことを証明せよ：
\[
\int_{-\infty}^{\infty} \frac{dx}{Ax^{2} + 2Bx + C}
  = \frac{\pi i}{\sqrt{B^{2} - AC}},
\]
\(\sqrt{B^{2} - AC}\) の符号は、実部が正となるように選ぶ。

\Appendix{I}{（III, IV, V 章への付録）}{方程式の根の存在証明}

初めに、$z$ の多項式 $P(z)$ を定義します。
\[
Z = P(z) = \alpha_{0} z^{n} + \alpha_{1} z^{n-1} + \dots + \alpha_{n}
\]
この式において、$\alpha_0, \alpha_1, \ldots, \alpha_n$ は実数または複素数の係数です。$z$ の値と $Z$ の値をそれぞれ $z$ 平面と $Z$ 平面上の点で表すことができます。もし $z$ が $z$ 平面上の閉曲線 $\gamma$ を描くならば、$Z$ も同様に $Z$ 平面上の閉曲線 $\Gamma$ を描きます。現在は、$\Gamma$ が原点を通らないことを仮定します。

任意の $Z$ の値に対して無数の異なる $\am Z$ の値があり、これらは $Z$ が $\Gamma$ 上を動くにつれて連続的に変化します。\footnote{ここで $\Gamma$ が原点を通らないという仮定が必要です。}各点に対応する特定の $\am Z$ の値を選択するために、まず初期の $Z$ の値に対応する特定の値を選び、その後 $Z$ が $\Gamma$ 上を動くにつれてこの値の連続的な変化を追います。以下の議論においては、「$Z$ の振幅」や式 $\am Z$ を用いて、選択した特定の振幅の値を指すことにします。したがって、$\am Z$ は実部と虚部 $X, Y$ の関数として一価かつ連続な関数を表します。
ページ分割点：434

当$Z$が$\Gamma$を描き終え元の位置に戻ったとき、その振幅は以前と同様である場合もある。これは$\Gamma$が原点を含まない場合、つまり図Bの(\ia)のような経路の場合に確実に成り立つ。また、その振幅は元の値から任意の$2\pi$の倍数だけ異なる場合も存在する。例えば、図Bの(\ib)のように原点周りに1周した場合、その振幅は$2\pi$増加する。これらの言及は$\Gamma$だけでなく、$Z$平面上の任意の閉曲線についても成り立つ。これら全てが原点を通らない限りである。どのような経路に対しても、その経路を描くときの振幅の変化量がある数値とみなすことができる。この数値は特定の振幅の初期選択に依存しない。

次に証明するが、次の結果が成り立つ：もし$Z$が元の位置に戻ったときにその振幅が異なるなら、$z$の経路は少なくとも1点を含むか、もしくはその上にある点において$Z = 0$である。

γをいくつかの小さなContoursに分割するために、軸に対して距離δ₁で平行線を引くことができる（図C参照）。これらの平行線の構築には明確な規則がある：最も自然な方法は、すべての直線x = kδ₁, y = kδ₁を描くことである。kは整数であり、正でも負でもよい。

もしこれらのContoursのいずれかの境界線上にZ = 0となる点が存在する場合、証明したいことは既に成り立っている。したがって、この状況では、そのような点が存在しないと仮定できる。

zがγを描くときのam Zの増分は、それぞれの小さなContoursを順番に描くときのam Zの増分の合計に等しい。もしzがγと同じ向きでそれぞれの小さなContoursを順番に描く場合、最終的には（図D参照）γの境界線を1回描き、各分割平行線は2回描かれ、反対方向で描かれる。

したがって、PQは2回描かれる：一度はPからQへ、もう一度はQからPへ。zがPからQへ移動するとき、am Zは連続的に変化し、Zは原点を通らないため；この場合、am Zの増分θがある。そして、zがQからPへ移動するときのam Zの増分は-θである。したがって、それぞれの小さなContoursの各部分を描くときのam Zの増分の合計は相互に打ち消され、γ自身の部分を描くときの増分だけが残る。

435ページへ続く。

したがって、$z$ が $\gamma$ を描くとき $\mathrm{arg}\,Z$ が変化するならば、少なくとも一つの小さな contour，すなわち $\gamma_1$ について $\mathrm{arg}\,Z$ が $z$ が $\gamma_1$ を描くときに変化する。

この contour は、補助平行線の一部で構成される正方形であるか、または補助平行線と $\gamma$ の境界からなるものである可能性がある。いずれの場合も、contour 上の任意の点は長さ $\delta_1$ の補助平行線の一部で構成される正方形 $\Delta_1$ 内に存在するか、その境界上にある。

次に、軸に対して距離 $\delta_2$ である小さな平行線を用いて $\gamma_1$ をさらに分割し、$\gamma_2$ を $\Delta_2$ の内部に完全に含む contour として定義できる。$\Delta_2$ の辺の長さは $\delta_2$ であり、また $\Delta_1$ 内に含まれている。

今後、零に収束する無限数列 $\delta_1, \delta_2, \dots, \delta_m, \dots$ を取る。例えば、$\delta_m = \delta_1 / 2^{m-1}$ とすることができる。
上記の議論を繰り返し使用することで、正方形の系列 $\Delta_1, \Delta_2, \dots, \Delta_m, \dots$ と contour の系列 $\gamma_1, \gamma_2, \dots, \gamma_m, \dots$ を決定できる。これらの条件を満たすように (i) $\Delta_{m+1}$ は完全に $\Delta_m$ 内に含まれている、(ii) $\gamma_m$ は完全に $\Delta_m$ 内に含まれている、(iii) $z$ が $\gamma_m$ を描くときに $\mathrm{arg}\,Z$ が変化する。

もし $(x_m, y_m)$ と $(x_m + \delta_m, y_m + \delta_m)$ が $\Delta_m$ の左下角と右上角であるとき、明らかに $x_1$, $x_2$, \dots, $x_m$, \dots は増加し、$x_1 + \delta_1$, $x_2 + \delta_2$, \dots, $x_m + \delta_m$, \dots は減少する数列であり、これらは共通の極限値 $x_0$ をもつ。同様に $y_m$ と $y_m + \delta_m$ も共通の極限値 $y_0$ をもつし、$(x_0, y_0)$ はすべての正方形 $\Delta_m$ 内に存在する唯一の点である。どの程度でも小さい $\delta$ についても、$(x_0, y_0)$ を含む辺が軸に平行で長さが $\delta$ の正方形を描くことができ、この正方形内には $z$ が回るとき $Z$ が変化する閉曲線を描くことができる。

今、
\[
P(x_0 + iy_0) = 0
\]
であることを示す。もし $P(x_0 + iy_0) = a$ であり、$|a| = \rho > 0$ のとき、$P(x + iy)$ は $x$ と $y$ の連続関数であるから、$(x_0, y_0)$ を中心とする正方形を描くことができ、その辺が軸に平行であり、かつ
\[
|P(x + iy) - P(x_0 + iy_0)| < \frac{1}{2}\rho
\]
となるような点 $x + iy$ 内部または境界上すべてにおいて。このような点について
\[
P(x + iy) = a + \phi
\]
であり、$|\phi| < \frac{1}{2}\rho$ である。今、この正方形内を完全に含む任意の閉曲線を考える。 $z$ がこの曲線を描くとき、$Z = a + \phi$ も同様の閉曲線を描く。しかし、後者の曲線は明らかに中心が $a$ で半径が $\frac{1}{2}\rho$ の円内に含まれており、この円は原点を含まない。したがって、$Z$ の振幅は変化しない。

しかし、これは上記で証明されたことを矛盾させる。すなわち、各正方形 $\Delta_m$ 内には $z$ が回るとき $Z$ の振幅が変化する閉曲線がある。したがって
\[
P(x_0 + iy_0) = 0
\]
である。

全てが残されているのは、あるContoursを選ぶことができることを示すことです。これにより \( Z \) は \( z \) が \(\gamma\) を描くときに変化します。

\[
Z = a_{0} z^{n} \left(1 + \frac{a_{1}}{a_{0}z} + \frac{a_{2}}{a_{0} z^{2}} + \dots
      + \frac{a_{n}}{a_{0} z^{n}}\right).
\]

任意の正の数 \(\delta\) でも何にせよ、次のように \( R \) を選ぶことができます：

\[
\frac{|a_{1}|}{|a_{0}| R} +
\frac{|a_{2}|}{|a_{0}| R^{2}} + \dots +
\frac{|a_{n}|}{|a_{0}| R^{n}} < \delta.
\]

そして、中心が原点で半径が \( R \) の円 \(\gamma\) 上のすべての点において、

\[
Z = a_{0} z^{n} (1 + \rho),
\]
ここで \( |\rho| < \delta \) です。次に、\( z \) が正方向で \(\gamma\) を描くとき、\( \am(1 + \rho) \) の値は変化しないことを示す方法として、先ほど使用した類似の議論を使用できます。一方で \( \am z^{n} \) は \( 2n\pi \) 増加します。したがって \( \am Z \) は \( 2n\pi \) 増加し、\( Z = 0 \) の根の存在を示す証明が完了します。

議論全体を通じて、\(\Gamma\) またはその分解された任意の小さなContoursが原点を通らないという仮定を行っています。この仮定は明らかに正当で、反対の仮定をした場合、定理の真偽を認めることになります。

読者には、前段落と第43節の議論から、次の結果を導くことを課題として残します：

\par
\text{任意のContours} \gamma \) において \( z \) が正方向で描くとき、\( \am Z \) の増分は \( 2k\pi \) であり、ここで \( k \) は \( Z = 0 \) の内部にある根の数です（多重根は重複して数えます）。


他にも異なる方法での証明がありますが、これは第102節以降の結果を多変数関数に拡張することに基づいています。

我々は、第102節と同様の方法で、関数$f(x, y)$の上界と下界を定義します。任意の平面$(x, y)$上の閉曲線$\gamma$で囲まれた領域内のすべての$x$および$y$の組に対応する点についてです。そして、連続な関数$f(x, y)$は、このような領域内でその上界と下界を達成することが証明できます。

\[
|Z| = |P(x + iy)|
\]
は$x$および$y$に対する正の連続関数です。点$\gamma$上のまたはその内部にある点で$|Z|$の下界$m$が存在する場合、ある点$z_{0}$が存在し、$|Z| = m$を満たします。これは$|Z|$の最小値をとる点です。もし$m = 0$であれば、$P(z_{0}) = 0$となり、求めたいことを証明しました。したがって、$m > 0$であると仮定できます。

点$z_{0}$は$\gamma$の内部または境界上に存在しなければなりませんが、もし$\gamma$が原点を中心とする円で、その半径$R$が十分大きい場合、最後の仮定は成り立たなくなります。すなわち、$|z| \to \infty$ として $|P(z)| \to \infty$ です。したがって、$z_{0}$は$\gamma$の内部に存在します。

ページ分割点437

もし$z = z_{0} + \zeta$と置き、$P(z)$を$\zeta$のべき乗に並び替えれば、

\[
P(z) = P(z_{0}) + A_{1}\zeta + A_{2}\zeta^{2} + \dots + A_{n}\zeta^{n}
\]

と表すことができます。$A_{k}$が最初の非零係数であり、$|A_{k}| = \mu$, $|\zeta| = \rho$とすることができます。$\rho$を十分小さく選べば、

\[
|A_{k+1}|\rho + |A_{k+2}|\rho^{2} + \dots + |A_{n}|\rho^{n-k} < \tfrac{1}{2}\mu
\]

が成り立ちます。すると、

\[
|P(z) - P(z_{0}) - A_{k}\zeta^{k}| < \tfrac{1}{2}\mu\rho^{k}
\]

かつ

\[
|P(z)| < |P(z_{0} + A_{k}\zeta^{k})| + \tfrac{1}{2}\mu\rho^{k}.
\]

今、$z$が中心$z_{0}$、半径$\rho$の円を一周するとき、
\[
P(z_{0}) + A_{k}\zeta^{k}
\]
は$k$回中心$P(z_{0})$、半径$|A_{k}\zeta^{k}| = \mu\rho^{k}$の円を一周し、この円と原点を通る直線との交点を通る。したがって、$z$が描く円上には$k$点で
\[
|P(z_{0}) + A_{k}\zeta^{k}| = |P(z_{0})| - \mu\rho^{k}
\]
となり、
\[
|P(z)| < |P(z_{0})| - \mu\rho^{k} + \frac{1}{2}\mu\rho^{k}
  = m - \frac{1}{2}\mu\rho^{k}
  < m;
\]
となる。これは$|P(z)|$の下限$m$が存在することと矛盾する。

したがって、$m$は$0$でなければならず、$P(z_{0}) = 0$である。

付録Iの例題

1. $f(z) = 0$の根の数を閉曲線内に含むことを示せ。この曲線を通るとき$\{\log f(z)\}/2\pi i$の増加量と一致する。

2. 任意の数$R$について、
\[
\frac{|a_{1}|}{R} + \frac{|a_{2}|}{R^{2}} + \dots + \frac{|a_{n}|}{R^{n}} < 1,
\]
が成り立つとき、$z^{n} + a_{1}z^{n-1} + \dots + a_{n} = 0$の根はすべて絶対値で$R$未満であることを示せ。特に、$z^{5} - 13z - 7 = 0$の根は絶対値で$2\frac{1}{67}$未満であることを示せ。

3. 実数$a$と$b$、奇数$p$に対して、方程式$z^{2p} + az + b = 0$の実部が正または負の根の数を求めよ。$a > 0$, $b > 0$のとき、その数は$p - 1$と$p + 1$である；$a < 0$, $b > 0$のとき、その数は$p + 1$と$p - 1$である；$b < 0$のとき、その数は$p$と$p$である。$a = 0$または$b = 0$の場合を議論せよ。特に$p = 1$のとき結果を検証せよ。

[大半径$R$の中心原点を持つ大きな半円と、その半円によって切断される虚軸部分を形成する曲線を通るとき$\am(z^{2p} + az + b)$の変化を追跡せよ。]

4. 同様に，以下の方程式を考える：
\[
z^{4q} + az + b = 0,\quad
z^{4q-1} + az + b = 0,\quad
z^{4q+1} + az + b = 0.
\]
\PageSep{438}

5. $\alpha$ と $\beta$ が実数であるとき，方程式 $z^{2n} + \alpha^{2} z^{2n-1} + \beta^{2} = 0$ の実部が正または負の根の個数は，$n$ が奇数であれば $n - 1$ と $n + 1$，偶数であれば $n$ と $n$ であることを示せ。
\MathTrip{1891.}

6. $z$ が点 $z = z_{1}$, $z = z_{2}$ を結ぶ直線上から近接する点から移動するとき，以下の式の増分はほぼ $\pi$ に等しいことを示せ：
\[
\am \left(\frac{1}{z - z_{1}} + \frac{1}{z - z_{2}}\right).
\]

7. $z = z_{1}$, $z = z_{2}$, $z = z_{3}$ を囲む contour は，三角形 $z_{1}$, $z_{2}$, $z_{3}$ 形成した辺の一部と，その三角形外側に位置する三つの小さな円の中心がそれぞれの点である三つの小円弧で定義される．$z$ が contour 上を動くとき，以下の式の増分は $-2\pi$ に等しいことを示せ：
\[
\am \left(\frac{1}{z - z_{1}} + \frac{1}{z - z_{2}} + \frac{1}{z - z_{3}}\right).
\]

8. 三次方程式 $f(z) = 0$ のすべての根を囲む閉曲線も，導出方程式 $f'(z) = 0$ の根を囲んでいることを証明せよ．[式
\[
f'(z) = f(z) \left(
  \frac{1}{z - z_{1}} + \frac{1}{z - z_{2}} + \frac{1}{z - z_{3}}
\right),
\]
ここで $z_{1}$, $z_{2}$, $z_{3}$ は $f(z) = 0$ の根であり，Ex.7 の結果を使用する．]

9. $f'(z) = 0$ の根は，三角形 $(z_{1}, z_{2}, z_{3})$ の辺の中点で接する楕円の焦点であることを示せ．[証明については Cesàro の \textit{Elementares Lehrbuch der algebraischen Analysis} (第 352 頁) を参照せよ．]

10. Ex.8 の結果を任意次数の方程式に拡張せよ．

\item 11. \( f(z) \)と\(\phi(z) \)が\( z \)の2つの多項式であり、\(\gamma\)が\( f(z) \)の根を通らない contour で、\(\gamma\)上のすべての点において\(|\phi(z)| < |f(z)|\)であるとき、方程式
\[
f(z) = 0,\quad
f(z) + \phi(z) = 0
\]
が\(\gamma\)内に含まれる根の数は同じです。

\item 12. \( a > e \)の場合、以下の方程式
\[
e^{z} = az,\quad
e^{z} = az^{2},\quad
e^{z} = az^{3}
\]
についてそれぞれ(i)正の根が一つ(ii)正と負の根が各一個(iii)正の根と複素根が二つ存在することが示されます。1910年数学オリンピック問題。

\appendix{II}{(第九章、第十章へ)}{重极限問題に関する一考察}

第九章、第十章の課程中、しばしば初学者を困惑させる種類の問題に遭遇しました。これらの問題は、最も一般的な形で扱われるとき、高度な数学において非常に困難でありながら最大限の興味と重要性を持つものです。

いくつかの具体的例を考えてみましょう。第213節では、
\[
\log(1 + x) = x - \frac{1}{2}x^{2} + \frac{1}{3}x^{3} - \dots,
\]
\(-1 < x \leq 1\)の範囲で、方程式
\[
\frac{1}{1 + t} = 1 - t + t^{2} - \dots
\]
を\(0\)から\(x\)まで積分することにより証明しました。証明したことは、
\[
\int_{0}^{x} \frac{dt}{1 + t}
  = \int_{0}^{x} dt - \int_{0}^{x} t\, dt + \int_{0}^{x} t^{2}\, dt - \dots;
\]
すなわち、無限級数 \(1 - t + t^{2} - \dots\) の和の積分（\(0\)から\(x\)まで）はその項の積分の和（同じ範囲で）に等しいということです。この事実を別の表現では、\((-1)^{n}t^{n}\)に対する無限級数の合算と、\(0\)から\(x\)までの積分が関数に対して交換可能であることを述べることもできます。つまり、これらの操作は関数に対してどの順序で行われても結果に影響を与えないということです。

また、§216で証明したように、指数関数
\[
\exp x = 1 + x + \frac{x^2}{2!} + \cdots
\]
の微分係数は自身に等しく、すなわち
\[
D_x \left(1 + x + \frac{x^2}{2!} + \cdots\right)
= D_x 1 + D_x x + D_x \frac{x^2}{2!} + \cdots;
\]
つまり、級数の和の微分係数はその項の微分係数の和に等しいということである，すなわち$x^n/n!$に対する無限和と$x$に関する微分が$x^n/n!$に対して適用されるとき交換律を満たす。

さらに、同じ節で偶然証明したように、関数$\exp x$は$x$の連続関数であり、言い換えれば
\[
\lim_{x \to \xi} \left(1 + x + \frac{x^2}{2!} + \cdots\right)
= 1 + \xi + \frac{\xi^2}{2!} + \cdots
= \lim_{x \to \xi} 1 + \lim_{x \to \xi} x + \lim_{x \to \xi} \frac{x^2}{2!} + \cdots;
\]
すなわち、級数の和の極限はその項の極限の和に等しいということである，すなわち$x = \xi$における無限和が連続であり、すなわち$x^n/n!$に対する無限和と$x$を$\xi$に近づける操作が適用されるとき交換律を満たす。

各々の場合において、私たちは結果の正しさについて特別な証明を与えました。本巻においては、どのような一つの定理からそれらのいずれかが即座に真であることを導く一般の定理を証明することはありませんでした。第xxxvii節1では有限数の連続項の和もまた連続であることを示しましたし、第113節では有限数の項の和の微分係数は各項の微分係数の和に等しいことを示しました。さらに、第160節では積分について対応する定理を述べました。したがって、特定の状況において、次の操作記号
\[
\lim_{x\to\xi} \dots,\quad
D_{x} \dots,\quad
\int \dots\, dx
\]
は有限数の項の和に対する操作と交換可能であることを証明しました。そして、特定の状況においても、無限数の項の和に対する操作との交換可能性があると考えるのは自然です。しかし、そのような状況を正確に定義することができるかどうかはまだ明らかではありません。

いくつかのさらなる交換可能および非交換可能な操作の例は、これらの点を明瞭にするのに役立つかもしれません。

(1) 2倍と3倍の乗算は常に交換可能です。すなわち
\[
2 × 3 × x = 3 × 2 × x
\]
がすべてのxに対して成り立ちます。

(2) zの実部を取り出す操作は、iで乗算する操作とは除くz=0の場合を除き交換不可能です。すなわち
\[
i × \Re(x + iy) = ix,\quad
\Re\{i × (x + iy)\} = -y.
\]

\Item{(3)} 二つの変数 $x$ および $y$ に対して零に収束する操作が関数 $f(x, y)$ に対して交換的であるかどうかは定かではありません。例えば
\[
\lim_{x \to 0} \left( \lim_{y \to 0} (x + y) \right) = \lim_{x \to 0} x = 0,\quad
\lim_{y \to 0} \left( \lim_{x \to 0} (x + y) \right) = \lim_{y \to 0} y = 0;
\]
が成り立ちますが、他方で
\begin{alignat*}{2}
\lim_{x \to 0} \left( \lim_{y \to 0} \frac{x - y}{x + y} \right)
  \&= \lim_{x \to 0} \frac{x}{x} \&\&= \lim_{x \to 0} 1 = 1,\\
\lim_{y \to 0} \left( \lim_{x \to 0} \frac{x - y}{x + y} \right)
  \&= \lim_{y \to 0} \frac{-y}{y} \&\&= \lim_{y \to 0} (-1) = -1.
\end{alignat*}

\Item{(4)} 操作 $\sum_{1}^{\infty} \dots$, $\lim_{x \to 1} \dots$ が交換的であるかどうかは定かではありません。例えば $x$ が $1$ より小さい値を取るとき
\begin{alignat*}{2}
\lim_{x \to 1} \left\{ \sum_{1}^{\infty} \frac{(-1)^n}{n} x^n \right\}
  \&= \lim_{x \to 1} \log(1 + x) \&\&= \log 2,\\
\sum_{1}^{\infty} \left\{ \lim_{x \to 1} \frac{(-1)^n}{n} x^n \right\}
  \&= \quad \sum\_{1}^{\infty} \frac{(-1)^n}{n} \&\&= \log 2;
\end{alignat*}
が成り立ちますが、他方で
\begin{align*}
\lim_{x \to 1} \left\{ \sum_{1}^{\infty} (x^n - x^{n+1}) \right\}
  &= \lim_{x \to 1} \{(1 - x) + (x - x^2) + \dots\}
   = \lim_{x \to 1} 1 = 1,\\
\sum_{1}^{\infty} \left\{ \lim_{x \to 1} (x^n - x^{n+1}) \right\}
  &= \sum_{1}^{\infty} (1 - 1) = 0 + 0 + 0 + \dots = 0.
\end{align*}

これらの例から、二つの与えられた操作が交換的である可能性は三つあることが示されます。すなわち、(1) 操作は常に交換的である；(2) 操作は例外的な状況を除いて決して交換的でない；(3) 一般的に現れる多くの普通の状況では操作が交換的である。

The本当に重要な場合（私たちは第 IX 章から引用した例によって示されるように）は、各操作が極限への移行を含むものであるときであり、そのような操作は微分や無限級数の和のようなものである。このような操作は「極限操作」と呼ばれる。二つの与えられた極限操作が交換可能であるか否かという一般的な問題は、数学全体において最も重要な問題の一つである。しかし、この種の質問を一般的な定理を使って扱おうとするのは、本書の範囲を超えるだろう。

ただし、上記の例によって示される通り、答えはそのような質問に対する一般的な解答の線に沿って進む。もし$L$と$L'$が二つの極限操作であるならば、数値$LL'z$と$L'Lz$は厳密理論的な意味での「一般」には等しくない。適切な工夫によって常に、$z$を選べば$LL'z$と$L'Lz$は互いに異なることができる。しかし、「普通」の意味で言えば、つまり「自然に起こりそうな多くの場合において」または「通常の場合において」、これらは等しいものである。

当然、純粋数学のような正確な科学では、このような答えには満足することはできない；そして、数学のより高度な分野では、これらの問題に対する詳細な調査は絶対的な必要性である。しかし、現在のところ、読者は我々が先ほど述べたことに注意を払うだけで十分である。実際において、二つの極限操作が交換可能であると仮定して得られた結果は「おそらく」真であり、少なくとも問題に対する答えへのヒントとして有用である。しかし、そのような答えは、一般の問題に対する更なる研究または特定の問題に対する特別な調査（第 IX 章で例示したように）がない限り、提案されたものとしてのみ受け入れられ、証明されるものではない。

大量の重要な重複極限問題に関する詳細な調査は、Bromwich の『無限級数』において見ることができます。
\PageSep{443}

付録{III}（第\SecNo[§]{158}節および第\okrickRef{Chapter}{IX}）

読者は、定義から理論を導く演習として、円関数の理論を計算することを試みることで有益な経験を得ることができます。
\CenterDef[\footnotemark]{\Item{(1)}}{$y = y(x) = \arctan x = \ds\int_{0}^{x} \frac{dt}{1 + t^{2}}$.}
\footnotetext{行末のこれらの文字は、その中に含まれる公式が定義であることを示します。}

式~\Eq{(1)} は、すべての実数 $x$ に対して一意的な $y$ の値を定義します。$y$ は連続かつ単調に増加するため、逆関数 $x = x(y)$ も存在し、連続かつ単調に増加します。私たちは
\CenterDef{\Item{(2)}}{$x = x(y) = \tan y$.}

を書き込みます。

$\pi$ を式~\Eq{(3)} によって定義します。
\CenterDef{\Item{(3)}}{$\frac{1}{2}\pi = \ds\int_{0}^{\infty} \frac{dt}{1 + t^{2}}$,}
この関数は $-\frac{1}{2}\pi < y < \frac{1}{2}\pi$ の範囲で定義されます。

さらに、
\CenterDef{\Item{(4)}}{$\cos y = \dfrac{1}{\sqrt{1 + x^{2}}},\quad \sin y = \dfrac{x}{\sqrt{1 + x^{2}}}$,}
ここで平方根は正です。そして $y$ が $-\frac{1}{2}\pi$ や $\frac{1}{2}\pi$ のとき、これらの関数を定義し直して、それらの値に対する連続性を保つようにします。
最後に、$\DPmod{(-\frac{1}{2}\pi, \frac{1}{2}\pi)}{[-\frac{1}{2}\pi, \frac{1}{2}\pi]}$ 以外の範囲で $\cos y$ と $\sin y$ を定義します。
\CenterDef{\Item{(5)}}{$\begin{alignedat}{2}\tan(y + \pi) &= &&\tan y,\\ \cos(y + \pi) &= -&&\cos y, \\ \sin(y + \pi) &= -&&\sin y.\end{alignedat}$}

これにより、すべての $y$ の値に対して $\cos y$ と $\sin y$ を定義し、$\tan y$ は奇数倍の $\frac{1}{2}\pi$ 以外のすべての $y$ の値に対して定義されます。$\cos y$ と $\sin y$ はすべての $y$ の値で連続であり、$\tan y$ はその定義が失われる点を除いて連続です。

进一步的发展僅依赖于加法定理。書き換えると、

\[
x = \frac{x_{1} + x_{2}}{1 - x_{1}x_{2}},
\]

方程式（1）を変形するには、次の置換を行う。

\[
t = \frac{x_{1} + u}{1 - x_{1}u},\quad
u = \frac{t - x_{1}}{1 + x_{1}t}.
\]

すると、

\begin{align*}
\arctan \frac{x_{1} + x_{2}}{1 - x_{1}x_{2}}
  &= \int_{-x_{1}}^{x_{2}} \frac{du}{1 + u^{2}}
   = \int_{0}^{x_{1}} \frac{du}{1 + u^{2}} + \int_{0}^{x_{2}} \frac{du}{1 + u^{2}} \\
  &= \arctan x_{1} + \arctan x_{2}.
\end{align*}

これを基に、

\[
\tan (y_{1} + y_{2}) = \dfrac{\tan y_{1} + \tan y_{2}}{1 - \tan y_{1}\tan y_{2}},
\]

という方程式が導かれる。この方程式は最初、$y_{1}$, $y_{2}$, および $y_{1} + y_{2}$ が $\DPmod{(-\frac{1}{2}\pi, \frac{1}{2}\pi)}{[-\frac{1}{2}\pi, \frac{1}{2}\pi]}$ 内にある場合に証明され、方程式（5）を用いてすべての $y_{1}$ および $y_{2}$ の値へ即座に拡張できる。

方程式（4）と方程式（6）から、

\[
\cos(y_{1} + y_{2}) = ±(\cos y_{1}\cos y_{2} - \sin y_{1}\sin y_{2}),
\]

を得る。符号を定めるためには $y_{2} = 0$ を置く。方程式は $\cos y_{1} = ±\cos y_{1}$ に簡約され、少なくとも一つの $y_{2}$ の値（すなわち $y_{2} = 0$）に対して正の符号を選ぶ必要があることを示す。連続性から考えると、すべての場合において正の符号を選ぶべきであることがわかる。同様の方法で $\sin(y_{1} + y_{2})$ の公式も導くことができる。

円関数の微分公式は通常の方法で導き出し、Taylorの定理から派生する冪級数も得られる。

無限級数理論に基づく円関数の別の理論がある。例えば $\cos x$ は方程式

\[
\cos x = 1 - \frac{x^{2}}{2!} + \frac{x^{4}}{4!} - \dots
\]

によって定義されるという理論は、WhittakerとWatsonの『Modern Analysis』（付録A）に見ることができる。

IV 附録：分析および幾何における無限について

あるシステムの解析幾何には「無限」な要素が含まれていますが、すべてのシステムには含まれていません。無限の点や無限の直線などがあります。この短い論文の目的は、これらの概念が極限の分析的教義に依存していないことを指摘することです。

通常のカートシアン幾何において、点とは実数のペア$(x, y)$です。直線とは、一次関係$ax + by + c = 0$を満たす点の集合であり、$a$と$b$がともにゼロでない限り存在します。無限の要素は存在せず、2つの直線は共有点を持たない場合があります。

実数ホモジネアス幾何において、点とは実数のトリオ$(x, y, z)$の集合であり、その成分が比例するトリオは同一視されます。直線とは、一次関係$ax + by + cz = 0$を満たす点の集合で、$a$、$b$、$c$がすべてゼロでない限り存在します。いくつかのシステムでは、1つの点や直線も他の点や直線と同じ地位に置かれています。他には、特定の「特別」な点と直線が特に区別され、それらの特別な要素に対する他の要素の関係性が強調されます。したがって、「実数ホモジネアスカートシアン幾何」と呼ばれるシステムでは、$z = 0$の点が特別であり、1つの特別な直線があります。この特別な直線は「無限の直線」と呼ばれています。

これは幾何学に関する論文ではなく、詳細に議論する機会もありません。重要な点は次のとおりです。分析における無限は「極限」的なものではなく、「実際の」ものではありません。「$\infty$」という記号は、この本全体を通じて「不完全な記号」として扱われ、独立した意味が付与されていませんが、その中に含まれる特定の表現には独立した意味が付けられています。しかし、幾何における無限は「実際のもの」であり、「極限」的なものではありません。「無限の直線」は他の直線と同じように直線であるという意味で存在します。

homogeneous と common 卡尔丹几何之间可以建立一种相关性，在これ种相关性中，第一系统的所有元素（特殊元素除外）在第二系统中有対应的元素。例如，直线 $ax + by + cz = 0$ 対应于直线 $ax + by + c = 0$。第一直线上的每一个点在第二直线上都有対应点，但有一个例外，即 $z = 0$ 的点。
当 $(x, y, z)$ 在第一直线上变化，并趋向于 $z = 0$ 的特殊点時，対应的第二直线上的点的变化使得它与原点的距离趋于无穷大。これ种相关性在历史上具有重要意义，因為该领域的词汇就是从これ里派生出来的，而且它经常用于说明目的。然而これ只是说明而已，不能基于此来建立几何无穷的理性解释。学生中普遍存在的混乱正是由于通常使用的解析几何教科书中将说明误认為真実情况而产生的。

創作物は公有領域の印刷版から作成し、これらの著作物には米国での著作権が存在しないため、財団（そしてあなた！）は米国内で許可なく複製・配布を行うことができます。著作権料を支払うことなく、Project Gutenberg-tmの概念と商標を保護する特別な規定は、このライセンスの利用規約の一部に記載されています。Project Gutenbergは登録商標であり、電子書籍を有偿で販売することは許可されていませんが、特定の場合には例外があります。この電子書籍を有償で配布しない場合、ルールを遵守するのは非常に簡単です。この電子書籍を使用して几乎任何目的，如衍生作品的创作、报告撰写、表演和研究都是允许的。您可以修改并打印分发これ些公共领域的电子书——您几乎可以做任何事情。但是，重新发布必须遵循商标许可，尤其是商业重新发布。

***著作権表示開始***

プロジェクト・ゲッutenberg-tmの使命を保護するために電子作品の無料配布を促進することを目的とするため、この作品（またはその他の関連する作品）を使用または配布する場合、または`Project Gutenberg`と関連する表現を使用する場合は、以下の全著作権表示規約に同意することに同意します。このファイルとともに提供されるか、またはhttp://gutenberg.net/licenseでオンラインで利用可能です。

第1条. 一般使用条件およびProject Gutenberg-tm電子作品の再配布

1.A. このProject Gutenberg-tm電子作品を読んだり使用したりすることにより、あなたはその著作権および知的財産（商標／著作権）に関する契約のすべての条項に同意し理解し、それに従うことを承認したと表明します。この契約のすべての条項に同意しない場合は、Project Gutenberg-tm電子作品の所有コピーを全て廃棄または破棄する必要があります。

1.B. 「Project Gutenberg」は登録商標です。この協定のすべての条項に同意することにより、電子作品に関連させることを許可される人だけが使用することができます。多くのProject Gutenberg-tm電子作品には、この協定の全部を遵守せずにでも以下の行為を行うことができます。1.C節をご覧ください。この協定を遵守し、Project Gutenberg-tm電子作品への将来アクセスを保護するための支援を行えば、多くのProject Gutenberg-tm電子作品に関連させることができます。1.E節をご覧ください。

1.C. ザ・プロジェクト・ジンクルベルト・リテラリー・アーカイブ・ファウンデーション（以下、「ファウンデーション」または「PGLAF」という）は、プロジェクト・ジンクルベルト-tm電子作品の集積著作権を所有しています。几乎所有的集合中的个别作品在美国均為公有领域。如果您位于美国并且所涉个别作品在美国為公有领域，我達不主张阻止您复制、分发、表演、展示或基于该作品创建衍生作品的权利，只要所有対项目・ジンクルベルトの参照被移除。当然，我達希望您能を通じて遵守本协议条款来支持项目・ジンクルベルト-tm的使命，即促进电子作品的免费访问，并且在不收费分享これ些作品時附带完整的项目・ジンクルベルト-tm许可証。

1.D. 你所在的国家或地区的版权法也规范了你可以対これ部作品做什么。大多数国家的版权法处于不断变化之中。如果您不在美国，请在下载、复制、展示、表演、分发或基于此作品或其他项目・ジンクルベルト-tm作品创建衍生作品之前，检查您所在国家的法律以及本协议的条款。

1.E. 除非你已经移除了所有対プロジェクト・ジンクルベルトの参照：

1.E.1. 次の文は、全著作権表示やその他の公開アクセスにリンクを設けるか、もしくはProject Gutenberg-tmライセンスに従うための直接的な方法が提供されている限り、Project Gutenberg-tm作品（「Project Gutenberg」という語が含まれるか、「Project Gutenberg」に関連しているすべての作品）が閲覧、表示、実行、コピー、または配布されるたびに目立つ位置に表示する必要があります：

この電子書籍は世界中どこでも無料で使用でき、制限もほとんどありません。あなたはこれを複製し、無償で贈呈したり再利用したりすることができます。この電子書籍と付属のProject Gutenberg-tmライセンスまたはwww.gutenberg.netでオンラインで提供されているものに従ってください。

1.E.2. 公共ドメインから派生した個別のProject Gutenberg-tm電子作品が著作権者の許可なしで作成された場合（著作権者への許可の表示がない場合）、その作品は米国の人々に対して無償で複製し、配布することができます。もし、この作品に「Project Gutenberg」が関連付けられたり表示されている場合、あなたは1.E.1から1.E.7までの要求に従うか、または1.E.8や1.E.9に記載されているように、著作権の使用とProject Gutenberg-tm商標を使用するための許可を得る必要があります。

1.E.3. 著作権者の許可を得て掲載された個別のProject Gutenberg-tm電子作品が掲載された場合、あなたの使用および配布は1.E.1から1.E.7までの要求と、著作権者がこの作品に付与した追加の条件に従う必要があります。追加の条件は、著作権者の許可を得て掲載されたすべての作品の最初に見つかるProject Gutenberg-tmライセンスにリンクされています。

1.E.4. 全体的なProject Gutenberg-tmライセンス条項をこの作品から削除したり、取り外したりしてはなりません。また、この作品の一部またはProject Gutenberg-tmに関連する他の作品が含まれているファイルも同様です。

1.E.5. この電子著作物やその一部を複製、表示、実行、配布または再配布しないでください。第1.E.1条に記載された文を目立つ位置で表示し、プロジェクト・ゲンティゥールス-tmライセンスの全文へのリンクまたは即時のアクセス方法を提供する必要があります。

1.E.6. この著作物は任意形式、包括的または非包括的二進数形式、圧縮形式、マークアップ形式、プロパティーや非プロパティーオービット形式、ワードプロセッシング形式やハイパーテキスト形式に変換し、配布することができます。ただし、プロジェクト・ゲンティゥールス-tm公式ウェブサイト（www.gutenberg.net）で掲載されている公式版とは異なる形式で著作物へのアクセスを提供または複製を配布する場合、ユーザーに対して無償で著作物の元の「 Plain Vanilla ASCII 」形式やその他の形式のコピーを提供する方法、または要求によりコピーを取得できる方法を提供する必要があります。代替形式には、第1.E.1条に記載されたプロジェクト・ゲンティゥールス-tmライセンスが含まれます。

1.E.7. 「 Plain Vanilla ASCII 」形式やその他の公式版で掲載されている著作物へのアクセス、表示、実行、複製または配布を有料で提供しないでください。ただし、第1.E.8条または1.E.9条に従っている場合を除きます。

1.E.8. プロジェクト・ゲンティゥールス-tm電子著作物のコピーを提供するか、またはアクセスを提供し、配布する料金を請求することができます。ただし、以下のようにする必要があります。

あなたは、Project Gutenberg-tm著作権料を総収益の20％とし、その計算方法で適用税率を計算する方法を使用して支払う必要があります。この費用は、Project Gutenberg-tm商標の所有者に負担されるが、彼はこの段落に基づいて著作権料をProject Gutenberg文学アーカイブ財団に寄付することに同意しています。著作権料は、あなたが確定申告書を作成（または法律で義務付けられている場合）する日から60日以内に支払う必要があります。著作権料は明確に「著作権料」と記載され、Project Gutenberg文学アーカイブ財団の指定した住所に送付されるべきです。

ユーザーが30日以内に書面（または電子メール）で通知し、彼または彼女は全Project Gutenberg-tmライセンスの条項に同意しない場合、あなたはその金銭を全額返還する必要があります。あなたは、そのようなユーザーに対し、物理媒体に保持されている作品のすべてのコピーを返却または破棄させ、他のProject Gutenberg-tm作品への使用とアクセスを停止させる必要があります。

電子作品に缺陷が見つかった場合、90日以内に通知した場合、あなたは第1.F.3段落に従って金銭全額または代替コピーを提供する責任があります。

あなたは、本契約による無料配布のすべての条件を遵守する必要があります。

1.E.9. あなたがProject Gutenberg-tm電子作品またはそのグループを異なる条件で料金を請求または配布したい場合、Project Gutenberg文学アーカイブ財団とMichael Hartから書面での許可が必要です。Foundationは第3節に記載されていますので、連絡先を確認してください。

1.F.

1.F.1. プロジェクト・ゲンデルバーグのボランティアと従業員は、著作権調査を行い、公開著作物を翻訳し、校正することでプロジェクト・ゲンデルバーグ収集を作成するために相当な努力を払っています。これらの努力にもかかわらず、プロジェクト・ゲンデルバーグ電子著作物及びその保存媒体には、不完全、正確でないまたは破損したデータ、翻訳の誤り、著作権その他の知的財産権侵害、破損または損傷したディスクその他の媒体、コンピュータウイルス、またはあなたの機器によって読み取れないコンピュータコードなどの「欠陥」が含まれる可能性があります。

1.F.2. 保証および損害に関する限定的な保証、損失の不承認 - この契約書第1.F.3条に記載されている「返品または返金権」以外については、プロジェクト・ゲンデルバーグ文学アーカイブ財団、プロジェクト・ゲンデルバーグ商標所有者およびこの契約に基づきプロジェクト・ゲンデルバーグ電子著作物を配布する他の当事者は、あなたに対する損害、費用及び経費、法律費用を含む一切の責任を免れることとします。あなたは、過失、厳格な責任、保証違反または契約違反による補填請求権については、第1.F.3条に記載されているもの以外には同意するものとします。また、財団、商標所有者及びこの契約に基づき配布する当事者は、あなたに対する実質的、直接的、間接的、結果的、懲罰的または偶然の損害について一切の責任を負わないことを同意するものとします。

1.F.3. 惟一の返品または返金の権利 - この電子書籍に90日以内に欠陥を発見した場合、あなたは支払った金銭（ある場合は）を返還してもらうために、受領元の人に対し文字で説明を行うことができます。物理媒体で受け取った場合、その媒体と共に文字による説明を提出する必要があります。欠陥のある作品を提供した人または法人は代替の複製物を提供する選択をすることができます。電子的に受け取った場合、作品を提供した人は再び電子的に受け取る機会を与えることを選択することができます。二回目の受け取りでも欠陥がある場合は、文字で返金を求めることができます。

1.F.4. この作品は第1.F.3条に定める惟一の返品または返金の権利を除き、「現状のまま」提供され、いかなる明示的または黙示的な保証も含みません。商品性や特定目的への適合性に関する保証も含みません。

1.F.5. 一部の州では、特定の黙示的保証の否定または一定の種類の損害に対する排除または制限を禁止する法律が存在します。この協定に記載された否定または制限が適用される州の法によって許可される最大の範囲で解釈される場合があります。本協定のどの条項も無効または執行不能であるとしても、その他の条項は依然として有効です。

1.F.6. 損賠 - あなたは、財団、商標権者、財団の代理人または従業員、Project Gutenberg-tm電子作品を本契約に従って提供する者、およびProject Gutenberg-tm電子作品の生産、推進及び配布に関連するボランティアから、直接的または間接的に次のいずれかが発生した全ての責任、費用及び経費、法律費用を免れるよう indemnify し、hold harmless することに同意します：(a) この作品又はProject Gutenberg-tm作品の配布、(b) Project Gutenberg-tm作品への変更、改訂、追加又は削除、(c) あなたが原因となった Defect。

2. Project Gutenberg-tm のミッションについて

Project Gutenberg-tmとは、最も多様な種類のコンピュータ（古きや新しいもの含む）で読み取れる電子作品の無料配布に特化したものです。その存在は数百人のボランティアと様々な分野の人々からの寄付によって支えられています。

ボランティアの支援を提供し、彼らが必要とする援助を与えることは、Project Gutenberg-tmの目標達成と、Project Gutenberg-tmコレクションが今後も自由に利用可能であることを確保するための鍵となります。2001年には、Project Gutenberg文学アーカイブ財団が設立され、Project Gutenberg-tmとその将来世代を安全で永久的な未来へと導くために設立されました。この財団についてさらに詳しく知るためには、セクション3及び4をご覧ください。また、財団のウェブサイト http://www.pglaf.org をご確認ください。

3. Project Gutenberg文学アーカイブ財団に関する情報

プロジェクト・ジンブルトン文学アーカイブ財団は、ミシシッピ州に設立され、内務省から税制上の非課税状態が認定された501(c)(3)の非営利教育法人です。財団のEIN（米国連邦税番号）は64-6221541で、その501(c)(3)の認定書はhttp://pglaf.org/fundraisingに掲載されています。米国および州法により許可される最大限の範囲で寄付金は税引額となります。

財団の本社事務所は、アラスカ州フェアバンクス市メランドロード4557番地にありますが、ボランティアや従業員は多数の場所に散在しています。事業事務所は、ユタ州サalt lake city市北1500西809番地にあり、電話（801）596-1887、メールbusiness@pglaf.orgで連絡できます。最新の連絡先情報は財団のウェブサイトと公式ページhttp://pglaf.orgで確認できます。

追加の連絡先情報：
グレゴリー・B・ニュービー博士
最高経営責任者および理事長
gbnewby@pglaf.org

第4条. プロジェクト・ジンブルトン文学アーカイブ財団への寄付について

プロジェクト・ジンブルトン-tmは広範な公衆の支援と寄付なくして存続することはできません。機械読取可能な形式で自由に配布できる公衆の領域およびライセンス作品の数を増やすという使命を達成するためには、古式的な機器も含む幅広い設備にアクセスできる最大限の範囲の機器で利用可能です。特に米国連邦税務署による非課税状態維持のために重要なのは、1ドルから5000ドルまでの小さな寄付です。

財団は、連邦合衆国の50州における慈善活動および寄付に関する法規制に準拠することをコミットしています。これらの法規制は一様ではなく、遵守するためには相当な努力と多くの書類作成、また費用が必要です。寄付の募集を行った州以外では寄付を募集していません。寄付を行うか、または特定の州の遵守状況を確認するには、http://pglaf.org を訪問してください。

寄付募集要件に達していない州から寄付を募ることはできませんが、そのような州からの無償の寄付を受ける禁止規定はないと知られています。

国際的な寄付は歓迎しますが、国外からの寄付による税制上の待遇については明確な声明をすることができません。米国の法規制だけでは、私たしの小さなスタッフには圧倒されます。

現在の寄付方法と連絡先についてProject Gutenberg Webページをご確認ください。寄付は現金払い、オンライン支払い、クレジットカードによる寄付など様々な方法で受け付けます。寄付を行うには、http://pglaf.org/donate を訪問してください。

第5条 Project Gutenberg-tm電子作品に関する一般情報

Project Gutenberg-tmの概念を提唱したのは、誰でも自由に共有できる電子版著作物のライブラリを作成し配布することを目的としたマシュー・S・ハート教授です。30年以上にわたり、彼は志願者による松散なネットワークのもとでProject Gutenberg-tm電子書籍を作成し、配布してきました。

Project Gutenberg-tm電子書籍は、米国では公衆の財産であることが確認された複数の印刷版から作成されることが多いですが、著作権表示が含まれている場合を除き、特定の紙面版に準拠しているとは限らないことに注意してください。

このウェブサイトには、プロジェクト・ゲンツベリューム-tmに関する情報が含まれており、その内容には著作権文学アーカイブ財団への寄付方法、新しい電子書籍の制作にどのように協力するか、および新しい電子書籍についてのメールニュースレターを購読する方法が記されている。

\end{document}